\documentclass{missingbook}
\usepackage{hydrogen}

\begin{document}

\frontmatter
\hypersetup{pdftitle={The Hydrogen Atoms of Deep Learning: Solvable Models and What They Get Wrong},pdfauthor={Written by Claude; conceived and edited by Xiaoyu Li}}
\mbtitlepage{The Hydrogen Atoms of Deep Learning}{Solvable Models and What They Get Wrong}{Written by Claude\\[2.5mm]Conceived and edited by Xiaoyu Li}{First edition \textperiodcentered{} October 2026}
\addchap{Preface}

This book began with a list of textbooks that should exist but do not. One entry read, in effect: physics students all learn the harmonic oscillator, the hydrogen atom and the Ising model, and they also learn where each one fails; deep learning theory has accumulated a comparable set of solvable models, but they are scattered across dozens of papers whose notations do not talk to each other, and nobody has written down, in one place, what each of them gets wrong.

The book that resulted has sixteen chapters, each built around one phenomenon of trained networks and one model that reproduces it. Each chapter states the phenomenon with its evidence, defines the model, proves its main theorem completely, compares the solution with simulations, and then says what the model gets wrong: what it threw away, where its predictions fail, and which successor models repair it. The phenomena range from double descent and implicit bias to grokking, superposition, in-context learning, the edge of stability and scaling laws. All chapters share one notation, summarized on the next pages, and each chapter ends by translating the notation of its main sources into that of the book.

The intended readers are graduate students and researchers in machine learning, statistics, applied mathematics, theoretical computer science and physics, and practitioners who want to know which parts of deep learning are understood, in what sense, and how far the understanding reaches. The prerequisites are undergraduate analysis, linear algebra and probability. Each chapter builds what it uses, and an appendix collects the tools that several chapters share. The prologue explains the plan of the book and solves its first model.

A solvable model is a lens, not a mirror. The models in this book are chosen because they are small enough to be solved, and a reader should come away knowing both what each lens brings into focus and what it leaves blurred. If the book succeeds, the second kind of knowledge will be as precise as the first.

\paragraph{A note on authorship.} This book was written by Claude (Opus 5.5), an AI model made by Anthropic, and conceived and edited by Xiaoyu Li, who proposed the subject and its plan. Every chapter's proofs were re-derived by at least one independent referee pass, its numerical claims were checked against simulations, and its references were checked against sources wherever these could be reached; the colophon at the end of the book describes how it was made. Mistakes surely remain. Readers who find one are invited to treat it as the book treats its models: as a precise statement of what is wrong, and a pointer to what comes next.

\mbtableofcontents
\addchap{Notation}

One notation is used throughout the book. Where a chapter needs a local symbol that clashes with this list, the chapter says so; where a source paper uses a different notation, the chapter's notes translate it.

\newcommand{\notationrow}[2]{#1 & #2\\}
\newcommand{\notationhead}[1]{\multicolumn{2}{@{}l}{\textsc{\color{sea}#1}}\\[0.3ex]}
\noindent{\small
\renewcommand{\arraystretch}{1.1}%
\begin{tabular}{@{}>{\raggedright\arraybackslash}p{30mm}>{\raggedright\arraybackslash}p{86mm}@{}}
\notationhead{sizes}
\notationrow{$n$}{number of training samples}
\notationrow{$d$}{input dimension}
\notationrow{$m$}{width of a hidden layer (number of hidden units)}
\notationrow{$p$}{number of features or parameters of a linear model (the prime modulus in \cref{ch:grokking})}
\notationrow{$K$}{number of classes}
\notationrow{$N$}{context length (\cref{ch:icl,ch:induction}); model size, with $D$ the data size (\cref{ch:scaling})}
\notationrow{$L$, $\ell$}{depth, and the index of a layer}
\notationrow{$k$}{a degree, a sparsity, or an information exponent}
\addlinespace
\notationhead{data and models}
\notationrow{$x$, $y$}{input and label}
\notationrow{$f(x;\theta)$}{the predictor, with parameters $\theta$}
\notationrow{$w$, $W$}{a weight vector and a weight matrix}
\notationrow{$\phi$}{an activation function}
\notationrow{$w_\star$, $f_\star$, $W_\star$}{a subscript star marks the target or the teacher}
\notationrow{$\lambda^\star$, $\Gamma^\star$}{a superscript star marks an optimal value}
\notationrow{$r=\ip{w,w_\star}$}{the overlap of two unit vectors; $Q$, $R$, $T$ are matrices of overlaps}
\notationrow{$\sigma$}{the standard deviation of label noise}
\notationrow{$\sigma_w^2$, $\sigma_b^2$}{the variances of weights and biases at initialization}
\notationrow{$\alpha$}{the scale of initialization}
\addlinespace
\notationhead{training}
\notationrow{$\Loss(\theta)$, $\Lhat(\theta)$}{population loss and training loss}
\notationrow{$\eta$}{step size}
\notationrow{$t$}{time: continuous for flows, an integer for iterations}
\notationrow{$\lambda$}{ridge parameter or weight decay; eigenvalues always carry a subscript, as in $\lambda_k$, $\lmax$, $\lmin$}
\end{tabular}

\medskip\noindent
\begin{tabular}{@{}>{\raggedright\arraybackslash}p{30mm}>{\raggedright\arraybackslash}p{86mm}@{}}
\notationhead{probability and functions}
\notationrow{$\cN(\mu,\Sigma)$, $\Unif$}{Gaussian and uniform laws}
\notationrow{$\Sph^{d-1}$}{the unit sphere in $\R^d$}
\notationrow{$\He_k$, $h_k$}{probabilists' Hermite polynomials and their normalization $h_k=\He_k/\sqrt{k!}$, orthonormal under $\cN(0,1)$}
\notationrow{$\hat g_k$}{the Hermite coefficient $\E[g(z)h_k(z)]$, $z\sim\cN(0,1)$}
\notationrow{$\ind\{\cdot\}$}{indicator}
\addlinespace
\notationhead{linear algebra}
\notationrow{$\ip{u,v}$, $\norm{u}$}{Euclidean inner product and norm}
\notationrow{$\norm{A}_{\op}$, $\norm{A}_F$}{operator and Frobenius norms}
\notationrow{$A^\top$, $I_d$, $\mathbf 1_d$}{transpose, identity, all-ones vector}
\addlinespace
\notationhead{asymptotics}
\notationrow{$a\lesssim b$, $a\asymp b$}{$a\le Cb$, and $a\lesssim b\lesssim a$, for a constant $C$ depending only on parameters declared fixed}
\notationrow{$O$, $\Omega$, $\Theta$, $o$, $\tilde O$}{the usual order symbols; $\tilde O$ hides polylogarithmic factors}
\addlinespace
\notationhead{status of an argument}
\notationrow{$\blacksquare$}{end of a complete proof}
\notationrow{proof sketch}{an argument that omits steps, with a pointer to where they are}
\notationrow{$\Diamond$}{end of a heuristic derivation: an argument that is not a proof}
\notationrow{observation}{an empirical regularity, with its evidence}
\end{tabular}\par}


\mainmatter
% Prologue of The Hydrogen Atoms of Deep Learning.
\specialchapter{Prologue}{The hydrogen atom, and what a solvable model is for}{4}\label{ch:prologue}
\begin{refsection}[chapters/00-prologue.bib]

Every physics student solves the same three problems before solving anything else: the harmonic oscillator, the hydrogen atom, and the Ising model. None of them describes a real system exactly, and every student learns, along with the solutions, the specific ways in which each one is wrong. This book proposes the same education for deep learning. It collects the models of deep learning that can be solved completely, solves each one in a single notation, and says, for each, exactly where it fails.

\addsec{Three models every physicist solves}

\paragraph{The hydrogen atom.} In 1885 Johann Balmer, a teacher at a girls' school in Basel, found that the wavelengths of the four visible lines of hydrogen fit the formula $\lambda = B\,n^2/(n^2-4)$ for $n=3,4,5,6$, with a single constant $B\approx364.56$ nanometers \parencite{Balmer1885}. His formula was a regularity without a reason. Johannes Rydberg generalized it to other elements and wrote it in terms of wave numbers, $1/\lambda = R(1/4-1/n^2)$ \parencite{Rydberg1890}. In 1913 Niels Bohr derived Rydberg's constant from a model in which the electron's angular momentum is quantized \parencite{Bohr1913}, and in 1926 Erwin Schr\"odinger derived the whole spectrum as the eigenvalues of a differential operator, in the first paper on his wave equation \parencite{Schrodinger1926}. The lines on the title page of this book stand at the wavelengths that Rydberg's formula predicts.\sidenote{The emblem draws the Balmer series, $n=3,4,\dots$, with the four visible lines labelled; the dotted line marks the series limit near $364.6$ nanometers, where the lines accumulate.}

The solution is exact, and it is wrong. Arnold Sommerfeld had already found in 1916 that relativistic corrections split the lines into closely spaced multiplets, the \term{fine structure}\index{fine structure} \parencite{Sommerfeld1916}; Paul Dirac's relativistic equation of 1928 accounted for it \parencite{Dirac1928}. Dirac's theory in turn predicted that two levels of the $n=2$ shell coincide. In 1947 Willis Lamb and Robert Retherford measured a small splitting between them \parencite{LambRetherford1947}, and within weeks Hans Bethe explained most of it by the interaction of the electron with the quantized electromagnetic field \parencite{Bethe1947}. The Lamb shift became one of the starting points of quantum electrodynamics. Each correction to the hydrogen atom names a piece of physics that the previous model had left out: relativity, then the field.

\paragraph{The Ising model.} Wilhelm Lenz proposed a lattice of spins that can point up or down and prefer to align with their neighbors; his student Ernst Ising solved it in one dimension in 1925, found no phase transition, and conjectured that there was none in higher dimensions either \parencite{Ising1925}. The conjecture was false. Rudolf Peierls argued in 1936 that the two-dimensional model does order at low temperature \parencite{Peierls1936}, and Lars Onsager computed its free energy exactly in 1944 \parencite{Onsager1944}. The cheaper solution, mean-field theory, is exact when every spin interacts with every other, and it predicts that the magnetization vanishes at the critical temperature like $(T_c-T)^{1/2}$. In two dimensions the exponent is $1/8$, as C.\,N.\ Yang derived from Onsager's solution \parencite{Yang1952}. Mean-field theory gets the phase diagram qualitatively right and the critical behavior wrong, and the reason it is wrong (fluctuations matter in low dimension) became the subject of the renormalization group.

\paragraph{The harmonic oscillator.} Expand a potential to second order around a minimum and every system becomes a set of independent oscillators, its normal modes. The approximation is the most useful in physics and it is wrong in instructive ways. A perfectly harmonic crystal does not expand when heated, and heat flows through it without resistance, because its normal modes never exchange energy; both thermal expansion and finite thermal conductivity come from the anharmonic terms that the expansion dropped \parencite{AshcroftMermin1976}.

\medskip

What these three have in common is not that they are true. Each is \emph{solvable}, because a symmetry or a limit makes it so (rotational invariance for hydrogen, a transfer matrix for Ising, quadratic energy for the oscillator). Each reproduces something real. Each makes quantitative predictions. And each fails in ways that are known, specific, and informative, so that learning the failure is learning the next piece of physics. George Box's remark that all models are wrong is often quoted as an excuse for a model \parencite{Box1979}. The physics curriculum treats it as an instruction: say how.

\addsec{The models of deep learning}

Deep learning has no standard model, and it may never have one: a trained network is a large, specific object shaped by its data, its architecture and its training procedure, and the phenomena it displays are many. But the theory of deep learning has, over the last decade, accumulated a surprising number of small models that can be solved exactly and that reproduce one phenomenon each. Least squares with Gaussian features shows a test error that falls, rises and falls again. A linear network with two layers learns in stages. A network whose weights are products of two numbers prefers sparse solutions when started small. Gradient descent on a single neuron in high dimension spends most of its time searching and then learns quickly. A small network trained on modular addition computes with Fourier modes. A linear attention layer trained on regression prompts performs one step of gradient descent on the prompt. A product of two scalars trained with a large step size jumps out of a sharp valley into a flat one.

These models are scattered across dozens of papers, written in different notations, by communities that rarely cite each other: statistical physics, high-dimensional statistics, optimization, theoretical computer science, and mechanistic interpretability. The places where the models fail are written down even less often, and are usually in other papers. A student who wants the deep learning analogue of the hydrogen atom has to assemble it from the papers, and then discover, unaided, what it gets wrong.

This book assembles sixteen of them. Each chapter takes one phenomenon observed in trained networks, gives the smallest model that reproduces it, proves the model's main theorem completely, and then says precisely what the model gets wrong, with references to the results that show it.\sidenote{The book is not a survey. Many important results of deep learning theory are absent because they are not solvable models of a phenomenon: generalization bounds, approximation theory, most of optimization theory. The monograph of \textcite{Zhao2026} covers the field's breadth, one topic per chapter.}

\addsec{What makes a model a hydrogen atom}

We hold every model in the book to four tests.

\begin{enumerate}
\item \emph{It reproduces a phenomenon.} Something observed in real networks, robustly enough that we can state it as an observation with evidence, appears in the model.
\item \emph{It is solvable.} Its main claim has a complete proof that a graduate student can check, short enough to teach in a lecture or two.
\item \emph{It predicts a number.} An exponent, a threshold, a time scale, a phase boundary, or a constant, which can be compared with an experiment.
\item \emph{It fails instructively.} The ways in which it is wrong are known and specific, and each failure points to a missing ingredient and to the next model.
\end{enumerate}

The tests exclude a great deal. A model that reproduces nothing observed is mathematics, and may be good mathematics, but it is not a model of deep learning. A heuristic that cannot be made rigorous may be right, and when it is the best available we include it, labelled as such; but the core of every chapter is a theorem. A model that predicts nothing quantitative cannot be tested, and a model whose failures are unknown has not been tested. The fourth test is the one most often skipped, and it is the one this book takes most seriously.

\addsec{A first hydrogen atom: gradient descent on a quadratic}

The simplest solvable model of training is the harmonic oscillator of optimization: gradient descent on a quadratic loss. It is worth solving carefully, because almost every chapter of this book is about one of its failures.

Let $H$ be a symmetric positive semidefinite $p\times p$ matrix with eigenvalues $\lambda_1\ge\dots\ge\lambda_p\ge0$ and orthonormal eigenvectors $e_1,\dots,e_p$, and let
\[
  \Loss(\theta)=\tfrac12\,(\theta-\theta_\star)^\top H(\theta-\theta_\star).
\]
Gradient descent with step size $\eta>0$ is the iteration $\theta_{t+1}=\theta_t-\eta\nabla\Loss(\theta_t)=\theta_t-\eta H(\theta_t-\theta_\star)$. Least squares is the main example: if $X\in\R^{n\times p}$ and $y\in\R^n$, then $\frac1{2n}\norm{X\theta-y}^2$ equals a quadratic of this form plus a constant, with $H=X^\top X/n$ and $\theta_\star$ a least-squares solution.

\begin{theorem*}[Normal modes of gradient descent]
For every initialization $\theta_0$ and every $t\ge0$, the error $\theta_t-\theta_\star$ has coordinates
\[
  \ip{\theta_t-\theta_\star,e_k}=(1-\eta\lambda_k)^t\,\ip{\theta_0-\theta_\star,e_k}\qquad(k=1,\dots,p),
\]
and the loss is $\Loss(\theta_t)=\frac12\sum_k\lambda_k(1-\eta\lambda_k)^{2t}\ip{\theta_0-\theta_\star,e_k}^2$. Consequently:
\begin{enumerate}
\item The component along $e_k$ converges to zero for every initialization if and only if $0<\eta\lambda_k<2$; it stays constant if $\lambda_k=0$, and it grows without bound, with alternating sign, if $\eta\lambda_k>2$ and it is nonzero initially.
\item If $H$ is positive definite, gradient descent converges for every initialization if and only if $\eta<2/\lambda_1$, and the step size $\eta=2/(\lambda_1+\lambda_p)$ gives $\norm{\theta_t-\theta_\star}\le\big(\frac{\kappa-1}{\kappa+1}\big)^t\norm{\theta_0-\theta_\star}$ with $\kappa=\lambda_1/\lambda_p$.
\end{enumerate}
\end{theorem*}

\begin{proof}
The error $\delta_t=\theta_t-\theta_\star$ satisfies $\delta_{t+1}=(I-\eta H)\delta_t$, so $\delta_t=(I-\eta H)^t\delta_0$. The matrix $I-\eta H$ has the eigenvectors $e_k$ with eigenvalues $1-\eta\lambda_k$, which gives the first formula, and the loss is $\frac12\sum_k\lambda_k\ip{\delta_t,e_k}^2$. A geometric sequence $(1-\eta\lambda_k)^tc$ with $c\ne0$ tends to zero exactly when $|1-\eta\lambda_k|<1$, that is, when $0<\eta\lambda_k<2$; it is constant when $\lambda_k=0$, and when $\eta\lambda_k>2$ its ratio is below $-1$, which is part (i). For part (ii), all $\lambda_k$ are positive, so (i) applies to every component, and the largest eigenvalue gives the binding constraint. With $\eta=2/(\lambda_1+\lambda_p)$ every factor satisfies $|1-\eta\lambda_k|\le(\lambda_1-\lambda_p)/(\lambda_1+\lambda_p)=(\kappa-1)/(\kappa+1)$, because $1-\eta\lambda_k$ decreases in $\lambda_k$ from $(\lambda_1-\lambda_p)/(\lambda_1+\lambda_p)$ at $\lambda_k=\lambda_p$ to its negative at $\lambda_k=\lambda_1$; orthogonality of the $e_k$ gives the bound on $\norm{\delta_t}$.
\end{proof}

Three predictions follow, and each will return.

\emph{A stability threshold.} The step size must stay below $2/\lambda_1$, twice the inverse of the largest curvature. Above it, the stiffest mode oscillates with growing amplitude (\cref{fig:prologue:gd}).

\emph{Learning by modes.} Different directions are learned at different speeds. In the limit of small steps, gradient flow $\dot\theta=-H(\theta-\theta_\star)$ gives $\ip{\theta(t)-\theta_\star,e_k}=e^{-\lambda_kt}\ip{\theta(0)-\theta_\star,e_k}$, so the mode $e_k$ is learned on the time scale $1/\lambda_k$: the large-eigenvalue directions first and the small ones last. In kernel regression the large-eigenvalue directions are the smooth functions, and this simple fact becomes the \emph{spectral bias} of \cref{ch:lazy}.

\emph{An implicit choice.} When $H$ is singular, the minimizers form an affine subspace, and gradient descent picks one of them without being told which. The components of the error along the null space of $H$ never change, and when $\eta$ is below $2/\lambda_1$ the other components tend to zero, so the limit is the minimizer closest to the initialization. For least squares with $n<p$, started at $\theta_0=0$, every iterate lies in the row space of $X$, because every gradient $X^\top(X\theta-y)/n$ does; the limit is therefore the minimum-norm interpolator $X^+y$. This interpolator is the protagonist of \cref{ch:doubledescent}, and \cref{ch:diagonal} shows how changing the parameterization changes the choice.

\begin{figure}[t]
\centering
\begin{fullwidth}
\centering
\begin{tikzpicture}
\begin{groupplot}[group style={group size=2 by 1,horizontal sep=16mm},hplot,width=80mm,height=62mm]
\nextgroupplot[xlabel={$\theta_1$},ylabel={$\theta_2$},xmin=-2.9,xmax=0.6,ymin=-1.15,ymax=1.15]
\pgfplotsinvokeforeach{0.3,1.2,2.7,4.9}{\addplot[seamist!80!sea,line width=0.4pt,domain=0:360,samples=90] ({sqrt(2*#1)*cos(x)},{sqrt(2*#1/10)*sin(x)});}
\addplot[sea,line width=0.9pt,mark=*,mark size=0.9pt] table[x=x,y=y] {figures/data/00-prologue-gd-a.dat};
\addplot[sandline!55!black,line width=0.8pt,mark=*,mark size=0.9pt] table[x=x,y=y,restrict expr to domain={\thisrow{t}}{0:14}] {figures/data/00-prologue-gd-b.dat};
\addplot[seadeep,line width=0.7pt,densely dotted,mark=*,mark size=0.9pt] table[x=x,y=y,restrict expr to domain={\thisrow{t}}{0:5}] {figures/data/00-prologue-gd-c.dat};
\nextgroupplot[xlabel={step $t$},ylabel={loss $\Loss(\theta_t)$},ymode=log,xmin=0,xmax=40,ymin=1e-7,ymax=1e3,legend pos=south west]
\addplot[theory] table[x=t,y=exact] {figures/data/00-prologue-gd-a.dat};
\addplot[theory3] table[x=t,y=exact] {figures/data/00-prologue-gd-b.dat};
\addplot[theory2] table[x=t,y=exact] {figures/data/00-prologue-gd-c.dat};
\addplot[simulation,each nth point=3] table[x=t,y=loss] {figures/data/00-prologue-gd-a.dat};
\addplot[simulation,each nth point=3] table[x=t,y=loss] {figures/data/00-prologue-gd-b.dat};
\addplot[simulation,each nth point=3] table[x=t,y=loss] {figures/data/00-prologue-gd-c.dat};
\legend{$\eta=0.05$,$\eta=0.17$,$\eta=0.205$}
\end{groupplot}
\end{tikzpicture}
\end{fullwidth}
\caption{Gradient descent on $\Loss(\theta)=\frac12(\theta_1^2+10\,\theta_2^2)$, for which $2/\lambda_1=0.2$. Left: iterates from $\theta_0=(-2.6,0.55)$ over level sets of the loss; with $\eta=0.05$ the path is smooth, with $\eta=0.17$ the stiff coordinate flips sign at every step but shrinks, and with $\eta=0.205$ it flips and grows (first five steps shown). Right: the loss, computed by running gradient descent (points) and from the normal-mode formula (curves); they agree to machine precision, which is what ``solvable'' means.}
\label{fig:prologue:gd}
\end{figure}

\addsec{What the quadratic gets wrong}

The quadratic is exact for linear regression. As a model of training a network it fails in at least five ways, and the list of its failures is close to a table of contents for this book.

\paragraph{The curvature is not constant.} The Hessian of a network's loss changes during training. The theorem forbids step sizes above $2/\lambda_1$, but networks are routinely trained with step sizes for which the curvature at initialization, or later in training, exceeds $2/\eta$, and the training does not diverge: the curvature adapts to the step size. A model with only two parameters already shows how (\cref{ch:stability}).

\paragraph{The loss is not convex.} A network's loss has saddle points, where the gradient vanishes but the loss is not minimal, and training slows down near them. A quadratic has no saddles, so it cannot produce the plateaus of training curves. Escaping from saddles is the subject of the chapters on deep linear networks (\cref{ch:deeplinear}), diagonal networks (\cref{ch:diagonal}), single-index models (\cref{ch:singleindex}), committee machines (\cref{ch:committee}), sparse parities (\cref{ch:parity}) and induction heads (\cref{ch:induction}).

\paragraph{The minimizer is not unique.} A network that can interpolate its training data has a manifold of minimizers, and which one training finds is decided by the parameterization, the initialization and the optimizer. In the quadratic this choice is the trivial projection described above. In networks it produces sparsity (\cref{ch:diagonal}), maximum margins (\cref{ch:margin}), Fourier features (\cref{ch:grokking}) and the simplex geometry of neural collapse (\cref{ch:collapse}).

\paragraph{The data have no structure.} In least squares the matrix $H$ is the covariance of the inputs, and its spectrum is all the quadratic knows about the data. That is enough for some purposes: the singular values of a random design produce double descent (\cref{ch:doubledescent}), and power-law spectra produce power-law learning curves (\cref{ch:scaling}). It is not enough for others: a target that depends on one hidden direction of a high-dimensional input (\cref{ch:singleindex}), or on the parity of a few bits (\cref{ch:parity}), is invisible to second-order statistics.

\paragraph{There are no features.} A quadratic in the parameters is a linear model, and a linear model has fixed features. The central fact about trained networks is that their features change. Wide networks can be made to behave like linear models (\cref{ch:lazy}), and the price of doing so is a provable failure to learn targets that networks with moving features learn easily. The chapters on what networks represent (\cref{ch:grokking,ch:superposition,ch:collapse}) and on transformers (\cref{ch:icl,ch:induction}) are about features that training builds.

\addsec{How each chapter is built}

Every chapter after this one has the same anatomy.

\emph{The phenomenon} states what is observed in trained networks, by whom and with what evidence, in an \emph{observation} environment that separates empirical regularities from theorems, and ends with the question the model will answer.

\emph{The model} defines the smallest solvable model of the phenomenon, says what was thrown away and why, and gives a \emph{dictionary} that translates the model's quantities into the quantities of a trained network.

\emph{The solution} proves the main theorem completely. Further results come with complete proofs where they fit, with proof sketches that say what is omitted, or with heuristic derivations. The three are marked differently: a proof ends with $\blacksquare$, a sketch is labelled as one, and a heuristic derivation, the kind of argument that physicists trust and mathematicians do not yet accept, ends with $\Diamond$. The difference matters. A heuristic derivation can be right and useful, and several of the most informative predictions in this book are derived rather than proved, but a reader should always know which is which.

\emph{Reading the solution} turns the theorem into predictions and compares them with simulations. Every chapter has at least one figure in which a curve is the model's exact prediction and the points are a simulation, produced by a script that is distributed with the book.

\emph{What the model gets wrong} has three parts: what was thrown away, where the model breaks (specific theorems or experiments that contradict it), and the fine structure (the corrections and successor models that repair it, and what is still open).

A \emph{model card} summarizes the chapter in six lines: phenomenon, model, solution, prediction, failure and successors. The exercises come in four kinds: \emph{consolidation} exercises check and extend the chapter's results; \emph{variation} exercises change one assumption of the model and ask for the new solution, which is the skill this book most wants to teach; \emph{simulation} exercises ask for a numerical experiment and say what it should show; \emph{research} exercises point to open problems. The notes at the end of each chapter close with a paragraph that translates the notation of the main source papers into the notation of this book.

\addsec{A map of the book}

\Cref{tab:prologue:map} lists the models. The six parts move from models that are linear in their parameters (Part I), through the two infinite-width limits (Part II), to models in which features are learned (Part III) and in which the learned structure can be read off the weights (Part IV), then to transformers (Part V) and to large step sizes and large models (Part VI). The epilogue looks across the chapters at what the models share and at what they get wrong together, and lists phenomena that still lack a hydrogen atom.

\begin{table}[p]
\begin{fullwidth}
\centering\small
\renewcommand{\arraystretch}{1.18}
\begin{tabular}{@{}r>{\raggedright\arraybackslash}p{28mm}>{\raggedright\arraybackslash}p{33mm}>{\raggedright\arraybackslash}p{41mm}>{\raggedright\arraybackslash}p{38mm}@{}}
\toprule
& \textsc{phenomenon} & \textsc{model} & \textsc{solution} & \textsc{gets wrong}\\
\midrule
1 & Double descent & Least squares with Gaussian features & Exact risk at every $n$ and $p$; ridge in the proportional limit & Features are fixed and isotropic\\
2 & Learning in stages & Deep linear networks & Exact decoupled mode dynamics, logistic for two layers & Only linear maps; aligned initialization\\
3 & Implicit bias of small initialization & Diagonal linear networks & A mirror flow interpolating between $\ell_2$ and $\ell_1$ & A linear predictor in disguise\\
4 & Training past zero error & Linear classifiers, exponential loss & The margin tends to the maximum margin & Logarithmic rates; networks reach only KKT points\\
\midrule
5 & Depth needs critical initialization & Wide random networks & An order--chaos phase diagram & Finite width matters when depth is comparable\\
6 & Wide networks as kernels & Linearization at initialization & Lazy training; a lower bound for every kernel & No feature learning\\
\midrule
7 & A search phase before learning & Single-index models & Online SGD needs about $d^{k-1}$ samples & Reusing data beats the exponent\\
8 & Plateaus and specialization & The soft committee machine & Exact ODEs; a symmetric saddle & Gaussian inputs and tiny steps\\
9 & Abrupt learning, hidden progress & Sparse parities & A statistical-query bound; one gradient step finds the support & Natural tasks have low-degree signal\\
\midrule
10 & Grokking & Modular addition & A Fourier network; margin-maximizing neurons are Fourier modes & The mechanism varies between runs\\
11 & Superposition & A ReLU toy autoencoder & Linear models never superpose; sparse features force it & Features are given and independent\\
12 & Neural collapse & Unconstrained features & Minimizers are simplex frames & Features are not free\\
\midrule
13 & In-context learning & Linear attention on regression & The optimum is one preconditioned gradient step & No softmax; Gaussian tasks\\
14 & Induction heads & Attention on random Markov chains & Bayes predictor; a two-layer circuit; a saddle & One circuit among many\\
\midrule
15 & The edge of stability & The product $uv$ & An exact law for the imbalance forces the catapult & No progressive sharpening\\
16 & Scaling laws & Power-law spectra and Zipf's law & Exponents from tails; compute-optimal allocation & The exponents are inputs\\
\bottomrule
\end{tabular}
\caption{The sixteen models, each with the phenomenon it reproduces, the main result proved in its chapter, and its most important failure.}
\label{tab:prologue:map}
\end{fullwidth}
\end{table}

\addsec{How to read this book}

The prerequisites are undergraduate analysis, linear algebra and probability, and familiarity with gradient descent and with neural networks at the level of an introductory course. The chapters are written to be read independently, and each builds what it uses; tools that several chapters share, such as Hermite polynomials, the Marchenko--Pastur law, martingale inequalities and the method of ordinary differential equations for stochastic recursions, are collected in \cref{ch:toolkit}. The notation is uniform across the book and is summarized before this prologue. A one-semester course can take Parts I through III in order; a course on interpretability can start with Part IV and continue with Part V; a reader interested in optimization can read this prologue and \cref{ch:deeplinear,ch:diagonal,ch:margin,ch:stability}.

Three cautions. First, a solvable model is not an explanation of deep learning. It is an explanation of one phenomenon in one idealized setting, and the claim that the same mechanism operates in a trained network is an empirical claim that each chapter examines separately. Second, the models in this book were chosen because they are solvable, and solvable models are biased toward phenomena that have simple mechanisms; the epilogue discusses what this bias hides. Third, the field moves quickly. The models here are the ones that have earned a place by 2026, and several chapters end with successor models that may soon deserve chapters of their own.

\notesandrefs

The history of the hydrogen spectrum, from Balmer's formula to quantum electrodynamics, is told in most textbooks of atomic physics; Balmer's formula, Bohr's model and Schr\"odinger's solution are in the original papers cited above, and Lamb and Retherford's measurement and Bethe's calculation appeared in the same volume of the \emph{Physical Review} in 1947. Onsager's solution of the two-dimensional Ising model \parencite{Onsager1944} and Yang's computation of its spontaneous magnetization \parencite{Yang1952} are the classic exact results; the failure of mean-field exponents in low dimension is the starting point of every account of critical phenomena. The role of anharmonicity in thermal expansion and thermal conductivity is in \textcite{AshcroftMermin1976}.

The convergence theory of gradient descent on quadratics, including the optimal constant step size, is standard; see \textcite{Polyak1987} and \textcite{Nesterov2018}. The observation that gradient descent on underdetermined least squares from zero converges to the minimum-norm solution is folklore in numerical analysis and became central to deep learning theory through the study of implicit regularization.

Four books are close to this one in spirit. \Textcite{EngelVandenBroeck2001} present the statistical mechanics of learning, including several models that reappear here in modern form. \Textcite{RobertsYaidaHanin2022} develop a single method, an effective field theory of wide networks organized as an expansion in the inverse width, and carry it much further than a single chapter could. \Textcite{Huang2021} covers the classical statistical-physics models of neural networks, such as spin glasses, the Hopfield model and the perceptron. \Textcite{Zhao2026}, a monograph posted on arXiv, gives a broad, proof-oriented account of deep learning theory in which each topic has about one chapter. The position paper of \textcite{SimonEtAl2026} argues that a scientific theory of deep learning is emerging and names solvable idealized settings as one of its five strands; this book is a textbook for that strand.

\printchapterbib
\end{refsection}


\casepart{Linear Models That Already Surprise}{The first models are linear in their inputs. They already fit noise and still predict well, learn in stages, choose sparse or maximum-margin solutions that nobody asked for, and show a test error that falls, rises and falls again. Linearity makes each of these behaviours solvable exactly.}{18}
% Chapter 1 of The Hydrogen Atoms of Deep Learning: double descent in least squares.
\casechapter{Double Descent}{Least squares in the proportional limit}{Ridgeless least squares, Gaussian features}{ch:doubledescent}
\begin{refsection}[chapters/01-doubledescent.bib]

\begin{chaptersummary}
As a model grows past the size at which it first fits its training data exactly, its test error can rise to a peak and fall again, sometimes below the bottom of the classical U. Least squares on isotropic Gaussian features, solved by the minimum-norm interpolator that gradient descent finds, does exactly this. We prove an exact formula for its test error at every sample size from two symmetries of the Gaussian design, and use the Marchenko--Pastur law to compute the risk of ridge regression, whose optimal tuning removes the peak. The model gets the shape right and the axis wrong: networks peak at an effective complexity rather than at a parameter count, the epoch-wise curve needs dynamics, and isotropic features never overfit benignly.
\end{chaptersummary}

\section{The phenomenon: a second descent}\label{sec:doubledescent:phenomenon}

Statistics teaches a U: as a model grows, its training error falls, while its test error first falls and then rises, because small models are biased and large ones fit the noise. A corollary is that a model that fits its training data exactly, noise included, should predict badly.

Deep networks violate the corollary routinely. They have more parameters than training examples, reach zero training error, and still predict well, although they can fit images with random labels perfectly \parencite{ZhangEtAl2017}. \Textcite{BelkinEtAl2019} reconciled the two pictures with one curve, which they named \term{double descent}\index{double descent}. As the capacity grows, the test error follows the classical U, peaks at the \term{interpolation threshold}\index{interpolation threshold}, the smallest capacity at which the model fits the training set exactly, and then falls again, in their experiments often below the bottom of the U. They showed this for random Fourier and random ReLU features fitted by minimum-norm least squares (on $10^4$ MNIST digits the peak sits at $10^4$ features), for random forests and $L_2$-boosting, and for networks with one hidden layer, whose threshold appeared near $nK$ parameters for $K=10$ classes: one parameter per fitted number.

\Textcite{NakkiranEtAl2020} found the same curve for ResNets and convolutional networks on CIFAR and for transformers trained to translate, and two more versions of it. In \term{epoch-wise double descent}\index{double descent!epoch-wise} the test error falls, rises, and falls again during training. In \term{sample-wise double descent}\index{double descent!sample-wise} more data hurts: for some transformer sizes, four and a half times more training sentences gave a worse test loss. They measured a training procedure by its \term{effective model complexity}\index{effective model complexity}, the largest number of examples on which it reaches near-zero training error, and found the peak where this equals the size of the training set; label noise, data augmentation, and more examples all move the peak toward larger models.

\begin{observation}[\textcite{BelkinEtAl2019}; \textcite{NakkiranEtAl2020}]\label[observation]{obs:doubledescent:dd}
On a fixed training set, train models of increasing capacity, each to (near) zero training error when it can. The test error follows the classical U until shortly before the interpolation threshold, peaks near the threshold, and decreases beyond it, sometimes below its classical minimum. A curve of the same shape appears as a function of training time, and, near the threshold, of the number of training examples.
\end{observation}

How robust is this? The location of the peak at the threshold is reliable; its height is not. In the experiments of \textcite{NakkiranEtAl2020} every form of double descent was strongest with label noise: on clean CIFAR-10 the model-wise curve showed only a plateau around the threshold, although clean CIFAR-100 and the transformers did show peaks. Regularization usually tames the peak (the cusp of \textcite{SpiglerEtAl2019} disappears with early stopping, and tuned weight decay makes the test error of the networks of \textcite{NakkiranEtAl2021} largely monotone), but in one setting of \textcite{NakkiranEtAl2020}, a ResNet on clean CIFAR-100, a peak survives optimal early stopping. Models trained for one pass over their data, as large language models are, never interpolate it; in language models double descent has been reported in training time when a small dataset is repeated for hundreds of epochs \parencite{MuennighoffEtAl2023}.

The curve is older than its name \parencite{LoogEtAl2020}.\sidenote{\Textcite{LoogEtAl2020} also list the peaking of the pseudo-Fisher linear discriminant in pattern recognition, analyzed by Raudys and Duin in 1998.} For linear models fitted with the pseudo-inverse, \textcite{ValletEtAl1989} saw a peak at $n=p$ in simulated learning curves, and \textcite{OpperEtAl1990} computed it exactly in the limit of large dimension. \Textcite{KroghHertz1992} showed that a weight decay matched to the noise suppresses the overfitting, \textcite{AdvaniEtAl2020} traced the peak to the small eigenvalues of the Marchenko--Pastur law, and \textcite{SpiglerEtAl2019} identified the threshold of deep networks with a jamming transition, the point at which a granular packing becomes rigid.

Three questions remain. Why is there a peak exactly where the model first fits its training data? Why does the test error come down past it, and how far? And is the second descent a property of the data, of the estimator, or of how capacity is counted? The model of this chapter answers all three exactly, for linear regression.

\section{The model: least squares on Gaussian features}\label{sec:doubledescent:model}

The simplest learner with a sharp interpolation threshold is linear least squares: with $p$ features and $n$ examples in general position, it cannot fit arbitrary labels when $p<n$, and it fits every set of labels exactly when $p\ge n$. We take the features to be independent standard Gaussians, the target linear, and the noise additive.

\begin{assumption}[Gaussian design, additive noise]\label[assumption]{ass:doubledescent:design}
The inputs $x_1,\dots,x_n$ are independent with law $\cN(0,I_p)$ and form the rows of the design matrix $X\in\R^{n\times p}$. The responses are $y=Xw_\star+\xi$, with a fixed target $w_\star\in\R^p$ of norm $r=\norm{w_\star}$ and a noise vector $\xi\in\R^n$ that is independent of $X$ and satisfies $\E\xi=0$ and $\E\xi\xi^\top=\sigma^2I_n$.
\end{assumption}

The noise need not be Gaussian nor have independent coordinates, which lets one theorem cover the second model below.

\paragraph{Test error.} A test point is a fresh pair $(x,y)$ with $x\sim\cN(0,I_p)$ and $y=\ip{w_\star,x}+\eps$, the noise $\eps$ having mean $0$ and variance $\sigma^2$ and being independent of $x$ and of the training data. Since $y-\ip{w,x}=\eps+\ip{w_\star-w,x}$ is a sum of two uncorrelated terms and $\E[xx^\top]=I_p$, the population loss is
\begin{equation}\label{eq:doubledescent:loss}
  \Loss(w)=\E\bigl(y-\ip{w,x}\bigr)^2=\sigma^2+\norm{w-w_\star}^2 .
\end{equation}
The \term{test error}\index{test error} of an estimator $\hat w$ is $\E\Loss(\hat w)$: the noise floor $\sigma^2$ plus the \term{risk}\index{risk} $\E\norm{\hat w-w_\star}^2$. The predictor $w=0$ has the \term{null risk}\index{null risk} $\Loss(0)=\sigma^2+r^2$.

\paragraph{Two estimators.} The training loss is $\Lhat(w)=\frac1{2n}\norm{y-Xw}^2$, and the \term{minimum-norm least-squares estimator}\index{minimum-norm interpolator} $\hat w=X^+y$ is its minimizer of smallest norm.\sidenote{The matrix $X^+$ is the Moore--Penrose pseudo-inverse of $X$, and in the physics literature $X^+y$ is the pseudo-inverse rule.} If $X$ has rank $p$ (so $p\le n$), then the minimizer is unique, $\hat w=(X^\top X)^{-1}X^\top y$. If $X$ has rank $n$ (so $p\ge n$), then the minimizers are the \term{interpolators}, the solutions of $Xw=y$; they differ from $X^\top(XX^\top)^{-1}y$, which lies in the row space of $X$, by vectors of $\ker X$, which are orthogonal to the row space, so $\hat w=X^\top(XX^\top)^{-1}y$. One of the two cases occurs with probability one (\cref{lem:doubledescent:wishart}).\sidenote{Near the threshold the mean of $\tr\bigl((X^\top X)^{-1}\bigr)$ in \cref{lem:doubledescent:wishart} is finite but its variance is not: by \cref{ex:doubledescent:wishart} it is infinite for $p\in\{n-3,n-2\}$, so simulated risks converge slowly there.} The \term{ridge estimator}\index{ridge regression} with $\lambda>0$ minimizes $\Lhat(w)+\frac\lambda2\norm w^2$; with the sample covariance $\hat\Sigma=\frac1nX^\top X$,
\begin{equation}\label{eq:doubledescent:ridge}
  \hat w_\lambda=(X^\top X+n\lambda I_p)^{-1}X^\top y=(\hat\Sigma+\lambda I_p)^{-1}\frac{X^\top y}{n},
\end{equation}
and $\hat w_\lambda\to X^+y$ as $\lambda\to0$, which is why $X^+y$ is also called \term{ridgeless}.\sidenote{Along an eigenvector of $\hat\Sigma$ with eigenvalue $t>0$, ridge multiplies the least-squares coefficient by $t/(t+\lambda)$, and gradient flow on $\Lhat$ from $w=0$, run for time $\tau$, multiplies it by $1-e^{-t\tau}$. The two filters are close when $\lambda\approx1/\tau$; this is the sense of the dictionary entry for $\lambda$.}

Why the minimum norm? Started at $w=0$, gradient descent on $\Lhat$ moves along gradients $-\frac1nX^\top(y-Xw)$, which lie in the row space of $X$; if it converges to a minimizer, then the limit is the only minimizer in the row space, $X^+y$ (\cref{ex:doubledescent:gd}). Past the threshold the training algorithm, not the loss, chooses the interpolator.\sidenote{Change the parameterization and the choice changes: gradient descent on $w=u\odot u-v\odot v$ from a small initialization finds an interpolator of small $\ell_1$ norm instead (\cref{ch:diagonal}).}

\paragraph{Weak features.} So far $p$ is both the number of parameters and the dimension of the truth, so changing $p$ changes the problem as well as the learner. \Textcite{BelkinEtAl2020} proposed a version in which only the learner changes.

\begin{assumption}[Weak features]\label[assumption]{ass:doubledescent:weak}
The input is $z\sim\cN(0,I_d)$, and $y=\ip{w_\star,z}+\eps$ with $w_\star\in\R^d$ and noise $\eps$ of mean $0$ and variance $\sigma^2$, independent of $z$. On a fixed set $S\subseteq\{1,\dots,d\}$ of $p$ coordinates, the learner fits $\hat w=X^+y$, the rows of $X\in\R^{n\times p}$ being the $z_{i,S}$ of $n$ independent examples, and predicts $\ip{\hat w,z_S}$.
\end{assumption}

The omitted coordinates are independent of the observed ones, so their contribution $\ip{w_{\star,S^c},z_{S^c}}$ to the label looks to the learner exactly like noise; ``weak'' means that each feature carries a small part of the signal.

\begin{dictionary}
number of features $p$ & number of parameters, or better an effective number of fitted directions\\
number of examples $n$ & training examples (times outputs, for several outputs)\\
threshold $p=n$ & smallest model that fits the training set exactly\\
noise variance $\sigma^2$ & label noise, plus the part of the target the model cannot represent\\
minimum-norm interpolator $X^+y$ & the interpolator that gradient descent finds from small initialization\\
ridge parameter $\lambda$ & weight decay, or early stopping at time of order $1/\lambda$\\
smallest singular value of $X$ & conditioning of the interpolation problem\\
Marchenko--Pastur law & spectrum of the feature (or tangent-kernel) covariance at finite sample size\\
identity covariance $I_p$ & a strongly anisotropic feature covariance (\cref{ch:scaling})\\
null risk $\sigma^2+r^2$ & test error of the constant predictor\\
\end{dictionary}

\paragraph{What was thrown away, and why first.} The features are fixed, so nothing is learned about the representation; but the threshold already appears with fixed random features \parencite{BelkinEtAl2019}. They are isotropic, which makes the risk depend on $w_\star$ only through its norm and gives exact formulas at finite sample size. The estimator is the end point of training, because the model-wise curve is a statement about end points. The loss is square, so the estimator is linear in the labels and every expectation is a trace. \Cref{sec:doubledescent:wrong} says what each choice costs.

\paragraph{What one had to notice.} Three things. The peak sits where the model first fits the training set, and least squares has the sharpest such threshold, at $p=n$, where its solution changes from the unique best fit to one of many exact fits. Past the threshold the training algorithm selects the solution, and gradient descent from zero selects the minimum norm. And with a rotation-invariant design the target matters only through its norm, and the spectrum of $X^\top X$ carries everything else.

\section{The solution: exact risk at every sample size}\label{sec:doubledescent:exact}

Two symmetries of the Gaussian design make the risk computable at every $n$ and $p$: rotations of the feature space reduce the error along $w_\star$ to a trace, and permutations of the columns reduce the trace to one entry of an inverse Wishart matrix.

\begin{lemma}[Rotational averaging]\label[lemma]{lem:doubledescent:rotation}
Let $X\in\R^{n\times p}$ have independent $\cN(0,1)$ entries, and let $F$ map almost every $n\times p$ matrix to a $p\times p$ matrix so that $F(XO)=O^\top F(X)O$ for every orthogonal $O\in\R^{p\times p}$. If $\E\norm{F(X)}_F<\infty$, then $\E F(X)=\frac1p\E\tr F(X)\cdot I_p$.
\end{lemma}

\begin{proof}
For an orthogonal $O$, the rows of $XO$ are the vectors $O^\top x_i$, again independent with law $\cN(0,I_p)$, so $XO$ has the law of $X$. Hence $M=\E F(X)$ satisfies $M=\E F(XO)=O^\top MO$ for every orthogonal $O$. Taking for $O$ the diagonal matrix that flips the sign of the $i$th coordinate gives $M_{ij}=-M_{ij}$ for $j\ne i$, so $M$ is diagonal; taking the permutation matrix that swaps coordinates $i$ and $j$ gives $M_{ii}=M_{jj}$. So $M$ is a multiple of $I_p$, and comparing traces identifies the multiple.
\end{proof}

\begin{lemma}[Inverse moments of a Wishart matrix]\label[lemma]{lem:doubledescent:wishart}
Let $X\in\R^{n\times p}$, with $p\le n$, have independent $\cN(0,1)$ entries. Then $X$ has rank $p$ almost surely, and
\[
  \E\tr\bigl((X^\top X)^{-1}\bigr)=\begin{cases}\dfrac{p}{n-p-1}, & p\le n-2,\\[1.2ex] +\infty, & p\in\{n-1,n\}.\end{cases}
\]
\end{lemma}

\begin{proof}
The plan is to show that each diagonal entry of $(X^\top X)^{-1}$ is the reciprocal of a chi-square variable with $n-p+1$ degrees of freedom.

\emph{The rank.} Conditionally on the other columns, the span of the columns other than the $j$th is a fixed subspace of dimension at most $p-1<n$, contained in a hyperplane $\{u\colon\ip{a,u}=0\}$ with $a\ne0$, and the $j$th column $x^{(j)}$ satisfies $\ip{a,x^{(j)}}\sim\cN(0,\norm a^2)$, which vanishes with probability $0$. So almost surely no column lies in the span of the others.

\emph{One diagonal entry.} Write $X=[x^{(1)},X']$, and let $Q$ be the orthogonal projection onto the orthogonal complement of the column space of $X'$, which has rank $n-p+1$ by the rank statement applied to $X'$ (in the trivial case $p=1$, $X'$ is empty and $Q=I_n$). The block-inversion formula for the positive definite matrix $X^\top X$ gives
\[
  \bigl[(X^\top X)^{-1}\bigr]_{11}=\frac1{\norm{x^{(1)}}^2-x^{(1)\top}X'(X'^\top X')^{-1}X'^\top x^{(1)}}=\frac1{\norm{Qx^{(1)}}^2}.
\]
Conditionally on $X'$, the column $x^{(1)}$ is still $\cN(0,I_n)$; writing $Q=UU^\top$ with $U$ an $n\times(n-p+1)$ matrix with orthonormal columns, $\norm{Qx^{(1)}}^2=\norm{U^\top x^{(1)}}^2$ is a sum of $k=n-p+1$ squares of independent standard Gaussians. So it has the $\chi^2_k$ law, conditionally on $X'$ and hence unconditionally.

\emph{The chi-square integral.} Integrating $1/t$ against the $\chi^2_k$ density gives
\[
  \E\Bigl[\frac1{\chi^2_k}\Bigr]=\frac1{2^{k/2}\Gamma(k/2)}\int_0^\infty t^{k/2-2}e^{-t/2}\dd t ,
\]
which is finite if and only if $k\ge3$, by the behavior $t^{k/2-2}$ near $t=0$. For $k\ge3$ the integral is $2^{k/2-1}\Gamma(k/2-1)$, and $\Gamma(k/2)=(k/2-1)\Gamma(k/2-1)$ turns the expectation into $1/(k-2)$. With $k=n-p+1$, the condition $k\ge3$ reads $p\le n-2$, and $k-2=n-p-1$.

\emph{The trace.} Permuting the columns of $X$ does not change its law, so every diagonal entry of $(X^\top X)^{-1}$ has the law of the first. Summing $p$ equal expectations gives $p/(n-p-1)$ for $p\le n-2$, and $+\infty$ for $p\in\{n-1,n\}$.
\end{proof}

\begin{theorem}[Exact risk of the minimum-norm interpolator]\label{thm:doubledescent:exact}
Under \cref{ass:doubledescent:design} with $\sigma>0$, the estimator $\hat w=X^+y$ satisfies
\begin{equation}\label{eq:doubledescent:exact}
  \E\norm{\hat w-w_\star}^2=\begin{cases}
  \sigma^2\,\dfrac{p}{n-p-1}, & p\le n-2,\\[1.4ex]
  +\infty, & \abs{p-n}\le1,\\[0.8ex]
  r^2\Bigl(1-\dfrac np\Bigr)+\sigma^2\,\dfrac{n}{p-n-1}, & p\ge n+2.
  \end{cases}
\end{equation}
Without noise ($\sigma=0$) the risk is $0$ for $p\le n$ and $r^2(1-n/p)$ for $p\ge n$. In every case the test error is $\E\Loss(\hat w)=\sigma^2+\E\norm{\hat w-w_\star}^2$.
\end{theorem}

\begin{proof}
The plan is to compute the error conditionally on $X$, where it is a quadratic form in the noise and in $w_\star$, and then to average over $X$ with the two lemmas.

\emph{Below the threshold, $p\le n$.} By \cref{lem:doubledescent:wishart}, $X$ has rank $p$ almost surely, so $\hat w=(X^\top X)^{-1}X^\top(Xw_\star+\xi)=w_\star+(X^\top X)^{-1}X^\top\xi$. Since $\xi$ is independent of $X$ with $\E\xi\xi^\top=\sigma^2I_n$,
\[
  \E\bigl[\norm{\hat w-w_\star}^2\bigm| X\bigr]=\tr\Bigl((X^\top X)^{-1}X^\top\,\E[\xi\xi^\top]\,X(X^\top X)^{-1}\Bigr)=\sigma^2\tr\bigl((X^\top X)^{-1}\bigr),
\]
and \cref{lem:doubledescent:wishart} gives $\sigma^2p/(n-p-1)$ for $p\le n-2$, and $+\infty$ for $p\in\{n-1,n\}$ because $\sigma>0$. Without noise, $\hat w=w_\star$.

\emph{Above the threshold, $p\ge n$.} \Cref{lem:doubledescent:wishart} applied to the $p\times n$ matrix $X^\top$ shows that $X$ has rank $n$ almost surely. Let $P=X^\top(XX^\top)^{-1}X$ be the orthogonal projection onto the row space of $X$. Then $\hat w=X^\top(XX^\top)^{-1}(Xw_\star+\xi)=Pw_\star+X^\top(XX^\top)^{-1}\xi$, and the error
\[
  \hat w-w_\star=-(I_p-P)w_\star+X^\top(XX^\top)^{-1}\xi
\]
is the sum of a vector of $\ker X$ and a vector of the row space, which are orthogonal, so $\norm{\hat w-w_\star}^2=w_\star^\top(I_p-P)w_\star+\xi^\top(XX^\top)^{-1}\xi$. The row space of $XO$ is $O^\top$ times the row space of $X$, so $P(XO)=O^\top P(X)O$, and $\tr P=\norm P_F^2=n$; \cref{lem:doubledescent:rotation} gives $\E P=\frac npI_p$, so the first term has mean $r^2(1-n/p)$. The second term has conditional mean $\sigma^2\tr\bigl((XX^\top)^{-1}\bigr)$, and \cref{lem:doubledescent:wishart}, with the roles of $n$ and $p$ exchanged, gives $\sigma^2n/(p-n-1)$ for $p\ge n+2$, and $+\infty$ for $p\in\{n,n+1\}$. Without noise only the first term remains.

\emph{Test error.} The estimator $\hat w$ is independent of the test pair, so \eqref{eq:doubledescent:loss} holds conditionally on the training data; take expectations.
\end{proof}

\begin{corollary}[Weak features; \textcite{BreimanFreedman1983}, \textcite{BelkinEtAl2020}]\label[corollary]{cor:doubledescent:weak}
Under \cref{ass:doubledescent:weak}, let $\sigma_S^2=\sigma^2+\norm{w_{\star,S^c}}^2>0$. The minimum-norm estimator on the features in $S$ has test error
\[
  \E\bigl(y-\ip{\hat w,z_S}\bigr)^2=\begin{cases}
  \sigma_S^2\,\dfrac{n-1}{n-p-1}, & p\le n-2,\\[1.4ex]
  +\infty, & \abs{p-n}\le1,\\[0.8ex]
  \norm{w_{\star,S}}^2\Bigl(1-\dfrac np\Bigr)+\sigma_S^2\,\dfrac{p-1}{p-n-1}, & p\ge n+2.
  \end{cases}
\]
\end{corollary}

\begin{proof}
Write the labels as $y_i=\ip{w_{\star,S},z_{i,S}}+\xi_i$ with $\xi_i=\ip{w_{\star,S^c},z_{i,S^c}}+\eps_i$. The coordinates of each $z_i$ are independent, so $\xi$ is independent of $X$, and its entries are independent with mean $0$ and variance $\sigma_S^2$. \Cref{ass:doubledescent:design} holds with target $w_{\star,S}$ and noise level $\sigma_S$, and at a test point the same decomposition gives the test error $\sigma_S^2+\E\norm{\hat w-w_{\star,S}}^2$. Insert \eqref{eq:doubledescent:exact} and use $1+\frac p{n-p-1}=\frac{n-1}{n-p-1}$ and $1+\frac n{p-n-1}=\frac{p-1}{p-n-1}$.
\end{proof}

Each added feature moves signal from $\sigma_S^2$ into $\norm{w_{\star,S}}^2$, where it can be learned; this creates a classical U below the threshold (\cref{sec:doubledescent:reading}).

The \term{proportional limit}\index{proportional limit} keeps $p/n$ fixed while both grow, so that bias and variance both stay of order one.\sidenote{It is the thermodynamic limit of the statistical mechanics of learning, in which the number of examples per weight, $\alpha=n/p=1/\gamma$, is held fixed.}

\begin{corollary}[Proportional limit]\label[corollary]{cor:doubledescent:limit}
Under \cref{ass:doubledescent:design}, let $n\to\infty$ and $p/n\to\gamma\in(0,\infty)$, with $r$ and $\sigma>0$ fixed (the target $w_\star\in\R^p$ changes with $p$ but keeps its norm). Then $\E\norm{\hat w-w_\star}^2\to\cR_0(\gamma)$, with
\begin{equation}\label{eq:doubledescent:limit}
  \cR_0(\gamma)=\begin{cases}\sigma^2\,\dfrac{\gamma}{1-\gamma}, & \gamma<1,\\[1.2ex] +\infty, & \gamma=1,\\[0.6ex] r^2\Bigl(1-\dfrac1\gamma\Bigr)+\dfrac{\sigma^2}{\gamma-1}, & \gamma>1.\end{cases}
\end{equation}
\end{corollary}

\begin{proof}
For $\gamma\ne1$, eventually $\abs{p-n}\ge2$, and $p<n$ exactly when $\gamma<1$; divide numerators and denominators in \eqref{eq:doubledescent:exact} by $n$ and let $p/n\to\gamma$. For $\gamma=1$, the bounds $\frac p{n-p-1}\ge\frac1{n/p-1}$ for $p\le n-2$ and $\frac n{p-n-1}\ge\frac1{p/n-1}$ for $p\ge n+2$ show that the risk tends to $+\infty$, and it is infinite when $\abs{p-n}\le1$.
\end{proof}

\section{The solution: ridge and the Marchenko--Pastur law}\label{sec:doubledescent:ridge}

The ridge risk depends on the whole spectrum of $\hat\Sigma$, which has a deterministic limit in the proportional regime.

\paragraph{The Marchenko--Pastur law.} For $\gamma>0$ let $t_\pm=(1\pm\sqrt\gamma)^2$, and let $\mu_\gamma$ be the probability measure on $[0,\infty)$ given by
\begin{equation}\label{eq:doubledescent:mp}
  \mu_\gamma(\dd t)=\Bigl(1-\frac1\gamma\Bigr)_+\delta_0(\dd t)+\frac{\sqrt{(t_+-t)(t-t_-)}}{2\pi\gamma\,t}\,\ind\{t_-\le t\le t_+\}\dd t .
\end{equation}
Its \term{Stieltjes transform}\index{Stieltjes transform} $m_\gamma(z)=\int\frac{\mu_\gamma(\dd t)}{t-z}$, defined for $z\in\C\setminus[0,\infty)$, is the solution of
\begin{equation}\label{eq:doubledescent:mpquad}
  \gamma z\,m^2+(z-1+\gamma)\,m+1=0
\end{equation}
with $\Im m>0$ when $\Im z>0$, and $m>0$ when $z<0$.\sidenote{We checked \eqref{eq:doubledescent:mp} and \eqref{eq:doubledescent:mpquad} in this normalization numerically, integrating the density against $1/(t-z)$ for $\gamma\in\{0.3,0.8,1,1.5,4\}$ and several $z$: the residuals are below $10^{-10}$.} The empirical spectral distribution of $\hat\Sigma$ is $\mu_p=\frac1p\sum_{i=1}^p\delta_{\lambda_i(\hat\Sigma)}$, so that $\frac1p\tr f(\hat\Sigma)=\int f\dd\mu_p$ for every function $f$ on $[0,\infty)$.

Gaussianity matters for the exact formulas of \cref{sec:doubledescent:exact}, but not for the limits: \textcite{HastieEtAl2022} prove the limits \eqref{eq:doubledescent:limit} for the risk conditional on $X$, with features whose independent entries have mean zero, unit variance, and a finite moment of order $4+\eta$.\sidenote{The Marchenko--Pastur law itself holds for independent entries with two moments; the Gaussian case, proved in the appendix, is all we need.}

\begin{theorem}[\textcite{MarchenkoPastur1967}; proved in \cref{thm:toolkit:mp}]\label{thm:doubledescent:mp}\index{Marchenko--Pastur law}
Let $X\in\R^{n\times p}$ have independent $\cN(0,1)$ entries, and let $n\to\infty$ with $p/n\to\gamma\in(0,\infty)$. Then, almost surely, $\mu_p$ converges weakly to $\mu_\gamma$.
\end{theorem}

We use the theorem through two spectral functions of $\lambda>0$,\sidenote{Mar\v{c}enko and Pastur published the law in 1967 in \emph{Matematicheskii Sbornik}, for a more general family of random matrices; the English translation appeared the same year.}
\[
  v_p(\lambda)=\int\frac{\lambda\,\mu_p(\dd t)}{t+\lambda}=\frac1p\tr\bigl(\lambda(\hat\Sigma+\lambda I)^{-1}\bigr),\quad
  v(\lambda)=\int\frac{\lambda\,\mu_\gamma(\dd t)}{t+\lambda}=\lambda m_\gamma(-\lambda).
\]
Since $1-v_p(\lambda)$ is the effective number of parameters of ridge regression per parameter, $v_p$ is the fraction of directions decided by the penalty rather than by the data. Multiplying \eqref{eq:doubledescent:mpquad} at $z=-\lambda$ by $-\lambda$ shows that $v(\lambda)$ solves $\gamma v^2+(1-\gamma+\lambda)v-\lambda=0$, whose roots have product $-\lambda/\gamma<0$; so $v(\lambda)$ is the positive root,
\begin{equation}\label{eq:doubledescent:v}
  v(\lambda)=\frac{\sqrt{(1-\gamma+\lambda)^2+4\gamma\lambda}-(1-\gamma+\lambda)}{2\gamma}.
\end{equation}

\begin{lemma}[Ridge risk in spectral form]\label[lemma]{lem:doubledescent:ridgespectral}
Under \cref{ass:doubledescent:design}, for every $n$, $p$, and $\lambda>0$,
\begin{equation}\label{eq:doubledescent:ridgefinite}
  \E\norm{\hat w_\lambda-w_\star}^2=r^2\,\E\!\int\frac{\lambda^2\,\mu_p(\dd t)}{(t+\lambda)^2}+\sigma^2\,\frac pn\,\E\!\int\frac{t\,\mu_p(\dd t)}{(t+\lambda)^2},
\end{equation}
and the two integrals are $v_p(\lambda)-\lambda v_p'(\lambda)$ and $v_p'(\lambda)$. In particular the risk depends on $w_\star$ only through $r$.
\end{lemma}

\begin{proof}
Insert $y=Xw_\star+\xi$ into \eqref{eq:doubledescent:ridge} and use $(\hat\Sigma+\lambda I)^{-1}\hat\Sigma-I=-\lambda(\hat\Sigma+\lambda I)^{-1}$:
\[
  \hat w_\lambda-w_\star=-\lambda(\hat\Sigma+\lambda I)^{-1}w_\star+(\hat\Sigma+\lambda I)^{-1}\frac{X^\top\xi}n .
\]
Conditionally on $X$ the second term has mean zero, so the cross term vanishes in expectation, and its squared norm has conditional mean $\frac{\sigma^2}{n^2}\tr\bigl((\hat\Sigma+\lambda I)^{-1}X^\top X(\hat\Sigma+\lambda I)^{-1}\bigr)=\frac{\sigma^2}n\tr\bigl(\hat\Sigma(\hat\Sigma+\lambda I)^{-2}\bigr)$. The first term contributes $\lambda^2w_\star^\top F(X)w_\star$ with $F(X)=(\hat\Sigma+\lambda I)^{-2}$, which satisfies $F(XO)=O^\top F(X)O$ and $\norm{F(X)}_{\op}\le\lambda^{-2}$, so \cref{lem:doubledescent:rotation} gives $\E F(X)=\frac1p\E\tr F(X)\cdot I_p$. Writing both traces as integrals against $\mu_p$ gives \eqref{eq:doubledescent:ridgefinite}. Finally $v_p'(\lambda)=\int\frac{t}{(t+\lambda)^2}\dd\mu_p$, and $\frac{\lambda}{t+\lambda}-\frac{\lambda t}{(t+\lambda)^2}=\frac{\lambda^2}{(t+\lambda)^2}$.
\end{proof}

\begin{theorem}[Ridge regression in the proportional limit]\label{thm:doubledescent:ridge}
Under \cref{ass:doubledescent:design}, let $n\to\infty$ and $p/n\to\gamma\in(0,\infty)$, with $r$, $\sigma$, and $\lambda>0$ fixed. Then
\begin{equation}\label{eq:doubledescent:ridgelimit}
  \E\norm{\hat w_\lambda-w_\star}^2\to\cR_\lambda(\gamma)=r^2\bigl(v(\lambda)-\lambda v'(\lambda)\bigr)+\sigma^2\gamma\,v'(\lambda),
\end{equation}
with $v$ given by \eqref{eq:doubledescent:v}. In terms of the Stieltjes transform, the bias is $r^2\lambda^2m_\gamma'(-\lambda)$ and the variance is $\sigma^2\gamma\bigl(m_\gamma(-\lambda)-\lambda m_\gamma'(-\lambda)\bigr)$. As $\lambda\to0$, $\cR_\lambda(\gamma)\to\cR_0(\gamma)$ for every $\gamma\ne1$.
\end{theorem}

\begin{proof}
The plan is to pass to the limit in \cref{lem:doubledescent:ridgespectral}: \cref{thm:doubledescent:mp} gives the limit of $\E v_p$, convexity gives the limit of $\E v_p'$, and the quadratic equation evaluates both.

\emph{The function $v$.} The integrand $t\mapsto\lambda/(t+\lambda)$ is bounded and continuous on $[0,\infty)$, so $v_p(\lambda)\to v(\lambda)$ almost surely by \cref{thm:doubledescent:mp}, and $\E v_p(\lambda)\to v(\lambda)$ by dominated convergence, since $0\le v_p\le1$.

\emph{Its derivative.} For each $t\ge0$ the function $\lambda\mapsto-\lambda/(t+\lambda)$ is convex, with derivative $-t/(t+\lambda)^2$ of absolute value at most $1/(4\lambda)$, because $(t+\lambda)^2\ge4t\lambda$. Hence $f_n=-\E v_p$ is convex and differentiable on $(0,\infty)$, with $f_n'=-\E v_p'$, and it converges pointwise to $f=-v$, which is differentiable by \eqref{eq:doubledescent:v}. Convexity places $f_n'(\lambda)$ between the difference quotients $\frac{f_n(\lambda)-f_n(\lambda-h)}h$ and $\frac{f_n(\lambda+h)-f_n(\lambda)}h$ for small $h>0$; letting $n\to\infty$ and then $h\to0$ gives $f_n'(\lambda)\to f'(\lambda)$, that is, $\E v_p'(\lambda)\to v'(\lambda)$. With $p/n\to\gamma$, \eqref{eq:doubledescent:ridgefinite} converges to \eqref{eq:doubledescent:ridgelimit}. The Stieltjes form follows from $v(\lambda)=\lambda m_\gamma(-\lambda)$ by the chain rule: $v'(\lambda)=m_\gamma(-\lambda)-\lambda m_\gamma'(-\lambda)$, and $v-\lambda v'=\lambda^2m_\gamma'(-\lambda)$.

\emph{The ridgeless limit.} Differentiating $\gamma v^2+(1-\gamma+\lambda)v-\lambda=0$ in $\lambda$, and using $2\gamma v+1-\gamma+\lambda=\sqrt{(1-\gamma+\lambda)^2+4\gamma\lambda}$ from \eqref{eq:doubledescent:v}, gives
\[
  v'(\lambda)=\frac{1-v(\lambda)}{\sqrt{(1-\gamma+\lambda)^2+4\gamma\lambda}} .
\]
As $\lambda\to0$ the square root tends to $\abs{1-\gamma}>0$, and $v(\lambda)\to\frac{\abs{1-\gamma}-(1-\gamma)}{2\gamma}=(1-\frac1\gamma)_+$. For $\gamma<1$ this gives $v\to0$ and $v'\to\frac1{1-\gamma}$, so $\cR_\lambda(\gamma)\to\sigma^2\frac\gamma{1-\gamma}$; for $\gamma>1$ it gives $v\to1-\frac1\gamma$ and $v'\to\frac1{\gamma(\gamma-1)}$, so $\cR_\lambda(\gamma)\to r^2(1-\frac1\gamma)+\frac{\sigma^2}{\gamma-1}$. Both agree with \eqref{eq:doubledescent:limit}.
\end{proof}

The limit $v(0+)=(1-\frac1\gamma)_+$ is the mass of $\mu_\gamma$ at zero, the fraction of directions that $n$ examples cannot see, and the bias of the ridgeless estimator per unit of $r^2$. How much ridge is best? The answer is explicit at every sample size.

\begin{proposition}[Optimal ridge]\label[proposition]{prop:doubledescent:optimal}
Under \cref{ass:doubledescent:design}, with $r>0$ and $\sigma>0$:
\begin{enumerate}
\item for all $n$ and $p$, the risk $\E\norm{\hat w_\lambda-w_\star}^2$ is minimized over $\lambda>0$ at $\lambda^\star=\sigma^2p/(nr^2)$, where it equals $r^2\,\E v_p(\lambda^\star)$;
\item for fixed $p$, the minimal risk does not increase with $n$;
\item as $n\to\infty$ and $p/n\to\gamma$, the minimal risk converges to
\begin{equation}\label{eq:doubledescent:optrisk}
  \cR^\star(\gamma)=\frac{r^2}2\Bigl(\sqrt{b^2+4s}-b\Bigr),\qquad s=\frac{\sigma^2}{r^2},\quad b=\frac1\gamma-1+s,
\end{equation}
which increases continuously from $0$ to $r^2$ as $\gamma$ increases from $0$ to $\infty$; in particular it is finite at $\gamma=1$.
\end{enumerate}
\end{proposition}

\begin{proof}
\emph{Part (i).} By \cref{lem:doubledescent:ridgespectral} the risk is $\E\int g_\lambda\dd\mu_p$ with $g_\lambda(t)=\frac{r^2\lambda^2+ct}{(t+\lambda)^2}$ and $c=\sigma^2p/n$. For each $t\ge0$,
\[
  \frac{\partial g_\lambda(t)}{\partial\lambda}=\frac{2r^2\lambda(t+\lambda)-2(r^2\lambda^2+ct)}{(t+\lambda)^3}=\frac{2t\,(r^2\lambda-c)}{(t+\lambda)^3},
\]
so $\lambda\mapsto g_\lambda(t)$ decreases on $(0,\lambda^\star]$ and increases on $[\lambda^\star,\infty)$, with $\lambda^\star=c/r^2$; integrating against $\mu_p$ and taking expectations preserves both monotonicities. At $\lambda^\star$ we have $c=r^2\lambda^\star$, so $g_{\lambda^\star}(t)=\frac{r^2\lambda^\star(\lambda^\star+t)}{(t+\lambda^\star)^2}=\frac{r^2\lambda^\star}{t+\lambda^\star}$, whose integral against $\mu_p$ is $r^2v_p(\lambda^\star)$.

\emph{Part (ii).} The idea is that ridge with $\lambda^\star$ is the best estimator linear in the labels when the target is random and isotropic, and that this class grows with $n$. Let $w_\star$ be uniform on the sphere of radius $r$, independent of $(X,\xi)$; by \cref{lem:doubledescent:ridgespectral} this does not change the risk of a ridge estimator. For $A\in\R^{p\times n}$ depending on $X$ only, $\E w_\star w_\star^\top=\frac{r^2}pI_p$ and $\E\xi\xi^\top=\sigma^2I_n$ give
\[
  \E\bigl[\norm{Ay-w_\star}^2\bigm| X\bigr]=\frac{r^2}p\norm{AX-I_p}_F^2+\sigma^2\norm A_F^2 ,
\]
a strictly convex quadratic function of $A$, whose gradient $\frac{2r^2}p(AX-I_p)X^\top+2\sigma^2A$ vanishes exactly at $A=X^\top(XX^\top+\frac{\sigma^2p}{r^2}I_n)^{-1}=(X^\top X+\frac{\sigma^2p}{r^2}I_p)^{-1}X^\top$: ridge with $n\lambda=\sigma^2p/r^2$, that is, with $\lambda=\lambda^\star$. Now append one example. Among the matrices allowed with $n+1$ examples is the one that applies the optimal $n$-example matrix to the first $n$ labels and ignores the last, with the conditional risk of the optimal $n$-example estimator. The optimal $(n+1)$-example estimator does at least as well for every $X$, hence on average.

\emph{Part (iii).} By the proof of \cref{thm:doubledescent:ridge}, $\E v_p(s\gamma)\to v(s\gamma)$, and $\abs{v_p'}\le\frac1{4\lambda}$ gives $\abs{\E v_p(\lambda^\star)-\E v_p(s\gamma)}\le\abs{\lambda^\star-s\gamma}/(4\min(\lambda^\star,s\gamma))\to0$, so the minimal risk tends to $r^2v(s\gamma)$. Substituting $\lambda=s\gamma$ in $\gamma v^2+(1-\gamma+\lambda)v-\lambda=0$ and dividing by $\gamma$ gives $v^2+bv-s=0$, whose positive root gives \eqref{eq:doubledescent:optrisk}. The map $b\mapsto\frac12(\sqrt{b^2+4s}-b)=\frac{2s}{\sqrt{b^2+4s}+b}$ is decreasing, and $b$ decreases with $\gamma$, so $\cR^\star$ increases; as $\gamma\to0$, $b\to\infty$ and $\cR^\star\to0$, and as $\gamma\to\infty$, $b\to s-1$ and $\cR^\star\to\frac{r^2}2\bigl((s+1)-(s-1)\bigr)=r^2$.
\end{proof}

In the unnormalized form $\norm{y-Xw}^2+n\lambda\norm w^2$ the optimal penalty $n\lambda^\star=\sigma^2p/r^2$ is the noise variance divided by the variance of one coordinate of the target, whatever $n$ is; ridge with this penalty is the posterior mean for the prior $w_\star\sim\cN(0,\frac{r^2}pI_p)$ and Gaussian noise. Part (ii) is the sample-wise monotonicity theorem of \textcite{NakkiranEtAl2021}, proved there for isotropic Gaussian covariates; the argument above uses only second moments of the noise and of the target.

\section{Reading the solution}\label{sec:doubledescent:reading}

\Cref{fig:doubledescent:risk} compares \cref{thm:doubledescent:exact} and \cref{cor:doubledescent:weak} with simulations at $n=40$; every simulated point lies within about two standard errors of its curve.

\begin{figure}[t]
\begin{fullwidth}
\centering
\begin{tikzpicture}
\begin{groupplot}[group style={group size=2 by 1,horizontal sep=15mm},hplot,
  width=0.5\linewidth,height=0.37\linewidth,ymode=log,unbounded coords=jump,
  xmin=0,log ticks with fixed point,
  ytick={0.25,0.5,1,2,4,8},yticklabels={$0.25$,$0.5$,$1$,$2$,$4$,$8$}]
\nextgroupplot[xlabel={number of features $p$},ylabel={test error},xmax=200,ymin=0.2,ymax=12,
  title={(a) well specified}]
\addplot[guide,domain=0:200,samples=2,forget plot] {1.25};
\addplot[guide,domain=0:200,samples=2,forget plot] {0.25};
\addplot[guide,forget plot] coordinates {(40,0.2) (40,12)};
\addplot[theory] table[x=p,y=wellspec] {figures/data/01-doubledescent-risk-theory.dat};
\addplot[simulation] table[x=p,y=wellspec] {figures/data/01-doubledescent-risk-mc.dat};
\node[font=\scriptsize,text=muted,anchor=south east] at (axis cs:198,1.27) {null risk};
\node[font=\scriptsize,text=muted,anchor=south east] at (axis cs:198,0.255) {noise floor};
\nextgroupplot[xlabel={number of features $p$},xmax=100,ymin=0.4,ymax=12,
  title={(b) weak features, $d=100$}]
\addplot[guide,domain=0:100,samples=2,forget plot] {1};
\addplot[guide,forget plot] coordinates {(40,0.4) (40,12)};
\addplot[theory] table[x=p,y=weakdecay] {figures/data/01-doubledescent-risk-theory.dat};
\addlegendentry{decaying signal}
\addplot[theory3] table[x=p,y=weakeven] {figures/data/01-doubledescent-risk-theory.dat};
\addlegendentry{evenly spread}
\addplot[simulation,forget plot] table[x=p,y=weakdecay] {figures/data/01-doubledescent-risk-mc.dat};
\addplot[simulation2,forget plot] table[x=p,y=weakeven] {figures/data/01-doubledescent-risk-mc.dat};
\end{groupplot}
\end{tikzpicture}
\caption{Test error of minimum-norm least squares against the number $p$ of features, for $n=40$ examples; the dashed vertical line marks $p=n$. Curves: \cref{thm:doubledescent:exact} and \cref{cor:doubledescent:weak}. Points: averages over $4000$ training sets (seed $20261017$), each fit scored by its exact test error \eqref{eq:doubledescent:loss}. (a)~Well specified, $r^2=1$, $\sigma^2=0.25$; the second descent bottoms out at $p^\star=82$. (b)~Weak features without label noise, $d=100$, $r=1$, the learner using the first $p$ coordinates; the signal is spread evenly ($w_{\star,j}^2=1/d$, squares) or decays ($w_{\star,j}^2\propto j^{-0.8}$, circles). Points with $\abs{p-n}\le3$, where the simulated average has infinite mean or variance, are omitted.}
\label{fig:doubledescent:risk}
\end{fullwidth}
\end{figure}

\paragraph{The anatomy of the peak.} Below the threshold the risk is pure variance, $\sigma^2\,\E\tr\bigl((X^\top X)^{-1}\bigr)=\sigma^2\frac pn\,\E\int t^{-1}\mu_p(\dd t)$: the noise is amplified by the inverse eigenvalues of $\hat\Sigma$, and the small ones dominate. In the limit, letting $z\to0$ in \eqref{eq:doubledescent:mpquad} gives $\int t^{-1}\mu_\gamma(\dd t)=\frac1{1-\gamma}$ for $\gamma<1$, the factor in \eqref{eq:doubledescent:limit}. As $\gamma$ increases to $1$ the lower edge $t_-=(1-\sqrt\gamma)^2$ of the spectrum moves down to $0$; at $\gamma=1$ the density \eqref{eq:doubledescent:mp} behaves like $1/(\pi\sqrt t)$ near $0$, and the inverse moment diverges.\sidenote{In the language of \textcite{SpiglerEtAl2019}, $p=n$ is the isostatic point of a jamming transition, at which the degrees of freedom equal the constraints and the response of the fitted weights to the noise diverges.} At finite size the smallest singular value of $X$ is close to $\sqrt n-\sqrt p$ when $p<n$ \parencite{Silverstein1985} but only of order $n^{-1/2}$ when $p=n$ \parencite{Edelman1988}.\sidenote{\Textcite{Edelman1988} also showed that the condition number of a square Gaussian matrix, divided by $n$, converges in distribution: at the threshold it grows linearly with $n$.} Past the threshold the bias $r^2(1-n/p)$ grows, because the interpolator lives in an $n$-dimensional subspace and sees a shrinking fraction of the target, while the variance falls like $\sigma^2n/(p-n)$, because with more directions available the $n$ noise values are interpolated by a shorter weight vector, whose squared length the test error pays. The second descent is the variance falling faster than the bias grows.

\paragraph{Where the second descent ends.} Treat $p$ in \eqref{eq:doubledescent:exact} as a real variable. On $p>n+1$ the derivative of the test error is $r^2n/p^2-\sigma^2n/(p-n-1)^2$, so if $\sigma<r$, then the test error is smallest at
\begin{equation}\label{eq:doubledescent:pstar}
  p^\star=\frac{(n+1)\,r}{r-\sigma},\qquad \E\Loss(\hat w)=\sigma^2+r^2-\frac n{n+1}(r-\sigma)^2 ,
\end{equation}
and if $\sigma\ge r$, then it decreases to the null risk without reaching it. In \cref{fig:doubledescent:risk}(a), $p^\star=82$ and the minimal test error is $1.006$, against a null risk of $1.25$. In the proportional limit the optimum is $\gamma^\star=r/(r-\sigma)$, with test error $2\sigma r$, below the null risk by $(r-\sigma)^2$. The same formulas give a window of harm: minimum-norm least squares predicts worse than the zero predictor exactly when
\begin{equation}\label{eq:doubledescent:window}
  \frac{r^2}{r^2+\sigma^2}<\gamma<\frac{r^2}{r^2-\sigma^2}
\end{equation}
(the upper end is $+\infty$ when $\sigma\ge r$; \cref{ex:doubledescent:optimum}). The window contains the threshold and widens with the noise. Here, however, a learner with more features also faces a truth with more coordinates, and $p^\star$ is only a local minimum: the test error $\sigma^2\frac{n-1}{n-p-1}$ below the threshold is smallest at $p=1$, and \textcite{HastieEtAl2022} note that in this well-specified model the global minimum always lies below the threshold. The weak-features model compares learners on one fixed truth.

\paragraph{Weak features.} When the features are ordered by importance and the signal is spread over many of them, the curve shows every phase of \cref{obs:doubledescent:dd}:\sidenote{The figures of \textcite{BelkinEtAl2020} use the same $d=100$ and $n=40$, with label noise.} the decaying profile of \cref{fig:doubledescent:risk}(b) descends to $0.755$ at $p=10$, peaks at the threshold, and descends again to $0.596$ at $p=90$, below the bottom of the U. If the signal is concentrated on few features, then the U wins instead: with $w_{\star,j}^2\propto j^{-2}$ the best model has $p=22<n$, with test error $0.048$ against $0.191$ at best past the threshold, in line with the classical analysis of \textcite{BreimanFreedman1983}. When the signal is spread evenly, $w_{\star,j}^2=r^2/d$, the first descent disappears: for $p\le n-2$ the test error $(\sigma^2+r^2(1-p/d))\frac{n-1}{n-p-1}$ increases with $p$ as soon as $d\ge n$, so the best classical model uses no features at all, while for $n+2\le p\le d$ \cref{cor:doubledescent:weak} gives (\cref{ex:doubledescent:even})
\begin{equation}\label{eq:doubledescent:even}
  \E\Loss=\sigma^2+r^2-\frac{2nr^2}d+\frac n{p-n-1}\Bigl(\sigma^2+r^2\Bigl(1-\frac{n+1}d\Bigr)\Bigr),
\end{equation}
which decreases in $p$. The best interpolating model uses all $d$ features, and it beats the null risk, hence every classical model, if and only if $r^2/\sigma^2>d/(d-n-1)$, the condition of \textcite{BelkinEtAl2020}. Each added feature brings signal $r^2/d$; below the threshold it costs more than that in variance, above it it lowers the variance too.

\paragraph{More data can hurt.} At fixed $p$ and $n\le p-2$, the risk $r^2(1-n/p)+\sigma^2\frac n{p-n-1}$ increases with $n$ as soon as $(p-n-1)^2<\frac{\sigma^2}{r^2}p(p-1)$: within about $p\sigma/r$ below $n=p$, each new example makes the interpolator worse, the sample-wise double descent of \textcite{Nakkiran2019} (\cref{ex:doubledescent:samplewise}).

\paragraph{Regularization removes the peak.} In \cref{fig:doubledescent:ridge}, a fixed ridge parameter lowers the peak and leaves a bump, while the optimal $\lambda^\star=\sigma^2\gamma/r^2$ leaves the increasing curve \eqref{eq:doubledescent:optrisk}, whose value at the threshold, $\frac{r^2}2(\sqrt{s^2+4s}-s)$, is below $\sigma r$. In this model the peak is a property of the ridgeless estimator, not of the data.\sidenote{A fixed penalty is not enough for monotonicity: the curve for $\lambda=0.01$ in \cref{fig:doubledescent:ridge} still falls from $1.43$ at $\gamma=1$ to $0.99$ at $\gamma=2$ before rising again.}

\begin{figure}[t]
\centering
\begin{tikzpicture}
\begin{axis}[hplot,height=0.5\linewidth,xlabel={$\gamma=p/n$},ylabel={test error},xmin=0,xmax=4,ymin=0.2,ymax=2,
  unbounded coords=jump,legend pos=north east,legend columns=1]
\addplot[muted,line width=0.8pt] table[x=gamma,y=ridgeless] {figures/data/01-doubledescent-ridge-theory.dat};
\addlegendentry{$\lambda=0$}
\addplot[theory3] table[x=gamma,y=lam0p01] {figures/data/01-doubledescent-ridge-theory.dat};
\addlegendentry{$\lambda=0.01$}
\addplot[theory2] table[x=gamma,y=lam0p1] {figures/data/01-doubledescent-ridge-theory.dat};
\addlegendentry{$\lambda=0.1$}
\addplot[theory] table[x=gamma,y=optimal] {figures/data/01-doubledescent-ridge-theory.dat};
\addlegendentry{$\lambda=\lambda^\star$}
\addplot[simulation,forget plot] table[x=gamma,y=lam0p01] {figures/data/01-doubledescent-ridge-mc.dat};
\addplot[simulation2,forget plot] table[x=gamma,y=lam0p1] {figures/data/01-doubledescent-ridge-mc.dat};
\addplot[simulation,mark=triangle*,mark size=2pt,mark options={fill=sea,draw=white,line width=0.2pt},color=sea,forget plot] table[x=gamma,y=optimal] {figures/data/01-doubledescent-ridge-mc.dat};
\end{axis}
\end{tikzpicture}
\caption{Test error of ridge regression against $\gamma=p/n$, for $r^2=1$ and $\sigma^2=0.25$. Curves: the limits \eqref{eq:doubledescent:ridgelimit}, \eqref{eq:doubledescent:limit}, and \eqref{eq:doubledescent:optrisk}. Points: simulations with $n=200$ and $p=\gamma n$, averaged over $500$ training sets per point (seed $1017$), for $\lambda=0.01$ (circles), $\lambda=0.1$ (squares), and $\lambda=\lambda^\star$ (triangles); the largest deviation from the curves is $2.0$ standard errors.}
\label{fig:doubledescent:ridge}
\end{figure}

\section{What the model gets wrong}\label{sec:doubledescent:wrong}

\paragraph{What was thrown away.} With fixed features the model cannot say what a trained network gains over random features, and its $p$ counts directions chosen before the data are seen. Without anisotropy the spectrum plays no role, yet the spectrum decides how much of the target an interpolator sees and where the fitted noise goes (\cref{ex:doubledescent:aniso,ex:doubledescent:spiked}); real feature spectra decay, often as power laws (\cref{ch:scaling}). With a linear target and independent noise, misspecification enters only as white noise, whereas what a small network cannot represent is structured and correlated with the input. Without time, early stopping enters only through its resemblance to ridge \parencite{AdvaniEtAl2020}, and the epoch-wise curve not at all. With square loss, interpolation is exact fitting; for classification with the logistic loss it is separation, and gradient descent selects a maximum-margin interpolator (\cref{ch:margin}).

\paragraph{Where it breaks.} \emph{The axis.} The model puts the peak at a parameter count, $p=n$; trained networks put it elsewhere. \Textcite{SpiglerEtAl2019} found that the jamming threshold of fully connected networks grows about linearly with the number of examples for random data, but much more slowly on a binary version of MNIST. \Textcite{NakkiranEtAl2020} found that the peak follows the effective model complexity, which moves with training time, label noise, and augmentation, so that one architecture sits on either side of the threshold depending on how long it is trained. The count of $nK$ parameters for $K$ outputs \parencite{BelkinEtAl2019} fits least squares with $K$ outputs and random features with a trained last layer, but for networks trained end to end it is an observation \parencite{AdvaniEtAl2020}, not a theorem. The model's $p$ is at best an effective number of fitted directions, which it does not tell us how to compute.

\emph{Time.} The model has no epoch-wise double descent: for isotropic features, the test error of gradient flow from zero, computed from the Marchenko--Pastur law, has a single minimum in time at every value of $\gamma$ and $\sigma/r$ that we checked numerically. \Textcite{HeckelYilmaz2021} showed that epoch-wise double descent arises when parts of a model are learned at different speeds, so that bias--variance trade-offs are superposed, and removed it from a convolutional network and a ResNet by changing the step sizes of some layers; \textcite{PezeshkiEtAl2022} derived it from fast and slow features in a teacher--student model (\cref{ex:doubledescent:epochwise}).

\emph{The tail.} As $\gamma\to\infty$ the test error of the isotropic interpolator tends to the null risk: the variance vanishes, but so does the fraction $n/p$ of the target that the interpolator sees. Isotropic features never overfit benignly, whereas interpolating networks, which fit noisy labels exactly \parencite{ZhangEtAl2017}, and many interpolating kernel machines overfit in a tempered rather than catastrophic way \parencite{MallinarEtAl2022}. In linear regression, benign overfitting needs structure in the covariance.

\begin{theorem}[\textcite{BartlettEtAl2020}, Gaussian case]\label{thm:doubledescent:bllt}\index{benign overfitting}
Let $x\sim\cN(0,\Sigma)$ in $\R^p$ with $p>n$ and $\Sigma$ positive definite with eigenvalues $\lambda_1\ge\dots\ge\lambda_p$, and let $y=\ip{w_\star,x}+\eps$ with $\eps\sim\cN(0,\sigma^2)$ independent of $x$. Define the effective ranks
\[
  r_k=\frac{\sum_{i>k}\lambda_i}{\lambda_{k+1}},\qquad R_k=\frac{\bigl(\sum_{i>k}\lambda_i\bigr)^2}{\sum_{i>k}\lambda_i^2},
\]
and let $k^*$ be the smallest $k\ge0$ with $r_k\ge bn$. There are universal constants $b,c,c_1>1$ such that, for $\delta\in(0,1)$ with $\log(1/\delta)<n/c$, the excess risk $\cE=\norm{\Sigma^{1/2}(X^+y-w_\star)}^2$ satisfies the following. If $k^*\ge n/c_1$, then $\E\cE\ge\sigma^2/c$. Otherwise $\E\cE\ge\frac{\sigma^2}c\bigl(\frac{k^*}n+\frac n{R_{k^*}}\bigr)$, and with probability at least $1-\delta$,
\begin{multline*}
  \cE\le c\norm{w_\star}^2\norm\Sigma_{\op}\max\Bigl\{\sqrt{\frac{r_0}n},\frac{r_0}n,\sqrt{\frac{\log(1/\delta)}n}\Bigr\}\\
  +c\log(1/\delta)\,\sigma^2\Bigl(\frac{k^*}n+\frac n{R_{k^*}}\Bigr).
\end{multline*}
\end{theorem}

The theorem names what the isotropic model lacks.\sidenote{\Textcite{BartlettEtAl2020} prove it for features with independent sub-Gaussian coordinates in the eigenbasis of $\Sigma$ and noise that is sub-Gaussian given $x$; we quote the Gaussian case.} The variance needs a \emph{small head}, $k^*\ll n$, few directions of large variance that cost $\sigma^2k^*/n$, and a \emph{heavy, flat tail}, $R_{k^*}\gg n$, many directions of small, comparable variance that absorb the fitted noise at the cost $\sigma^2n/R_{k^*}$, the analogue of $\sigma^2n/(p-n-1)$. The bias needs $r_0=\tr\Sigma/\norm\Sigma_{\op}\ll n$, so that the tail does not drown the head. For $\Sigma=I_p$ with $p\ge bn$, $k^*=0$ and $R_0=p$, so the variance vanishes as $p/n\to\infty$; but $r_0=p$, and the bias $r^2(1-n/p)$ of \cref{thm:doubledescent:exact} does not vanish. \Cref{ex:doubledescent:spiked} builds a benign covariance by hand.

\emph{The counting.} \Textcite{CurthEtAl2023} found that in the non-deep experiments of \textcite{BelkinEtAl2019} the horizontal axis switches at the threshold from one way of adding capacity to another, and that against an effective number of parameters measured at test points the curves fold back into a U. Our model shows the mechanism in closed form. The prediction of $X^+y$ at a test point $x_0$ is $\ip{\hat s(x_0),y}$ with smoothing weights $\hat s(x_0)=(X^+)^\top x_0$, and with $p_{\mathrm{eff}}=n\,\E\norm{\hat s(x_0)}^2$ the variance in \cref{thm:doubledescent:exact} is $\sigma^2p_{\mathrm{eff}}/n$ on both sides of the threshold. But $p_{\mathrm{eff}}=np/(n-p-1)$ below the threshold and $n^2/(p-n-1)$ above it: past $p=n$, every added feature makes the estimator effectively simpler.

\emph{The cure.} Optimal ridge removes the peak in this model, at a positive $\lambda^\star$. For anisotropic data the optimal ridge parameter can be zero or even negative \parencite{KobakEtAl2020,TsiglerBartlett2023}: many low-variance directions already act as a ridge, and an interpolator is then optimal, not merely harmless. The isotropic model cannot represent this.

\paragraph{Fine structure.} Random nonlinear features were solved by \textcite{MeiMontanari2022}: in the proportional limit they behave like linear features plus independent noise, the ridgeless peak stays at $N=n$ features, and the best test error is reached as $N/n\to\infty$, with a ridge parameter that vanishes above a critical signal-to-noise ratio. For anisotropic covariances the ridge risk is given by deterministic equivalents that generalize \cref{thm:doubledescent:ridge} \parencite{HastieEtAl2022,AtanasovEtAl2026}, which power-law spectra turn into the scaling laws of \cref{ch:scaling}. Between benign and catastrophic overfitting lies tempered overfitting, a finite excess risk bounded away from zero, as for kernels with power-law spectra and, empirically, for interpolating networks \parencite{MallinarEtAl2022}; \textcite{LiangRakhlin2020} analyze kernel interpolation in high dimension. Gradient flow in the random-features model gives model-wise, sample-wise, and epoch-wise descents in one exact solution \parencite{BodinMacris2021}, and in toy models of superposition double descent separates memorizing data points from learning features (\cref{ch:superposition}; \cite{HenighanEtAl2023}). Open: a solvable model with feature learning that predicts where the peak of a trained network sits, and how high it is.

\section{Model card}\label{sec:doubledescent:card}

\modelcard{Test error falls, peaks where the model first fits the training set, and falls again, in model size, training time, and sample size.}
{Least squares on $p$ isotropic Gaussian features from $n$ examples, minimum-norm or ridge; weak features: $p$ of the $d$ features of a fixed truth.}
{Exact risk at every $(n,p)$ (\cref{thm:doubledescent:exact}); ridge risk from the Marchenko--Pastur law (\cref{thm:doubledescent:ridge}); optimal ridge $\lambda^\star=\sigma^2p/(nr^2)$ (\cref{prop:doubledescent:optimal}).}
{An infinite peak at $p=n$; the second descent ends at $p^\star=(n+1)r/(r-\sigma)$, near test error $2\sigma r$; a window of harm; more data hurting below the threshold; no peak under optimal ridge.}
{The axis (networks peak at an effective complexity); no epoch-wise descent; no benign overfitting with isotropic features; one peak only.}
{Random features; anisotropic deterministic equivalents; tempered overfitting; several time scales.}

The rotational invariance that made the risk computable returns in the kernel lower bound of \cref{ch:lazy}, in the single-index and committee models of \cref{ch:singleindex,ch:committee}, and in the in-context regression of \cref{ch:icl}, whose optimal one-step learner is compared there with ridge. The Marchenko--Pastur law is the first spectrum of the book; kernel spectra follow in \cref{ch:lazy}, and power-law spectra, which turn deterministic equivalents into scaling laws, in \cref{ch:scaling}. The minimum-norm interpolator is the first implicit bias of the book; a reparameterization changes the norm to $\ell_1$ in \cref{ch:diagonal}, and the logistic loss gives the maximum margin in \cref{ch:margin}. Epoch-wise double descent needs features learned at different speeds, the staged learning of \cref{ch:deeplinear}.

\begin{exercises}
\item \consolidation (Inverse Wishart matrices.) Let $X\in\R^{n\times p}$ have independent $\cN(0,1)$ entries. (a) For $n\ge p+2$, show that $\E[(X^\top X)^{-1}]=I_p/(n-p-1)$ with \cref{lem:doubledescent:rotation}, using $\norm A_F\le\tr A$ for positive semidefinite $A$. (b) Show that $\E[T^{-2}]=\frac1{(k-2)(k-4)}$ for $T\sim\chi^2_k$ with $k\ge5$, and $\E[T^{-2}]=+\infty$ for $k\le4$. (c) Deduce that $\tr\bigl((X^\top X)^{-1}\bigr)$ has finite variance if and only if $p\le n-4$, and that with Gaussian noise the simulated test error of $X^+y$ has infinite variance for $p\in\{n-3,n-2,n+2,n+3\}$.\label{ex:doubledescent:wishart}

\item \consolidation (Gradient descent finds the minimum norm.) Run gradient descent $w_{t+1}=w_t-\eta\nabla\Lhat(w_t)$ from $w_0$ with $0<\eta<2/\lmax(\hat\Sigma)$. (a) Show that $w_t\to X^+y+(I_p-P)w_0$, with $P$ the orthogonal projection onto the row space of $X$. (b) Let $p>n$ and $w_0\sim\cN(0,\alpha^2I_p)$, independent of the data. Show that the limit has risk $\E\norm{X^+y-w_\star}^2+\alpha^2(p-n)$.\label{ex:doubledescent:gd}

\item \consolidation (The overparameterized optimum.) (a) Derive \eqref{eq:doubledescent:pstar}, and show that the test error decreases on $p\ge n+2$ when $\sigma\ge r$. (b) Show that in the proportional limit the optimum is $\gamma^\star=r/(r-\sigma)$, with test error $2\sigma r$. (c) Prove that $\sigma^2+\cR_0(\gamma)>\sigma^2+r^2$ exactly in the window \eqref{eq:doubledescent:window}.\label{ex:doubledescent:optimum}

\item \consolidation (More data can hurt.) Fix $p$, $r$, and $\sigma>0$, and treat the risk $r^2(1-n/p)+\sigma^2n/(p-n-1)$ as a function of a real variable $n<p-1$. (a) Show that it increases with $n$ exactly when $(p-n-1)^2<\frac{\sigma^2}{r^2}p(p-1)$. (b) For $p=100$, $r=1$, and $\sigma=\frac12$, show that among $n\le98$ the risk is smallest at $n=49$, where it equals $0.755$, and that at least $135$ examples are needed to beat this value. (c) Why does this not contradict \cref{prop:doubledescent:optimal}(ii)?\label{ex:doubledescent:samplewise}

\item \consolidation (Moments of the Marchenko--Pastur law.) (a) Show that $\E\frac1p\tr\hat\Sigma=1$ and $\E\frac1p\tr\hat\Sigma^2=1+\frac{p+1}n$. (b) Substitute the expansion $m_\gamma(z)=-\sum_{k\ge0}z^{-k-1}\int t^k\mu_\gamma(\dd t)$, valid for large $\abs z$, in \eqref{eq:doubledescent:mpquad}, and show that the first two moments of $\mu_\gamma$ are $1$ and $1+\gamma$. (c) Show that $\int_{(0,\infty)}t^{-1}\mu_\gamma(\dd t)$ equals $\frac1{1-\gamma}$ for $\gamma<1$ and $\frac1{\gamma(\gamma-1)}$ for $\gamma>1$. \hint{in (c), let $\lambda\to0$ in $v'(\lambda)$.}\label{ex:doubledescent:moments}

\item \variation (Weak features with evenly spread signal.) Under \cref{ass:doubledescent:weak}, let $w_{\star,j}^2=r^2/d$ for every $j$ and $S=\{1,\dots,p\}$. (a) Derive \eqref{eq:doubledescent:even}, show that it decreases in $p$ on $n+2\le p\le d$, and show that its value at $p=d$ is below the null risk exactly when $r^2/\sigma^2>d/(d-n-1)$. (b) Let $p/n\to\gamma$ and $d/n\to\delta>1$. Show that on $1<\gamma\le\delta$ the test error tends to $\sigma^2+r^2-\frac{2r^2}\delta+\frac{\sigma^2+r^2(1-1/\delta)}{\gamma-1}$. (c) Let the learner use ridge on its $p$ features with the best $\lambda$ for that $p$. Show that the test error then does not increase with $p$, at every $n$. \hint{adapt \cref{prop:doubledescent:optimal}(ii), with a random sign pattern for $w_\star$; a learner with $p+1$ features may ignore one.}\label{ex:doubledescent:even}

\item \variation (Anisotropic design.) Let the rows of $X$ and the test inputs be independent $\cN(0,\Sigma)$ with $\Sigma$ positive definite, so that $\Loss(w)=\sigma^2+\norm{\Sigma^{1/2}(w-w_\star)}^2$. (a) Show that for $p\le n-2$ the test error of least squares is $\sigma^2\frac{n-1}{n-p-1}$, whatever $\Sigma$ is. (b) For $p\ge n+2$, show that $\Sigma^{-1/2}(X\Sigma^{-1/2})^+y$ has the risk of \cref{thm:doubledescent:exact} with $r=\norm{\Sigma^{1/2}w_\star}$, and check numerically that $X^+y$ does not.\label{ex:doubledescent:aniso}

\item \variation\hard (Benign overfitting with a spiked covariance.) Let $x=(z,\sqrt\epsilon\,u)$ with independent standard Gaussian $z\in\R^k$ and $u\in\R^q$, let $y=\ip{w_\star,z}+\eps$ with $w_\star\in\R^k$ and noise variance $\sigma^2$, and let $\hat w=(\hat w_z,\hat w_u)$ be the minimum-norm interpolator of $n<q$ examples, with data matrices $Z$ and $U$. (a) Show that $\hat w_z=Z^\top(ZZ^\top+\epsilon UU^\top)^{-1}y$ and that the excess risk is $\norm{\hat w_z-w_\star}^2+\epsilon\norm{\hat w_u}^2$. (b) Suppose for a moment that $UU^\top=qI_n$. Show that $\hat w_z$ is then ridge regression on the head with $n\lambda=\epsilon q$, and that $\epsilon\norm{\hat w_u}^2\le\norm y^2/q$; this is the implicit ridge of \textcite{KobakEtAl2020}. (c) Check numerically, with $k=5$, $\norm{w_\star}=1$, $\sigma^2=\frac14$, $n\in\{100,200,400\}$, $q=n^2/10$, and $\epsilon q/n=0.5/\sqrt n$, that the excess risk is about $0.04$, $0.02$, and $0.01$. (d) Compute the effective ranks of \cref{thm:doubledescent:bllt} for this covariance and check its conditions when $k$ is fixed, $q/n\to\infty$, and $\epsilon q/n\to0$. \hint{in (b), $Z^\top(ZZ^\top+cI_n)^{-1}=(Z^\top Z+cI_k)^{-1}Z^\top$.}\label{ex:doubledescent:spiked}

\item \simulation (Random ReLU features.) Take $d=100$, $n=300$, a unit vector $w_\star$, $x\sim\cN(0,I_d)$, and $y=\ip{w_\star,x}+\eps$ with $\eps\sim\cN(0,0.1)$. For $N$ from $30$ to $4000$, draw $v_1,\dots,v_N\sim\cN(0,I_d/d)$, fit minimum-norm least squares on the features $\relu(\ip{v_j,x})$, and average the test error on $2000$ fresh points over $20$ draws. (a) Observe a minimum of about $0.8$ for $N$ between $100$ and $150$, a peak above $7$ near $N=n$, and a second descent to about $0.19$ at $N=4000$. (b) Write $\relu(g)=\mu_0+\mu_1g+\mu_\star\zeta(g)$ for $g\sim\cN(0,1)$, with $\zeta$ of unit variance and uncorrelated with $1$ and $g$; show that $\mu_0=1/\sqrt{2\pi}$, $\mu_1=\frac12$, and $\mu_\star^2=\frac14-\frac1{2\pi}$. (c) Explain why the peak sits at $N=n$ rather than $N=d$, using the picture of \textcite{MeiMontanari2022}: linear features plus independent noise of variance $\mu_\star^2$.\label{ex:doubledescent:relu}

\item \simulation (Epoch-wise double descent from two time scales.) Take $n=100$ examples, $40$ fast features of variance $1$ and $20$ slow features of variance $\epsilon=0.01$, a target $w_\star$ that is constant on each group, with $\norm{\Sigma^{1/2}w_\star}^2=1$ on each group, and noise variance $\sigma^2=0.3$. Run gradient descent from zero with step $\eta=0.05$ for $2\times10^5$ steps, recording the test error $\sigma^2+\norm{\Sigma^{1/2}(w_t-w_\star)}^2$ and averaging over $400$ draws. (a) Observe a first minimum near step $25$ (test error about $1.69$), a bump near step $200$ (about $1.89$), a second minimum near step $1.3\times10^4$ (about $0.72$), and a slow rise toward the least-squares value $\sigma^2\frac{n-1}{n-p-1}\approx0.76$ of \cref{ex:doubledescent:aniso}. (b) Explain each phase with the time scales $1/\eta$ and $1/(\eta\epsilon)$, treating the unlearned slow signal as noise, as in the weak-features model. (c) Give the slow features the step size $\eta/\epsilon$, a preconditioning in the spirit of \textcite{HeckelYilmaz2021}, and observe a single minimum, about $0.63$ near step $100$.\label{ex:doubledescent:epochwise}

\item \research (A solvable model of epoch-wise double descent.) In \cref{ex:doubledescent:epochwise} the time scales are put in by hand, as in the linear models of \textcite{HeckelYilmaz2021} and \textcite{PezeshkiEtAl2022} and the random-feature model of \textcite{BodinMacris2021}. Open: a model in which training creates the time scales, for instance a two-layer network that learns features in stages as in \cref{ch:deeplinear}, and which predicts when the epoch-wise peak occurs and how its height depends on the label noise.\label{ex:doubledescent:researchepoch}

\item \research (Where is the peak of a network?) \Textcite{NakkiranEtAl2020} place the peak of trained networks where the effective model complexity equals the size of the training set. Open: compute the effective model complexity of a two-layer network trained by gradient descent from its width, input dimension, data distribution, and training time, in some solvable limit.\label{ex:doubledescent:researchpeak}
\end{exercises}

\notesandrefs
\Cref{thm:doubledescent:exact} is elementary, and its two halves were found decades apart. Below the threshold, the Gaussian weak-features model is the one that \textcite{BreimanFreedman1983} used to ask how many variables to enter in a regression, and they found the best number to be a small fraction of $n$. \Textcite{BelkinEtAl2020} took the model past the threshold and proved the formula and the weak-features corollary, while \textcite{HastieEtAl2022} and \textcite{BartlettEtAl2020}, also first posted in 2019, analyzed interpolation with general covariances. The inverse Wishart mean behind \cref{lem:doubledescent:wishart} is classical (\cref{thm:toolkit:inv-wishart}).

\Textcite{DobribanWager2018} give the proportional asymptotics of ridge for general covariances; we followed the isotropic form of \textcite{HastieEtAl2022}. \Textcite{KroghHertz1992} found the optimal weight decay, the $n\lambda^\star$ of \cref{prop:doubledescent:optimal}, and \textcite{AdvaniEtAl2020} showed that optimal early stopping is nearly as good as optimal ridge, with neither showing a peak.

\paragraph{Translating the literature.} \Textcite{HastieEtAl2022} use our $\gamma$, $r$, $\sigma$, and ridge normalization; they write $\beta$ for $w_\star$, and their risk $R_X(\hat\beta;\beta)$ is our risk conditional on $X$, without the noise floor. \Textcite{DobribanWager2018} write $w$ for $w_\star$ and $\alpha^2$ for $r^2$, set the noise variance to $1$, include the noise in their risk, and write $v$ for the companion Stieltjes transform, of the spectrum of $\frac1nXX^\top$, which is not our $v$; their $\lambda^*=\gamma/\alpha^2$ is our $\lambda^\star$ at $\sigma=1$. \Textcite{BelkinEtAl2020} write $D$, $T$, and $\beta$ for our $d$, $S$, and $w_\star$, and their risk includes the noise. \Textcite{NakkiranEtAl2021} penalize the unnormalized squared loss, so their $\lambda$ and $d$ are our $n\lambda$ and $p$, and their $\lambda_{\mathrm{opt}}=d\sigma^2/\norm{\beta^*}^2$ is our $n\lambda^\star$. \Textcite{MeiMontanari2022} write $N$ for the number of features, $\psi_1=N/d$, and $\psi_2=n/d$, so that our $\gamma$ is their $\psi_1/\psi_2$.

\printchapterbib
\end{refsection}

% Chapter 2 of The Hydrogen Atoms of Deep Learning: exact dynamics of deep linear networks.
\newcommand{\dlS}{\Sigma_{yx}}
\newcommand{\dlE}{\mathcal{E}}
\newcommand{\dlhalf}{t_{1/2}}
\newcommand{\dlLzero}{\Loss_0}

\casechapter{Learning in Stages}{Exact dynamics of deep linear networks}{Deep linear networks, whitened inputs}{ch:deeplinear}
\begin{refsection}[chapters/02-deeplinear.bib]

\begin{chaptersummary}
Networks trained from small initial weights often learn in stages: the loss sits on a plateau, drops quickly, sits on a lower plateau, and drops again, and what is learned first is the coarsest structure of the task. A deep linear network trained by gradient flow on whitened inputs reproduces this. We prove that the differences $W_{\ell+1}^\top W_{\ell+1}-W_\ell W_\ell^\top$ are conserved, that balanced layers make the end-to-end map follow a closed equation, and that from an aligned initialization the network learns the singular modes of the input--output correlation matrix one at a time, each along a logistic curve whose midpoint is at $\log(s/a_0)/(2s)$ for a mode of strength $s$ and initial strength $a_0$; depth $L\ge3$ stretches the plateau to order $a_0^{-(1-2/L)}/((L-2)s)$. The model cannot learn what has no linear correlation with the input, it orders the stages of nonlinear networks by the wrong statistic, and its implicit bias toward low rank is not described by a norm.
\end{chaptersummary}

\section{The phenomenon: plateaus, drops, and coarse before fine}\label{sec:deeplinear:phenomenon}

Training curves of neural networks are often not smooth. A network started from small weights may show a long stretch in which the loss barely moves, then a rapid drop, then another stretch of near-stagnation, and so on: a \term{loss staircase}\index{loss staircase}. Plateaus were a familiar obstacle in training small sigmoidal networks by backpropagation, and the committee machine of \cref{ch:committee} is the classical model of one of them. They return whenever the initialization is small.

The steps are not arbitrary. A small network introduced by David Rumelhart learns to answer questions such as ``a canary can \dots'' about a handful of living things (birds, fish, trees, and flowers), and \textcite{RogersMcClelland2004} made it the centerpiece of a theory of semantic cognition. Its hidden representations \term{differentiate progressively}\index{progressive differentiation}: early in training all items look alike; then plants separate from animals; later birds separate from fish and trees from flowers; and the distinctions between individual items come last. Developmental psychologists report a similar coarse-to-fine order in children's concepts, and \textcite{SaxeEtAl2019} set out to explain why a network trained by gradient descent should follow it.\sidenote{They also treat illusory correlations during learning, typicality, category coherence, and inductive generalization, and they contrast the learning trajectories of deep networks with those of shallow ones.}

The same pattern appears at a larger scale and in other guises.
\begin{itemize}
\item \Textcite{KalimerisEtAl2019} trained convolutional networks on image classification and measured, with a conditional mutual information, how much of a network's predictions a linear classifier explains. In the first epochs nearly all of the gain is explained by a linear classifier, and later epochs add functions of increasing complexity while retaining what the linear classifier predicted.
\item \Textcite{AroraEtAl2019} and \textcite{GissinEtAl2020} trained deep factorizations $W=W_L\cdots W_1$ of a matrix from small initialization and found that the singular values of $W$ are learned one at a time, more sharply separated as the depth grows: \term{incremental learning}\index{incremental learning}.
\item \Textcite{SimonEtAl2023} found, in a linearized model and in experiments, that self-supervised methods (Barlow Twins, SimCLR, VICReg) trained on ResNets learn their embeddings in a sequence of discrete, well-separated steps, one eigendirection of the embedding at a time.
\end{itemize}

\begin{observation}[\textcite{SaxeEtAl2014,SaxeEtAl2019}; \textcite{AroraEtAl2019}; \textcite{SimonEtAl2023}]\label[observation]{obs:deeplinear:stages}
Networks trained by gradient descent from small initialization often learn in stages. The training loss alternates between plateaus and rapid drops; each drop adds one direction to the learned map (a singular direction of the input--output map, or an eigendirection of a learned embedding); and directions that carry more of the input--output correlation are learned first. The effect is sharper for deeper networks and smaller initializations and fades for large initializations.
\end{observation}

The status of this observation is mixed. Stagewise learning is robust in linear and nearly linear settings (matrix factorization, linearized self-supervised learning) and visible in some nonlinear ones, but large-scale training with standard initialization and adaptive optimizers shows smooth loss curves in most cases, and where stages do appear in nonlinear networks their order need not follow linear statistics (\cref{sec:deeplinear:wrong}). The order ``coarse before fine'' is itself a hypothesis about what ``coarse'' means. The model of this chapter makes it precise: coarse means a large singular value of the input--output correlation matrix.

The question is therefore: \emph{why should gradient descent learn in stages, which stages come first, and how long does each plateau last?}

\section{The model: a deep linear network on whitened inputs}\label{sec:deeplinear:model}

\paragraph{Data and network.} Inputs $x\in\R^d$ and outputs $y\in\R^{d_y}$ have a joint law with finite second moments. Write
\[
  \Sigma_{xx}=\E[xx^\top]\in\R^{d\times d},\qquad \dlS=\E[yx^\top]\in\R^{d_y\times d}.
\]
A \term{deep linear network}\index{deep linear network} of depth $L\ge1$ is the map $f(x;\theta)=W_L\cdots W_1x$ with $W_\ell\in\R^{d_\ell\times d_{\ell-1}}$, $d_0=d$ and $d_L=d_y$; the hidden widths $d_1,\dots,d_{L-1}$ play the role of $m$. The parameters are $\theta=(W_1,\dots,W_L)$, and the \term{end-to-end matrix}\index{end-to-end matrix} is $W=W_L\cdots W_1$. The network computes only linear maps, so its expressive power is that of the matrix $W$ (of rank at most $\min_\ell d_\ell$). What depth changes is the geometry of the parameterization, and with it the dynamics of training.

\begin{assumption}[Whitened inputs]\label[assumption]{ass:deeplinear:white}
The inputs are whitened: $\Sigma_{xx}=I_d$.
\end{assumption}

\paragraph{The loss.} We train on the square loss, whose population version is $\Loss(\theta)=\frac12\E\norm{y-W_L\cdots W_1x}^2$. It depends on $\theta$ only through $W$, and only through two moments of the data.

\begin{lemma}[Reduction to a matrix problem]\label[lemma]{lem:deeplinear:reduction}
Under \cref{ass:deeplinear:white}, $\Loss(\theta)=\dlE(W)$ with
\[
  \dlE(W)=\frac12\norm{\dlS-W}_F^2+\dlLzero,\qquad \dlLzero=\frac12\bigl(\E\norm{y}^2-\norm{\dlS}_F^2\bigr)\ge0 .
\]
\end{lemma}

\begin{proof}
Expand the square and take expectations term by term, using $\E\ip{y,Wx}=\tr(W\E[xy^\top])=\ip{\dlS,W}_F$ and $\E\norm{Wx}^2=\tr(W\Sigma_{xx}W^\top)=\norm{W}_F^2$:
\[
  \Loss(\theta)=\frac12\E\norm{y}^2-\ip{\dlS,W}_F+\frac12\norm W_F^2=\frac12\norm{\dlS-W}_F^2+\dlLzero .
\]
The constant is the minimum of $\dlE$ over all matrices, which is the error of the best linear predictor, so it is nonnegative.
\end{proof}

The same computation applies to the training loss on $n$ samples, with $\Sigma_{xx}$ and $\dlS$ replaced by their empirical versions; whitening then means $\frac1n\sum_ix_ix_i^\top=I_d$, which requires $n\ge d$. Everything below holds verbatim for the training loss of whitened data.\sidenote{Whitening a dataset by $x\mapsto\Sigma_{xx}^{-1/2}x$ is always possible when $\Sigma_{xx}$ is invertible, but it changes the first layer to $W_1\Sigma_{xx}^{1/2}$, and gradient flow in the new variables is not gradient flow in the old ones; \cref{ex:deeplinear:nonwhite} treats the case in which the old flow still decouples.}

\paragraph{Singular modes.} Let $q=\min(d,d_y)$ and write a singular value decomposition
\begin{equation}\label{eq:deeplinear:svd}
  \dlS=USV^\top=\sum_{i=1}^qs_iu_iv_i^\top,\qquad s_1\ge s_2\ge\dots\ge s_q\ge0,
\end{equation}
with $U=[u_1,\dots,u_q]\in\R^{d_y\times q}$ and $V=[v_1,\dots,v_q]\in\R^{d\times q}$ having orthonormal columns and $S=\diag(s_1,\dots,s_q)$. We call $(s_i,u_i,v_i)$ the \term{singular mode}\index{singular mode} $i$ of the task: the input direction $v_i$ predicts the output pattern $u_i$ with strength $s_i$. In this chapter the index $i$ labels modes (Saxe and coauthors use $\alpha$, which in this book is an initialization scale), and $q$ is the number of modes.

\paragraph{Training.} We train by gradient flow, $\dot W_\ell=-\nabla_{W_\ell}\Loss(\theta)$ for $\ell=1,\dots,L$; gradient descent with step $\eta$ follows it with $t=\eta\times(\text{number of steps})$ when $\eta$ is small. To write the gradient, put
\[
  A_\ell=W_L\cdots W_{\ell+1},\qquad C_\ell=W_{\ell-1}\cdots W_1,
\]
with $A_L=I_{d_y}$ and $C_1=I_d$, so that $W=A_\ell W_\ell C_\ell$ for each $\ell$, and write $G=\nabla\dlE(W)$; under \cref{ass:deeplinear:white}, $G=W-\dlS$.

\begin{lemma}[Layer gradients]\label[lemma]{lem:deeplinear:grad}
For every differentiable $\dlE$ and $\Loss(\theta)=\dlE(W_L\cdots W_1)$, $\nabla_{W_\ell}\Loss=A_\ell^\top GC_\ell^\top$.
\end{lemma}

\begin{proof}
Perturbing $W_\ell$ by $H$ perturbs $W$ by $A_\ell HC_\ell$, so the derivative of $\Loss$ in the direction $H$ is $\ip{G,A_\ell HC_\ell}_F=\tr(G^\top A_\ell HC_\ell)=\ip{A_\ell^\top GC_\ell^\top,H}_F$.
\end{proof}

\paragraph{What was thrown away, and why.} The \emph{nonlinearity}: without it the network can represent only linear maps, and the loss is a polynomial in the weights; yet the training dynamics remain nonlinear, because the gradient for one layer is multiplied by all the others. This is the point of the model: it separates the nonlinearity of \emph{learning} from the nonlinearity of the \emph{function}. \emph{Whitening} makes the loss see the data only through $\dlS$, so that the singular modes of one matrix organize everything. \emph{Population loss and gradient flow} remove sampling noise and step-size effects. Finally the solution in \cref{sec:deeplinear:modes} needs an \emph{aligned initialization}, which no practitioner uses; \cref{sec:deeplinear:random} explains why small random initializations end up nearly aligned.

\begin{dictionary}
end-to-end matrix $W=W_L\cdots W_1$ & the input--output map of a trained network, linearized\\
singular mode $(s_i,u_i,v_i)$ of $\dlS$ & a distinction in the data, such as plant versus animal: the input direction $v_i$ that signals it and the output pattern $u_i$ it predicts\\
singular value $s_i$ & the statistical strength of a distinction, which sets how fast it is learned\\
mode strength $a_i(t)=u_i^\top W(t)v_i$ & how much of the distinction the network has learned at time $t$\\
hidden directions $R_\ell$ & the internal representation of a distinction in layer $\ell$\\
initial strength $a_i(0)$, or the per-layer scale $\alpha$ & the initialization scale; small values put the network in the regime of stages\\
plateau, then logistic rise of $a_i$ & a plateau of the training loss, then a drop\\
balancedness $W_{\ell+1}^\top W_{\ell+1}=W_\ell W_\ell^\top$ & layers of equal size, automatically maintained by gradient training\\
\end{dictionary}

\paragraph{What one had to notice.} That a network computing a linear map can have nonlinear training dynamics, so that the simplest model of a plateau is a product of matrices rather than a nonlinearity; that the square loss with whitened inputs sees only $\dlS$, whose singular value decomposition supplies natural coordinates in which the problem might decouple; and that the scalar version of the problem, $\dot b=c(s-bc)$ and $\dot c=b(s-bc)$ for a product $bc$, already has a saddle at the origin and a plateau before escape.

\begin{figure}[t]
\centering
\begin{tikzpicture}[x=0.86cm,y=0.95cm,
  nd/.style={circle,draw=sea,fill=seafoam,line width=0.5pt,minimum size=3.2mm,inner sep=0pt},
  lf/.style={font=\scriptsize,text=ink}]
  \node[nd] (r) at (3.5,3) {};
  \node[nd] (a) at (1.5,2) {}; \node[nd] (b) at (5.5,2) {};
  \foreach \i/\x in {0/0.5,1/2.5,2/4.5,3/6.5}{\node[nd] (c\i) at (\x,1) {};}
  \foreach \i in {0,...,7}{\node[nd] (l\i) at (\i,0) {};}
  \draw[sea] (r)--(a) (r)--(b) (a)--(c0) (a)--(c1) (b)--(c2) (b)--(c3);
  \foreach \i in {0,...,3}{\pgfmathtruncatemacro{\j}{2*\i}\pgfmathtruncatemacro{\k}{2*\i+1}\draw[sea] (c\i)--(l\j) (c\i)--(l\k);}
  \foreach \i/\n in {0/canary,1/robin,2/salmon,3/sunfish,4/oak,5/pine,6/rose,7/daisy}{\node[lf,anchor=north,rotate=0] at (\i,-0.22) {\n};}
  \node[lf,anchor=west] at (8.0,3) {level $0$: $s^2=15/8$, one mode};
  \node[lf,anchor=west] at (8.0,2) {level $1$: $s^2=7/8$, one mode};
  \node[lf,anchor=west] at (8.0,1) {level $2$: $s^2=3/8$, two modes};
  \node[lf,anchor=west] at (8.0,0) {level $3$: $s^2=1/8$, four modes};
\end{tikzpicture}
\caption{A hierarchical dataset in the style of \textcite{RogersMcClelland2004} and \textcite{SaxeEtAl2019}. Each of the $P=8$ items is an input $x=\sqrt P\,e_p$, so that $\Sigma_{xx}=I_8$, and its output $y\in\{0,1\}^{15}$ has one coordinate per node of the tree (``is an animal'', ``is a bird'', \dots), equal to $1$ for the nodes on the path from the root to the item (the leaf included). The labels on the right group the modes into levels, each level learned in one stage. The input directions $v_i$ of the singular modes of $\dlS$ can be taken to be the Haar functions of the tree: the constant function, then one function per internal node that is $+1$ on its left subtree, $-1$ on its right, and $0$ elsewhere. A mode at a lower level is stronger (\cref{ex:deeplinear:tree}).}\label{fig:deeplinear:tree}
\end{figure}

\section{The solution: conservation and balance}\label{sec:deeplinear:balance}

The first results hold for every loss of the form $\dlE(W_L\cdots W_1)$, whitened or not. They come from a symmetry: for an invertible $d_\ell\times d_\ell$ matrix $M$, replacing $(W_{\ell+1},W_\ell)$ by $(W_{\ell+1}M^{-1},MW_\ell)$ leaves $W$ unchanged. By Noether's principle a continuous symmetry of a gradient flow comes with a conserved quantity, and here it is a matrix.\sidenote{For a gradient flow the mechanism is simple: a loss that is constant along the orbits $M\mapsto(W_{\ell+1}M^{-1},MW_\ell)$ has its gradient orthogonal to the orbits, and orthogonality to the tangent directions $(-W_{\ell+1}X,XW_\ell)$, for all $d_\ell\times d_\ell$ matrices $X$, is exactly the equality of the two products in the proof below.}

\begin{theorem}[Conservation law; \textcite{AroraCohenHazan2018}, \textcite{DuHuLee2018}]\label{thm:deeplinear:conservation}
Let $\dlE\colon\R^{d_y\times d}\to\R$ be continuously differentiable and let $\theta(t)$ solve the gradient flow of $\Loss(\theta)=\dlE(W_L\cdots W_1)$ on an interval. Then for each $\ell\in\{1,\dots,L-1\}$ the $d_\ell\times d_\ell$ matrix
\[
  D_\ell(t)=W_{\ell+1}(t)^\top W_{\ell+1}(t)-W_\ell(t)W_\ell(t)^\top
\]
is constant in $t$. In particular $\norm{W_{\ell+1}}_F^2-\norm{W_\ell}_F^2=\tr D_\ell$ is constant.
\end{theorem}

\begin{proof}
The plan is to compute the two halves of $\dot D_\ell$ and to see that they are equal. By \cref{lem:deeplinear:grad}, $\dot W_\ell=-A_\ell^\top GC_\ell^\top$. Two identities between the partial products follow from their definitions: $W_\ell C_\ell=C_{\ell+1}$ and $A_{\ell+1}W_{\ell+1}=A_\ell$. Using the first,
\[
  W_\ell\dot W_\ell^\top=-W_\ell C_\ell G^\top A_\ell=-C_{\ell+1}G^\top A_\ell ,
\]
and using the second,
\[
  \dot W_{\ell+1}^\top W_{\ell+1}=-C_{\ell+1}G^\top A_{\ell+1}W_{\ell+1}=-C_{\ell+1}G^\top A_\ell .
\]
The two products are equal; call them $N$. Then $\frac{\dd}{\dd t}(W_\ell W_\ell^\top)=N+N^\top=\frac{\dd}{\dd t}(W_{\ell+1}^\top W_{\ell+1})$, so $\dot D_\ell=0$. Taking the trace gives the last statement.
\end{proof}

The scalar case $L=2$, $d=d_1=d_y=1$ says that $b^2-c^2$ is conserved for the product $bc$: the flow moves along hyperbolas, the picture behind the catapult of \cref{ch:stability} and the diagonal networks of \cref{ch:diagonal}. The same argument applies neuron by neuron to ReLU networks (\cref{ex:deeplinear:relu}).

We call $\theta$ \term{balanced}\index{balanced initialization} if $D_\ell=0$ for every $\ell$. By \cref{thm:deeplinear:conservation}, a flow that starts balanced stays balanced; a flow that starts from small weights stays nearly balanced, since $\norm{D_\ell}_F$ is at most of the order of the squared initialization scale while the weights themselves grow to order one. For balanced weights the layers have the same nonzero singular values, and the end-to-end matrix obeys an equation of its own.

\begin{lemma}[Powers of balanced products]\label[lemma]{lem:deeplinear:powers}
Let $\theta$ be balanced. Then for $1\le\ell\le L$,
\[
  A_\ell A_\ell^\top=(WW^\top)^{(L-\ell)/L},\qquad C_\ell^\top C_\ell=(W^\top W)^{(\ell-1)/L},
\]
where powers of positive semidefinite matrices are taken in their eigenbasis.
\end{lemma}

\begin{proof}
The plan is an induction that carries a power of a Gram matrix through the product. Let $P_j=W_L\cdots W_j$ for $1\le j\le L$, so that $P_1=W$ and $P_{\ell+1}=A_\ell$ (with $P_{L+1}=I$). We show, by downward induction on $j$, that for every integer $k\ge1$
\begin{equation}\label{eq:deeplinear:claim}
  P_j\,(W_j^\top W_j)^{k-1}P_j^\top=(W_LW_L^\top)^{L-j+k}.
\end{equation}
For $j=L$ the left side is $W_L(W_L^\top W_L)^{k-1}W_L^\top=(W_LW_L^\top)^k$, by regrouping the factors. For $j<L$, write $P_j=P_{j+1}W_j$; regroup $W_j(W_j^\top W_j)^{k-1}W_j^\top=(W_jW_j^\top)^k$; replace $W_jW_j^\top$ by $W_{j+1}^\top W_{j+1}$, which is balancedness; and apply \eqref{eq:deeplinear:claim} for $j+1$ with $k+1$ in place of $k$:
\[
  P_j(W_j^\top W_j)^{k-1}P_j^\top=P_{j+1}(W_{j+1}^\top W_{j+1})^{k}P_{j+1}^\top=(W_LW_L^\top)^{L-j+k}.
\]
With $k=1$, \eqref{eq:deeplinear:claim} gives $P_jP_j^\top=(W_LW_L^\top)^{L-j+1}$. For $j=1$ this is $WW^\top=(W_LW_L^\top)^L$, and since $W_LW_L^\top$ is positive semidefinite it is the unique positive semidefinite $L$-th root, $W_LW_L^\top=(WW^\top)^{1/L}$. For $j=\ell+1$ it gives $A_\ell A_\ell^\top=(W_LW_L^\top)^{L-\ell}=(WW^\top)^{(L-\ell)/L}$. The statement for $C_\ell$ is the same argument applied to the transposed network $W^\top=W_1^\top\cdots W_L^\top$, which is balanced as well.
\end{proof}

\begin{theorem}[End-to-end dynamics; \textcite{AroraCohenHazan2018}]\label{thm:deeplinear:e2e}
Let $\dlE$ be continuously differentiable, and let $\theta(t)$ solve the gradient flow of $\dlE(W_L\cdots W_1)$ from a balanced initialization. Then $W(t)=W_L(t)\cdots W_1(t)$ satisfies
\begin{equation}\label{eq:deeplinear:e2e}
  \dot W=-\sum_{j=1}^L(WW^\top)^{(L-j)/L}\,\nabla\dlE(W)\,(W^\top W)^{(j-1)/L} .
\end{equation}
\end{theorem}

\begin{proof}
By the product rule, $\dot W=\sum_{j=1}^LA_j\dot W_jC_j$. By \cref{lem:deeplinear:grad}, $\dot W_j=-A_j^\top GC_j^\top$, so $\dot W=-\sum_jA_jA_j^\top G\,C_j^\top C_j$. By \cref{thm:deeplinear:conservation} the flow stays balanced, and \cref{lem:deeplinear:powers} turns the two Gram matrices into the stated powers.
\end{proof}

For $L=1$ this is plain gradient flow, $\dot W=-\nabla\dlE(W)$. For $L\ge2$ the gradient is multiplied on both sides by powers of $WW^\top$ and $W^\top W$, which vanish when $W$ is small and are of order one when $W$ is: depth acts as a preconditioner that slows the directions in which $W$ is small. Near $W=0$ the right side of \eqref{eq:deeplinear:e2e} is of order $\norm W^{2-2/L}$, so $W=0$ is a stationary point, and $\theta=0$ is a saddle of $\Loss$ for every $L\ge2$ whenever $\dlS\ne0$ (with a vanishing Hessian when $L\ge3$). \Textcite{AroraCohenHazan2018} prove further that, for $L\ge2$, the effect of this preconditioning cannot be obtained by adding a regularizer to $\dlE$ and following its gradient.

\section{The solution: learning one mode at a time}\label{sec:deeplinear:modes}

We now return to whitened inputs and solve the flow exactly from initializations that respect the singular modes of $\dlS$.

\begin{definition}[Aligned, balanced initialization]\label[definition]{def:deeplinear:aligned}
Assume $d_\ell\ge q$ for $1\le\ell\le L-1$. Fix matrices $R_\ell\in\R^{d_\ell\times q}$ with orthonormal columns, $R_\ell^\top R_\ell=I_q$, for $1\le\ell\le L-1$, and put $R_0=V$ and $R_L=U$ from \eqref{eq:deeplinear:svd}. For a diagonal matrix $B=\diag(\beta_1,\dots,\beta_q)$ with $\beta_i\ge0$, the \term{aligned}\index{aligned initialization} weights are
\[
  W_\ell=R_\ell B R_{\ell-1}^\top,\qquad \ell=1,\dots,L .
\]
\end{definition}

Aligned weights are balanced: $W_{\ell+1}^\top W_{\ell+1}=R_\ell B^2R_\ell^\top=W_\ell W_\ell^\top$. Their partial products telescope, because $R_\ell^\top R_\ell=I_q$ (the first identity for $\ell\le L-1$, the second for $\ell\ge2$):
\begin{equation}\label{eq:deeplinear:telescope}
  W_L\cdots W_{\ell+1}=UB^{L-\ell}R_\ell^\top,\qquad W_{\ell-1}\cdots W_1=R_{\ell-1}B^{\ell-1}V^\top,\qquad W=UB^LV^\top .
\end{equation}
So the end-to-end matrix has the singular vectors of $\dlS$, and its strength on mode $i$ is $a_i=\beta_i^L$. Each layer passes mode $i$ from the direction $R_{\ell-1}e_i$ to the direction $R_\ell e_i$ with gain $\beta_i$.

\begin{theorem}[Exact learning dynamics; \textcite{SaxeEtAl2014}]\label{thm:deeplinear:exact}
Assume \cref{ass:deeplinear:white} and $d_\ell\ge q$ for $1\le\ell\le L-1$, with $L\ge2$. Let $\theta(0)$ be aligned with $B(0)=\diag(\beta_i(0))$, $\beta_i(0)\ge0$. For each $i$ let $\beta_i(t)$ be the solution of
\begin{equation}\label{eq:deeplinear:beta}
  \dot\beta_i=\beta_i^{L-1}\bigl(s_i-\beta_i^L\bigr).
\end{equation}
Then the gradient flow of $\Loss$ exists for all $t\ge0$, and $\theta(t)$ is aligned with $B(t)=\diag(\beta_i(t))$ and the same $R_\ell$, for every $t\ge0$. Consequently $W(t)=\sum_ia_i(t)u_iv_i^\top$ with $a_i=\beta_i^L$, the mode strengths evolve independently by
\begin{equation}\label{eq:deeplinear:aode}
  \dot a_i=L\,a_i^{2-2/L}\,(s_i-a_i),
\end{equation}
and the loss is $\Loss(\theta(t))=\dlLzero+\frac12\sum_{i=1}^q\bigl(s_i-a_i(t)\bigr)^2$.
\end{theorem}

\begin{proof}
The plan is to show that the aligned matrices built from the scalar solutions solve the gradient flow, and then to invoke uniqueness of solutions.

\emph{The scalar equations.} The right side of \eqref{eq:deeplinear:beta} is a polynomial, so a unique solution exists on some interval. If $\beta_i(0)=0$, then $\beta_i\equiv0$ is the solution. If $\beta_i(0)>0$, then the interval $J_i=(0,\max\{\beta_i(0),s_i^{1/L}\}]$ contains $\beta_i(0)$, and the solution cannot leave it: it cannot reach $0$, because $0$ is a stationary point and solutions are unique, and it cannot exceed the upper end $c$, because $\dot\beta_i\le0$ whenever $\beta_i^L\ge s_i$, and $c^L\ge s_i$. A solution confined to a bounded set exists for all $t\ge0$.

\emph{The candidate solves the flow.} Let $\tilde\theta(t)$ be the aligned weights with $B(t)$. By \eqref{eq:deeplinear:telescope}, $\tilde W=UB^LV^\top$, so $G=\tilde W-\dlS=U(B^L-S)V^\top$. By \cref{lem:deeplinear:grad} and \eqref{eq:deeplinear:telescope}, using $U^\top U=V^\top V=I_q$ and that diagonal matrices commute (for $\ell=L$ the factor $A_L=I_{d_y}$ stands where \eqref{eq:deeplinear:telescope} would give $UU^\top$, and for $\ell=1$ the factor $C_1=I_d$ stands for $VV^\top$; nothing changes, since $U^\top UU^\top=U^\top$ and $VV^\top V=V$),
\[
  \nabla_{W_\ell}\Loss(\tilde\theta)=R_\ell B^{L-\ell}U^\top\,U(B^L-S)V^\top\,VB^{\ell-1}R_{\ell-1}^\top=R_\ell\,B^{L-1}(B^L-S)\,R_{\ell-1}^\top .
\]
On the other hand $\frac{\dd}{\dd t}\tilde W_\ell=R_\ell\dot B R_{\ell-1}^\top$, and \eqref{eq:deeplinear:beta} says $\dot B=B^{L-1}(S-B^L)$. Hence $\frac{\dd}{\dd t}\tilde W_\ell=-\nabla_{W_\ell}\Loss(\tilde\theta)$ for every $\ell$ and every $t\ge0$.

\emph{Uniqueness.} The gradient of $\Loss$ is a polynomial in the entries of $\theta$, hence locally Lipschitz, so the gradient flow from $\theta(0)$ has at most one solution on each interval. Since $\tilde\theta$ is a solution on $[0,\infty)$ with $\tilde\theta(0)=\theta(0)$, the flow exists for all time and equals $\tilde\theta$.

\emph{Consequences.} The formula for $W$ is \eqref{eq:deeplinear:telescope}. For $a_i=\beta_i^L$, the chain rule and \eqref{eq:deeplinear:beta} give $\dot a_i=L\beta_i^{L-1}\dot\beta_i=L\beta_i^{2L-2}(s_i-a_i)=La_i^{2-2/L}(s_i-a_i)$. The loss follows from \cref{lem:deeplinear:reduction}, since $\norm{\dlS-W}_F^2=\norm{S-B^L}_F^2$ by the orthonormality of the columns of $U$ and $V$.
\end{proof}

We worked with $\beta_i$ rather than $a_i$ because $a\mapsto a^{2-2/L}$ is not Lipschitz at $0$ when $L\ge3$: equation \eqref{eq:deeplinear:aode} alone would allow solutions that leave $0$ spontaneously, and the layer equation \eqref{eq:deeplinear:beta} rules them out. The mode equation also follows from \cref{thm:deeplinear:e2e}: inserting $W=UD V^\top$ with $D=\diag(a_i)$ into \eqref{eq:deeplinear:e2e} gives $\dot D=-\sum_jD^{2(L-j)/L}(D-S)D^{2(j-1)/L}=LD^{2-2/L}(S-D)$.

\paragraph{Two layers: the logistic curve.} For $L=2$, \eqref{eq:deeplinear:aode} is the logistic equation $\dot a=2a(s-a)$, and everything is explicit.

\begin{corollary}[Two layers]\label[corollary]{cor:deeplinear:logistic}
Let $L=2$, and consider one mode with $a(0)=a_0>0$. If $s>0$, then
\begin{equation}\label{eq:deeplinear:logistic}
  a(t)=\frac{s}{1+(s/a_0-1)\,e^{-2st}},
\end{equation}
and the time that $a$ takes to pass from $a_1$ to $a_2$, for $a_0\le a_1<a_2<s$, is $\frac1{2s}\log\frac{a_2(s-a_1)}{a_1(s-a_2)}$. In particular, if $a_0<s/2$, then the mode is half learned at
\begin{equation}\label{eq:deeplinear:thalf}
  \dlhalf=\frac1{2s}\log\Bigl(\frac s{a_0}-1\Bigr),
\end{equation}
and $a$ rises from $s/10$ to $9s/10$ in time $\frac{\log 81}{2s}=\frac{2\log3}{s}$, independent of $a_0$ when $a_0\le s/10$. If $s=0$, then the mode decays: $a(t)=a_0/(1+2a_0t)$.
\end{corollary}

\begin{proof}
Put $E(t)=(s/a_0-1)e^{-2st}$, so that the right side of \eqref{eq:deeplinear:logistic} is $s/(1+E)$, equal to $a_0$ at $t=0$. Since $\dot E=-2sE$, its derivative is $2s^2E/(1+E)^2$, which equals $2a(s-a)$ because $s-a=sE/(1+E)$; by uniqueness \eqref{eq:deeplinear:logistic} is the solution. Solving \eqref{eq:deeplinear:logistic} for $t$ gives $t(a)=\frac1{2s}\log\frac{a(s-a_0)}{a_0(s-a)}$, and the travel time is $t(a_2)-t(a_1)$. With $a_1=a_0$ and $a_2=s/2$ this is \eqref{eq:deeplinear:thalf}; with $a_1=s/10$ and $a_2=9s/10$ it is $\frac1{2s}\log81$. For $s=0$ the equation is $\dot a=-2a^2$, solved by $a_0/(1+2a_0t)$.
\end{proof}

The logistic curve is the signature of a saddle. Near $a=0$ the mode grows exponentially, $a\approx a_0e^{2st}$, with rate $2s$; the time to grow from $a_0$ to order $s$ is the logarithm of the ratio divided by the rate; and the final approach to $s$ is again exponential, with the same rate. The plateau has length $\frac1{2s}\log(s/a_0)$ and the drop has width of order $1/s$: small initialization lengthens the plateau only logarithmically, which is why stages separate cleanly only when $a_0$ is many orders of magnitude below $s$.

\paragraph{Depth.} For $L\ge3$ the origin is a degenerate saddle: near $0$, $\dot a\approx Ls\,a^{2-2/L}$ with an exponent above $1$, so growth is polynomial and the plateau is polynomially long in $1/a_0$.

\begin{proposition}[Plateau length in depth $L\ge3$]\label[proposition]{prop:deeplinear:depth}
Let $L\ge3$, $s>0$ and $0<a_0<s/2$, and let $\dlhalf$ be the time at which the solution of \eqref{eq:deeplinear:aode} from $a_0$ reaches $s/2$. With $\kappa=1-2/L$,
\[
  \frac{a_0^{-\kappa}-(s/2)^{-\kappa}}{(L-2)\,s}\;\le\;\dlhalf\;\le\;\frac{a_0^{-\kappa}-(s/2)^{-\kappa}}{(L-2)\,s}+\frac{(s/2)^{2/L}}{s^2}.
\]
In particular $\dlhalf=a_0^{-(1-2/L)}/((L-2)s)+O(s^{-(2-2/L)})$ as $a_0\to0$. The time to rise from $s/10$ to $9s/10$ is $c_L\,s^{-(2-2/L)}$ with $c_L=\int_{1/10}^{9/10}\frac{\dd z}{Lz^{2-2/L}(1-z)}$, independent of $a_0$.
\end{proposition}

\begin{proof}
The plan is to write $\dlhalf$ as an integral and to bound its integrand. On $(0,s)$ the right side of \eqref{eq:deeplinear:aode} is positive, so $a$ increases, and separating variables,
\[
  \dlhalf=\int_{a_0}^{s/2}\frac{\dd a}{L\,a^{2-2/L}(s-a)} .
\]
For $0\le a\le s/2$ we have $\frac1s\le\frac1{s-a}=\frac1s+\frac{a}{s(s-a)}\le\frac1s+\frac{2a}{s^2}$. The first term contributes
\[
  \frac1{Ls}\int_{a_0}^{s/2}a^{-1-\kappa}\,\dd a=\frac{a_0^{-\kappa}-(s/2)^{-\kappa}}{L\kappa s}=\frac{a_0^{-\kappa}-(s/2)^{-\kappa}}{(L-2)s},
\]
which is the lower bound. The second term contributes at most
\[
  \frac2{Ls^2}\int_0^{s/2}a^{-1+2/L}\,\dd a=\frac2{Ls^2}\cdot\frac L2\Bigl(\frac s2\Bigr)^{2/L}=\frac{(s/2)^{2/L}}{s^2},
\]
which gives the upper bound. Since $(s/2)^{-\kappa}/s$ and $(s/2)^{2/L}/s^2$ are both of order $s^{-(2-2/L)}$, the asymptotic statement follows. For the rise time, separate variables between $s/10$ and $9s/10$ and substitute $a=sz$.
\end{proof}

With a per-layer initialization scale $\alpha$, so that $a_0=\alpha^L$, the plateau lasts about $\alpha^{-(L-2)}/((L-2)s)$, against $\log(s/\alpha^2)/(2s)$ for two layers: \emph{depth lengthens the plateau and sharpens the transition}, in the sense that the drop takes a time independent of $\alpha$ while the plateau grows like a power of $1/\alpha$.\sidenote{\Textcite{GissinEtAl2020} make the same point about incremental learning: shallow models show it only from exponentially small initializations, deep ones already from polynomially small ones.} The layer equation \eqref{eq:deeplinear:beta} reads $\dot\beta\approx s\beta^{L-1}$ near $0$, the escape equation of an order-$L$ saddle; it is the equation of the search phase in \cref{ch:singleindex} with depth $L$ in the role of the information exponent, and the escape time $\beta_0^{-(L-2)}/((L-2)s)$ has the same form.

In both cases \emph{the order of learning is the order of the singular values}: for equal initial strengths, a larger $s_i$ gives an earlier $\dlhalf$, by \eqref{eq:deeplinear:thalf} and \cref{prop:deeplinear:depth}. The loss therefore descends a staircase. Between the learning of mode $k$ and of mode $k+1$ it rests near
\begin{equation}\label{eq:deeplinear:staircase}
  \Loss_k=\dlLzero+\frac12\sum_{i>k}s_i^2 ,
\end{equation}
the loss of the end-to-end matrix $W_k=\sum_{i\le k}s_iu_iv_i^\top$, the best approximation of $\dlS$ of rank $k$.

\paragraph{The landscape.} The points $W_k$ are not only resting places of one trajectory: they are the critical points of the loss. \Textcite{BaldiHornik1989} proved this for two layers with a bottleneck, in the language of principal component analysis.

\begin{proposition}[Critical points of two-layer networks; after \textcite{BaldiHornik1989}]\label[proposition]{prop:deeplinear:critical}
Let $L=2$, assume \cref{ass:deeplinear:white}, and let $s_1>s_2>\dots>s_q>0$. Let $(W_1,W_2)$ be a critical point of $\Loss$ with $W_2\in\R^{d_y\times m}$ of full column rank $m$. Then
\begin{enumerate}
\item $W=W_2W_1=\sum_{i\in I}s_iu_iv_i^\top$ for a set $I\subseteq\{1,\dots,q\}$ with $\abs I\le m$, and $\Loss=\dlLzero+\frac12\sum_{i\notin I}s_i^2$;
\item if $m\le q$ and $I\ne\{1,\dots,m\}$, then the critical point is a saddle: every neighborhood contains points of lower loss and points of higher loss. Hence every local minimum with $W_2$ of full column rank has $W=\sum_{i\le m}s_iu_iv_i^\top$, and it is a global minimum over the networks of width $m$.
\end{enumerate}
\end{proposition}

\begin{proof}
The plan for part (i) is to solve the two critical-point equations in turn. By \cref{lem:deeplinear:grad} they read $W_2^\top(\dlS-W)=0$ and $(\dlS-W)W_1^\top=0$. Since $W_2^\top W_2$ is invertible, the first gives $W_1=(W_2^\top W_2)^{-1}W_2^\top\dlS$, so that $W=\Pi\dlS$ with $\Pi=W_2(W_2^\top W_2)^{-1}W_2^\top$, the orthogonal projection onto the range of $W_2$, of rank $m$. Inserting $W_1^\top=\dlS^\top W_2(W_2^\top W_2)^{-1}$ into the second equation and multiplying it on the right by $W_2^\top$ gives
\[
  (I-\Pi)\,\dlS\dlS^\top\,\Pi=0 ,
\]
that is, $\dlS\dlS^\top\Pi=\Pi\dlS\dlS^\top\Pi$. The right side is symmetric, so transposing gives $\Pi\dlS\dlS^\top=\dlS\dlS^\top\Pi$: the projection commutes with $\dlS\dlS^\top=\sum_is_i^2u_iu_i^\top$. Its range is therefore invariant under $\dlS\dlS^\top$, and because the eigenvalues $s_1^2>\dots>s_q^2>0$ are simple, $\Pi=\sum_{i\in I}u_iu_i^\top+\Pi_0$ for a set $I$ and a projection $\Pi_0$ onto a subspace of the kernel of $\dlS^\top$. Then $\Pi_0\dlS=0$ and $W=\Pi\dlS=\sum_{i\in I}s_iu_iv_i^\top$; the rank count gives $\abs I\le m$, and the loss follows from \cref{lem:deeplinear:reduction}.
\end{proof}

\begin{proof}[Proof of part (ii)]
The plan is to rotate the range of $W_2$ slightly, from a weak direction toward a strong one, and to compare losses with the identity $\Loss=\dlLzero+\frac12\norm{\dlS}_F^2-\frac12\tr(P\dlS\dlS^\top)$, valid whenever $W=P\dlS$ for an orthogonal projection $P$ (expand $\frac12\norm{(I-P)\dlS}_F^2$ and use $P^2=P$). First we find two directions. If $\abs I=m$, then some $j\le m$ lies outside $I$ and some $i\in I$ exceeds $m$; put $s_i$ for its singular value. If $\abs I<m$, then $\Pi_0\ne0$; let $u_i$ be a unit vector in its range, which lies in the kernel of $\dlS^\top$, put $s_i=0$, and pick some $j\notin I$, which exists because $\abs I<m\le q$. In both cases $s_j>s_i$, $u_i$ lies in the range of $\Pi$, and $\Pi u_j=0$, because $u_j$ is orthogonal to the $u_k$ with $k\in I$ and to the kernel of $\dlS^\top$.

Let $Q_\eps$ be the rotation by the angle $\eps$ in the plane of $u_i$ and $u_j$ that takes $u_i$ to $\cos\eps\,u_i+\sin\eps\,u_j$ and fixes the orthogonal complement of the plane. Put $W_2(\eps)=Q_\eps W_2$ and $W_1(\eps)=(W_2^\top W_2)^{-1}W_2^\top Q_\eps^\top\dlS$; both depend continuously on $\eps$ and equal $(W_2,W_1)$ at $\eps=0$. Their product is $Q_\eps\Pi Q_\eps^\top\dlS=\Pi_\eps\dlS$, with $\Pi_\eps=Q_\eps\Pi Q_\eps^\top$ an orthogonal projection. Write $\Pi=u_iu_i^\top+\Pi'$; the range of $\Pi'$ is orthogonal to $u_i$ and to $u_j$, so $Q_\eps$ fixes it and $\tr(\Pi'Q_\eps^\top\dlS\dlS^\top Q_\eps)=\tr(\Pi'\dlS\dlS^\top)$. The remaining term changes because $u_i$ and $u_j$ are eigenvectors of $\dlS\dlS^\top$ with eigenvalues $s_i^2$ and $s_j^2$:
\[
  (Q_\eps u_i)^\top\dlS\dlS^\top(Q_\eps u_i)=\cos^2\!\eps\,s_i^2+\sin^2\!\eps\,s_j^2 .
\]
Hence $\Loss$ at $(W_1(\eps),W_2(\eps))$ is lower than at $(W_1,W_2)$ by $\frac12(s_j^2-s_i^2)\sin^2\eps>0$ for every $\eps\in(0,\pi/2)$. For points of higher loss, move $W_1$ alone to $W_1+\eps X$ with $W_2X\ne0$: the first-order change of $\Loss$ vanishes by the first critical-point equation, and the loss rises by $\frac12\eps^2\norm{W_2X}_F^2$. For the last statement, a local minimum is not a saddle, so $I=\{1,\dots,m\}$, and by the Eckart--Young theorem $\sum_{i\le m}s_iu_iv_i^\top$ minimizes $\norm{\dlS-W}_F$ over all matrices of rank at most $m$, which contains every product $W_2W_1$ of width $m$.
\end{proof}

Critical points with $W_2$ of deficient rank require a separate argument. \Textcite{BaldiHornik1989} give the full classification, and \textcite{Kawaguchi2016} extends ``every local minimum is global'' to deep linear networks, and shows that for $L\ge3$ some saddles have no direction of negative curvature.

The trajectory of \cref{thm:deeplinear:exact} from a small aligned initialization therefore passes from the saddle $W_0=0$ near the saddles $W_1,W_2,\dots$ in turn, lingering near each until the next mode escapes, and the plateaus of the staircase are visits to saddles.

\section{The solution: small random initialization}\label{sec:deeplinear:random}

No one initializes a network in alignment with the singular vectors of the task. What happens from a small \emph{random} initialization? The answer has two phases: first the weights align without the loss moving, and then the aligned solution takes over.

\begin{heuristic}[Heuristic derivation, $L=2$]
Let $W_1(0)$ and $W_2(0)$ be small, with entries of order $\alpha$. While $W=W_2W_1$ is negligible against $\dlS$, the flow is approximately linear: $\dot W_1\approx W_2^\top\dlS$ and $\dot W_2\approx\dlS W_1^\top$. Project onto the modes: with $x_i=W_1v_i\in\R^{m}$ and $z_i=W_2^\top u_i\in\R^m$, the linearized flow gives $\dot x_i=s_iz_i$ and $\dot z_i=s_ix_i$, because $\dlS v_i=s_iu_i$ and $\dlS^\top u_i=s_iv_i$. The modes decouple: $x_i+z_i$ grows like $e^{s_it}$ and $x_i-z_i$ decays like $e^{-s_it}$. After a time of a few multiples of $1/s_i$,
\[
  x_i(t)\approx z_i(t)\approx\xi_ie^{s_it},\qquad \xi_i=\tfrac12\bigl(W_1(0)v_i+W_2(0)^\top u_i\bigr),
\]
and the mode strength $a_i=u_i^\top Wv_i=\ip{z_i,x_i}$ is close to $\norm{\xi_i}^2e^{2s_it}$. That is the early part of the logistic curve \eqref{eq:deeplinear:logistic} with initial strength
\begin{equation}\label{eq:deeplinear:aeff}
  \hat a_i=\norm{\xi_i}^2=\tfrac14\norm{W_1(0)v_i+W_2(0)^\top u_i}^2 .
\end{equation}
At this point the weights are aligned and balanced up to errors that decay: $x_i\approx z_i$, and the hidden directions $\xi_i/\norm{\xi_i}$ of different modes are nearly orthogonal when $m$ is much larger than the number of modes, since independent random vectors in $\R^m$ are. The aligned solution with initial strengths $\hat a_i$ then describes the rest of the trajectory. For a mode with multiplicity, $\hat a_i$ is replaced by the eigenvalues of the Gram matrix of the $\xi_i$ in its eigenspace. With i.i.d.\ $\cN(0,\alpha^2)$ entries, $\hat a_i\approx\frac12m\alpha^2$.\qedhere
\end{heuristic}

A network whose second layer starts at zero still escapes: then $z_i(0)=0$, but $x_i+z_i$ starts at $x_i(0)$ and grows all the same, and $\hat a_i=\frac14\norm{W_1(0)v_i}^2$.\sidenote{This is why the common practice of initializing a last layer at zero does not stall training, while initializing \emph{every} layer at zero does: the origin is a stationary point of the flow.} The first phase is the \term{silent alignment}\index{silent alignment} of \textcite{AtanasovEtAl2022}: the directions of the weights change by order one while the loss, and the function, hardly move. It explains why the aligned solution describes random initializations at all; \cref{fig:deeplinear:tree-dyn} shows how accurately. For deeper networks the first phase is itself nonlinear ($\dot\beta\approx s\beta^{L-1}$), and we know no closed formula for the effective initial strength; the scaling $\dlhalf\propto\alpha^{-(L-2)}$ survives, with a prefactor (\cref{fig:deeplinear:depth}).

\begin{remark}[Saddle to saddle]\label[remark]{rem:deeplinear:saddle}
For losses other than $\frac12\norm{\dlS-W}_F^2$, such as matrix sensing or completion, the modes do not decouple, but the staircase survives in a weaker form. \Textcite{LiLuoLyu2021} show that for two layers, gradient flow from an infinitesimal initialization is equivalent, under assumptions, to \emph{greedy low-rank learning}: escape from the current saddle along its top eigendirection, which raises the rank by one, then minimize at that rank, and repeat. \Textcite{JacotEtAl2021} (an arXiv preprint) conjecture the same \term{saddle-to-saddle}\index{saddle-to-saddle dynamics} picture for deep linear networks, with saddles at maps of increasing rank, and prove a theorem for the dynamics between the first two saddles. \Textcite{GidelEtAl2019} show that for two layers and vanishing initialization the iterates of gradient descent pass through the solutions of reduced-rank regression of increasing rank.
\end{remark}

\section{Reading the solution}\label{sec:deeplinear:reading}

\paragraph{Predictions.} The model predicts five things. (i) \emph{A staircase}: the loss rests near the values \eqref{eq:deeplinear:staircase}, one plateau per distinct singular value. (ii) \emph{Order and timing}: mode $i$ is half learned at $\dlhalf\approx\frac1{2s_i}\log(s_i/\hat a_i)$ for two layers, so strong modes come first, and the delay between modes $i$ and $i+1$ is about $\log(s/\hat a)$ times the difference of $\frac1{2s_{i+1}}$ and $\frac1{2s_i}$. (iii) \emph{Sharpness}: each drop takes time $2\log3/s_i$ for $L=2$, independent of the initialization. (iv) \emph{Depth}: with per-layer scale $\alpha$ the plateau scales like $\log(1/\alpha)$ for $L=2$ and like $\alpha^{-(L-2)}$ for $L\ge3$. (v) \emph{The endpoint}: when every hidden width is at least $q$, $W(t)\to\dlS$, the least-squares solution; with a bottleneck of width $m<q$, the aligned trajectory with $R_\ell$ of $m$ columns reaches the best rank-$m$ approximation, the solution of reduced-rank regression. A linear autoencoder ($y=x$) needs inputs that are not whitened, since otherwise $\dlS=I_d$; with $\Sigma_{xx}=V\Lambda V^\top$ the decoupled flow of \cref{ex:deeplinear:nonwhite} applies, and a bottleneck of width $m$ learns the projection onto the top $m$ principal components, in the order of the eigenvalues.\sidenote{This is the connection \textcite{BaldiHornik1989} made: a linear autoencoder with a bottleneck, trained on the squared error, has a unique minimum up to invertible transformations of the hidden layer, the projection onto the leading principal subspace.}

\paragraph{Hierarchies.} For the tree of \cref{fig:deeplinear:tree} the modes are the levels of the tree, and their strengths decrease with depth in the tree: in a complete binary tree of depth $D$ the mode that splits a node at depth $j$ (the root has depth $0$) has $s^2=(2^{D-j}-1)/2^D$, and the constant mode has $s^2=(2^{D+1}-1)/2^D$ (\cref{ex:deeplinear:tree}). In \cref{fig:deeplinear:tree}, level $0$ is the constant mode and level $k\ge1$ consists of the $2^{k-1}$ modes that split the nodes at depth $k-1$, that is, the modes that tell apart the nodes of the $k$-th row. Broad distinctions are stronger because they are supported by more items and more output features, and the network learns ``animal versus plant'' first and ``canary versus robin'' last. This is the progressive differentiation of \cref{sec:deeplinear:phenomenon}, explained by one statistic: the singular value.

\begin{figure}[t]
\begin{fullwidth}
\centering
\begin{tikzpicture}
\begin{groupplot}[group style={group size=2 by 1,horizontal sep=17mm},hplot,width=0.475\linewidth,height=0.37\linewidth]
\nextgroupplot[xmin=0,xmax=20,ymin=-0.03,ymax=1.05,xlabel={time $t$},ylabel={fraction learned $a/s$},
  legend style={at={(0.97,0.05)},anchor=south east}]
\addplot[theory] table[x=t,y=th1] {figures/data/02-deeplinear-tree.dat};
\addplot[simulation,mark repeat=4] table[x=t,y=sim1] {figures/data/02-deeplinear-tree.dat};
\addplot[theory] table[x=t,y=th2] {figures/data/02-deeplinear-tree.dat};
\addplot[simulation,mark repeat=4] table[x=t,y=sim2] {figures/data/02-deeplinear-tree.dat};
\addplot[theory] table[x=t,y=th3] {figures/data/02-deeplinear-tree.dat};
\addplot[simulation,mark repeat=4] table[x=t,y=sim3] {figures/data/02-deeplinear-tree.dat};
\addplot[theory] table[x=t,y=th4] {figures/data/02-deeplinear-tree.dat};
\addplot[simulation,mark repeat=4] table[x=t,y=sim4] {figures/data/02-deeplinear-tree.dat};
\node[font=\scriptsize,text=ink] at (axis cs:2.0,0.62) {$0$};
\node[font=\scriptsize,text=ink] at (axis cs:4.7,0.38) {$1$};
\node[font=\scriptsize,text=ink] at (axis cs:7.3,0.38) {$2$};
\node[font=\scriptsize,text=ink] at (axis cs:11.5,0.38) {$3$};
\legend{logistic prediction,gradient descent}
\nextgroupplot[xmin=0,xmax=20,ymin=0,ymax=2.15,xlabel={time $t$},ylabel={loss $\Loss$},
  legend style={at={(0.97,0.95)},anchor=north east}]
\addplot[guide] coordinates {(0,1.0625) (20,1.0625)};
\addplot[guide] coordinates {(0,0.625) (20,0.625)};
\addplot[guide] coordinates {(0,0.25) (20,0.25)};
\addplot[theory] table[x=t,y=lossth] {figures/data/02-deeplinear-tree.dat};
\addplot[simulation,mark repeat=3] table[x=t,y=losssim] {figures/data/02-deeplinear-tree.dat};
\legend{staircase \eqref{eq:deeplinear:staircase},,,prediction,gradient descent}
\end{groupplot}
\end{tikzpicture}
\caption{Learning the tree of \cref{fig:deeplinear:tree} with a two-layer linear network of width $m=64$, by full-batch gradient descent with step $\eta=0.01$ (time $t=\eta\times$steps) from a random initialization with i.i.d.\ $\cN(0,\alpha^2)$ entries, $\alpha=0.003$, seed $0$. Left: the fraction learned, $\tr(U_k^\top WV_k)/(q_ks_k)$ for the $q_k$ modes $U_k,V_k$ of level $k$ (labels $0$ to $3$), against the logistic curves \eqref{eq:deeplinear:logistic} with the effective initial strengths \eqref{eq:deeplinear:aeff} read off from the initialization (between $1.8\times10^{-4}$ and $3.8\times10^{-4}$; for the levels with several modes, the curves are averaged). The largest deviation is $0.017$. Right: the loss, $\frac12\norm{\dlS-W}_F^2$, descends the staircase $2$, $1.0625$, $0.625$, $0.25$, $0$ of \eqref{eq:deeplinear:staircase}.}\label{fig:deeplinear:tree-dyn}
\end{fullwidth}
\end{figure}

\Cref{fig:deeplinear:tree-dyn} tests predictions (i) and (ii) from a random initialization. The four levels are half learned at $t=3.1$, $4.2$, $6.3$, and $10.0$, as predicted by \eqref{eq:deeplinear:thalf} with the effective initial strengths, and the agreement holds over the whole trajectory, not only at the midpoints, to within $0.02$. The loss sits on each step of the staircase for a time that grows as the singular value shrinks, because both the rate $2s$ and the log-ratio $\log(s/\hat a)$ favor strong modes. The residual deviation comes from the cross-talk between hidden directions that the heuristic neglects. With width $m=16$ instead of $64$, where random hidden directions are further from orthogonal, it grows to about $0.1$, for the weakest level.

\begin{figure}[t]
\centering
\begin{tikzpicture}
\begin{axis}[hplot,xmode=log,ymode=log,xmin=8e-4,xmax=0.7,ymin=0.25,ymax=800,xlabel={per-layer initialization scale $\alpha$},ylabel={half-learning time $\dlhalf$},
  legend style={at={(0.5,1.03)},anchor=south},legend columns=3]
\addplot[theory] table[x=alpha,y=L2] {figures/data/02-deeplinear-depth-theory.dat};
\addplot[theory2] table[x=alpha,y=L3] {figures/data/02-deeplinear-depth-theory.dat};
\addplot[theory3] table[x=alpha,y=L4] {figures/data/02-deeplinear-depth-theory.dat};
\addplot[simulation2] table[x=alpha,y=aligned] {figures/data/02-deeplinear-depth-sim.dat};
\addplot[simulation] table[x=alpha,y=random] {figures/data/02-deeplinear-depth-sim.dat};
\legend{$L=2$ exact,$L=3$ exact,$L=4$ exact,aligned (GD),random (GD)}
\end{axis}
\end{tikzpicture}
\caption{Depth lengthens the plateau. One mode of strength $s=1$ in $\R^{4\times4}$, hidden width $32$, full-batch gradient descent with step $0.02$; $\dlhalf$ is the first time $u_1^\top Wv_1\ge\frac12$. Curves: the exact half-learning time of \cref{thm:deeplinear:exact} with $a_0=\alpha^L$, from \eqref{eq:deeplinear:thalf} for $L=2$ and by quadrature for $L=3,4$. Squares: aligned, balanced initialization with per-layer gain $\alpha$ (seed $0$), within about $2\%$ of the curves. Dots: Gaussian initialization with entries $\cN(0,\alpha^2/d_{\ell-1})$, averaged over seeds $0$, $1$, $2$. For $L=2$ these are shifted down by about $0.4$, the shift $\frac12\log2.25$ predicted by \eqref{eq:deeplinear:aeff}, which gives $\hat a\approx2.25\alpha^2$ here; for $L=3$ and $L=4$ they follow the slopes $-(L-2)$ with a factor of about $0.7$.}\label{fig:deeplinear:depth}
\end{figure}

\Cref{fig:deeplinear:depth} tests prediction (iv). On logarithmic axes the two-layer curve is nearly flat, a logarithm, while the curves for $L=3$ and $L=4$ are straight lines of slopes $-1$ and $-2$, as \cref{prop:deeplinear:depth} predicts: with $\alpha=0.04$, two layers learn the mode at $t\approx3.2$ and four layers at $t\approx310$. Gradient descent from aligned initializations reproduces the exact curves, which checks \cref{thm:deeplinear:exact} together with the small error of a finite step. Random initializations learn somewhat faster, because their effective initial strength exceeds $\alpha^L$, but the exponent is the same.

\paragraph{Limits and regimes.} The staircase needs two separations: of the singular values (distinct steps) and of scales ($\hat a_i\ll s_i$, long plateaus). When the initialization is large, $\hat a_i$ is comparable with $s_i$, the logistic curves start near their midpoints, and all modes are learned together at rate about $2s_i$, which is the lazy, kernel-like regime of \cref{ch:lazy}; for two layers and aligned weights, the flow is then close to the linear flow $\dot W=-2a_0(W-\dlS)$ with a fixed preconditioner. Small initialization is the rich regime, and the staircase is its signature: the order parameters are the mode strengths $a_i$, and the plateau times are their escape times from the saddle at zero.

\section{What the model gets wrong}\label{sec:deeplinear:wrong}

\paragraph{What was thrown away.} \emph{The nonlinearity.} A deep linear network represents only linear maps, so it learns no features that a linear regression could not, and its hidden layers are rotations and rescalings of the modes of $\dlS$; whatever nonlinear networks gain from depth in expressive power is absent. \emph{Whitening.} With $\Sigma_{xx}\ne I$ the gradient is $W\Sigma_{xx}-\dlS$, and the modes decouple only when $\Sigma_{xx}$ and $\dlS$ share right singular vectors (\cref{ex:deeplinear:nonwhite}); in general they do not, and the input correlations mix the modes. \emph{Aligned and balanced initialization.} The heuristic of \cref{sec:deeplinear:random} recovers it from small random weights, but only approximately, and only for two layers in closed form; large or unbalanced initializations change the dynamics qualitatively \parencite{KuninEtAl2024}. \emph{Population loss.} With $n$ samples, $\dlS$ is replaced by a noisy estimate whose weak modes are dominated by noise, and the late stages fit noise: \textcite{AdvaniEtAl2020} and \textcite{LampinenGanguli2019} compute the resulting generalization dynamics, in which early stopping helps. \emph{Gradient flow.} With large steps, gradient descent is not close to the flow; for the scalar product this is the catapult of \cref{ch:stability}.

\paragraph{Where it breaks.} \emph{Tasks without linear correlation.} If $\dlS=0$, every mode has $s_i=0$, and by \cref{thm:deeplinear:exact} and \cref{cor:deeplinear:logistic} a linear network learns nothing: every aligned $a_i$ decays like $1/t$. XOR, sparse parities and single-index targets of information exponent at least $2$ have $\dlS=0$ (\cref{ch:singleindex,ch:parity}), yet nonlinear networks learn them, after plateaus whose length the linear model cannot predict. \emph{Stage order in nonlinear networks.} \Textcite{RefinettiEtAl2023} show that networks trained with SGD on images first use the mean and covariance of each class and only later higher-order statistics, a distributional simplicity bias that the linear model shares only in its first stage. \Textcite{AbbeEtAl2023} prove, under assumptions on the training procedure, a saddle-to-saddle dynamics for two-layer nonlinear networks learning functions of a few Gaussian coordinates, in which the order and the duration of the stages are set by the hierarchical structure of the target (its leap complexity; they conjecture times of order $d^{\max(\text{leap},2)}$ in dimension $d$, up to logarithmic factors), not by singular values. \Textcite{KalimerisEtAl2019} see the linear stage first in convolutional networks, as the model predicts, but what follows is not a linear map. \emph{Implicit bias.} From small initialization the staircase suggests a bias toward low rank. \Textcite{GunasekarEtAl2017} conjectured that for matrix factorization the bias is the minimum nuclear norm. \Textcite{AroraEtAl2019} found that depth strengthens the bias toward low rank and argued that it may not be captured by the nuclear norm or other norms, and \textcite{RazinCohen2020} constructed matrix factorization problems in which gradient flow drives every norm and quasi-norm to infinity while the rank is minimized. \Textcite{LiLuoLyu2021} give counterexamples to the nuclear norm conjecture. The model's stages are therefore a better description of the bias than a norm, but the model with whitened, fully observed data cannot exhibit this: its endpoint is always $\dlS$.

\paragraph{Fine structure.} \emph{Unbalanced and large initializations.} \Textcite{TarmounEtAl2021} analyze two-layer flows with an arbitrary conserved $D_1$, without small, balanced or aligned initialization, and relate them to Riccati equations; \textcite{BraunEtAl2022} give exact solutions from initializations that are not aligned with the task, building on \textcite{Fukumizu1998}, and use them to study prior knowledge and continual learning; \textcite{DomineEtAl2025} solve the family of initializations they call $\lambda$-balanced, with $D_1$ a multiple of the identity, which interpolates between the rich regime of this chapter and the lazy regime; and \textcite{KuninEtAl2024} show that unbalanced initializations that make earlier layers learn faster can speed up feature learning in nonlinear networks. \emph{Gates.} \Textcite{SaxeSodhaniLewallen2022} add input-dependent gates to linear pathways, which yields solvable models of nonlinear networks in which pathways race and shared representations win. \emph{Saddle-to-saddle.} \Textcite{LiLuoLyu2021}, \textcite{JacotEtAl2021} and \textcite{GidelEtAl2019} extend the staircase to general losses, \textcite{PesmeFlammarion2023} make it exact for diagonal networks (\cref{ch:diagonal}), and \textcite{AbbeEtAl2023} for nonlinear networks on hierarchical targets. \emph{Convergence.} \Textcite{AroraEtAl2019conv} prove linear convergence of gradient descent for deep linear networks from approximately balanced initializations whose loss is below that of every rank-deficient solution. Open: a theory that predicts the order and timing of stages in nonlinear networks on natural data, and a description of the implicit bias of deep factorizations from small initialization for general measurement operators.

\section{Model card}\label{sec:deeplinear:card}

\modelcard{Plateaus and sudden drops in training loss; coarse distinctions learned before fine ones; incremental, low-rank learning from small initialization.}{Deep linear network $W_L\cdots W_1$, square loss, whitened inputs, gradient flow, aligned (or small random) initialization.}{Conserved $W_{\ell+1}^\top W_{\ell+1}-W_\ell W_\ell^\top$; closed end-to-end dynamics for balanced weights; exact decoupled mode dynamics $\dot a=La^{2-2/L}(s-a)$, logistic for two layers; critical points are partial singular value decompositions.}{A loss staircase $\dlLzero+\frac12\sum_{i>k}s_i^2$; mode $i$ learned at $\log(s_i/a_0)/(2s_i)$ ($L=2$) or $a_0^{-(1-2/L)}/((L-2)s_i)$ ($L\ge3$); drop width independent of initialization; order set by singular values of $\dlS$.}{Nothing is learned when $\dlS=0$ (XOR, parity); nonlinear stages follow higher-order structure; the implicit bias toward low rank is not a norm; no finite-sample or step-size effects.}{Unbalanced and $\lambda$-balanced exact solutions; gated linear networks; saddle-to-saddle dynamics for general losses, diagonal networks, and nonlinear networks.}

This chapter opens the book's longest thread: \emph{escape from a saddle at small initialization}. The logistic escape here, with time $\log(1/\text{initialization})/\text{strength}$, reappears coordinate by coordinate in the diagonal networks of \cref{ch:diagonal}, as the search phase of \cref{ch:singleindex} (where depth is replaced by the information exponent), as the plateaus of the committee machine in \cref{ch:committee}, as hidden progress in sparse parities (\cref{ch:parity}), and in the formation of induction heads (\cref{ch:induction}); the quadratic of \hyperref[ch:prologue]{the prologue} is the degenerate case with no saddle. The scalar case of the conservation law is the invariant of the two-parameter model in \cref{ch:stability}.

The second thread is \emph{lazy versus rich}. Large initializations turn the staircase into simultaneous exponential relaxation, the regime of \cref{ch:lazy}; small ones produce the rich dynamics of this chapter, and \cref{ch:diagonal} shows how the same choice selects an implicit bias.

\begin{exercises}
\item \consolidation (The logistic curve.) (a) Integrate $\dot a=2a(s-a)$ by partial fractions and recover \cref{cor:deeplinear:logistic}. (b) Show that the loss of one mode, $\frac12(s-a)^2$, decreases fastest in time when $a=s/3$, and compute the time at which it occurs. (c) For the tree of \cref{fig:deeplinear:tree-dyn}, with mean effective initial strengths $2.8\times10^{-4}$, $3.5\times10^{-4}$, $2.7\times10^{-4}$ and $2.9\times10^{-4}$ for the four levels, predict the four half-learning times and compare with the values $3.1$, $4.2$, $6.3$ and $10.0$ quoted in the text.\label{ex:deeplinear:logistic}

\item \consolidation (The tree.) Let $P=2^D$ items sit at the leaves of a complete binary tree of depth $D$, with $x=\sqrt Pe_p$ and one output coordinate per node, as in \cref{fig:deeplinear:tree}. (a) Show that $(\dlS^\top\dlS)_{pp'}$ is $\frac1P$ times the number of common ancestors of $p$ and $p'$ (the item itself counts if $p=p'$). (b) Show that the constant vector is an eigenvector with eigenvalue $(2^{D+1}-1)/P$, and that the Haar function of a node at depth $j<D$ (the root has depth $0$) is an eigenvector with eigenvalue $(2^{D-j}-1)/P$. (c) Conclude that the $2^j$ modes that split the nodes at depth $j$ have $s^2=(2^{D-j}-1)/2^D$, and that the loss staircase for $D=3$ is $2$, $1.0625$, $0.625$, $0.25$, $0$.\label{ex:deeplinear:tree}

\item \consolidation (Neurons of a ReLU network are balanced.) Let $f(x)=\sum_{r=1}^ma_r\phi(\ip{w_r,x})$ with $\phi(z)=\max(z,0)$, and let the parameters follow the gradient flow of $\E[\ell(f(x),y)]$ for a differentiable loss $\ell$, with the convention $\phi'(0)=0$. Using $\phi'(z)z=\phi(z)$, show that $\norm{w_r}^2-a_r^2$ is constant in time for every neuron $r$. What does this say about a neuron whose output weight is driven to zero?\label{ex:deeplinear:relu}

\item \consolidation (End-to-end dynamics for two layers.) (a) For $L=2$ and balanced weights, show directly that $W_2W_2^\top=(WW^\top)^{1/2}$. (b) For a scalar output, $W=w^\top$ with $w\in\R^d$ and gradient $g=\nabla\dlE$ written as a vector, show that \eqref{eq:deeplinear:e2e} becomes $\dot w=-\norm w^{2-2/L}\bigl(g+(L-1)\ip{g,\hat w}\hat w\bigr)$ with $\hat w=w/\norm w$: depth speeds up motion along $w$ by the factor $L$. (c) Show that for $W=UDV^\top$ with $D$ diagonal, \eqref{eq:deeplinear:e2e} gives back \eqref{eq:deeplinear:aode}.\label{ex:deeplinear:e2e}

\item \consolidation (Escape from a degenerate saddle.) (a) Solve $\dot\beta=s\beta^{L-1}$ exactly for $L\ge3$, and show that the solution blows up at time $\beta_0^{-(L-2)}/((L-2)s)$. (b) Explain why this time is the leading term of \cref{prop:deeplinear:depth} with $a_0=\beta_0^L$. (c) Compare with the escape time of the search phase in \cref{ch:singleindex}, and identify the depth with the information exponent.\label{ex:deeplinear:escape}

\item \variation (Unbalanced scalar network.) Let $L=2$ and $d=d_1=d_y=1$, with weights $b=W_1$ and $c=W_2$ and $\Loss=\frac12(s-cb)^2$. (a) Show that $\delta=b^2-c^2$ is conserved and that $a=cb$ satisfies $\dot a=\sqrt{4a^2+\delta^2}\,(s-a)$. (b) Show that from $a(0)=0$ and $\delta\ne0$ the mode reaches $s/2$ by time $\log2/\abs\delta$, so a strongly unbalanced network has no plateau. (c) Show that as $\delta\to0$ with $a(0)=a_0>0$ fixed, the solution converges to the logistic curve. (d) Interpret $\abs\delta$ as an initialization scale that makes training lazy.\label{ex:deeplinear:unbalanced}

\item \variation (Non-whitened inputs.) Let $\Sigma_{xx}=V\Lambda V^\top$ and $\dlS=USV^\top$ share the matrix $V$, with $\Lambda=\diag(\lambda_1,\dots,\lambda_q)$ positive (in this exercise $\lambda_i$ are eigenvalues). Show that the aligned ansatz still solves the flow for $L=2$, with $\dot a_i=2a_i(s_i-\lambda_ia_i)$. Which modes are learned first, and to what values do they converge? Show that the order of learning is set by $s_i$ while the final value is $s_i/\lambda_i$, and that in general the two orders differ.\label{ex:deeplinear:nonwhite}

\item \variation (Weight decay.) Add $\frac\lambda2(\norm{W_1}_F^2+\norm{W_2}_F^2)$ to the loss of a two-layer network with whitened inputs. (a) Show that $D_1(t)=e^{-2\lambda t}D_1(0)$, so weight decay balances the layers. (b) Show that the aligned ansatz still solves the flow, with $\dot a_i=2a_i(s_i-\lambda-a_i)$. (c) Conclude that the network converges to the soft-thresholded matrix $\sum_i(s_i-\lambda)_+u_iv_i^\top$ when every $a_i(0)>0$, and that modes with $s_i\le\lambda$ are never learned.\label{ex:deeplinear:wd}

\item \variation\hard (Depth-independent learning time.) For $L\ge3$ and fixed $0<a_0<s/2$, show that $L\,\dlhalf\to\int_{a_0}^{s/2}\frac{\dd a}{a^2(s-a)}=\frac1{sa_0}-\frac2{s^2}+\frac1{s^2}\log\frac{s-a_0}{a_0}$ as $L\to\infty$. Interpret: if the end-to-end initial strength is held fixed (each layer starts near the identity, with gain $a_0^{1/L}\to1$) and the step size is scaled like $1/L$, the learning time does not grow with depth. Compare with the per-layer scaling $\alpha$ of \cref{fig:deeplinear:depth}, and with the depth-independent learning times that \textcite{SaxeEtAl2014} report for special initializations. \hint{$a^{-2+2/L}\le\max(a^{-2},1)$ for $a>0$; use dominated convergence.}\label{ex:deeplinear:deep}

\item \simulation (Depth on the tree.) Modify the script of \cref{fig:deeplinear:tree-dyn} to train a three-layer network ($d_1=d_2=64$) from i.i.d.\ $\cN(0,\alpha^2/d_{\ell-1})$ entries with $\alpha=0.1$, and plot the fraction learned for each level. With step $0.01$, the reader should observe four stages, half learned near $t\approx5$, $7$, $10$ to $12$, and $18$ (the exact values vary with the seed by about $10\%$), each drop narrow compared with the plateau before it; and with $\alpha=1$ all levels are learned together, $90\%$ learned by $t\approx0.2$. Explain both observations with \cref{prop:deeplinear:depth}.\label{ex:deeplinear:sim-depth}

\item \simulation (XOR.) Let $x\in\{\pm1\}^2$ be uniform, augmented by a constant coordinate, and $y=x_1x_2$. (a) Compute $\Sigma_{xx}$ and $\dlS$, and show that $\dlS=0$. (b) Train a two-layer linear network and a two-layer ReLU network of width $16$ from small random initialization ($\cN(0,10^{-4})$ entries) by gradient descent with step $0.05$. The reader should observe that the linear network's loss stays at $\frac12$, while for most seeds the ReLU network reaches zero loss within about $500$ steps (some seeds stall at loss $\frac18$, with one input misfit): the ReLU network learns a function that the linear model says carries no signal.\label{ex:deeplinear:xor}

\item \research (Stage order in nonlinear networks.) For a two-layer ReLU network trained from small initialization on the tree data of \cref{fig:deeplinear:tree}, with inputs given by random features of the items rather than one-hot vectors, determine whether the stages still follow the singular values of $\dlS$, and how the answer changes when the outputs depend nonlinearly on the features. \Textcite{SaxeSodhaniLewallen2022} solve gated networks with fixed gates, and \textcite{AbbeEtAl2023} treat functions of a few Gaussian coordinates; a theory of stage order for nonlinear networks on structured, non-Gaussian data is open.\label{ex:deeplinear:research-order}

\item \research (Saddle to saddle in depth.) Prove that, for $L\ge3$ and small random initialization, gradient flow on $\frac12\norm{\dlS-W}_F^2$ with distinct singular values visits neighborhoods of the saddles $W_1,W_2,\dots$ in order, with explicit times. For two layers, \textcite{LiLuoLyu2021} and \textcite{GidelEtAl2019} establish versions of this; \textcite{JacotEtAl2021} conjecture it in depth and prove a theorem for the passage between the first two saddles. A complete proof for all saddles, with the prefactor of \cref{fig:deeplinear:depth}, is open.\label{ex:deeplinear:research-saddle}
\end{exercises}

\notesandrefs
Deep linear networks entered the theory of learning with \textcite{BaldiHornik1989}, who showed that a two-layer linear autoencoder or regressor has no spurious local minima and that its critical points are projections onto sets of principal components. \Textcite{Fukumizu1998} analyzed batch gradient descent in linear networks with one hidden layer, through a matrix Riccati equation, and used it to explain overtraining when the network has more hidden units than the target needs. The exact solution of \cref{thm:deeplinear:exact}, the logistic mode curves, the depth-$L$ mode equation, and the observation that linear networks reproduce the plateaus and sudden drops of nonlinear ones are from \textcite{SaxeEtAl2014}, who also found that special initializations, and random orthogonal ones, give learning times that do not grow with depth.\sidenote{The 2014 paper also connected depth-independent learning to the faithful propagation of gradients through deep nonlinear networks at the edge of chaos, the subject of \cref{ch:signalprop}.} \Textcite{SaxeEtAl2019} applied the solution to semantic development, building on the networks of \textcite{RogersMcClelland2004}.

The conservation law was found independently by \textcite{AroraCohenHazan2018}, for linear networks, and by \textcite{DuHuLee2018}, who proved that gradient flow preserves the differences of squared layer norms in homogeneous networks with linear, ReLU and leaky ReLU activations, the per-neuron statement of \cref{ex:deeplinear:relu}. \Cref{thm:deeplinear:e2e} is from \textcite{AroraCohenHazan2018}, who used it to argue that depth can accelerate optimization; \textcite{AroraEtAl2019conv} turned it into a convergence proof for gradient descent. \Textcite{Kawaguchi2016} extended the landscape results to deep linear networks.

Incremental learning and implicit low-rank bias were studied by \textcite{GidelEtAl2019}, \textcite{GissinEtAl2020}, \textcite{AroraEtAl2019}, \textcite{RazinCohen2020} and \textcite{LiLuoLyu2021}, and saddle-to-saddle dynamics by \textcite{JacotEtAl2021}. \Textcite{AtanasovEtAl2022} named the silent alignment of \cref{sec:deeplinear:random} and used it to show that networks trained from small initialization behave like kernel machines with a learned kernel. Generalization dynamics with finite data are in \textcite{AdvaniEtAl2020} and \textcite{LampinenGanguli2019}. Exact solutions beyond the aligned case are in \textcite{TarmounEtAl2021}, \textcite{BraunEtAl2022}, \textcite{DomineEtAl2025} and \textcite{KuninEtAl2024}.

\paragraph{Translating the literature.} \Textcite{SaxeEtAl2014} write a three-layer network as $W^{32}W^{21}$: their $W^{21}$ and $W^{32}$ are our $W_1$ and $W_2$, their $\Sigma^{31}$ and $\Sigma^{11}$ are our $\dlS$ and $\Sigma_{xx}$, their modes are indexed by $\alpha$ (our $i$), with $\Sigma^{31}=\sum_\alpha s_\alpha u^\alpha v^{\alpha\top}$, and their time constant $\tau=1/\lambda$ (with $\lambda$ their learning rate) is our unit of time, so their $t/\tau$ is our $t$. \Textcite{SaxeEtAl2019} write $\Sigma^{yx}=USV^\top$ for our $\dlS$ and collect the mode strengths in a diagonal matrix $A(t)$, our $\diag(a_i(t))$. \Textcite{AroraCohenHazan2018} write $N$ for our $L$, $W_e$ for our end-to-end $W$, and $L^1(W_e)$ for our $\dlE(W)$; they state the end-to-end equation with weight decay, and index its sum in the opposite order, with $[W_eW_e^\top]^{(j-1)/N}$ on the left, which is \eqref{eq:deeplinear:e2e} with $j$ replaced by $L+1-j$. \Textcite{BaldiHornik1989} write the network as $AB$, with $B$ the first layer and $A$ the second (our $W_1$ and $W_2$), and $\Sigma_{YX}\Sigma_{XX}^{-1}\Sigma_{XY}$ for the matrix whose eigenvectors define the critical points (our $\dlS\dlS^\top$ when $\Sigma_{xx}=I$).

\printchapterbib
\end{refsection}

% Chapter 3 of The Hydrogen Atoms of Deep Learning: diagonal linear networks and the bias of small initialization.
\newcommand{\diagG}{G}
\newcommand{\diagwl}[1]{w_{\ell_{#1}}}
\newcommand{\diagwphi}{w^\Phi_\alpha}
\newcommand{\diagbar}{\bar w}
\newcommand{\diagH}{\mathrm{H}}

\casechapter{The Bias of Small Initialization}{Diagonal linear networks, from kernels to sparsity}{Diagonal linear networks}{ch:diagonal}
\begin{refsection}[chapters/03-diagonal.bib]

\begin{chaptersummary}
An overparameterized network can fit its training data in many ways, and which fit gradient descent finds depends on how large the weights are at the start: large initializations behave like a kernel method, small ones learn sparse or low-rank structure, one piece at a time. A linear predictor written as $w=u\odot u-v\odot v$ and trained by gradient flow from $u=v=\alpha\mathbf 1$ shows the whole story. We prove that the flow converges to the interpolator of smallest hyperbolic entropy $\Phi_\alpha$, a potential that is a squared Euclidean norm for large $\alpha$ and an $\ell_1$ norm for small $\alpha$; that for an orthogonal design the coordinates are learned one at a time, at times $\log(1/\alpha)/(2\abs{w_{\star,i}})$; and we sketch how, for a general design, the flow jumps from saddle to saddle. Step size, stochastic gradients, classification losses, and parametrizations that are not diagonal all change the answer.
\end{chaptersummary}

\section{The phenomenon: one network, many fits}\label{sec:diagonal:phenomenon}

A network with more parameters than training examples has, as a rule, many parameter vectors that fit the training set exactly. \Textcite{ZhangEtAl2017} made the point vivid in 2016: standard convolutional networks reach zero training error on CIFAR-10 even when the labels are replaced by random ones, and then of course predict the test labels at chance. The architecture alone therefore cannot explain why the same networks generalize on real labels; among the many interpolators, the training algorithm must be choosing a good one. \Textcite{NeyshaburTomiokaSrebro2015} had argued this earlier: in their experiments the test error of two-layer networks did not deteriorate as the width grew past what was needed to fit the data, and they proposed that gradient descent regularizes implicitly, toward small norms of some kind; the selection is now called the \term{implicit bias}\index{implicit bias} of the algorithm.\sidenote{The phrase \emph{implicit regularization} is in their title. The analogy they drew was with matrix factorization, where a factorization of unbounded rank still generalizes when its norm is controlled.}

Which interpolator, then? For matrix factorization, \textcite{GunasekarEtAl2017} observed that gradient descent from an initialization close to the origin, with small steps, finds a solution of nearly minimal nuclear norm, and they conjectured that it does so exactly in the limit.\sidenote{They proved the conjecture when the measurement matrices commute, a case that reduces to the diagonal networks of this chapter. In general it fails; see \cref{sec:diagonal:wrong}.} The size of the initialization turned out to be the decisive parameter. \Textcite{ChizatOyallonBach2019} showed that a model rescaled so that its output at initialization is large trains in a \term{lazy regime}\index{lazy training}: the parameters hardly move, the model stays close to its linearization at initialization, and training is kernel regression with the tangent kernel of \cref{ch:lazy}. They reported that deep convolutional networks trained this way perform worse. \Textcite{WoodworthEtAl2020} named the two ends of the scale the \term{kernel regime}\index{kernel regime} and the \term{rich regime}\index{rich regime}, and found the transition between them in matrix factorization and in multilayer networks as the initialization scale varies. \Textcite{GeigerEtAl2020} located the crossover for networks of width $h$ at an output scale of order $h^{-1/2}$, and found that the better regime depends on the architecture: fully connected networks generalized better in the lazy regime, convolutional networks in the feature-learning one.

A second, dynamical signature accompanies small initializations. \Textcite{GissinShalevShwartzDaniely2020} observed \term{incremental learning}\index{incremental learning} in deep linear models, quadratic networks and convolutional linear networks: started near zero, the model acquires its components one after another, in order of decreasing strength, with long plateaus between them. They noted that in shallow models the effect needs an exponentially small initialization, and in deeper ones only a polynomially small one.

\begin{observation}[\textcite{ChizatOyallonBach2019}; \textcite{WoodworthEtAl2020}; \textcite{GeigerEtAl2020}]\label[observation]{obs:diagonal:regimes}
In overparameterized models trained by gradient methods to near-zero training loss, the interpolator reached depends on the scale of the initialization. Large scales give the predictor of the tangent-kernel model at initialization. Small scales give predictors with structure that the kernel does not favor (sparse vectors, low-rank matrices), acquired incrementally, with plateaus whose length grows as the scale shrinks. Which regime generalizes better depends on the task and the architecture: for convolutional networks on image data the small-scale regime has been better, for fully connected networks the comparison has gone the other way.
\end{observation}

The two regimes, and a crossover between them, have been seen in every family of models where they were looked for; but the crossover scale depends on width, depth, learning rate and loss, and in practice initialization is set by conventions (\cref{ch:signalprop}) rather than tuned. The question for this chapter is: \emph{what exactly does gradient descent select among the interpolators, as a function of the initialization scale, and how does it get there?} We want the answer in a model where it is a theorem.

\section{The model: a linear predictor written as a difference of squares}\label{sec:diagonal:model}

\paragraph{Data and loss.} We are given a design matrix $X\in\R^{n\times d}$ with rows $x_1,\dots,x_n$ and labels $y\in\R^n$, and we fit linear predictors $x\mapsto\ip{w,x}$ by the training loss
\begin{equation}\label{eq:diagonal:loss}
  \Lhat(w)=\frac1{2n}\norm{Xw-y}^2,\qquad \nabla\Lhat(w)=\frac1nX^\top(Xw-y).
\end{equation}
The case of interest is $n<d$ with $X$ of rank $n$, when the \term{interpolators}, the solutions of $Xw=y$, form an affine subspace of dimension $d-n$. Two of them will matter: the \emph{minimum-$\ell_2$-norm interpolator} $\diagwl2=X^+y$, and a \emph{minimum-$\ell_1$-norm interpolator} $\diagwl1\in\argmin\{\norm w_1:Xw=y\}$, the \term{basis pursuit}\index{basis pursuit} solution of \textcite{ChenDonohoSaunders1998}.

\paragraph{The network.} We write the weight vector as
\begin{equation}\label{eq:diagonal:param}
  w=u\odot u-v\odot v,\qquad u,v\in\R^d,
\end{equation}
with $\odot$ the coordinatewise product, so that $w_i=u_i^2-v_i^2$, and we run gradient flow on $F(u,v)=\Lhat(u\odot u-v\odot v)$ from the balanced initialization $u(0)=v(0)=\alpha\mathbf 1_d$, with initialization scale $\alpha>0$. Then $w(0)=0$. With $g(t)=\nabla\Lhat(w(t))$, the chain rule gives $\partial F/\partial u_i=2u_ig_i$ and $\partial F/\partial v_i=-2v_ig_i$, so the flow is
\begin{equation}\label{eq:diagonal:flow}
  \dot u=-2\,u\odot g,\qquad \dot v=2\,v\odot g .
\end{equation}
The model is a \term{diagonal linear network}\index{diagonal linear network}: a two-layer linear network whose hidden layer has $2d$ units, unit $i$ seeing only coordinate $i$ of the input with weight $u_i$ (or $v_i$) and sending it to the output with the same weight, $+u_i$ (or $-v_i$). The second factor of $u_i$ is the output weight; tying it to the input weight is the balanced initialization of \cref{ch:deeplinear}, preserved here by the flow.

\paragraph{What was thrown away, and why.} Nearly everything. The function class is the class of linear predictors, the same as for the plain linear model; only the \emph{parametrization} changes, and with it the geometry in which gradient flow moves. That is the point: in a linear model the set of interpolators is known exactly, so the question ``which interpolator?'' can be asked cleanly, and every answer must come from the parametrization and the initialization, since nothing else differs from ordinary least squares, whose gradient flow from $w(0)=0$ converges to $\diagwl2$. The parametrization is the simplest one that is not linear in the parameters and keeps the coordinates separate; the difference of squares lets $w$ take both signs, and the product $u\odot v$ would serve as well (\cref{ex:diagonal:product}). Infinitesimal steps, no noise, and the square loss are restrictions that the literature lifts, each changing the answer (\cref{sec:diagonal:wrong}).

\begin{dictionary}
weight vector $w=u\odot u-v\odot v$ & the function computed by a network; linear in the input, quadratic in the parameters\\
pair $(u_i,v_i)$ of a coordinate & a neuron, with tied input and output weights\\
initialization scale $\alpha$ & standard deviation of the initial weights, or the output multiplier of \textcite{ChizatOyallonBach2019}\\
conserved products $u_iv_i=\alpha^2$ & the balancedness invariants of homogeneous networks (\cref{ch:deeplinear})\\
potential $\Phi_\alpha$, limit $w^\Phi_\alpha$ & the implicit regularizer of training; the interpolator it selects\\
large $\alpha$: minimum-$\ell_2$ interpolator & the lazy or kernel regime: regression with the tangent kernel at initialization\\
small $\alpha$: minimum-$\ell_1$ interpolator & the rich regime: few active features, learned from the data\\
accelerated time $t/\log(1/\alpha)$; saddles & plateaus and sudden drops of the loss; incremental learning\\
\end{dictionary}

\paragraph{The kernel at initialization.} The gradient of the predictor $\ip{u\odot u-v\odot v,x}$ with respect to $(u,v)$ at initialization is $(2\alpha x,-2\alpha x)$, so the tangent kernel of \cref{ch:lazy} is $K(x,x')=8\alpha^2\ip{x,x'}$, the linear kernel. Kernel regression with it, run to interpolation from zero, gives $\diagwl2$. A lazy diagonal network is therefore an ordinary least-squares solver, and anything else it does is feature learning.

\paragraph{What one had to notice.} Three things. First, that a reparameterization which leaves the model class unchanged changes the implicit bias; this was clear for mirror descent, where the choice of potential decides the limit \parencite{GunasekarEtAl2018}, and the observation of \textcite{WoodworthEtAl2020} was that a reparameterization can \emph{be} a mirror descent in disguise. Second, that the coordinates of a diagonal network share nothing but the gradient, so each one can be integrated in closed form given the integrated gradient. Third, that this closed form makes the initialization scale appear inside the potential, which interpolates between the two norms of sparse regression.

\section{The solution: gradient flow is a mirror flow}\label{sec:diagonal:mirror}

\paragraph{Integrating one coordinate.} Write $\diagG(t)=\int_0^tg(s)\,\mathrm ds\in\R^d$ for the integrated gradient. Each equation in \eqref{eq:diagonal:flow} is linear in its own variable once $g$ is known.

\begin{proposition}[Closed form of the parameters]\label[proposition]{prop:diagonal:closed}
On every interval $[0,T)$ on which the flow \eqref{eq:diagonal:flow} exists,
\begin{equation}\label{eq:diagonal:closed}
  \begin{gathered}
  u=\alpha e^{-2\diagG},\qquad v=\alpha e^{2\diagG},\qquad w=-2\alpha^2\sinh(4\diagG),\\
  u_iv_i=\alpha^2,\qquad u_i^2+v_i^2=\sqrt{w_i^2+4\alpha^4},
  \end{gathered}
\end{equation}
with the exponential and the hyperbolic sine taken coordinatewise.
\end{proposition}

\begin{proof}
The scalar equation $\dot u_i=-2g_i(t)u_i$ with $u_i(0)=\alpha$ has the unique solution $u_i=\alpha e^{-2\diagG_i}$, and in the same way $v_i=\alpha e^{2\diagG_i}$. Then $u_iv_i=\alpha^2$, and $w_i=\alpha^2(e^{-4\diagG_i}-e^{4\diagG_i})=-2\alpha^2\sinh(4\diagG_i)$. Finally $(u_i^2+v_i^2)^2=(u_i^2-v_i^2)^2+4u_i^2v_i^2=w_i^2+4\alpha^4$, and $u_i^2+v_i^2>0$.
\end{proof}

Inverting the hyperbolic sine, $\frac14\arcsinh\bigl(w_i/(2\alpha^2)\bigr)=-\diagG_i$. Differentiating in $t$ gives the central identity of the chapter,
\begin{equation}\label{eq:diagonal:mirrorflow}
  \frac{\mathrm d}{\mathrm dt}\nabla\Phi_\alpha\bigl(w(t)\bigr)=-\nabla\Lhat\bigl(w(t)\bigr),\qquad \nabla\Phi_\alpha(w)_i=\frac14\arcsinh\frac{w_i}{2\alpha^2},
\end{equation}
for the potential
\begin{equation}\label{eq:diagonal:potential}
  \Phi_\alpha(w)=\frac{\alpha^2}4\sum_{i=1}^dq\Bigl(\frac{w_i}{\alpha^2}\Bigr),\qquad q(z)=2-\sqrt{4+z^2}+z\arcsinh\frac z2 .
\end{equation}
Indeed $q'(z)=-z/\sqrt{4+z^2}+\arcsinh(z/2)+z/\sqrt{4+z^2}=\arcsinh(z/2)$, so $\partial_i\Phi_\alpha(w)=\frac14q'(w_i/\alpha^2)=\frac14\arcsinh(w_i/2\alpha^2)$. Equation \eqref{eq:diagonal:mirrorflow} says that $w(t)$ follows \term{mirror descent}\index{mirror descent} in continuous time, with the mirror map $\nabla\Phi_\alpha$: the flow is a gradient flow in the dual coordinates $\nabla\Phi_\alpha(w)$, not in $w$.\sidenote{Mirror descent is due to Nemirovsky and Yudin; see \textcite{NemirovskyYudin1983}. With the entropy as potential it becomes the exponentiated gradient method of \textcite{KivinenWarmuth1997}.} In the parameters the flow is an ordinary gradient flow; in the predictor it is a mirror flow, and the reparameterization has chosen the mirror. The function $\Phi_\alpha$ is, up to a factor and a constant, the \term{hyperbolic entropy}\index{hyperbolic entropy} (``hypentropy'') of \textcite{GhaiHazanSinger2020}: with $\beta=2\alpha^2$, $\frac{\alpha^2}4q(w_i/\alpha^2)=\frac14\bigl(w_i\arcsinh(w_i/\beta)-\sqrt{w_i^2+\beta^2}+\beta\bigr)$.

\begin{lemma}[The potential]\label[lemma]{lem:diagonal:q}
The function $q$ is even, smooth and strictly convex, with $q(0)=q'(0)=0$, $q'(z)=\arcsinh(z/2)$ and $q''(z)=1/\sqrt{4+z^2}$. For every real $z$,
\begin{equation}\label{eq:diagonal:qquad}
  \frac{z^2}{2\sqrt{4+z^2}}\le q(z)\le\frac{z^2}4,
\end{equation}
and
\begin{equation}\label{eq:diagonal:qlog}
  0\le q(z)-\abs z\bigl(\log\abs z-1\bigr)<2,
\end{equation}
with the convention $0\log0=0$.
\end{lemma}

\begin{proof}
The formulas for $q'$ and $q''$ were checked above, and $q''>0$; evenness and $q(0)=0$ are read off from \eqref{eq:diagonal:potential}. For \eqref{eq:diagonal:qquad}, Taylor's formula with integral remainder gives $q(z)=\int_0^z(z-s)q''(s)\,\mathrm ds$. For $z>0$ the weight $q''(s)$ decreases from $q''(0)=\frac12$ to $q''(z)$ on $[0,z]$, so $\frac12q''(z)z^2\le q(z)\le\frac14z^2$; evenness covers $z<0$.

For \eqref{eq:diagonal:qlog}, let $r(z)=q(z)-\abs z(\log\abs z-1)$, an even function with $r(0)=0$, continuous at $0$ because $z\log z\to0$. For $z>0$,
\[
  r'(z)=\arcsinh\frac z2-\log z=\log\Bigl(\frac12+\sqrt{\frac14+\frac1{z^2}}\Bigr)>\log1=0,
\]
using $\arcsinh x=\log(x+\sqrt{1+x^2})$. So $r$ increases on $(0,\infty)$ and $r\ge0$. As $z\to\infty$, $\sqrt{4+z^2}-z\to0$ and, since $\sqrt{\frac14+\frac1{z^2}}\le\frac12+\frac1{z^2}$, $0<z\,r'(z)\le z\log(1+z^{-2})\le1/z\to0$; since
\[
  r(z)=2-\bigl(\sqrt{4+z^2}-z\bigr)+z\,r'(z),
\]
$r(z)\to2$, and an increasing function stays below its limit.
\end{proof}

The two bounds describe the two faces of $\Phi_\alpha$. By \eqref{eq:diagonal:qquad}, where $\abs{w_i}\ll\alpha^2$ the potential is $\Phi_\alpha(w)\approx\norm w^2/(16\alpha^2)$, a squared Euclidean norm. By \eqref{eq:diagonal:qlog}, writing $\eps_\alpha=\log(1/\alpha^2)$ and $\diagH(w)=\sum_i\abs{w_i}(\log\abs{w_i}-1)$,
\begin{equation}\label{eq:diagonal:l1form}
  \Phi_\alpha(w)=\frac{\eps_\alpha}4\norm w_1+\frac14\diagH(w)+R_\alpha(w),\qquad 0\le R_\alpha(w)<\frac{d\alpha^2}2,
\end{equation}
for every $w\in\R^d$: substitute $z=w_i/\alpha^2$ in \eqref{eq:diagonal:qlog}, multiply by $\alpha^2/4$, and note $\frac{\alpha^2}4\abs z(\log\abs z-1)=\frac14\abs{w_i}(\log\abs{w_i}+\eps_\alpha-1)$. When $\alpha$ is small, $\eps_\alpha$ is large and $\Phi_\alpha$ is an $\ell_1$ norm with an entropy-like correction.

\paragraph{The main theorem.} We now prove that gradient flow converges, and to what.

\begin{theorem}[Implicit bias of diagonal linear networks]\label{thm:diagonal:bias}
Let $X\in\R^{n\times d}$ and $y\in\R^n$ be such that some $w^\circ\in\R^d$ satisfies $Xw^\circ=y$, and let $\alpha>0$. Then:
\begin{enumerate}
\item The flow \eqref{eq:diagonal:flow} from $u(0)=v(0)=\alpha\mathbf 1_d$ exists for all $t\ge0$, and $u(t)$ and $v(t)$ have positive coordinates and stay bounded.
\item For every interpolator $w^\circ$ and every $t>0$, $\Lhat(w(t))\le\Phi_\alpha(w^\circ)/(2t)$.
\item The problem of minimizing $\Phi_\alpha(w)$ subject to $Xw=y$ has a unique solution $\diagwphi$, and $w(t)\to\diagwphi$ as $t\to\infty$.
\end{enumerate}
\end{theorem}

\begin{proof}
The plan is to use the Bregman divergence of $\Phi_\alpha$ to an interpolator as a Lyapunov function, which gives boundedness, the rate, and existence for all time at once; then to identify the limit by the optimality conditions, using that $\nabla\Phi_\alpha(w(t))$ never leaves the row space of $X$.

\emph{A Lyapunov function.} Let $[0,T)$ be the maximal interval of existence of the flow; the right side of \eqref{eq:diagonal:flow} is a polynomial in $(u,v)$, so a unique maximal solution exists, and if $T<\infty$, then $(u,v)$ is unbounded on $[0,T)$. Fix an interpolator $w^\circ$ and let
\[
  D(t)=\Phi_\alpha(w^\circ)-\Phi_\alpha(w(t))-\ip{\nabla\Phi_\alpha(w(t)),w^\circ-w(t)},
\]
the Bregman divergence\index{Bregman divergence} of $\Phi_\alpha$ from $w(t)$ to $w^\circ$. It is nonnegative because $\Phi_\alpha$ is convex (\cref{lem:diagonal:q}). Differentiating, the terms $\mp\ip{\nabla\Phi_\alpha(w),\dot w}$ cancel, and \eqref{eq:diagonal:mirrorflow} gives
\[
  \dot D=-\Bigl\langle\frac{\mathrm d}{\mathrm dt}\nabla\Phi_\alpha(w),\,w^\circ-w\Bigr\rangle=\ip{\nabla\Lhat(w),w^\circ-w}=-2\Lhat(w),
\]
where the last step is $\ip{\nabla\Lhat(w),w^\circ-w}=\frac1n\ip{Xw-y,\,Xw^\circ-Xw}$ with $Xw^\circ=y$. So $D$ is nonincreasing, and $D(0)=\Phi_\alpha(w^\circ)$ because $w(0)=0$, $\Phi_\alpha(0)=0$ and $\nabla\Phi_\alpha(0)=0$.

\emph{Boundedness and existence.} We show that the sublevel sets of $D$, as a function of $w$, are bounded. Write $\varphi(x)=\frac{\alpha^2}4q(x/\alpha^2)$ for one coordinate of $\Phi_\alpha$, so that $D=\sum_iD_i$ with $D_i=\varphi(w_i^\circ)-\varphi(w_i)-\varphi'(w_i)(w_i^\circ-w_i)\ge0$. A computation from \eqref{eq:diagonal:potential} gives $x\varphi'(x)-\varphi(x)=\frac14\bigl(\sqrt{x^2+4\alpha^4}-2\alpha^2\bigr)\ge\frac14(\abs x-2\alpha^2)$, and $\abs{\varphi'(x)}=\frac14\arcsinh(\abs x/2\alpha^2)\le\frac14\log(1+\abs x/\alpha^2)$, because $\arcsinh s\le\log(1+2s)$ for $s\ge0$. Hence
\[
  4D_i\ge4\varphi(w_i^\circ)+\abs{w_i}-2\alpha^2-\abs{w^\circ_i}\log\bigl(1+\abs{w_i}/\alpha^2\bigr),
\]
which tends to infinity with $\abs{w_i}$. Since $D(t)\le\Phi_\alpha(w^\circ)$ and every $D_j\ge0$, each $w_i(t)$ stays in a bounded set, uniformly in $t<T$. By \cref{prop:diagonal:closed}, $u_i^2+v_i^2=\sqrt{w_i^2+4\alpha^4}$, so $(u,v)$ is bounded on $[0,T)$, and therefore $T=\infty$. Positivity is \eqref{eq:diagonal:closed}. This proves (i).

\emph{The rate.} Along the flow $\frac{\mathrm d}{\mathrm dt}\Lhat(w(t))=\frac{\mathrm d}{\mathrm dt}F(u,v)=-\norm{\nabla F(u,v)}^2\le0$, so $\Lhat(w(t))$ is nonincreasing, and integrating $\dot D=-2\Lhat$,
\[
  2t\,\Lhat(w(t))\le2\int_0^t\Lhat(w(s))\,\mathrm ds=D(0)-D(t)\le\Phi_\alpha(w^\circ),
\]
which is (ii).

\emph{The minimizer.} By \eqref{eq:diagonal:l1form}, $\Phi_\alpha(w)\ge\frac14\sum_i\abs{w_i}(\log(\abs{w_i}/\alpha^2)-1)$, which tends to infinity with $\norm w$; so $\Phi_\alpha$ attains its minimum on the closed nonempty set $\{Xw=y\}$, and the minimizer is unique because $\Phi_\alpha$ is strictly convex. A point $\bar w$ with $X\bar w=y$ is this minimizer if and only if $\nabla\Phi_\alpha(\bar w)$ lies in the row space $\operatorname{range}(X^\top)$. Indeed, if $\nabla\Phi_\alpha(\bar w)=X^\top\nu$, then for every interpolator $w'$ convexity gives $\Phi_\alpha(w')\ge\Phi_\alpha(\bar w)+\ip{X^\top\nu,w'-\bar w}=\Phi_\alpha(\bar w)$, since $X(w'-\bar w)=0$; and conversely the gradient at a minimizer is orthogonal to $\ker X$, whose orthogonal complement is $\operatorname{range}(X^\top)$.

\emph{The limit.} By \eqref{eq:diagonal:mirrorflow} and $\nabla\Phi_\alpha(w(0))=0$,
for every $t$,
\[
  \nabla\Phi_\alpha(w(t))=-\int_0^t\nabla\Lhat(w(s))\,\mathrm ds=-\frac1nX^\top\int_0^t\bigl(Xw(s)-y\bigr)\,\mathrm ds\in\operatorname{range}(X^\top).
\]
Let $\bar w$ be a limit point of $w(t)$ along a sequence $t_k\to\infty$; one exists because $w$ is bounded. By (ii), $\Lhat(\bar w)=\lim_k\Lhat(w(t_k))=0$, so $X\bar w=y$. The subspace $\operatorname{range}(X^\top)$ is closed and $\nabla\Phi_\alpha$ is continuous, so $\nabla\Phi_\alpha(\bar w)\in\operatorname{range}(X^\top)$, and $\bar w=\diagwphi$ by the optimality condition. A bounded trajectory whose only limit point is $\diagwphi$ converges to it, which proves (iii).
\end{proof}

Two remarks on the hypotheses. The theorem does not need $n<d$: if the interpolator is unique, then $\diagwphi$ is that interpolator. And it needs an interpolator; when $y\notin\operatorname{range}(X)$ the same argument, with $w^\circ$ a least-squares solution, shows that the flow converges to the least-squares solution of smallest $\Phi_\alpha$ (\cref{ex:diagonal:noisy}). The rate (ii) is slow for small $\alpha$, since $\Phi_\alpha(w^\circ)\approx\frac14\log(1/\alpha^2)\norm{w^\circ}_1$: the time to fit grows like $\log(1/\alpha)$, the length of the plateaus we meet in \cref{sec:diagonal:dynamics}.

\section{The solution: from kernel to sparsity}\label{sec:diagonal:limits}

Which interpolator is $\diagwphi$? The two faces of $\Phi_\alpha$ suggest the answer at the two ends of the scale.

\begin{theorem}[The two limits]\label{thm:diagonal:limits}
Let $X$ and $y$ be as in \cref{thm:diagonal:bias}, and let $\diagwl2=X^+y$ be the minimum-$\ell_2$-norm interpolator.
\begin{enumerate}
\item \emph{Kernel regime.} If $\alpha^4\ge\norm{\diagwl2}^2$, then $\displaystyle\norm{\diagwphi-\diagwl2}\le\frac{\norm{\diagwl2}^2}{2^{5/4}\alpha^2}$. In particular $\diagwphi\to\diagwl2$ as $\alpha\to\infty$.
\item \emph{Rich regime.} Let $m_1=\min\{\norm w_1:Xw=y\}$ and let $\diagwl1$ attain it. For $0<\alpha<1$,
\[
  m_1\le\norm{\diagwphi}_1\le m_1+\frac{\diagH(\diagwl1)+d+2d\alpha^2}{\log(1/\alpha^2)} .
\]
Every limit point of $\diagwphi$ as $\alpha\to0$ is a minimum-$\ell_1$-norm interpolator, and if the minimum-$\ell_1$-norm interpolator is unique, then $\diagwphi\to\diagwl1$.
\end{enumerate}
\end{theorem}

\begin{proof}
\emph{Part (i).} The plan is to sandwich $16\alpha^2\Phi_\alpha$ between the squared Euclidean norm and a slightly smaller multiple of it, and to use that the excess squared norm of an interpolator over $\diagwl2$ is its squared distance to $\diagwl2$. Let $\Psi(w)=16\alpha^2\Phi_\alpha(w)=4\alpha^4\sum_iq(w_i/\alpha^2)$, which has the same minimizer $\diagwphi$ on $\{Xw=y\}$. By \eqref{eq:diagonal:qquad}, with $z=w_i/\alpha^2$,
\[
  \frac{w_i^2}{\sqrt{1+w_i^2/(4\alpha^4)}}=\frac{4\alpha^4z^2}{2\sqrt{4+z^2}}\le4\alpha^4q(z)\le w_i^2,
\]
and summing, $\norm w^2/\sqrt{1+\norm w^2/(4\alpha^4)}\le\Psi(w)\le\norm w^2$, where we used $w_i^2\le\norm w^2$ in the denominators. Let $\rho=\norm{\diagwphi}$ and $m=\norm{\diagwl2}$. Minimality gives $\Psi(\diagwphi)\le\Psi(\diagwl2)\le m^2$, so
\begin{equation}\label{eq:diagonal:sandwich}
  \frac{\rho^2}{\sqrt{1+\rho^2/(4\alpha^4)}}\le m^2 .
\end{equation}
If $\rho^2>4\alpha^4$, then the left side of \eqref{eq:diagonal:sandwich} exceeds $\rho^2/\sqrt{2\rho^2/(4\alpha^4)}=\sqrt2\,\alpha^2\rho>2\sqrt2\,\alpha^4\ge2\sqrt2\,m^2$, using the hypothesis $\alpha^4\ge m^2$; this contradicts \eqref{eq:diagonal:sandwich}. So $\rho^2\le4\alpha^4$, the left side is at least $\rho^2/\sqrt2$, and $\rho^2\le\sqrt2\,m^2$. Feeding this back into \eqref{eq:diagonal:sandwich} and using $\sqrt{1+x}\le1+x/2$,
\[
  \rho^2\le m^2\sqrt{1+\frac{\sqrt2\,m^2}{4\alpha^4}}\le m^2+\frac{\sqrt2\,m^4}{8\alpha^4} .
\]
Finally $\diagwl2$ lies in $\operatorname{range}(X^\top)=(\ker X)^\perp$ and $\diagwphi-\diagwl2\in\ker X$, so $\rho^2=m^2+\norm{\diagwphi-\diagwl2}^2$ by Pythagoras, and $\norm{\diagwphi-\diagwl2}^2\le\sqrt2\,m^4/(8\alpha^4)$. Taking square roots, $\sqrt{\sqrt2/8}=2^{-5/4}$.

\emph{Part (ii).} The plan is to compare $\Phi_\alpha$ at $\diagwphi$ and at $\diagwl1$ through the expansion \eqref{eq:diagonal:l1form}. The lower bound $m_1\le\norm{\diagwphi}_1$ holds because $\diagwphi$ interpolates. For $0<\alpha<1$ we have $\eps_\alpha=\log(1/\alpha^2)>0$. Minimality of $\diagwphi$ and \eqref{eq:diagonal:l1form} at both points give
\[
  \frac{\eps_\alpha}4\norm{\diagwphi}_1+\frac14\diagH(\diagwphi)\le\Phi_\alpha(\diagwphi)\le\Phi_\alpha(\diagwl1)<\frac{\eps_\alpha}4m_1+\frac14\diagH(\diagwl1)+\frac{d\alpha^2}2 .
\]
The function $x\mapsto x\log x-x$ is at least $-1$ on $[0,\infty)$, with its minimum at $x=1$, so $\diagH(\diagwphi)\ge-d$. Multiplying by $4/\eps_\alpha$ and rearranging gives the upper bound. For the limits, let $\alpha_k\to0$ be a sequence along which $w^\Phi_{\alpha_k}$ converges, to $\bar w$ say; such sequences exist because the bound keeps $\norm{\diagwphi}_1$ bounded for $\alpha\le\frac12$. Then $X\bar w=y$, and the bound gives $\norm{\bar w}_1\le m_1$, so $\bar w$ is a minimum-$\ell_1$-norm interpolator. If that interpolator is unique, then every convergent subsequence has the same limit, and the bounded family $\diagwphi$ converges to it.
\end{proof}

The kernel limit is no surprise once \eqref{eq:diagonal:qquad} is known: a large initialization makes every coordinate move little relative to $\alpha^2$, where the potential is quadratic. The rate in (i) is not sharp; since $q(z)=\frac{z^2}4-\frac{z^4}{192}+O(z^6)$, the potential differs from the quadratic by a relative amount of order $\alpha^{-4}$, and $\norm{\diagwphi-\diagwl2}$ is in fact of order $\alpha^{-4}$ (\cref{ex:diagonal:taylor}). The rich limit is the surprise: no $\ell_1$ penalty appears anywhere in the training procedure, yet small initialization produces basis pursuit. The bound in (ii) is slow, through $1/\log(1/\alpha)$, but it concerns norms; the vector itself converges much faster when the $\ell_1$ problem is well conditioned (\cref{sec:diagonal:reading}). Ties between $\ell_1$ minimizers are broken by the entropy term of \eqref{eq:diagonal:l1form} (\cref{ex:diagonal:tie}).

\paragraph{Why this matters for generalization.} Suppose $y=Xw_\star$ for a $k$-sparse $w_\star$ and $X$ has independent standard Gaussian entries. The minimum-$\ell_2$ interpolator is the projection of $w_\star$ on the row space, a uniformly random $n$-dimensional subspace, so $\E\norm{\diagwl2-w_\star}^2=(1-n/d)\norm{w_\star}^2$, the overparameterized branch of \cref{ch:doubledescent} without noise: the kernel regime learns nothing until $n$ approaches $d$. Basis pursuit, in contrast, recovers $w_\star$ exactly with high probability as soon as $n$ exceeds the statistical dimension of the descent cone of the $\ell_1$ norm at $w_\star$, which is of order $k\log(d/k)$, and fails with high probability below it; the transition is sharp, with width of order $\sqrt d$ \parencite{AmelunxenEtAl2014}.\sidenote{For $d=100$ and $k=5$ the statistical dimension is $20.4$. Exact recovery by $\ell_1$ minimization is the subject of compressed sensing, where dual certificates are the standard tool.} Between the two, the initialization scale interpolates between a sample complexity of order $d$ and one of order $k\log(d/k)$.

\section{The solution: incremental learning and saddles}\label{sec:diagonal:dynamics}

\Cref{thm:diagonal:bias} says where the flow ends, not how it gets there. When the coordinates decouple, the path can be computed exactly.

\begin{assumption}[Orthogonal design]\label[assumption]{ass:diagonal:orth}
The design satisfies $X^\top X=nI_d$ (so $n\ge d$). We write $\diagbar=X^\top y/n$ for the least-squares solution.
\end{assumption}

Under \cref{ass:diagonal:orth}, $\nabla\Lhat(w)=w-\diagbar$ and $\Lhat(w)=\frac12\norm{w-\diagbar}^2+c_0$ with $c_0=\frac1{2n}\norm y^2-\frac12\norm{\diagbar}^2\ge0$; for noiseless labels $y=Xw_\star$, $\diagbar=w_\star$ and $c_0=0$. The design is not overparameterized, so there is no implicit bias to speak of; what remains is the dynamics. By \cref{prop:diagonal:closed}, each coordinate obeys an autonomous equation,
\begin{equation}\label{eq:diagonal:coord}
  \dot w_i=2u_i\dot u_i-2v_i\dot v_i=-4\sqrt{w_i^2+4\alpha^4}\,(w_i-\diagbar_i),
\end{equation}
a gradient flow on $\frac12(w_i-\diagbar_i)^2$ with the state-dependent rate $4\sqrt{w_i^2+4\alpha^4}$. Near $w_i=0$ the rate is $8\alpha^2$, tiny: this is why small initializations wait. Once $\abs{w_i}\gg\alpha^2$ the rate is $4\abs{w_i}$, and the coordinate grows exponentially, like the modes of \cref{ch:deeplinear}.

\begin{theorem}[Incremental learning under an orthogonal design]\label{thm:diagonal:orth}
Assume \cref{ass:diagonal:orth} and let $s=\diagbar_i$ for one coordinate $i$.
\begin{enumerate}
\item If $s=0$, then $w_i\equiv0$. If $s>0$, let $\rho=\sqrt{s^2+4\alpha^4}$, $p_\pm=(s\pm\rho)/2$, and $\kappa_0=(p_+-\alpha^2)/(\alpha^2-p_-)$; then
\begin{equation}\label{eq:diagonal:riccati}
  w_i(t)=P(t)-\frac{\alpha^4}{P(t)},\qquad P(t)=\frac{p_++\kappa_0e^{-4\rho t}p_-}{1+\kappa_0e^{-4\rho t}},
\end{equation}
and $w_i(t)$ increases from $0$ to $s$. If $s<0$, then $w_i(t)$ is minus the solution for $-s$.
\item If $s\ne0$, then for every $\tau\ge0$ with $\tau\ne1/(2\abs s)$,
\[
  \lim_{\alpha\to0}w_i\bigl(\tau\log(1/\alpha)\bigr)=\begin{cases}0,&\tau<1/(2\abs s),\\ s,&\tau>1/(2\abs s).\end{cases}
\]
\item Consequently, for every $\tau\ge0$ different from all the times $1/(2\abs{\diagbar_j})$,
\[
  \lim_{\alpha\to0}\Lhat\bigl(w(\tau\log(1/\alpha))\bigr)=c_0+\frac12\sum_{j:\,2\abs{\diagbar_j}\tau<1}\diagbar_j^2 .
\]
\end{enumerate}
\end{theorem}

\begin{proof}
The plan is to reduce \eqref{eq:diagonal:flow} for one coordinate to a Riccati equation for $P=u_i^2$ with constant coefficients and to solve it.

\emph{Part (i).} By \eqref{eq:diagonal:flow} with $g_i=w_i-s$ and \cref{prop:diagonal:closed}, $w_i=P-\alpha^4/P$ and
\[
  \dot P=2u_i\dot u_i=-4P\Bigl(P-\frac{\alpha^4}P-s\Bigr)=-4\bigl(P^2-sP-\alpha^4\bigr)=-4(P-p_+)(P-p_-),
\]
because $p_\pm$ are the roots of $P^2-sP-\alpha^4$ (their sum is $s$ and their product $(s^2-\rho^2)/4=-\alpha^4$). If $s=0$, then $P\equiv\alpha^2$ solves this equation with the right initial value, so $w_i\equiv0$. Let $s>0$. Then $p_-<0<\alpha^2<p_+$, the middle inequality because $P^2-sP-\alpha^4$ equals $-s\alpha^2<0$ at $P=\alpha^2$. The constants $p_\pm$ are solutions, so by uniqueness $P(t)$ stays in $(p_-,p_+)$, where $\dot P>0$. On this interval
\[
  \frac{\mathrm d}{\mathrm dt}\log\frac{p_+-P}{P-p_-}=-\dot P\,\frac{p_+-p_-}{(p_+-P)(P-p_-)}=-4(p_+-p_-)=-4\rho,
\]
so $(p_+-P)/(P-p_-)=\kappa_0e^{-4\rho t}$, and solving for $P$ gives \eqref{eq:diagonal:riccati}. Since $P$ increases to $p_+$ and $x\mapsto x-\alpha^4/x$ increases on $(0,\infty)$, $w_i$ increases from $0$ to $p_+-\alpha^4/p_+=p_++p_-=s$. For $s<0$, the map $(u_i,v_i,s)\mapsto(v_i,u_i,-s)$ preserves \eqref{eq:diagonal:flow} for coordinate $i$ and the initial condition, and changes the sign of $w_i$.

\emph{Part (ii).} By (i) we may take $s>0$. Put $t=\tau\log(1/\alpha)$ and $\kappa=\kappa_0e^{-4\rho t}$. As $\alpha\to0$, $p_+\to s$, $p_-=-\alpha^4/p_+\to0$ and $\alpha^2\kappa_0=\alpha^2(p_+-\alpha^2)/(\alpha^2-p_-)\to s$, while $\rho-s=O(\alpha^4)$, so $(\rho-s)\log(1/\alpha)\to0$. Hence
\[
  \log\kappa=\log(\alpha^2\kappa_0)+2\log(1/\alpha)-4\rho\tau\log(1/\alpha)=(2-4s\tau)\log(1/\alpha)+\log s+o(1),
\]
which tends to $-\infty$ if $\tau>1/(2s)$ and to $+\infty$ if $\tau<1/(2s)$. In the first case $P\to p_+\to s$ by \eqref{eq:diagonal:riccati}, so $w_i=P-\alpha^4/P\to s$. In the second case $P=p_+/(1+\kappa)+p_-\kappa/(1+\kappa)\to0$, and $P\ge\alpha^2$ because $P$ increases from $\alpha^2$; so $\alpha^4/P\le\alpha^2\to0$ and $w_i\to0$.

\emph{Part (iii).} Apply (ii) to every coordinate in $\Lhat(w)=c_0+\frac12\sum_j(w_j-\diagbar_j)^2$.
\end{proof}

In the \term{accelerated time}\index{accelerated time} $\tau=t/\log(1/\alpha)$ the loss converges, as $\alpha\to0$, to a staircase: coordinate $j$ is learned at $\tau_j=1/(2\abs{\diagbar_j})$, the largest first, and the loss drops by $\frac12\diagbar_j^2$ at each step. In the original time the transitions are logistic. With $t_{1/2}$ the time at which $w_i=s/2$, \eqref{eq:diagonal:riccati} gives, for $\abs{w_i}\gg\alpha^2$, $w_i/s\approx1/(1+e^{-4\abs s(t-t_{1/2})})$: a transition of width of order $1/\abs s$ after a plateau of length
\begin{equation}\label{eq:diagonal:thalf}
  t_{1/2}=\frac{\log(\abs s/\alpha^2)}{4\abs s}-\frac{\alpha^2}{2s^2}+O(\alpha^4)
\end{equation}
(\cref{ex:diagonal:half}), which is the time at which the exponential growth $\alpha^2e^{4\abs st}$ of the linearized equation reaches $\abs s$. This is the logistic escape of \cref{ch:deeplinear},\sidenote{There the mode strength enters through a rate $2s$; here the rate is $4\abs s$ because the gradient enters each factor through $\dot u=-2u\odot g$, so that $u_i^2$ grows like $e^{4\abs st}$. The constants depend on how the loss is normalized, the logarithm of the initialization does not.} with the same dependence on the strength $\abs s$ of the mode and on the logarithm of the initialization; it is also the shape of the search phase of \cref{ch:singleindex} when the information exponent is $2$.

\paragraph{General designs: saddle to saddle.} When $X^\top X$ is not diagonal, coordinates interact through the gradient, and a coordinate that has been learned can be unlearned. The closed form \eqref{eq:diagonal:closed} still holds, and it suggests how to take the limit $\alpha\to0$. Put $L_\alpha=\log(1/\alpha)$ and define the rescaled dual variable
\[
  \zeta(\tau)=-\frac{2\diagG(\tau L_\alpha)}{L_\alpha}=-2\int_0^\tau\nabla\Lhat\bigl(w(\tau'L_\alpha)\bigr)\,\mathrm d\tau' .
\]
By \eqref{eq:diagonal:closed}, $w_i=2\alpha^2\sinh(2L_\alpha\zeta_i)$, so $\abs{w_i}=\alpha^{2-2\abs{\zeta_i}}(1+o(1))$ as $\alpha\to0$ when $\zeta_i\ne0$ is fixed. A coordinate whose dual variable satisfies $\abs{\zeta_i}<1$ is exponentially small in $L_\alpha$, a coordinate of order one has $\abs{\zeta_i}=1+O(1/L_\alpha)$, and $\abs{\zeta_i}>1$ is impossible for a bounded trajectory. The dual variable moves at speed $2\nabla\Lhat(w)$ in accelerated time while the predictor sits still, and a new coordinate wakes up when its $\abs{\zeta_i}$ reaches $1$. This leads to the following algorithm.\index{saddle-to-saddle dynamics}

\begin{definition}[The saddle-to-saddle path]\label[definition]{def:diagonal:path}
Start from $\tau_0=0$, $w^{(0)}=0$ and $\zeta(0)=0$. Given the saddle $w^{(k)}$ and $\zeta(\tau_k)$, let $g^{(k)}=\nabla\Lhat(w^{(k)})$; if $g^{(k)}=0$, stop. Otherwise let $\zeta(\tau)=\zeta(\tau_k)-2(\tau-\tau_k)g^{(k)}$ for $\tau\ge\tau_k$, let $\tau_{k+1}$ be the first $\tau>\tau_k$ at which some coordinate with $\abs{\zeta_i(\tau_k)}<1$ reaches $\abs{\zeta_i(\tau)}=1$, write $\zeta^+=\zeta(\tau_{k+1})$, let $A$ be the set of coordinates with $\abs{\zeta^+_i}=1$, and let
\[
  w^{(k+1)}=\argmin\bigl\{\Lhat(w):\ w_i=0\text{ for }i\notin A,\ \ \zeta^+_iw_i\ge0\text{ for }i\in A\bigr\}.
\]
The limit path is $\tilde w(\tau)=w^{(k)}$ for $\tau_k\le\tau<\tau_{k+1}$.
\end{definition}

The constraint set is a cone of sign-constrained vectors on the active coordinates, and the saddles $w^{(k)}$ are critical points of the loss restricted to coordinate subspaces, hence saddle points of $F$ when they are not global minima. The procedure resembles the \textsc{lars} algorithm of \textcite{EfronEtAl2004},\sidenote{\textsc{lars} stands for least angle regression. With a small modification it computes the whole Lasso path, one breakpoint at a time.} which computes the regularization path of the Lasso: active sets grow, sometimes lose a coordinate, and each step solves a least-squares problem on the active set with sign constraints.

\begin{theorem}[\textcite{PesmeFlammarion2023}, in our normalization]\label{thm:diagonal:saddles}
Under nondegeneracy assumptions on $(X,y)$ that hold for data in general position (see their paper for the precise conditions), the procedure of \cref{def:diagonal:path} stops after finitely many steps, at the minimum-$\ell_1$-norm interpolator when $y\in\operatorname{range}(X)$, and $w(\tau\log(1/\alpha))\to\tilde w(\tau)$ as $\alpha\to0$, uniformly on compact sets of accelerated times that avoid the jump times $\tau_k$.
\end{theorem}

\begin{proofsketch}
We describe the mechanism and omit the uniform control of the transitions, which is the technical heart of the paper of Pesme and Flammarion. \emph{Plateaus.} Between jumps, the active coordinates are of order one and the inactive ones are of size $\alpha^{2-2\abs{\zeta_i}}$, exponentially small in $L_\alpha$. On the active coordinates, the flow in the original time $t$ runs at rate of order one, so on the time scale $L_\alpha$ it has converged to a minimizer of $\Lhat$ over vectors supported on the active set with the signs of $\zeta$: the saddle $w^{(k)}$. The gradient there vanishes on the active coordinates with nonzero value and pushes $\abs{\zeta_i}$ inward on the active coordinates at zero, so $\zeta$ moves linearly as in \cref{def:diagonal:path}, and the inactive coordinates grow like $\alpha^{2-2\abs{\zeta_i(\tau)}}$. \emph{Jumps.} When an inactive $\abs{\zeta_i}$ reaches $1$, coordinate $i$ reaches order one; the flow then leaves the saddle along a heteroclinic orbit in time $O(1)$ in $t$, which is $O(1/L_\alpha)$ in $\tau$, and settles at the next constrained minimizer. \emph{End.} When the path reaches an interpolator $\tilde w$, $\zeta$ stops, with $\zeta\in\operatorname{range}(X^\top)$, $\norm\zeta_\infty\le1$, and $\zeta_i=\sgn(\tilde w_i)$ on the support of $\tilde w$: the optimality conditions of basis pursuit. Pesme and Flammarion make the transitions rigorous with an arc-length reparametrization of time, and prove that the number of jumps is finite.
\end{proofsketch}

Under \cref{ass:diagonal:orth}, $g^{(0)}=-\diagbar$, the coordinate $j$ reaches $\abs{\zeta_j}=1$ at $\tau=1/(2\abs{\diagbar_j})$, and \cref{def:diagonal:path} reproduces the staircase of \cref{thm:diagonal:orth}, a useful check of the normalization. For general designs the limit path visits saddles that are not on the way to $\diagwl1$ in a monotone sense: coordinates can enter with one sign and leave, and the loss decreases in steps whose heights depend on the correlations of the design.

\section{Reading the solution}\label{sec:diagonal:reading}

The theorems make three kinds of prediction: \emph{where} the flow ends, as a function of $\alpha$; \emph{when} each piece is learned; and \emph{in which order}.

\begin{figure}[t]
\centering
\begin{tikzpicture}
\begin{axis}[hplot,xmode=log,ymode=log,xmin=8e-5,xmax=12,ymin=1e-7,ymax=3,
  xlabel={initialization scale $\alpha$},ylabel={$\norm{w-w_\star}^2$},
  legend style={at={(0.97,0.03)},anchor=south east},unbounded coords=jump]
\addplot[guide] table[x=alpha,y=errl2] {figures/data/03-diagonal-sparse-theory.dat};
\addplot[theory2] table[x=alpha,y=guide] {figures/data/03-diagonal-sparse-theory.dat};
\addplot[theory] table[x=alpha,y=err] {figures/data/03-diagonal-sparse-theory.dat};
\addplot[simulation] table[x=alpha,y=err] {figures/data/03-diagonal-sparse-sim.dat};
\legend{minimum-$\ell_2$ interpolator,slope $4\gamma^\star$,$w^\Phi_\alpha$ (theory),gradient flow}
\end{axis}
\end{tikzpicture}
\caption{From kernel to sparse recovery. Gaussian design with $n=40$, $d=100$, a $5$-sparse target with entries $\pm1/\sqrt5$ on a random support, noiseless labels (seed $0$). Curve: the squared error of $w^\Phi_\alpha=\argmin\{\Phi_\alpha(w):Xw=y\}$, computed by Newton's method on the dual problem. Points: the gradient flow \eqref{eq:diagonal:flow}, integrated in the variables $(\log u,\log v)$ until $\Lhat<10^{-20}$; at all eleven values of $\alpha$ the flow ends within $5\times10^{-10}$ of $w^\Phi_\alpha$. Horizontal dashed line: the error $0.604$ of the minimum-$\ell_2$ interpolator. Sloped dashed line: the power law $\alpha^{4\gamma^\star}$ of \eqref{eq:diagonal:richrate}, with $\gamma^\star=0.516$ computed from the data.}\label{fig:diagonal:sparse}
\end{figure}

\paragraph{Where: the crossover in $\alpha$.} \Cref{fig:diagonal:sparse} shows the squared error of the limit against $\alpha$ on a sparse recovery problem with $n=40$ samples in dimension $d=100$, well above the $\ell_1$ threshold $20.4$ but far below $d$. For $\alpha\ge1$ the flow ends at the minimum-$\ell_2$ interpolator, with error $0.604$, close to the average $1-n/d=0.6$. Between $\alpha=1$ and $\alpha=0.03$ the error falls by more than an order of magnitude, and below that it falls as a power of $\alpha$. The flow and the static minimizer agree to within $5\times10^{-10}$ at every scale, which is \cref{thm:diagonal:bias} seen at work. The crossover happens where $\Phi_\alpha$ changes face, that is, where the quadratic cost $\norm{\diagwl2}^2/(16\alpha^2)$ of the dense interpolator becomes comparable with the $\ell_1$ cost $\frac14\log(1/\alpha^2)\norm{w_\star}_1$ of the sparse one. This gives the estimate $\alpha_c^2\log(1/\alpha_c^2)\approx\norm{\diagwl2}^2/(4\norm{w_\star}_1)$, used again in \cref{ch:grokking}; here it gives $\alpha_c\approx0.098$, while the error has halved at $\alpha\approx0.125$, so the estimate is right to within a factor of about $1.3$.

\paragraph{Where: the rich-regime exponent.} Below the crossover the error is a power law. The exponent comes from the dual description of $\diagwphi$.

\begin{heuristic}
By the optimality condition in the proof of \cref{thm:diagonal:bias}, $\nabla\Phi_\alpha(\diagwphi)=X^\top\nu$ for some $\nu\in\R^n$, that is, $\diagwphi=2\alpha^2\sinh(4X^\top\nu)$. Write $4\nu=\eps_\alpha\xi$ with $\eps_\alpha=\log(1/\alpha^2)$. Then $\abs{(\diagwphi)_i}\approx\alpha^{2(1-\abs{(X^\top\xi)_i})}$ whenever $(X^\top\xi)_i\ne0$. On the support $S$ of $w_\star$ the coordinates are of order one, which forces $(X^\top\xi)_i=\sgn(w_{\star,i})+O(1/\eps_\alpha)$; off the support $\abs{(X^\top\xi)_i}<1$, and the coordinate is of size $\alpha^{2(1-\abs{(X^\top\xi)_i})}$. A vector $\xi$ with $X_S^\top\xi=\sgn(w_{\star,S})$ and $\norm{X_{S^c}^\top\xi}_\infty<1$ is a \term{dual certificate}\index{dual certificate} for basis pursuit; it exists when $\ell_1$ minimization recovers $w_\star$ with some room to spare. Which certificate does $\diagwphi$ pick? The dual objective of the problem defining $\diagwphi$ contains the terms $\frac{\alpha^2}2\sum_i(\cosh(4(X^\top\nu)_i)-1)\approx\frac14\sum_i\alpha^{2(1-\abs{(X^\top\xi)_i})}$, dominated off the support by the largest $\abs{(X^\top\xi)_i}$, so for small $\alpha$ the dual picks the certificate minimizing $\max_{i\notin S}\abs{(X^\top\xi)_i}=1-\gamma^\star$. The largest off-support coordinate is then of size $\alpha^{2\gamma^\star}$, the on-support coordinates are perturbed by the same order to keep $Xw=y$, and
\begin{equation}\label{eq:diagonal:richrate}
  \norm{\diagwphi-w_\star}^2\asymp\alpha^{4\gamma^\star},\qquad \gamma^\star=1-\min\bigl\{\norm{X_{S^c}^\top\xi}_\infty:X_S^\top\xi=\sgn(w_{\star,S})\bigr\}.\qedhere
\end{equation}
\end{heuristic}

The margin $\gamma^\star$ of the best certificate is the value of a linear program,\sidenote{Minimize $m$ over $(\xi,m)$ subject to $X_S^\top\xi=\sgn(w_{\star,S})$ and $\abs{(X^\top\xi)_i}\le m$ for $i\notin S$; then $\gamma^\star=1-m$. The figure script solves it with the HiGHS solver of SciPy.} $0.516$ for the instance of \cref{fig:diagonal:sparse}; the measured slope between $\alpha=10^{-3}$ and $10^{-4}$ is $2.09$, against $4\gamma^\star=2.06$. The convergence of the \emph{vector} is thus polynomial in $\alpha$, much faster than the logarithmic convergence of the norms in \cref{thm:diagonal:limits}, and it degrades as the certificate's margin shrinks, that is, as $n$ approaches the $\ell_1$ threshold. \Textcite{MattStoger2025} (unrefereed) prove matching upper and lower bounds of this kind; for depth two the exponent is given by a null space constant of $X$.

\begin{figure}[t]
\begin{fullwidth}
\centering
\begin{tikzpicture}
\begin{groupplot}[group style={group size=2 by 1,horizontal sep=17mm},hplot,width=0.475\linewidth,height=0.37\linewidth]
\nextgroupplot[xmin=0,xmax=20,ymin=-0.05,ymax=1.55,ytick={0,0.5,1},xlabel={time $t$},ylabel={$w_i(t)/w_{\star,i}$},
  legend columns=3,legend style={at={(0.5,0.98)},anchor=north,font=\scriptsize}]
\addplot[theory] table[x=t,y=r1] {figures/data/03-diagonal-orthogonal-theory.dat};
\addplot[theory2] table[x=t,y=r2] {figures/data/03-diagonal-orthogonal-theory.dat};
\addplot[theory3] table[x=t,y=r3] {figures/data/03-diagonal-orthogonal-theory.dat};
\addplot[muted,line width=0.9pt,dashdotted,no marks] table[x=t,y=r4] {figures/data/03-diagonal-orthogonal-theory.dat};
\addplot[simulation] table[x=t,y=r1] {figures/data/03-diagonal-orthogonal-sim.dat};
\addplot[simulation,forget plot] table[x=t,y=r2] {figures/data/03-diagonal-orthogonal-sim.dat};
\addplot[simulation,forget plot] table[x=t,y=r3] {figures/data/03-diagonal-orthogonal-sim.dat};
\addplot[simulation,forget plot] table[x=t,y=r4] {figures/data/03-diagonal-orthogonal-sim.dat};
\legend{{$w_{\star,i}=1$},{$-0.6$},{$0.35$},{$0.2$},gradient descent}
\nextgroupplot[xmin=0,xmax=3,ymin=-0.03,ymax=0.8,xlabel={accelerated time $\tau=t/\log(1/\alpha)$},ylabel={training loss},
  legend style={at={(0.97,0.97)},anchor=north east}]
\addplot[guide] table[x=tau,y=stair] {figures/data/03-diagonal-staircase.dat};
\addplot[theory3] table[x=tau,y=a2] {figures/data/03-diagonal-staircase.dat};
\addplot[theory2] table[x=tau,y=a4] {figures/data/03-diagonal-staircase.dat};
\addplot[theory] table[x=tau,y=a8] {figures/data/03-diagonal-staircase.dat};
\legend{$\alpha\rightarrow0$,$\alpha=10^{-2}$,$\alpha=10^{-4}$,$\alpha=10^{-8}$}
\end{groupplot}
\end{tikzpicture}
\caption{Incremental learning under an orthogonal design ($d=4$, $n=8$, $X=\sqrt n\,Q$ with $Q$ a random $8\times4$ matrix with orthonormal columns, seed $1$; $w_\star=(1,-0.6,0.35,0.2)$, noiseless). Left: $\alpha=10^{-3}$; curves are the exact solution \eqref{eq:diagonal:riccati}, points are gradient descent on $(u,v)$ with step $10^{-3}$, recorded every $0.5$ time units, within $3\times10^{-3}$ of the curves. The half times are $3.45$, $5.54$, $9.12$ and $15.26$, equal to $\log(\abs{w_{\star,i}}/\alpha^2)/(4\abs{w_{\star,i}})$ to within $2\times10^{-5}$. Right: the training loss from \eqref{eq:diagonal:riccati} in accelerated time for three scales, and its limit, the staircase of \cref{thm:diagonal:orth} with steps at $\tau=0.5$, $0.83$, $1.43$ and $2.5$.}\label{fig:diagonal:orthogonal}
\end{fullwidth}
\end{figure}

\paragraph{When: plateaus of length $\log(1/\alpha)$.} \Cref{fig:diagonal:orthogonal} shows the orthogonal design. Each coordinate follows a logistic curve, the larger targets first; plain gradient descent with a small step tracks the exact solution. In accelerated time the loss curves for $\alpha=10^{-2}$, $10^{-4}$ and $10^{-8}$ approach the staircase, and the steps sharpen as $\alpha$ decreases because the transition width $1/\abs s$ is fixed in $t$ while the plateaus grow like $\log(1/\alpha)$. The time to learn a coordinate is inversely proportional to its size: a feature twice as weak takes twice as long, in units of $\log(1/\alpha)$. The first plateau, before anything is learned, lasts $\log(1/\alpha)/(2\max_i\abs{\diagbar_i})$; a small initialization buys sparsity with time.

\begin{figure}[t]
\begin{fullwidth}
\centering
\begin{tikzpicture}
\begin{groupplot}[group style={group size=2 by 1,horizontal sep=17mm},hplot,width=0.475\linewidth,height=0.37\linewidth]
\nextgroupplot[xmin=0,xmax=3.6,ymin=-0.95,ymax=0.5,xlabel={accelerated time $\tau$},ylabel={coordinates of $w$},
  legend style={at={(0.97,0.97)},anchor=north east}]
\addplot[theory,const plot] table[x=tau,y=w1] {figures/data/03-diagonal-saddles-limit.dat};
\addplot[theory,const plot,forget plot] table[x=tau,y=w2] {figures/data/03-diagonal-saddles-limit.dat};
\addplot[theory,const plot,forget plot] table[x=tau,y=w3] {figures/data/03-diagonal-saddles-limit.dat};
\addplot[theory,const plot,forget plot] table[x=tau,y=w4] {figures/data/03-diagonal-saddles-limit.dat};
\addplot[simulation,mark repeat=6] table[x=tau,y=a24w1] {figures/data/03-diagonal-saddles-flow.dat};
\addplot[simulation,mark repeat=6,forget plot] table[x=tau,y=a24w2] {figures/data/03-diagonal-saddles-flow.dat};
\addplot[simulation,mark repeat=6,forget plot] table[x=tau,y=a24w3] {figures/data/03-diagonal-saddles-flow.dat};
\addplot[simulation,mark repeat=6,forget plot] table[x=tau,y=a24w4] {figures/data/03-diagonal-saddles-flow.dat};
\legend{limit path,flow $\alpha=10^{-24}$}
\nextgroupplot[xmin=0,xmax=3.6,ymin=-0.03,ymax=0.9,xlabel={accelerated time $\tau$},ylabel={training loss},
  legend style={at={(0.97,0.97)},anchor=north east}]
\addplot[guide] table[x=tau,y=loss] {figures/data/03-diagonal-saddles-limit.dat};
\addplot[theory3] table[x=tau,y=a4loss] {figures/data/03-diagonal-saddles-flow.dat};
\addplot[theory2] table[x=tau,y=a12loss] {figures/data/03-diagonal-saddles-flow.dat};
\addplot[theory] table[x=tau,y=a24loss] {figures/data/03-diagonal-saddles-flow.dat};
\legend{limit path,$\alpha=10^{-4}$,$\alpha=10^{-12}$,$\alpha=10^{-24}$}
\end{groupplot}
\end{tikzpicture}
\caption{Saddle to saddle for a general design: $n=3$, $d=6$, Gaussian design and labels (seed $281$); four coordinates become active, the other two stay at zero. Curves on the left and the dashed staircase on the right: the limit path of \cref{def:diagonal:path}, with jumps at $\tau=0.319$, $0.582$, $1.485$ and $2.907$. Points on the left and curves on the right: the gradient flow, integrated in $(\log u,\log v)$ with an adaptive Runge--Kutta method (relative tolerance $10^{-10}$). At the third jump one coordinate enters and another, active since the second jump, returns to zero. At distance more than $0.15$ from the jumps the flow is within $0.27$, $0.12$ and $0.024$ of the limit path for $\alpha=10^{-4}$, $10^{-12}$ and $10^{-24}$.}\label{fig:diagonal:saddles}
\end{fullwidth}
\end{figure}

\paragraph{In which order: saddles.} \Cref{fig:diagonal:saddles} shows a general design with three samples in dimension six. The flow visits four saddles before reaching the minimum-$\ell_1$ interpolator, and at the third jump a coordinate that had been learned is abandoned: learning is not monotone in the support. The limit predicts the plateau \emph{heights}, losses of least-squares fits on subsets of coordinates, and the jump \emph{times} in units of $\log(1/\alpha)$; the approach to it is slow, through $1/\log(1/\alpha)$, so that at $\alpha=10^{-4}$ the plateaus are visible but rounded, and the limit is crisp only at scales that no network uses.

\paragraph{Limits of the predictions.} All of these are statements about $\alpha\to0$ or $\alpha\to\infty$ at fixed data, or about $t\to\infty$ at fixed $\alpha$, and the time to reach the limit grows quickly as $\alpha$ shrinks: in the run of \cref{fig:diagonal:sparse} the training loss falls below $10^{-20}$ at $t\approx10^4$ for $\alpha=10^{-2}$ but only at $t\approx3\times10^6$ for $\alpha=10^{-4}$, because the off-support coordinates, which carry the last $\alpha^{4\gamma^\star}$ of the error, move at rates proportional to their own tiny size.

\section{What the model gets wrong}\label{sec:diagonal:wrong}

\paragraph{What was thrown away.} \emph{Nonlinearity.} The predictor is linear in $x$; the network has no features to learn except the choice of coordinates, and ``sparse'' means sparse in a basis fixed in advance. In a real network the analogue of the coordinate basis is itself learned. \emph{Infinitesimal steps.} The mirror-flow identity \eqref{eq:diagonal:mirrorflow} holds for the flow; for gradient descent with step $\eta$ the products $u_iv_i$ are no longer conserved but multiplied at each step by $1-4\eta^2g_i^2$, and the potential drifts. \emph{Noise.} There is no label noise and no stochastic gradient; both change the implicit bias. \emph{The loss and the task.} The square loss on a regression problem attains its minimum, on the affine set of interpolators; logistic losses on classification problems do not. \emph{The parametrization.} Diagonal networks are among the few models whose gradient flow is a mirror flow; for others no potential need exist. \emph{The limit.} The theorems describe $t\to\infty$, and then $\alpha\to0$ or $\alpha\to\infty$; training stops at finite time, at finite $\alpha$.

\paragraph{Where it breaks.} \emph{Step size and stochasticity change the bias.} \Textcite{PesmePillaudVivienFlammarion2021} model stochastic gradient descent by a stochastic gradient flow and prove that it converges to the minimizer of a potential of the same form with an effective initialization scale smaller than $\alpha$, so stochasticity pushes toward the rich regime; the slower the loss converges, the stronger the effect. \Textcite{NacsonEtAl2022} show that large step sizes on non-centered data move gradient descent from the kernel to the rich regime, through a dynamical-stability argument and a weighted $\ell_1$ bound. \Textcite{EvenEtAl2023} prove convergence with macroscopic step sizes and characterize the limits: large steps help stochastic gradient descent recover sparse vectors, but for full-batch gradient descent they can hurt, most strongly near the edge of stability (\cref{ch:stability}). On the instance of \cref{fig:diagonal:sparse} the effect is small but visible (\cref{ex:diagonal:stepsize}). \emph{Label noise} is stronger still: \textcite{HaoChenEtAl2021} show that stochastic gradient descent with label noise recovers a sparse target from every initialization, where Gaussian parameter noise does not, and \textcite{PillaudVivienReygnerFlammarion2022} prove that label-noise dynamics solve a Lasso problem, a bias independent of $\alpha$. \emph{Classification.} With the exponential loss on separable data, \textcite{MoroshkoEtAl2020} find that whether the limit is the maximum-$\ell_1$ or the maximum-$\ell_2$ margin classifier depends on the initialization scale \emph{and} on the training accuracy reached, and that some limits appear only at losses far below $10^{-100}$: the asymptotic theory of \cref{ch:margin} and the rich limit of this chapter do not commute. \emph{Beyond diagonal networks.} \Textcite{LiWangLeeArora2022} show that gradient flow under a commuting parametrization is a mirror flow, and that conversely every mirror flow arises this way; their necessary condition for such an equivalence is weaker than commutation, and outside this class no potential need exist. For matrix factorization, where the analogue of the $\ell_1$ norm would be the nuclear norm, the conjecture of \textcite{GunasekarEtAl2017} fails: \textcite{RazinCohen2020} construct matrix completion problems in which gradient flow from small initialization drives every norm and quasi-norm to infinity while lowering the rank, and \textcite{LiLuoLyu2021} show that from infinitesimal initialization the flow performs greedy low-rank learning, which differs from nuclear norm minimization. For a single ReLU neuron with the square loss, \textcite{VardiShamir2021} prove that no function of the parameters describes the implicit bias exactly. \emph{The real phenomenon.} The model cannot explain why fully connected networks do better in the lazy regime \parencite{GeigerEtAl2020}: it has no notion of a target that a kernel fits well.

\paragraph{Fine structure.} \emph{Depth.} \Textcite{WoodworthEtAl2020} solve diagonal networks $w=u^{\odot L}-v^{\odot L}$ of every depth $L$; the potential changes, the rich limit is still $\ell_1$, and it is reached at polynomially rather than exponentially small scales (\cref{ex:diagonal:depth}), as \textcite{GissinShalevShwartzDaniely2020} observed; \textcite{ChouMalyRauhut2023} give quantitative $\ell_1$ bounds for deeper factorizations, and \textcite{MattStoger2025} matching rates in $\alpha$. \emph{Initialization shape.} \Textcite{AzulayEtAl2021} show that the relative scales of the layers at initialization enter the potential, and that the infinitesimal mirror-descent picture breaks down beyond diagonal models. \emph{Dynamics.} \Textcite{Berthier2023} proves incremental learning, one coordinate at a time, in an underparameterized setting with anticorrelated features; \textcite{PesmeFlammarion2023} handle general designs (\cref{thm:diagonal:saddles}); \textcite{Berthier2025} (unrefereed) relates the whole trajectory to the Lasso regularization path, with time playing the role of an inverse regularization parameter, exactly under a monotonicity assumption on the path. \Textcite{NishiyamaImaizumi2026} derive the high-dimensional dynamics on random designs by dynamical mean-field theory. \emph{Statistics.} With noisy labels, early-stopped gradient descent on the Hadamard parametrization attains the minimax rate of sparse regression \parencite{VaskeviciusKanadeRebeschini2019,ZhaoYangHe2022}, and the rich limit, which interpolates the noise, is the minimum-$\ell_1$ interpolator, whose prediction error is of order $\sigma^2/\log(d/n)$ \parencite{WangDonhauserYang2022}: consistent, but slowly. Open are a description of the implicit bias of gradient descent with large steps on general data, and a comparable theory for parametrizations without a potential, matrix factorization first.

\section{Model card}\label{sec:diagonal:card}

\modelcard{Among the interpolators of an overparameterized model, gradient descent picks one that depends on the initialization scale: kernel-like for large scales, sparse or low rank and learned incrementally for small ones.}{A linear predictor $w=u\odot u-v\odot v$, square loss, gradient flow from $u=v=\alpha\mathbf1$.}{Gradient flow is a mirror flow with the hyperbolic entropy $\Phi_\alpha$; it converges to $\argmin\{\Phi_\alpha:Xw=y\}$; this tends to the minimum-$\ell_2$ interpolator as $\alpha\to\infty$ and to the minimum-$\ell_1$ one as $\alpha\to0$; closed-form incremental learning for orthogonal designs; saddle-to-saddle limit for general ones.}{Error $\asymp\alpha^{4\gamma^\star}$ in the rich regime; crossover at $\alpha_c^2\log(1/\alpha_c^2)\approx\norm{\diagwl2}^2/(4\norm{w_\star}_1)$; feature $i$ learned at $t\approx\log(1/\alpha)/(2\abs{w_{\star,i}})$; $\Lhat(w(t))\le\Phi_\alpha(w_\star)/(2t)$.}{Step size, stochastic gradients and label noise change the bias; classification limits depend on the training accuracy; no potential exists for matrix factorization or ReLU networks; fully connected networks can be better in the lazy regime.}{Depth-$L$ diagonal networks; initialization shape; stochastic and large-step dynamics; label-noise SGD and the Lasso; greedy low-rank learning in matrix factorization; dynamical mean-field theory.}

The chapter sits at the crossing of three threads of the book. \emph{Escape from small initialization}: the plateaus of length $\log(1/\alpha)/\text{strength}$ are those of deep linear networks in \cref{ch:deeplinear}, and reappear in the search phase of \cref{ch:singleindex}, the committee-machine plateaus of \cref{ch:committee} and the hidden progress of \cref{ch:parity}. \emph{Implicit bias}: the minimum-$\Phi_\alpha$ interpolator is the regression counterpart of the maximum-margin classifier of \cref{ch:margin}, and \cref{ch:grokking} shows that weight decay acts on a diagonal network as an initialization scale shrinking like $e^{-\lambda t}$, which turns the crossover of \cref{fig:diagonal:sparse} into grokking. \emph{Lazy versus rich}: the kernel limit is the lazy training of \cref{ch:lazy} with a linear kernel, whose risk is the overparameterized branch of \cref{ch:doubledescent}.

\begin{exercises}
\item \consolidation (Two expansions of $q$.) (a) Show that $q(z)=\frac{z^2}4-\frac{z^4}{192}+O(z^6)$ as $z\to0$, and $q(z)=\abs z\log\abs z-\abs z+2-1/\abs z+O(\abs z^{-3})$ as $\abs z\to\infty$. (b) Use (a) and the optimality condition of \cref{thm:diagonal:bias} to show that $\norm{\diagwphi-\diagwl2}=O(\alpha^{-4})$ as $\alpha\to\infty$, and compare with \cref{thm:diagonal:limits}(i). \label{ex:diagonal:taylor}

\item \consolidation (The half time.) From \eqref{eq:diagonal:riccati}, show that $w_i(t_{1/2})=s/2$ for $t_{1/2}=\frac1{4\rho}\log\bigl(\kappa_0\frac{P_h-p_-}{p_+-P_h}\bigr)$ with $P_h=\frac14(s+\sqrt{s^2+16\alpha^4})$, and deduce \eqref{eq:diagonal:thalf}. Check it against the numbers in the caption of \cref{fig:diagonal:orthogonal}. \label{ex:diagonal:half}

\item \consolidation\hard (Ties in the rich limit.) Suppose that the set $M$ of minimum-$\ell_1$-norm interpolators has more than one point. (a) Show that two points of $M$ have no coordinate of opposite signs, and deduce that $w\mapsto\sum_i\abs{w_i}\log\abs{w_i}$ is strictly convex on $M$. (b) Show that $\diagwphi$ converges, as $\alpha\to0$, to the unique minimizer of $\sum_i\abs{w_i}\log\abs{w_i}$ on $M$. \hint{for (a), if $w,w'\in M$ then $\frac12(w+w')\in M$, which forces $\abs{w_i+w'_i}=\abs{w_i}+\abs{w'_i}$. For (b), compare $\Phi_\alpha(\diagwphi)$ with $\Phi_\alpha(w')$ for $w'\in M$ through \eqref{eq:diagonal:l1form}, and use $\norm{\diagwphi}_1\ge\norm{w'}_1$ to show $\diagH(\diagwphi)\le\diagH(w')+2d\alpha^2$.}\label{ex:diagonal:tie}

\item \consolidation (Noisy labels.) Drop the assumption that $y\in\operatorname{range}(X)$. Rerun the proof of \cref{thm:diagonal:bias} with $w^\circ$ a minimizer of $\Lhat$, and show that $\Lhat(w(t))-\min\Lhat\le\Phi_\alpha(w^\circ)/(2t)$ and that $w(t)$ converges to the minimizer of $\Phi_\alpha$ among the minimizers of $\Lhat$. \label{ex:diagonal:noisy}

\item \consolidation (Lazy means small motion.) For $\alpha\ge\norm{\diagwl2}^{1/2}$, use \cref{prop:diagonal:closed} and \cref{thm:diagonal:limits} to show that $\max_i\abs{u_i(\infty)-\alpha}/\alpha=O(\norm{\diagwl2}/\alpha^2)$: each parameter moves by a vanishing fraction of its initial value, the defining property of lazy training in \cref{ch:lazy}. \label{ex:diagonal:lazy}

\item \variation (Initialization shape.) Start from $u(0)=a$ and $v(0)=b$ with positive vectors $a,b$. Show that $u_iv_i=a_ib_i$ is conserved and that $w(t)$ converges to the minimizer, among interpolators, of
\[
  \sum_{i=1}^d\Bigl(\frac{a_ib_i}4\,q\Bigl(\frac{w_i}{a_ib_i}\Bigr)-\frac14\log\frac{a_i}{b_i}\,w_i\Bigr).
\]
Which signs does $a_i>b_i$ favor, and what happens to the rich limit when $a_ib_i\to0$ with $a_i/b_i$ fixed? Compare with \textcite{AzulayEtAl2021}. \label{ex:diagonal:shape}

\item \variation (The product parametrization.) Let $w=u\odot v$, $u(0)=c\mathbf1$ and $v(0)=0$, with $c>0$. Show that $u_i^2-v_i^2=c^2$ is conserved, that $w=-\frac{c^2}2\sinh(2\diagG)$, and that the trajectory is that of \eqref{eq:diagonal:param} with $\alpha^2=c^2/4$ run at half speed: $w_{\mathrm{prod}}(2t)=w(t)$. \label{ex:diagonal:product}

\item \variation\hard (Depth.) Let $w=u^{\odot L}-v^{\odot L}$ with $L\ge3$ and $u(0)=v(0)=\alpha\mathbf1$. (a) Show that $u_i^{2-L}=\alpha^{2-L}+L(L-2)\diagG_i$ and $v_i^{2-L}=\alpha^{2-L}-L(L-2)\diagG_i$ as long as both are positive, and that $w_i=-\alpha^L\psi_L\bigl(L(L-2)\alpha^{L-2}\diagG_i\bigr)$ with $\psi_L(z)=(1-z)^{-L/(L-2)}-(1+z)^{-L/(L-2)}$ on $(-1,1)$. (b) Deduce a mirror flow with $\nabla\Phi(w)_i=\psi_L^{-1}(w_i/\alpha^L)/(L(L-2)\alpha^{L-2})$. (c) Show that $L(L-2)\alpha^{L-2}\,\Phi(w)\to\norm w_1$ as $\alpha\to0$, with a relative error that is a positive power of $\alpha$ for $w$ in a compact set avoiding the coordinate hyperplanes, instead of the $1/\log(1/\alpha)$ of depth two. \label{ex:diagonal:depth}

\item \variation (Unequal feature scales.) Replace \cref{ass:diagonal:orth} by $X^\top X/n=\diag(\lambda_1,\dots,\lambda_d)$ with $\lambda_i>0$. Solve the flow coordinate by coordinate and show that coordinate $i$ is learned at accelerated time $1/(2\lambda_i\abs{\diagbar_i})=n/(2\abs{\ip{X_{:i},y}})$: the order of learning is the order of the correlations of the features with the labels, not of the coefficients. \label{ex:diagonal:scales}

\item \simulation (The phase transition.) With $d=100$ and $k=5$, Gaussian designs and $\pm1$ targets, compute the limit $w^\Phi_\alpha$ of the flow (by \cref{thm:diagonal:bias}, through Newton's method on the dual, as in the script of \cref{fig:diagonal:sparse}) for $\alpha\in\{10,10^{-3},10^{-8}\}$, $n=10,15,\dots,60$ and $20$ random instances each, and plot the fraction of instances with $\norm{w-w_\star}\le10^{-2}\norm{w_\star}$. The reader should observe that $\alpha=10$ never succeeds; that for $\alpha=10^{-3}$ the success rate rises from $0$ at $n=25$ to nearly $1$ at $n=50$; and that for $\alpha=10^{-8}$ the rise comes earlier and is steeper, crossing one half between $n=20$ and $n=25$, close to the statistical dimension $20.4$ of the $\ell_1$ descent cone \parencite{AmelunxenEtAl2014}. Explain the shift with \eqref{eq:diagonal:richrate}: near the threshold the certificate margin $\gamma^\star$ is small. \label{ex:diagonal:transition}

\item \simulation (Step size.) On the instance of \cref{fig:diagonal:sparse} with $\alpha=0.1$, run gradient descent on $(u,v)$ with the steps $0.0004$, $0.0087$, $0.0216$ and $0.0346$ until $\Lhat<10^{-11}$. The reader should observe errors $\norm{w-w_\star}^2$ of $0.2284$, $0.2289$, $0.2297$ and $0.2306$: the limit moves away from $w^\Phi_\alpha$ as $\eta$ grows, toward denser solutions on this centered design, as \textcite{EvenEtAl2023} predict for full-batch gradient descent. Find the step at which gradient descent diverges, and explain the trend using the factor $1-4\eta^2g_i^2$ by which a step multiplies $u_iv_i$. \label{ex:diagonal:stepsize}

\item \simulation (Saddles.) With the script of \cref{fig:diagonal:saddles}, change the seed, and plot the accelerated times at which the training loss of the flow crosses the midpoints of the predicted steps, for $\alpha$ from $10^{-4}$ down to $10^{-24}$. The reader should observe convergence to the jump times of \cref{def:diagonal:path}, with an error decreasing roughly like $1/\log(1/\alpha)$. Show from \cref{def:diagonal:path} that if a coordinate has opposite signs at two saddles, then it vanishes at some saddle in between. \label{ex:diagonal:saddlesim}

\item \research (Stochastic gradients with large steps.) \Textcite{PesmePillaudVivienFlammarion2021} treat a stochastic gradient flow and \textcite{EvenEtAl2023} stochastic gradient descent with large steps, each with conditions on the data. A description of the limit of stochastic gradient descent on diagonal networks for general designs, as an explicit function of the step size, the batch size and $\alpha$, uniform in the dimension, is open; so is the high-dimensional test error of that limit on Gaussian designs. \label{ex:diagonal:sgd}

\item \research (A potential for matrices.) For matrix sensing with a symmetric factorization $W=UU^\top$ from $U(0)=\alpha I$, the flow is a mirror flow when the measurement matrices commute \parencite{GunasekarEtAl2017}, and \textcite{LiLuoLyu2021} describe the limit $\alpha\to0$ as greedy low-rank learning under assumptions. Characterize the limit of the flow at fixed small $\alpha$ for general measurements, or prove that no function of the end-to-end matrix describes it. \label{ex:diagonal:matrix}
\end{exercises}

\notesandrefs
The implicit-bias program grew out of the observations of \textcite{NeyshaburTomiokaSrebro2015} and \textcite{ZhangEtAl2017}. \Textcite{GunasekarEtAl2018} showed that mirror descent in underdetermined regression converges to the interpolator closest to the initialization in the Bregman divergence of its potential, which is the template of \cref{thm:diagonal:bias}. The diagonal network and the potential $\Phi_\alpha$ are from \textcite{WoodworthEtAl2020}, first posted in 2019, who also treated depth $L$ and the transition between the regimes; the equivalent product parametrization $w=u\odot v$ (\cref{ex:diagonal:product}) had been used in statistics by \textcite{Hoff2017}, who showed that the Lasso can be computed by alternating ridge regressions in the factors, and \textcite{VaskeviciusKanadeRebeschini2019} and \textcite{ZhaoYangHe2022} used Hadamard-product parametrizations with early stopping for sparse regression. The potential is the hyperbolic entropy of \textcite{GhaiHazanSinger2020}, introduced to interpolate between gradient descent and the exponentiated gradient method of \textcite{KivinenWarmuth1997}, which does well when few inputs are relevant.

The quantitative forms of \cref{thm:diagonal:limits} are ours; Woodworth and coauthors give the limits and discuss how small $\alpha$ must be. \Cref{thm:diagonal:orth} is a direct integration in the spirit of \textcite{WoodworthEtAl2020} and \textcite{GissinShalevShwartzDaniely2020}; the algorithm of \textcite{PesmeFlammarion2023} is restated in our normalization and checked numerically in \cref{fig:diagonal:saddles}.

\paragraph{Translating the literature.} \Textcite{WoodworthEtAl2020} write $\beta$ for our $w$ and $Q_\alpha(\beta)=\alpha^2\sum_iq(\beta_i/\alpha^2)$ with the same $q$, so their $Q_\alpha$ is our $4\Phi_\alpha$ and has the same minimizers; their depth is $D$, our $L$. \Textcite{PesmeFlammarion2023}, \textcite{PesmePillaudVivienFlammarion2021} and \textcite{EvenEtAl2023} write $\beta$ for our $w$; Pesme and Flammarion use $\beta=u\odot v$ from $u(0)=\sqrt2\,\alpha\mathbf1$, $v(0)=0$, which by \cref{ex:diagonal:product} is our network with $\alpha^2/2$ in place of $\alpha^2$, run at half speed. Their normalizations of the loss and of accelerated time may differ from ours by further constant factors, which the reader should take from the paper at hand rather than from us. We have restated their saddle-to-saddle algorithm with our dual variable $\zeta$, whose coordinates reach $\pm1$ at the jumps. \Textcite{GhaiHazanSinger2020} write $\phi_\beta(x)=x\arcsinh(x/\beta)-\sqrt{x^2+\beta^2}$ for the hypentropy; our $\Phi_\alpha(w)$ is $\frac14\sum_i(\phi_\beta(w_i)-\phi_\beta(0))$ with $\beta=2\alpha^2$. \Textcite{ChizatOyallonBach2019} write $\alpha$ for a multiplier of the model's output, which plays the role of our $\alpha^2$: the output of a diagonal network scales as $\alpha^2$ when its parameters scale as $\alpha$.

\printchapterbib
\end{refsection}

% Chapter 4 of The Hydrogen Atoms of Deep Learning: training past zero error.
\newcommand{\marginu}{u^\star}
\newcommand{\marginwhat}{\hat w}
\newcommand{\marginwt}{\tilde w}
\newcommand{\marginsm}{\tilde\gamma}
\newcommand{\marginS}{\mathcal S}

\casechapter{Training Past Zero Error}{Logistic regression on separable data}{Linear classifiers, exponential-tailed losses}{ch:margin}
\begin{refsection}[chapters/04-margin.bib]

\begin{chaptersummary}
Networks trained with the cross-entropy reach zero training error long before training stops, and what happens afterwards matters: the loss keeps falling like $1/t$, the weights grow like $\log t$, and the classifier keeps changing. The hydrogen atom of this regime is a linear classifier trained by gradient flow on the exponential loss over linearly separable data. We prove that the loss is at most $1/(1+\gamma^{\star2}t)$, that the norm grows like $\log t/\gamma^\star$, that the normalized margin is within a factor $1-\log n/\log(\gamma^{\star2}t/n)$ of the maximal margin $\gamma^\star$, and that the direction of the weights converges to the hard-margin support vector machine; the same holds for gradient descent, and a monotone smoothed margin carries the argument to homogeneous networks. The model is right about the direction and wrong about the speed: every rate is logarithmic, and for nonlinear networks the limit is a stationary point of the margin problem that need not be a maximizer.
\end{chaptersummary}

\section{The phenomenon: the loss keeps falling after the errors stop}\label{sec:margin:phenomenon}

A classifier trained with the cross-entropy never finishes. Once every training point is classified correctly, the loss is still positive, and it can be lowered further by scaling up the logits; no finite set of weights attains its infimum. Large networks reach this regime early: \textcite{ZhangEtAl2017} showed that standard architectures for image classification drive the training error to zero even when the labels are replaced by random ones, and on real labels zero training error is the ordinary state of affairs long before the end of training. What does gradient descent do in the time that remains?

\Textcite{SoudryEtAl2018} gave a detailed answer for a convolutional network on CIFAR-10 trained by stochastic gradient descent with momentum at a constant learning rate. After the training error reached zero, the training loss decayed roughly like $t^{-1}$ in the number of iterations, the Euclidean norm of the last weight layer grew roughly like $\log t$, the validation loss eventually started to \emph{increase}, and the validation error, in contrast, kept slowly improving.\sidenote{The opposite movements of validation loss and validation error matter in practice: early stopping on the validation loss would have stopped training while the error was still improving.} A year earlier \textcite{HofferHubaraSoudry2017} had found, across several architectures, that the distance of the weights from their initialization grows logarithmically with the number of updates, and that the generalization gap of large-batch training closes when the number of updates is matched to that of small-batch training: train longer, generalize better.

\begin{observation}[\textcite{SoudryEtAl2018}; \textcite{HofferHubaraSoudry2017}]\label[observation]{obs:margin:late}
For networks trained with the cross-entropy past zero training error, at a constant learning rate: the training loss decays roughly like $1/t$ and the weight norm grows roughly like $\log t$; the classifier keeps changing; the test loss can increase while the test error decreases.
\end{observation}

The status of the observation is uneven. The rates of the loss and of the norm have been seen repeatedly and are what the model below predicts. ``The test error keeps improving'' is weaker: it was reported for particular runs, it depends on the data, and it fails in simple settings. In a linear model trained on a high-dimensional Gaussian mixture we found the test error drifting slightly \emph{upward} after the training error reached zero (\cref{sec:margin:wrong}). What is robust is that the training does not stop: the predictor keeps moving, slowly, and in a definite direction.

Boosting saw the same phenomenon first. \Textcite{SchapireEtAl1998} observed that AdaBoost, which greedily minimizes an exponential loss, keeps lowering its test error after the training error reaches zero, and explained it by the growth of the training margins.\sidenote{For a voting classifier the margin of a training example is the weighted vote for its correct label minus the largest weighted vote for an incorrect label, normalized by the total weight.} The explanation was disputed for boosting, and margin-based explanations of deep networks are disputed today (\cref{sec:margin:wrong}). But the mechanism they point to is real, and in the model of this chapter it is exact.

Three questions follow. Where does the predictor go once the training set is fit, given that it never converges? How fast? And does the answer depend on the loss, the step size and the architecture?

\section{The model: a linear classifier and an exponential tail}\label{sec:margin:model}

The data are $n$ pairs $(x_i,y_i)$ with $x_i\in\R^d$ and labels $y_i\in\{-1,+1\}$. A linear classifier predicts $\sgn\ip{w,x}$, and the \term{margin} of the training point $i$ is $q_i(w)=y_i\ip{w,x_i}$\index{margin}; the point is classified correctly when $q_i(w)>0$. It is convenient to fold the labels into the data, $z_i=y_ix_i$, so that $q_i(w)=\ip{w,z_i}$. The training loss is the \term{exponential loss}\index{exponential loss}
\begin{equation}\label{eq:margin:loss}
  \Lhat(w)=\frac1n\sum_{i=1}^ne^{-\ip{w,z_i}},\qquad \nabla\Lhat(w)=-\frac1n\sum_{i=1}^ne^{-\ip{w,z_i}}z_i ,
\end{equation}
and we train it by gradient flow, $\dot w(t)=-\nabla\Lhat(w(t))$, from $w(0)=0$; gradient descent comes in \cref{sec:margin:steps}. The logistic loss $\frac1n\sum_i\log(1+e^{-q_i})$ has the same tail, $\log(1+e^{-q})=e^{-q}(1+o(1))$ as $q\to\infty$, and gives the same results with more bookkeeping (\cref{ex:margin:logistic}).\sidenote{With $K$ classes the cross-entropy of logits $f_1,\dots,f_K$ with correct class $c$ is $\log\bigl(1+\sum_{k\ne c}e^{-(f_c-f_k)}\bigr)$, again exponential in the logit gaps $f_c-f_k$; \textcite{SoudryEtAl2018} treat the multiclass case.}

\begin{assumption}[Separable, bounded data]\label[assumption]{ass:margin:data}
Every $z_i$ satisfies $\norm{z_i}\le1$, and some $u\in\R^d$ satisfies $\ip{u,z_i}>0$ for every $i$.
\end{assumption}

The bound on the norms only fixes units: rescaling the data by $c$ rescales time by $c^{-2}$ and margins by $c$. Separability is the substance.

\begin{definition}[Normalized margin and maximal margin]\label[definition]{def:margin:margin}
The \term{normalized margin}\index{margin!normalized} of $w\ne0$ is $\gamma(w)=\min_iq_i(w)/\norm w$. The \term{maximal margin} is $\gamma^\star=\max_{\norm u=1}\min_i\ip{u,z_i}$, and a unit vector attaining it is a \term{max-margin direction}\index{max-margin direction}. The \term{support vectors}\index{support vector} are the points with $\ip{\marginu,z_i}=\gamma^\star$.
\end{definition}

The normalized margin is invariant under $w\mapsto cw$ for $c>0$, as the predictions are. The maximum exists because $u\mapsto\min_i\ip{u,z_i}$ is continuous on the compact unit sphere, and by \cref{ass:margin:data} it is positive; by the Cauchy--Schwarz inequality $\gamma^\star\le\min_i\norm{z_i}\le1$. In the language of support vector machines, $\gamma^\star$ is the largest width of a slab around a hyperplane through the origin that separates the two classes, and $\marginwhat=\marginu/\gamma^\star$ is the solution of the hard-margin problem\index{hard-margin SVM} $\min\norm w$ subject to $\ip{w,z_i}\ge1$ for all $i$ \parencite{BoserGuyonVapnik1992,CortesVapnik1995}.

\paragraph{What was thrown away.} Everything below the last layer: the features $x_i$ are fixed, and only a linear readout is trained. Stochastic gradients, momentum, adaptive methods, weight decay and normalization layers are gone, and the step size is infinitesimal. Two simplifications are deliberate and right to make first. Separability is the defining property of the regime ``past zero error'': a network that fits its training set has, in particular, last-layer features on which the training set is separable. And the exponential loss isolates the one property of the cross-entropy that matters here, the exponential tail, without the curvature near $q=0$ that complicates the beginning of training and is irrelevant at its end.

\begin{dictionary}
linear predictor $\ip{w,x}$ on fixed inputs $x_i$ & the last layer of a network, on the features of the layer below\\
exponential loss $e^{-q}$ & cross-entropy, whose tail is $e^{-q}$ in the logit gap\\
separable data, $\gamma^\star>0$ & zero training error, reached early by overparameterized networks\\
$\Lhat(w(t))\approx1/(\gamma^{\star2}t)$ & training loss decaying like $1/t$\\
$\norm{w(t)}\approx\log t/\gamma^\star$ & weight norms and logit scales growing like $\log t$\\
direction $w/\norm w$, normalized margin $\gamma(w)$ & the classifier itself, and the quantity it slowly improves\\
max-margin direction $\marginu$, support vectors & hard-margin SVM on the features; Karush--Kuhn--Tucker points for homogeneous networks\\
time $t$ & learning rate times number of steps\\
\end{dictionary}

\paragraph{What one had to notice.} That a loss without a minimizer still determines a limit, the limit of the \emph{direction}; and that the exponential tail turns the gradient into a soft minimum. Write $w=\beta u$ with $\norm u=1$. Then
\[
  -\log\bigl(n\Lhat(\beta u)\bigr)=-\log\sum_ie^{-\beta\ip{u,z_i}}
\]
is a free energy at inverse temperature $\beta$, and as $\beta\to\infty$ it tends to the ground-state energy $\beta\min_i\ip{u,z_i}$: the Boltzmann weights $e^{-\beta\ip{u,z_i}}/\sum_je^{-\beta\ip{u,z_j}}$, which are the weights of the points in the gradient \eqref{eq:margin:loss}, concentrate on the points of smallest margin. Gradient descent on the loss at large norm is gradient ascent on the smallest margin.\sidenote{The function $-\log\sum_ie^{-q_i}$ is the log-sum-exp, or soft minimum, of the margins: it lies between $\min_iq_i-\log n$ and $\min_iq_i$.}\index{free energy} Training past zero error is a slow quench to zero temperature, and its ground state is the support vector machine.

\section{The solution: the loss, the norm, and the margin}\label{sec:margin:solution}

We first describe the target, the max-margin direction, by a lemma that collects the convex geometry we need. Write $\Delta_n=\{p\in\R^n:p_i\ge0,\ \sum_ip_i=1\}$ for the probability simplex and $C=\{\sum_ip_iz_i:p\in\Delta_n\}$ for the convex hull of the points $z_i$.

\begin{lemma}[The max-margin direction]\label[lemma]{lem:margin:svm}
Under \cref{ass:margin:data}:
\begin{enumerate}
\item The max-margin direction $\marginu$ is unique.
\item The point $z^\star$ of $C$ closest to the origin is $z^\star=\gamma^\star\marginu$; equivalently, $\gamma^\star=\min_{p\in\Delta_n}\norm{\sum_ip_iz_i}$.
\item Write $z^\star=\sum_ip^\star_iz_i$ with $p^\star\in\Delta_n$. Then $p^\star_i>0$ only for support vectors, and $\marginwhat=\marginu/\gamma^\star=\sum_i\alpha_iz_i$ with $\alpha_i=p^\star_i/\gamma^{\star2}\ge0$ and $\sum_i\alpha_i=1/\gamma^{\star2}$.
\item If a unit vector $u$ satisfies $\min_i\ip{u,z_i}\ge\gamma^\star-\delta$ with $\delta\ge0$, then $\norm{u-\marginu}^2\le2\delta/\gamma^\star$.
\end{enumerate}
\end{lemma}

\begin{proof}
\emph{Part (i).} Suppose that unit vectors $u_1\ne u_2$ both attain $\gamma^\star$. Their midpoint $v=\frac12(u_1+u_2)$ satisfies $\ip{v,z_i}\ge\gamma^\star$ for every $i$, and $\norm v<1$ by strict convexity of the Euclidean ball. Then $v\ne0$, and $v/\norm v$ has margin at least $\gamma^\star/\norm v>\gamma^\star$, a contradiction.

\emph{Part (ii).} The set $C$ is compact and convex, so it has a unique point $z^\star$ closest to the origin, characterized by $\ip{z-z^\star,z^\star}\ge0$ for every $z\in C$. By \cref{ass:margin:data} every $z\in C$ has $\ip{u,z}>0$ for the separating $u$, so $z^\star\ne0$. For every unit vector $u$, the minimum of the linear function $\ip{u,\cdot}$ over $C$ is attained at a vertex, so
\[
  \min_i\ip{u,z_i}=\min_{z\in C}\ip{u,z}\le\ip{u,z^\star}\le\norm{z^\star},
\]
and hence $\gamma^\star\le\norm{z^\star}$. Conversely, for $u=z^\star/\norm{z^\star}$ the characterization of $z^\star$ gives $\ip{u,z}\ge\ip{u,z^\star}=\norm{z^\star}$ for every $z\in C$, in particular for every $z_i$. So $\gamma^\star=\norm{z^\star}$, this $u$ attains it, and by part (i) $u=\marginu$.

\emph{Part (iii).} By part (ii), $\gamma^\star=\ip{\marginu,z^\star}=\sum_ip^\star_i\ip{\marginu,z_i}$, a convex combination of numbers that are all at least $\gamma^\star$. So every $i$ with $p^\star_i>0$ has $\ip{\marginu,z_i}=\gamma^\star$. Dividing $z^\star=\gamma^\star\marginu$ by $\gamma^{\star2}$ gives the formula for $\marginwhat$, and $\sum_i\alpha_i=1/\gamma^{\star2}$ because $p^\star\in\Delta_n$.

\emph{Part (iv).} Because $z^\star\in C$, the first display of part (ii) gives $\ip{u,z^\star}\ge\min_i\ip{u,z_i}\ge\gamma^\star-\delta$, that is, $\ip{u,\marginu}\ge1-\delta/\gamma^\star$. Since both vectors have norm one, $\norm{u-\marginu}^2=2-2\ip{u,\marginu}\le2\delta/\gamma^\star$.
\end{proof}

Part (iii) is the Karush--Kuhn--Tucker condition\index{Karush--Kuhn--Tucker point} of the hard-margin problem: $\marginwhat$ is a nonnegative combination of the support vectors. Part (ii) is the classical duality between the hard-margin support vector machine and the nearest point of a convex hull.\sidenote{The weights $p^\star_i$ are the Lagrange multipliers of the hard-margin problem, normalized to sum to one.} Part (iv) turns margins into angles: a direction whose margin is within $\delta$ of the best is within $\sqrt{2\delta/\gamma^\star}$ of $\marginu$. \Cref{fig:margin:geometry} shows the picture.

\begin{figure}[t]
\centering
\begin{tikzpicture}[scale=1.55,font=\small]
\fill[seamist] (1.6,0.3) -- (1.9,1.4) -- (1.3,2.0) -- (0.6,1.3) -- cycle;
\draw[sea,line width=0.6pt] (1.6,0.3) -- (1.9,1.4) -- (1.3,2.0) -- (0.6,1.3) -- cycle;
\draw[muted,densely dashed] (-0.55,0.55) -- (1.35,-1.35);
\draw[muted,densely dotted] (2.15,-0.25) -- (0.05,1.85);
\draw[->,ink,line width=0.7pt] (0,0) -- (0.95,0.95);
\foreach \p in {(1.6,0.3),(1.9,1.4),(1.3,2.0),(0.6,1.3),(1.3,1.0),(1.5,1.5),(1.1,1.6),(1.6,0.9),(1.0,1.25)}
  \fill[sandline!55!black] \p circle (1.3pt);
\draw[seadeep,line width=0.9pt] (1.6,0.3) circle (2.6pt) (0.6,1.3) circle (2.6pt);
\fill[ink] (0,0) circle (0.9pt) node[below left] {$0$};
\fill[ink] (0.95,0.95) circle (1.1pt);
\draw[muted,line width=0.4pt] (0.98,0.9) -- (1.16,0.58);
\node[anchor=north] at (1.18,0.58) {$z^\star$};
\node[anchor=north,text=muted] at (1.25,-1.25) {$\{w:\ip{\marginu,w}=0\}$};
\node[anchor=west,text=sea] at (1.95,1.75) {$C$};
\node[anchor=west,text=seadeep] at (1.7,0.25) {support vector};
\end{tikzpicture}
\caption{The geometry of \cref{lem:margin:svm} in the plane. The points are the $z_i=y_ix_i$; their convex hull $C$ is shaded. The point $z^\star$ of $C$ closest to the origin has norm $\gamma^\star$ and direction $\marginu$. The decision boundary (dashed) is orthogonal to $\marginu$, and the support vectors (circled) lie on the dotted line $\ip{\marginu,z}=\gamma^\star$; $z^\star$ is a convex combination of them.}\label{fig:margin:geometry}
\end{figure}

Two inequalities connect the loss to the geometry. The first bounds the gradient in terms of the loss; the second bounds the smallest margin.

\begin{lemma}[Gradient and soft minimum]\label[lemma]{lem:margin:grad}
Under \cref{ass:margin:data}, every $w\in\R^d$ satisfies
\begin{gather}
  \gamma^\star\,\Lhat(w)\le\ip{-\nabla\Lhat(w),\marginu}\le\norm{\nabla\Lhat(w)}\le\Lhat(w),\label{eq:margin:grad}\\
  -\log\bigl(n\Lhat(w)\bigr)\le\min_iq_i(w)\le\gamma^\star\norm w .\label{eq:margin:softmin}
\end{gather}
\end{lemma}

\begin{proof}
For \eqref{eq:margin:grad}: by \eqref{eq:margin:loss}, $\ip{-\nabla\Lhat(w),\marginu}=\frac1n\sum_ie^{-q_i(w)}\ip{z_i,\marginu}$, and every $\ip{z_i,\marginu}\ge\gamma^\star$; the middle inequality is the Cauchy--Schwarz inequality; and the triangle inequality with $\norm{z_i}\le1$ gives the last. For \eqref{eq:margin:softmin}: the sum $n\Lhat(w)=\sum_ie^{-q_i(w)}$ is at least its largest term $e^{-\min_iq_i(w)}$, which gives the left inequality; the right one holds for $w=0$, and for $w\ne0$ it is the definition of $\gamma^\star$ applied to $w/\norm w$.
\end{proof}

The lower bound in \eqref{eq:margin:grad} says that the loss has no flat directions while it is positive: the gradient always has a component of size at least $\gamma^\star\Lhat$ along $\marginu$. This is what drives everything.

\begin{theorem}[Loss, norm, and margin under gradient flow]\label{thm:margin:main}
Under \cref{ass:margin:data}, gradient flow on \eqref{eq:margin:loss} from $w(0)=0$ exists for all $t\ge0$, and for every $t\ge0$:
\begin{enumerate}
\item $\dfrac1{1+t}\le\Lhat(w(t))\le\dfrac1{1+\gamma^{\star2}t}$; in particular every training point is classified correctly for $t>(n-1)/\gamma^{\star2}$;
\item $\dfrac1{\gamma^\star}\log\dfrac{1+\gamma^{\star2}t}n\le\norm{w(t)}\le\dfrac1{\gamma^{\star2}}\log\bigl(1+\gamma^{\star2}t\bigr)$;
\item $\min_iq_i(w(t))\ge\gamma^\star\norm{w(t)}-\log n$;
\item for $t>(n-1)/\gamma^{\star2}$ the normalized margin satisfies
\begin{equation}\label{eq:margin:rate}
  \gamma\bigl(w(t)\bigr)\ge\gamma^\star\Bigl(1-\frac{\log n}{\log\bigl((1+\gamma^{\star2}t)/n\bigr)}\Bigr).
\end{equation}
\end{enumerate}
\end{theorem}

\begin{proof}
The plan is to show that the loss decays at least like $1/t$, because the gradient is at least $\gamma^\star$ times the loss, and then to follow the free energy $N(t)=-\log(n\Lhat(w(t)))$, which by \eqref{eq:margin:softmin} sits below the smallest margin and grows at least $\gamma^\star$ times as fast as the path length of $w$.

\emph{Existence.} The loss \eqref{eq:margin:loss} is infinitely differentiable, so the flow exists on a maximal interval $[0,T)$. Along it $\frac{\dd}{\dd t}\Lhat(w(t))=-\norm{\nabla\Lhat(w(t))}^2\le0$, so $\Lhat(w(t))\le\Lhat(0)=1$, and by \eqref{eq:margin:grad} $\norm{\dot w}\le\Lhat\le1$. Then $\norm{w(t)}\le t$ stays bounded on bounded intervals, the solution cannot blow up, and $T=\infty$. Write $\Lhat(t)$ for $\Lhat(w(t))$.

\emph{Part (i).} By \eqref{eq:margin:grad}, $\gamma^{\star2}\Lhat^2\le\norm{\nabla\Lhat}^2\le\Lhat^2$, so the derivative $-\norm{\nabla\Lhat}^2$ of $\Lhat(t)$ lies between $-\Lhat^2$ and $-\gamma^{\star2}\Lhat^2$. The loss is positive, so we may divide by $\Lhat^2$: the derivative of $1/\Lhat(t)$ lies between $\gamma^{\star2}$ and $1$. Integrating from $0$, where $1/\Lhat=1$,
\[
  1+\gamma^{\star2}t\le\frac1{\Lhat(t)}\le1+t .
\]
If $1+\gamma^{\star2}t>n$, then $\Lhat(t)<1/n$, so every term $e^{-q_i}/n$ is less than $1/n$ and every margin is positive.

\emph{Part (iii).} Differentiate $N(t)=-\log(n\Lhat(t))$ and use the lower bound of \eqref{eq:margin:grad}, $\norm{\nabla\Lhat}\ge\gamma^\star\Lhat$, once:
\[
  \dot N=\frac{\norm{\nabla\Lhat}^2}{\Lhat}\ge\gamma^\star\norm{\nabla\Lhat}=\gamma^\star\norm{\dot w}.
\]
Integrate from $0$ to $t$. The length of the path is at least the distance between its endpoints, so with $N(0)=-\log n$ and $w(0)=0$,
\begin{equation}\label{eq:margin:Ngrowth}
  N(t)\ge-\log n+\gamma^\star\int_0^t\norm{\dot w(s)}\,\dd s\ge\gamma^\star\norm{w(t)}-\log n .
\end{equation}
The left inequality of \eqref{eq:margin:softmin} gives $\min_iq_i(w(t))\ge N(t)$, which proves part (iii).

\emph{Part (ii).} For the lower bound, chain the right inequality of \eqref{eq:margin:softmin}, the left one, and part (i):
\[
  \gamma^\star\norm{w(t)}\ge\min_iq_i(w(t))\ge-\log\bigl(n\Lhat(t)\bigr)\ge\log\frac{1+\gamma^{\star2}t}n .
\]
For the upper bound, $\norm{\dot w}=\norm{\nabla\Lhat}\le\Lhat$ by \eqref{eq:margin:grad}, so by part (i)
\[
  \norm{w(t)}\le\int_0^t\Lhat(s)\,\dd s\le\int_0^t\frac{\dd s}{1+\gamma^{\star2}s}=\frac1{\gamma^{\star2}}\log\bigl(1+\gamma^{\star2}t\bigr).
\]
\emph{Part (iv).} For $t>(n-1)/\gamma^{\star2}$ the lower bound of part (ii) is positive, so $w(t)\ne0$. Dividing part (iii) by $\norm{w(t)}$ gives $\gamma(w(t))\ge\gamma^\star-\log n/\norm{w(t)}$, and the lower bound of part (ii) turns this into \eqref{eq:margin:rate}.
\end{proof}

The time $(n-1)/\gamma^{\star2}$ to zero training error is a cousin of Novikoff's bound: the perceptron makes at most $1/\gamma^{\star2}$ mistakes on data with $\norm{z_i}\le1$ \parencite{Novikoff1962}.\sidenote{Novikoff presented the bound at a symposium held at the Polytechnic Institute of Brooklyn in April 1962; the paper is cited with the year 1962 or 1963.} The factor $n$ comes from the normalization of the loss: zero training error needs $\Lhat<1/n$, not $\Lhat<1$.

\begin{corollary}[Directional convergence]\label[corollary]{cor:margin:direction}
Under the hypotheses of \cref{thm:margin:main}, for $t>(n-1)/\gamma^{\star2}$,
\[
  \Bigl\|\frac{w(t)}{\norm{w(t)}}-\marginu\Bigr\|^2\le\frac{2\log n}{\log\bigl((1+\gamma^{\star2}t)/n\bigr)} .
\]
In particular $w(t)/\norm{w(t)}\to\marginu$ as $t\to\infty$.
\end{corollary}

\begin{proof}
\index{directional convergence}By \eqref{eq:margin:rate} the unit vector $u=w(t)/\norm{w(t)}$ has $\min_i\ip{u,z_i}\ge\gamma^\star-\delta$ with $\delta=\gamma^\star\log n/\log((1+\gamma^{\star2}t)/n)$. Part (iv) of \cref{lem:margin:svm} gives $\norm{u-\marginu}^2\le2\delta/\gamma^\star$.
\end{proof}

The rate of \eqref{eq:margin:rate} is $O(\log n/\log t)$ in the margin, and \textcite{JiTelgarsky2021} prove a rate $O(\ln n/\ln t)$ of convergence to the max-margin direction for gradient descent with a constant step size and show that it is tight in both $n$ and $t$. The rate of the corollary, $O((\log t)^{-1/2})$ in the angle, is not: the direction converges at the rate $O(1/\log t)$ for almost every data set, and at the rate $O(\log\log t/\log t)$ for some degenerate ones \parencite{SoudryEtAl2018}. The heuristic derivation below shows where the faster rate comes from, and what its constant is.

\begin{heuristic}[Heuristic derivation of the asymptotic expansion]
Guess that $w(t)=\marginwhat\log t+\marginwt+o(1)$ for fixed vectors $\marginwhat$ and $\marginwt$, and that the error term may be differentiated, so that $\dot w(t)=\marginwhat/t+o(1/t)$. Insert this into the flow. The point $i$ contributes
\[
  \frac1ne^{-\ip{w(t),z_i}}z_i=\frac1nt^{-\ip{\marginwhat,z_i}}e^{-\ip{\marginwt,z_i}}z_i\bigl(1+o(1)\bigr)
\]
to $\dot w(t)$. A point with $\ip{\marginwhat,z_i}<1$ would contribute a term decaying more slowly than $1/t$, which nothing could balance; so $\ip{\marginwhat,z_i}\ge1$ for all $i$. The points with $\ip{\marginwhat,z_i}>1$ contribute $o(1/t)$, and matching the terms of order $1/t$ leaves
\[
  \marginwhat=\sum_{i\in\marginS}\alpha_iz_i,\qquad \alpha_i=\frac1ne^{-\ip{\marginwt,z_i}}>0,\qquad \marginS=\{i:\ip{\marginwhat,z_i}=1\}.
\]
These are the Karush--Kuhn--Tucker conditions of the convex problem $\min\norm w$ subject to $\ip{w,z_i}\ge1$, which are sufficient for optimality; so $\marginwhat=\marginu/\gamma^\star$ is the hard-margin solution, and $\marginS$ is the set of support vectors. If the support vectors are linearly independent, then the coefficients $\alpha_i$ are determined by $\marginwhat$, and the equations
\begin{equation}\label{eq:margin:wtilde}
  \ip{\marginwt,z_i}=-\log(n\alpha_i)\qquad\text{for }i\in\marginS
\end{equation}
fix the component of $\marginwt$ in the span of the support vectors; the orthogonal component is set during the transient. Three predictions follow. The loss is dominated by the support vectors, and by part (iii) of \cref{lem:margin:svm},
\[
  \Lhat(w(t))\approx\frac1t\sum_{i\in\marginS}\alpha_i=\frac1{\gamma^{\star2}t}.
\]
The norm is $\norm{w(t)}=\log t/\gamma^\star+\ip{\marginu,\marginwt}+o(1)$. And the normalized margin is
\[
  \frac{\gamma(w(t))}{\gamma^\star}=\frac{\log t+a}{\log t+b}+o\Bigl(\frac1{\log t}\Bigr),\qquad a=\min_{i\in\marginS}\ip{\marginwt,z_i},\quad b=\gamma^\star\ip{\marginu,\marginwt},
\]
a gap of order $(b-a)/\log t$, while the angle to $\marginu$ is about $\gamma^\star\norm{P^\perp\marginwt}/\log t$, with $P^\perp$ the projection orthogonal to $\marginu$. A point with $r_i=\ip{\marginwhat,z_i}>1$ perturbs the expansion by a term of order $t^{1-r_i}$, so the expansion is accurate once $t^{r_i-1}$ is large for every non-support vector.
\end{heuristic}

\Textcite{SoudryEtAl2018} made this rigorous for gradient descent with small constant steps: for almost every separable data set the residual $w(t)-\marginwhat\log t$ stays bounded, and in degenerate cases it may grow like $\log\log t$ in directions orthogonal to the support vectors. The heuristic also explains the sharpness of part (i) of \cref{thm:margin:main}: the upper bound $1/(1+\gamma^{\star2}t)$ is attained asymptotically, while the lower bound $1/(1+t)$ is off by the factor $\gamma^{\star2}$.

\section{The solution: steps and networks}\label{sec:margin:steps}

Two extensions bring the model closer to practice: finite steps, and networks. The first is a short induction. The second needs one new idea, a smoothed margin that increases along the flow.

\begin{proposition}[Gradient descent]\label[proposition]{prop:margin:gd}
Under \cref{ass:margin:data}, let $w_{k+1}=w_k-\eta\nabla\Lhat(w_k)$ with $w_0=0$ and $0<\eta\le\frac12$, and write $\Lhat_k=\Lhat(w_k)$. For every $k\ge0$:
\begin{enumerate}
\item $\Lhat_k\le1/(1+\frac12\eta\gamma^{\star2}k)$;
\item $\min_iq_i(w_k)\ge\gamma^\star\norm{w_k}-\log n-\eta\bigl(\gamma^\star\eta+2/\gamma^\star\bigr)$;
\item $\norm{w_k}\ge\frac1{\gamma^\star}\log\bigl((1+\frac12\eta\gamma^{\star2}k)/n\bigr)$.
\end{enumerate}
In particular $\gamma(w_k)\to\gamma^\star$ and $w_k/\norm{w_k}\to\marginu$ as $k\to\infty$.
\end{proposition}

The proof rests on a one-step inequality, which replaces the identity $\frac{\dd}{\dd t}\Lhat=-\norm{\nabla\Lhat}^2$ of the flow.

\begin{lemma}[One step]\label[lemma]{lem:margin:step}
Let $w\in\R^d$, $g=\nabla\Lhat(w)$ and $w^+=w-\eta g$, and suppose that $\eta\Lhat(w)\le\frac12$. Then
\[
  \Lhat(w^+)\le\Lhat(w)-\eta\norm g^2\bigl(1-\eta\Lhat(w)\bigr).
\]
\end{lemma}

\begin{proof}
Put $v=-g=\frac1n\sum_ie^{-q_i(w)}z_i$ and $b_i=\eta\ip{v,z_i}$, so that $q_i(w^+)=q_i(w)+b_i$. By \eqref{eq:margin:grad} and $\norm{z_i}\le1$, $\abs{b_i}\le\eta\norm v\le\eta\Lhat(w)\le\frac12$. On $[-1,1]$ the function $e^{-b}$ lies below $1-b+b^2$: the difference $h(b)=1-b+b^2-e^{-b}$ has $h(0)=h'(0)=0$ and $h''(b)=2-e^{-b}$, which is positive for $b>-\log2$, so $h\ge0$ on $[-\log2,1]$; on $[-1,-\log2]$ the function $h$ is concave, with $h(-1)=3-e>0$ and $h(-\log2)=\log2+(\log2)^2-1>0$, so $h>0$ there too. Hence
\[
  \Lhat(w^+)=\frac1n\sum_ie^{-q_i(w)}e^{-b_i}\le\frac1n\sum_ie^{-q_i(w)}\bigl(1-b_i+b_i^2\bigr).
\]
Expand the right side term by term. The constant terms give $\Lhat(w)$. The linear terms give $-\frac\eta n\sum_ie^{-q_i(w)}\ip{v,z_i}=-\eta\norm v^2$, by the definition of $v$. In the quadratic terms $b_i^2\le\eta^2\norm v^2$, so they give at most $\eta^2\Lhat(w)\norm v^2$.
\end{proof}

\begin{proof}[Proof of \cref{prop:margin:gd}]
The plan is to repeat the proof of \cref{thm:margin:main} with sums instead of integrals, using \cref{lem:margin:step} and absorbing the second-order terms into a constant. Write $g_k=\nabla\Lhat(w_k)$.

\emph{Part (i).} We show by induction that $\Lhat_k$ is nonincreasing, so that $\eta\Lhat_k\le\eta\Lhat_0=\eta\le\frac12$ for all $k$. If $\eta\Lhat_k\le\frac12$, then \cref{lem:margin:step} and \eqref{eq:margin:grad} give
\[
  \Lhat_{k+1}\le\Lhat_k-\frac\eta2\norm{g_k}^2\le\Lhat_k\bigl(1-c\Lhat_k\bigr),\qquad c=\frac{\eta\gamma^{\star2}}2 ,
\]
so $\Lhat_{k+1}\le\Lhat_k$ and the induction proceeds. Since $0\le c\Lhat_k<1$ and $1/(1-x)\ge1+x$ for $x\in[0,1)$, inverting gives $1/\Lhat_{k+1}\ge1/\Lhat_k+c$, and summing from $\Lhat_0=1$ gives $1/\Lhat_k\ge1+ck$.

\emph{Part (ii).} Let $N_k=-\log(n\Lhat_k)$ and $a_k=\eta\norm{g_k}^2(1-\eta\Lhat_k)/\Lhat_k$. \Cref{lem:margin:step} says $\Lhat_{k+1}\le\Lhat_k(1-a_k)$, and $a_k\le\eta\Lhat_k\le\frac12$ by \eqref{eq:margin:grad}; since $-\log(1-x)\ge x$ for $x\in[0,1)$, we get $N_{k+1}-N_k\ge a_k$. Now bound $a_k$ below, using $\norm{g_k}\ge\gamma^\star\Lhat_k$ from \eqref{eq:margin:grad} and then, with $\Delta w_k=w_{k+1}-w_k$, the bound $\norm{\Delta w_k}=\eta\norm{g_k}\le\eta\Lhat_k$:
\[
  a_k\ge\gamma^\star\eta\norm{g_k}\bigl(1-\eta\Lhat_k\bigr)=\gamma^\star\norm{\Delta w_k}-\gamma^\star\eta\Lhat_k\norm{\Delta w_k}\ge\gamma^\star\norm{\Delta w_k}-\gamma^\star\eta^2\Lhat_k^2 .
\]
Summing over $k<K$, using the triangle inequality $\sum_{k<K}\norm{\Delta w_k}\ge\norm{w_K}$ and $N_0=-\log n$,
\[
  N_K\ge\gamma^\star\norm{w_K}-\log n-\gamma^\star\eta^2\sum_{k\ge0}\Lhat_k^2 .
\]
By part (i) and the comparison of a decreasing sum with an integral, $\sum_{k\ge0}\Lhat_k^2\le\sum_{k\ge0}(1+ck)^{-2}\le1+\int_0^\infty(1+cs)^{-2}\dd s=1+1/c$, and $\gamma^\star\eta^2(1+1/c)=\gamma^\star\eta^2+2\eta/\gamma^\star$. Finally $\min_iq_i(w_K)\ge N_K$ by \eqref{eq:margin:softmin}.

\emph{Part (iii).} As in the proof of \cref{thm:margin:main}, \eqref{eq:margin:softmin} and part (i) give $\gamma^\star\norm{w_k}\ge\min_iq_i(w_k)\ge N_k\ge\log((1+ck)/n)$.

\emph{The limits.} Parts (ii) and (iii) give $\gamma(w_k)\ge\gamma^\star-(\log n+\gamma^\star\eta^2+2\eta/\gamma^\star)/\norm{w_k}$ with $\norm{w_k}\to\infty$, and part (iv) of \cref{lem:margin:svm} turns this into directional convergence.
\end{proof}

The constant $2\eta/\gamma^\star$ is the price of discreteness, and it is harmless: it enters only through $1/\norm{w_k}$. The restriction $\eta\le\frac12$ comes from $\Lhat_0=1$ and the descent lemma; larger steps work too, but by a different mechanism (\cref{sec:margin:wrong}).

\paragraph{Homogeneous networks.} Let $f(x;\theta)$ be a network with parameters $\theta\in\R^p$, and suppose that it is \term{positively homogeneous}\index{homogeneous network} of degree $L\ge1$ in $\theta$: $f(x;c\theta)=c^Lf(x;\theta)$ for every $c>0$. A ReLU network without biases and with $L$ layers is homogeneous of degree $L$; a linear classifier has $L=1$.\sidenote{A bias in the first layer can be absorbed by appending a constant coordinate to $x$, which keeps the network homogeneous; biases in later layers cannot.} Training uses the margins $q_i(\theta)=y_if(x_i;\theta)$ in the loss $\Lhat(\theta)=\frac1n\sum_ie^{-q_i(\theta)}$, and the normalized margin is now $\min_iq_i(\theta)/\norm\theta^L$. The free energy of the proof of \cref{thm:margin:main}, divided by the matching power of the norm, becomes the \term{smoothed margin}\index{margin!smoothed}
\[
  \marginsm(\theta)=\frac{-\log\bigl(n\Lhat(\theta)\bigr)}{\norm\theta^L},
\]
which by \eqref{eq:margin:softmin}, whose proof uses only $n\Lhat=\sum_ie^{-q_i}$, satisfies $\marginsm(\theta)\le\min_iq_i(\theta)/\norm\theta^L\le\marginsm(\theta)+\log n/\norm\theta^L$.

\begin{proposition}[The smoothed margin increases; \textcite{LyuLi2020}]\label[proposition]{prop:margin:smoothed}
Let $\theta\mapsto f(x_i;\theta)$ be continuously differentiable and positively homogeneous of degree $L\ge1$ for every $i$, and let $\theta(t)$ solve $\dot\theta=-\nabla\Lhat(\theta)$ on an interval $[t_0,T)$. If $\Lhat(\theta(t_0))<1/n$, then on $[t_0,T)$ the loss stays below $1/n$, the norm $\norm{\theta(t)}$ is strictly increasing, and $\marginsm(\theta(t))$ is nondecreasing.
\end{proposition}

\begin{proof}
The plan is to compute the growth rates of the numerator $\nu=-\log(n\Lhat)$ and of the denominator $\rho^L$, with $\rho=\norm\theta$, and to compare them through Euler's identity for homogeneous functions. Write $\Lhat$, $q_i$, $\nu$, $\rho$ for their values at $\theta(t)$.

\emph{Preliminaries.} Along the flow $\Lhat$ is nonincreasing, so $\Lhat<1/n$ and $\nu>0$ on $[t_0,T)$. Because $e^{-q_i}\le n\Lhat$ for each $i$, every margin satisfies $q_i\ge\nu>0$. The parameter $\theta$ is never $0$, because $f(x;0)=0$ by homogeneity and $\Lhat(0)=1$. Differentiating $c\mapsto f(x;c\theta)=c^Lf(x;\theta)$ at $c=1$ gives Euler's identity $\ip{\theta,\nabla_\theta q_i(\theta)}=Lq_i(\theta)$.

\emph{The norm.} By Euler's identity and $q_i\ge\nu$,
\begin{equation}\label{eq:margin:norm}
  \rho\dot\rho=\ip{\theta,\dot\theta}=\frac1n\sum_ie^{-q_i}\ip{\theta,\nabla q_i}=\frac Ln\sum_ie^{-q_i}q_i\ge L\nu\Lhat>0 .
\end{equation}
\emph{The numerator.} The flow gives $\dot\nu=-\dot{\Lhat}/\Lhat=\norm{\nabla\Lhat}^2/\Lhat$. The component of $\nabla\Lhat$ along $\theta$ is $\ip{\nabla\Lhat,\theta}/\rho=-\dot\rho$, so $\norm{\nabla\Lhat}\ge\dot\rho$ and $\dot\nu\ge\dot\rho^2/\Lhat$.

\emph{The ratio.} Both $\nu$ and $\rho$ are positive, so $\log\marginsm=\log\nu-L\log\rho$ is differentiable, and
\[
  \frac{\dd}{\dd t}\log\marginsm=\frac{\dot\nu}\nu-\frac{L\dot\rho}\rho\ge\frac{\dot\rho^2}{\Lhat\nu}-\frac{L\dot\rho}\rho=\frac{\dot\rho}{\Lhat\nu\rho}\bigl(\rho\dot\rho-L\nu\Lhat\bigr)\ge0,
\]
where the last step uses $\dot\rho>0$ and \eqref{eq:margin:norm} twice.
\end{proof}

For $L=1$ the proposition is a second proof of the margin part of \cref{thm:margin:main}, without its constants (\cref{ex:margin:linearsm}). For networks it is the key step of a theory that we state without its full proof.

\begin{theorem}[\textcite{LyuLi2020}; \textcite{JiTelgarsky2020}; stated with a sketch]\label{thm:margin:lyuli}
Let $f(x;\theta)$ be locally Lipschitz, admit a chain rule for Clarke subgradients, and be positively homogeneous of degree $L\ge1$ in $\theta$, such as a ReLU network without biases, possibly with max-pooling and convolutions; let $\theta(t)$ be a gradient flow of the exponential or logistic loss, in the sense of Clarke subgradients; and suppose that $\Lhat(\theta(t_0))<1/n$ at some time $t_0$. Then:
\begin{enumerate}
\item $\Lhat(\theta(t))=\Theta\bigl(1/(t(\log t)^{2-2/L})\bigr)$ and $\norm{\theta(t)}=\Theta\bigl((\log t)^{1/L}\bigr)$ as $t\to\infty$;
\item every limit point $\bar\theta$ of $\theta(t)/\norm{\theta(t)}$ is, after rescaling, a Karush--Kuhn--Tucker point of
\begin{equation}\label{eq:margin:kkt}
  \min\ \tfrac12\norm\theta^2\qquad\text{subject to}\quad y_if(x_i;\theta)\ge1\quad\text{for all }i ,
\end{equation}
and the normalized margin $\min_iq_i(\theta(t))/\norm{\theta(t)}^L$ converges, its limit being the normalized margin of every such limit point;
\item if the network is definable in an o-minimal structure, which covers ReLU networks with the exponential and logistic losses, then $\theta(t)/\norm{\theta(t)}$ converges \parencite{JiTelgarsky2020}.
\end{enumerate}
\end{theorem}

\begin{proofsketch}
Part (i) follows from \cref{prop:margin:smoothed} by comparing $\nu$ and $\rho$ in the differential inequalities of its proof, as in \cref{thm:margin:main}; the logarithmic factor comes from $\nu\asymp\rho^L$. For part (ii), the Karush--Kuhn--Tucker conditions of \eqref{eq:margin:kkt} ask that $\theta$ be a nonnegative combination $\sum_i\lambda_i\nabla q_i(\theta)$ with $\lambda_i=0$ whenever $q_i(\theta)>1$. The negative gradient $\frac1n\sum_ie^{-q_i}\nabla q_i$ is such a combination with weights that, after normalization, vanish on the points of larger margin; what remains is to show that $\theta$ and $-\nabla\Lhat$ become parallel along a sequence of times, which \textcite{LyuLi2020} deduce from the integrability of $\frac{\dd}{\dd t}\log\marginsm$ and a quantitative notion of approximate Karush--Kuhn--Tucker points. Part (iii) needs a Kurdyka--{\L}ojasiewicz inequality for definable functions, unbounded and nonsmooth, which bounds the length of the curve $\theta/\norm\theta$ in a logarithmic clock. \Textcite{JiTelgarsky2020} also show that, when the network has locally Lipschitz gradients, $-\nabla\Lhat(\theta(t))$ and $\theta(t)$ become aligned. The full proofs are in the two papers.
\end{proofsketch}

For a linear classifier \eqref{eq:margin:kkt} is the hard-margin problem, which is convex, so its Karush--Kuhn--Tucker points are its solutions and the theorem returns \cref{cor:margin:direction}. For networks it is not convex, and part (ii) gives a stationary point, not a maximizer (\cref{sec:margin:wrong}).

\section{Reading the solution}\label{sec:margin:reading}

\Cref{thm:margin:main} and the heuristic expansion make four quantitative predictions. (i) The loss decays like $1/(\gamma^{\star2}t)$, with the constant exactly $1/\gamma^{\star2}$ at late times. (ii) The norm grows like $\log t/\gamma^\star$, with an offset $\ip{\marginu,\marginwt}$ fixed by the support vectors. (iii) The normalized margin approaches $\gamma^\star$ with a gap of order $1/\log t$: halving the gap takes squaring the training time.\sidenote{A gap $c/\log t$ at time $t$ becomes $c/(2\log t)$ at time $t^2$.} (iv) The time scale is $1/\gamma^{\star2}$: the data are fit by time $(n-1)/\gamma^{\star2}$, and the asymptotic regime begins once the non-support vectors are suppressed, when $t^{r_i-1}$ is large for every $r_i=\ip{\marginwhat,z_i}>1$.

\begin{figure}[t]
\begin{fullwidth}
\centering
\begin{tikzpicture}
\begin{groupplot}[group style={group size=3 by 1,horizontal sep=12mm},hplot,width=0.34\linewidth,height=0.31\linewidth,
  xmode=log,xmin=0.1,xmax=1e9,xtick={1,1e3,1e6,1e9},xlabel={time $t$}]
\nextgroupplot[ymin=0,ymax=1.08,ylabel={$\gamma^{\star2}t\,\Lhat(w(t))$},legend style={at={(0.97,0.03)},anchor=south east}]
\addplot[theory] table[x=t,y=lossbd] {figures/data/04-margin-flow.dat};
\addplot[simulation,mark repeat=3] table[x=t,y=lossx] {figures/data/04-margin-flow.dat};
\legend{bound,flow}
\nextgroupplot[ymin=0,ymax=1.04,ylabel={$\gamma(w(t))/\gamma^\star$},legend style={at={(0.03,0.5)},anchor=west}]
\addplot[theory3] table[x=t,y=marglo] {figures/data/04-margin-flow.dat};
\addplot[theory,restrict x to domain=9.2:21] table[x=t,y=margas] {figures/data/04-margin-flow.dat};
\addplot[simulation,mark repeat=5] table[x=t,y=marg] {figures/data/04-margin-flow.dat};
\legend{bound \eqref{eq:margin:rate},expansion,flow}
\nextgroupplot[ymin=0,ymax=80,ylabel={$\norm{w(t)}$},legend style={at={(0.03,0.97)},anchor=north west}]
\addplot[theory3] table[x=t,y=normlo] {figures/data/04-margin-flow.dat};
\addplot[theory,restrict x to domain=9.2:21] table[x=t,y=normas] {figures/data/04-margin-flow.dat};
\addplot[simulation,mark repeat=5] table[x=t,y=norm] {figures/data/04-margin-flow.dat};
\legend{lower bound,expansion,flow}
\end{groupplot}
\end{tikzpicture}
\caption{Gradient flow of the exponential loss on $n=40$ points in the plane: the vectors $z_i=y_ix_i$ are drawn from $\cN(\mu,0.2^2I_2)$ with $\mu=(0.55,0.25)$ and rescaled to largest norm $1$ (seed $31$), so that $\gamma^\star=0.2034$, with two support vectors and smallest non-support ratio $r_i=1.49$. The flow is integrated with a stiff solver in the clock $\log(1+t)$. Left: $\gamma^{\star2}t\,\Lhat$ against the upper bound $\gamma^{\star2}t/(1+\gamma^{\star2}t)$ of \cref{thm:margin:main}; it reaches $0.999$ at $t=10^9$. Middle: the normalized margin, the bound \eqref{eq:margin:rate} and the heuristic expansion $\marginwhat\log t+\marginwt$ with $\marginwt$ from \eqref{eq:margin:wtilde}, drawn for $t\ge10^4$. Right: the norm, the lower bound of part (ii) and the expansion $\log t/\gamma^\star+\ip{\marginu,\marginwt}$.}\label{fig:margin:flow}
\end{fullwidth}
\end{figure}

\Cref{fig:margin:flow} tests these predictions on a data set small enough to see. In the left panel, $\gamma^{\star2}t\Lhat$ climbs to the bound and meets it: the upper bound of part (i) is the late-time law, not just a bound. The climb is slow because early on all forty points carry weight in the loss, and $\Lhat$ exceeds the single-support-vector value. In the middle panel the normalized margin starts at $0.91\gamma^\star$, the margin of the direction $\sum_iz_i$ that the flow takes first, rises through the transient, and then creeps: from $0.975\gamma^\star$ at $t=10^6$ to $0.988\gamma^\star$ at $t=10^9$. The expansion, which has no free parameter, matches it to within $0.003$ from $t=10^5$ on and to within $5\times10^{-4}$ from $t=10^7$ on. The rigorous bound \eqref{eq:margin:rate} is valid but loose, at $0.73\gamma^\star$ at $t=10^9$: it pays $\log n$ for not knowing which points are support vectors. In the right panel the norm is a straight line in $\log t$, with slope $1/\gamma^\star=4.92$ per unit of $\log t$, as the lower bound of part (ii) and the expansion say. The angle to $\marginu$ is still $0.020$ at $t=10^9$; the expansion predicts $\norm{P^\perp\marginwt}/\norm{w(t)}=1.452/71.4=0.020$.

\begin{figure}[t]
\centering
\begin{tikzpicture}
\begin{axis}[hplot,xmode=log,xmin=10,xmax=1e9,xtick={10,1e3,1e5,1e7,1e9},ymin=-0.004,ymax=0.012,
  xlabel={time $t$},ylabel={margin$/\norm\theta^2$},
  yticklabel style={/pgf/number format/fixed,/pgf/number format/precision=3},legend style={at={(0.97,0.97)},anchor=north east}]
\addplot[guide] table[x=t,y=up] {figures/data/04-margin-relu.dat};
\addplot[theory] table[x=t,y=sm] {figures/data/04-margin-relu.dat};
\addplot[simulation,mark repeat=2] table[x=t,y=nm] {figures/data/04-margin-relu.dat};
\addplot[muted,line width=0.5pt,densely dotted] coordinates {(1000,-0.004) (1000,0.012)};
\legend{$\marginsm+\log n/\norm\theta^2$,smoothed margin $\marginsm$,normalized margin}
\end{axis}
\end{tikzpicture}
\caption{A ReLU network $f(x)=\sum_{r=1}^{50}a_r\relu(\ip{w_r,x})$ without biases, homogeneous of degree $L=2$, trained by gradient flow of the exponential loss on $n=30$ points of the unit circle labeled by $\sgn(x_1x_2)$, which no linear classifier separates; weights initialized $\cN(0,0.3^2)$, seed $0$. The flow is integrated by the Euler method in the clock $\dd\tau=\Lhat\,\dd t$ with $\dd\tau=10^{-3}$. The loss drops below $1/n$ at $t\approx10^3$ (dotted line); from then on the smoothed margin never decreases, as \cref{prop:margin:smoothed} asserts, and the normalized margin lies between it and the dashed upper bound.}\label{fig:margin:relu}
\end{figure}

\Cref{fig:margin:relu} shows the network version. Before the loss reaches $1/n$ the smoothed margin is negative and the normalized margin wanders; after, the smoothed margin increases monotonically, from $2.6\times10^{-4}$ at $t=10^3$ to $4.1\times10^{-3}$ at $t=10^9$, and the normalized margin settles near $4.27\times10^{-3}$, squeezed between the two curves as $\log n/\norm\theta^2$ shrinks. The squared norm grows by about $600$ per decade of $t$, linear in $\log t$, and $\Lhat\,t\log t$ stays between $75$ and $100$ over the last three decades, as part (i) of \cref{thm:margin:lyuli} predicts for $L=2$ (\cref{ex:margin:relurate}).

\paragraph{The test loss rises.} Let $(x,y)$ be a test point that the max-margin classifier gets wrong, $\ip{\marginwhat,yx}<0$. Its exponential loss grows like $t^{\abs{\ip{\marginwhat,yx}}}$ and its logistic loss like $\abs{\ip{\marginwhat,yx}}\log t$; so the average test loss eventually increases, while the test error converges to that of the max-margin classifier. This is the decoupling of \cref{obs:margin:late}, and \textcite{SoudryEtAl2018} derived it in this form. The model says that the test loss is the wrong thing to monitor when training past zero error; it does not say that the test error improves.

\paragraph{Weights in the loss do not matter in the limit.} If the loss weights the points, $\Lhat(w)=\sum_ic_ie^{-\ip{w,z_i}}$ with $c_i>0$ and $\sum_ic_i=1$, then \cref{thm:margin:main} holds with $\log n$ replaced by $\log(1/\min_ic_i)$ (\cref{ex:margin:weights}): the limit direction is $\marginu$ whatever the weights. Importance weighting\index{importance weighting} for class imbalance or distribution shift therefore has an effect that fades as training continues past zero error. \Textcite{ByrdLipton2019} observed exactly this in deep networks: the effect of importance weights diminishes over training once the network separates the training data, and $\ell_2$ regularization restores part of it. Here the model's prediction holds for networks.

\paragraph{Limits of the model.} Every statement is about $t\to\infty$ at fixed data. The parameter that sets the time scale is $\gamma^\star$, which can only decrease as points are added, and on hard data sets $1/\gamma^{\star2}$ is large. The logarithm makes the asymptotic regime expensive even when $\gamma^\star$ is not small. In \cref{fig:margin:flow} the transient lasts until about $t=10^5$, some four thousand times the natural unit $1/\gamma^{\star2}\approx24$; with a non-support ratio $r_i=1.15$ instead of $1.49$, we found the expansion still $0.6\%$ off at $t=10^9$.

\section{What the model gets wrong}\label{sec:margin:wrong}

\paragraph{What was thrown away.} \emph{Fixed features.} The model trains a linear readout on frozen inputs; in a network the features move, and the geometry that defines the margin changes during training. \Cref{thm:margin:lyuli} restores feature learning for homogeneous networks, but only at the level of stationary points, and biases, residual connections and normalization layers break homogeneity. \emph{Separability.} Real training sets contain mislabeled and ambiguous points, which a network may or may not fit; for a linear model on nonseparable data the max-margin story applies only to a separable part (\cref{ex:margin:nonsep}). \emph{Infinite time.} The theorems are statements about $t\to\infty$, and the rates are logarithmic. \emph{No regularization.} With weight decay $\lambda$ the norm stops growing at about $\log(1/\lambda)/\gamma^\star$, and the limit becomes a regularized solution, which approaches the max-margin direction only as $\lambda\to0$ \parencite{RossetZhuHastie2003,WeiEtAl2019}. \emph{The optimizer.} Gradient flow, or gradient descent with $\eta\le\frac12$; no stochasticity, momentum, adaptive preconditioning, or learning-rate schedules.

\paragraph{Where it breaks.} \emph{The limit is never reached.} A gap of order $1/\log t$ means that the margin of a trained network reflects its training time and initialization as much as its limit: in \cref{fig:margin:flow} the direction is still $0.02$ away from $\marginu$ after $10^9$ time units, on forty points in the plane. Statements about ``the'' implicit bias of a network trained for $10^5$ steps should be read as statements about a point on a slow curve. \emph{Stationary is not maximal.} \Textcite{VardiShamirSrebro2022} show that the Karush--Kuhn--Tucker points of \cref{thm:margin:lyuli} are in many cases not even local maximizers of the margin, for example for depth-two ReLU networks and depth-two diagonal linear networks, while for deep linear fully connected networks and several other settings they identify conditions that guarantee local or global optimality. The physicist's picture of a quench to the ground state fails: the network can freeze in a metastable state. \emph{Margins do not predict generalization in deep networks.} Margin-based bounds such as the spectrally normalized bound of \textcite{BartlettFosterTelgarsky2017} are far from tight; \textcite{NagarajanKolter2019} found that such norm-based bounds can grow with the training set size while the test error falls, and that uniform convergence itself can be vacuous for gradient descent; and in a large-scale comparison over thousands of networks trained on CIFAR-10 and SVHN, \textcite{JiangEtAl2020} found sharpness-based measures among the most promising predictors of generalization, while several norm-based measures, spectral ones included, correlated poorly or even negatively with it. The model explains the direction of the weights, not the quality of the classifier. \emph{The test error need not improve.} In a linear model on a high-dimensional Gaussian mixture ($x_i=y_i\mu+g_i$ with $\mu=0.25\,\mathbf 1_d$, $g_i\sim\cN(0,I_d)$, $d=200$, $n=100$) we found, for each of four seeds, the exact test error of the flow rising from about $0.05\%$ at zero training error to about $0.13\%$ at $t=4\times10^6$: early-stopped gradient descent is close to the mean-difference classifier, which is near-optimal for isotropic noise, and the max-margin classifier is worse. When $d\gg n$ every training point becomes a support vector and the max-margin classifier coincides with the minimum-norm interpolator of the labels \parencite{MuthukumarEtAl2021}, which brings the double descent of \cref{ch:doubledescent} to classification. \emph{Large steps.}\index{gradient descent!large step size} \Cref{prop:margin:gd} needs $\eta\le\frac12$ to make the loss decrease monotonically. \Textcite{WuEtAl2024} show that gradient descent on logistic regression with separable data and a much larger constant step size $\eta$ first goes through a phase in which the loss oscillates, leaves it within $O(\eta)$ steps, and then decreases at the rate $\tilde O(1/(\eta t))$; with $\eta$ of the order of the step budget $T$, the loss after $T$ steps is $\tilde O(1/T^2)$, much faster than the small-step analysis allows. The monotone picture of this chapter is the wrong picture of the beginning of training at large learning rates, the regime of \cref{ch:stability}. \emph{Other optimizers, other norms.} The Euclidean norm is a property of gradient descent. Steepest descent with respect to another norm converges to the max-margin direction for that norm \parencite{GunasekarEtAl2018a}, and \textcite{ZhangZouCao2024} prove that Adam, in linear logistic regression with diminishing learning rates, converges in direction to the maximum-$\ell_\infty$-margin classifier. Networks trained with Adam, which is most of them, are not described by the $\ell_2$ geometry of this chapter. \emph{Label noise.} With noisy labels, fitting the training set means fitting the noise. Whether the max-margin classifier still generalizes depends on the dimension: \textcite{ChatterjiLong2021} prove that, with sufficient overparameterization, the max-margin classifier trained on noisy data achieves nearly optimal population risk, and \textcite{FreiChatterjiBartlett2022} prove the analogue for two-layer networks trained by gradient descent on the logistic loss. The model, which has no noise and no population, is silent on this benign overfitting.

\paragraph{Fine structure.} \emph{Architecture changes the norm.} For deep linear fully connected networks the end-to-end predictor still converges to the $\ell_2$ max-margin direction, and the layers align \parencite{JiTelgarsky2019a}; for linear convolutional networks of depth $L$ the limit is a stationary point of an $\ell_{2/L}$ bridge penalty on the Fourier coefficients of the predictor \parencite{GunasekarEtAl2018b}; for diagonal linear networks it is an $\ell_1$-type margin, as the potential of \cref{ch:diagonal} suggests (\cref{ex:margin:diagonal}); and for infinitely wide two-layer networks \textcite{ChizatBach2020} characterize the limit as a max-margin classifier in a non-Hilbertian function space, which adapts to low-dimensional structure. \emph{Nonseparable data.} \Textcite{JiTelgarsky2019b} show that on general data gradient descent on logistic regression converges in direction to the max-margin predictor of a maximal separable subset, at a rate of about $\log\log t/\log t$, and recovers the offset on the nonseparable part at a rate of about $(\log t)^2/\sqrt t$. \emph{Faster rates.} The $1/\log t$ is a property of constant steps on a loss whose gradient vanishes. Gradient descent with aggressive step sizes that compensate for the vanishing gradient converges to the max-margin direction at the rate $O(\log t/\sqrt t)$ for the exponential loss \parencite{NacsonEtAl2019}, and \textcite{JiTelgarsky2021} reach $O(1/t)$ for the margin with an aggressive schedule, by reading gradient descent on the primal as mirror descent on the dual distribution $p\in\Delta_n$ of \cref{lem:margin:svm}. \emph{Boosting.} The same mathematics appeared first for boosting: \textcite{RossetZhuHastie2004} showed that boosting with infinitesimal steps follows a regularization path to an $\ell_1$ max-margin classifier, and \textcite{Telgarsky2013} that shrinkage yields approximate max-margin classifiers. \emph{Open.} Rates of directional convergence for nonlinear homogeneous networks are unknown in general, as is a description of which Karush--Kuhn--Tucker point gradient descent selects from a given initialization; and a theory of the margin at the finite training times used in practice, rather than in the limit, is missing.

\section{Model card}\label{sec:margin:card}

\modelcard{After zero training error, the cross-entropy keeps falling like $1/t$, the weights grow like $\log t$, and the classifier keeps changing; test loss and test error decouple.}{Linear classifier on fixed separable data with $\norm{z_i}\le1$, exponential loss, gradient flow or gradient descent with $\eta\le\frac12$ from $w=0$; homogeneous networks as an extension.}{$\Lhat\le1/(1+\gamma^{\star2}t)$; $\norm w\approx\log t/\gamma^\star$; normalized margin within $\gamma^\star\log n/\log(\gamma^{\star2}t/n)$ of $\gamma^\star$; direction converges to the hard-margin SVM; the smoothed margin of a homogeneous network increases once $\Lhat<1/n$.}{$\gamma^{\star2}t\Lhat\to1$; margin gap of order $1/\log t$; time scale $1/\gamma^{\star2}$; importance weights fade; test loss of misclassified points grows like $\log t$.}{Logarithmic rates never reached in practice; stationary points instead of maximizers for ReLU networks; margins do not predict deep-network generalization; large steps and Adam change the dynamics and the norm; label noise.}{Steepest descent and Adam ($\ell_\infty$ margins); linear convolutional and diagonal networks; Lyu--Li and Ji--Telgarsky for homogeneous networks; nonseparable data; normalized and primal--dual methods with fast rates.}

The chapter is the hub of the thread on implicit bias. The diagonal networks of \cref{ch:diagonal} trained on classification data are homogeneous of degree two, and the potential that selects the sparse interpolator there becomes the $\ell_1$ margin here (\cref{ex:margin:diagonal}). The max-margin argument of \cref{ch:grokking} is \cref{thm:margin:lyuli} for a network homogeneous of degree three, with the margin problem solved in Fourier space; and the unconstrained-features model of neural collapse in \cref{ch:collapse} is a max-margin problem for the last layer and the features together, whose solution is a simplex. In all three, ``max-margin solutions as attractors'' means the limit of the dynamics studied here.

Two contrasts are instructive. In \cref{ch:doubledescent} the least-squares loss has a minimizer, and gradient descent from zero converges to the minimum-norm interpolator in finite effective time; here the loss has none, and the bias appears only in the limit of the direction. And the lazy regime of \cref{ch:lazy} is the case in which the features hardly move: then the network trained past zero error behaves like the linear model of this chapter on its tangent features, and converges, in the kernel sense, to a max-margin classifier of the neural tangent kernel.

\begin{exercises}
\item \consolidation (Three points.) Let $z_1=(1,0)$, $z_2=(0,1)$ and $z_3=(0.6,0.6)$. (a) Show that $\gamma^\star=1/\sqrt2$, $\marginu=(1,1)/\sqrt2$, that the support vectors are $z_1$ and $z_2$, and that $\alpha=(1,1,0)$ in \cref{lem:margin:svm}. (b) Show that gradient flow from $0$ stays on the diagonal, $w(t)=s(t)(1,1)$, with $\dot s=\frac13(e^{-s}+0.6e^{-1.2s})$, so that the normalized margin equals $\gamma^\star$ for every $t>0$. (c) Show that $s(t)=\log(t/3)+O(t^{-1/5})$, and compare with \eqref{eq:margin:wtilde}. \hint{in (c), $u=e^s$ satisfies $\dot u=\frac13(1+0.6u^{-1/5})$.}\label{ex:margin:three}

\item \consolidation (The convex bound, and the logistic loss.) Let $\ell$ be the exponential or the logistic loss and $\Lhat(w)=\frac1n\sum_i\ell(\ip{w,z_i})$, which is convex. (a) Show that gradient flow from $0$ satisfies $\Lhat(w(t))\le\Lhat(v)+\norm v^2/(2t)$ for every $v\in\R^d$ and $t>0$. \hint{differentiate $\frac12\norm{w(t)-v}^2$ and use convexity and the monotonicity of $\Lhat(w(t))$.} (b) Taking $v=\marginu\log t/\gamma^\star$, deduce $\Lhat(w(t))\le\bigl(1+(\log t)^2/(2\gamma^{\star2})\bigr)/t$ for $t\ge1$, for both losses. (c) For the logistic loss, show that $e^{-q}\le2\ell(q)$ and $\abs{\ell'(q)}\ge\ell(q)(1-2\ell(q))$ for $q\ge0$. From the first time at which $\Lhat<\frac{\log2}n$, adapt the proof of part (iii) of \cref{thm:margin:main} with $N=-\log(2n\Lhat)$, bounding the error terms by $\int\Lhat^2\,\dd t<\infty$ from (b), to prove that the normalized margin converges to $\gamma^\star$. \label{ex:margin:logistic}

\item \consolidation (The smoothed margin of a linear classifier.) For $L=1$, \cref{prop:margin:smoothed} says that $\marginsm(w(t))=-\log(n\Lhat)/\norm w$ is nondecreasing after the first time $t_0$ at which $\Lhat<1/n$. (a) Show that $\marginsm(w(t))\le\gamma^\star$ for all $t$, and that $\marginsm(w(t))\to\gamma^\star$. (b) Show that monotonicity alone does not give a rate, and recover one from \eqref{eq:margin:Ngrowth}. (c) Check the monotonicity numerically for gradient descent with $\eta=\frac12$ on the data of \cref{fig:margin:flow}. \label{ex:margin:linearsm}

\item \consolidation (The descent lemma, another way.) (a) Show that $\nabla^2\Lhat(w)=\frac1n\sum_ie^{-\ip{w,z_i}}z_iz_i^\top$, so that $0\preceq\nabla^2\Lhat(w)\preceq\Lhat(w)I_d$ under \cref{ass:margin:data}, and that $\Lhat$ has no global Lipschitz gradient. (b) Show that $\Lhat(w+\tau v)\le e^{\tau\norm v}\Lhat(w)$ for $\tau\ge0$, and deduce from Taylor's formula a version of \cref{lem:margin:step} with $1-\eta\Lhat(w)$ replaced by $1-\frac\eta2e^{\eta\Lhat(w)}\Lhat(w)$. (c) For which step sizes does this version give a rate from $w_0=0$? \label{ex:margin:hessian}

\item \consolidation (Norm and loss from the expansion.) Assume the expansion of the heuristic derivation. (a) Show that $\gamma^\star\ip{\marginu,\marginwt}\ge\min_{i\in\marginS}\ip{\marginwt,z_i}$, so that the predicted gap $(b-a)/\log t$ is nonnegative. (b) Show that the angle between $w(t)$ and $\marginu$ is $\norm{P^\perp\marginwt}/\norm{w(t)}+o(1/\log t)$. (c) A misclassified test point has $\ip{\marginwhat,z}<0$. Show that its logistic loss is $\abs{\ip{\marginwhat,z}}\log t+O(1)$. \hint{in (a), $\marginu$ is a convex combination of the support vectors divided by $\gamma^\star$.}\label{ex:margin:expansion}

\item \variation (Other norms.) Let $\norm\cdot$ be a norm on $\R^d$ with dual norm $\norm\cdot_*$, assume $\norm{z_i}_*\le1$, and let $\gamma^\star=\max_{\norm u\le1}\min_i\ip{u,z_i}>0$. Consider an absolutely continuous solution of the steepest-descent flow, in which $\dot w(t)$ maximizes $\ip{-\nabla\Lhat(w),v}$ over $\norm v\le\norm{\nabla\Lhat(w)}_*$. (a) Show that $\frac{\dd}{\dd t}\Lhat=-\norm{\nabla\Lhat}_*^2$ and $\norm{\nabla\Lhat}_*\ge\gamma^\star\Lhat$. (b) Prove parts (i)--(iv) of \cref{thm:margin:main} with $\norm\cdot$ in place of the Euclidean norm. (c) For $\norm\cdot=\norm\cdot_\infty$, show that the flow is sign gradient descent rescaled by $\norm{\nabla\Lhat}_1$, and relate the result to \textcite{ZhangZouCao2024}. (d) Show by an example in $\R^2$ that the max-margin direction for $\norm\cdot_\infty$ need not be unique, so that part (iv) of \cref{lem:margin:svm} fails. \label{ex:margin:norms}

\item \variation (Weighted losses.) Let $\Lhat(w)=\sum_ic_ie^{-\ip{w,z_i}}$ with $c_i>0$ and $\sum_ic_i=1$. Show that \cref{thm:margin:main} holds with $\log n$ replaced by $\log(1/\min_ic_i)$ and $n$ in part (i) replaced by $1/\min_ic_i$. Conclude that the limit direction does not depend on the weights, and that up-weighting a rare class only delays the approach to $\marginu$. Compare with \textcite{ByrdLipton2019}. \label{ex:margin:weights}

\item \variation (Nonseparable data.) Let $z_1=(1,0)$, $z_2=(-1,0)$ and $z_3=(0.6,0.8)$, so that no direction classifies all three points. (a) Show that $\inf\Lhat=\frac23$ and that it is not attained. (b) Show that gradient flow from $0$ satisfies $w_1(t)\to0$ and $w_2(t)-\frac54\log t\to\frac54\log\frac{16}{75}$, so that $w(t)/\norm{w(t)}\to(0,1)$, the max-margin direction of the separable part restricted to the orthogonal complement of the span of $z_1,z_2$. Compare with \textcite{JiTelgarsky2019b}. \hint{show first that $w_1\ge0$ and $w_1(t)=O(1/t)$, then solve for $w_2$ with $w_1=0$.}\label{ex:margin:nonsep}

\item \variation\hard (Diagonal networks and the $\ell_1$ margin.) Let $f(x;u,v)=\ip{u\odot u-v\odot v,x}$, the diagonal network of \cref{ch:diagonal}, homogeneous of degree two. (a) Show that $\min\{\norm u^2+\norm v^2:u\odot u-v\odot v=\beta\}=\norm\beta_1$. (b) Deduce that $(u,v)$ is a global solution of \eqref{eq:margin:kkt} if and only if $\beta=u\odot u-v\odot v$ solves $\min\norm\beta_1$ subject to $\ip{\beta,z_i}\ge1$, the $\ell_1$ max-margin problem. (c) For $z_1=(1,0)$ and $z_2=(0.6,0.6)$, show that the $\ell_2$ max-margin solution is $\marginwhat=(1,\frac23)$, and that the solutions of the $\ell_1$ problem form the segment $\{\beta_1+\beta_2=\frac53,\ 1\le\beta_1\le\frac53\}$, which contains $\marginwhat$. The limit of the diagonal network is not determined by the $\ell_1$ problem alone. (d) By \textcite{VardiShamirSrebro2022}, the Karush--Kuhn--Tucker points of this network need not be local maximizers. Explain why part (b) does not contradict this. \hint{in (a), $u_j^2+v_j^2\ge\abs{u_j^2-v_j^2}$ for each coordinate, with equality when $u_j=0$ or $v_j=0$.}\label{ex:margin:diagonal}

\item \simulation (Gradient descent and its bounds.) On the data of \cref{fig:margin:flow}, run gradient descent on the exponential loss with $\eta\in\{0.25,0.5\}$ for $2\times10^5$ steps and check parts (i) and (ii) of \cref{prop:margin:gd} at every step. The reader should find both bounds satisfied, with $\min_iq_i(w_k)-\gamma^\star\norm{w_k}$ ending near $-0.2$ for both step sizes: the constants $\log n=3.7$ and $\eta(\gamma^\star\eta+2/\gamma^\star)$ are almost entirely slack. Then try $\eta\in\{1,2,5\}$, outside the proposition: on these data the loss still decreases monotonically, and the normalized margin at time $\eta k$ is close to that of the flow. \label{ex:margin:gdsim}

\item \simulation (Rates for a ReLU network.) Use the output of the script of \cref{fig:margin:relu}. (a) Plot $\norm{\theta}^2$ against $\log t$ and check that it is asymptotically linear. (b) Assuming $\marginsm\approx\bar\gamma$ constant and $\min_iq_i\approx\bar\gamma\norm\theta^2$, use \eqref{eq:margin:norm} to derive $u=\bar\gamma\norm\theta^2$ from $e^u/u\approx4\bar\gamma t/n$, and $\Lhat\approx e^{-u}/n$. The reader should find that this predicts $\Lhat(10^9)\approx3.7\times10^{-9}$ with $\bar\gamma=0.00427$, against $3.69\times10^{-9}$ simulated, and that $\Lhat\,t\log t$ varies slowly, as part (i) of \cref{thm:margin:lyuli} says for $L=2$. (c) Repeat with several initialization seeds, keeping the data fixed, and record the normalized margin at $t=10^7$; we found $0.00425$ for five seeds, so on these data the flow appears to select the same stationary point. \label{ex:margin:relurate}

\item \research (Rates for networks.) For a two-layer ReLU network on separable data, bound the rate at which $\theta(t)/\norm{\theta(t)}$ converges, or the normalized margin approaches its limit. \Textcite{LyuLi2020} give rates for the loss and the norm but not for the direction; \textcite{JiTelgarsky2020} prove convergence without a rate. Even an example with a matching lower bound worse than $1/\log t$ would be new to us. \label{ex:margin:rates}

\item \research (Which stationary point?) Characterize, for a two-layer ReLU network and a given data set, the Karush--Kuhn--Tucker points of \eqref{eq:margin:kkt} that gradient flow from small random initialization can reach, and decide when it reaches a global maximizer of the margin. \Textcite{VardiShamirSrebro2022} give settings with and without local optimality; a general criterion is open. \label{ex:margin:which}
\end{exercises}

\notesandrefs
The implicit bias\index{implicit bias} of gradient descent toward the max-margin direction was proved for linear models on separable data by \textcite{SoudryEtAl2018}, whose conference version appeared at ICLR 2018; they gave the expansion $\marginwhat\log t+\rho(t)$, the rates $1/\log t$ and $\log\log t/\log t$, and the CIFAR-10 experiments of \cref{sec:margin:phenomenon}. The ideas were older in boosting, where AdaBoost minimizes an exponential loss by coordinate descent: \textcite{SchapireEtAl1998} tied its test error to the margins, \textcite{RossetZhuHastie2003} showed that regularized solutions of losses with exponential tails approach margin-maximizing separators as the regularization vanishes, \textcite{RossetZhuHastie2004} connected boosting to an $\ell_1$ max-margin classifier, and \textcite{Telgarsky2013} analyzed shrinkage. The max-margin classifier itself goes back to the support vector machines of \textcite{BoserGuyonVapnik1992} and \textcite{CortesVapnik1995}, and the margin as the quantity that controls learning to the perceptron analysis of \textcite{Novikoff1962}.

Our proof of \cref{thm:margin:main} follows the free-energy argument used for homogeneous networks by \textcite{LyuLi2020}, specialized to $L=1$, with the convex geometry of \cref{lem:margin:svm} supplying the direction; the quantitative directional bound of \cref{cor:margin:direction} is a short consequence of part (iv) of that lemma. The primal--dual view, in which the normalized weights $e^{-q_i}/\sum_je^{-q_j}$ are a dual variable converging to $p^\star$, is developed by \textcite{JiTelgarsky2021}, who also show that the rate $\ln n/\ln t$ is tight. \Cref{prop:margin:gd} is a standard descent argument; the constant $\frac12$ in the step size is not optimized.

For networks, \textcite{LyuLi2020} proved the monotonicity of the smoothed margin, the rates of the loss and the norm, and convergence to Karush--Kuhn--Tucker points, for homogeneous networks with locally Lipschitz outputs; \textcite{JiTelgarsky2020} proved directional convergence and alignment using o-minimal geometry. Earlier, \textcite{JiTelgarsky2019a} treated deep linear networks and \textcite{GunasekarEtAl2018a,GunasekarEtAl2018b} worked out the effect of the optimization geometry and of convolutional architectures. \Textcite{WeiEtAl2019} connected the regularized problem with small weight decay to the max-margin problem for networks, a link used in \cref{ch:grokking}.

\paragraph{Translating the literature.} \Textcite{SoudryEtAl2018} absorb the labels into the inputs, writing $\mathbf x_n$ for our $z_i$ and $N$ for our $n$; their loss is the sum $\mathcal L(\mathbf w)=\sum_n\ell(\mathbf w^\top\mathbf x_n)$, which is $n$ times our $\Lhat$; they run gradient descent with step $\eta$, so that their iteration $t$ corresponds to our time $n\eta t$; their $\hat{\mathbf w}$ is our $\marginwhat=\marginu/\gamma^\star$, and their residual $\boldsymbol\rho(t)$ is our $w(t)-\marginwhat\log t$ up to the constant $\marginwhat\log(n\eta)$ that the change of clock produces. \Textcite{LyuLi2020} write $\Phi(\boldsymbol\theta;\mathbf x)$ for our $f(x;\theta)$, $q_n$ for our $q_i$, $L$ for the degree of homogeneity, as we do, and $\rho$ for $\norm\theta$; their loss is again a sum, so for the exponential loss their smoothed margin $\log(1/\mathcal L)/\rho^L$ is our $\marginsm$. The dual variable of \textcite{JiTelgarsky2021} is our vector of normalized weights $e^{-q_i}/\sum_je^{-q_j}$, whose limit is $p^\star$.

\printchapterbib
\end{refsection}


\casepart{Infinite Width}{Two limits make wide networks simple. At initialization a wide network is a Gaussian process whose behaviour in depth has phases, as a magnet does. Under small updates it is a linear model in its parameters, a kernel machine. Both limits are exact, and both lose what makes networks interesting.}{32}
% Chapter 5 of The Hydrogen Atoms of Deep Learning: signal propagation in wide random networks.
\newcommand{\signalpropV}{\mathcal{V}}
\newcommand{\signalpropC}{\mathcal{C}}
\newcommand{\signalpropG}{\widehat{G}}

\casechapter{Order and Chaos at Initialization}{Signal propagation in wide random networks}{Infinitely wide random networks}{ch:signalprop}
\begin{refsection}[chapters/05-signalprop.bib]

\begin{chaptersummary}
A deep network trains only if its random initialization is scaled with care, and the deepest trainable network is set by a depth scale that diverges on a curve in the plane of weight and bias variances. Infinitely wide random networks explain this. We prove that their pre-activations are Gaussian processes whose covariance obeys a layer-to-layer recursion, and that the recursion has two phases: an ordered phase, where all inputs become perfectly correlated at a geometric rate $\chi_1$ per layer, and a chaotic phase, where nearby inputs decorrelate; the depth scale is $-1/\log$ of the slope of the correlation map at its stable fixed point, and it diverges like $1/\abs{\chi_1-1}$. At He initialization ReLU networks sit exactly at the boundary and decorrelate only polynomially, $1-c^\ell\sim 9\pi^2/(2\ell^2)$. The model is mean-field theory: it is exact when width is infinite at fixed depth, it has no training and no architecture, and at finite width its fluctuations grow with depth over width, so that criticality on average says little about the typical network once $L$ is comparable to $m$.
\end{chaptersummary}

\section{The phenomenon: a narrow window of trainable initializations}\label{sec:signalprop:phenomenon}

Before a network learns anything, its weights are random, and signals must already pass through it: forward, so that the output depends on the input, and backward, so that every layer receives a gradient. Each layer multiplies the signal by a random matrix and bends it through a nonlinearity. Composing $L$ such maps either shrinks or amplifies, and generically it does so exponentially in $L$.

For gradients through time in recurrent networks the problem was identified early. \Textcite{Hochreiter1991}, in a diploma thesis written in German, analyzed how the error signal propagated backward through a recurrent network decays or blows up exponentially with the number of time steps, and \textcite{BengioEtAl1994} showed that a recurrent network that stores information robustly must have vanishing gradients over long time lags. These are the \term{vanishing and exploding gradients}\index{vanishing and exploding gradients}.\sidenote{An unrolled recurrent network is a deep feedforward network whose layers share their weights, so the two problems are one problem, with depth playing the role of time.}

For depth the practice came first. \Textcite{GlorotBengio2010} found that deep sigmoid networks with the then-standard initialization drove their top hidden layer into saturation early in training, and proposed to draw the weights of a layer with $m_{\mathrm{in}}$ inputs and $m_{\mathrm{out}}$ outputs with variance $2/(m_{\mathrm{in}}+m_{\mathrm{out}})$, a compromise between keeping the variance of the activations and that of the back-propagated gradients constant in a linear network. \Textcite{HeEtAl2015} redid the calculation for rectifiers: a ReLU passes on half of the second moment of its input, so the weight variance must be $2/m_{\mathrm{in}}$ to keep the variance of the responses constant. With this initialization they trained a $30$-layer rectifier network from scratch, while with the variance of \textcite{GlorotBengio2010} the same network stalled.

\Textcite{PooleEtAl2016} turned the variance-counting arguments into a theory. Following two inputs through a deep random network of large width, they found two phases. In the \term{ordered phase}\index{ordered phase} every pair of inputs is mapped to nearly the same representation; in the \term{chaotic phase}\index{chaotic phase} nearby inputs are driven apart, and a smooth curve of inputs is mapped to a curve whose length and curvature grow exponentially with depth. \Textcite{SchoenholzEtAl2017} read off depth scales from the same theory: the number of layers over which the correlation between two inputs relaxes, and over which gradient norms grow or decay. The correlation depth scale diverges on the boundary between the phases, the \term{edge of chaos}\index{edge of chaos}, and their experiments put the limit of trainability at a fixed multiple of it.

\begin{observation}[\textcite{SchoenholzEtAl2017}]\label[observation]{obs:signalprop:depth}
Fully connected $\tanh$ networks with Gaussian weights of variance $\sigma_w^2/m$ and biases of variance $\sigma_b^2$, were trained by stochastic gradient methods on MNIST and on CIFAR-10 for many values of the depth $L$ and of $(\sigma_w^2,\sigma_b^2)$. Networks trained when $L\lesssim6\xi_c$ and failed when $L$ was much larger, where $\xi_c$ is the correlation depth scale of the infinitely wide random network with the same variances. Close to the critical line, where $\xi_c$ diverges, the trainable depth grew accordingly.
\end{observation}

How robust is this? The divergence of the trainable depth at the critical line has been reproduced many times; \textcite{XiaoEtAl2018} trained vanilla convolutional networks with $10{,}000$ layers, without residual connections or normalization, by initializing at criticality with orthogonal convolution kernels. The constant $6$ is empirical and depends on the setting, the training budget, and the criterion for success.\sidenote{The theory says nothing about the constant: it predicts a depth scale, not a training outcome.} Trainable is also not the same as trains well: \textcite{PenningtonEtAl2017} found that critical ReLU networks learn more slowly than critical sigmoidal networks with orthogonal weights, because criticality fixes the mean squared singular value of the input--output Jacobian and not its spread. Finally, modern networks escape the constraint by architecture rather than by tuning: residual connections and normalization layers make depth easy for reasons that the plain model does not contain (\cref{sec:signalprop:wrong}).

The questions for the model are therefore: What computes $\xi_c$? Why is there a sharp boundary, and why does the depth scale diverge on it? Where does the ReLU network at He initialization sit?

\section{The model: infinitely wide random networks}\label{sec:signalprop:model}

\begin{assumption}[Random fully connected network]\label[assumption]{ass:signalprop:net}
The input is $x\in\R^d$, the widths are $m_0=d$ and $m_1,\dots,m_L$, and the activation $\phi\colon\R\to\R$ is continuous and polynomially bounded: $\abs{\phi(t)}\le C_\phi(1+\abs t)^{k_\phi}$ for constants $C_\phi,k_\phi$. The \term{pre-activations}\index{pre-activation} are
\[
  h^{1}(x)=W^{1}x+b^{1},\qquad h^{\ell+1}(x)=W^{\ell+1}\phi\bigl(h^\ell(x)\bigr)+b^{\ell+1}\quad(1\le\ell<L),
\]
with $\phi$ applied coordinatewise, and all weights and biases are independent, with $W^{\ell}_{ij}\sim\cN(0,\sigma_w^2/m_{\ell-1})$ and $b^\ell_i\sim\cN(0,\sigma_b^2)$, where $\sigma_w^2>0$ and $\sigma_b^2\ge0$.
\end{assumption}

We write $a^0(x)=x$ and $a^\ell(x)=\phi(h^\ell(x))$ for the post-activations, so that $h^{\ell+1}=W^{\ell+1}a^\ell+b^{\ell+1}$ for $0\le\ell<L$. The widths are the hidden widths $m$ of the global notation, one per layer; the limit of the model sends them to infinity at fixed depth.

\paragraph{The kernel recursion.} The infinitely wide network is described by one deterministic function per layer, its \term{kernel}\index{kernel!of a random network}:
\begin{equation}\label{eq:signalprop:kernel}
  K^1(x,x')=\sigma_b^2+\sigma_w^2\frac{\ip{x,x'}}d,\qquad K^{\ell+1}(x,x')=\sigma_b^2+\sigma_w^2\,\E\bigl[\phi(z)\phi(z')\bigr],
\end{equation}
where $(z,z')$ is a centered Gaussian pair with $\E z^2=K^\ell(x,x)$, $\E z'^2=K^\ell(x',x')$, and $\E zz'=K^\ell(x,x')$. The polynomial bound makes the expectation finite. The two \term{order parameters}\index{order parameters!of signal propagation} are the length $q^\ell(x)=K^\ell(x,x)$, the limit of $\frac1{m_\ell}\norm{h^\ell(x)}^2$, and the correlation
\[
  c^\ell(x,x')=\frac{K^\ell(x,x')}{\sqrt{q^\ell(x)\,q^\ell(x')}},
\]
the limit of the cosine of the angle between $h^\ell(x)$ and $h^\ell(x')$. With $z\sim\cN(0,1)$, and $(u,v)$ standard Gaussians with correlation $c$, the length obeys the one-dimensional \term{length map}\index{length map}
\begin{equation}\label{eq:signalprop:length}
  q^{\ell+1}=\signalpropV(q^\ell),\qquad \signalpropV(q)=\sigma_b^2+\sigma_w^2\,\E\bigl[\phi(\sqrt q\,z)^2\bigr],
\end{equation}
and if both inputs sit at a fixed point $\bar q>0$ of $\signalpropV$, then the correlation obeys the \term{correlation map}\index{correlation map}
\begin{equation}\label{eq:signalprop:cmap}
  c^{\ell+1}=\signalpropC(c^\ell),\qquad \signalpropC(c)=\frac{\sigma_b^2+\sigma_w^2\,\E\bigl[\phi(\sqrt{\bar q}\,u)\phi(\sqrt{\bar q}\,v)\bigr]}{\bar q}.
\end{equation}
The fixed point is written $\bar q$, not $q^\star$, because a superscript star marks an optimum in this book. The quantity that will decide everything is the slope
\begin{equation}\label{eq:signalprop:chi}
  \chi_1=\sigma_w^2\,\E\bigl[\phi'(\sqrt{\bar q}\,z)^2\bigr].
\end{equation}

\begin{dictionary}
width $m_\ell\to\infty$ at fixed depth & width much larger than depth, $L/m\ll1$\\
weight and bias variances $\sigma_w^2,\sigma_b^2$ & the initialization scheme (LeCun, Glorot, He)\\
length $q^\ell$ & mean squared pre-activation of layer $\ell$\\
correlation $c^\ell(x,x')$ & cosine similarity of the representations of two inputs\\
fixed point $\bar c$ of $\signalpropC$ & how similar all inputs look to a deep layer\\
slope $\chi_1$ & growth factor per layer of squared gradient norms\\
depth scale $\xi$ & depth beyond which training fails (about $6\xi$ for $\tanh$ in \cref{obs:signalprop:depth})\\
critical line $\chi_1=1$ & the family of initializations that train deep networks\\
kernel $K^L$ & prior covariance of a wide Bayesian network; the NNGP kernel\\
random weights & weights at initialization only; training moves them\\
\end{dictionary}

\paragraph{What was thrown away, and why first.} The network is never trained: trainability is decided at initialization, at least in the sense of \cref{obs:signalprop:depth}, because a signal that does not reach the output cannot be learned from. The width is infinite, which turns sums over neurons into expectations and makes the theory deterministic, as the field of a magnet becomes deterministic when each spin interacts with infinitely many others. The architecture is a plain stack of dense layers, because it is the case in which depth is hardest. The data enter only through the inner products $\ip{x,x'}/d$, so two inputs are as good as a dataset.

\paragraph{What one had to notice.} Three things. Conditionally on one layer, the next is exactly Gaussian, because it is linear in Gaussian weights; only its covariance is random, and that covariance is an average over the neurons of the previous layer. Two inputs suffice: their lengths and their correlation are closed under the layer map, like the magnetization of a mean-field magnet. And the correlation map is a power series with nonnegative coefficients, so its behavior on $[0,1]$ is decided by its slope at the fixed point $c=1$.

\section{The solution: Gaussian processes at infinite width}\label{sec:signalprop:gp}

Let $\cX$ be a finite set of inputs, and write $h^\ell_i(\cX)=(h^\ell_i(x))_{x\in\cX}\in\R^{\cX}$ for the values of neuron $i$ of layer $\ell$ on these inputs. Let $\cF_\ell$ be the $\sigma$-algebra generated by $W^1,b^1,\dots,W^\ell,b^\ell$, with $\cF_0$ trivial.

\begin{lemma}[Exact conditional Gaussianity]\label[lemma]{lem:signalprop:conditional}
Under \cref{ass:signalprop:net}, for $0\le\ell<L$, conditionally on $\cF_\ell$, the vectors $h^{\ell+1}_1(\cX),\dots,h^{\ell+1}_{m_{\ell+1}}(\cX)$ are independent with law $\cN(0,\signalpropG^\ell)$, where
\[
  \signalpropG^\ell(x,x')=\sigma_b^2+\frac{\sigma_w^2}{m_\ell}\sum_{j=1}^{m_\ell}a^\ell_j(x)\,a^\ell_j(x'),\qquad x,x'\in\cX.
\]
In particular $\signalpropG^0=K^1$ on $\cX$, and the vectors $h^1_i(\cX)$ are independent with law $\cN(0,K^1)$.
\end{lemma}

\begin{proof}
The post-activations $a^\ell(x)$ are $\cF_\ell$-measurable, and the row $(W^{\ell+1}_{i1},\dots,W^{\ell+1}_{im_\ell},b^{\ell+1}_i)$ is a Gaussian vector independent of $\cF_\ell$ and of the other rows. Conditionally on $\cF_\ell$, the vector $h^{\ell+1}_i(\cX)$ is therefore a fixed linear image of this row, hence centered Gaussian, and the vectors for different $i$ are independent. Since the entries of the row are independent with variances $\sigma_w^2/m_\ell$ and $\sigma_b^2$, the conditional covariance of $h^{\ell+1}_i(x)=\sum_jW^{\ell+1}_{ij}a^\ell_j(x)+b^{\ell+1}_i$ and $h^{\ell+1}_i(x')$ is $\frac{\sigma_w^2}{m_\ell}\sum_ja^\ell_j(x)a^\ell_j(x')+\sigma_b^2$. For $\ell=0$, $a^0(x)=x$ and $m_0=d$ give $K^1$.
\end{proof}

The lemma is exact at every width. It says that a layer acts on a finite set of inputs through the $\abs\cX\times\abs\cX$ matrix $\signalpropG^\ell$ alone: in law, $h^{\ell+1}(\cX)$ is an $m_{\ell+1}\times\abs\cX$ matrix of independent standard Gaussians times $(\signalpropG^\ell)^{1/2}$. On two inputs, a network of width $m$ is therefore a Markov chain on $m\times2$ matrices, which can be simulated in time $O(mL)$ instead of $O(m^2L)$; all the finite-width simulations of this chapter use this reduction.\sidenote{We checked the reduction against simulations with full $m\times m$ weight matrices at $m=200$; the averaged correlations agree to within the Monte Carlo error.}

\begin{theorem}[Gaussian-process limit, sequential widths]\label{thm:signalprop:gp}\index{Gaussian process limit}
Under \cref{ass:signalprop:net}, let $\cX$ be a finite set of inputs, $1\le\ell\le L$, and $J\ge1$, with $m_\ell\ge J$. In the iterated limit in which $m_1\to\infty$, then $m_2\to\infty$, and so on up to $m_{\ell-1}\to\infty$, the law of $(h^\ell_1(\cX),\dots,h^\ell_J(\cX))$ converges weakly to the law of $J$ independent $\cN(0,K^\ell)$ vectors, with $K^\ell$ restricted to $\cX$. That is, for every bounded continuous $F\colon(\R^{\cX})^J\to\R$,
\[
  \lim_{m_{\ell-1}\to\infty}\cdots\lim_{m_1\to\infty}\E F\bigl(h^\ell_1(\cX),\dots,h^\ell_J(\cX)\bigr)=\E F(z_1,\dots,z_J),\qquad z_1,\dots,z_J\overset{\text{i.i.d.}}{\sim}\cN(0,K^\ell).
\]
For $\ell=1$ no limit is needed, and the law of $(h^\ell_1(\cX),\dots,h^\ell_J(\cX))$ does not depend on $m_\ell,\dots,m_L$.
\end{theorem}

\begin{proof}
The plan is induction on $\ell$: by \cref{lem:signalprop:conditional}, layer $\ell+1$ is a Gaussian whose covariance is a continuous function of the neurons of layer $\ell$; the induction hypothesis makes those neurons i.i.d.\ Gaussian, and the law of large numbers makes the covariance deterministic.

\emph{Two facts.} First, centered Gaussian laws depend continuously on their covariance: if $\Sigma_n\to\Sigma$, then the characteristic functions $\exp(-\frac12t^\top\Sigma_nt)$ converge to $\exp(-\frac12t^\top\Sigma t)$, and L\'evy's continuity theorem gives $\cN(0,\Sigma_n)\to\cN(0,\Sigma)$ weakly. Second, the last claim of the theorem holds because $h^\ell$ is a function of $W^1,b^1,\dots,W^\ell,b^\ell$, and the first $J$ neurons of layer $\ell$ use only the first $J$ rows of $W^\ell$ and $b^\ell$.

\emph{Base case.} By \cref{lem:signalprop:conditional} with $\ell=0$, the vectors $h^1_i(\cX)$ are i.i.d.\ $\cN(0,K^1)$ at every width.

\emph{Induction step.} Assume the statement for layer $\ell$ and every number of neurons, and fix $J$, $F$, and $m_\ell$. For $y=(y_1,\dots,y_{m_\ell})\in(\R^{\cX})^{m_\ell}$ define the $\cX\times\cX$ matrix $G(y)=\sigma_b^2\mathbf 1\mathbf 1^\top+\frac{\sigma_w^2}{m_\ell}\sum_j\phi(y_j)\phi(y_j)^\top$, and let $\Phi_F(y)=\E F(g_1,\dots,g_J)$ with $g_1,\dots,g_J$ i.i.d.\ $\cN(0,G(y))$. The function $\Phi_F$ is bounded by $\sup\abs F$, and it is continuous: $\phi$ is continuous, so $G$ is continuous, and the first fact applies. By \cref{lem:signalprop:conditional}, conditionally on $\cF_\ell$ the first $J$ neurons of layer $\ell+1$ are i.i.d.\ $\cN(0,G(h^\ell_1(\cX),\dots,h^\ell_{m_\ell}(\cX)))$, so
\[
  \E F\bigl(h^{\ell+1}_1(\cX),\dots,h^{\ell+1}_J(\cX)\bigr)=\E\,\Phi_F\bigl(h^\ell_1(\cX),\dots,h^\ell_{m_\ell}(\cX)\bigr).
\]
The induction hypothesis with $m_\ell$ neurons in place of $J$, applied to the bounded continuous function $\Phi_F$, shows that after the limits $m_1\to\infty,\dots,m_{\ell-1}\to\infty$ the right side becomes $\E\,\Phi_F(z_1,\dots,z_{m_\ell})$ with $z_j$ i.i.d.\ $\cN(0,K^\ell)$ (for $\ell=1$ this is the base case, with no limit).

Now let $m_\ell\to\infty$. Each entry of $\phi(z_j)\phi(z_j)^\top$ is integrable, since $\abs{\phi(s)\phi(t)}\le\frac12(\phi(s)^2+\phi(t)^2)$ and $\phi(z_j(x))^2$ is bounded by a polynomial in a Gaussian variable. By the strong law of large numbers, $G(z_1,\dots,z_{m_\ell})\to\sigma_b^2\mathbf 1\mathbf 1^\top+\sigma_w^2\E[\phi(z)\phi(z)^\top]$ almost surely, and by \eqref{eq:signalprop:kernel} the limit is $K^{\ell+1}$ on $\cX$. By the first fact, $\Phi_F(z_1,\dots,z_{m_\ell})\to\E F(z'_1,\dots,z'_J)$ almost surely, with $z'_i$ i.i.d.\ $\cN(0,K^{\ell+1})$, and bounded convergence gives the same limit for the expectation. This is the statement for layer $\ell+1$.
\end{proof}

The proof shows more: in the same iterated limit, now including $m_\ell\to\infty$, the empirical kernel $\signalpropG^\ell$ converges in probability to $K^{\ell+1}$, and in particular $\frac1{m_\ell}\norm{h^\ell(x)}^2\to q^\ell(x)$ (\cref{ex:signalprop:empirical}). This is the precise sense of \emph{mean-field exactness}: each neuron of layer $\ell+1$ feels the other neurons of layer $\ell$ only through an average over infinitely many of them, and the average does not fluctuate.

Taking the widths to infinity one at a time is a convenience of the proof. The theorem holds when they grow together.

\begin{theorem}[Simultaneous limits; \textcite{MatthewsEtAl2018}, \textcite{Hanin2023}]\label{thm:signalprop:simultaneous}
The conclusion of \cref{thm:signalprop:gp} holds when $\min(m_1,\dots,m_{\ell-1})\to\infty$ in an arbitrary manner, for activations satisfying a growth condition: \textcite{MatthewsEtAl2018} assume a linear envelope $\abs{\phi(t)}\le c+M\abs t$, and \textcite{Hanin2023} allows general nonlinearities and non-Gaussian weights under moment conditions.
\end{theorem}

For one hidden layer the statement is the central limit theorem, and \textcite{Neal1996} used it to show that Bayesian networks with a Gaussian prior on many hidden units define Gaussian-process priors; \textcite{LeeEtAl2018} and \textcite{MatthewsEtAl2018} extended it to depth, and \textcite{Yang2019} to most modern architectures.

\paragraph{The length map.} The kernel recursion needs a fixed length to define the correlation map. For odd activations such as $\tanh$ the length map has a unique attracting fixed point, and the proof is a Fourier argument.

\begin{proposition}[Length map of $\tanh$]\label[proposition]{prop:signalprop:length}
Let $\phi=\tanh$. The length map $\signalpropV$ of \eqref{eq:signalprop:length} is increasing and strictly concave on $[0,\infty)$, with $\signalpropV(0)=\sigma_b^2$, $\signalpropV'(0)=\sigma_w^2$, and $\signalpropV<\sigma_b^2+\sigma_w^2$.
\begin{enumerate}
\item If $\sigma_b^2>0$, then $\signalpropV$ has a unique fixed point $\bar q>0$, and $q^\ell\to\bar q$ monotonically from every $q^1\ge0$.
\item If $\sigma_b^2=0$ and $\sigma_w^2\le1$, then $q^\ell\to\bar q=0$ from every $q^1\ge0$.
\item If $\sigma_b^2=0$ and $\sigma_w^2>1$, then $\signalpropV$ has a unique positive fixed point $\bar q$, and $q^\ell\to\bar q$ monotonically from every $q^1>0$.
\end{enumerate}
In cases (i) and (iii), $\signalpropV'(\bar q)<1$, so the convergence is geometric.
\end{proposition}

\begin{proof}
The plan is to write $1-\tanh^2=\operatorname{sech}^2$ as a mixture of cosines with positive weights, which turns $q\mapsto\E[\operatorname{sech}^2(\sqrt qz)]$ into a mixture of decaying exponentials.

\emph{The representation.} The function $\operatorname{sech}^2$ is integrable, with Fourier transform $\int_\R\operatorname{sech}^2(t)\cos(\omega t)\dd t=g(\omega)$, where $g(\omega)=\pi\omega/\sinh(\pi\omega/2)$ for $\omega\ne0$ and $g(0)=2$ (\cref{ex:signalprop:fourier}). The function $g$ is positive and decays exponentially, so Fourier inversion gives $\operatorname{sech}^2(t)=\frac1{2\pi}\int_\R g(\omega)\cos(\omega t)\dd\omega$. Since $\E\cos(\omega\sqrt q z)=e^{-q\omega^2/2}$, Fubini's theorem gives
\[
  f(q)\defeq\E\bigl[\tanh^2(\sqrt q\,z)\bigr]=1-\frac1{2\pi}\int_\R g(\omega)\,e^{-q\omega^2/2}\dd\omega .
\]
\emph{Monotonicity and concavity.} The functions $\omega^2g(\omega)$ and $\omega^4g(\omega)$ are integrable, so we may differentiate under the integral for $q\ge0$:
\[
  f'(q)=\frac1{4\pi}\int_\R\omega^2g(\omega)e^{-q\omega^2/2}\dd\omega>0,\qquad f''(q)=-\frac1{8\pi}\int_\R\omega^4g(\omega)e^{-q\omega^2/2}\dd\omega<0 .
\]
At $q=0$, differentiating the inversion formula twice at $t=0$ gives $\frac1{2\pi}\int\omega^2g=-(\operatorname{sech}^2)''(0)=2$, because $\operatorname{sech}^2t=1-t^2+O(t^4)$; so $f'(0)=1$. With $\signalpropV=\sigma_b^2+\sigma_w^2f$ and $0\le f<1$, this proves the first sentence.

\emph{Case (ii).} For $q>0$, strict concavity gives $\signalpropV(q)<\signalpropV(0)+\signalpropV'(0)q=\sigma_w^2q\le q$. So $q^\ell$ decreases, its limit is a fixed point of the continuous map $\signalpropV$, and $0$ is the only one.

\emph{Cases (i) and (iii).} Let $g_0(q)=\signalpropV(q)-q$, strictly concave, with $g_0(q)\to-\infty$ because $\signalpropV$ is bounded. In case (i), $g_0(0)=\sigma_b^2>0$; in case (iii), $g_0(0)=0$ and $g_0'(0)=\sigma_w^2-1>0$; in both cases $g_0$ is positive somewhere on $(0,\infty)$ and has a zero $\bar q>0$ there by the intermediate value theorem. If $0<q_1<q_2$ were two zeros, then strict concavity on $[0,q_2]$ would give $g_0(q_1)>(1-q_1/q_2)g_0(0)\ge0$, a contradiction; so $\bar q$ is unique, $g_0>0$ on $(0,\bar q)$ and $g_0<0$ on $(\bar q,\infty)$. For $0<q<\bar q$, monotonicity gives $q<\signalpropV(q)<\signalpropV(\bar q)=\bar q$, and for $q>\bar q$ it gives $\bar q<\signalpropV(q)<q$: the iterates move monotonically toward $\bar q$, and their limit, a positive fixed point, is $\bar q$. In case (i) the start $q^1=0$ is mapped to $\sigma_b^2>0$. Finally, strict concavity gives $g_0'(\bar q)<\bigl(g_0(\bar q)-g_0(0)\bigr)/\bar q\le0$, that is, $\signalpropV'(\bar q)<1$.
\end{proof}

In case (ii) the network forgets the scale of its input: every length decays to zero, slowly when $\sigma_w^2=1$, since $f(q)=q-2q^2+O(q^3)$ gives $q^\ell\sim1/(2\ell)$. ReLU is the opposite extreme: $\E[\relu(\sqrt qz)^2]=q/2$, so with $\sigma_b=0$ the length map is the multiplication by $\sigma_w^2/2$, and every length is a fixed point exactly when $\sigma_w^2=2$.

\section{The solution: order, chaos, and the critical line}\label{sec:signalprop:phase}

Fix a fixed point $\bar q>0$ of the length map and two inputs with $q^1(x)=q^1(x')=\bar q$, so that $q^\ell=\bar q$ in every layer. Write $\bar\phi(t)=\phi(\sqrt{\bar q}\,t)$ and expand it in the normalized Hermite polynomials, with coefficients $a_k=\E[\bar\phi(z)h_k(z)]$ (\cref{lem:toolkit:hermite}). By \cref{cor:toolkit:dual}, the correlation map \eqref{eq:signalprop:cmap} is the power series
\begin{equation}\label{eq:signalprop:series}
  \signalpropC(c)=\sum_{k\ge0}\beta_kc^k,\qquad \beta_0=\frac{\sigma_b^2+\sigma_w^2a_0^2}{\bar q},\quad \beta_k=\frac{\sigma_w^2a_k^2}{\bar q}\ (k\ge1),
\end{equation}
with nonnegative coefficients, convergent on $[-1,1]$, with $\sum_k\beta_k=\signalpropC(1)=\signalpropV(\bar q)/\bar q=1$.\sidenote{The dual-activation view of \textcite{DanielyEtAl2016} is exactly this: a layer acts on correlations through the power series of its activation's squared Hermite coefficients, and composition of layers is composition of power series.}

\begin{theorem}[Order and chaos]\label{thm:signalprop:phase}\index{order--chaos transition}
Under \cref{ass:signalprop:net}, let $\phi$ be absolutely continuous and not constant, with $\phi'$ polynomially bounded, let $\bar q>0$ be a fixed point of $\signalpropV$, and let $\chi_1$ be given by \eqref{eq:signalprop:chi}.
\begin{enumerate}
\item The map $\signalpropC$ is continuous on $[-1,1]$, maps $[-1,1]$ into itself and $[0,1]$ into $[0,1]$, is nondecreasing and convex on $[0,1]$, satisfies $\signalpropC(1)=1$, and has left derivative $\signalpropC'(1)=\chi_1$ at $1$. For every $c\in[-1,1]$,
\begin{equation}\label{eq:signalprop:contraction}
  0\le1-\signalpropC(c)\le\chi_1\,(1-c).
\end{equation}
\item \emph{Ordered phase, $\chi_1<1$.} For every $c^1\in[-1,1]$, $1-c^\ell\le2\chi_1^{\ell-1}$, so $c^\ell\to1$. If $c^1\in(-1,1)$, then $c^\ell<1$ for every $\ell$ and $\frac1\ell\log(1-c^\ell)\to\log\chi_1$.
\item \emph{Critical line, $\chi_1=1$.} If $\phi$ is not affine and $c^1\in[0,1)$, then $c^\ell$ increases to $1$, and $(1-c^{\ell+1})/(1-c^\ell)\to1$: the decay is slower than every geometric rate.
\item \emph{Chaotic phase, $\chi_1>1$.} The map $\signalpropC$ has a unique fixed point $\bar c\in[0,1)$, and $\signalpropC'(\bar c)<1$. For every $c^1\in[0,1)$, $c^\ell\to\bar c$ monotonically, and $\abs{c^\ell-\bar c}=O(\rho^\ell)$ for every $\rho\in(\signalpropC'(\bar c),1)$.
\end{enumerate}
\end{theorem}

\begin{proof}
The plan is to read every claim off the power series \eqref{eq:signalprop:series}, using the identity
\begin{equation}\label{eq:signalprop:telescope}
  1-\signalpropC(c)=\sum_{k\ge1}\beta_k(1-c^k)=(1-c)\sum_{k\ge1}\beta_k\bigl(1+c+\dots+c^{k-1}\bigr),
\end{equation}
which follows from $\sum_k\beta_k=1$.

\emph{Part (i).} The coefficients are nonnegative and summable, so the series converges uniformly on $[-1,1]$ and $\signalpropC$ is continuous there; on $[0,1]$ every term is nondecreasing and convex. Since $\abs{c^k}\le1$, the first expression in \eqref{eq:signalprop:telescope} is nonnegative, and $\signalpropC(c)\ge\beta_0-\sum_{k\ge1}\beta_k\ge-1$; for $c\in[0,1]$ also $\signalpropC(c)\ge\beta_0\ge0$. For the slope, the Hermite coefficients of $\bar\phi'$ are $\sqrt{k}\,a_k$ by Gaussian integration by parts (\cref{lem:toolkit:hermite}, as in the proof of \cref{lem:toolkit:price}), so Parseval's identity and $\bar\phi'(t)=\sqrt{\bar q}\,\phi'(\sqrt{\bar q}\,t)$ give
\[
  \sum_{k\ge1}k\beta_k=\frac{\sigma_w^2}{\bar q}\sum_{k\ge1}ka_k^2=\frac{\sigma_w^2}{\bar q}\E\bigl[\bar\phi'(z)^2\bigr]=\sigma_w^2\E\bigl[\phi'(\sqrt{\bar q}\,z)^2\bigr]=\chi_1 ,
\]
which is finite because $\phi'$ is polynomially bounded. For $c\in[0,1)$ the sum in \eqref{eq:signalprop:telescope} increases to $\sum_kk\beta_k$ as $c\uparrow1$, by monotone convergence, so the left derivative at $1$ is $\chi_1$. For $c\in[-1,1]$, $\abs{1+c+\dots+c^{k-1}}\le k$ bounds the sum by $\chi_1$, which proves \eqref{eq:signalprop:contraction}.

\emph{Part (ii).} Iterating \eqref{eq:signalprop:contraction} gives $1-c^\ell\le\chi_1^{\ell-1}(1-c^1)\le2\chi_1^{\ell-1}$. Since $\phi$ is not constant, $\beta_k>0$ for some $k\ge1$, so for $c\in(-1,1)$ the first expression in \eqref{eq:signalprop:telescope} is positive and $\signalpropC(c)\ge\beta_0-\sum_{k\ge1}\beta_k\abs c^k>-1$: the interval $(-1,1)$ is invariant, and $c^\ell<1$ for every $\ell$. The ratio $(1-c^{\ell+1})/(1-c^\ell)$ is the sum in \eqref{eq:signalprop:telescope} at $c=c^\ell$. As $c^\ell\to1$, each term $\beta_k(1+c^\ell+\dots+(c^\ell)^{k-1})$ tends to $k\beta_k$ and is bounded by $k\beta_k$, which is summable; by dominated convergence the ratio tends to $\chi_1$, and the logarithms of the ratios have Ces\`aro means converging to $\log\chi_1$, which is the claim.

\emph{Part (iii).} If $\phi$ is not affine, then $\beta_k>0$ for some $k\ge2$ (if $a_k=0$ for all $k\ge2$, then $\bar\phi$ is affine), so for $c\in[0,1)$ the sum in \eqref{eq:signalprop:telescope} is strictly less than $\sum_kk\beta_k=1$, and $\signalpropC(c)>c$. Since $\signalpropC$ maps $[0,1]$ into itself, $c^\ell$ increases to a fixed point in $(c^1,1]$, and $\signalpropC(c)>c$ on $[0,1)$ leaves only $1$. The ratio tends to $\chi_1=1$ by the argument of part (ii).

\emph{Part (iv).} Let $g_1(c)=\signalpropC(c)-c$, convex and continuous on $[0,1]$, with $g_1(0)=\beta_0\ge0$, $g_1(1)=0$, and left derivative $\chi_1-1>0$ at $1$; so $g_1<0$ on some interval $(1-\delta,1)$. By the intermediate value theorem $g_1$ has a zero in $[0,1-\delta]$; let $\bar c$ be the largest one. On $(\bar c,1)$, $g_1$ has no zero and takes negative values, so $g_1<0$ there. If $0\le c_0<\bar c$ were another zero, then convexity on $[c_0,1]$ with $g_1(c_0)=g_1(1)=0$ would give $g_1\le0$ on $[c_0,1]$, and the convex function $g_1$ would have the interior maximum $g_1(\bar c)=0$ on this interval, forcing $g_1\equiv0$ on $[c_0,1]$, which contradicts $g_1<0$ near $1$. So $\bar c$ is the unique fixed point in $[0,1)$, $g_1>0$ on $[0,\bar c)$, and $g_1<0$ on $(\bar c,1)$. Because $\signalpropC$ is nondecreasing on $[0,1]$, $c<\bar c$ implies $c<\signalpropC(c)\le\bar c$, and $c>\bar c$ implies $\bar c\le\signalpropC(c)<c$; so $c^\ell$ moves monotonically toward $\bar c$ and converges to a fixed point in $[0,1)$, which is $\bar c$. By convexity, $\signalpropC'(\bar c)\le\bigl(\signalpropC(c)-\signalpropC(\bar c)\bigr)/(c-\bar c)<1$ for each $c\in(\bar c,1)$, where the last inequality is $g_1(c)<0$. Finally, $\signalpropC$ is differentiable at $\bar c<1$, so for $\rho\in(\signalpropC'(\bar c),1)$ there is $\eta>0$ with $\abs{\signalpropC(c)-\bar c}\le\rho\abs{c-\bar c}$ for $\abs{c-\bar c}\le\eta$; once $c^\ell$ is within $\eta$ of $\bar c$ the distance decreases by the factor $\rho$ per layer.
\end{proof}

\Cref{thm:signalprop:phase} says that a deep random network has two fates. In the ordered phase all inputs end up with the same representation at a geometric rate, and the network computes, at depth, a function that ignores its input. In the chaotic phase two inputs, however close, end up with correlation $\bar c<1$, a correlation that the network assigns to every pair of distinct inputs: nearby points are pulled apart, and at depth the network behaves like a random function of its input. By Price's theorem (\cref{lem:toolkit:price}), $\signalpropC'(c)=\sigma_w^2\E[\phi'(\sqrt{\bar q}u)\phi'(\sqrt{\bar q}v)]$, so the slope at $\bar c$ is a correlation of derivatives, as $\chi_1$ is a mean square of derivatives.

The rate of approach to the stable fixed point defines the \term{depth scale}\index{depth scale}
\begin{equation}\label{eq:signalprop:xi}
  \xi=-\frac1{\log\signalpropC'(\bar c)},\qquad \bar c=1\ \text{in the ordered phase, so that}\ \xi=-\frac1{\log\chi_1}\ \text{there}.
\end{equation}
It is the $\xi_c$ of \cref{obs:signalprop:depth}. The \term{critical line}\index{critical line} is the curve $\chi_1=1$ in the plane $(\sigma_w^2,\sigma_b^2)$; for $\tanh$ it starts at $(1,0)$, where $\bar q=0$, and $\chi_1$ is evaluated at the fixed length $\bar q(\sigma_w^2,\sigma_b^2)$ of \cref{prop:signalprop:length}. How does $\xi$ diverge as the line is approached?

\begin{heuristic}[Critical exponent of the depth scale]
Let $\signalpropC''(1)=\sum_kk(k-1)\beta_k=\sigma_w^2\bar q\,\E[\phi''(\sqrt{\bar q}\,z)^2]$ be finite and positive, and expand near $c=1$: $\signalpropC(1-\eps)=1-\chi_1\eps+\frac12\signalpropC''(1)\eps^2+O(\eps^3)$. In the ordered phase $\xi=-1/\log\chi_1\approx1/(1-\chi_1)$ as $\chi_1\uparrow1$. In the chaotic phase, the fixed point $\bar c=1-\bar\eps$ solves $\chi_1\bar\eps\approx\bar\eps+\frac12\signalpropC''(1)\bar\eps^2$, so $\bar\eps\approx2(\chi_1-1)/\signalpropC''(1)$, and $\signalpropC'(\bar c)\approx\chi_1-\signalpropC''(1)\bar\eps\approx2-\chi_1$. Both sides give
\[
  \xi\approx\frac1{\abs{\chi_1-1}}\qquad(\chi_1\to1),
\]
a divergence with exponent $1$, the value of mean-field theory.\qedhere
\end{heuristic}

The derivation assumes that $\bar c\to1$ as $\chi_1\downarrow1$, which fails at the end point $\sigma_b=0$ of the $\tanh$ line: there $\signalpropC$ is odd, $\bar c=0$ throughout the chaotic phase, and the transition is degenerate (\cref{ex:signalprop:exponent}). The numerics of \cref{fig:signalprop:phase}(b) confirm the formula on both sides at $\sigma_b^2=0.05$: at $\chi_1=0.998$, $\xi=496$ against $1/(1-\chi_1)=497$, and at $\chi_1=1.0029$, $\xi=349$ against $345$.

\paragraph{ReLU at He initialization.} For $\phi=\relu$ and $\sigma_b=0$, the length map is $q\mapsto\frac{\sigma_w^2}2q$, so lengths neither explode nor vanish exactly when $\sigma_w^2=2$, the initialization of \textcite{HeEtAl2015}. Then $\chi_1=2\,\E[\ind\{z>0\}]=1$: He initialization is critical. The correlation map of the arc-cosine kernel of \textcite{ChoSaul2009} decides how a critical ReLU network forgets.

\begin{theorem}[ReLU at He initialization]\label{thm:signalprop:relu}\index{ReLU network!He initialization}
Let $\phi=\relu$, $\sigma_w^2=2$, and $\sigma_b=0$. Then $\signalpropV(q)=q$ for every $q\ge0$, and for every $\bar q>0$ the correlation map is
\begin{equation}\label{eq:signalprop:arccos}
  \signalpropC(c)=\frac{\sqrt{1-c^2}+(\pi-\arccos c)\,c}{\pi},\qquad c\in[-1,1],
\end{equation}
with $\signalpropC'(c)=1-\frac{\arccos c}\pi$. For $\eps\in[0,1]$,
\begin{equation}\label{eq:signalprop:relu-expansion}
  \signalpropC(1-\eps)=1-\eps+\frac{2\sqrt2}{3\pi}\eps^{3/2}+R(\eps),\qquad 0\le R(\eps)\le\frac{\pi-2}{5\sqrt2\,\pi}\eps^{5/2}.
\end{equation}
For every $c^1\in[-1,1)$, $c^\ell$ increases to $1$ from $\ell=2$ on, and
\begin{equation}\label{eq:signalprop:ninepi}
  \lim_{\ell\to\infty}\ell^2\,(1-c^\ell)=\frac{9\pi^2}2 .
\end{equation}
\end{theorem}

\begin{proof}
The plan is to compute $\signalpropC$ by integrating its derivative from $c=0$, then to expand near $c=1$, and finally to linearize the recursion in the variable $(1-c)^{-1/2}$.

\emph{Lengths.} By symmetry of $z$, $\E[\relu(\sqrt qz)^2]=\frac q2\E z^2=\frac q2$, so $\signalpropV(q)=2\cdot\frac q2=q$.

\emph{The map.} By homogeneity $\E[\relu(\sqrt{\bar q}u)\relu(\sqrt{\bar q}v)]=\bar q\,\kappa(c)$ with $\kappa(c)=\E[\relu(u)\relu(v)]$, so $\signalpropC=2\kappa$. Price's theorem (\cref{lem:toolkit:price}) applies to $\relu$, which is absolutely continuous with derivative $\ind\{t>0\}$, and gives $\kappa'(c)=\Pr[u>0,v>0]$ for $c\in(-1,1)$. Write $u=\ip{e,g}$ and $v=\ip{e',g}$ for a standard Gaussian $g\in\R^2$ and unit vectors $e,e'$ at angle $\theta=\arccos c$. The direction of $g$ is uniform on the circle, and the set of directions with both inner products positive is an arc of length $\pi-\theta$, so $\kappa'(c)=(\pi-\arccos c)/(2\pi)$. At $c=0$, independence gives $\kappa(0)=(\E\relu(u))^2=1/(2\pi)$. Integrating with the antiderivative $s\arccos s-\sqrt{1-s^2}$ of $\arccos s$,
\[
  \kappa(c)=\frac1{2\pi}+\frac1{2\pi}\int_0^c(\pi-\arccos s)\dd s=\frac{\sqrt{1-c^2}+(\pi-\arccos c)\,c}{2\pi},
\]
first for $c\in(-1,1)$ and then on $[-1,1]$ by continuity (\cref{cor:toolkit:dual}). This is \eqref{eq:signalprop:arccos}, and $\signalpropC'=2\kappa'$.

\emph{The expansion.} For $s\in[0,1]$ the series $\arcsin s=s+\sum_{n\ge1}\gamma_ns^{2n+1}$ has nonnegative coefficients, so $(\arcsin s-s)/s^3$ is nondecreasing and $0\le\arcsin s-s\le(\frac\pi2-1)s^3$. Since $\arccos(1-t)=2\arcsin\sqrt{t/2}$ for $t\in[0,2]$,
\[
  \arccos(1-t)=\sqrt{2t}+\varrho(t),\qquad 0\le\varrho(t)\le\frac{\pi-2}{2\sqrt2}\,t^{3/2}.
\]
Integrating $\signalpropC'$ from $1-\eps$ to $1$ and substituting $s=1-t$,
\[
  1-\signalpropC(1-\eps)=\int_0^\eps\Bigl(1-\frac{\arccos(1-t)}\pi\Bigr)\dd t=\eps-\frac{2\sqrt2}{3\pi}\eps^{3/2}-\frac1\pi\int_0^\eps\varrho(t)\dd t ,
\]
which is \eqref{eq:signalprop:relu-expansion} with $R(\eps)=\frac1\pi\int_0^\eps\varrho$, and the bound on $\varrho$ integrates to the bound on $R$.

\emph{Convergence.} For $c<1$, $\signalpropC'<1$ on $(c,1)$, so $1-\signalpropC(c)<1-c$, that is, $\signalpropC(c)>c$; also $\signalpropC\le1$, and $\signalpropC(c)\ge\signalpropC(-1)=0$ because $\signalpropC$ is nondecreasing. So $c^2\in[0,1)$, the sequence $(c^\ell)_{\ell\ge2}$ increases, and its limit is a fixed point of $\signalpropC$, of which $1$ is the only one.

\emph{The rate.} Let $\eps_\ell=1-c^\ell\to0$, $\eps_\ell>0$, and $u_\ell=\eps_\ell^{-1/2}$; write $a=\frac{2\sqrt2}{3\pi}$. By \eqref{eq:signalprop:relu-expansion}, $\eps_{\ell+1}=\eps_\ell(1-\delta_\ell)$ with $\delta_\ell=a\eps_\ell^{1/2}+R(\eps_\ell)/\eps_\ell$, and $0\le\delta_\ell-a\eps_\ell^{1/2}\le\eps_\ell^{3/2}$. For $\delta\in[0,\frac12]$, Taylor's theorem gives $(1-\delta)^{-1/2}=1+\frac\delta2+r(\delta)$ with $0\le r(\delta)\le2\delta^2$. Once $\delta_\ell\le\frac12$,
\[
  u_{\ell+1}-u_\ell=u_\ell\Bigl(\frac{\delta_\ell}2+r(\delta_\ell)\Bigr)=\frac a2+O\bigl(\eps_\ell^{1/2}\bigr),
\]
because $u_\ell\delta_\ell=a+O(\eps_\ell)$ and $u_\ell\delta_\ell^2=O(\eps_\ell^{1/2})$. Hence $u_{\ell+1}-u_\ell\to a/2$, and by the Stolz--Ces\`aro theorem $u_\ell/\ell\to a/2$. So $\ell^2\eps_\ell=(\ell/u_\ell)^2\to4/a^2=\frac{4\cdot9\pi^2}{8}=\frac{9\pi^2}2$.
\end{proof}

ReLU networks at He initialization decorrelate polynomially: no depth scale, but a power law. The critical $\tanh$ network also decays as a power, but a different one: with $\signalpropC''(1)$ finite, $\eps_{\ell+1}=\eps_\ell-\frac12\signalpropC''(1)\eps_\ell^2+O(\eps_\ell^3)$ gives $1-c^\ell\sim2/(\signalpropC''(1)\ell)$ (\cref{ex:signalprop:smooth}). The exponent at the critical point depends on the smoothness of $\phi$ at the origin, through the first non-analytic term of $\signalpropC$ at $1$: $\eps^2$ for $\tanh$, $\eps^{3/2}$ for ReLU.

\paragraph{Gradients.} Forward correlations are half of the story; the other half is the size of the gradient. Let $\Loss$ be a loss that depends on the network through $h^L$, and let $\delta^\ell=\partial\Loss/\partial h^\ell$. The chain rule gives $\delta^\ell=D^\ell(W^{\ell+1})^\top\delta^{\ell+1}$ with $D^\ell=\diag(\phi'(h^\ell))$.

\begin{heuristic}[Gradient norms]
Assume the \term{gradient independence assumption}\index{gradient independence assumption}: in second moments, the matrix $W^{\ell+1}$ of the backward pass may be replaced by an independent copy, and $\delta^{\ell+1}$ is independent of $h^\ell$. Then, conditionally on $\delta^{\ell+1}$ and $h^\ell$, each coordinate $\sum_iW^{\ell+1}_{ij}\delta^{\ell+1}_i$ has variance $\frac{\sigma_w^2}{m_\ell}\norm{\delta^{\ell+1}}^2$, and summing over $j$,
\[
  \E\norm{\delta^\ell}^2=\sigma_w^2\,\E\Bigl[\frac1{m_\ell}\sum_{j}\phi'(h^\ell_j)^2\Bigr]\,\E\norm{\delta^{\ell+1}}^2\ \longrightarrow\ \chi_1\,\E\norm{\delta^{\ell+1}}^2
\]
at infinite width and length $\bar q$, by the law of large numbers of \cref{thm:signalprop:gp}. Iterating, $\E\norm{\delta^\ell}^2\approx\chi_1^{L-\ell}\,\E\norm{\delta^L}^2$: gradients vanish geometrically in the ordered phase, explode in the chaotic phase, and keep their mean squared size on the critical line, with depth scale $\xi_\nabla=-1/\log\chi_1$.\qedhere
\end{heuristic}

The assumption is false at finite width, where the forward and backward passes share the same matrices, but it gives the correct answer in the infinite-width limit for this network and many others: \textcite{Yang2020} gave a simple checkable condition under which computations with it are exact in the limit, and showed that they can be wrong when the condition fails.\sidenote{For $\tanh$ networks of depth $30$ at $\sigma_b^2=0.05$ we measured the ratio $\norm{\delta^0}^2/\norm{\delta^{30}}^2$. At criticality, averaged over $40$ networks, it was $0.89$ at width $800$ and $1.00$ at width $1600$, against $\chi_1^{30}=1$, with a spread across networks of about $25\%$. Over $10$ networks of width $800$ it was $0.0071$ against $0.0064$ at $\sigma_w^2=1.2$, and $265$ against $297$ at $\sigma_w^2=3$.} Two cases can be done exactly at every width, and they show what the mean hides.

\begin{theorem}[Jacobian norms at finite width]\label{thm:signalprop:finite}\index{finite-width fluctuations}
Let all widths equal $m$, $\sigma_b=0$, and fix an input $x\ne0$ and a direction $\dot x\in\R^d$. Let $u^\ell=\partial_th^\ell(x+t\dot x)|_{t=0}$, so that $u^1=W^1\dot x$ and $u^{\ell+1}=W^{\ell+1}D^\ell u^\ell$, and let $Y_\ell=\norm{D^\ell u^\ell}^2$, the squared norm of the derivative of $a^\ell=\phi(h^\ell)$ along $\dot x$.
\begin{enumerate}
\item If $\phi(t)=t$ (a deep linear network), then $Y_L=Y_1\prod_{\ell=1}^{L-1}F_\ell$, where $F_1,\dots,F_{L-1}$ are i.i.d., independent of $Y_1$, with law $\frac{\sigma_w^2}m\chi^2_m$.
\item If $\phi=\relu$ with $D^\ell=\diag(\ind\{h^\ell>0\})$ and $\sigma_w^2=2$, then the same holds with $F_\ell$ of law $\frac2m\chi^2_{B}$, where $B\sim\mathrm{Bin}(m,\frac12)$ and $\chi^2_0=0$.
\end{enumerate}
In both cases $\E Y_L=\chi_1^{L-1}\E Y_1$, and
\[
  \frac{\E Y_L^2}{(\E Y_L)^2}=\frac{\E Y_1^2}{(\E Y_1)^2}\Bigl(1+\frac\kappa m\Bigr)^{L-1},\qquad \kappa=2\ \text{(linear)},\quad \kappa=5\ \text{(ReLU)}.
\]
\end{theorem}

\begin{proof}
The plan is to show that, conditionally on $\cF_\ell$, the ratio $Y_{\ell+1}/Y_\ell$ has a fixed law that does not depend on $\cF_\ell$; independence of the factors and the moment formulas follow.

\emph{Linear case.} Here $D^\ell=I_m$ and $u^{\ell+1}=W^{\ell+1}u^\ell$. Conditionally on $\cF_\ell$, the coordinates of $W^{\ell+1}u^\ell$ are independent $\cN(0,\frac{\sigma_w^2}m\norm{u^\ell}^2)$, so $Y_{\ell+1}=Y_\ell\,F_\ell$ with $F_\ell\sim\frac{\sigma_w^2}m\chi^2_m$ independent of $\cF_\ell$ (if $Y_\ell=0$, then $Y_{\ell+1}=0$ and $F_\ell$ is drawn independently). Since $Y_1,F_1,\dots,F_{\ell-1}$ are $\cF_\ell$-measurable, the factors are independent of each other and of $Y_1$.

\emph{ReLU case.} Conditionally on $\cF_\ell$, \cref{lem:signalprop:conditional} (with $\cX=\{x\}$, extended to the pair $(h,u)$ by the same argument) shows that the pairs $(h^{\ell+1}_i,u^{\ell+1}_i)$, $i\le m$, are i.i.d.\ centered Gaussian, with $\Var(u^{\ell+1}_i)=\frac2m Y_\ell$ and $\Var(h^{\ell+1}_i)=\frac2m\norm{a^\ell}^2$. If $a^\ell=0$, then $D^\ell=0$, so $Y_\ell=0$ and $Y_{\ell+1}=0$. Otherwise $\Var(h^{\ell+1}_i)>0$, and the law of a centered Gaussian pair is invariant under $(h,u)\mapsto(-h,-u)$, so for every Borel set $A$, $\Pr[h>0,u^2\in A]=\Pr[h<0,u^2\in A]=\frac12\Pr[u^2\in A]$: the sign of $h^{\ell+1}_i$ is a fair coin independent of $(u^{\ell+1}_i)^2$. Hence $Y_{\ell+1}=\sum_i\ind\{h^{\ell+1}_i>0\}(u^{\ell+1}_i)^2$ is, conditionally on $\cF_\ell$, a sum of $B\sim\mathrm{Bin}(m,\frac12)$ independent squares of $\cN(0,\frac2mY_\ell)$ variables, that is, $Y_\ell$ times a variable of law $\frac2m\chi^2_B$. The case $\ell=1$ uses $x\ne0$, and independence follows as in the linear case.

\emph{Moments.} For the linear case $\E F=\sigma_w^2$ and $\E F^2=\frac{\sigma_w^4}{m^2}(m^2+2m)=\sigma_w^4(1+\frac2m)$, and $\chi_1=\sigma_w^2$. For ReLU, $\E[\chi^2_B\mid B]=B$ and $\E[(\chi^2_B)^2\mid B]=B^2+2B$, with $\E B=\frac m2$ and $\E B^2=\frac m4+\frac{m^2}4$; so $\E F=1=\chi_1$ and $\E F^2=\frac4{m^2}\bigl(\frac{m^2}4+\frac m4+m\bigr)=1+\frac5m$. Independence gives $\E Y_L=\E Y_1(\E F)^{L-1}$ and $\E Y_L^2=\E Y_1^2(\E F^2)^{L-1}$.
\end{proof}

At criticality the mean of $Y_L/Y_1$ is exactly $1$ at every width, as the heuristic predicts, but $\log(Y_L/Y_1)$ is a sum of $L-1$ independent terms, each with variance close to $\kappa/m$: the typical network has $\log(Y_L/Y_1)\approx-\frac{\kappa}{2m}L\pm\sqrt{\kappa L/m}$. The mean is carried by rare networks with large gradients, and the parameter that controls the typical behavior is $L/m$, which the infinite-width limit sets to zero. For ReLU this is the result of \textcite{Hanin2018} and \textcite{HaninNica2020}, who prove the analogous statements for the input--output Jacobian with variable widths, where $\sum_\ell1/m_\ell$ replaces $L/m$.

\paragraph{Dynamical isometry.} Criticality controls the mean squared singular value of the input--output Jacobian; trainability also depends on its spread. A network has \term{dynamical isometry}\index{dynamical isometry} if all singular values of its Jacobian concentrate near $1$. \Textcite{SaxeEtAl2014} observed that deep linear networks with orthogonal weight matrices have this property exactly, while Gaussian weights give a spread that grows with depth. \Textcite{PenningtonEtAl2017} computed the limiting singular value distribution with free probability: ReLU networks cannot achieve dynamical isometry with Gaussian or orthogonal weights, whereas $\tanh$ networks can, with orthogonal weights on the critical line near $\bar q=0$, and then learn faster.

\section{Reading the solution}\label{sec:signalprop:reading}

\paragraph{The phase diagram.} \Cref{fig:signalprop:phase}(a) shows the critical line of $\tanh$ networks, computed by solving $\chi_1=1$ with $\bar q$ from \cref{prop:signalprop:length}. It starts at $(\sigma_w^2,\sigma_b^2)=(1,0)$ and bends to larger weight variances as the bias variance grows: a bias adds a common component to all pre-activations, which raises $\bar q$, pushes the neurons into the saturated part of $\tanh$, and lowers $\E[\phi'^2]$, so more weight variance is needed to keep $\chi_1=1$. Left of the line is order and right of it is chaos, and the contours of $\xi$ crowd toward the line from both sides. Panel (b) shows $\xi$ along the cut $\sigma_b^2=0.05$, where the critical point is $\sigma_w^2=1.761$: the divergence is the critical slowing down of a second-order transition, and the curve $1/\abs{\chi_1-1}$ of the heuristic derivation is indistinguishable from it near the line. Combined with \cref{obs:signalprop:depth}, the plot is a map of trainable depth: at $\sigma_w^2=1.2$, where $\xi=5.9$, networks much deeper than $6\xi\approx36$ layers should fail; at $\sigma_w^2=1.7$, where $\xi=75$, several hundred layers should train.

\begin{figure}[t]
\begin{fullwidth}
\centering
\begin{tikzpicture}
\begin{groupplot}[group style={group size=2 by 1,horizontal sep=17mm},hplot,
  width=0.5\linewidth,height=0.37\linewidth,unbounded coords=jump]
\nextgroupplot[xlabel={$\sigma_w^2$},ylabel={$\sigma_b^2$},xmin=0,xmax=8,ymin=0,ymax=2,
  title={(a) phase diagram of $\tanh$},legend pos=north west,legend style={font=\scriptsize}]
\addplot[theory] table[x=crit,y=sb2] {figures/data/05-signalprop-phase.dat};
\addlegendentry{critical line $\chi_1=1$}
\addplot[theory2] table[x=ord20,y=sb2] {figures/data/05-signalprop-phase.dat};
\addlegendentry{$\xi=20$}
\addplot[theory3] table[x=ord5,y=sb2] {figures/data/05-signalprop-phase.dat};
\addlegendentry{$\xi=5$}
\addplot[theory2,forget plot] table[x=cha20,y=sb2] {figures/data/05-signalprop-phase.dat};
\addplot[theory3,forget plot] table[x=cha5,y=sb2] {figures/data/05-signalprop-phase.dat};
\node[font=\small,text=muted] at (axis cs:1.0,1.2) {ordered};
\node[font=\small,text=muted] at (axis cs:6.2,0.5) {chaotic};
\nextgroupplot[xlabel={$\sigma_w^2$},ylabel={depth scale $\xi$},xmin=0.5,xmax=8,ymode=log,ymin=1,ymax=2000,
  title={(b) $\xi$ along $\sigma_b^2=0.05$},legend pos=north east,legend style={font=\scriptsize},
  log ticks with fixed point]
\addplot[guide,forget plot] coordinates {(1.761,1) (1.761,2000)};
\addplot[theory] table[x=sw2,y=xi] {figures/data/05-signalprop-xi.dat};
\addlegendentry{$-1/\log\signalpropC'(\bar c)$}
\addplot[theory3,densely dashed] table[x=sw2,y expr={1/abs(\thisrow{chi1}-1)}] {figures/data/05-signalprop-xi.dat};
\addlegendentry{$1/\abs{\chi_1-1}$}
\end{groupplot}
\end{tikzpicture}
\caption{Infinitely wide $\tanh$ networks, computed by quadrature (theory only). (a)~The critical line $\chi_1=1$ in the plane of weight and bias variances, and the contours $\xi=5$ and $\xi=20$ of the depth scale \eqref{eq:signalprop:xi} on both sides of it. (b)~The depth scale along $\sigma_b^2=0.05$ (critical point $\sigma_w^2=1.761$, dotted), with the approximation $1/\abs{\chi_1-1}$ of the heuristic derivation.}
\label{fig:signalprop:phase}
\end{fullwidth}
\end{figure}

\paragraph{Correlations, layer by layer.} \Cref{fig:signalprop:corr} follows two inputs with initial correlation $0.6$ through $\tanh$ networks of width $1000$ in the three phases, against the iterated correlation map. In the ordered phase ($\sigma_w^2=1.2$, $\chi_1=0.845$) the correlation approaches $1$ geometrically, the gap $1-c^\ell$ shrinking by about $\chi_1$ per layer; at criticality it approaches $1$ like $6.2/\ell$, the power law of \cref{ex:signalprop:smooth}; in the chaotic phase ($\sigma_w^2=4$) it falls to $\bar c=0.166$. Every average lies within three standard errors of its curve. The spread across networks is the finite-width correction: it is small in the ordered phase, where the map contracts strongly, and largest in the chaotic phase, where the representation is a random function of the input.

\begin{figure}[t]
\centering
\begin{tikzpicture}
\begin{axis}[hplot,height=0.5\linewidth,xlabel={layer $\ell$},ylabel={correlation $c^\ell$},xmin=0,xmax=51,ymin=0,ymax=1.02,
  legend pos=south east,legend columns=1,legend style={font=\scriptsize}]
\addplot[guide,forget plot,domain=0:51,samples=2] {0.165578};
\addplot[theory] table[x=layer,y=ordth] {figures/data/05-signalprop-corr.dat};
\addlegendentry{ordered, $\sigma_w^2=1.2$}
\addplot[theory2] table[x=layer,y=critth] {figures/data/05-signalprop-corr.dat};
\addlegendentry{critical, $\sigma_w^2=1.761$}
\addplot[theory3] table[x=layer,y=chath] {figures/data/05-signalprop-corr.dat};
\addlegendentry{chaotic, $\sigma_w^2=4$}
\addplot[simulation,forget plot,each nth point=2] table[x=layer,y=ordmc] {figures/data/05-signalprop-corr.dat};
\addplot[simulation2,forget plot,each nth point=2] table[x=layer,y=critmc] {figures/data/05-signalprop-corr.dat};
\addplot[simulation,mark=triangle*,mark size=2pt,forget plot,each nth point=2] table[x=layer,y=chamc] {figures/data/05-signalprop-corr.dat};
\end{axis}
\end{tikzpicture}
\caption{Correlation of the pre-activations of two inputs in $\tanh$ networks with $\sigma_b^2=0.05$, starting from $c^1=0.6$ at the fixed length $\bar q$. Curves: iterates of the correlation map \eqref{eq:signalprop:cmap}. Marks: networks of width $m=1000$, averages over $50$ networks of the empirical cosine $\ip{h^\ell(x),h^\ell(x')}/(\norm{h^\ell(x)}\norm{h^\ell(x')})$, every second layer (seed $20261017$); the standard deviation across networks is about $0.03$ at criticality and $0.07$ in the chaotic phase. The dotted line is the chaotic fixed point $\bar c=0.166$.}
\label{fig:signalprop:corr}
\end{figure}

\paragraph{The ReLU constant.} \Cref{fig:signalprop:relu}(a) shows $\ell^2(1-c^\ell)$ for two orthogonal inputs of a ReLU network at He initialization. The limit $9\pi^2/2=44.41$ is approached slowly: the proof gives $u_{\ell+1}-u_\ell=\frac a2+\frac{3a^2}{8u_\ell}+O(u_\ell^{-2})$, hence $u_\ell=\frac a2\ell+\frac{3a}4\log\ell+O(1)$ and a relative correction of about $3\log\ell/\ell$ to $\ell^2(1-c^\ell)$; at $\ell=1000$ the exact iterate gives $43.2$, at $\ell=10^4$ it gives $44.26$. Simulations at width $1000$ agree with the iterated map up to depth $181$, within two standard errors.

\paragraph{Depth over width.} \Cref{fig:signalprop:relu}(b) tests \cref{thm:signalprop:finite} at width $m=50$ up to depth $151$: the variance of $\log(Y_L/Y_1)$ grows linearly in $L/m$, with slope close to $5$ for ReLU and $2$ for linear networks. At $L=m$ the typical gradient differs from its mean by a factor of order $e^{\pm\sqrt5}$, ten either way, although every network is critical on average. The same parameter controls the forward correlations: with the exact two-input chain at widths $25$ to $200$, the mean of $1-c^L$ at $L=m$ still matches $\signalpropC$, but its median is only about a quarter of the mean, at every width (\cref{ex:signalprop:fluct}).

\begin{figure}[t]
\begin{fullwidth}
\centering
\begin{tikzpicture}
\begin{groupplot}[group style={group size=2 by 1,horizontal sep=17mm},hplot,
  width=0.5\linewidth,height=0.37\linewidth,unbounded coords=jump]
\nextgroupplot[xlabel={layer $\ell$},ylabel={$\ell^2(1-c^\ell)$},xmode=log,xmin=1.5,xmax=1100,ymin=0,ymax=50,
  title={(a) ReLU, He initialization},legend pos=south east,legend style={font=\scriptsize},log ticks with fixed point]
\addplot[guide,domain=1.5:1100,samples=2] {44.4132};
\addlegendentry{$9\pi^2/2$}
\addplot[theory] table[x=layer,y=theory] {figures/data/05-signalprop-relu.dat};
\addlegendentry{iterated map}
\addplot[simulation,error bars/.cd,y dir=both,y explicit] table[x=layer,y=mc,y error=mcse] {figures/data/05-signalprop-relu.dat};
\addlegendentry{width $1000$}
\nextgroupplot[xlabel={depth over width $L/m$},ylabel={$\Var\log(Y_L/Y_1)$},xmin=0,xmax=3.1,ymin=0,ymax=17,
  title={(b) fluctuations at width $m=50$},legend pos=north west,legend style={font=\scriptsize}]
\addplot[theory] table[x=Lm,y=relutheory] {figures/data/05-signalprop-fluct.dat};
\addlegendentry{ReLU, theorem}
\addplot[theory2] table[x=Lm,y=lintheory] {figures/data/05-signalprop-fluct.dat};
\addlegendentry{linear, theorem}
\addplot[simulation] table[x=Lm,y=relumc] {figures/data/05-signalprop-fluct.dat};
\addlegendentry{ReLU, simulation}
\addplot[simulation2] table[x=Lm,y=linmc] {figures/data/05-signalprop-fluct.dat};
\addlegendentry{linear, simulation}
\end{groupplot}
\end{tikzpicture}
\caption{Critical networks at large and at finite width. (a)~ReLU at $\sigma_w^2=2$, $\sigma_b=0$, two orthogonal inputs: the iterated arc-cosine map \eqref{eq:signalprop:arccos} and its limit \eqref{eq:signalprop:ninepi}; points are averages over $400$ networks of width $1000$, with standard errors (seed $20261017$). (b)~Variance of the log squared Jacobian norm $\log(Y_L/Y_1)$ of \cref{thm:signalprop:finite}. Lines: $(L-1)\Var\log F$, computed exactly from the laws of the factors ($\Var\log F=0.108$ for ReLU and $0.041$ for linear networks at $m=50$). Points: $2000$ networks of width $50$ with full Gaussian weight matrices.}
\label{fig:signalprop:relu}
\end{fullwidth}
\end{figure}

\paragraph{The physics.} The order parameters $q$ and $c$ play the role of the magnetization of a mean-field magnet: at infinite width each neuron couples to infinitely many others, the coupling is through an average, and the average obeys a closed deterministic equation. Depth plays the role of time, or of the distance from a boundary at which a perturbation is applied, and $\xi$ is a correlation length. The ordered phase is the high-temperature phase, in which perturbations of the input die out, and the chaotic phase is the analogue of a positive Lyapunov exponent: infinitesimally close inputs separate by a factor $\sqrt{\chi_1}$ per layer, because $1-c^{\ell+1}\approx\chi_1(1-c^\ell)$ is the squared distance of the two representations in units of $\bar q$. At the critical line the correlation length diverges with the mean-field exponent $1$, and at the critical point the approach to equilibrium is a power law, as in critical slowing down. Mean-field theory is exact here for the same reason as in the Curie--Weiss model: the interaction range, the width, is infinite.

\section{What the model gets wrong}\label{sec:signalprop:wrong}

\paragraph{What was thrown away.} \emph{Finite width.} The limit sets $L/m=0$, and every fluctuation of order $\sqrt{L/m}$ disappears with it (\cref{thm:signalprop:finite}). \emph{Training.} The theory describes the network at time zero; it predicts where gradient descent can start, not what it finds, how fast, or whether the weights move far from their initial values (\cref{ch:lazy}). \emph{Architecture.} No residual connections, no normalization, no convolutions, no attention; each changes the recursion, and some remove the phase transition altogether. \emph{Data.} Only pairwise inner products of inputs enter, so the theory cannot distinguish a dataset from a set of points with the same Gram matrix, and it says nothing about which functions of the data are easy to learn. \emph{Weights.} Gaussian weights are not essential for the forward recursion \parencite{Hanin2023}, but orthogonal weights change the Jacobian spectrum and with it the trainability \parencite{PenningtonEtAl2017}.

\paragraph{Where it breaks.} \emph{Depth comparable to width.} Gradient statistics at finite width are controlled by $\sum_\ell1/m_\ell$, which is $L/m$ for constant widths: for ReLU networks the variance of the squared entries of the input--output Jacobian grows exponentially in this sum \parencite{Hanin2018}, and the logarithm of the squared norm of a Jacobian--vector product is approximately Gaussian, with mean and variance of order $\sum_\ell1/m_\ell$, when depth and width grow together \parencite{HaninNica2020}. \Cref{thm:signalprop:finite} is the exactly solvable shadow of this result. When $L\gtrsim m$ the infinite-width kernel is the mean of a broad distribution and not the typical value: the network that is actually initialized is not the Gaussian process. \Textcite{RobertsEtAl2022} organize the finite-width corrections as an expansion in $L/m$, and find that criticality is necessary to keep the corrections from growing exponentially, but that even at criticality the expansion parameter is $L/m$.

\emph{Criticality is necessary, not sufficient.} Critical ReLU networks train more slowly than critical sigmoidal networks with orthogonal weights, because the spread of the Jacobian singular values, which the model does not compute, grows with depth for ReLU \parencite{PenningtonEtAl2017}. Inside the trainable region, which initialization trains fastest depends on the activation beyond $\chi_1$ \parencite{HayouEtAl2019}. And training can leave the region in which the initialization analysis applies.

\emph{Kernels are not networks.} The kernel $K^L$ is the covariance of the infinitely wide Bayesian network \parencite{LeeEtAl2018}. In the large study of \textcite{LeeEtAl2020}, these kernels (and the tangent kernels of \cref{ch:lazy}) outperformed finite fully connected networks trained in the standard way, but underperformed finite convolutional networks; large learning rates and weight decay, which help finite networks, break the correspondence. The infinite-width kernel misses what trained convolutional networks gain by learning features.

\emph{Normalization changes the answer.} The model suggests that a well-chosen activation and variance cure exploding gradients. \Textcite{YangEtAl2019} showed that in fully connected networks with batch normalization and no skip connections, gradients explode exponentially with depth at initialization whatever the activation and the weight variance: the normalization itself causes the explosion, and such networks are untrainable at large depth.

\emph{Attention collapses.} A stack of self-attention layers without skip connections and feedforward blocks converges to a rank-one output doubly exponentially in depth \parencite{DongEtAl2021}, an ordered phase that variance tuning cannot leave. In transformers at initialization, this rank collapse makes the gradients of queries and keys vanish, and \textcite{NociEtAl2022} prevent it by scaling the residual branches with depth.

\paragraph{Fine structure.} The $1/m$ corrections are computed systematically by the effective theory of \textcite{RobertsEtAl2022}, in which the four-point vertex of the pre-activations grows linearly in $L/m$ at criticality. When depth and width go to infinity together with $L/m$ fixed, the covariance of a network whose activation is shaped toward the identity at a rate tied to the width converges to the solution of a stochastic differential equation, the \emph{neural covariance SDE} of \textcite{LiEtAl2022}: the deterministic recursion of this chapter becomes a diffusion, and $L/m$ is its time. Residual networks do not have the exponential depth scales of plain networks, and their signal propagation is subexponential or polynomial \parencite{YangSchoenholz2017}; scaling the residual branches by $1/\sqrt L$ gives well-defined infinite-depth limits in which hyperparameters transfer across depth \parencite{YangEtAl2024,BordelonEtAl2024}. Deep kernel shaping \parencite{MartensEtAl2021} controls the whole correlation map, not only its slope at $1$, by transforming the activation, and trains deep networks without skip connections or normalization. Open: a theory of trainability, rather than of signal propagation, that predicts the constant in $L\lesssim6\xi$, and a description of the regime $L\sim m$ during training.

\section{Model card}\label{sec:signalprop:card}

\modelcard{Deep networks train only for well-scaled initializations; the maximal trainable depth is a multiple of a depth scale that diverges on an order--chaos boundary; He initialization for ReLU.}
{Fully connected network with Gaussian weights $\cN(0,\sigma_w^2/m)$ and biases $\cN(0,\sigma_b^2)$, widths to infinity at fixed depth; order parameters $q^\ell$ and $c^\ell$.}
{Gaussian-process limit (\cref{thm:signalprop:gp}); length map (\cref{prop:signalprop:length}); order and chaos from the dual power series (\cref{thm:signalprop:phase}); ReLU at He initialization (\cref{thm:signalprop:relu}); exact Jacobian fluctuations at finite width (\cref{thm:signalprop:finite}).}
{The critical line $\chi_1=1$; depth scales $\xi=-1/\log\signalpropC'(\bar c)\approx1/\abs{\chi_1-1}$; gradient growth $\chi_1^{L}$; $1-c^\ell\sim9\pi^2/(2\ell^2)$ for ReLU and $\sim2/(\signalpropC''(1)\ell)$ for smooth critical activations.}
{Finite width (fluctuations of order $\sqrt{L/m}$); training; residual connections, normalization, and attention; the spread of the Jacobian spectrum; the constant in $L\lesssim6\xi$.}
{$1/m$ effective field theory; neural covariance SDE at fixed $L/m$; dynamical isometry; residual and depth-$\mu$P scalings; kernel shaping.}

The kernel $K^L$ of this chapter is the starting point of \cref{ch:lazy}: the neural tangent kernel is built from the same recursion together with the derivative map $c\mapsto\signalpropC'(c)$, and the order--chaos analysis decides whether that kernel is degenerate at depth. The Hermite expansion of the activation, which turns a layer into a power series in the correlation, returns in \cref{ch:singleindex}, where the lowest nonzero Hermite coefficient of a link function sets the information exponent; in both chapters Gaussian inputs make a nonlinearity act diagonally on Hermite polynomials, as characters diagonalize convolution in \cref{ch:grokking}. The spectrum of a product of random matrices in \cref{thm:signalprop:finite} is a relative of the Marchenko--Pastur law of \cref{ch:doubledescent}.

\begin{exercises}
\item \consolidation (The empirical kernel concentrates.) In the setting of \cref{thm:signalprop:gp}, show that in the iterated limit $m_1\to\infty,\dots,m_\ell\to\infty$ the empirical kernel $\signalpropG^\ell(x,x')$ converges in probability to $K^{\ell+1}(x,x')$, and deduce that $\frac1{m_\ell}\norm{h^\ell(x)}^2\to q^\ell(x)$. \hint{apply the theorem to bounded continuous functions of $\frac1J\sum_{j\le J}\phi(h^\ell_j(x))\phi(h^\ell_j(x'))$, truncated, and use the exchangeability of the neurons.}\label{ex:signalprop:empirical}

\item \consolidation (The Fourier transform of $\operatorname{sech}^2$.) (a) Integrate $e^{i\omega t}\operatorname{sech}^2t$, for $\omega>0$, over the boundary of the rectangle with corners $\pm R$ and $\pm R+i\pi$, and show that $\int_\R\operatorname{sech}^2(t)e^{i\omega t}\dd t=\pi\omega/\sinh(\pi\omega/2)$. (b) Deduce that $\E[\tanh^2(\sqrt qz)]=q-2q^2+O(q^3)$, and that for $\sigma_b=0$ and $\sigma_w^2=1$ the length of a $\tanh$ network satisfies $\ell q^\ell\to\frac12$. \hint{in (a), the only pole inside is the double pole at $i\pi/2$, where $\operatorname{sech}^2t\approx-(t-i\pi/2)^{-2}$, and the top side of the rectangle is $-e^{-\pi\omega}$ times the bottom one.}\label{ex:signalprop:fourier}

\item \consolidation\hard (The critical exponent, rigorously.) Assume $\signalpropC''(1)<\infty$. (a) Show that $\xi(1-\chi_1)\to1$ as $\chi_1\uparrow1$. (b) For $\tanh$ with a fixed $\sigma_b^2>0$, let $\sigma_w^2$ decrease to its critical value, and show that $1-\bar c\sim2(\chi_1-1)/\signalpropC''(1)$ and $\xi(\chi_1-1)\to1$ as $\chi_1\downarrow1$. (c) For $\tanh$ with $\sigma_b=0$, show that $\bar c=0$ for every $\sigma_w^2>1$, and that $\chi_1-1$ and $1-\signalpropC'(0)$ are both of order $(\sigma_w^2-1)^2$. Why does (b) not apply?\label{ex:signalprop:exponent}

\item \consolidation (Smooth activations at criticality.) Let $\chi_1=1$ and $0<\signalpropC''(1)<\infty$, and let $\signalpropC(1-\eps)=1-\eps+\frac12\signalpropC''(1)\eps^2+O(\eps^3)$. Show that $\ell(1-c^\ell)\to2/\signalpropC''(1)$, and that $\signalpropC''(1)=\sigma_w^2\bar q\,\E[\phi''(\sqrt{\bar q}z)^2]$. For $\tanh$ at $\sigma_b^2=0.05$, check numerically that $\signalpropC''(1)=0.324$, so that $1-c^\ell\approx6.17/\ell$, as in \cref{fig:signalprop:corr}. \hint{use $1/\eps_{\ell+1}-1/\eps_\ell$.}\label{ex:signalprop:smooth}

\item \consolidation\hard (The logarithmic correction.) In the proof of \cref{thm:signalprop:relu}, include the next term of the expansion \eqref{eq:signalprop:relu-expansion} and show that $u_\ell=\frac a2\ell+\frac{3a}4\log\ell+O(1)$. Deduce that $\ell^2(1-c^\ell)=\frac{9\pi^2}2\bigl(1-\frac{3\log\ell}\ell+O(\frac1\ell)\bigr)$, and compare with the values $43.2$ at $\ell=1000$ and $44.26$ at $\ell=10^4$.\label{ex:signalprop:log}

\item \variation (The error function.) Let $\phi=\erf$. (a) Show that for a centered Gaussian pair with variances $q_1,q_2$ and covariance $q_{12}$, $\E[\erf(z)\erf(z')]=\frac2\pi\arcsin\frac{2q_{12}}{\sqrt{(1+2q_1)(1+2q_2)}}$, the kernel of \textcite{Williams1997}. (b) Show that $\chi_1=\frac{4\sigma_w^2}{\pi\sqrt{1+4\bar q}}$, that the critical point at $\sigma_b=0$ is $\sigma_w^2=\pi/4$, and that the critical line is the curve $\sigma_w^2=\frac\pi4\sqrt{1+4\bar q}$, $\sigma_b^2=\bar q-\frac12\sqrt{1+4\bar q}\arcsin\frac{2\bar q}{1+2\bar q}$, $\bar q>0$. \hint{in (a), use Price's theorem and $\E[\phi'(z)\phi'(z')]$ for Gaussian $\phi'$.}\label{ex:signalprop:erf}

\item \variation (Leaky ReLU.) Let $\phi(t)=\max(t,\alpha t)$ with $\alpha\in[0,1)$ and $\sigma_b=0$. Show that lengths are preserved exactly when $\sigma_w^2=2/(1+\alpha^2)$, that this initialization is critical, and that
\[
  \signalpropC(c)=\frac{(1+\alpha)^2c+(1-\alpha)^2\frac2\pi\bigl(\sqrt{1-c^2}+c\arcsin c\bigr)}{2(1+\alpha^2)} .
\]
Deduce that $\ell^2(1-c^\ell)\to\frac{9\pi^2}2\bigl(\frac{1+\alpha^2}{(1-\alpha)^2}\bigr)^2$. Why does the constant blow up as $\alpha\to1$?\label{ex:signalprop:leaky}

\item \variation (ReLU off the He point.) Let $\phi=\relu$. (a) For $\sigma_w^2<2$ and $\sigma_b^2>0$, show that $\bar q=\sigma_b^2/(1-\sigma_w^2/2)$ and $\chi_1=\sigma_w^2/2<1$: ReLU networks with a fixed length are never chaotic. (b) For $\sigma_w^2>2$, show that lengths grow geometrically, so that \cref{thm:signalprop:phase} does not apply; compute the correlation map at $\sigma_b=0$ and show that it does not depend on $\sigma_w^2$.\label{ex:signalprop:relubias}

\item \variation (A residual network.) Let $h^{\ell+1}=h^\ell+\frac{\gamma}{\sqrt L}W^{\ell+1}\relu(h^\ell)$ with $W^{\ell+1}_{ij}\sim\cN(0,2/m)$ and no biases. (a) Derive the length recursion $q^{\ell+1}=(1+\gamma^2/L)q^\ell$ at infinite width, and show that $q^L/q^1\to e^{\gamma^2}$ as $L\to\infty$, while without the factor $1/\sqrt L$ the length grows like $(1+\gamma^2)^L$. (b) Derive the correlation recursion and show that, with the factor $1/\sqrt L$, the correlation of two inputs converges as $L\to\infty$ to the solution at time $1$ of an ordinary differential equation in $c$; write it down.\label{ex:signalprop:resnet}

\item \simulation (Gradients in three phases.) Build $\tanh$ networks of width $800$ and depth $30$ with $\sigma_b^2=0.05$ and $\sigma_w^2\in\{1.2,1.761,3\}$, start from a random input of length $\bar q$, and backpropagate $\Loss=\ip{v,h^{30}}$ with $v\sim\cN(0,I_m)$. Average $\norm{\delta^0}^2/\norm{\delta^{30}}^2$ over $10$ networks and compare with $\chi_1^{30}$ ($0.0064$, $1$, and $297$). Observe agreement within the spread across networks, which is $30$--$60\%$ of the mean. Repeat at width $200$ and observe the spread grow, as \cref{thm:signalprop:finite} suggests.\label{ex:signalprop:grad}

\item \simulation (Typical versus mean.) Use the two-input Markov chain of \cref{lem:signalprop:conditional} for ReLU at He initialization, with orthogonal inputs, widths $m\in\{25,50,100,200\}$, depths $L\in\{m/4,m/2,m\}$, and $1000$ networks each. (a) Observe that the mean of $1-\hat c^L$ agrees with the iterated map within a few percent in every case. (b) Observe that the ratio of the median of $1-\hat c^L$ to its mean is about $0.75$, $0.5$, and $0.23$ for $L/m=\frac14,\frac12,1$, independently of $m$, and that the relative standard deviation grows from about $0.8$ to about $2.2$. Explain why the controlling parameter is $L/m$.\label{ex:signalprop:fluct}

\item \research (Trainability beyond signal propagation.) \Cref{obs:signalprop:depth} relates the trainable depth to $6\xi$, with a constant fitted to experiments. Open: a solvable model of the first steps of training of a deep network near the critical line that predicts the trainable depth, including the constant and its dependence on the learning rate and the training budget.\label{ex:signalprop:researchtrain}

\item \research (The regime $L\sim m$.) The neural covariance SDE of \textcite{LiEtAl2022} describes shaped networks at initialization when $L/m$ is fixed. Open: an analogous limit for unshaped critical networks (ReLU at He initialization) that describes the law of the Jacobian spectrum and of $c^L$ at $L/m$ fixed, and its consequences for the first steps of gradient descent.\label{ex:signalprop:researchLm}
\end{exercises}

\notesandrefs
Wide random networks became Gaussian processes with \textcite{Neal1996}, who studied priors for Bayesian networks with one hidden layer, and \textcite{Williams1997} computed the covariance functions for error-function and Gaussian units. \Textcite{ChoSaul2009} introduced the arc-cosine kernels as the covariances of infinitely wide rectifier layers and composed them into multilayer kernels, and \textcite{DanielyEtAl2016} developed the duality between activations and their kernels, including the power-series form of the correlation map. The deep Gaussian-process limit was stated by \textcite{LeeEtAl2018}, who used it for Bayesian inference on MNIST and CIFAR-10, and proved by \textcite{MatthewsEtAl2018}; the sequential proof given here is the argument behind both, made exact by \cref{lem:signalprop:conditional}. \Textcite{Hanin2023} gives a short proof for general weight distributions, and \textcite{Yang2019} extends the limit to most modern architectures.

The order--chaos picture is due to \textcite{PooleEtAl2016}, who also showed that in the chaotic phase the length of a curve of inputs grows exponentially with depth, a measure of expressivity, and the depth scales and the link to trainability to \textcite{SchoenholzEtAl2017}, who also showed that dropout destroys the critical point. \Textcite{PenningtonEtAl2017} and \textcite{XiaoEtAl2018} added the spectrum of the Jacobian and orthogonal initializations, building on the orthogonal initializations of deep linear networks of \textcite{SaxeEtAl2014}. \Textcite{HayouEtAl2019} analyze the edge of chaos for general activations, including ReLU-like ones.

The exact product structure of \cref{thm:signalprop:finite} is elementary for linear networks; for ReLU it rests on the sign symmetry used by \textcite{Hanin2018}, whose analysis, and that of \textcite{HaninNica2020}, covers the full Jacobian and variable widths. The book of \textcite{RobertsEtAl2022} develops the whole subject, criticality included, as a perturbation theory in $1/m$.\sidenote{The Hermite form of the correlation map makes the order--chaos theorem a statement about power series with nonnegative coefficients summing to one, generating functions of probability distributions; \cref{thm:signalprop:phase}(ii)--(iv) is then the extinction and survival dichotomy of a branching process, with $\chi_1$ as the mean offspring number.}

\paragraph{Translating the literature.} \Textcite{PooleEtAl2016} and \textcite{SchoenholzEtAl2017} write $N_l$ for our width $m_\ell$, $q^l_{aa}$ and $q^l_{ab}$ for our $q^\ell$ and $K^\ell(x,x')$, $\mathcal V$ for the length map, $q^*$ and $c^*$ for our fixed points $\bar q$ and $\bar c$, $c^l_{ab}$ for our $c^\ell$, and $\mathcal Dz$ for the standard Gaussian measure; their $\chi_1$ is ours, their $\chi_c$ is our $\signalpropC'(\bar c)$, their $\xi_c=-1/\log\chi_c$ is our $\xi$, and their $\xi_\nabla$ is our gradient depth scale $-1/\log\chi_1$. \Textcite{LeeEtAl2018} write $K^l(x,x')$ as we do, with $n_l$ for the width. \Textcite{RobertsEtAl2022} write $n_\ell$ for the width, $C_W$ and $C_b$ for our $\sigma_w^2$ and $\sigma_b^2$, $K^{(\ell)}$ for the kernel, and use parallel and perpendicular susceptibilities $\chi_\parallel$ and $\chi_\perp$, which at a fixed point are our $\signalpropV'(\bar q)$ and our $\chi_1$; their criticality condition $\chi_\parallel=\chi_\perp=1$ is the critical line together with marginal stability of the length. \Textcite{HaninNica2020} write $n_j$ for the widths and $\beta=\sum_j1/n_j$, which is $L/m$ for constant widths.

\printchapterbib
\end{refsection}

% Chapter 6 of The Hydrogen Atoms of Deep Learning: lazy training and the neural tangent kernel.
\casechapter{Lazy Training}{The neural tangent kernel and what it cannot learn}{Linearization around initialization}{ch:lazy}
\begin{refsection}[chapters/06-lazy.bib]

\begin{chaptersummary}
Very wide networks, trained with small steps from a standard random initialization, fit their training data while each weight barely moves, and their predictions follow those of kernel regression with a fixed kernel, the neural tangent kernel. The model is the first-order Taylor expansion of the network around its initialization. We prove that a smooth model whose Jacobian is well conditioned at initialization and varies slowly converges to zero training loss inside a small ball, that a two-layer network of width of order $n^2$ satisfies these conditions, and that training then learns the eigenmodes of the kernel in order of their eigenvalues. We then prove the cost: no method that predicts from a fixed $n$-dimensional space of functions can learn a degree-$k$ function of one unknown direction on the sphere from fewer than about $d^k/k!$ samples. Trained networks do far better because their features move.
\end{chaptersummary}

\section{The phenomenon: weights that do not move}\label{sec:lazy:phenomenon}

A network with millions of parameters is trained by gradient descent on a loss that is not convex in them. For very wide networks trained with small steps, the optimization is nonetheless simple. The training loss decreases to zero at an exponential rate, and the weights at the end are close to the weights at the start: the \emph{relative} change of each weight matrix decreases as the width grows. \Textcite{DuEtAl2019} proved that gradient descent drives the training loss of a two-layer ReLU network to zero at a linear rate when the width is a large polynomial in the number of examples and no two inputs are parallel, and \textcite{AllenZhuEtAl2019} extended such guarantees to deep fully connected, convolutional, and residual networks.

\Textcite{JacotEtAl2018} gave the phenomenon its name and its mathematics. During gradient descent the function computed by the network moves along the gradient of the loss in function space, preconditioned by a kernel, the \term{neural tangent kernel}\index{neural tangent kernel}
\[
  K_\theta(x,x')=\ip{\nabla_\theta f(x;\theta),\nabla_\theta f(x';\theta)} .
\]
In a suitable scaling, as the widths of all hidden layers go to infinity the kernel at initialization converges to a deterministic limit $K^\infty$, and it stays equal to $K^\infty$ throughout training. The trained network is then kernel regression with $K^\infty$. \Textcite{LeeEtAl2019} showed that a wide network trained by gradient descent stays close to its first-order Taylor expansion around the initialization, in theory and in experiments with fully connected, convolutional, and residual networks. Because $K^\infty$ is explicit, the infinitely wide network can be computed without training it; \textcite{AroraEtAl2019a} computed the kernel of an $11$-layer convolutional network with global average pooling and reached $77\%$ test accuracy on CIFAR-10 with kernel regression, about ten points above the best result reported before for a Gaussian-process kernel.

The kernel picture also predicts \emph{which} functions are learned fast: the error along each eigenfunction of the kernel decays at a rate equal to its eigenvalue, and these decrease with frequency. \Textcite{RahamanEtAl2019} named the resulting tendency of networks to fit low frequencies first the \term{spectral bias}\index{spectral bias}, and \textcite{BasriEtAl2019} predicted the time needed to learn each frequency for data uniform on a sphere, where the eigenfunctions are spherical harmonics, and confirmed the prediction in shallow and deep networks.

\begin{observation}[\textcite{JacotEtAl2018,LeeEtAl2019,DuEtAl2019,BasriEtAl2019}]\label[observation]{obs:lazy:lazy}
Train a network with $m$ units per hidden layer, initialized with independent Gaussian weights of variance proportional to $1/(\text{fan-in})$, by gradient descent with a small step on the square loss. As $m$ grows, the training loss converges to zero exponentially fast, the relative change $\norm{\theta_t-\theta_0}/\norm{\theta_0}$ of the parameters decreases to zero, the tangent kernel changes less and less during training, and the outputs on new inputs approach those of the network linearized at its initialization. Components of the target along eigenfunctions of the tangent kernel with larger eigenvalues are learned first.
\end{observation}

The observation is robust in the regime it describes, and it is a regime, not a law of deep learning: convolutional networks forced to stay lazy perform worse than the same networks trained normally, and the tangent kernel of networks trained in the usual way changes fast early in training (\cref{sec:lazy:wrong}).\sidenote{The term \emph{lazy training} is due to \textcite{ChizatEtAl2019}, whose note first appeared in December 2018 under the title ``A note on lazy training in supervised differentiable programming''.}

Two questions follow. Why does a nonlinear, nonconvex training problem behave like a linear one when the network is wide, and how wide is wide enough? And what does the linear behavior cost: which targets can a network that does not move its features never learn efficiently? The model of this chapter answers both.

\section{The model: linearization around initialization}\label{sec:lazy:model}

Strip the network down to what gradient descent sees. Let $x_1,\dots,x_n$ be the training inputs and $y\in\R^n$ the labels, and collect the outputs of the network on the training set into one map
\[
  F\colon\R^P\to\R^n,\qquad F(\theta)=\bigl(f(x_1;\theta),\dots,f(x_n;\theta)\bigr),
\]
with $P$ the number of parameters. The training loss is $\Lhat(\theta)=\frac12\norm{F(\theta)-y}^2$ (we do not divide by $n$, so that the kernel below has entries of order one), and we study its \term{gradient flow}\index{gradient flow}
\begin{equation}\label{eq:lazy:flow}
  \dot\theta_t=-\nabla\Lhat(\theta_t)=-J(\theta_t)^\top u_t,\qquad u_t=F(\theta_t)-y ,
\end{equation}
where $J(\theta)\in\R^{n\times P}$ is the Jacobian of $F$ and $u_t$ is the \term{residual}.\sidenote{We write $u$ for the residual because $r$ is reserved for overlaps throughout the book.} The chain rule gives the equation of motion of the outputs,
\begin{equation}\label{eq:lazy:residual}
  \dot u_t=J(\theta_t)\dot\theta_t=-K_t\,u_t,\qquad K_t=J(\theta_t)J(\theta_t)^\top\in\R^{n\times n},
\end{equation}
and the \term{tangent kernel}\index{tangent kernel} $K_t$ is the matrix of the kernel $K_{\theta_t}$ of \cref{sec:lazy:phenomenon} on the training inputs. If $K_t$ did not depend on $t$, then \eqref{eq:lazy:residual} would be a linear equation with the solution $u_t=\exp(-tK_0)u_0$, and every question about training would be a question about the spectrum of one fixed positive semidefinite matrix. This is exactly the dynamics of the linear model
\begin{equation}\label{eq:lazy:lin}
  f^{\mathrm{lin}}(x;\theta)=f(x;\theta_0)+\ip{\nabla_\theta f(x;\theta_0),\theta-\theta_0},
\end{equation}
the \term{linearized network}\index{linearized network}, whose tangent kernel is $K_{\theta_0}$ at every $\theta$. Lazy training is the statement that the network and its linearization stay close.

The general theorem of \cref{sec:lazy:convergence} is about $F$; to see when its hypotheses hold we specialize to the simplest network.

\begin{assumption}[Two-layer network]\label[assumption]{ass:lazy:twolayer}
The inputs are unit vectors $x_1,\dots,x_n\in\Sph^{d-1}$, no two of which are parallel ($x_i\ne\pm x_j$ for $i\ne j$), and the labels satisfy $\abs{y_i}\le1$. The network is
\[
  f(x;W)=\frac1{\sqrt m}\sum_{r=1}^ma_r\,\phi(\ip{w_r,x}),
\]
with first-layer weights $w_1,\dots,w_m\in\R^d$, the rows of $W\in\R^{m\times d}$, and output weights $a_r\in\{-1,1\}$ that are not trained. At initialization the $w_r$ are independent $\cN(0,I_d)$ and the $a_r$ are independent uniform signs, independent of $W$. The activation $\phi$ is twice continuously differentiable, is not affine, and satisfies $\abs{\phi'}\le\beta$ and $\abs{\phi''}\le\beta$.
\end{assumption}

The activation $\tanh$ qualifies with $\beta=1$; ReLU does not, because its derivative jumps, although the conclusions survive with more work \parencite{DuEtAl2019}. The factor $m^{-1/2}$ is the \term{NTK parameterization}\index{NTK parameterization}: with $\ip{w_r,x}\sim\cN(0,1)$ each term is of order one, the output at initialization is a sum of $m$ independent terms of mean zero, and the scale makes it of order one.\sidenote{In the standard parameterization the factor $m^{-1/2}$ sits in the variance of the output weights instead. The two give the same network at initialization but gradients of different sizes, so the same step size means different things; \textcite{YangHu2021} classify these choices.} For this network $\partial f(x;W)/\partial w_r=m^{-1/2}a_r\phi'(\ip{w_r,x})\,x$ and $a_r^2=1$, so the tangent kernel and its expectation at initialization are
\begin{equation}\label{eq:lazy:kinf}
\begin{gathered}
  K_W(x,x')=\frac{\ip{x,x'}}m\sum_{r=1}^m\phi'(\ip{w_r,x})\phi'(\ip{w_r,x'}),\\
  K^\infty(x,x')=\ip{x,x'}\,\kappa(\ip{x,x'}),
\end{gathered}
\end{equation}
where $\kappa(\rho)=\E[\phi'(g)\phi'(g')]$ for standard Gaussians $g,g'$ with correlation $\rho$: for unit $x,x'$ the pair $(\ip{w,x},\ip{w,x'})$ is such a pair with $\rho=\ip{x,x'}$. By the law of large numbers $K_{W_0}\to K^\infty$ as $m\to\infty$. We write $K^\infty$ also for the $n\times n$ matrix $(K^\infty(x_i,x_j))_{ij}$, and $\lmin=\lmin(K^\infty)$.

\begin{dictionary}
parameters $\theta\in\R^P$, map $F$ & all weights; the outputs on the training set\\
Jacobian $J(\theta_0)$ & per-example gradients at initialization\\
tangent kernel $K_t=J_tJ_t^\top$ & empirical neural tangent kernel on the training set\\
limit kernel $K^\infty$ & kernel of the infinitely wide architecture (fully connected, convolutional, \dots)\\
$\lmin(K^\infty)$ & how distinguishable the training inputs are to the architecture\\
residual $u_t$ & training errors\\
linearized model \eqref{eq:lazy:lin} & the network in the lazy regime\\
eigenvectors of $K^\infty$ & modes learned in order of their eigenvalues: spectral bias\\
width $m$, scale $m^{-1/2}$ & hidden-layer width; NTK versus standard or mean-field parameterization\\
a fixed space of $n$ functions & a kernel method, random features, or a lazy network\\
\end{dictionary}

\paragraph{What was thrown away, and why first.} The features: the model is linear in the parameters, so the functions $x\mapsto\nabla_\theta f(x;\theta_0)$ are fixed before training, and learning is regression on them. Finite step sizes and stochastic gradients: we study the gradient flow. The second layer and depth: in \cref{ass:lazy:twolayer} only one layer is trained (\cref{ex:lazy:bothlayers} trains both, and \cref{thm:lazy:deepntk} states the kernel of deep networks). Each simplification moves toward a model in which the parameters move as little as the data allow; that is the right first model because there a nonconvex problem becomes solvable, and \cref{sec:lazy:kernels} computes the price.

\paragraph{What one had to notice.} Two facts, one about scale and one about geometry. The gradient of a wide network with respect to each single weight is small, of order $m^{-1/2}$, while the outputs depend on all $m$ weights; so the outputs can change by order one while every weight changes by order $m^{-1/2}$. And the Hessian of the outputs, which is what makes $K_t$ move, is block diagonal with one block of norm $O(m^{-1/2})$ per neuron, so it vanishes in the limit while the kernel, a sum of $m$ terms of order $1/m$, does not. A third fact makes the model useful: the limit kernel is computable in closed form, by the Gaussian recursions of signal propagation (\cref{ch:signalprop}).

\section{The solution: convergence without moving}\label{sec:lazy:convergence}

The first theorem isolates what is needed of $F$: a Jacobian that is well conditioned at initialization and changes slowly over a ball whose radius is set by the initial error.

\begin{theorem}[Lazy training]\label{thm:lazy:lazy}\index{lazy training}
Let $F\colon\R^P\to\R^n$ be continuously differentiable with Jacobian $J$, let $y\in\R^n$ and $\theta_0\in\R^P$, and let $u_0=F(\theta_0)-y$. Let $\mu>0$, let $R=\norm{u_0}/\mu$, and let $B$ be the closed ball of radius $R$ around $\theta_0$. Assume
\begin{enumerate}
\item $\lmin\bigl(J(\theta_0)J(\theta_0)^\top\bigr)\ge4\mu^2$;
\item $\norm{J(\theta)-J(\theta')}_{\op}\le L_J\norm{\theta-\theta'}$ for all $\theta,\theta'\in B$;
\item $L_JR\le\mu$.
\end{enumerate}
Then the gradient flow \eqref{eq:lazy:flow} started at $\theta_0$ exists for all $t\ge0$, and for every $t\ge0$,
\begin{equation}\label{eq:lazy:lazy}
  \norm{u_t}\le e^{-\mu^2t}\norm{u_0},\qquad \norm{\theta_t-\theta_0}\le\frac{\norm{u_0}-\norm{u_t}}\mu .
\end{equation}
In particular $\theta_t$ stays in $B$, and it converges to a point $\theta_\infty\in B$ with $F(\theta_\infty)=y$.
\end{theorem}

\begin{proof}
The plan is to show that the kernel $K_t$ stays at least $\mu^2$ on the ball, which makes the residual decay, and that the path length is controlled by the decrease of the residual, which keeps the flow inside the ball.

\emph{The degenerate case.} If $u_0=0$, then $\theta_0$ is a stationary point of \eqref{eq:lazy:flow}, the flow is constant, and \eqref{eq:lazy:lazy} holds with $R=0$. From now on $u_0\ne0$, so $R>0$.

\emph{The Jacobian stays well conditioned on $B$.} Let $\theta\in B$ and $v\in\R^n$. Hypothesis (i) gives $\norm{J(\theta_0)^\top v}^2=v^\top J(\theta_0)J(\theta_0)^\top v\ge4\mu^2\norm v^2$, and hypotheses (ii) and (iii) give $\norm{J(\theta)-J(\theta_0)}_{\op}\le L_JR\le\mu$. By the triangle inequality,
\begin{equation}\label{eq:lazy:cond}
  \norm{J(\theta)^\top v}\ge\norm{J(\theta_0)^\top v}-\norm{(J(\theta)-J(\theta_0))^\top v}\ge2\mu\norm v-\mu\norm v=\mu\norm v .
\end{equation}

\emph{Local existence.} On $B$ the map $J$ is Lipschitz by (ii), hence bounded, say $\norm{J(\theta)}_{\op}\le b$, and $F$ is then Lipschitz with constant $b$. So the vector field $\theta\mapsto-J(\theta)^\top(F(\theta)-y)$ is Lipschitz on $B$, and by the Picard--Lindel\"of theorem the flow has a unique solution on a maximal interval $[0,T^*)$ during which it stays in the open ball $B^\circ$; if $T^*<\infty$, then $\theta_t$ leaves every compact subset of $B^\circ$ as $t\uparrow T^*$.

\emph{The residual decays.} For $t<T^*$, \eqref{eq:lazy:residual} and \eqref{eq:lazy:cond} give
\[
  \frac{\dd}{\dd t}\norm{u_t}^2=-2u_t^\top K_tu_t=-2\norm{J(\theta_t)^\top u_t}^2\le-2\mu^2\norm{u_t}^2 ,
\]
so $\norm{u_t}\le e^{-\mu^2t}\norm{u_0}$. The same identity and $\norm{J(\theta_t)^\top u_t}\le b\norm{u_t}$ give $\frac{\dd}{\dd t}\norm{u_t}^2\ge-2b^2\norm{u_t}^2$, hence $\norm{u_t}\ge e^{-b^2t}\norm{u_0}>0$.

\emph{The path length is bounded by the decrease of the residual.} Since $u_t\ne0$, the norm $\norm{u_t}$ is differentiable, and \eqref{eq:lazy:cond} with $v=u_t$ gives
\[
  \frac{\dd}{\dd t}\norm{u_t}=-\frac{\norm{J(\theta_t)^\top u_t}^2}{\norm{u_t}}\le-\mu\norm{J(\theta_t)^\top u_t}=-\mu\norm{\dot\theta_t}.
\]
Integrating from $0$ to $t$ gives $\norm{\theta_t-\theta_0}\le\int_0^t\norm{\dot\theta_s}\dd s\le(\norm{u_0}-\norm{u_t})/\mu$, which is \eqref{eq:lazy:lazy} on $[0,T^*)$.

\emph{The flow never leaves.} Suppose $T^*<\infty$. By the last two steps, for $t<T^*$,
\[
  \norm{\theta_t-\theta_0}\le R-\frac{\norm{u_t}}\mu\le R-\frac{e^{-b^2T^*}\norm{u_0}}\mu ,
\]
so $\theta_t$ stays in a closed ball of radius strictly less than $R$, a compact subset of $B^\circ$, which contradicts the maximality of $T^*$. Hence $T^*=\infty$ and \eqref{eq:lazy:lazy} holds for all $t$.

\emph{Convergence.} The total path length $\int_0^\infty\norm{\dot\theta_s}\dd s$ is at most $\norm{u_0}/\mu$, so $\theta_t$ is Cauchy as $t\to\infty$ and converges to some $\theta_\infty\in B$. By continuity of $F$, $F(\theta_\infty)-y=\lim_tu_t=0$.
\end{proof}

The hypothesis is not convexity but $\norm{\nabla\Lhat}^2\ge2\mu^2\Lhat$ on $B$, a Polyak--{\L}ojasiewicz inequality, on which \textcite{LiuEtAl2022} build a theory of overparameterized optimization. The radius $R=\norm{u_0}/\mu$ does not involve $L_J$: the parameters travel a distance set by the initial error and the conditioning, while the curvature only has to be small \emph{on that distance}. Wide networks will make $R$ independent of the width and $L_J$ of order $m^{-1/2}$.

Under the same hypotheses the network stays close to its linearization. Let $\theta^{\mathrm{lin}}_t$ be the gradient flow of $\frac12\norm{F^{\mathrm{lin}}(\theta)-y}^2$ with $F^{\mathrm{lin}}(\theta)=F(\theta_0)+J(\theta_0)(\theta-\theta_0)$, from $\theta_0$, and $u^{\mathrm{lin}}_t$ its residual. Its kernel is the constant $K_0=J(\theta_0)J(\theta_0)^\top$, so $u^{\mathrm{lin}}_t=\exp(-tK_0)u_0$.

\begin{corollary}[Closeness to the linearized model]\label[corollary]{cor:lazy:linear}
Under the hypotheses of \cref{thm:lazy:lazy}, let $\beta_0=\norm{J(\theta_0)}_{\op}$. For every $t\ge0$,
\[
  \norm{u_t-u^{\mathrm{lin}}_t}\le\frac{\beta_0L_J\norm{u_0}^2}{\mu^3}\,e^{-\mu^2t},\qquad
  \norm{\theta_t-\theta^{\mathrm{lin}}_t}\le\frac{2\beta_0^2L_J\norm{u_0}^2}{\mu^5}.
\]
If, moreover, a test input $x$ has $\norm{\nabla_\theta f(x;\theta)-\nabla_\theta f(x;\theta')}\le L_x\norm{\theta-\theta'}$ on $B$, then $\abs{f(x;\theta_t)-f^{\mathrm{lin}}(x;\theta^{\mathrm{lin}}_t)}\le\frac12L_xR^2+\norm{\nabla_\theta f(x;\theta_0)}\,\norm{\theta_t-\theta^{\mathrm{lin}}_t}$.
\end{corollary}

\begin{proof}
The plan is to write the difference of the two residuals by variation of constants around the linear equation and to bound the forcing term by the motion of the kernel.

\emph{The kernel moves little.} Hypothesis (i) gives $4\mu^2\le\lmin(K_0)\le\beta_0^2$, so $\mu\le\beta_0/2$. By \cref{thm:lazy:lazy}, $\norm{\theta_s-\theta_0}\le R$, so $\epsilon\defeq\sup_s\norm{J(\theta_s)-J(\theta_0)}_{\op}\le L_JR\le\mu$. Writing $J_s=J(\theta_s)$ and $J_0=J(\theta_0)$,
\[
  \norm{K_s-K_0}_{\op}=\norm{(J_s-J_0)J_s^\top+J_0(J_s-J_0)^\top}_{\op}\le\epsilon(2\beta_0+\epsilon),
\]
which is at most $3\beta_0L_JR$.

\emph{Variation of constants.} The difference $\delta_t=u_t-u^{\mathrm{lin}}_t$ satisfies $\delta_0=0$ and $\dot\delta_t=-K_0\delta_t-(K_t-K_0)u_t$, so $\delta_t=-\int_0^t\exp(-(t-s)K_0)(K_s-K_0)u_s\dd s$. Since $K_0\succeq4\mu^2I_n$, $\norm{\exp(-(t-s)K_0)}_{\op}\le e^{-4\mu^2(t-s)}$, and with $\norm{u_s}\le e^{-\mu^2s}\norm{u_0}$,
\begin{align*}
  \norm{\delta_t}&\le3\beta_0L_JR\norm{u_0}\int_0^te^{-4\mu^2(t-s)}e^{-\mu^2s}\dd s\\
  &\le3\beta_0L_JR\norm{u_0}\,\frac{e^{-\mu^2t}}{3\mu^2}=\frac{\beta_0L_J\norm{u_0}^2}{\mu^3}e^{-\mu^2t}.
\end{align*}

\emph{Parameters.} Integrating $\dot\theta_s-\dot\theta^{\mathrm{lin}}_s=-(J_s-J_0)^\top u_s-J_0^\top\delta_s$ and using the bounds above,
\begin{align*}
  \norm{\theta_t-\theta^{\mathrm{lin}}_t}&\le\int_0^\infty\Bigl(L_JR\norm{u_s}+\beta_0\norm{\delta_s}\Bigr)\dd s\\
  &\le\frac{L_J\norm{u_0}^2}{\mu^3}+\frac{\beta_0^2L_J\norm{u_0}^2}{\mu^5}\le\frac{2\beta_0^2L_J\norm{u_0}^2}{\mu^5},
\end{align*}
the last step by $\mu\le\beta_0$.

\emph{A test input.} Taylor's theorem with a Lipschitz gradient bounds the difference between $f(x;\theta_t)$ and $f^{\mathrm{lin}}(x;\theta_t)$ by $\frac12L_x\norm{\theta_t-\theta_0}^2\le\frac12L_xR^2$. The linearized model changes by $\ip{\nabla_\theta f(x;\theta_0),\theta_t-\theta^{\mathrm{lin}}_t}$ between $\theta^{\mathrm{lin}}_t$ and $\theta_t$.
\end{proof}

The linearized model is solved in closed form, and its solution is kernel regression.

\begin{proposition}[Lazy training is kernel regression]\label[proposition]{prop:lazy:kernelreg}
Suppose $K_0=J(\theta_0)J(\theta_0)^\top$ is invertible, and write $K_0(x,X)$ for the row vector $(K_{\theta_0}(x,x_i))_{i\le n}$. The linearized flow satisfies, for every $t\ge0$ and every input $x$,
\begin{equation}\label{eq:lazy:krr}
  f^{\mathrm{lin}}(x;\theta^{\mathrm{lin}}_t)=f(x;\theta_0)-K_0(x,X)\,K_0^{-1}\bigl(I_n-\exp(-tK_0)\bigr)u_0 .
\end{equation}
As $t\to\infty$ it converges to the kernel interpolant $f(x;\theta_0)+K_0(x,X)K_0^{-1}(y-F(\theta_0))$. The parameters move by $\norm{\theta^{\mathrm{lin}}_t-\theta_0}^2=u_0^\top K_0^{-1}(I_n-\exp(-tK_0))^2u_0$, which increases to $u_0^\top K_0^{-1}u_0$, the smallest $\norm{\theta-\theta_0}^2$ among all $\theta$ with $F^{\mathrm{lin}}(\theta)=y$.
\end{proposition}

\begin{proof}
With $J_0=J(\theta_0)$, the linearized flow is $\dot\theta=-J_0^\top\exp(-tK_0)u_0$, so $\theta^{\mathrm{lin}}_t-\theta_0=-J_0^\top\int_0^t\exp(-sK_0)\dd s\,u_0=-J_0^\top K_0^{-1}(I_n-\exp(-tK_0))u_0$. Since $\nabla_\theta f(x;\theta_0)^\top J_0^\top=K_0(x,X)$, inserting this in \eqref{eq:lazy:lin} gives \eqref{eq:lazy:krr}. The squared norm is $v^\top K_0^{-1}J_0J_0^\top K_0^{-1}v=v^\top K_0^{-1}v$ with $v=(I_n-\exp(-tK_0))u_0$; in an eigenbasis of $K_0$ with eigenvalues $\lambda_k>0$ it is $\sum_k(1-e^{-\lambda_kt})^2c_k^2/\lambda_k$, with $c_k$ the coordinates of $u_0$, which increases in $t$. Finally every $\Delta$ with $J_0\Delta=-u_0$ splits as $\Delta=J_0^\top c+\Delta^\perp$ with $J_0\Delta^\perp=0$; then $K_0c=-u_0$, and $\norm\Delta^2=u_0^\top K_0^{-1}u_0+\norm{\Delta^\perp}^2$.
\end{proof}

The interpolant is the kernel analogue of the minimum-norm least-squares estimator of \cref{ch:doubledescent}: among all functions in the span of the tangent features that fit the residual labels, gradient flow from $\theta_0$ picks the one of smallest norm, here the norm of the reproducing kernel Hilbert space of $K_0$.

We now check the hypotheses of \cref{thm:lazy:lazy} for the two-layer network. Everything rests on one fact about the limit kernel.

\begin{lemma}[The limit kernel is positive definite]\label[lemma]{lem:lazy:pd}
Under \cref{ass:lazy:twolayer}, the matrix $K^\infty=(K^\infty(x_i,x_j))_{i,j\le n}$ is positive definite.
\end{lemma}

\begin{proof}
The plan is to expand $K^\infty$ in tensor powers of the inputs and to show that a vector annihilated by infinitely many of them vanishes.

\emph{Expansion.} Since $\phi'$ is bounded, it has a Hermite expansion $\phi'=\sum_kb_kh_k$ in $L^2(\cN(0,1))$ with $\sum_kb_k^2=\E\phi'(g)^2\le\beta^2$, and \cref{cor:toolkit:dual} gives $\kappa(\rho)=\sum_kb_k^2\rho^k$. Hence $K^\infty_{ij}=\sum_kb_k^2\ip{x_i,x_j}^{k+1}$, and since $\ip{x_i,x_j}^{p}=\ip{x_i^{\otimes p},x_j^{\otimes p}}$, for every $v\in\R^n$
\[
  v^\top K^\infty v=\sum_{k\ge0}b_k^2\,\Bigl\|\sum_{i=1}^nv_i\,x_i^{\otimes(k+1)}\Bigr\|^2\ge0 .
\]

\emph{Infinitely many terms.} If only finitely many $b_k$ were nonzero, then $\phi'$ would agree almost everywhere, hence everywhere by continuity, with a polynomial; a bounded polynomial is constant, and $\phi$ would be affine, which \cref{ass:lazy:twolayer} excludes. So the set $S=\{k\colon b_k\ne0\}$ is infinite.

\emph{A vector in the kernel vanishes.} Suppose $v^\top K^\infty v=0$. Then $\sum_iv_ix_i^{\otimes(k+1)}=0$ for each $k\in S$, and pairing with $z^{\otimes(k+1)}$ gives $\sum_iv_i\ip{x_i,z}^{k+1}=0$ for every $z\in\R^d$. The vectors $z$ with $\ip{x_i,z}=0$ for some $i$, or with $\ip{x_i,z}=\pm\ip{x_j,z}$ for some $i\ne j$, lie in finitely many hyperplanes, the orthogonal complements of the nonzero vectors $x_i$ and $x_i\mp x_j$ (nonzero because no two inputs are parallel); these do not cover $\R^d$. Choose $z$ outside them and relabel the inputs so that $c_i=\abs{\ip{x_i,z}}$ satisfies $c_1>c_2>\dots>c_n>0$, with signs $\epsilon_i=\sgn\ip{x_i,z}$. If $v\ne0$, let $i_0$ be the first index with $v_{i_0}\ne0$; dividing by $c_{i_0}^{k+1}$,
\[
  v_{i_0}\epsilon_{i_0}^{k+1}=-\sum_{i>i_0}v_i\epsilon_i^{k+1}\Bigl(\frac{c_i}{c_{i_0}}\Bigr)^{k+1}\qquad(k\in S),
\]
and the right side tends to $0$ as $k\to\infty$ in $S$, while the left side has absolute value $\abs{v_{i_0}}>0$. This contradiction shows $v=0$.
\end{proof}

The hypothesis on parallel inputs is needed: for $\tanh$ the function $\phi'$ is even, $\kappa(-1)=\kappa(1)$, and an input $x$ together with $-x$ gives two proportional rows of $K^\infty$ (\cref{ex:lazy:pd}).\sidenote{\Textcite{DuEtAl2019} prove positive definiteness for ReLU under the same hypothesis that no two inputs are parallel.} Positivity is qualitative. How large $\lmin$ is depends on the data and the activation, and it can be very small: for $32$ equally spaced points on a circle with a bias coordinate and $\phi=\tanh$, the setting of \cref{fig:lazy:modes}, $\lmin\approx2.6\times10^{-8}$, because the eigenvalues of a smooth kernel decay fast with frequency.

\begin{lemma}[The network at initialization]\label[lemma]{lem:lazy:init}
Under \cref{ass:lazy:twolayer}, let $J(W)\in\R^{n\times md}$ be the Jacobian of $W\mapsto(f(x_i;W))_{i\le n}$, let $c_\phi=\E\phi(g)^2$ for $g\sim\cN(0,1)$, and let $\delta\in(0,1)$.
\begin{enumerate}
\item For all $W,W'$, $\norm{J(W)}_{\op}\le\beta\sqrt n$ and $\norm{J(W)-J(W')}_{\op}\le\beta\sqrt{n/m}\,\norm{W-W'}_F$.
\item If $m\ge\frac{8\beta^4n^2}{\lmin^2}\log\frac{4n^2}\delta$, then with probability at least $1-\frac\delta2$, $\norm{K_0-K^\infty}_{\op}\le\frac{\lmin}2$, and hence $\lmin(K_0)\ge\frac{\lmin}2$.
\item With probability at least $1-\frac\delta2$, $\norm{u_0}^2\le2n(c_\phi+1)/\delta$.
\end{enumerate}
\end{lemma}

\begin{proof}
\emph{Part (i).} The entry of $J(W)$ in row $i$ and in the block of coordinates of $w_r$ is $m^{-1/2}a_r\phi'(\ip{w_r,x_i})x_i^\top$. With $\abs{\phi'}\le\beta$, $a_r^2=1$, and $\norm{x_i}=1$, the squared Frobenius norm is at most $\sum_i\sum_r\beta^2/m=n\beta^2$, which bounds the operator norm. For the difference, $\abs{\phi'(s)-\phi'(s')}\le\beta\abs{s-s'}$ and the Cauchy--Schwarz inequality $\abs{\ip{w_r-w_r',x_i}}\le\norm{w_r-w_r'}$ give
\[
  \norm{J(W)-J(W')}_F^2=\frac1m\sum_{i=1}^n\sum_{r=1}^m\bigl(\phi'(\ip{w_r,x_i})-\phi'(\ip{w_r',x_i})\bigr)^2\le\frac{n\beta^2}m\norm{W-W'}_F^2 .
\]

\emph{Part (ii).} By \eqref{eq:lazy:kinf}, $K_0=\frac1m\sum_rZ_r$, where the matrix $Z_r$ has the entries $(Z_r)_{ij}=\ip{x_i,x_j}\phi'(\ip{w_r,x_i})\phi'(\ip{w_r,x_j})$. The matrices $Z_r$ are independent, have mean $K^\infty$, and have entries in $[-\beta^2,\beta^2]$. Hoeffding's inequality (\cref{thm:toolkit:hoeffding}) for the entry $(i,j)$, with deviation $ms$ and $m$ summands of range $2\beta^2$, gives
\[
  \Pr\Bigl(\abs{(K_0-K^\infty)_{ij}}\ge s\Bigr)\le2\exp\Bigl(-\frac{2m^2s^2}{m(2\beta^2)^2}\Bigr)=2\exp\Bigl(-\frac{ms^2}{2\beta^4}\Bigr).
\]
Take $s=\lmin/(2n)$ and a union bound over the $n^2$ entries: with probability at least $1-2n^2\exp(-m\lmin^2/(8n^2\beta^4))\ge1-\frac\delta2$, every entry is below $s$, so $\norm{K_0-K^\infty}_{\op}\le\norm{K_0-K^\infty}_F\le ns=\lmin/2$. Weyl's inequality then gives $\lmin(K_0)\ge\lmin-\lmin/2$.

\emph{Part (iii).} The output $f(x_i;W_0)=m^{-1/2}\sum_ra_r\phi(\ip{w_r,x_i})$ is a normalized sum of $m$ independent terms; each has mean zero, because $a_r$ is a symmetric sign independent of $w_r$, and second moment $c_\phi$, which is finite because $\abs{\phi(s)}\le\abs{\phi(0)}+\beta\abs s$. Hence $\E f(x_i;W_0)=0$, $\E f(x_i;W_0)^2=c_\phi$, and $\E\norm{u_0}^2=\sum_i(c_\phi+y_i^2)\le n(c_\phi+1)$. Markov's inequality gives (iii).
\end{proof}

\begin{theorem}[Wide two-layer networks train lazily]\label{thm:lazy:twolayer}
Under \cref{ass:lazy:twolayer}, let $\delta\in(0,1)$ and suppose
\begin{equation}\label{eq:lazy:width}
  m\ge\frac{n^2}{\lmin^2}\max\Bigl\{8\beta^4\log\frac{4n^2}\delta,\ \frac{128\beta^2(c_\phi+1)}\delta\Bigr\}.
\end{equation}
With probability at least $1-\delta$ over the initialization, the gradient flow of $\Lhat(W)=\frac12\sum_i(f(x_i;W)-y_i)^2$ from $W_0$ satisfies, for all $t\ge0$,
\begin{align*}
  \norm{u_t}&\le e^{-\lmin t/8}\norm{u_0},\\
  \norm{W_t-W_0}_F&\le\sqrt{8/\lmin}\,\norm{u_0}\le4\sqrt{n(c_\phi+1)/(\lmin\delta)},\\
  \max_{r\le m}\norm{w_r(t)-w_r(0)}&\le8\beta\sqrt n\,\norm{u_0}/(\lmin\sqrt m),\\
  \norm{K_t-K_0}_{\op}&\le3\sqrt8\,\beta^2n\norm{u_0}/\sqrt{\lmin m} .
\end{align*}
\end{theorem}

\begin{proof}
The plan is to apply \cref{thm:lazy:lazy} with $\mu^2=\lmin/8$ on the event where parts (ii) and (iii) of \cref{lem:lazy:init} hold, which has probability at least $1-\delta$ by a union bound.

\emph{The hypotheses.} The first term of \eqref{eq:lazy:width} and \cref{lem:lazy:init}(ii) give $\lmin(K_0)\ge\lmin/2=4\mu^2$, hypothesis (i). \Cref{lem:lazy:init}(i) gives hypothesis (ii) on all of $\R^{m\times d}$ with $L_J=\beta\sqrt{n/m}$. For hypothesis (iii), $L_JR=\beta\sqrt{n/m}\,\norm{u_0}/\mu\le\mu$ is equivalent to $m\ge n\beta^2\norm{u_0}^2/\mu^4=64n\beta^2\norm{u_0}^2/\lmin^2$, which holds by \cref{lem:lazy:init}(iii) and the second term of \eqref{eq:lazy:width}.

\emph{The conclusions.} \Cref{thm:lazy:lazy} gives the decay of $u_t$ and $\norm{W_t-W_0}_F\le R=\norm{u_0}/\mu$, and \cref{lem:lazy:init}(iii) bounds $\norm{u_0}$. For one neuron, $\dot w_r=-m^{-1/2}a_r\sum_iu_{t,i}\phi'(\ip{w_r,x_i})x_i$, so $\norm{\dot w_r}\le\beta m^{-1/2}\sum_i\abs{u_{t,i}}\le\beta\sqrt{n/m}\,\norm{u_t}$; integrating $e^{-\mu^2t}\norm{u_0}$ over $t\ge0$ gives the factor $1/\mu^2=8/\lmin$. For the kernel, the first step of the proof of \cref{cor:lazy:linear} gives $\norm{K_t-K_0}_{\op}\le3\beta_0L_JR$ with $\beta_0\le\beta\sqrt n$ by \cref{lem:lazy:init}(i).
\end{proof}

Read the theorem as a statement about scales. The displacement $\norm{W_t-W_0}_F$ is bounded independently of $m$, while $\norm{W_0}_F\approx\sqrt{md}$, so the relative change of the weights is of order $m^{-1/2}$; so is the largest move of a single neuron, and so is the change of the kernel. The exponents in the width $n^2/\lmin^2$ are not sharp: \textcite{DuEtAl2019} needed a larger power of $n$ for ReLU, and \textcite{OymakSoltanolkotabi2020} reduced the requirement for smooth activations to a number of parameters of order $n^2$. \Cref{sec:lazy:reading} shows that the kernel in fact moves by order $1/m$, not $m^{-1/2}$.

\section{The solution: what a kernel can and cannot learn}\label{sec:lazy:kernels}

\paragraph{Learning by modes.} On the training set the answer is in \eqref{eq:lazy:krr}: if $v_k$ is an eigenvector of $K_0$ with eigenvalue $\lambda_k$, then $\ip{v_k,u^{\mathrm{lin}}_t}=e^{-\lambda_kt}\ip{v_k,u_0}$, the normal modes of the quadratic of \hyperref[ch:prologue]{the prologue}. For the population, let $x\sim\Unif(\Sph^{d-1})$, let $K(x,x')=\psi(\ip{x,x'})$ be a bounded rotation-invariant kernel, and let $T_K$ be its integral operator, $(T_Kg)(x)=\E_{x'}[K(x,x')g(x')]$. The \term{population linearized flow} is the gradient flow of $\frac12\norm{f-f_\star}^2$ in the geometry of the kernel,
\begin{equation}\label{eq:lazy:popflow}
  \partial_tf_t=-T_K(f_t-f_\star),
\end{equation}
the limit $n\to\infty$ of the function-space flow $\partial_tf_t(x)=-\frac1n\sum_iK(x,x_i)(f_t(x_i)-y_i)$ for the loss $\frac1{2n}\sum_i(f(x_i)-y_i)^2$ with $y_i=f_\star(x_i)$.

\begin{proposition}[Spectral bias on the sphere]\label[proposition]{prop:lazy:spectral}\index{spectral bias}
Let $\Pi_k$ be the orthogonal projection of $L^2(\Sph^{d-1})$ onto the spherical harmonics $\cH_k$ of degree $k$, and let $\lambda_k=\frac{\abs{\Sph^{d-2}}}{\abs{\Sph^{d-1}}}\int_{-1}^1\psi(t)P_{d,k}(t)(1-t^2)^{(d-3)/2}\dd t$, with $P_{d,k}$ the Gegenbauer polynomial normalized by $P_{d,k}(1)=1$. The solution of \eqref{eq:lazy:popflow} with $f_0-f_\star\in L^2$ satisfies $\Pi_k(f_t-f_\star)=e^{-\lambda_kt}\,\Pi_k(f_0-f_\star)$ for every $k\ge0$ and $t\ge0$.
\end{proposition}

\begin{proof}
The operator $T_K$ is bounded on $L^2$, with norm at most $\sup\abs\psi$, so \eqref{eq:lazy:popflow} is a linear equation with a bounded operator and the unique solution $f_t-f_\star=\exp(-tT_K)(f_0-f_\star)$. By the Funk--Hecke formula (\cref{thm:toolkit:addition}(iii)), $T_KY=\lambda_kY$ for every $Y\in\cH_k$, so $T_K$ commutes with $\Pi_k$ and $\exp(-tT_K)\Pi_k=e^{-\lambda_kt}\Pi_k$.
\end{proof}

For the two-layer network, $K^\infty(x,x')=\rho\,\kappa(\rho)$ with $\rho=\ip{x,x'}$ and $\kappa(\rho)=\sum_jb_j^2\rho^j$, so the eigenvalues are determined by the Hermite coefficients of $\phi'$, and their decay by the smoothness of $\phi$. For ReLU, $\kappa(\rho)=(\pi-\arccos\rho)/(2\pi)$, a classical arc-cosine kernel \parencite{ChoSaul2009}, and the eigenvalues decay like $k^{-d}$ \parencite{BiettiMairal2019};\sidenote{For $d=3$ we computed $\lambda_4/\lambda_6=3.47$ by quadrature, against $(6/4)^3=3.375$.} for smooth activations they decay faster. The parity of $\phi'$ matters too. For ReLU, $\rho\kappa(\rho)=\frac\rho4+\frac{\rho\arcsin\rho}{2\pi}$ is a degree-one term plus an even function, so $\lambda_k=0$ for every odd $k\ge3$ (\cref{ex:lazy:relu}): a two-layer ReLU network without bias cannot learn odd harmonics of degree three or more in the lazy regime at all \parencite{BasriEtAl2019}. A bias, implemented as a constant input coordinate, removes these zeros.

\paragraph{What no fixed space of functions can learn.} High-degree harmonics are not only slow but expensive in samples, for every kernel, by a counting argument.

\begin{theorem}[What a kernel cannot learn]\label{thm:lazy:lower}\index{kernel lower bound}
Let $d\ge3$, $k\ge0$, and $N=N_{d,k}=\dim\cH_k$. For $w\in\Sph^{d-1}$ let
\[
  f_w(x)=\frac{Z_k(x,w)}{\sqrt N}=\sqrt N\,P_{d,k}(\ip{w,x}),
\]
with $Z_k$ the reproducing kernel of $\cH_k$ (\cref{thm:toolkit:addition}(ii)), so that $\norm{f_w}_{L^2}=1$. Let $w\sim\Unif(\Sph^{d-1})$, and let $V\subset L^2(\Sph^{d-1})$ be a random subspace of dimension at most $n$, spanned by a random orthonormal family independent of $w$. Then
\[
  \E\norm{f_w-\Pi_Vf_w}^2=1-\frac{\E\tr(\Pi_V\Pi_k)}N\ge1-\frac nN .
\]
Consequently, every estimator $\hat f$ that takes values in $V$, whatever it does with the labels, satisfies $\E\norm{\hat f-f_w}^2\ge1-n/N$.
\end{theorem}

\begin{proof}
The plan is to compute the energy of $f_w$ in one fixed direction, averaged over $w$, with the reproducing property, and to add over an orthonormal basis of $V$.

\emph{Normalization.} The reproducing property gives $\norm{Z_k(\cdot,w)}^2=Z_k(w,w)=N$, so $\norm{f_w}=1$; the second expression for $f_w$ is $Z_k(x,w)=N\,P_{d,k}(\ip{x,w})$ from \cref{thm:toolkit:addition}(ii).

\emph{One direction.} Let $e\in L^2$. Since $Z_k$ is the kernel of $\Pi_k$, $\ip{e,f_w}=N^{-1/2}\int e(x)Z_k(x,w)\dd x=N^{-1/2}(\Pi_ke)(w)$, the integral being against the uniform probability measure. Averaging over $w\sim\Unif(\Sph^{d-1})$,
\[
  \E_w\ip{e,f_w}^2=\frac1N\norm{\Pi_ke}^2\le\frac{\norm e^2}N .
\]

\emph{The subspace.} Condition on $V$, which is independent of $w$, and let $e_1,\dots,e_{n'}$, $n'\le n$, be its orthonormal basis. Then $\norm{f_w-\Pi_Vf_w}^2=1-\sum_j\ip{e_j,f_w}^2$, and the previous step gives $\E_w\sum_j\ip{e_j,f_w}^2=\frac1N\sum_j\norm{\Pi_ke_j}^2=\frac1N\tr(\Pi_V\Pi_k)\le\frac{n'}N$. Take expectations over $V$.

\emph{Estimators.} If $\hat f\in V$, then $\hat f-\Pi_Vf_w\in V$ is orthogonal to $f_w-\Pi_Vf_w$, so $\norm{\hat f-f_w}^2\ge\norm{f_w-\Pi_Vf_w}^2$.
\end{proof}

The theorem applies to every method that predicts from a fixed space of at most $n$ functions. Kernel ridge regression with \emph{every} bounded kernel predicts in the span of $K(\cdot,x_i)$, which depends on the inputs only; the linearized network of \cref{prop:lazy:kernelreg}, started independently of $w$, predicts in $f(\cdot;\theta_0)+\mathrm{span}\{K_0(\cdot,x_i)\}$; a random-features model with $p$ features predicts in a space of dimension $p$ whatever $n$ is. Each therefore needs $n$ (or $p$) of order $N_{d,k}=\frac{d^k}{k!}(1+O(1/d))$ to reach error $\frac12$ on a target that depends on one unknown direction through a polynomial of degree $k$. For $d=30$ that is $N_{30,2}=464$ for $k=2$ and $N_{30,3}=4930$ for $k=3$.\sidenote{For Gaussian inputs the same counting works with Hermite polynomials in place of spherical harmonics, up to one radial direction (\cref{ex:lazy:gauss}).}

A network whose first layer moves can first find $w$ and then fit a function of one variable. \Cref{ch:singleindex} shows that one neuron trained by online stochastic gradient descent finds the direction of a target of information exponent $k$ from about $d^{k-1}$ samples (up to logarithms), and \textcite{DamianEtAl2022} prove that gradient descent on a two-layer network outside the kernel regime learns functions of $r$ hidden directions of degree $p$ from about $d^2r+dr^p$ samples, against $d^p$ for kernels. The kernel pays for not knowing the direction: it must be ready for all $N_{d,k}$ of them.

\paragraph{The kernel of a deep network.} The limit kernel of a deep fully connected network obeys the same Gaussian recursion as its forward signals (\cref{ch:signalprop}), with one more term.

\begin{theorem}[\textcite{JacotEtAl2018}; nonasymptotic form by \textcite{AroraEtAl2019a}]\label{thm:lazy:deepntk}
Consider $h^1(x)=\sigma_wd^{-1/2}W^1x+\sigma_bb^1$ and $h^{\ell+1}=\sigma_wm^{-1/2}W^{\ell+1}\phi(h^\ell)+\sigma_bb^{\ell+1}$ for $\ell<L$, with independent standard Gaussian entries in all $W^\ell$ and $b^\ell$, hidden widths $m$, one output, and $\phi$ Lipschitz and twice differentiable with bounded second derivative. Define $\Sigma^1(x,x')=\sigma_w^2\ip{x,x'}/d+\sigma_b^2$ and, with $(g,g')$ centered Gaussian with covariance $\Sigma^\ell$ restricted to $\{x,x'\}$,
\begin{gather*}
  \Sigma^{\ell+1}=\sigma_w^2\E[\phi(g)\phi(g')]+\sigma_b^2,\qquad \dot\Sigma^{\ell+1}=\sigma_w^2\E[\phi'(g)\phi'(g')],\\
  \Theta^{\ell+1}=\Theta^\ell\,\dot\Sigma^{\ell+1}+\Sigma^{\ell+1},
\end{gather*}
with $\Theta^1=\Sigma^1$. As the hidden widths tend to infinity one after the other, the tangent kernel of the output with respect to all weights and biases converges in probability to $\Theta^L$ at initialization, and, along gradient flow on the square loss over a fixed horizon, stays equal to $\Theta^L$ in the limit.
\end{theorem}

\begin{proofsketch}
By the chain rule the tangent kernel of $h^{\ell+1}$ is the contribution $\Sigma^{\ell+1}$ of the weights and biases of layer $\ell+1$, which converges by the law of large numbers as in \cref{ch:signalprop}, plus the tangent kernel $\Theta^\ell$ of each unit of layer $\ell$ (by induction) multiplied by $\frac{\sigma_w^2}m\sum_j\phi'(h^\ell_j(x))\phi'(h^\ell_j(x'))\to\dot\Sigma^{\ell+1}$; the cross terms average out. During training the weights of each layer move by $O(m^{-1/2})$ in operator norm, as in \cref{thm:lazy:lazy}. The details are in \textcite[Theorems~1 and~2]{JacotEtAl2018}; \textcite{AroraEtAl2019a} give explicit widths.
\end{proofsketch}

For the two-layer network of \cref{ass:lazy:twolayer} ($L=2$, $\sigma_w=1$, $\sigma_b=0$, unit inputs scaled by $\sqrt d$) the part of $\Theta^2$ coming from the first layer is $\Sigma^1\dot\Sigma^2=\ip{x,x'}\kappa(\ip{x,x'})$, which is \eqref{eq:lazy:kinf}; training the second layer adds $\Sigma^2$ (\cref{ex:lazy:bothlayers}).

\paragraph{The other infinite width.} Laziness comes from the factor $m^{-1/2}$, not from the width itself \parencite{ChizatEtAl2019}. The following derivation shows how.

\begin{heuristic}[Lazy and mean-field scalings]
Write $f(x)=m^{-\gamma}\sum_ra_r\phi(\ip{w_r,x})$ with $\frac12\le\gamma\le1$, and train with the step (or time scale) $\eta$. The function moves at speed $\partial_tf=-\eta\sum_r\nabla_{w_r}f\cdot\nabla_{w_r}\Lhat$, a sum of $m$ terms each of order $m^{-2\gamma}$, so it changes by order one per unit time only if $\eta\asymp m^{2\gamma-1}$. Each neuron then moves at speed $\eta\norm{\nabla_{w_r}\Lhat}\asymp m^{2\gamma-1}m^{-\gamma}=m^{\gamma-1}$. For $\gamma=\frac12$, the NTK parameterization with step of order one, each neuron moves by $O(m^{-1/2})$ while the function changes by $O(1)$: lazy. For $\gamma=1$, the \term{mean-field parameterization}\index{mean-field limit} with step of order $m$, each neuron moves by $O(1)$, and the kernel $\eta\sum_r\nabla_{w_r}f(x)\cdot\nabla_{w_r}f(x')$, an average of $m$ terms of order one, changes by order one: the features are learned. Between the two, the motion of one neuron is $m^{\gamma-1}$ and laziness sets in as $m\to\infty$ for every $\gamma<1$.
\end{heuristic}

The mean-field limit is not a linear model but a nonlinear partial differential equation for the distribution of neurons. Let $\rho^m_t=\frac1m\sum_r\delta_{(a_r(t),w_r(t))}$ for a network $f(x)=\frac1m\sum_ra_r\phi(\ip{w_r,x})$ trained in both layers, and write $f_\rho(x)=\int a\phi(\ip{w,x})\dd\rho(a,w)$.

\begin{theorem}[\textcite{MeiEtAl2018,ChizatBach2018,RotskoffVandenEijnden2022,SirignanoSpiliopoulos2020}; stated without proof]\label{thm:lazy:meanfield}
Under regularity conditions on $\phi$ and on the data, and with initial neurons independent with law $\rho_0$, the empirical measure $\rho^m_t$ of gradient flow on the population loss $\frac12\E(f(x)-f_\star(x))^2$, run with time scale $m$, converges weakly as $m\to\infty$, uniformly on bounded time intervals, to the solution of
\[
  \partial_t\rho_t=\nabla\cdot\bigl(\rho_t\nabla\Psi(\cdot;\rho_t)\bigr),
\]
with $\Psi(a,w;\rho)=\E_x[(f_\rho(x)-f_\star(x))\,a\,\phi(\ip{w,x})]$: a Wasserstein gradient flow of the loss as a function of $\rho$.
\end{theorem}

\Textcite{ChizatBach2018} prove that, under homogeneity conditions and for an initialization of full support, if the flow converges, then its limit is a global minimizer. \Textcite{YangHu2021} classify the scalings of deep networks: each infinite-width limit either learns features or follows kernel gradient descent, and in their maximal-update parameterization ($\mu$P) every layer learns features.

\section{Reading the solution}\label{sec:lazy:reading}

\paragraph{Modes.} \Cref{fig:lazy:modes} tests the linearized solution where it can be solved exactly. With equally spaced inputs on the circle the infinite-width kernel is circulant, its eigenvectors on the training set are the Fourier modes, and the linearized model predicts that the Fourier coefficient of frequency $k$ of the residual is multiplied by $1-\eta\lambda_k$ at each step of gradient descent. The eigenvalues $\lambda_1=3.62$, $\lambda_2=0.317$, and $\lambda_3=0.086$ of $K^\infty$ decrease by factors of $11$ and $4$, and the three frequencies are learned one after the other on time scales separated by the same factors. A network of width $16384$ follows the three predicted curves until each coefficient has decayed to about $10^{-2}$; below that, the residual leaks between modes, because the kernel of a finite network is not exactly circulant and drifts slightly during training. A network of width $64$ learns all three frequencies more slowly than predicted, and its first coefficient stalls near $2\times10^{-2}$: at that width the kernel at initialization already differs from $K^\infty$ by about $15\%$ in operator norm.

\begin{figure}[t]
\centering
\begin{tikzpicture}
\begin{axis}[hplot,height=0.44\linewidth,xlabel={training time $t$},ylabel={$\ip{\cos(k\vartheta),u_t}/\ip{\cos(k\vartheta),u_0}$},
  xmode=log,xmin=0.02,xmax=30,ymode=log,ymin=1e-3,ymax=1.5,unbounded coords=discard,legend pos=south west,legend columns=1]
\addplot[theory] table[x=t,y=k1] {figures/data/06-lazy-modes-theory.dat};
\addlegendentry{$k=1$}
\addplot[theory2] table[x=t,y=k2] {figures/data/06-lazy-modes-theory.dat};
\addlegendentry{$k=2$}
\addplot[theory3] table[x=t,y=k3] {figures/data/06-lazy-modes-theory.dat};
\addlegendentry{$k=3$}
\addlegendimage{simulation}
\addlegendentry{$m=16384$}
\addplot[simulation,forget plot] table[x=t,y=m16384_k1] {figures/data/06-lazy-modes-sim.dat};
\addplot[simulation,forget plot] table[x=t,y=m16384_k2] {figures/data/06-lazy-modes-sim.dat};
\addplot[simulation,forget plot] table[x=t,y=m16384_k3] {figures/data/06-lazy-modes-sim.dat};
\addplot[simulation2,mark size=1.5pt,forget plot] table[x=t,y=m64_k1] {figures/data/06-lazy-modes-sim.dat};
\addplot[simulation2,mark size=1.5pt,forget plot] table[x=t,y=m64_k2] {figures/data/06-lazy-modes-sim.dat};
\addplot[simulation2,mark size=1.5pt,forget plot] table[x=t,y=m64_k3] {figures/data/06-lazy-modes-sim.dat};
\addlegendimage{simulation2,mark size=1.5pt}
\addlegendentry{$m=64$}
\end{axis}
\end{tikzpicture}
\caption{Learning by modes. Training set: $n=32$ equally spaced angles $\vartheta_i$ on the circle, inputs $(\cos\vartheta_i,\sin\vartheta_i,1)/\sqrt2$, labels $y_i=\cos\vartheta_i+\cos2\vartheta_i+\cos3\vartheta_i$. Network: \cref{ass:lazy:twolayer} with $\phi=\tanh$ and an antisymmetric initialization ($w_{r+m/2}=w_r$, $a_{r+m/2}=-a_r$), so that $f(\cdot;W_0)=0$ and $u_0=-y$; full-batch gradient descent with step $\eta=0.02$ (seed $6$). Curves: the linearized prediction $(1-\eta\lambda_k)^{t/\eta}$ with the eigenvalues of $K^\infty$, computed by Gauss--Hermite quadrature. Circles: width $m=16384$; squares: $m=64$. Ratios below $10^{-3}$ are not shown.}
\label{fig:lazy:modes}
\end{figure}

\paragraph{How far things move.} \Cref{fig:lazy:width} measures the motion of the parameters and of the kernel. \Cref{prop:lazy:kernelreg} predicts the displacement exactly in the limit: after $s$ steps the linearized model has moved by $D^2=u_0^\top(K^\infty)^{-1}(I-(I-\eta K^\infty)^s)^2u_0$, (here $D=7.40$) whatever the width, while $\E\norm{W_0}_F^2=3m$. The predicted relative change $D/\sqrt{3m}$ matches the simulation to within $2\%$ from $m=512$ on. The largest move of a single neuron also decreases like $m^{-1/2}$ (from $0.45$ at $m=1024$ to $0.11$ at $m=16384$). The kernel, however, moves by order $1/m$, not $m^{-1/2}$: its relative change falls by a factor of about $2$ for each doubling of the width. \Cref{thm:lazy:twolayer} bounds it by $m^{-1/2}$, but the leading-order change cancels.

\begin{heuristic}[Why the kernel moves by $1/m$]
Let $\chi(w)=\phi'(\ip{w,x})\phi'(\ip{w,x'})$ and expand the kernel to first order in the displacements:
\[
  K_t(x,x')-K_0(x,x')\approx\frac{\ip{x,x'}}m\sum_r\nabla\chi(w_r)\cdot(w_r(t)-w_r(0)).
\] The displacement of neuron $r$ is $-m^{-1/2}a_r\int_0^t\sum_iu_{s,i}\phi'(\ip{w_r(s),x_i})x_i\dd s$, a sign $a_r$ times a vector $m^{-1/2}g_r$ that depends on $w_r$ and on the common residual. The first-order change is therefore $m^{-3/2}\sum_ra_rh_r$ with $h_r$ of order one. The signs $a_r$ are independent of $w_r$ and enter the residual only through averages over all neurons, so the sum behaves like a sum of $m$ terms of random sign, of order $\sqrt m$, and the first-order change is of order $m^{-1}$. The second-order terms are $\frac1m\sum_rO(\norm{w_r(t)-w_r(0)}^2)=O(1/m)$. For the antisymmetric initialization of the figure, the first-order terms of each pair of neurons cancel to leading order.
\end{heuristic}

The $1/m$ rate agrees with the analysis of \textcite{DyerGurAri2020}, in which the tangent kernel is constant during training up to corrections of order $1/m$.\sidenote{\Textcite{DyerGurAri2020} derive this from a conjecture on the scaling of correlation functions with width, which they analyze completely for deep linear networks and support with evidence for nonlinear ones.} We checked that the same rate holds with independent random signs (\cref{ex:lazy:simwidth}).

\begin{figure}[t]
\centering
\begin{tikzpicture}
\begin{axis}[hplot,height=0.4\linewidth,xmode=log,ymode=log,xlabel={width $m$},ylabel={relative change},
  xmin=20,xmax=26000,ymin=2e-5,ymax=1.5,log basis x=2,legend pos=south west,legend style={font=\scriptsize},
  xtick={32,128,512,2048,8192},xticklabels={$2^5$,$2^7$,$2^9$,$2^{11}$,$2^{13}$}]
\addplot[theory] table[x=m,y=relWtheory] {figures/data/06-lazy-width.dat};
\addlegendentry{$D/\sqrt{3m}$, linearized}
\addplot[simulation] table[x=m,y=relW] {figures/data/06-lazy-width.dat};
\addlegendentry{$\norm{W_T-W_0}_F/\norm{W_0}_F$}
\addplot[simulation2] table[x=m,y=relK] {figures/data/06-lazy-width.dat};
\addlegendentry{$\norm{K_T-K_0}_{\op}/\norm{K_0}_{\op}$}
\addplot[guide,domain=20:26000,samples=2] {1.82/sqrt(x)};
\addlegendentry{slope $-\frac12$}
\addplot[muted,line width=0.5pt,dotted,domain=20:26000,samples=2] {9.67/x};
\addlegendentry{slope $-1$}
\end{axis}
\end{tikzpicture}
\caption{Motion of the weights and of the tangent kernel during training, against the width. Same training set, network, and initialization as \cref{fig:lazy:modes}, labels $y_i=\cos\vartheta_i+\cos2\vartheta_i$, gradient descent with $\eta=0.02$ for time $T=25$, at which the training residual is between $0.2\%$ and $2.3\%$ of $\norm y$. Points: means over $4$ seeds. Curve: the exact linearized displacement $D=7.40$ of \cref{prop:lazy:kernelreg} divided by $\sqrt{3m}$. Dashed and dotted lines: reference slopes through the kernel change at $m=32$ and $m=1024$.}
\label{fig:lazy:width}
\end{figure}

\paragraph{The price.} \Cref{fig:lazy:gap} compares kernel ridge regression with a trained network on the target $f_w$ of \cref{thm:lazy:lower} with $k=2$ in dimension $d=30$, where $N_{30,2}=464$. Kernel ridge regression with a two-layer ReLU tangent kernel, given the best ridge parameter in hindsight, learns nothing below $n=100$, has error $0.6$ at $n=400$, and needs about $n=1000$ samples to reach error $0.2$; the theorem's bound $1-n/464$ lies below it, as it must. A two-layer network with $256$ units, trained in both layers from a small output scale,\sidenote{No theorem of this chapter covers this network, trained by Adam with weight decay; it is there to show what moving features buy.} reaches error $0.2$ at $n=240$ and $0.1$ at $n=480$. At this dimension the gap is a factor of four to five in samples, and it grows like $d$: by the theorem the kernel needs about $d^2/2$ samples, while a network that finds the direction needs order $d$ up to logarithms (\cref{ch:singleindex}). The kernel curve also shows what the bound does not: $n\ge N$ samples are not enough, and even at $n=600$ the error is $0.4$; \textcite{GhorbaniEtAl2021} prove that kernel ridge regression learns the degree-$2$ part only once $n\gg d^2$.

\begin{figure}[t]
\centering
\begin{tikzpicture}
\begin{axis}[hplot,height=0.4\linewidth,xmode=log,xlabel={number of samples $n$},ylabel={relative test error},
  xmin=20,xmax=2200,ymin=0,ymax=1.4,legend pos=north east]
\addplot[theory,domain=20:464,samples=60] {1-x/464};
\addlegendentry{bound $1-n/N_{30,2}$}
\addplot[simulation2] table[x=n,y=krr] {figures/data/06-lazy-gap-kernel.dat};
\addlegendentry{kernel ridge regression}
\addplot[simulation] table[x=n,y=net] {figures/data/06-lazy-gap-net.dat};
\addlegendentry{trained network}
\addplot[guide] coordinates {(464,0) (464,1.4)};
\end{axis}
\end{tikzpicture}
\caption{A kernel against a network that learns its features, on the target $f_w=\sqrt N\,P_{d,2}(\ip{w,x})$ with $x\sim\Unif(\Sph^{29})$ and no noise; dashed line at $n=N_{30,2}=464$. Squares: kernel ridge regression with the two-layer ReLU tangent kernel with a bias coordinate (both layers), ridge parameter chosen from $\{10^{-8},\dots,10\}$ by test error. Circles: a network $m^{-1/2}\sum_ra_r\relu(\sqrt d\ip{w_r,x}+b_r)$ with $m=256$ trained in all parameters by full-batch Adam (step $0.02$, $2000$ steps, weight decay $10^{-4}$), from $a_r\sim\cN(0,10^{-2})$. Relative errors $\norm{\hat f-f_w}^2/\norm{f_w}^2$ on $4000$ test points, averaged over $3$ draws of $w$ and of the data.}
\label{fig:lazy:gap}
\end{figure}

\section{What the model gets wrong}\label{sec:lazy:wrong}

\paragraph{What was thrown away.} The features: in the lazy regime the representation is the random one at initialization, and everything learned is a linear combination of fixed tangent features. Width that is large compared with the data: \cref{thm:lazy:twolayer} needs $m$ of order $n^2/\lmin^2$, and $\lmin$ shrinks with $n$ and is tiny for smooth activations (\cref{sec:lazy:convergence}), so the theorem covers widths that no practical network has. Finite step sizes: gradient descent with a large step leaves the linear regime even at large width, and its catapult dynamics are the subject of \cref{ch:stability}. Depth and architecture enter only through the kernel, and the interaction of layers, in which one layer learns features that the next one uses, is absent. And the initialization is data independent, so the kernel knows nothing about the task.

\paragraph{Where it breaks.} \emph{Performance.} The CNTK's $77\%$ on CIFAR-10 \parencite{AroraEtAl2019a} is well below what the same kind of architecture reaches when trained in the usual way. In the comparisons of \textcite{LeeEtAl2020} infinite-width kernels beat finite fully connected networks but lose to finite convolutional networks trained with large steps and weight decay, the very ingredients that break the correspondence. Deep convolutional networks forced into the lazy regime by a large output scale test worse \parencite{ChizatEtAl2019}, and \textcite{GeigerEtAl2020} found feature learning better for convolutional networks, though not for fully connected ones.

\emph{The kernel moves.} In networks trained in the usual way the tangent kernel is not constant. \Textcite{FortEtAl2020} found that it changes rapidly during the first few epochs and then slowly, and that kernel regression with the kernel reached partway through training comes close to the performance of the fully trained network, which the initial kernel does not. \Textcite{AtanasovEtAl2022} observed, in networks trained from small initializations, a \emph{silent alignment}: early in training, while it is still small, the kernel rotates toward the task, and the network then behaves like kernel regression with the final kernel; the network is a kernel machine, but with a learned kernel. \Textcite{HaninNica2020} showed that the tangent kernel of a deep ReLU network at initialization is not deterministic when depth and width grow together, its fluctuations being exponential in depth over width, and that its first gradient step moves it by an amount set by the same ratio, so a wide network can learn features if it is deep enough.

\emph{Sample complexity.} \Cref{thm:lazy:lower} is the cleanest instance of a family of separations. \Textcite{GhorbaniEtAl2021} show that kernel ridge regression with a rotation-invariant kernel learns, in the limit of infinite width, only the projection of the target on polynomials of degree $\ell$ when $d^\ell\ll n\ll d^{\ell+1}$, and that the linearized two-layer models with finite width are limited by the number of their features in the same way; random features cannot learn even a single ReLU neuron with Gaussian inputs unless their number or their weights are exponential in $d$ \parencite{YehudaiShamir2019}; and three-layer residual networks learn some functions efficiently through hierarchical feature learning that kernels cannot \parencite{AllenZhuLi2019}.

\paragraph{Fine structure.} \emph{Corrections in $1/m$.} \Textcite{DyerGurAri2020} computed the corrections to the infinite-width dynamics order by order in $1/m$ with Feynman diagrams, \textcite{HuangYau2020} wrote the evolution of the kernel as an infinite hierarchy of equations whose truncations approximate it, and \textcite{BaiLee2020} replaced the linearization by a quadratic model that can escape the kernel regime and gains a factor of $d$ in sample complexity. \emph{The other limits.} The mean-field limit of \cref{thm:lazy:meanfield} keeps feature learning at infinite width, and $\mu$P extends it to deep networks \parencite{YangHu2021}. \emph{Crossovers} are the subject of \cref{ex:lazy:crossover}. \emph{One step of feature learning.} \Textcite{BaEtAl2022} and \textcite{DamianEtAl2022} analyze networks whose first layer takes one or a few large gradient steps, enough to align the features with the target's directions and to beat every kernel. Open: a quantitative theory, for a realistic deep architecture and data, of how much the tangent kernel moves during training, in which directions, and how much the motion improves the test error.

\section{Model card}\label{sec:lazy:card}

\modelcard{Very wide networks fit their training data while each weight barely moves; their predictions track kernel regression with the neural tangent kernel; low frequencies are learned first.}
{A smooth map $F$ from parameters to outputs, trained by gradient flow on the square loss; specialized to a two-layer network $m^{-1/2}\sum_ra_r\phi(\ip{w_r,x})$ with Gaussian initialization.}
{Convergence inside a ball of radius $\norm{u_0}/\mu$ when $L_J\norm{u_0}\le\mu^2$ (\cref{thm:lazy:lazy}); closeness to the linearization (\cref{cor:lazy:linear}); width $n^2/\lmin^2$ suffices for two-layer networks (\cref{thm:lazy:twolayer}); kernel lower bound $1-n/N_{d,k}$ (\cref{thm:lazy:lower}).}
{Exponential decay of each kernel mode at rate $\lambda_k$; relative weight change of order $m^{-1/2}$, with the exact displacement $\sqrt{u_0^\top K^{-1}u_0}$; kernel methods need about $d^k/k!$ samples for a degree-$k$ function of one direction.}
{The features never move; the widths covered are unrealistic; kernels underperform trained networks; the kernel of practical networks changes early and aligns with the task.}
{$1/m$ corrections and the tangent hierarchy; quadratic models; mean-field limits and $\mu$P; one-step feature learning; lazy-to-rich crossovers.}

The chapter is the lazy end of a thread that runs through the book. Large initializations made deep linear networks learn all modes at once instead of in stages (\cref{ch:deeplinear}) and made diagonal networks find the minimum-$\ell_2$ interpolator instead of a sparse one (\cref{ch:diagonal}); grokking is read in \cref{ch:grokking} as a transition from an early lazy phase that memorizes to a late rich phase that generalizes; and the single-index model of \cref{ch:singleindex} shows what the network gains by moving its features: $d^{k-1}$ samples instead of the kernel's $d^k$. Spectra are a second thread: the time scale $1/\lambda_k$ of each mode here and the power-law spectra that turn kernel regression into scaling laws in \cref{ch:scaling} say that a linear model learns its target eigenvalue by eigenvalue. On the Boolean cube the same counting argument, with characters instead of spherical harmonics, shows why kernels cannot learn sparse parities (\cref{ch:parity}).

\begin{exercises}
\item \consolidation (Gradient descent on the linearized model.) Run gradient descent $\theta_{s+1}=\theta_s-\eta J_0^\top(F^{\mathrm{lin}}(\theta_s)-y)$ with $K_0=J_0J_0^\top$ invertible. (a) Show that $u_s=(I_n-\eta K_0)^su_0$ and that $u_s\to0$ for every $u_0$ if and only if $0<\eta<2/\lmax(K_0)$. (b) Show that $\norm{\theta_s-\theta_0}^2=u_0^\top K_0^{-1}(I_n-(I_n-\eta K_0)^s)^2u_0$, and that for $\eta\le1/\lmax(K_0)$ it increases with $s$ to $u_0^\top K_0^{-1}u_0$.\label{ex:lazy:gd}

\item \consolidation (Arc-cosine kernels and odd harmonics.) Let $g,g'$ be standard Gaussians with correlation $\rho$. (a) Show that $\Pr(g>0,g'>0)=(\pi-\arccos\rho)/(2\pi)$ and $\E[\relu(g)\relu(g')]=(\sqrt{1-\rho^2}+(\pi-\arccos\rho)\rho)/(2\pi)$. (b) For $\phi=\relu$ in \eqref{eq:lazy:kinf}, show that $K^\infty(x,x')=\frac\rho4+\frac{\rho\arcsin\rho}{2\pi}$ with $\rho=\ip{x,x'}$, and deduce from the parity of $P_{d,k}$ and the orthogonality of the Gegenbauer polynomials that $\lambda_k=0$ for every odd $k\ge3$. (c) On $\Sph^2$, where $P_{3,k}$ is the Legendre polynomial and the weight in \cref{prop:lazy:spectral} is $\frac12$, show that $\lambda_0=\frac1{16}$ and $\lambda_1=\frac1{12}$, and check numerically that $\lambda_2=\frac7{256}$. \hint{in (a), the event $\{g>0,g'>0\}$ is a wedge of angle $\pi-\arccos\rho$ in the plane; in (c), $\int_0^1t\arcsin t\dd t=\frac\pi8$.}\label{ex:lazy:relu}

\item \consolidation (Parallel inputs.) (a) For $\phi=\tanh$ and inputs $x_1$ and $x_2=-x_1$, show that the two rows of $K^\infty$ are proportional, so $\lmin(K^\infty)=0$. (b) Show that appending a bias coordinate, $\tilde x=(x,1)/\sqrt2$, makes distinct unit inputs pairwise non-parallel, so that \cref{lem:lazy:pd} applies to every set of distinct inputs.\label{ex:lazy:pd}

\item \consolidation (Counting dimensions.) (a) Check that $N_{30,2}=464$ and $N_{30,3}=4930$. (b) A two-layer network of width $m$ on inputs in $\R^d$ has $md$ tangent features. Show that its linearization, trained on $n$ examples from an initialization independent of $w$, has expected error at least $1-(\min\{n,md\}+1)/N_{d,k}$ on the target of \cref{thm:lazy:lower}.\label{ex:lazy:count}

\item \variation\hard (Training both layers.) In \cref{ass:lazy:twolayer} let the $a_r$ be trained too, starting from $\pm1$. (a) Show that the tangent kernel becomes $K_W+K_a$ with $K_a(x,x')=\frac1m\sum_r\phi(\ip{w_r,x})\phi(\ip{w_r,x'})$, whose expectation is the Gaussian-process kernel $\E[\phi(g)\phi(g')]$, so that $\lmin$ can only increase. (b) Show that on the ball of radius $R$ around the initialization, $\norm{J(\theta)-J(\theta')}_{\op}\le\beta\sqrt{n/m}\,\sqrt{1+2(1+R)^2}\,\norm{\theta-\theta'}$. (c) Deduce a version of \cref{thm:lazy:twolayer} with $m$ of order $n^3/(\lmin^3\delta^2)$, up to a factor depending on $\beta$ and $c_\phi$. Where does the extra factor $n/(\lmin\delta)$ come from, and how would a bound on the motion of each $a_r$ remove it? \hint{in (b), $\abs{a_r}\le1+R$ on the ball, and $a\phi'(s)-a'\phi'(s')=(a-a')\phi'(s)+a'(\phi'(s)-\phi'(s'))$.}\label{ex:lazy:bothlayers}

\item \variation (Laziness is a matter of scale.) Let $F$ satisfy hypotheses (i) and (ii) of \cref{thm:lazy:lazy} with constants $\mu$ and $L_J$ on every ball around $\theta_0$, with $F(\theta_0)=0$, and train $\alpha F$ on $\frac12\norm{\alpha F(\theta)-y}^2$ for a scale $\alpha>0$. (a) Show that the theorem applies to $\alpha F$ as soon as $\alpha\ge L_J\norm y/\mu^2$. (b) Show that the parameters then move by at most $\norm y/(\alpha\mu)$, which tends to $0$ as $\alpha\to\infty$, and that the linearization error of \cref{cor:lazy:linear} is of order $1/\alpha$. (c) Compare with the initialization scale $\alpha$ of diagonal linear networks in \cref{ch:diagonal}, where $F(\theta_0)\ne0$ is avoided by the parameterization $u\odot u-v\odot v$.\label{ex:lazy:scale}

\item \variation\hard (Gaussian inputs.) Let $x\sim\cN(0,I_d)$, let $w\sim\Unif(\Sph^{d-1})$, and let $f_w(x)=h_2(\ip{w,x})$. (a) With $H_\alpha=\prod_ih_{\alpha_i}(x_i)$, show that
\[
  h_2(\ip{w,x})=\sum_{\abs\alpha=2}\sqrt{2/\alpha!}\,w^\alpha H_\alpha(x).
\] (b) For $g$ in the span of the $H_\alpha$ with $\abs\alpha=2$, write $\ip{g,f_w}=w^\top Aw$ with $A$ symmetric and show that $\norm g=\norm A_F$. (c) Using $\E_w(w^\top Aw)^2=\frac{2\tr A^2+(\tr A)^2}{d(d+2)}$, show that the operator $g\mapsto\E_w[\ip{g,f_w}f_w]$ has the eigenvalue $\frac1d$ once (on $g\propto\norm x^2-d$) and $\frac2{d(d+2)}$ on the rest of its range. (d) Deduce that every $n$-dimensional $V$ independent of $w$ has $\E\norm{f_w-\Pi_Vf_w}^2\ge1-\frac1d-\frac{2(n-1)}{d(d+2)}$, and explain the term $\frac1d$, which the sphere does not have. \hint{for (d), the sum of $\ip{e_j,Ce_j}$ over an orthonormal family of size $n$ is at most the sum of the $n$ largest eigenvalues of $C$.}\label{ex:lazy:gauss}

\item \simulation (Weights and kernel against width.) Repeat the experiment of \cref{fig:lazy:width} with independent uniform signs $a_r$ and independent $w_r\sim\cN(0,I_3)$, for $m\in\{256,1024,4096\}$ and $3$ seeds. Observe a relative weight change of about $0.26$, $0.14$, and $0.07$, and a relative kernel change of about $0.04$, $0.011$, and $0.003$ (seeds $0$ to $2$; other seeds differ by up to $20\%$). (a) Which quantity follows $m^{-1/2}$ and which $m^{-1}$? (b) Explain both with \cref{prop:lazy:kernelreg} and the heuristic derivation of \cref{sec:lazy:reading}.\label{ex:lazy:simwidth}

\item \simulation (A kernel that cannot see odd harmonics.) On $\Sph^2$, fit kernel interpolation with $K(x,x')=\rho\,(\pi-\arccos\rho)/(2\pi)$, $\rho=\ip{x,x'}$ (a regularization of $10^{-10}$ suffices), to $n\in\{50,200,800,2000\}$ uniform samples of the targets $\sqrt5P_2(\ip{w,x})$ and $\sqrt7P_3(\ip{w,x})$, with $P_k$ the Legendre polynomials. (a) Observe that the relative test error for degree $2$ falls below $10^{-4}$ by $n=200$, and that for degree $3$ it stays above $1$ (between about $1.4$ and $2.2$ at every $n$, near $1.45$ at $n=2000$). (b) Explain both with \cref{ex:lazy:relu}. (c) Add a bias coordinate as in \cref{fig:lazy:modes} and observe that the degree-$3$ target is learned.\label{ex:lazy:simodd}

\item \research (The lazy-to-rich crossover.) Train the network of \cref{fig:lazy:gap} with output $\alpha(f(x;\theta)-f(x;\theta_0))$ and loss $\frac1{\alpha^2}$ times the square loss, for a range of $\alpha$ and $m$. For large $\alpha$ it approaches kernel regression with its initial tangent kernel; for small $\alpha$ it learns $w$. \Textcite{GeigerEtAl2020} map such crossovers empirically, and \textcite{WoodworthEtAl2020} solve one for diagonal linear networks. Open: for a nonlinear network on high-dimensional data, predict the scale $\alpha_c(m,n,d)$ at which the test error leaves the kernel curve, and the shape of the transition.\label{ex:lazy:crossover}

\end{exercises}

\notesandrefs
The neural tangent kernel was introduced by \textcite{JacotEtAl2018}, first posted in June 2018; the plan of \cref{thm:lazy:twolayer} follows \textcite{DuEtAl2019}, and \textcite{AroraEtAl2019b} used the kernel to explain why random labels are fitted more slowly than true ones and to bound the generalization error. The abstract form of \cref{thm:lazy:lazy}, with a condition on the Jacobian at initialization and its Lipschitz constant, follows \textcite{OymakSoltanolkotabi2020} and \textcite{ChizatEtAl2019}.

\Cref{thm:lazy:lower} is a simple instance of the lower bounds of \textcite{GhorbaniEtAl2021}, who also prove matching upper bounds; our proof uses only the addition theorem. The mean-field limit was found independently by \textcite{MeiEtAl2018}, \textcite{ChizatBach2018}, \textcite{RotskoffVandenEijnden2022}, and \textcite{SirignanoSpiliopoulos2020}, all first posted in 2018.

\paragraph{Translating the literature.} \Textcite{JacotEtAl2018} write $\Theta$ for our $K^\infty$ and $\Theta^{(L)}$ for our $\Theta^L$, $n_\ell$ for the widths, $\beta$ for our $\sigma_b$, and $\Sigma^{(\ell)}$ and $\dot\Sigma^{(\ell)}$ as we do; they train with the loss normalized by the number of examples, which rescales time by $n$. \Textcite{DuEtAl2019} write $u(t)$ for our $F(\theta_t)$ (so $y-u(t)$ is our $-u_t$), $\mathbf H(t)$, $\mathbf H^\infty$, and $\lambda_0$ for our $K_t$, $K^\infty$, and $\lmin$, and $\sigma$ for $\phi$. \Textcite{ChizatEtAl2019} write $h$ for our $F$, $Dh(w_0)$ for $J(\theta_0)$, $\alpha$ for the output scale of \cref{ex:lazy:scale}, and measure laziness by $\kappa_h=\norm{h(w_0)-y^\star}\,\norm{D^2h(w_0)}/\norm{Dh(w_0)}^2$, the analogue of our $L_J\norm{u_0}/\mu^2$.

\printchapterbib
\end{refsection}


\casepart{Feature Learning in High Dimensions}{Real networks find the few directions that matter. In high dimension that search is the hard part, and its cost is set by how much signal the target leaks at low order. Plateaus, sudden drops and hidden progress are three faces of the same search.}{46}
% Chapter 7 of The Hydrogen Atoms of Deep Learning: single-index models and the information exponent.
\casechapter{The Search Phase}{Single-index models and the information exponent}{Single-index model, online SGD}{ch:singleindex}
\begin{refsection}[chapters/07-singleindex.bib]

\begin{chaptersummary}
When a network must find a few relevant directions in a high-dimensional input, its loss often stays flat for a long time and then drops. The single-index model, a target $g(\ip{w_\star,x})$ of a Gaussian input learned by one neuron with online stochastic gradient descent on the sphere, isolates this search. Its state is one number, the overlap $r=\ip{w,w_\star}$, which leaves the equator $r=0$ at a rate proportional to $r^{k-1}$, where $k$ is the information exponent, the lowest Hermite degree that the target shares with the neuron. We prove that the population dynamics escapes in time of order $d^{(k-2)/2}$ and that online SGD, whose step size the noise at the equator limits to order $d^{-k/2}$, finds the direction with of order $d^{k-1}$ samples when $k\ge3$, with explicit constants; a theorem we quote shows that this is sharp up to logarithms. The model is wrong about the task itself: reusing samples, smoothing the loss, or transforming the labels finds the direction with far fewer samples, so the information exponent measures one algorithm, not the difficulty of the problem.
\end{chaptersummary}

\section{The phenomenon: plateaus whose length is set by the target}\label{sec:singleindex:phenomenon}

A trained network computes features, functions of the input that make the task easy, and it has to find them. When the task depends on a few directions of a high-dimensional input, finding them is a search over a sphere whose dimension is that of the input, and in the training curve the search shows up as a plateau: the loss stays near its initial value for a long time and then falls in a comparatively short one. Plateaus of this kind are among the oldest observations about gradient training. Committee machines trained online sit on plateaus while their hidden units are interchangeable (\cref{ch:committee}), networks trained on sparse parities show no progress in the loss for a long time and then learn abruptly (\cref{ch:parity}), and transformers acquire induction heads in a sudden transition (\cref{ch:induction}). This chapter isolates the simplest setting in which the length of the plateau can be computed, and the answer depends on the target through one integer.

The integer comes from the Hermite expansion (\cref{lem:toolkit:hermite}).\sidenote{The polynomials carry the name of Charles Hermite, who studied them in 1864; Laplace and Chebyshev had used them earlier.} A target $y=g(\ip{w_\star,x})$ with $x\sim\cN(0,I_d)$ and $\E g(z)^2<\infty$ for $z\sim\cN(0,1)$ expands as $g=\sum_{j\ge0}\hat g_jh_j$, and the first index $j\ge1$ with $\hat g_j\ne0$ says how much of the target is visible to low-degree statistics of the input. The targets $\ip{w_\star,x}^3$ and $\He_3(\ip{w_\star,x})=\ip{w_\star,x}^3-3\ip{w_\star,x}$ differ by a linear function. The first is correlated with $\ip{w_\star,x}$, so the empirical average of $yx$ already points at $w_\star$; the second is uncorrelated with every polynomial of degree at most $2$ in $x$, so no statistic of that kind sees it. Experiments show that this difference controls how long gradient training searches.

\begin{observation}[\textcite{BenArousGheissariJagannath2021}; \textcite{AbbeBoixAdseraMisiakiewicz2023}; \textcite{ArnaboldiEtAl2023}]\label[observation]{obs:singleindex:plateau}
Train a single neuron or a two-layer network by one-pass stochastic gradient descent on inputs with independent coordinates, Gaussian or uniform on $\{-1,1\}^d$, with a target that depends on one or a few directions. The loss stays near its initial value for a stretch of training and then falls quickly; most of the samples are spent in the flat stretch, and the stretch is longer when the lowest-degree component of the target has higher degree. For a two-layer network on $\{-1,1\}^{100}$ learning $z_1+z_1\cdots z_5+z_1\cdots z_9+z_1\cdots z_{14}$, the plateaus before the four drops last, by the authors' estimate, about $d$, $d^3$, $d^3$, and $d^4$ steps, where $1$, $4$, $4$, and $5$ new coordinates enter. For targets with many flat directions at initialization, such as phase retrieval, two-layer networks escape after a number of samples of order $d\log d$, and overparameterization shortens the escape only by a constant factor.
\end{observation}

These are controlled experiments on synthetic data, and they are robust: many groups have reproduced them, with different architectures, losses, and step sizes. Their limits are as clear. Most use one or two dimensions, so the growth of the plateau as a power of $d$ is read off from theory rather than measured. In practical networks plateaus are common, but none has been traced to the Hermite degree of a target. The model of this chapter supplies the missing dimension sweep (\cref{fig:singleindex:escape}). The question it answers: how long does the search last, as a function of the dimension and of the target, and why does the answer depend on the target only through one integer?

\section{The model: one neuron searching for one direction}\label{sec:singleindex:model}

Write $P_w^\perp=I_d-ww^\top$ for the orthogonal projection onto the tangent space of the sphere $\Sph^{d-1}$ at a unit vector $w$.

\begin{definition}[The single-index model and spherical online SGD]\label[definition]{def:singleindex:model}
Let $d\ge2$ and $w_\star\in\Sph^{d-1}$, and let $g$ and $\phi$ be polynomials. Inputs are $x\sim\cN(0,I_d)$ and labels are $y=g(\ip{w_\star,x})$. A neuron with weight $w$ predicts $\phi(\ip{w,x})$ and pays the \term{correlation loss}\index{correlation loss} $\ell(w;x,y)=-y\,\phi(\ip{w,x})$; its population loss is $\Loss(w)=\E\,\ell(w;x,y)$ for $w\in\R^d$. Given independent samples $(x_t,y_t)$ for $t=0,1,2,\dots$, a start $w_0\in\Sph^{d-1}$, and a step size $\eta>0$, \term{spherical online SGD}\index{online SGD} is the recursion
\begin{equation}\label{eq:singleindex:sgd}
\begin{aligned}
  w_{t+1}&=\frac{w_t-\eta G_t}{\norm{w_t-\eta G_t}},\\
  G_t&=P_{w_t}^\perp\nabla_w\ell(w_t;x_t,y_t)=-y_t\,\phi'(\ip{w_t,x_t})\,P_{w_t}^\perp x_t .
\end{aligned}
\end{equation}
The \term{overlap}\index{overlap} is $r_t=\ip{w_t,w_\star}$.
\end{definition}

Each step uses a new sample, so after $t$ steps the algorithm has seen $t$ samples: in this model the sample complexity and the running time are the same number. The update moves along the tangent vector $-G_t$ and returns to the sphere; because $G_t$ is orthogonal to $w_t$, the normalization divides by $\sqrt{1+\eta^2\norm{G_t}^2}\ge1$, never by zero. By \cref{lem:singleindex:hermite} below, the population loss on the sphere depends on $w$ only through the overlap, $\Loss(w)=-\Psi(\ip{w,w_\star})$, with
\begin{equation}\label{eq:singleindex:Psi}
  \Psi(r)=\sum_{j\ge0}\hat g_j\hat\phi_j\,r^j,\qquad -1\le r\le1 .
\end{equation}
For the matched pure link $g=\phi=h_k$ this is $\Psi(r)=r^k$. That case is the hydrogen atom of the chapter, and the reader may keep it in mind throughout.

\begin{definition}[Information exponent]\label[definition]{def:singleindex:ie}
The \term{information exponent}\index{information exponent} of a function $g$ with $\E g(z)^2<\infty$ is $k^\star(g)=\min\{j\ge1:\hat g_j\ne0\}$. For a target $g$ learned by a neuron $\phi$ on the correlation loss, the relevant integer is $k^\star=\min\{j\ge1:\hat g_j\hat\phi_j\ne0\}$, the lowest degree at which $\Psi$ is not constant. It is at least $k^\star(g)$, with equality when the Hermite coefficients of the student do not vanish.
\end{definition}

The name is due to \textcite{BenArousGheissariJagannath2021}, who define it through the derivatives of the population loss at the equator; \textcite{DudejaHsu2018} had already organized the learning of single-index models by the Hermite expansion of the link.\sidenote{Dudeja and Hsu designed a semiparametric estimator rather than analyzing SGD; their sample sizes are polynomial in $d$, with an exponent set by the smoothness of the link and the strength of its signal.} Linear functions, the ReLU, $\tanh$, the sign, and $z^3$ have information exponent $1$, because each is correlated with $z$; even functions such as $z^2$ and $\abs z$ have information exponent $2$; and $\He_k$ has information exponent $k$ (\cref{ex:singleindex:exponents}). Near the equator $r=0$ the population loss is $-\Psi(0)-\hat g_k\hat\phi_kr^k+O(r^{k+1})$ with $k=k^\star$, so for $k\ge2$ the equator is a critical set of the loss, flat to order $k$. A random start lies almost on it, because a uniformly random unit vector has overlap of order $d^{-1/2}$.\sidenote{Concentration of measure: all but a fraction $2e^{-d\eps^2/2}$ of the sphere lies within distance $\eps$ of each equator.} The search phase\index{search phase} is the escape from this degenerate saddle; for $k=2$ it is the logistic escape of \cref{ch:deeplinear}, and for larger $k$ it is the higher-order saddle that recurs in \cref{ch:committee,ch:parity}.

Each simplification removes one difficulty so that the search can be seen alone. Gaussian inputs make the joint law of $\ip{w_\star,x}$ and $\ip{w,x}$ depend only on $r$, so the overlap is the only order parameter, and the Hermite basis diagonalizes the correlation of the two (\cref{lem:singleindex:hermite}), as the characters diagonalize noise on the cube in \cref{ch:parity}. This rotational invariance is the thread that also runs through \cref{ch:doubledescent,ch:lazy,ch:committee,ch:icl}. One direction and a known link separate the search for $w_\star$ from the fitting of a function. The sphere removes the dynamics of the norm, which in a network is set by the other layers and by weight decay; the price is the normalization step, which pulls the overlap toward the equator and is one of the two effects that limit the step size. The correlation loss changes nothing essential: on the sphere $\frac12\E(y-\phi(\ip{w,x}))^2=\frac12\E g(z)^2+\frac12\E\phi(z)^2-\Psi(r)$, so the square loss has the same landscape, and its stochastic gradient differs by a term of mean zero (\cref{ex:singleindex:square}). Fresh samples make the gradient noise a martingale difference sequence and make samples and steps the same thing.

To invent the model one has to notice two facts. Rotational invariance collapses the $d$-dimensional dynamics to one number, so that a high-dimensional search becomes a one-dimensional random walk with drift. And near the equator only the lowest common Hermite degree matters: the drift of the overlap is of order $r^{k-1}$ while the noise is of order one, so the search is a contest between a signal that vanishes at the starting point and a noise that does not.

\begin{dictionary}
weight $w\in\Sph^{d-1}$ & incoming weights of one hidden neuron, normalized\\
direction $w_\star$ & a feature on which the target depends\\
overlap $r=\ip{w,w_\star}$ & alignment of a neuron with the feature: the progress of feature learning\\
equator $\abs r\lesssim d^{-1/2}$ & where every neuron starts under random initialization\\
Hermite coefficients $\hat g_j$ & how the target depends on the feature, degree by degree\\
information exponent $k^\star$ & how much of the target leaks to low-degree statistics\\
step size $\eta$ & learning rate\\
one fresh sample per step & one pass over a large data set\\
normalization after each step & weight normalization, weight decay, or the scale set by other layers\\
search phase, $r$ near $0$ & a loss plateau\\
descent phase, $r$ near $1$ & the drop that ends the plateau\\
\end{dictionary}

\section{The solution: the overlap and its drift}\label{sec:singleindex:drift}

Everything in this section rests on one identity for Gaussian vectors, which turns the $d$-dimensional problem into a problem about one number.

\begin{lemma}[Hermite correlation]\label[lemma]{lem:singleindex:hermite}
Let $U$ and $V$ be jointly Gaussian standard variables with $\E[UV]=\rho$. Then $\E[h_j(U)h_k(V)]=\delta_{jk}\rho^k$ for all $j,k\ge0$, and all functions $f_1$ and $f_2$ with $\E f_1(z)^2<\infty$ and $\E f_2(z)^2<\infty$ satisfy $\E[f_1(U)f_2(V)]=\sum_{k\ge0}\hat f_{1,k}\hat f_{2,k}\rho^k$, with $\hat f_{i,k}=\E[f_i(z)h_k(z)]$. In particular $\Loss(w)=-\Psi(\ip{w,w_\star})$ for every $w\in\Sph^{d-1}$.
\end{lemma}

\begin{proof}
The first statement is \cref{thm:toolkit:hermite-corr}; we recall the proof because it takes three lines. The generating function $e^{sz-s^2/2}=\sum_{k\ge0}s^kh_k(z)/\sqrt{k!}$ (\cref{lem:toolkit:hermite}) and the Gaussian moment generating function $\E e^{sU+tV}=e^{(s^2+2\rho st+t^2)/2}$ give, for real $s$ and $t$,
\[
  \sum_{j,k\ge0}\frac{s^jt^k}{\sqrt{j!\,k!}}\,\E[h_j(U)h_k(V)]=\E\bigl[e^{sU-s^2/2}e^{tV-t^2/2}\bigr]=e^{\rho st}=\sum_{k\ge0}\frac{(\rho st)^k}{k!} .
\]
Taking the expectation term by term is legitimate because $\E\abs{h_j(U)h_k(V)}\le1$ by the Cauchy--Schwarz inequality, so the double series converges absolutely; comparing the coefficients of $s^jt^k$ proves the first statement. For the second, expand $f_1$ and $f_2$ in the orthonormal basis $(h_k)$ and pass to the limit in the bilinear form $(f_1,f_2)\mapsto\E[f_1(U)f_2(V)]$, which the Cauchy--Schwarz inequality bounds by $(\E f_1(z)^2\,\E f_2(z)^2)^{1/2}$. For the last statement, $\ip{w_\star,x}$ and $\ip{w,x}$ are standard Gaussians with correlation $\ip{w,w_\star}$.
\end{proof}

For a unit vector $w$ with overlap $r$ and a fresh input $x$, write $u=\ip{w_\star,x}$, $v=\ip{w,x}$, and $F=g(u)\,\phi'(v)$, so that the spherical gradient is $G=-F\,P_w^\perp x$. Two scalars govern one step: the component of $-G$ along $w_\star$, which moves the overlap, and the squared length of $G$, which drives the normalization,
\[
  a=-\ip{G,w_\star}=F\,(u-rv),\qquad N=\norm{G}^2=F^2\norm{P_w^\perp x}^2 .
\]
One function of the overlap controls the size of both. For $-1\le r\le1$ let $v$ and $z$ be independent standard Gaussians, and put
\begin{equation}\label{eq:singleindex:lambda}
  \lambda(r)=\E\Bigl[g\bigl(rv+\sqrt{1-r^2}\,z\bigr)^2\,\phi'(v)^2\,(1+z^2)\Bigr].
\end{equation}

\begin{lemma}[One step in terms of the overlap]\label[lemma]{lem:singleindex:onestep}
Let $w\in\Sph^{d-1}$ have overlap $r$, let $x\sim\cN(0,I_d)$, and define $u$, $v$, $F$, $a$, and $N$ as above. Then
\begin{enumerate}
\item[\textup{(a)}] $\E[F\,P_w^\perp x]=\Psi'(r)\,P_w^\perp w_\star$, and consequently $\E a=(1-r^2)\,\Psi'(r)$;
\item[\textup{(b)}] $\E a^2\le\lambda(r)$ and $\E N\le d\,\lambda(r)$.
\end{enumerate}
\end{lemma}

\begin{proof}
For (a), let $e$ be a unit vector orthogonal to $w$, and let $c(\theta)=w\cos\theta+e\sin\theta$, a curve on the sphere through $w$ with velocity $e$. By \cref{lem:singleindex:hermite}, $\E[g(u)\phi(\ip{c(\theta),x})]=\Psi(\ip{c(\theta),w_\star})$ for every $\theta$. Differentiate at $\theta=0$, under the expectation on the left, which dominated convergence allows because the integrand is a polynomial in Gaussian variables: $\E[F\ip{e,x}]=\Psi'(r)\ip{e,w_\star}$. The vector $m=\E[F\,P_w^\perp x]$ is orthogonal to $w$ and satisfies $\ip{m,e}=\E[F\ip{e,x}]$ for every such $e$, so $m=\Psi'(r)\,P_w^\perp w_\star$. Since $a=F\ip{P_w^\perp x,w_\star}$ and $\ip{P_w^\perp w_\star,w_\star}=1-r^2$, the inner product of $m$ with $w_\star$ is $\E a=(1-r^2)\Psi'(r)$.

For (b), suppose first that $\abs r<1$. Put $e=(w_\star-rw)/\sqrt{1-r^2}$ and $z=\ip{e,x}$, and let $x'$ be the projection of $x$ onto the orthogonal complement of $w$ and $w_\star$. Then $v$, $z$, and $x'$ are independent, $v$ and $z$ are standard Gaussians, $u=rv+\sqrt{1-r^2}\,z$, $P_w^\perp x=ze+x'$, and $\E\norm{x'}^2=d-2$. Hence $a=\sqrt{1-r^2}\,Fz$, so that $\E a^2\le\E[F^2z^2]\le\lambda(r)$, and
\[
  \E N=\E\bigl[F^2\bigl(z^2+\norm{x'}^2\bigr)\bigr]=\E[F^2z^2]+(d-2)\,\E[F^2]\le d\,\lambda(r).
\]
In the degenerate case $w=\pm w_\star$ we have $u=\pm v$, so $a=0$; and $P_w^\perp x$ is independent of $v$ with $\E\norm{P_w^\perp x}^2=d-1$, so $\E N=(d-1)\E[F^2]\le d\,\lambda(\pm1)$.
\end{proof}

\begin{remark}[The overlap is a Markov chain]\label[remark]{rem:singleindex:markov}
The proof of (b) shows more. Given the past, the triple $(u_t,v_t,\norm{P_{w_t}^\perp x_t}^2)$ that enters step $t$ has the law of $(r_tv+\sqrt{1-r_t^2}\,z,\;v,\;z^2+Q)$, with $v$ and $z$ standard Gaussians and $Q$ a $\chi^2_{d-2}$ variable, all independent; and $r_{t+1}$ is a function of $r_t$ and this triple (\cref{lem:singleindex:ineq} below). So $(r_t)$ is a Markov chain on $[-1,1]$ whose transition depends on $d$ only through $Q$. Simulating it costs a few operations per step in every dimension, and we use it in \cref{sec:singleindex:reading} (\cref{ex:singleindex:markov}).
\end{remark}

The population flow is the limit of \eqref{eq:singleindex:sgd} as $\eta\to0$ with $\eta t$ fixed.\sidenote{The reduction to one order parameter is exact here, in every dimension, because the student has one neuron; the committee machine of \cref{ch:committee} needs a matrix of overlaps.} Its overlap obeys a closed equation, and the time it needs to leave the equator already shows the role of $k$.

\begin{theorem}[Population dynamics and escape times]\label{thm:singleindex:flow}
Suppose $\hat g_j\hat\phi_j\ge0$ for every $j\ge1$ and that $\Psi$ is not constant, and let $k=k^\star$. The spherical gradient flow $\dot w=-P_w^\perp\nabla\Loss(w)$ started at $w_0\in\Sph^{d-1}$ stays on the sphere, and its overlap solves
\begin{equation}\label{eq:singleindex:ode}
  \dot r=(1-r^2)\,\Psi'(r).
\end{equation}
If $0<r_0<\frac12$, then $r$ reaches $\frac12$ at the time $t_{1/2}=\int_{r_0}^{1/2}\dd s/((1-s^2)\Psi'(s))$, which satisfies
\[
  \frac{I_k(r_0)}{2^{k-1}\Psi'(\frac12)}\le t_{1/2}\le\frac{4\,I_k(r_0)}{3k\,\hat g_k\hat\phi_k},\qquad I_k(r_0)=\int_{r_0}^{1/2}s^{1-k}\dd s .
\]
For $r_0=d^{-1/2}$ with $d>4$, the integral is $I_1=\frac12-d^{-1/2}$, $I_2=\frac12\log d-\log2$, and $I_k=(d^{(k-2)/2}-2^{k-2})/(k-2)$ for $k\ge3$: the escape takes a time of order $1$, $\log d$, and $d^{(k-2)/2}$ respectively.
\end{theorem}

\begin{proof}
By \cref{lem:singleindex:onestep}(a), every unit vector $w$ satisfies
\[
  -P_w^\perp\nabla\Loss(w)=\E[F\,P_w^\perp x]=\Psi'(\ip{w,w_\star})\,(w_\star-\ip{w,w_\star}w).
\]
Read the right side as a polynomial vector field on $\R^d$. Along its solutions $q=\norm w^2-1$ satisfies the linear equation $\dot q=-2\Psi'(r)\,r\,q$, so $q$ stays $0$ when it starts at $0$: the flow stays on the compact sphere and therefore exists for all time. The inner product of the flow with $w_\star$ is \eqref{eq:singleindex:ode}.

Nonnegative coefficients make $s\mapsto\Psi'(s)/s^{k-1}=\sum_{j\ge k}j\hat g_j\hat\phi_js^{j-k}$ nondecreasing on $(0,1]$, so for $0<s\le\frac12$
\[
  k\,\hat g_k\hat\phi_k\,s^{k-1}\le\Psi'(s)\le2^{k-1}\Psi'(\tfrac12)\,s^{k-1},
\]
and $\frac34\le1-s^2\le1$. The right side of \eqref{eq:singleindex:ode} is therefore bounded below by a positive number on $[r_0,\frac12]$, so $r$ increases and reaches $\frac12$ in finite time, and the substitution $s=r(t)$ gives the formula for $t_{1/2}$. The two displayed bounds on $\Psi'$ bound the integrand $1/((1-s^2)\Psi'(s))$ above and below, which gives the bounds on $t_{1/2}$; the values of $I_k$ are elementary integrals.
\end{proof}

For $k\ge2$ the population flow is not slow because the landscape is rough; on the hemisphere $r>0$ it has no critical points at all. It is slow because the landscape is flat where the search starts. Online SGD adds noise to this weak drift, and the noise sets the step size.

\begin{heuristic}[Why the step size must be of order $d^{-k/2}$]
Take $\Psi(r)=r^k$ and expand one step of \eqref{eq:singleindex:sgd} to second order in $\eta$ (the exact form is \cref{lem:singleindex:ineq}): at overlap $\rho\ll1$ the overlap changes by a drift $\eta k\rho^{k-1}$, a centered noise $\eta\xi$ with $\E\xi^2$ of order one, and a normalization term $-\frac12\eta^2N\rho$ with $\E N$ of order $d$, by \cref{lem:singleindex:onestep}. To double the overlap, the drift needs $L\approx\rho^{2-k}/(\eta k)$ steps. Over these steps the noise accumulates to about $\eta\sqrt L\approx(\eta\rho^{2-k}/k)^{1/2}$, which is small compared with $\rho$ only if $\eta\ll\rho^k$; the normalization accumulates to about $\eta^2d\rho L\approx\eta d\rho^{3-k}/k$, which is small compared with $\rho$ only if $\eta\ll\rho^{k-2}/d$. For $k\ge3$ both conditions are hardest at the start, $\rho=d^{-1/2}$, where both read $\eta\ll d^{-k/2}$. The doubling times form a geometric series dominated by its first term, about $d^{(k-2)/2}/(\eta k)$, which is of order $d^{k-1}$ when $\eta$ is of order $d^{-k/2}$. In words: the number of samples is the population escape time divided by the largest step that the noise at the equator allows. For $k=1$ the binding condition is the normalization at $\rho\approx\frac12$, which gives $\eta\approx1/d$ and $d$ samples; for $k=2$ every doubling takes about $1/(2\eta)$ steps and both conditions read $\eta\ll1/d$, which gives about $d\log d$ samples.
\end{heuristic}

\section{The solution: online SGD needs \texorpdfstring{$d^{k-1}$}{d\^{}(k-1)} samples}\label{sec:singleindex:sgd}

We now make the heuristic rigorous for $k\ge3$. The proof uses only the following properties of the pair $(g,\phi)$.

\begin{assumption}[The search-phase hypotheses]\label[assumption]{ass:singleindex:search}
The functions $g$ and $\phi$ are polynomials, $k\ge3$, and there are constants $c_0>0$ and $\Lambda<\infty$ such that $\Psi'(r)\ge c_0r^{k-1}$ and $\lambda(r)\le\Lambda$ for $0\le r\le\frac12$.
\end{assumption}

\begin{example}[Matched Hermite links]\label[example]{exm:singleindex:hermite}
For $g=\phi=h_k$ we have $\Psi(r)=r^k$, so the drift condition holds with $c_0=k$. The function $\lambda$ is a polynomial with $\lambda(0)=2k(k+1)$ (\cref{ex:singleindex:constants}). For $k=3$ a direct expansion gives $\lambda(r)=6(4+45r^2+12r^4-30r^6)$, which increases on $[0,\frac12]$, so $\Lambda=\lambda(\frac12)=1491/16<94$; for $k=4$ one finds $\Lambda=\lambda(\frac12)=3633/8<455$. More generally, every matched link $g=\phi$ with $\hat g_1=\dots=\hat g_{k-1}=0$ satisfies the drift condition with $c_0=k\hat g_k^2$.
\end{example}

The first step is an exact form of the expansion used in the heuristic. With $(r)_+=\max(r,0)$ and $\cF_t$ the $\sigma$-algebra generated by $x_0,\dots,x_{t-1}$ (the start $w_0$ is fixed, so $w_t$ and $r_t$ are $\cF_t$-measurable), write $a_t$ and $N_t$ for the quantities of \cref{lem:singleindex:onestep} computed with $w=w_t$ and the sample $x_t$.

\begin{lemma}[One step of SGD, bounded below]\label[lemma]{lem:singleindex:ineq}
For every $t$,
\begin{equation}\label{eq:singleindex:ineq}
  r_{t+1}\ge r_t+\eta a_t-\eta^2E_t,\qquad E_t=\tfrac12N_t\,(r_t)_++\abs{a_t}\sqrt{N_t/2},
\end{equation}
and $\E[E_t\mid\cF_t]\le\lambda(r_t)\bigl(\frac d2(r_t)_++\sqrt{d/2}\bigr)$.
\end{lemma}

\begin{proof}
Because $G_t$ is orthogonal to $w_t$, $\norm{w_t-\eta G_t}^2=1+\eta^2N_t$, and the inner product of \eqref{eq:singleindex:sgd} with $w_\star$ gives the exact recursion
\[
  r_{t+1}=\frac{r_t+\eta a_t}{\sqrt{1+\eta^2N_t}}=(r_t+\eta a_t)(1-\theta_t),\qquad\theta_t=1-(1+\eta^2N_t)^{-1/2}.
\]
Convexity of $q\mapsto(1+q)^{-1/2}$ gives $(1+q)^{-1/2}\ge1-q/2$ for $q\ge0$, so $0\le\theta_t\le\min(1,\eta^2N_t/2)$; and $\min(1,p)\le\sqrt p$ for $p\ge0$ gives also $\theta_t\le\eta\sqrt{N_t/2}$. Bound $r_t+\eta a_t$ above by $(r_t)_++\eta\abs{a_t}$, then use the first bound on $\theta_t$ in the third term below and the second bound in the fourth:
\[
  r_{t+1}\ge r_t+\eta a_t-\theta_t(r_t)_+-\theta_t\,\eta\abs{a_t}\ge r_t+\eta a_t-\tfrac12\eta^2N_t(r_t)_+-\eta^2\abs{a_t}\sqrt{N_t/2} .
\]
For the conditional expectation, $w_t$ is $\cF_t$-measurable and $x_t$ is independent of $\cF_t$, so \cref{lem:singleindex:onestep}(b) applies conditionally with $w=w_t$: $\E[N_t\mid\cF_t]\le d\lambda(r_t)$, and by the Cauchy--Schwarz inequality $\E[\abs{a_t}\sqrt{N_t}\mid\cF_t]\le(\lambda(r_t)\cdot d\lambda(r_t))^{1/2}$.
\end{proof}

The proof of the theorem cuts the search into epochs in which the overlap doubles. The next lemma treats one epoch; it is stated for a general band of overlaps because the descent phase will use it too. For a unit vector $w$ and a fresh sample, $E$ denotes $\frac12N(r)_++\abs a\sqrt{N/2}$, the quantity $E_t$ of \eqref{eq:singleindex:ineq} at $w_t=w$.

\begin{lemma}[One epoch]\label[lemma]{lem:singleindex:epoch}
Let $h>0$, $\ell<\ell'$, $m\ge0$, $v\ge0$, and $e\ge0$ be numbers such that every $w\in\Sph^{d-1}$ with $\ell-h\le\ip{w,w_\star}<\ell'$ satisfies $\E a\ge m$, $\E a^2\le v$, and $\E E\le e$. Let $\tau$ be a stopping time with $r_\tau\ge\ell$ on the event $\{\tau<\infty\}$, let $L\ge1$ be an integer, and let $\sigma$ be the first time $t\ge\tau$ with $r_t\notin[\ell-h,\ell')$. Then, outside an event of probability at most
\begin{equation}\label{eq:singleindex:epochfail}
  \frac{16\eta^2vL}{h^2}+\frac{2\eta^2eL}{h},
\end{equation}
either $\tau=\infty$ or $r_t>\ell-h+\eta m(t-\tau)$ for every integer $t$ with $\tau\le t\le\min(\sigma,\tau+L)$. Consequently, outside the same event and if $\tau<\infty$:
\begin{enumerate}
\item[\textup{(i)}] if $\eta mL\ge\ell'-\ell+h$, then $r_t\ge\ell'$ for some $t\le\tau+L$, and $r_s>\ell-h$ for $\tau\le s\le t$;
\item[\textup{(ii)}] if $\ell'>1$, then $r_t>\ell-h$ for $\tau\le t\le\tau+L$.
\end{enumerate}
\end{lemma}

\begin{proof}
The plan is to split the change of the overlap after $\tau$ into a drift, a martingale, and a nonnegative loss, and to bound the martingale by Doob's inequality and the loss by Markov's inequality. Let $\mu_s=\E[a_s\mid\cF_s]$ and $\xi_s=a_s-\mu_s$, and let $\chi_s$ be the indicator of the event $\{\tau\le s<\min(\sigma,\tau+L)\}$. This event belongs to $\cF_s$, because $\tau$ and $\sigma$ are stopping times and $\{s<\tau+L\}$ is the complement of $\{\tau\le s-L\}$; on it $r_s\in[\ell-h,\ell')$, so that $\mu_s\ge m$, $\E[\xi_s^2\mid\cF_s]\le\E[a_s^2\mid\cF_s]\le v$, and $\E[E_s\mid\cF_s]\le e$. Consider
\[
  M_t=\sum_{s<t}\chi_s\xi_s,\qquad S=\sum_{s\ge0}\chi_sE_s .
\]
The process $(M_t)$ is a martingale with orthogonal increments, so $\E M_t^2=\sum_{s<t}\E\bigl[\chi_s\E[\xi_s^2\mid\cF_s]\bigr]\le vL$ for every $t$, since at most $L$ of the indicators equal $1$. Doob's $L^2$ inequality (\cref{lem:toolkit:doob}) gives $\E\max_{t\le n}M_t^2\le4vL$ for every $n$, hence $\E\sup_tM_t^2\le4vL$ by monotone convergence, and Markov's inequality gives $\Pr(\sup_t\abs{\eta M_t}\ge h/2)\le16\eta^2vL/h^2$. In the same way $\E S\le eL$, so $\Pr(\eta^2S\ge h/2)\le2\eta^2eL/h$. Let $\cE$ be the event that $\sup_t\abs{\eta M_t}<h/2$ and $\eta^2S<h/2$; its complement has probability at most \eqref{eq:singleindex:epochfail}.

On $\cE\cap\{\tau<\infty\}$, let $\tau\le t\le\min(\sigma,\tau+L)$. Every $s$ with $\tau\le s<t$ has $\chi_s=1$ and no $s<\tau$ has, so summing \eqref{eq:singleindex:ineq} over $\tau\le s<t$ and using $\mu_s\ge m$ gives
\[
  r_t\ge r_\tau+\eta\sum_{s=\tau}^{t-1}\mu_s+\eta M_t-\eta^2\sum_{s=\tau}^{t-1}E_s>\ell+\eta m(t-\tau)-\frac h2-\frac h2 ,
\]
which is the main claim. For (i), suppose $\sigma>\tau+L$; then $t=\tau+L$ is admissible, so $r_t>\ell-h+\eta mL\ge\ell'$, whereas $t<\sigma$ forces $r_t<\ell'$. Hence $\sigma\le\tau+L$, and since $r_\sigma>\ell-h$, the overlap leaves the band through the top: $r_\sigma\ge\ell'$. For (ii), the overlap never reaches $\ell'>1$, so an exit at a time $\sigma\le\tau+L$ would have $r_\sigma<\ell-h$, contradicting the main claim; hence $\sigma>\tau+L$, and the main claim with $m\ge0$ gives (ii).
\end{proof}

\begin{theorem}[Online SGD finds the direction]\label{thm:singleindex:sgd}
Suppose \cref{ass:singleindex:search} holds, let $0<\eps\le1$, and let $r_0\ge d^{-1/2}$. Run \eqref{eq:singleindex:sgd} with step size and horizon
\[
  \eta=\frac{c_0\,\eps}{2^{k+8}\Lambda}\,d^{-k/2},\qquad T=\frac{2^{2k+10}\Lambda}{c_0^2\,\eps}\,d^{k-1}.
\]
Then with probability at least $1-\eps$ some $t\le T$ has $r_t\ge\frac12$. In particular $T$ samples suffice to leave the search phase.
\end{theorem}

\begin{proof}
The plan is to cut the climb from $r_0$ to $\frac12$ into epochs in which the overlap doubles, to apply \cref{lem:singleindex:epoch}(i) to each, and to add up the failure probabilities, which form a geometric series because $k\ge3$. If $r_0\ge\frac12$, there is nothing to prove. Otherwise let $J\ge1$ be the least integer with $2^Jr_0\ge\frac12$, put $\rho_j=2^jr_0$ for $0\le j<J$ and $\rho_J=\frac12$, so that $\rho_{j+1}=\min(2\rho_j,\frac12)$, and let $\tau_j$ be the first time $t$ with $r_t\ge\rho_j$. In particular $\tau_0=0$.

\emph{The constants of an epoch.} Fix $j<J$ and use \cref{lem:singleindex:epoch} with $\tau=\tau_j$, $\ell=\rho_j$, $h=\rho_j/2$, and $\ell'=\rho_{j+1}$. Every overlap $r$ in the band $[\rho_j/2,\rho_{j+1})$ lies in $(0,\frac12)$, so \cref{lem:singleindex:onestep}(a) and \cref{ass:singleindex:search} give $\E a=(1-r^2)\Psi'(r)\ge\frac34c_0(\rho_j/2)^{k-1}=:m_j$, and \cref{lem:singleindex:onestep}(b) gives $\E a^2\le\Lambda$. \Cref{lem:singleindex:ineq} gives $\E E\le\Lambda(\frac d2\cdot2\rho_j+\sqrt{d/2})\le2\Lambda d\rho_j$, because $r<2\rho_j$ in the band and $\sqrt{d/2}\le d\rho_j$ (here $\rho_j\ge r_0\ge d^{-1/2}$ is used). Let $L_j$ be the least integer at least
\[
  x_j=\frac{3\rho_j}{2\eta m_j}=\frac{2^k}{c_0\,\eta\,\rho_j^{k-2}} .
\]
Since $\ell'-\ell+h\le\frac32\rho_j$, the hypothesis $\eta m_jL_j\ge\ell'-\ell+h$ of (i) holds. So outside an event $\cB_j$ whose probability is at most \eqref{eq:singleindex:epochfail}, if $\tau_j<\infty$, then $\tau_{j+1}\le\tau_j+L_j$.

\emph{The failure probability.} The number $x_j$ is at least $1$, so $L_j\le2x_j$. To see this, compare the drift at overlap $\frac12$ with its noise: $\frac34c_02^{1-k}\le\frac34\Psi'(\frac12)=\E a\le(\E a^2)^{1/2}\le\Lambda^{1/2}$, so $c_0^2<4^k\Lambda$ and $x_j\ge2^k/(c_0\eta)\ge2^{2k+8}\Lambda/c_0^2>1$. Inserting $v=\Lambda$, $e=2\Lambda d\rho_j$, $h=\rho_j/2$, and $L_j\le2x_j$ into \eqref{eq:singleindex:epochfail},
\[
  \Pr(\cB_j)\le\frac{64\eta^2\Lambda L_j}{\rho_j^2}+8\eta^2\Lambda dL_j\le\frac{2^{k+1}\Lambda\eta}{c_0}\Bigl(\frac{64}{\rho_j^k}+\frac{8d}{\rho_j^{k-2}}\Bigr).
\]
Since $\rho_j=2^jr_0$ and $r_0\ge d^{-1/2}$, the geometric series give $\sum_{j<J}\rho_j^{-k}\le\frac87d^{k/2}$ and, because $k\ge3$, $\sum_{j<J}\rho_j^{2-k}\le2d^{(k-2)/2}$. Hence
\[
  \sum_{j<J}\Pr(\cB_j)\le\frac{2^{k+1}\Lambda\eta\,d^{k/2}}{c_0}\Bigl(\frac{512}{7}+16\Bigr)<\frac{2^{k+8}\Lambda\eta\,d^{k/2}}{c_0}=\eps .
\]

\emph{The time.} Outside the union of the events $\cB_j$, induction on $j$ gives $\tau_j\le L_0+\dots+L_{j-1}$, and therefore
\[
  \tau_J\le\sum_{j<J}L_j\le\frac{2^{k+1}}{c_0\eta}\sum_{j<J}\rho_j^{2-k}\le\frac{2^{k+2}d^{(k-2)/2}}{c_0\,\eta}=T .\qedhere
\]
\end{proof}

Three remarks place the theorem. The constants are explicit but not good: for $k=3$ and $\eps=0.1$ the step size is $1.6\times10^{-6}d^{-3/2}$, while the simulations of \cref{sec:singleindex:reading} escape reliably with steps about $10^4$ times larger. The losses come from Chebyshev-type bounds, the margins of $\rho_j/4$ allowed to the noise and to the normalization, the factor $2^{k-1}$ lost by bounding the drift at the bottom of each band, and the use of $\lambda(\frac12)$ instead of $\lambda(0)$; the exponents are what the proof gets right. Second, the warm start $r_0\ge d^{-1/2}$ holds with probability about $0.16$ for a uniformly random $w_0$, because $\sqrt d\,r_0$ is approximately standard Gaussian; for even links the symmetry $w\mapsto-w$ makes $\abs{r_0}\ge d^{-1/2}$ enough, and a start at $\gamma d^{-1/2}$ with $0<\gamma<1$ costs a factor $\gamma^k$ in the step size (\cref{ex:singleindex:otherk}). For odd $k$ a start with $r_0<0$ is genuinely bad: the population flow then tends to the equator from below and never crosses it (\cref{ex:singleindex:sign}). Third, the failure probability enters as $1/\eps$, the price of using only second moments; \textcite{BenArousGheissariJagannath2021} obtain probability $1-o(1)$ from random starts at the cost of factors $\log^2d$, and \textcite{ZweigEtAl2023} prove bounds of the same type for a class of non-Gaussian inputs.\sidenote{The doubling epochs, the martingale bound, and the treatment of the normalization are standard; we have arranged them so that every constant can be checked by hand.}

After the search, the overlap must still climb from $\frac12$ to nearly $1$. This descent phase\index{descent phase} is fast, and its accuracy is limited by the step size.

\begin{proposition}[The descent phase]\label[proposition]{prop:singleindex:descent}
Let $g=\phi=h_k$ with $k\ge3$, let $d\ge16$, let $\bar\Lambda=\max_{0\le r\le1}\lambda(r)$, and let $\eta$ and $T$ be as in \cref{thm:singleindex:sgd} with $c_0=k$ and $\Lambda=\bar\Lambda$. If $r_0\ge d^{-1/2}$, then for every $\delta\in(0,1)$ and every integer $t$ with $T\le t\le2T$,
\[
  \Pr\bigl(r_t<1-\delta\bigr)\le2\eps+\frac1\delta\Bigl(e^{-\kappa\eta(t-T)}+\frac{2^{k-10}\,\eps}{d^{(k-2)/2}}\Bigr),\qquad\kappa=k\,4^{1-k}.
\]
\end{proposition}

\begin{proof}
The plan is to wait for the search to end, to show with \cref{lem:singleindex:epoch}(ii) that the overlap then stays above $\frac14$, and to show that above $\frac14$ the distance $D_t=1-r_t$ contracts in expectation. The hypotheses of \cref{thm:singleindex:sgd} hold with $c_0=k$ and the larger constant $\bar\Lambda$, so the first time $\tau$ with $r_\tau\ge\frac12$ satisfies $\Pr(\tau>T)\le\eps$.

\emph{The overlap stays above $\frac14$.} For overlaps $r\in[\frac14,1]$, \cref{lem:singleindex:onestep} gives $\E a=kr^{k-1}(1-r^2)\ge0$ and $\E a^2\le\bar\Lambda$, and \cref{lem:singleindex:ineq} gives $\E E\le\bar\Lambda(\frac d2+\sqrt{d/2})\le\bar\Lambda d$. Apply \cref{lem:singleindex:epoch}(ii) with $\ell=\frac12$, $h=\frac14$, $\ell'=2$, $m=0$, and $L=\floor{2T}$: outside an event of probability at most $8\eta^2\bar\Lambda L(32+d)$, if $\tau<\infty$, then $r_s>\frac14$ for $\tau\le s\le\tau+L$, in particular for $\tau\le s\le2T$. The definitions of $\eta$ and $T$ give $\eta^2\bar\Lambda T=\eps/(64d)$, so this probability is at most $\eps(32+d)/(4d)\le\frac34\eps$ when $d\ge16$.

\emph{Contraction.} Let $\nu$ be the first time $s\ge\tau$ with $r_s<\frac14$, and for $s\ge T$ let $V_s=D_s\,\ind\{\tau\le T,\ s<\nu\}$. The event in the indicator belongs to $\cF_s$, and on it $\frac14\le r_s\le1$, so that $\E[a_s\mid\cF_s]=kr_s^{k-1}(1+r_s)(1-r_s)\ge\kappa D_s$ and $\E[E_s\mid\cF_s]\le\bar\Lambda d$. By \eqref{eq:singleindex:ineq}, $D_{s+1}\le D_s-\eta a_s+\eta^2E_s$, and $D_{s+1}\ge0$. The indicator can only switch off as $s$ grows, so
\[
  \E V_{s+1}\le\E\bigl[\ind\{\tau\le T,\ s<\nu\}\,\E[D_{s+1}\mid\cF_s]\bigr]\le(1-\kappa\eta)\,\E V_s+\eta^2\bar\Lambda d .
\]
Starting from $\E V_T\le\frac34$ and iterating, then inserting the value of $\eta$,
\[
  \E V_t\le\tfrac34(1-\kappa\eta)^{t-T}+\frac{\eta\bar\Lambda d}{\kappa}\le e^{-\kappa\eta(t-T)}+\frac{2^{k-10}\,\eps}{d^{(k-2)/2}} .
\]

\emph{Conclusion.} If $\tau\le T$ and $\nu>t$, then $V_t=D_t$. So the event $\{r_t<1-\delta\}$ is contained in the union of $\{\tau>T\}$, the event that the overlap falls below $\frac14$ between $\tau$ and $2T$, and $\{V_t>\delta\}$, whose probabilities are at most $\eps$, $\frac34\eps$, and $\E V_t/\delta$ by Markov's inequality.
\end{proof}

At $t=2T$ the exponential term equals $\exp(-2^{4-k}d^{(k-2)/2})$, which is negligible: with probability at least $1-2\eps$, up to terms that vanish as $d$ grows, the overlap is within $\delta$ of $1$ for every fixed $\delta>0$. The descent takes about $\kappa^{-1}\eta^{-1}\log(1/(\eps\delta))$ steps, of order $d^{k/2}\log(1/(\eps\delta))/\eps$, a vanishing fraction of $T$: the samples are spent in the search.

\subsection{What is known to be sharp}\label{sec:singleindex:sharp}

\Cref{thm:singleindex:sgd} is an upper bound for one algorithm. Two kinds of lower bound complete the picture: one for online SGD, and one for every algorithm that learns from correlations. In the next statement, $a_d\ll b_d$ means $a_d/b_d\to0$ as $d\to\infty$, and the square loss replaces the correlation loss; on the sphere the two have the same population landscape.

\begin{theorem}[Sharp thresholds for online SGD; \textcite{BenArousGheissariJagannath2021}]\label{thm:singleindex:bagj}
Let $g$ be a polynomial with information exponent $k$, let $\phi=g$, and run \eqref{eq:singleindex:sgd} for the square loss $\frac12(y-\phi(\ip{w,x}))^2$ from $w_0$ uniform on the hemisphere $\{w:\ip{w,w_\star}\ge0\}$, with $n$ samples and step size $\eta$ depending on $d$, $n$ at most polynomial in $d$.
\begin{enumerate}
\item[\textup{(a)}] If $1/n\ll\eta\ll1/\sqrt{nd}$, and $n\gg d$ for $k=1$, $n\gg d\log^2d$ for $k=2$, or $n\gg d^{k-1}\log^2d$ for $k\ge3$, then $r_n\to1$ in probability.
\item[\textup{(b)}] If $n\ll d$ for $k=1$ or $n\ll d\log d$ for $k=2$, with $\eta=O(1/d)$, or if $n\ll d^{k-1}$ for $k\ge3$, with $\eta=O(1/\sqrt{nd})$, then $\max_{t\le n}\abs{r_t}\to0$ in probability.
\end{enumerate}
\end{theorem}

The thresholds $d$, $d\log d$, and $d^{k-1}$ are sharp up to the factors $\log^2d$. The lower bound for $k\ge3$ is proved only for step sizes $\eta=O(1/\sqrt{nd})$, of order $d^{-k/2}$ at the threshold; it says nothing about larger steps, which \cref{sec:singleindex:reading} shows to be harmful for a different reason. \Textcite{BarzilaiShamir2026} prove lower bounds for standard SGD, without the spherical constraint and without very small steps, by an argument outside the statistical-query framework described next. For networks that sum many single-index components with an even link of information exponent $k\ge3$, \textcite{RenEtAl2025} locate the sharp transition times at which online SGD picks up each direction.

The second lower bound concerns \term{correlational statistical queries}\index{correlational statistical query}: an algorithm chooses functions $\psi$ of the input with $\E\psi(x)^2\le1$ and receives, for each, a number within a tolerance $\tau$ of $\E[y\,\psi(x)]$, with the errors chosen by an adversary. The model is the correlational restriction \parencite{BshoutyFeldman2002} of the statistical queries of \textcite{Kearns1998}; a gradient of the square or correlation loss is such a query, estimated from samples.

\begin{theorem}[Correlational lower bound, informal; \textcite{DamianLeeSoltanolkotabi2022}, \textcite{DamianEtAl2023}]\label{thm:singleindex:csq}
Let $g$ satisfy $\E g(z)=0$ and $\E g(z)^2=1$ and have information exponent $k\ge1$. If an algorithm makes $q$ correlational statistical queries of tolerance $\tau$ and recovers $w_\star$ to a fixed accuracy for every $w_\star\in\Sph^{d-1}$, then $\tau\le C_k(\log(qd)/d)^{k/4}$, with $C_k$ depending only on $k$ and the accuracy.
\end{theorem}

An empirical average of $n$ samples estimates $\E[y\psi(x)]$ to accuracy about $n^{-1/2}$. Reading $\tau\approx n^{-1/2}$ turns the theorem into the heuristic bound $n\gtrsim(d/\log(qd))^{k/2}$ for gradient methods; it is a heuristic, not a theorem about SGD, because the errors of SGD are random rather than adversarial, as both papers stress. The bound has the same form for odd and even $k$; parity matters only on the algorithmic side, where the simplest matching estimator, a partial trace of the empirical Hermite tensor, needs about $d^{\lceil k/2\rceil}$ samples, and odd $k$ needs a further power iteration to reach $d^{k/2}$ \parencite{DamianEtAl2023,DamianEtAl2024}. For $k\le2$ the correlational bound is weaker than the trivial requirement $n\gtrsim d$, and online SGD is optimal up to logarithms. For $k\ge3$ it is $d^{k/2}$, and online SGD is off by the factor $d^{k/2-1}$. Smoothing the loss closes this gap (\cref{sec:singleindex:wrong}).

\section{Reading the solution}\label{sec:singleindex:reading}

The solution has three clocks. The population flow needs time of order $d^{(k-2)/2}$ to leave the equator. The step size is limited by what happens at the equator, where the drift is weakest: to order $d^{-k/2}$ for $k\ge3$, by the noise and the normalization together, and to order $1/d$ for $k\le2$, by the normalization. The number of samples is the first clock divided by the second: $d^{k-1}$ for $k\ge3$, $d\log d$ for $k=2$, and $d$ for $k=1$, the thresholds of \cref{thm:singleindex:bagj}. The samples are spent almost entirely near the equator: the epoch lengths in the proof of \cref{thm:singleindex:sgd} are proportional to $\rho_j^{2-k}$, so for $k\ge3$ the first doubling of the overlap costs at least as much as all later ones together, and the descent costs a vanishing fraction (\cref{prop:singleindex:descent}).

\begin{figure}[t]
\centering
\begin{tikzpicture}
\begin{axis}[hplot,xmode=log,ymode=log,log basis x=2,xmin=12,xmax=5600,ymin=100,ymax=3e6,
  xlabel={dimension $d$},ylabel={samples to reach $\abs{r_t}\ge\frac12$},legend pos=south east,
  xtick={16,64,256,1024,4096},xticklabels={$16$,$64$,$256$,$1024$,$4096$}]
\addplot[theory] table[x=d,y=k3] {figures/data/07-singleindex-escape-theory.dat};
\addplot[theory2] table[x=d,y=k2] {figures/data/07-singleindex-escape-theory.dat};
\addplot[theory3] table[x=d,y=k1] {figures/data/07-singleindex-escape-theory.dat};
\addplot[simulation,error bars/.cd,y dir=both,y explicit] table[x=d,y=median,
  y error plus expr=\thisrow{q75}-\thisrow{median},y error minus expr=\thisrow{median}-\thisrow{q25}]
  {figures/data/07-singleindex-escape-k3.dat};
\addplot[simulation,forget plot,error bars/.cd,y dir=both,y explicit] table[x=d,y=median,
  y error plus expr=\thisrow{q75}-\thisrow{median},y error minus expr=\thisrow{median}-\thisrow{q25}]
  {figures/data/07-singleindex-escape-k2.dat};
\addplot[simulation,forget plot,error bars/.cd,y dir=both,y explicit] table[x=d,y=median,
  y error plus expr=\thisrow{q75}-\thisrow{median},y error minus expr=\thisrow{median}-\thisrow{q25}]
  {figures/data/07-singleindex-escape-k1.dat};
\legend{population flow $k=3$,population flow $k=2$,population flow $k=1$,SGD median and quartiles}
\end{axis}
\end{tikzpicture}
\caption{Escape times of spherical online SGD for the matched links $g=\phi=h_k$, from $r_0=d^{-1/2}$, with $\eta=0.02/d$ for $k\le2$ and $\eta=0.02\,d^{-3/2}$ for $k=3$. Curves: the population prediction $t_{1/2}/\eta$ of \cref{thm:singleindex:flow}, with no fitted constant. Points: medians and quartiles of the first time $t$ with $r_t\ge\frac12$ (with $\abs{r_t}\ge\frac12$ for $k=2$), over $200$ runs for $k\le2$ and $100$ runs for $k=3$, simulated with the exact one-dimensional chain of \cref{rem:singleindex:markov} (seed $20261017$). One run in a hundred at $k=3$, $d=256$, had not escaped after three times the predicted time.}\label{fig:singleindex:escape}
\end{figure}

\Cref{fig:singleindex:escape} tests the first two clocks. The step sizes are those of the heuristic, $\eta=0.02/d$ for $k\le2$ and $\eta=0.02\,d^{-3/2}$ for $k=3$, and the curves are the population prediction with no fitted constant. The slopes are $1$ for $k=1$, slightly more than $1$ for $k=2$ because of the logarithm, and $2$ for $k=3$. The medians agree with the curves within $1\%$ for $k=1$, within $7\%$ for $k=2$, and within $12\%$ for $k=3$, the worst cases at $d=16$, where the start $r_0=\frac14$ is already halfway to the target. The spread is the other prediction. For $k=1$ the quartiles narrow from $\pm7\%$ of the median at $d=16$ to $\pm0.3\%$ at $d=4096$, for $k=2$ from $\pm16\%$ to $\pm5\%$, while for $k=3$ the interquartile range stays between $25\%$ and $35\%$ of the median for $d\ge32$. At the step size $\eta\approx d^{-k/2}$ the noise is as large as it can be without stopping the search, so the escape time stays random as $d\to\infty$. We checked the one-dimensional chain against $d$-dimensional SGD for $k=3$ at $d=16$ and $d=32$, with $1000$ runs each: the medians agree within $2\%$.

\begin{figure}[t]
\centering
\begin{tikzpicture}
\begin{groupplot}[group style={group size=2 by 1,horizontal sep=1.45cm},hplot,width=0.51\linewidth,height=0.47\linewidth,
  xlabel={samples$/d^2$},xmin=0,xmax=36]
\nextgroupplot[ymode=log,ymin=0.05,ymax=1.2,ylabel={overlap $r_t$},ytick={0.0625,0.125,0.25,0.5,1},
  yticklabels={$\frac1{16}$,$\frac18$,$\frac14$,$\frac12$,$1$}]
\pgfplotsinvokeforeach{1,...,8}{\addplot[muted!70,line width=0.35pt] table[x=tau,y=r#1] {figures/data/07-singleindex-traj.dat};}
\addplot[theory] table[x=tau,y=ode] {figures/data/07-singleindex-traj.dat};
\nextgroupplot[ymin=0,ymax=1.08,ylabel={loss $1-r_t^3$}]
\pgfplotsinvokeforeach{1,...,8}{\addplot[muted!70,line width=0.35pt] table[x=tau,y=l#1] {figures/data/07-singleindex-traj.dat};}
\addplot[theory] table[x=tau,y=lode] {figures/data/07-singleindex-traj.dat};
\end{groupplot}
\end{tikzpicture}
\caption{Eight runs of spherical online SGD in dimension $d=128$ for $g=\phi=h_3$, with $\eta=0.02\,d^{-3/2}$ and $r_0=d^{-1/2}$ (grey), and the population flow (curve). Left: the overlap, on a logarithmic scale. Right: the population square loss $\frac12\E(y-h_3(\ip{w_t,x}))^2=1-r_t^3$. The runs simulate \eqref{eq:singleindex:sgd} in $\R^{128}$ directly (seed $7$).}\label{fig:singleindex:trajectories}
\end{figure}

\Cref{fig:singleindex:trajectories} shows single runs. On a logarithmic scale the overlap grows from the start, as the population flow predicts; for $k=3$ the flow makes $1/r_t$ decrease almost linearly in $t$ while $r_t$ is small. Yet the loss $1-r_t^3$ stays within $10^{-2}$ of its initial value until the overlap passes $0.2$: progress is hidden from the loss because the loss sees $r^k$, not $r$, the same mechanism as the hidden progress of \cref{ch:parity}. The runs escape at different times, between $0.8$ and $1.4$ times the population prediction, and then descend in a short and nearly deterministic time. They settle at $r\approx0.985$, not at $1$. Near $r=1$ the normalization pulls the overlap down by about $\frac12\eta^2\E N\approx\frac12\eta^2(d-1)\E[F^2]$ per step, and the drift pulls it up by about $\eta kr^{k-1}(1-r^2)\approx2k\eta(1-r)$; they balance at $1-r\approx\eta(d-1)\E[F^2]/(4k)$, which is $0.014$ here, since $\E[F^2]=93$ at $r=1$ for $k=3$. A smaller step after the search removes this floor, as in the two-phase schedules of \textcite{DamianEtAl2023}.

The equator also has a picture of its own. Near $r=0$, \cref{lem:singleindex:onestep} and \cref{lem:singleindex:ineq} give, to leading order, $r_{t+1}-r_t\approx\eta kr_t^{k-1}-\frac12\eta^2kd\,r_t+\eta\xi_t$ with $\E\xi_t^2\approx k(2k+1)$ for the matched link $h_k$. Without the drift this is an autoregressive process whose stationary variance is $(2k+1)/d$ (\cref{ex:singleindex:stepsize}): SGD noise by itself keeps the overlap at the scale $d^{-1/2}$ of a random start. The drift beats the restoring force $\frac12\eta^2kd\,r$ only where $r^{k-2}>\eta d/2$, and this threshold lies inside the band of width $\sqrt{(2k+1)/d}$ exactly when $\eta\lesssim d^{-k/2}$. Larger steps do not stop the search outright but make it a rare event, a large deviation of the equatorial noise (\cref{ex:singleindex:stepsize}); the lower bound of \cref{thm:singleindex:bagj} does not cover them, and \textcite{Rangriz2026} proves limit theorems at the critical step size in which, near a fixed point, the rescaled summary statistics follow an Ornstein--Uhlenbeck process.\sidenote{At this step size the equator is a stable fixed point of the effective dynamics in the sense of \textcite{BenArousGheissariJagannath2022}, who prove limit theorems for summary statistics of SGD on time scales of order $1/\eta$.}

\section{What the model gets wrong}\label{sec:singleindex:wrong}

\paragraph{What was thrown away.}
Isotropic Gaussian inputs, which made the overlap a sufficient statistic and the Hermite polynomials the right basis; with structured or non-Gaussian inputs neither holds. One relevant direction and a known link: a network must find several directions and fit an unknown function of them, and the order in which directions are found becomes part of the problem. One neuron with nothing on top of it: width changes only the start, since the best of $m$ random neurons has overlap about $\sqrt{2\log m/d}$, which shortens the search by a factor of order $(\log m)^{(k-2)/2}$ (\cref{ex:singleindex:width}); for $k=2$, \textcite{ArnaboldiEtAl2023} find that overparameterization helps only by a constant factor. One pass over fresh samples, which is what makes samples and steps the same number and the noise a martingale. The sphere, which fixes the norm that real networks must learn. The correlation loss with a matched student, small constant steps, and no momentum or adaptive preconditioning. And a warm start.

\paragraph{Where it breaks.}
The model's central prediction, that a target with information exponent $k$ costs about $d^{k-1}$ samples, is a statement about one algorithm, and other gradient-based algorithms do much better (\cref{tab:singleindex:ladder}). \emph{Smoothing the loss.} \Textcite{DamianEtAl2023} run online SGD on a loss averaged over a cap of radius $\lambda$ around $w$; with $\lambda=d^{1/4}$ the signal near the equator grows faster than the noise, and $d^{k/2}$ samples suffice up to logarithms, which matches \cref{thm:singleindex:csq}. So $d^{k-1}$ is not a limit of correlational methods, only of plain online SGD. \emph{Reusing samples.} Two steps on the same batch let the update depend nonlinearly on the label, which acts as an implicit transformation of $y$. \Textcite{DandiEtAl2024} show that gradient descent that reuses a batch of $n\asymp d$ samples finds, within two steps, directions that one-pass SGD with as many samples cannot find, and \textcite{LeeEtAl2024} show that a two-layer network trained by SGD with reused minibatches learns every polynomial single-index target with $d$ samples up to polylogarithmic factors; \textcite{ArnaboldiEtAl2024a}, in a workshop paper, reach a similar conclusion for multi-index targets. \emph{The information exponent is not the hardness of the problem.} \Textcite{DamianEtAl2024} define the generative exponent\index{generative exponent}, the least information exponent of $T(y)$ over all square-integrable transformations $T$ of the label, prove that statistical-query and low-degree algorithms need $d^{k_{\mathrm g}/2}$ samples when it equals $k_{\mathrm g}$, and match this with a partial-trace algorithm. Every polynomial link has $k_{\mathrm g}\le2$; for $y=\He_3(\ip{w_\star,x})$ the generative exponent is $1$, so the hydrogen atom of this chapter is easy for the right algorithm. \emph{The loss and the student matter.} \Textcite{JoshiMisiakiewiczSrebro2024} show for sparse functions that gradient queries of the square loss are exactly as weak as correlational queries, while those of losses such as the absolute loss are as strong as general statistical queries; and a student whose Hermite coefficients vanish where the target's do not can lose the signal altogether (\cref{ex:singleindex:exponents}). \emph{Anisotropy.} When the input covariance has a spike aligned with $w_\star$, the effective dimension drops and fewer samples suffice \parencite{BaEtAl2023}; but \textcite{MousaviHosseiniEtAl2023} show that spherical gradient flow that ignores the covariance can then stay at the equator, while a normalization adapted to the covariance exploits the spike.

\begin{table}[t]
\centering\small
\caption{The cost, in samples and up to logarithmic factors, of finding $w_\star$ for $y=\He_k(\ip{w_\star,x})$ with $k\ge3$, by method. Statistical queries reach $d$ because the generative exponent of every polynomial link is at most $2$; the information-theoretic cost is of order $d$ for every single-index model.}\label{tab:singleindex:ladder}
\begin{tabular}{@{}lll@{}}
\toprule
method & samples & source\\
\midrule
kernels, random features & $d^k$ & \parencite{GhorbaniEtAl2021}; \cref{ch:lazy}\\
online SGD, one neuron & $d^{k-1}$ & \cref{thm:singleindex:sgd,thm:singleindex:bagj}\\
correlational queries & $d^{k/2}$ & \cref{thm:singleindex:csq}\\
smoothed online SGD & $d^{k/2}$ & \parencite{DamianEtAl2023}\\
statistical queries & $d$ & \parencite{DamianEtAl2024}\\
SGD with reused samples & $d$ & \parencite{LeeEtAl2024}\\
\bottomrule
\end{tabular}
\end{table}

\paragraph{Fine structure.}
The successors repair one omission each. For targets with several directions, the information exponent becomes the leap complexity\index{leap complexity}: \textcite{AbbeBoixAdseraMisiakiewicz2022} characterize the sparse targets learnable with order $d$ samples by a merged-staircase property, and \textcite{AbbeBoixAdseraMisiakiewicz2023} conjecture that online SGD needs about $d^{\max(L-1,1)}$ samples for a target of leap $L$, proving it in restricted Gaussian settings; the generative version of the leap, with a matching spectral algorithm that needs about $d^{\max(1,k/2)}$ samples for a leap $k$, is due to \textcite{DamianLeeBruna2025}. With large steps and $n\asymp d$ samples, a single gradient step already learns the linear part of a target \parencite{BaEtAl2022}. When the link is unknown, a shallow network can learn it together with the direction, with about $d^{k}$ samples for gradient flow on the empirical loss \parencite{BiettiEtAl2022}, and two-timescale gradient flow on multi-index targets passes a sequence of saddles predicted by the Hermite structure \parencite{BiettiBrunaPillaudVivien2025}. Beyond Gaussian inputs, a local growth condition on the population drift replaces the information exponent \parencite{ZweigEtAl2023}, and for spherically symmetric inputs spherical harmonics replace Hermite polynomials, with one estimator that is optimal in samples and another, based on online SGD, that is optimal in running time \parencite{JoshiEtAl2025}. The step size itself is a knob: \textcite{TsiolisMousaviHosseiniErdogdu2025} show a phase transition, as the learning rate grows, from the regime governed by the information exponent to the one governed by the generative exponent, for one-pass SGD and for SGD with batch reuse. The survey of \textcite{BrunaHsu2025} organizes these results. Open are whether plain minibatch SGD, with neither smoothing nor reuse, can reach $d^{k/2}$; the limit law of the escape time at the critical step size for odd $k$ (\cref{ex:singleindex:law}); and lower bounds for online SGD that hold for every step size.

\section{Model card}\label{sec:singleindex:card}

\modelcard{A loss plateau while a network searches for the relevant directions of a high-dimensional input; longer for targets whose low-degree Hermite components vanish.}
{$y=g(\ip{w_\star,x})$, $x\sim\cN(0,I_d)$; one neuron $\phi(\ip{w,x})$ on the sphere; online SGD on the correlation loss; order parameter $r=\ip{w,w_\star}$.}
{Population drift $(1-r^2)\Psi'(r)\approx k\hat g_k\hat\phi_kr^{k-1}$; flow escapes in time of order $d^{(k-2)/2}$; with $\eta\asymp\eps d^{-k/2}$, online SGD reaches $r=\frac12$ within $O(d^{k-1}/\eps)$ samples for $k\ge3$ (proved), then descends in $O(d^{k/2}\log(1/\eps\delta)/\eps)$ samples.}
{Sample complexities $d$, $d\log d$, $d^{k-1}$ for $k=1$, $2$, $\ge3$; a step size limited by the equator; flat loss with hidden progress in $r$; an accuracy floor proportional to $\eta d$.}
{It measures one algorithm: smoothing reaches $d^{k/2}$, sample reuse and label transforms reach about $d$ for polynomial links; anisotropic and non-Gaussian inputs change the exponent; one direction, known link, one neuron.}
{Generative exponent; correlational and statistical-query lower bounds; smoothing and partial traces; leap complexity for multi-index targets; effective dynamics of SGD.}

The search phase is the third appearance in this book of an escape from a saddle at small initialization. In \cref{ch:deeplinear} the escape is logistic, with time $\log(1/\text{initialization})/\text{strength}$, which is the case $k=2$ here; in \cref{ch:diagonal} the same mechanism orders the coordinates; here a degenerate saddle of order $k$ turns the logarithm into the power $d^{(k-2)/2}$. The committee machine of \cref{ch:committee} adds a symmetric saddle between students, analyzed with the same order parameters and the ODE method (\cref{thm:toolkit:ode-method}), and \cref{ch:parity} is the Boolean version of this chapter: a parity on $k$ coordinates has no correlation with low-degree functions, its hidden progress is the growth of $\log r$ under a flat loss, and its statistical-query bound is the counterpart of \cref{thm:singleindex:csq}.

Against \cref{ch:lazy} the chapter gives the first quantitative separation between feature learning and kernels: a kernel method needs about $d^k$ samples for a degree-$k$ target, one neuron trained by online SGD needs $d^{k-1}$, and better gradient methods need $d^{k/2}$ or fewer. Rotational invariance of Gaussian data, which reduced everything to one overlap here, reduces the committee machine to a few overlaps in \cref{ch:committee} and makes the spectra of \cref{ch:doubledescent} computable.

\begin{exercises}
\item \consolidation (Information exponents.) (a) Show that $z\mapsto\relu(z)$, $\tanh z$, $\sgn z$, and $z^3$ have information exponent $1$, that $z^2$ and $\abs z$ have information exponent $2$ with $\widehat{\abs{z}}_2=1/\sqrt\pi$, and that $\He_k$ has information exponent $k$. (b) Show that the Hermite coefficients of the ReLU vanish at every odd degree $j\ge3$. Conclude that a ReLU neuron trained on the correlation loss for the target $h_3(\ip{w_\star,x})$ has $\Psi\equiv0$: it receives no signal at all, at every overlap. \hint{$\relu(z)=\frac z2+\frac{\abs z}2$.}\label{ex:singleindex:exponents}

\item \consolidation (Noise at the equator.) Let $g=\phi=h_k$. (a) Use $zh_k=\sqrt{k+1}\,h_{k+1}+\sqrt k\,h_{k-1}$ to show that $\lambda(0)=2k(k+1)$, that $\E a^2=k(2k+1)$ at overlap $0$, and that $\E N=k(2k+1)+(d-2)k$ there. (b) Check the value $\lambda(\frac12)=1491/16$ of \cref{exm:singleindex:hermite} by expanding $h_3(\frac12v+\frac{\sqrt3}2z)$ in monomials.\label{ex:singleindex:constants}

\item \consolidation (The overlap is a Markov chain.) Prove the claim of \cref{rem:singleindex:markov}: given $\cF_t$, the law of $r_{t+1}$ depends only on $r_t$. Write the transition as an explicit function of $r_t$ and of independent $v,z\sim\cN(0,1)$ and $Q\sim\chi^2_{d-2}$.\label{ex:singleindex:markov}

\item \consolidation (Square loss.) For the square loss $\frac12(y-\phi(\ip{w,x}))^2$, show that the spherical gradient is that of the correlation loss plus $\phi(v)\phi'(v)P_w^\perp x$, that this extra term has mean zero, and that \cref{thm:singleindex:sgd} holds for the square loss with $g(u)$ replaced by $g(u)-\phi(v)$ in the definition \eqref{eq:singleindex:lambda} of $\lambda$.\label{ex:singleindex:square}

\item \consolidation\hard (Other exponents and colder starts.) (a) Adapt the proof of \cref{thm:singleindex:sgd} to $k=2$: show that $\eta$ of order $\eps/(d\log d)$ and $O(d\log^2d/\eps)$ samples suffice, and locate the logarithm that \cref{thm:singleindex:bagj} suggests should be removable. (b) Do the same for $k=1$, with $\eta$ of order $\eps/d$ and $O(d/\eps)$ samples. (c) For $k\ge3$ and $r_0\ge\gamma d^{-1/2}$ with $0<\gamma<1$, show that the conclusion holds with $\eta$ multiplied by $\gamma^k$ and $T$ by $\gamma^{2-2k}$. \hint{in (a) every epoch contributes the same failure probability, and there are about $\log_2\sqrt d$ epochs.}\label{ex:singleindex:otherk}

\item \variation (Label noise.) Let $y=g(\ip{w_\star,x})+\sigma\zeta$ with $\zeta\sim\cN(0,1)$ independent of $x$. Show that the drift of \cref{lem:singleindex:onestep}(a) is unchanged and that $\lambda(r)$ increases by $2\sigma^2\E\phi'(z)^2$, which is $2k\sigma^2$ for $\phi=h_k$. By what factor does the bound of \cref{thm:singleindex:sgd} grow?\label{ex:singleindex:noise}

\item \variation (A little leakage.) Let $g=\phi=h_3+\beta h_1$ with $0<\beta<1$, so that $\Psi(r)=r^3+\beta^2r$ and the information exponent is $1$. Solve the population flow near the equator and show that the escape from $r_0=d^{-1/2}$ takes time of order $\min(\sqrt d,1/\beta)$. For which $\beta$ does the degree-$1$ leakage dominate the search?\label{ex:singleindex:leak}

\item \variation (Starting below the equator.) Let $g=\phi=h_k$ and $-1<r_0<0$. (a) For even $k$, show that the population flow converges to $-w_\star$, and explain why that is as good as $w_\star$. (b) For odd $k$, show that $r(t)$ increases to $0$ and never becomes positive. What must the noise of SGD do in that case?\label{ex:singleindex:sign}

\item \variation (Width.) Start $m$ neurons independently and uniformly on the sphere, and train each by \eqref{eq:singleindex:sgd}. Show that the largest initial overlap is about $\sqrt{2\log m/d}$ when $\log m\ll d$, and use \cref{thm:singleindex:flow} to show that the fastest neuron escapes earlier than a typical one by a factor of order $(2\log m)^{(k-2)/2}$.\label{ex:singleindex:width}

\item \simulation (Escape times for $k=4$.) Using the chain of \cref{rem:singleindex:markov} with $g=\phi=h_4$, $r_0=d^{-1/2}$, and $\eta=0.02\,d^{-2}$, measure the median escape time to $\abs r\ge\frac12$ over $50$ runs for $d=16$, $32$, $64$, $128$. You should observe a slope approaching $3$ on a log-log plot, and medians about $30\%$ above $t_{1/2}/\eta$ from \cref{thm:singleindex:flow} at $d=16$ but within $10\%$ for $d\ge32$; explain why the agreement is worse at small $d$.\label{ex:singleindex:k4}

\item \simulation (The equator and the step size.) (a) Show that the recursion $r_{t+1}=(1-\frac12\eta^2kd)r_t+\eta\xi_t$ with independent $\xi_t$ of variance $k(2k+1)$ has stationary variance close to $(2k+1)/d$ when $\eta^2kd\ll1$. (b) For $k=3$ and $d=256$, simulate the chain of \cref{rem:singleindex:markov} with $\eta=c\,d^{-3/2}$ for $c$ between $0.25$ and $16$, and plot the median escape time divided by $d^2$ against $c$. You should observe a decrease like $1/(3c)$ for small $c$, a minimum for $c$ between $2$ and $4$, and a steep increase beyond (at $c=16$ no run escapes within $30d^2$ steps). Relate the increase to the equatorial band of (a).\label{ex:singleindex:stepsize}

\item \research (The law of the search time.) Let $g=\phi=h_k$ with $k\ge3$, $\eta=c\,d^{-k/2}$, and $r_0=d^{-1/2}$. Show heuristically that $\bar r_s=\sqrt d\,r_{\floor{sd^{k-1}}}$ approximately solves $\dd\bar r=c\bigl(k\bar r^{k-1}-\frac{ck}2\bar r\bigr)\dd s+c\sqrt{k(2k+1)}\,\dd W_s$, so that the escape time divided by $d^{k-1}$ should converge in law to the blow-up time of this equation started at $\bar r=1$. A proof is open as far as we know: the limit theorems of \textcite{BenArousGheissariJagannath2022} cover time scales of order $1/\eta$, shorter by the factor $d^{(k-2)/2}$.\label{ex:singleindex:law}

\item \research (Plain SGD at the correlational rate.) Can minibatch SGD without explicit smoothing or sample reuse, with suitably tuned batch size and step size, learn targets of information exponent $k\ge3$ with about $d^{k/2}$ samples? Test this numerically for $h_3$, and if it succeeds, try to identify an implicit smoothing. Open as far as we know; \textcite{ArnaboldiEtAl2024b} analyze the effect of the batch size on online SGD, and \textcite{DamianEtAl2023} the effect of explicit smoothing.\label{ex:singleindex:minibatch}
\end{exercises}

\notesandrefs
The rigorous theory of the search phase starts with \textcite{BenArousGheissariJagannath2021}, who named the information exponent, proved the thresholds of \cref{thm:singleindex:bagj} for a general class of losses, and described the search and descent phases; a precursor is their analysis of Langevin dynamics and gradient descent for spiked tensor models \parencite{BenArousGheissariJagannath2020}, where the order of the tensor plays the role of $k$ and sets the signal strength needed to leave the equator.\sidenote{In tensor PCA the equator is the set of vectors uncorrelated with the planted spike, and a random start lies on it for the same reason as here.} Before them, \textcite{DudejaHsu2018} had used the Hermite expansion of the link to learn single-index models with an estimator tailored to one harmonic at a time, and \textcite{TanVershynin2023} had shown, in work first posted in 2019, that online SGD solves phase retrieval, an example with $k=2$, from an arbitrary start.\sidenote{Phase retrieval, $y=\abs{\ip{w_\star,x}}$, is even, so its information exponent is $2$, and the escape is the logistic escape of \cref{ch:deeplinear}.} On the information-theoretic side, \textcite{MondelliMontanari2019} located the weak-recovery threshold of single-index models at a number of samples proportional to $d$.

The correlational lower bound is due to \textcite{DamianLeeSoltanolkotabi2022}, who also showed that two-layer networks trained by gradient descent beat kernel methods on such targets; \textcite{DamianEtAl2023} extended it to every link and matched it with smoothed online SGD. \Textcite{AbbeBoixAdseraMisiakiewicz2023} stated the corresponding bounds for multi-index targets in terms of the leap. The generative exponent of \textcite{DamianEtAl2024} explained why the correlational picture is not the whole story, and the sample-reuse papers \parencite{DandiEtAl2024,LeeEtAl2024,ArnaboldiEtAl2024a} showed that ordinary gradient methods can realize the label transformations it calls for. \Textcite{BrunaHsu2025} survey algorithms for single- and multi-index models.

The effective-dynamics theorems of \textcite{BenArousGheissariJagannath2022} turn the heuristic expansion of \cref{sec:singleindex:drift} into limit theorems for summary statistics on time scales of order $1/\eta$, and the ODE method of \cref{thm:toolkit:ode-method} is the special case without noise in the limit. For sharp transition times in networks with many neurons see \textcite{RenEtAl2025}; for the role of the batch size and the step size see \textcite{ArnaboldiEtAl2024b} and \textcite{TsiolisMousaviHosseiniErdogdu2025}.

\paragraph{Translating the literature.}
\Textcite{BenArousGheissariJagannath2021} write $N$ for our $d$, $M=\alpha N$ for our $n$, $\delta/N$ for our $\eta$, and the latitude $m$, the inner product of the parameter with the planted direction, for our $r$; their parameter lives on the sphere of dimension $N-1$, as ours does. Their population loss is our $-\Psi(r)$ up to a constant, and their information exponent, the order of the first nonvanishing derivative of the population loss at $m=0$, agrees with our $k^\star$ for matched links. \Textcite{DamianEtAl2023} write $\sigma$ for the known link, our $g=\phi$; $w^\star$ for $w_\star$; $\alpha=w\cdot w^\star$ for $r$; $\He_k$ for the monic Hermite polynomials, with coefficients $c_k=\E[\sigma\He_k]=\sqrt{k!}\,\hat g_k$; $k^\star$ for the information exponent; and $L(w;x;y)=1-\sigma(w\cdot x)y$ for our $\ell+1$, with $\lambda$ the smoothing radius, unrelated to our $\lambda(r)$. In \textcite{DamianEtAl2024} the letter $k^\star$ changes role and denotes the generative exponent, our $k_{\mathrm g}$. \Textcite{DudejaHsu2018} write $p$ for our $d$.

\printchapterbib
\end{refsection}

% Chapter 8 of The Hydrogen Atoms of Deep Learning: the soft committee machine and the ODE method.
\newcommand{\committeeM}{\mathsf{M}}
\newcommand{\committeeOm}{\bar\omega}
\newcommand{\committeelu}{\lambda_{\mathrm u}}

\casechapter{Plateaus and Specialization}{The committee machine and the ODE method}{Soft committee machine, teacher and student}{ch:committee}
\begin{refsection}[chapters/08-committee.bib]

\begin{chaptersummary}
A network with a few hidden units, trained online from small weights on a target computed by a network of the same form, often learns in two steps: the error drops to an intermediate value, stays there while the hidden units are interchangeable copies of one another, and drops again when they specialize. The soft committee machine, a sum of $\erf$ units learning from a teacher of the same kind on Gaussian inputs, reproduces this. We prove that as the dimension grows the overlaps between student and teacher weights follow a deterministic system of ordinary differential equations, with closed-form coefficients and a term of order $\eta^2$ that comes from the gradient noise. From it we derive a learning-rate threshold $\pi\sqrt{5/3}$ for a single unit, and, for $K$ students and $K$ orthonormal teachers, an exact symmetric fixed point with an explicit unstable eigenvalue, which gives a plateau whose length grows like $\log d$. The model is silent about why real networks seldom show such plateaus: structured data, overparameterization, preconditioned gradients, and large steps all change the picture, and the plateau's length is set by fluctuations that the limiting equations leave out.
\end{chaptersummary}

\section{The phenomenon: redundant units and a sudden division of labor}\label{sec:committee:phenomenon}

A hidden layer has a symmetry that a single neuron lacks. Permuting the hidden units of a two-layer network, together with their outgoing weights, does not change the function it computes, so every function the network can represent is represented many times, and the loss landscape repeats itself under the permutation group. Points fixed by a permutation, where two units carry the same weights, are special: the gradient there respects the symmetry, so two identical units receive identical updates and stay identical. Small random initializations drive the network toward such points, because at the start every unit sees roughly the same average target and moves the same way.

The consequence shows up in training curves as a plateau that has nothing to do with a search for a hidden direction. The chapter on single-index models (\cref{ch:singleindex}) produced plateaus because the target was invisible to low-degree statistics; here every hidden unit sees its signal at once, and the delay comes instead from the need to break a symmetry. Physicists who analyzed online learning in small sigmoidal networks during the 1990s called the redundant stretch the \term{symmetric phase}\index{symmetric phase} and the end of it \term{specialization}\index{specialization}.

\begin{observation}[\textcite{SaadSolla1995b}; \textcite{BiehlRieglerWohler1996}; \textcite{FukumizuAmari2000}]\label[observation]{obs:committee:plateau}
Train a two-layer network with a few sigmoidal hidden units, by online gradient descent from small random weights, on inputs with independent coordinates and labels produced by a network of the same form. The error falls quickly to an intermediate value and then stays nearly constant for a long stretch, during which the hidden units have nearly the same incoming weights, each a blend of all the teacher's units. The stretch ends with a second drop, in which the units differentiate and each aligns with one teacher unit. With several teacher units the dynamics can pass through several such plateaus in a cascade. The plateaus are associated with critical points at which units are duplicated, which exist for every target, loss, and activation function.
\end{observation}

The evidence is of an unusual kind. The first two sources combine an exact theory, the one this chapter develops, with simulations of the stochastic algorithm; the third is a structural statement about the loss landscape: critical points of a network with fewer hidden units reappear in a larger network as lines of critical points, along which training can stall \parencite{FukumizuAmari2000}.\sidenote{Fukumizu and Amari show that such a line can contain local minima as well as saddles, and that the conclusion does not depend on the target, the loss, or the activation, only on the hierarchical structure of the model.} Plateaus of this kind were a recognized problem of backpropagation: \textcite{Amari1998} suggested that the natural gradient would make it less serious. In current practice they are seldom seen: \textcite{YoshidaOkada2019} note that plateaus were studied intensively in the 1990s but rarely appear in modern deep learning, and they show that data whose covariance has small and dispersed eigenvalues make the plateau less noticeable. The phenomenon is robust in its native setting, small networks, small initializations, and isotropic inputs, and fragile outside it; \cref{sec:committee:wrong} returns to the reasons. A modern descendant, the passage of SGD through a sequence of saddles while two-layer networks learn hierarchical targets, is the subject of \textcite{AbbeBoixAdseraMisiakiewicz2023}.

The question the model answers: what sets the level of the plateau, why does it last as long as it does, and what in the dynamics decides when it ends?

\section{The model: the soft committee machine}\label{sec:committee:model}

Write $g(z)=\erf(z/\sqrt2)=2\Pr(\cN(0,1)\le z)-1$.\sidenote{So $g$ is the Gaussian distribution function rescaled to $(-1,1)$, the probit counterpart of $\tanh$.} This is an odd sigmoid with values in $(-1,1)$, and
\begin{equation}\label{eq:committee:gprime}
  g'(z)=\sqrt{\tfrac2\pi}\,e^{-z^2/2},\qquad g''(z)=-z\,g'(z).
\end{equation}
Every derivative of $g$ is a polynomial times $e^{-z^2/2}$, so all derivatives are bounded.

\begin{definition}[Soft committee machine, teacher and student]\label[definition]{def:committee:model}
Let $d\ge m+m_\star$. A \term{teacher}\index{teacher--student setup} has weights $w_{\star1},\dots,w_{\star m_\star}\in\R^d$ with Gram matrix $T=(\ip{w_{\star k},w_{\star k'}})$, and a \term{student} has weights $W=(w_1,\dots,w_m)$. For an input $x\sim\cN(0,I_d)$ put $z_j=\ip{w_j,x}$ and $z_{\star k}=\ip{w_{\star k},x}$; the label and the prediction are
\[
  y=\sum_{k=1}^{m_\star}g(z_{\star k}),\qquad f(x;W)=\sum_{j=1}^mg(z_j),
\]
and the population loss is $\Loss(W)=\frac12\E\bigl(y-f(x;W)\bigr)^2$. Online SGD with step size $\eta/d$ draws a fresh input $x^s$ at step $s=0,1,2,\dots$ and updates
\begin{equation}\label{eq:committee:sgd}
  w_j^{s+1}=w_j^s+\frac\eta d\,\delta_j^s\,x^s,\qquad \delta_j^s=g'(z_j^s)\bigl(y^s-f(x^s;W^s)\bigr),
\end{equation}
which is a gradient step on $\frac12(y^s-f(x^s;W^s))^2$. The \term{order parameters}\index{order parameters} are the overlaps
\[
  Q_{jl}=\ip{w_j,w_l},\qquad R_{jk}=\ip{w_j,w_{\star k}},\qquad T_{kk'}=\ip{w_{\star k},w_{\star k'}}.
\]
\end{definition}

This is the \term{soft committee machine}\index{soft committee machine}: a two-layer network whose second layer is fixed to ones, so that the units vote with equal weight, and whose units are soft versions of the sign. Its weights have norm of order one, so the preactivations $z_j$ are of order one, and the gradient step $\frac\eta d\delta_jx$ has norm of order $\eta d^{-1/2}$, which is the scaling at which each step changes the overlaps by $O(1/d)$. Time is measured in samples per dimension, $t=s/d$.

Rotational invariance does the first reduction. Given the current weights, the vector $(z_1,\dots,z_m,z_{\star1},\dots,z_{\star m_\star})$ is a centered Gaussian with covariance
\begin{equation}\label{eq:committee:M}
  \committeeM=\begin{pmatrix}Q&R\\R^\top&T\end{pmatrix},
\end{equation}
so the population loss, and the mean of every update of the overlaps, depend on the $d$-dimensional weights only through $Q$ and $R$. The order parameters are to this model what the single overlap $r$ was to the single-index model of \cref{ch:singleindex}: the reason the high-dimensional dynamics collapses to a few numbers. The difference is that now there are several of them, and permutations of the student units act on them.

What was thrown away is, first, every layer but one: the readout weights are fixed, so the student can only move its hidden units. Inputs are isotropic Gaussian, which makes the overlaps sufficient statistics. The width is a fixed number while $d\to\infty$, which is the opposite of the infinite-width limits of \cref{ch:signalprop,ch:lazy}. The labels are noiseless and realizable when $m\ge m_\star$, and each sample is used once. Each simplification isolates the mechanism of interest, the interplay between the permutation symmetry and the gradient noise, and the activation $g$ is chosen because every Gaussian integral the dynamics needs is then elementary (\cref{lem:committee:erf}). The target has information exponent $1$ in each direction, since $\hat g_1=\E[zg(z)]=\E g'(z)=1/\sqrt\pi\ne0$, so no unit has to search; whatever delay we find is due to symmetry, not to the flatness studied in \cref{ch:singleindex}.

To invent the model one had to notice that two facts compose. Rotational invariance reduces the state to the Gram matrix of the student and teacher weights; and the gradient noise of online learning in high dimension averages out over $d$ steps, except for a deterministic contribution of order $\eta^2$ that a naive gradient-flow limit misses. Together they turn a stochastic process in $\R^{md}$ into an autonomous ODE in $m(m+1)/2+mm_\star$ variables.

\begin{dictionary}
student weights $w_j$ & incoming weights of hidden unit $j$\\
teacher weights $w_{\star k}$ & the features the target is built from\\
fixed readout $f=\sum_jg(z_j)$ & second layer, frozen at equal weights\\
overlap $R_{jk}$ & how much unit $j$ has learned feature $k$\\
overlap $Q_{jl}$, $j\ne l$ & redundancy between units $j$ and $l$\\
symmetric fixed point & a plateau: units are copies, each a blend of all features\\
specialization & the drop that ends the plateau: each unit takes one feature\\
unstable eigenvalue $\committeelu$ & the rate at which redundancy breaks\\
asymmetry of order $d^{-1/2}$ & the accidental preferences of a random initialization\\
$\eta^2$ term & gradient noise, which limits the step size\\
time $t=s/d$ & samples per input dimension\\
\end{dictionary}

\section{The solution: Gaussian integrals and the ODE limit}\label{sec:committee:ode}

Three lemmas prepare the main theorem. The first is the computational engine: a Gaussian weight $e^{-z^2/2}$, which is what $g'$ contributes, can be absorbed into the covariance.

\begin{lemma}[Gaussian tilting]\label[lemma]{lem:committee:tilt}
Let $X\sim\cN(0,C)$ in $\R^n$ with $C$ positive semidefinite, let $S=(s_1,\dots,s_r)$ be a list of indices in $\{1,\dots,n\}$, repetitions allowed, and write $X_S=(X_{s_1},\dots,X_{s_r})$, $C_{SS}=(C_{s_is_{i'}})$, and $C_{\cdot S}$ for the $n\times r$ matrix with columns $C_{\cdot s_i}$. For every measurable $H$ of polynomial growth,
\[
  \E\bigl[e^{-\norm{X_S}^2/2}H(X)\bigr]=\det(I_r+C_{SS})^{-1/2}\,\E H(\tilde X),
\]
where $\tilde X\sim\cN(0,\tilde C)$ with $\tilde C=C-C_{\cdot S}(I_r+C_{SS})^{-1}C_{\cdot S}^\top$.
\end{lemma}

\begin{proof}
Write $C=LL^\top$ with $L\in\R^{n\times n}$ and $X=LZ$ with $Z\sim\cN(0,I_n)$, and let $L_S$ be the $r\times n$ matrix whose rows are the rows $s_1,\dots,s_r$ of $L$, so that $X_S=L_SZ$ and $C_{SS}=L_SL_S^\top$. The weight times the standard Gaussian density is
\[
  e^{-\norm{L_Sz}^2/2}(2\pi)^{-n/2}e^{-\norm z^2/2}=(2\pi)^{-n/2}\exp\bigl(-\tfrac12z^\top(I_n+L_S^\top L_S)z\bigr),
\]
an unnormalized Gaussian density with covariance $\Sigma=(I_n+L_S^\top L_S)^{-1}$ and total mass $\det(I_n+L_S^\top L_S)^{-1/2}=\det(I_r+L_SL_S^\top)^{-1/2}$, by Sylvester's determinant identity. Hence the left side equals $\det(I_r+C_{SS})^{-1/2}\E H(L\tilde Z)$ with $\tilde Z\sim\cN(0,\Sigma)$. The Woodbury identity $\Sigma=I_n-L_S^\top(I_r+L_SL_S^\top)^{-1}L_S$ and $LL_S^\top=C_{\cdot S}$ give $\operatorname{Cov}(L\tilde Z)=L\Sigma L^\top=C-C_{\cdot S}(I_r+C_{SS})^{-1}C_{\cdot S}^\top$.
\end{proof}

The covariance of $\tilde X$ is a Schur complement of $I+C$ minus the identity, hence positive semidefinite. The second lemma says how a Gaussian expectation changes with the covariance; it contains Price's theorem (\cref{lem:toolkit:price}) for variables that are not standardized, and it will also supply the Lipschitz bounds that the ODE method needs.

\begin{lemma}[Gaussian interpolation]\label[lemma]{lem:committee:interp}
Let $\psi\colon\R^n\to\R$ be twice continuously differentiable with $\abs\psi$, $\abs{\partial_a\psi}$, and $\abs{\partial_a\partial_b\psi}$ bounded by $\kappa(1+\norm\xi^p)$. For positive semidefinite $C$ write $\E_C\psi=\E\psi(X)$ with $X\sim\cN(0,C)$. For positive semidefinite $C,C'$, with $C_\theta=(1-\theta)C+\theta C'$,
\[
  \E_{C'}\psi-\E_C\psi=\frac12\int_0^1\sum_{a,b=1}^n(C'-C)_{ab}\,\E_{C_\theta}[\partial_a\partial_b\psi]\dd\theta .
\]
Consequently, if $\norm C_F,\norm{C'}_F\le B$, then $\abs{\E_{C'}\psi-\E_C\psi}\le L_B\norm{C'-C}_F$ with $L_B=\frac12n\kappa(1+\mu_pB^{p/2})$, where $\mu_p=\E\norm{\cN(0,I_n)}^p$.
\end{lemma}

\begin{proof}
The plan is to join the two laws by a path of Gaussians and to differentiate along it, converting the result with Stein's lemma. Let $X\sim\cN(0,C)$ and $X'\sim\cN(0,C')$ be independent and $X_\theta=\sqrt{1-\theta}\,X+\sqrt\theta\,X'$, which has law $\cN(0,C_\theta)$; put $\varphi(\theta)=\E\psi(X_\theta)$. For $0<\theta<1$ the growth bounds justify differentiating under the expectation:
\[
  \varphi'(\theta)=\sum_a\E\Bigl[\partial_a\psi(X_\theta)\Bigl(\frac{X'_a}{2\sqrt\theta}-\frac{X_a}{2\sqrt{1-\theta}}\Bigr)\Bigr].
\]
Condition on $X$ and apply Stein's lemma (\cref{thm:toolkit:stein}) to $X'$ and the function $h(X')=\partial_a\psi(\sqrt{1-\theta}X+\sqrt\theta X')$, whose derivatives are $\partial_bh=\sqrt\theta\,\partial_b\partial_a\psi$: this gives $\E[X'_ah(X')]=\sqrt\theta\sum_bC'_{ab}\E[\partial_a\partial_b\psi(X_\theta)]$. The same computation for $X$ gives the second term, and
\[
  \varphi'(\theta)=\frac12\sum_{a,b}(C'_{ab}-C_{ab})\,\E[\partial_a\partial_b\psi(X_\theta)].
\]
The right side is bounded on $(0,1)$, and $\varphi$ is continuous on $[0,1]$ by dominated convergence, so integrating proves the identity. For the bound, $\norm{C_\theta}_F\le B$ gives $\E\norm{X_\theta}^p\le\norm{C_\theta}_{\op}^{p/2}\mu_p\le B^{p/2}\mu_p$, so every $\abs{\E_{C_\theta}\partial_a\partial_b\psi}$ is at most $\kappa(1+\mu_pB^{p/2})$, and the Cauchy--Schwarz inequality over the $n^2$ pairs $(a,b)$ finishes the proof.
\end{proof}

Taking $C'$ near $C$ shows that $C\mapsto\E_C\psi$ is continuously differentiable on the cone of positive semidefinite matrices, with partial derivative $\E_C[\partial_a\partial_b\psi]$ with respect to an off-diagonal pair $C_{ab}=C_{ba}$ and $\frac12\E_C[\partial_a^2\psi]$ with respect to a diagonal entry. The third lemma evaluates every integral the committee machine needs.\sidenote{These are the integrals $I_2$, $I_3$, and $I_4$ of \textcite{SaadSolla1995b}; the matrix form of $I_4$ through a Schur complement is our arrangement of their formula.}

\begin{lemma}[The $\erf$ integrals]\label[lemma]{lem:committee:erf}
Let $X\sim\cN(0,C)$ with $C$ positive semidefinite, and let $a,b,c,c'$ be indices, not necessarily distinct. Then
\begin{align*}
  I_2(a,b)&\defeq\E[g(X_a)g(X_b)]=\frac2\pi\arcsin\frac{C_{ab}}{\sqrt{(1+C_{aa})(1+C_{bb})}},\\
  I_3(a,b,c)&\defeq\E[g'(X_a)\,X_b\,g(X_c)]\\
  &=\frac2\pi\,\frac{C_{bc}(1+C_{aa})-C_{ab}C_{ac}}{(1+C_{aa})\sqrt{(1+C_{aa})(1+C_{cc})-C_{ac}^2}},\\
  I_4(a,b;c,c')&\defeq\E[g'(X_a)g'(X_b)g(X_c)g(X_{c'})]\\
  &=\frac{4}{\pi^2}\,\frac{1}{\sqrt{\det(I_2+C_{SS})}}\arcsin\frac{\tilde C_{cc'}}{\sqrt{(1+\tilde C_{cc})(1+\tilde C_{c'c'})}},
\end{align*}
where $S=(a,b)$ and $\tilde C=C-C_{\cdot S}(I_2+C_{SS})^{-1}C_{\cdot S}^\top$. All denominators are at least $1$, and all arguments of $\arcsin$ lie in $(-1,1)$.
\end{lemma}

\begin{proof}
For $I_2$, only the covariance of the pair $(X_a,X_b)$ matters; write $\alpha=C_{aa}$, $\beta=C_{bb}$, and $\gamma=C_{ab}$, so that $\abs\gamma\le\sqrt{\alpha\beta}$. Move the off-diagonal entry from $0$ to $\gamma$ with the diagonal fixed; every matrix on the way is positive semidefinite. By \cref{lem:committee:interp} with $\psi(u,v)=g(u)g(v)$, whose derivatives are bounded, the derivative with respect to the off-diagonal entry $c$ is $\E[g'(X_a)g'(X_b)]$, which is $\frac2\pi\E[e^{-(X_a^2+X_b^2)/2}]=\frac2\pi((1+\alpha)(1+\beta)-c^2)^{-1/2}$ by \cref{lem:committee:tilt} with $H=1$. At $c=0$ the two variables are independent and $g$ is odd, so the expectation vanishes. Integrating,
\[
  I_2(a,b)=\frac2\pi\int_0^\gamma\frac{\dd c}{\sqrt{(1+\alpha)(1+\beta)-c^2}}=\frac2\pi\arcsin\frac{\gamma}{\sqrt{(1+\alpha)(1+\beta)}} .
\]
For $I_3$, apply \cref{lem:committee:tilt} with $S=(a)$: $I_3=\sqrt{2/\pi}\,(1+C_{aa})^{-1/2}\E[\tilde X_bg(\tilde X_c)]$. Stein's lemma (\cref{thm:toolkit:stein}) gives $\E[\tilde X_bg(\tilde X_c)]=\tilde C_{bc}\E g'(\tilde X_c)$, and \cref{lem:committee:tilt} with $H=1$ gives $\E g'(\tilde X_c)=\sqrt{2/\pi}\,(1+\tilde C_{cc})^{-1/2}$. Substituting $\tilde C_{bc}=C_{bc}-C_{ab}C_{ac}/(1+C_{aa})$ and $1+\tilde C_{cc}=((1+C_{aa})(1+C_{cc})-C_{ac}^2)/(1+C_{aa})$ gives the formula. For $I_4$, apply \cref{lem:committee:tilt} with $S=(a,b)$ and then the formula for $I_2$ to the Gaussian vector $\tilde X$.

For the last sentence: $(1+C_{aa})(1+C_{cc})-C_{ac}^2\ge1+C_{aa}+C_{cc}\ge1$ because $C_{ac}^2\le C_{aa}C_{cc}$; $\det(I_2+C_{SS})\ge1$ because $C_{SS}$ is positive semidefinite; and for a positive semidefinite matrix $D$, $\abs{D_{ab}}\le\sqrt{D_{aa}D_{bb}}<\sqrt{(1+D_{aa})(1+D_{bb})}$, which applies to $C$ and to $\tilde C$.
\end{proof}

With the lemma, the population loss is explicit. Expanding the square and using $I_2$ with the covariance \eqref{eq:committee:M},
\begin{equation}\label{eq:committee:loss}
\begin{aligned}
  \pi\Loss={}&\sum_{k,k'}\arcsin\frac{T_{kk'}}{\sqrt{(1+T_{kk})(1+T_{k'k'})}}+\sum_{j,l}\arcsin\frac{Q_{jl}}{\sqrt{(1+Q_{jj})(1+Q_{ll})}}\\
  &-2\sum_{j,k}\arcsin\frac{R_{jk}}{\sqrt{(1+Q_{jj})(1+T_{kk})}}.
\end{aligned}
\end{equation}
For an orthonormal teacher, $T=I_{m_\star}$, the loss at $W=0$ is $\frac1\pi m_\star\arcsin\frac12=m_\star/6$.

\subsection{The order parameters follow an ODE}\label{sec:committee:odelimit}

Collect the order parameters in a vector $\omega=((Q_{jl})_{j\le l},(R_{jk}))\in\R^D$ with $D=m(m+1)/2+mm_\star$, and write $\committeeM(\omega)$ for the matrix \eqref{eq:committee:M}, with the teacher Gram matrix $T$ fixed. Call $\omega$ \term{realizable} if $\committeeM(\omega)$ is positive semidefinite; these are exactly the order parameters of some weights, in every dimension $d\ge m+m_\star$. Taking the inner product of \eqref{eq:committee:sgd} with $w_{\star k}$ and with $w_l$ gives the exact recursion
\begin{equation}\label{eq:committee:recursion}
\begin{aligned}
  R_{jk}^{s+1}&=R_{jk}^s+\frac\eta d\,\delta_j^sz_{\star k}^s,\\
  Q_{jl}^{s+1}&=Q_{jl}^s+\frac\eta d\bigl(\delta_j^sz_l^s+\delta_l^sz_j^s\bigr)+\frac{\eta^2}{d^2}\,\delta_j^s\delta_l^s\,\norm{x^s}^2 .
\end{aligned}
\end{equation}
The last term is the one that a gradient-flow limit would drop. It carries a factor $\norm{x^s}^2/d^2\approx1/d$, the same order as the others, because the step moves the weights along $x^s$, a vector of length $\sqrt d$. For realizable $\omega$ define $F_\eta(\omega)\in\R^D$ by Gaussian expectations under $(z,z_\star)\sim\cN(0,\committeeM(\omega))$, with $\delta_j=g'(z_j)(\sum_kg(z_{\star k})-\sum_lg(z_l))$:
\begin{equation}\label{eq:committee:field}
\begin{aligned}
  F_\eta(\omega)_{R,jk}&=\eta\,\E[\delta_jz_{\star k}],\\
  F_\eta(\omega)_{Q,jl}&=\eta\,\E[\delta_jz_l+\delta_lz_j]+\eta^2\,\E[\delta_j\delta_l].
\end{aligned}
\end{equation}
Expanding $\delta_j$ and applying \cref{lem:committee:erf} with the covariance $\committeeM(\omega)$ makes the field explicit:
\begin{align*}
  \E[\delta_jz_{\star k}]&=\sum_{k'}I_3(z_j,z_{\star k},z_{\star k'})-\sum_{l'}I_3(z_j,z_{\star k},z_{l'}),\\
  \E[\delta_jz_l]&=\sum_{k'}I_3(z_j,z_l,z_{\star k'})-\sum_{l'}I_3(z_j,z_l,z_{l'}),\\
  \E[\delta_j\delta_l]&=\sum_{k,k'}I_4(z_j,z_l;z_{\star k},z_{\star k'})-2\sum_{k,l'}I_4(z_j,z_l;z_{\star k},z_{l'})\\
  &\quad+\sum_{l',l''}I_4(z_j,z_l;z_{l'},z_{l''}),
\end{align*}
where each $I$ is evaluated by reading the entries of its arguments' covariance from $Q$, $R$, and $T$. This is the \term{Saad--Solla system}\index{Saad--Solla equations}. To apply the ODE method we extend $F_\eta$ to all of $\R^D$: let $\Pi$ be the Frobenius-norm projection onto the closed convex cone of positive semidefinite matrices of size $m+m_\star$, and put $\bar F_\eta(\omega)=F_\eta(\Pi\committeeM(\omega))$, where $F_\eta(\committeeM)$ means the expectations \eqref{eq:committee:field} under $\cN(0,\committeeM)$. On realizable points $\bar F_\eta=F_\eta$.

\begin{theorem}[The ODE limit]\label{thm:committee:ode}
Fix $m$, $m_\star$, $\eta>0$, a positive semidefinite teacher Gram matrix $T$, a realizable starting point $\committeeOm_0$, and a horizon $t_1>0$. The equation $\dot{\committeeOm}=\bar F_\eta(\committeeOm)$, $\committeeOm(0)=\committeeOm_0$, has a unique solution on $[0,t_1]$, the solution is realizable for every $t$, and so it solves $\dot{\committeeOm}=F_\eta(\committeeOm)$ with the explicit field above. There is a constant $C$, depending only on $m,m_\star,\eta,T,\committeeOm_0,t_1$, such that the following holds. Let $d\ge m+m_\star$, let the teacher have Gram matrix $T$, run \eqref{eq:committee:sgd} from deterministic weights with order parameters $\omega_0$, and let $\omega_s$ be the order parameters after $s$ steps. For every $a\in(0,1]$ with $\Delta\defeq C(\norm{\omega_0-\committeeOm_0}+a+d^{-1/2})<1$,
\[
  \Pr\Bigl(\max_{0\le s\le t_1d}\norm{\omega_s-\committeeOm(s/d)}>\Delta\Bigr)\le\frac{C}{a^2d}.
\]
\end{theorem}

With $\omega_0\to\committeeOm_0$ and $a=d^{-1/4}$, the order parameters stay within $O(\norm{\omega_0-\committeeOm_0}+d^{-1/4})$ of the ODE over $t_1d$ steps, with probability $1-O(d^{-1/2})$. A random initialization $w_j^0\sim\cN(0,\frac{q_0}dI_d)$ has $\norm{\omega_0-\committeeOm_0}=O(d^{-1/2})$ with high probability for $\committeeOm_0=(q_0I_m,0)$, and the theorem applies conditionally on the initialization.

\begin{proof}
The plan is to verify the hypotheses of the ODE method (\cref{thm:toolkit:ode-method}): a field that is Lipschitz and bounded on a neighborhood of the path, increments whose noise has bounded conditional variance, and a small bias.

\emph{Bounds on the integrands.} Since $\abs g<1$ and $0<g'\le\sqrt{2/\pi}$, every $\delta_j$ satisfies $\abs{\delta_j}\le\bar\delta\defeq\sqrt{2/\pi}\,(m+m_\star)$. The integrands $\delta_jz_{\star k}$, $\delta_jz_l$, and $\delta_j\delta_l$ are smooth functions of $(z,z_\star)$ whose derivatives up to order two are bounded by $\kappa(1+\norm{(z,z_\star)})$ for a constant $\kappa$ depending on $m$ and $m_\star$, by \eqref{eq:committee:gprime}.

\emph{The field is Lipschitz on bounded sets and grows sublinearly.} The map $\omega\mapsto\committeeM(\omega)$ is affine with $\norm{\committeeM(\omega)-\committeeM(\omega')}_F\le\sqrt2\norm{\omega-\omega'}$, and $\Pi$ is nonexpansive with $\Pi0=0$, so $\norm{\Pi\committeeM(\omega)}_F\le\norm{\committeeM(\omega)}_F$. By \cref{lem:committee:interp}, $\bar F_\eta$ is Lipschitz on every ball. By the Cauchy--Schwarz inequality, $\abs{\E[\delta_jz_l]}\le\bar\delta\,(\Pi\committeeM)_{ll}^{1/2}\le\bar\delta\norm{\committeeM(\omega)}_F^{1/2}$, and similarly for the other entries, so $\norm{\bar F_\eta(\omega)}\le c_0(1+\norm\omega)$. A locally Lipschitz field with linear growth has a unique solution on $[0,t_1]$; call it $\committeeOm$, let $r_1=\max_t\norm{\committeeOm(t)}$, and let $U$ be the open ball of radius $r_1+2$. On $U$ the field is bounded by some $B$ and $L$-Lipschitz, and every point within distance $\rho=1$ of the path lies in $U$.

\emph{Noise and bias.} Let $\cF_s$ be generated by $x^0,\dots,x^{s-1}$. Given $\cF_s$, the vector $(z^s,z_\star^s)$ is $\cN(0,\committeeM(\omega_s))$, and $\committeeM(\omega_s)$ is positive semidefinite, so $\Pi\committeeM(\omega_s)=\committeeM(\omega_s)$. Write $d(\omega_{s+1}-\omega_s)=\bar F_\eta(\omega_s)+\xi_{s+1}+\epsilon_{s+1}$ with $\epsilon_{s+1}=\E[d(\omega_{s+1}-\omega_s)\mid\cF_s]-\bar F_\eta(\omega_s)$ and $\xi_{s+1}$ the centered remainder. By \eqref{eq:committee:recursion}, the $R$-entries of $\epsilon_{s+1}$ vanish and its $Q$-entries are $\eta^2\E[\delta_j\delta_l(\norm{x^s}^2/d-1)\mid\cF_s]$, which are at most $\eta^2\bar\delta^2\E\abs{\norm x^2/d-1}\le\eta^2\bar\delta^2\sqrt{2/d}$ in absolute value, because $\Var\norm x^2=2d$. So $\norm{\epsilon_{s+1}}\le\epsilon\defeq\sqrt D\,\eta^2\bar\delta^2\sqrt{2/d}$. The conditional second moment of $\xi_{s+1}$ is at most that of $d(\omega_{s+1}-\omega_s)$, whose entries are bounded in mean square by $\eta^2\bar\delta^2T_{kk}$, by $3\eta^2\bar\delta^2(Q_{jj}+Q_{ll})+3\eta^4\bar\delta^4\E\norm x^4/d^2$, and $\E\norm x^4=d^2+2d\le3d^2$. On the event $\{\omega_s\in U\}$ these are bounded by a constant $\varsigma^2$.

\emph{Conclusion.} \Cref{thm:toolkit:ode-method} applies with these constants: with $\Delta'=(\norm{\omega_0-\committeeOm_0}+a+t_1\epsilon+t_1LB/(2d))e^{Lt_1}$, if $\Delta'<1$, the probability that $\norm{\omega_s-\committeeOm(s/d)}>\Delta'$ for some $s\le t_1d$ is at most $\varsigma^2t_1/(a^2d)$. Since $\epsilon$ and $1/d$ are at most constants times $d^{-1/2}$, $\Delta'\le\Delta$ for a suitable $C$, which proves the bound.

\emph{Realizability.} Fix $t\in[0,t_1]$. For each $d\ge m+m_\star$ choose initial weights with $\omega_0=\committeeOm_0$ exactly, which is possible because $\committeeOm_0$ is realizable, and take $a=d^{-1/4}$. For large $d$ the bound shows that with positive probability $\norm{\omega_{\floor{td}}-\committeeOm(\floor{td}/d)}\le C(2d^{-1/4})$, so some realizable point lies within $2Cd^{-1/4}+Bd^{-1}$ of $\committeeOm(t)$. The realizable set is closed, so $\committeeOm(t)$ is realizable, and there $\bar F_\eta=F_\eta$.
\end{proof}

\begin{remark}[The overlaps form a Markov chain]\label[remark]{rem:committee:chain}
The proof shows more. Given $\cF_s$, the update \eqref{eq:committee:recursion} is a function of $\omega_s$ and of the triple $(z^s,z_\star^s,\norm{x^s}^2)$, whose conditional law depends only on $\omega_s$: if $P$ is the orthogonal projection onto the span of the current student and teacher weights, of dimension $\varrho$, then $\norm{x}^2=\norm{Px}^2+\chi^2_{d-\varrho}$ with an independent chi-square variable, and $\norm{Px}^2$ is a function of $(z,z_\star)$. So $(\omega_s)$ is a Markov chain whose law can be simulated exactly in every dimension at a cost independent of $d$ per step. All simulations in this chapter use it; for $m=m_\star=2$, $d=100$, and $400$ runs, its median loss agrees with that of SGD in $\R^{100}$ within $1\%$ at $t=10,20,40,80$ (\cref{ex:committee:chain}).
\end{remark}

\begin{remark}[Two time scales]\label[remark]{rem:committee:tau}
The $\eta$-term of $F_\eta$ is the image of gradient flow on the population loss: $-\nabla_{w_j}\Loss(W)=\E[\delta_jx]$, and its inner products with $w_{\star k}$ and $w_l$ are the entries of $F_\eta/\eta$ without the $\eta^2$ term. So for small $\eta$, in the time $\tau=\eta t$, the order parameters follow
\begin{equation}\label{eq:committee:G}
  \frac{\dd\omega}{\dd\tau}=G(\omega),\qquad G(\omega)_{R,jk}=\E[\delta_jz_{\star k}],\quad G(\omega)_{Q,jl}=\E[\delta_jz_l+\delta_lz_j],
\end{equation}
the order-parameter shadow of $\dot w_j=-\nabla_{w_j}\Loss(W)$, and the $\eta^2$ term is a correction of relative size $\eta$. Nothing below requires $\eta$ to be small except where we say so.
\end{remark}

\subsection{One unit: the learning-rate threshold}\label{sec:committee:single}

The simplest case already shows what the $\eta^2$ term does. Let $m=m_\star=1$ and $T=1$, write $R=R_{11}$ and $Q=Q_{11}$, and use the coordinates $u=R-1$ and $v=Q-R^2$, the squared distance of $w$ from the line through $w_\star$; the realizable set is $v\ge0$, and $\norm{w-w_\star}^2=u^2+v$. Put $c_1=\E[g'(z)^2z^2]$, $c_2=\E[g'(z)^2]$, and $c_3=\E[g'(z)^4]$ for $z\sim\cN(0,1)$, so that
\[
  c_1=\frac{2}{3\sqrt3\,\pi},\qquad c_2=\frac{2}{\sqrt3\,\pi},\qquad c_3=\frac{4}{\sqrt5\,\pi^2};
\]
these values follow from $g'(z)^2=\frac2\pi e^{-z^2}$, $\E e^{-z^2}=3^{-1/2}$, $\E z^2e^{-z^2}=3^{-3/2}$, and $\E e^{-2z^2}=5^{-1/2}$.

\begin{proposition}[A single unit and its learning-rate threshold]\label[proposition]{prop:committee:single}
For every $\eta>0$ the point $w=w_\star$, that is $(u,v)=(0,0)$, is a fixed point of the ODE of \cref{thm:committee:ode}. In the coordinates $(u,v)$ the Jacobian of the field there is upper triangular, with diagonal entries $-\eta c_1$ and $-(2\eta c_2-\eta^2c_3)$. Consequently, with
\[
  \eta_{\max}=\frac{2c_2}{c_3}=\pi\sqrt{5/3}\approx4.056,
\]
the fixed point is linearly stable for $0<\eta<\eta_{\max}$: solutions starting near it converge exponentially, and the loss decays at the rate $\min(2\eta c_1,\,2\eta c_2-\eta^2c_3)$, which is largest, $8\sqrt5/27\approx0.663$, at $\eta=\frac23\eta_{\max}$. For $\eta>\eta_{\max}$ the fixed point is unstable, and an eigenvector of the positive eigenvalue points into the realizable region $v>0$.
\end{proposition}

\begin{proof}
The plan is to write the ODE in the coordinates $(u,v)$, to compute its derivatives at $(0,0)$ with the interpolation lemma, and to read off the eigenvalues. Represent $z=Rz_\star+\sqrt v\,\zeta$ with $z_\star,\zeta$ independent standard Gaussians, and let $\delta=g'(z)(g(z_\star)-g(z))$. At $w=w_\star$ we have $z=z_\star$ and $\delta=0$ identically, so every term of \eqref{eq:committee:field} vanishes: the point is fixed for every $\eta$.

\emph{The equation for $v$.} By the chain rule and \eqref{eq:committee:field},
\[
  \dot v=\dot Q-2R\dot R=2\eta\,\E[\delta(z-Rz_\star)]+\eta^2\E[\delta^2]=2\eta\sqrt v\,\E[\delta\zeta]+\eta^2\E[\delta^2].
\]
Stein's lemma in $\zeta$ gives $\E[\delta\zeta]=\sqrt v\,\E[\partial_z\delta]$, so $\dot v=2\eta v\,A+\eta^2B$ with $A=\E[\partial_z\delta]$ and $B=\E[\delta^2]$, both continuously differentiable functions of $(R,Q)$ by \cref{lem:committee:interp}. At the fixed point, $\partial_z\delta=g''(z)(g(z_\star)-g(z))-g'(z)^2=-g'(z_\star)^2$, so $A=-c_2$. For $B$, use the derivative rules after \cref{lem:committee:interp}, with $Q$ the variance of $z$ and $R$ its covariance with $z_\star$. At $z=z_\star$, $\delta=0$ and $\partial_z\delta=-\partial_{z_\star}\delta=-g'^2$, so $\partial_Q B=\frac12\E[\partial_z^2\delta^2]=\E[(\partial_z\delta)^2]=c_3$ and $\partial_RB=\E[\partial_z\partial_{z_\star}\delta^2]=2\E[\partial_z\delta\,\partial_{z_\star}\delta]=-2c_3$. Since $\dd Q=\dd v+2R\,\dd R$ and $\dd R=\dd u$, at $R=1$ this reads $\dd B=c_3\,\dd v$. Hence $\partial_u\dot v=0$ and $\partial_v\dot v=2\eta A+\eta^2c_3=-(2\eta c_2-\eta^2c_3)$ at the fixed point.

\emph{The equation for $u$.} Here $\dot u=\eta\phi(R,Q)$ with $\phi=\E[\delta z_\star]$. At the fixed point, writing $z_\star=z$ after differentiating,
\begin{align*}
  \partial_R\phi&=\E[\partial_z\partial_{z_\star}(\delta z_\star)]=\E[-g'^2-z_\star^2g'^2]=-c_2-c_1,\\
  \partial_Q\phi&=\tfrac12\E[z_\star\partial_z^2\delta]=\tfrac12\E[-3z_\star g'g'']=\tfrac32c_1,
\end{align*}
using $\partial_z^2\delta=-3g'g''$ at $z=z_\star$ and $g''=-zg'$. Moving along $u$ at fixed $v$ changes $Q$ by $2\,\dd u$, so $\partial_u\dot u=\eta(\partial_R\phi+2\partial_Q\phi)=\eta(2c_1-c_2)=-\eta c_1$, because $c_2=3c_1$.

\emph{Conclusion.} The Jacobian is $\begin{psmallmatrix}-\eta c_1&*\\0&-(2\eta c_2-\eta^2c_3)\end{psmallmatrix}$. For $\eta<\eta_{\max}$ both eigenvalues are negative and the linearization theorem gives local exponential convergence. The loss is $\frac12\E[g'(z_\star)^2(z-z_\star)^2]$ to second order, that is $\frac12(c_1u^2+c_2v)$, and $u^2$ decays at rate $2\eta c_1$ while $v$, forced by terms of order $u^2$, decays at rate $\min(2\eta c_1,2\eta c_2-\eta^2c_3)$; so does the loss. The two rates are equal at $\eta=2(c_2-c_1)/c_3=\frac23\eta_{\max}$, where the first is increasing and the second decreasing in $\eta$, which gives the maximum $\frac43\eta_{\max}c_1=8\sqrt5/27$. For $\eta>\eta_{\max}$ the eigenvalue $-(2\eta c_2-\eta^2c_3)$ is positive, and its eigenvector in a triangular matrix has a nonzero second coordinate, which can be taken positive.
\end{proof}

The threshold is a stability condition of the kind met in \cref{ch:prologue} and \cref{ch:stability}: gradient descent on a quadratic with curvature $\lambda_1$ diverges for steps above $2/\lambda_1$. Here the curvature in the directions orthogonal to $w_\star$ is $c_2$ per unit time, and the noise term $\eta^2c_3v$ plays the role of the overshoot. The variance of the gradient, not the curvature alone, sets the threshold. Beyond it the ODE converges to a different fixed point with positive error: numerically, $\Loss\approx0.0056$ at $\eta=4.2$ and $0.036$ at $\eta=5$, with $Q>1$; the student overshoots in norm.

\section{The solution: the symmetric plateau}\label{sec:committee:plateau}

Now take $K\ge2$ students and $K$ orthonormal teachers: $m=m_\star=K$ and $T=I_K$. The permutations of the students, and separately of the teachers, leave the dynamics unchanged, and they single out two linear subspaces of order parameters,
\begin{align*}
  \cS&=\{\omega:\ R_{jk}=R\text{ for all }j,k;\ Q_{jj}=Q\text{ for all }j;\ Q_{jl}=C\text{ for all }j\ne l\},\\
  \cS_0&=\cS\cap\{C=Q\}.
\end{align*}
Points of $\cS_0$ are the order parameters of identical students, $w_1=\dots=w_K$, since then $\norm{w_j-w_l}^2=2Q-2C=0$.

\begin{theorem}[The symmetric plateau]\label{thm:committee:plateau}
Let $m=m_\star=K\ge2$ and $T=I_K$.
\begin{enumerate}
\item[\textup{(a)}] For every $\eta>0$, solutions of the ODE of \cref{thm:committee:ode} that start in $\cS$ stay in $\cS$, and solutions that start in $\cS_0$ stay in $\cS_0$. For SGD itself, students that start identical stay identical.
\item[\textup{(b)}] The point $P_\star\in\cS_0$ with
\[
  R=a\defeq\frac{1}{\sqrt{K(2K-1)}},\qquad Q=C=\frac1{2K-1},
\]
the order parameters of $w_j=a\sum_kw_{\star k}$ for every $j$, is a fixed point of the gradient-flow field $G$ of \eqref{eq:committee:G}, and the only one of the form $w_j=a'\sum_kw_{\star k}$ for all $j$ with $a'>0$. Its loss is
\[
  \Loss_\star=\frac K6-\frac{K^2}\pi\arcsin\frac1{2K}.
\]
\item[\textup{(c)}] Let $\Xi$ be the direction with $\Xi_{R,11}=\Xi_{R,22}=1$, $\Xi_{R,12}=\Xi_{R,21}=-1$, and all other entries $0$. Then $\Xi$ is an eigenvector of the Jacobian of $G$ at $P_\star$ with eigenvalue
\[
  \committeelu=\frac{2K}{\pi\,(2K+1)^{3/2}(2K-1)^{1/2}}>0 .
\]
So $P_\star$ is a saddle, and its unstable direction makes student $1$ prefer teacher $1$ and student $2$ prefer teacher $2$. For $K=2$, $\committeelu=4\sqrt{15}/(75\pi)\approx0.0657$ and $\Loss_\star=\frac13-\frac4\pi\arcsin\frac14\approx0.01161$.
\item[\textup{(d)}] Perturbations of $P_\star$ in which the students split into directions orthogonal to all teachers, $w_j=w+\varepsilon u_j$ with $\sum_ju_j=0$ and each $u_j$ orthogonal to $w_{\star1},\dots,w_{\star K}$, have velocity $O(\varepsilon^2)$ under the gradient flow: these directions are marginal.
\end{enumerate}
\end{theorem}

The plateau is the time spent near $P_\star$: the dynamics is drawn into $\cS_0$ toward $P_\star$ along stable directions and leaves it along $\Xi$. In $\cS_0$ itself and for $K=2$, the two eigenvalues at $P_\star$ are $-0.592$ and $-1.972$ (computed numerically), so $P_\star$ attracts identical students exponentially. The full $7\times7$ Jacobian of $G$ at $P_\star$ for $K=2$ has the eigenvalues
\[
  0.0657,\quad0,\quad0,\quad-0.592,\quad-0.920,\quad-0.986,\quad-1.972,
\]
where the zeros belong to directions that are of second order in the weights, among them the marginal directions of (d), and $-0.920$ to identical students turning toward one teacher (\cref{ex:committee:modes}). Within $\cS$ the fixed point therefore attracts only algebraically along the coordinate $Q-C$; numerically $Q-C$ decays like $4/\tau$ along the flow.

\begin{proof}
The plan is to prove (a) by checking that the field is tangent to the two subspaces, (b) with the integrals of \cref{lem:committee:erf}, and (c) and (d) in weight space, where the gradient flow is a gradient flow and the Hessian of the loss can be computed with Stein's lemma.

\emph{(a) Invariance.} We use one fact about ODEs: if a field $F$ is Lipschitz on bounded sets and $F(V)\subseteq V$ for a linear subspace $V$, then solutions starting in $V$ stay in $V$, because the ODE restricted to $V$ has a solution, which solves the full equation, and solutions are unique. For $\omega\in\cS$, the covariance $\committeeM(\omega)$ is unchanged when the students are permuted, and also when the teachers are permuted, so the joint law of $(z,z_\star)$ is invariant under both kinds of relabeling, while the label and the prediction are symmetric functions. Hence $\E[\delta_jz_{\star k}]$ does not depend on $(j,k)$, $\E[\delta_jz_j]$ and $\E[\delta_j^2]$ do not depend on $j$, and $\E[\delta_jz_l+\delta_lz_j]$ and $\E[\delta_j\delta_l]$ for $j\ne l$ do not depend on the pair. So $F_\eta(\omega)\in\cS$, and the same holds for $\bar F_\eta$, because the projection $\Pi$ commutes with simultaneous permutations of rows and columns. For $\omega\in\cS_0$ the variables $z_1,\dots,z_K$ coincide almost surely, because $\E(z_j-z_l)^2=2Q-2C=0$; so are the $\delta_j$, and $F_\eta(\omega)_{Q,jl}$ does not depend on whether $j=l$: $F_\eta(\omega)\in\cS_0$. For SGD, \eqref{eq:committee:sgd} gives identical students identical updates.

\emph{(b) The fixed point.} At $w_j=w=a'\sum_kw_{\star k}$, all preactivations equal $z=a'\sum_kz_{\star k}$, with $\Var z=Q=Ka'^2$ and $\operatorname{Cov}(z,z_{\star k})=a'$. Put $\Lambda=\sqrt{2+(2K-1)a'^2}$. By \cref{lem:committee:erf}(ii), the three integrals $\E[g'(z)z_{\star1}g(z_{\star1})]$, $\E[g'(z)z_{\star1}g(z_{\star2})]$, and $\E[g'(z)z_{\star1}g(z)]$ are $\frac2\pi(1+Q)^{-1}$ times, respectively,
\[
  \frac{1+Q-a'^2}{\Lambda},\qquad\frac{-a'^2}{\Lambda},\qquad\frac{a'}{\sqrt{1+2Q}},
\]
where we used $2(1+Q)-a'^2=\Lambda^2$, $(1+Q)^2-Q^2=1+2Q$, and $C_{bc}(1+C_{aa})-C_{ab}C_{ac}=a'(1+Q)-a'Q$ in the third. Summing with weights $1$, $K-1$, and $-K$, and using $1+Q-Ka'^2=1$,
\[
  G_{R,jk}=\frac{2}{\pi(1+Q)}\Bigl(\frac1{\sqrt{2+(2K-1)a'^2}}-\frac{Ka'}{\sqrt{1+2Ka'^2}}\Bigr).
\]
The bracket vanishes exactly when $a'>0$ and $1+2Ka'^2=K^2a'^2(2+(2K-1)a'^2)$, a quadratic equation in $a'^2$ with product of roots $-1/(K^2(2K-1))<0$, hence with exactly one positive root, and $a'^2=1/(K(2K-1))$ satisfies it: $\frac{1}{2K-1}+\frac{2(K-1)}{2K-1}=1$. The $Q$-entries of $G$ vanish too, because $z=a\sum_kz_{\star k}$ makes $\E[\delta_jz]=a\sum_k\E[\delta_jz_{\star k}]=0$. For the loss, \eqref{eq:committee:loss} with $Q/(1+Q)=\frac1{2K}$ and $a/\sqrt{2(1+Q)}=\frac1{2K}$ gives $\Loss_\star=\frac1\pi(K\arcsin\frac12+K^2\arcsin\frac1{2K}-2K^2\arcsin\frac1{2K})$.

\emph{The Hessian at $P_\star$.} Let $W_\star=(w,\dots,w)$ with $w=a\sum_kw_{\star k}$, and let $z=\ip{w,x}$, $\delta=g'(z)(y-Kg(z))$. The point $W_\star$ is a critical point of $\Loss$: the gradient $-\E[\delta x]$ lies in the span of the teachers, because $\delta$ depends on $x$ only through $z_{\star1},\dots,z_{\star K}$, and its inner products with the teachers vanish by (b). Differentiating $\nabla_{w_j}\Loss=-\E[g'(z_j)(y-f)x]$ once more at $W_\star$, the blocks of the Hessian are $H_{jl}=\cB+\ind\{j=l\}\cA$ with
\[
  \cB=\E[g'(z)^2xx^\top],\qquad \cA=-\E[(y-Kg(z))g''(z)\,xx^\top]=\E[\delta z\,xx^\top],
\]
using $g''(z)=-zg'(z)$. For a perturbation $W=W_\star+\varepsilon(p_1,\dots,p_K)$ with $\sum_jp_j=0$, the sum $\sum_lH_{jl}p_l$ equals $\cA p_j$, so the gradient flow gives $\dot p_j=-\cA p_j+O(\varepsilon)$.

\emph{(d) The marginal directions.} Write $x=\sum_kz_{\star k}w_{\star k}+x_\perp$, with $x_\perp$ independent of the $z_{\star k}$. Since $\delta z$ depends only on the $z_{\star k}$, $\cA u=\E[\delta z]\,u$ for $u$ orthogonal to the teachers, and $\E[\delta z]=0$ at $P_\star$ by (b). So $\cA u_j=0$, which is (d).

\emph{(c) The specialization eigenvalue.} Let $v=(w_{\star1}-w_{\star2})/\sqrt2$ and $V=\ip{v,x}=(z_{\star1}-z_{\star2})/\sqrt2$, a standard Gaussian independent of $z$ and of $z_{\star3},\dots,z_{\star K}$. Exchanging teachers $1$ and $2$ preserves the law of $x$ and the function $\delta z$ but changes the sign of $V$; so $\ip{w_{\star k},\cA v}=\E[\delta z\,z_{\star k}V]$ vanishes for $k\ge3$ and is antisymmetric in $k\in\{1,2\}$, and with $\cA$ vanishing off the teacher span this means $\cA v=\mu v$ with $\mu=\E[\delta zV^2]$. For a standard Gaussian $V$ and a smooth $h$ of polynomial growth, two applications of Stein's lemma give $\E[V^2h(V)]=\E h(V)+\E h''(V)$. Write $z_{\star1}=U+V/\sqrt2$ and $z_{\star2}=U-V/\sqrt2$ with $U=(z_{\star1}+z_{\star2})/2$; then $V$ is independent of $(U,z_{\star3},\dots,z_{\star K})$, and $z$ is a function of these variables. Apply the identity conditionally on them. Only $g(z_{\star1})+g(z_{\star2})$ depends on $V$, with second derivative $\frac12(g''(z_{\star1})+g''(z_{\star2}))$, so
\[
  \mu=\E[\delta z]+\tfrac12\E\bigl[g'(z)z\,(g''(z_{\star1})+g''(z_{\star2}))\bigr]=-\E\bigl[g'(z)g'(z_{\star1})\,z\,z_{\star1}\bigr],
\]
using $\E[\delta z]=0$, the symmetry between teachers $1$ and $2$, and $g''=-zg'$. The pair $(z,z_{\star1})$ has covariance $C_2=\begin{psmallmatrix}Q&a\\a&1\end{psmallmatrix}$, with $\det(I_2+C_2)=2(1+Q)-a^2=(2K+1)/K$. \Cref{lem:committee:tilt} with $S$ the pair gives
\[
  \E[g'(z)g'(z_{\star1})zz_{\star1}]=\frac{2}{\pi}\,\frac{\bigl(C_2(I_2+C_2)^{-1}\bigr)_{12}}{\det(I_2+C_2)^{1/2}}=\frac2\pi\,\frac{a}{\det(I_2+C_2)^{3/2}},
\]
because $C_2-C_2(I_2+C_2)^{-1}C_2=C_2(I_2+C_2)^{-1}$, whose off-diagonal entry is $(a(1+Q)-aQ)/\det(I_2+C_2)$. Inserting $a$ and the determinant gives $-\mu=\committeelu$.

\emph{From weights to order parameters.} Take $p_1=v$, $p_2=-v$, and $p_j=0$ for $j\ge3$. Along $W_\star+\varepsilon(p_j)$ the order parameters are $\omega(\varepsilon)=P_\star+\frac\varepsilon{\sqrt2}\Xi+O(\varepsilon^2)$, because $\ip{v,w_{\star1}}=-\ip{v,w_{\star2}}=1/\sqrt2$ and $\ip{v,w}=0$. By \cref{rem:committee:tau}, $G(\omega(\varepsilon))$ is the time derivative of the order parameters under the gradient flow, whose first-order part is $\dot w_j=\varepsilon\committeelu p_j$. Its $R$-entries are $\varepsilon\committeelu\ip{p_j,w_{\star k}}=\frac{\varepsilon}{\sqrt2}\committeelu\Xi_{R,jk}$, and its $Q$-entries are $\varepsilon\committeelu(\ip{p_j,w}+\ip{w,p_l})=0$. Differentiating at $\varepsilon=0$ gives $\frac1{\sqrt2}J\Xi=\frac1{\sqrt2}\committeelu\Xi$ for the Jacobian $J$ of $G$, which is continuously differentiable near $P_\star$ by \cref{lem:committee:interp}.
\end{proof}

The proof of (d) uses the identity $g''=-zg'$, and it is special to this activation: for a general odd $g$ the orthogonal split has curvature $-\E[(y-Kg(z))g''(z)]$, which need not vanish (\cref{ex:committee:activation}).

\begin{heuristic}[The plateau lasts $\frac1{2\committeelu(\eta)}\log d$]
Let $\committeelu(\eta)$ be the eigenvalue of the full field $F_\eta$ in the direction $\Xi$, at the symmetric fixed point $P_\star(\eta)$ of $F_\eta$ in $\cS_0$; by (c), $\committeelu(\eta)=\eta\committeelu+O(\eta^2)$. Start from small random weights. The initial order parameters are symmetric up to fluctuations of order $d^{-1/2}$, and in time $O(1)$ the ODE carries them close to $P_\star(\eta)$ while the component along $\Xi$, call it $S$, stays of order $d^{-1/2}$. Near $P_\star(\eta)$ the recursion \eqref{eq:committee:recursion} for $S$ reads $S_{s+1}=S_s+\frac1d(\committeelu(\eta)S_s+\xi_s)$ to first order, with centered noise $\xi_s$ of variance of order one, so in time $t$,
\[
  S(t)\approx e^{\committeelu(\eta)t}\Bigl(S(0)+\frac1d\sum_{s<td}e^{-\committeelu(\eta)s/d}\xi_s\Bigr),
\]
and the bracket converges to a random variable $\Gamma d^{-1/2}$ with $\Gamma$ of order one: the initial asymmetry and the asymmetry injected by the noise are of the same order. The plateau ends when $S$ becomes of order one, at
\[
  t_{\mathrm{esc}}\approx\frac{1}{2\committeelu(\eta)}\log d+\frac1{\committeelu(\eta)}\log\frac{1}{\abs\Gamma}+O(1),
\]
a time of order $\log d$ with fluctuations of order one that do not shrink as $d$ grows. In steps, the plateau costs about $\frac{d}{2\committeelu(\eta)}\log d$ samples.
\end{heuristic}

The heuristic leaves the range of \cref{thm:committee:ode}, which controls the dynamics only up to a fixed horizon $t_1$, while $t_{\mathrm{esc}}$ grows with $d$; and the escape is decided by fluctuations of order $d^{-1/2}$, exactly the scale the theorem neglects. Without them the ODE started on $\cS$ never leaves it, and SGD started from identical students never specializes, at every $d$.\sidenote{In floating point an exact tie can be broken by rounding errors, but these are many orders of magnitude below $d^{-1/2}$, so such a run would sit on the plateau far longer.} This is the escape from a saddle at small initialization that recurs in this book, with the logistic time $\log(1/\text{initialization})/\text{strength}$ of \cref{ch:deeplinear}: the initialization is $d^{-1/2}$ and the strength is $\committeelu(\eta)$.

\section{Reading the solution}\label{sec:committee:reading}

The solution makes four quantitative predictions. The plateau height $\Loss_\star=K/6-(K^2/\pi)\arcsin\frac1{2K}$ is $0.0116$, $0.0203$, and $0.0284$ for $K=2,3,4$, a small fraction of the initial loss $K/6$: the symmetric student, which sends every unit to the average teacher direction, already captures most of the variance of the label, and what remains is the part that requires a division of labor. The plateau length is $\frac1{2\committeelu}\log d$ in units of $\tau=\eta t$, and since $\committeelu\approx\frac1{2\pi K}$ for large $K$, it grows linearly in the number of units as well as logarithmically in the dimension.\sidenote{The exact values are $\committeelu=0.0657$, $0.0461$, and $0.0356$ for $K=2,3,4$, against $1/(2\pi K)=0.0796$, $0.0531$, and $0.0398$.} The level and the escape rate depend on $\eta$ through the $\eta^2$ term: at $\eta=0.5$ and $K=2$ the plateau rises from $0.0116$ to $0.0147$, and $\committeelu(\eta)/\eta$ is $0.0662$, $0.0670$, $0.0687$, and $0.0748$ for $\eta=0.1$, $0.25$, $0.5$, and $1$. Finally, every unit's step size is limited by the threshold of \cref{prop:committee:single}, and the committee's threshold is lower: for $K=2$ the specialized fixed point $R=Q=I$ loses stability at $\eta\approx2.31$, against $4.06$ for one unit (\cref{ex:committee:etamax}).

\begin{figure}[t]
\centering
\begin{tikzpicture}
\begin{groupplot}[group style={group size=2 by 1,horizontal sep=1.35cm},hplot,width=0.48\linewidth,height=0.46\linewidth,
  xlabel={$t=$ steps$/d$},xmin=0,xmax=300]
\nextgroupplot[ymode=log,ymin=2e-5,ymax=0.5,ylabel={population loss}]
\addplot[guide] coordinates {(0,0.01468) (300,0.01468)};
\addplot[theory] table[x=t,y=eg] {figures/data/08-committee-plateau-ode.dat};
\addplot[simulation2] table[x=t,y=eg300] {figures/data/08-committee-plateau-sgd.dat};
\addplot[simulation] table[x=t,y=eg3000] {figures/data/08-committee-plateau-sgd.dat};
\nextgroupplot[ymin=0,ymax=1.05,ylabel={overlap},legend pos=north west]
\addplot[theory] table[x=t,y=R11] {figures/data/08-committee-plateau-ode.dat};
\addplot[theory2] table[x=t,y=R12] {figures/data/08-committee-plateau-ode.dat};
\addplot[simulation] table[x=t,y=R11_3000] {figures/data/08-committee-plateau-sgd.dat};
\addplot[simulation2] table[x=t,y=R12_3000] {figures/data/08-committee-plateau-sgd.dat};
\legend{ODE $R_{11}$,ODE $R_{12}$,SGD $R_{11}$,SGD $R_{12}$}
\end{groupplot}
\end{tikzpicture}
\caption{Plateau and specialization for $K=2$ students and orthonormal teachers, $\eta=0.5$. Left: the population loss \eqref{eq:committee:loss} along the ODE of \cref{thm:committee:ode} (curve), and medians over $16$ runs of SGD in dimension $d=300$ (squares) and $d=3000$ (dots); the dashed line is the plateau level $\Loss_\star(\eta)=0.0147$ of the full field. Right: the overlaps of student $1$ with teachers $1$ and $2$, ODE (curves) and medians at $d=3000$ (points). All runs start from the same order parameters, $R=0.003\,I$ and $Q-RR^\top=0.01\,I$, and are simulated with the exact chain of \cref{rem:committee:chain} (seed $8$).}\label{fig:committee:plateau}
\end{figure}

\Cref{fig:committee:plateau} tests the ODE limit across a plateau. The runs start with a slight preference of each student for its own teacher, $R(0)=0.003\,I$, so that the ODE itself escapes; its loss falls below half the plateau level at $t=166$. In the first $20$ time units the loss falls by a factor $23$ to the plateau, and both students move to the blend $R_{jk}\approx0.41$, close to $a=1/\sqrt6\approx0.408$ (the full field at $\eta=0.5$ puts the fixed point at $0.410$). On the plateau the medians agree with the ODE to within $1\%$ in loss, for both dimensions. The escape is where the noise matters: the medians cross the half-plateau level at $t=170$ in both dimensions, about $4$ time units after the ODE, and individual runs scatter around it, because the noise injects asymmetry of order $d^{-1/2}$ in both directions, which shortens some plateaus and lengthens others. After the escape the descent is again deterministic: $R_{11}$ rises to $0.98$ and $R_{12}$ falls to $0.01$ by $t=300$, with the loss decaying exponentially at the rate of the specialized fixed point.

\begin{figure}[t]
\centering
\begin{tikzpicture}
\begin{axis}[hplot,xmode=log,log basis x=2,xmin=80,xmax=4000,ymin=110,ymax=230,
  xlabel={dimension $d$},ylabel={escape time $t_{\mathrm{esc}}$},legend pos=north west,
  xtick={100,200,400,800,1600,3200},xticklabels={$100$,$200$,$400$,$800$,$1600$,$3200$}]
\addplot[theory] table[x=d,y=theory] {figures/data/08-committee-escape.dat};
\addplot[simulation,error bars/.cd,y dir=both,y explicit] table[x=d,y=median,
  y error plus expr=\thisrow{q75}-\thisrow{median},y error minus expr=\thisrow{median}-\thisrow{q25}]
  {figures/data/08-committee-escape.dat};
\legend{slope $1/(2\committeelu(\eta))$,SGD median and quartiles}
\end{axis}
\end{tikzpicture}
\caption{Length of the symmetric plateau from random initialization, $K=2$, $\eta=0.5$, $w_j(0)\sim\cN(0,\frac{0.01}{d}I_d)$. Points: median and quartiles over $40$ runs of the first time at which the loss falls below half the plateau level, simulated with the exact chain of \cref{rem:committee:chain} (seed $9$). Line: slope $1/(2\committeelu(\eta))=14.6$ per unit of $\log d$, computed from the Jacobian of the full field at $P_\star(\eta)$, with the intercept fitted to the medians.}\label{fig:committee:escape}
\end{figure}

\Cref{fig:committee:escape} tests the heuristic from random initialization, where the only asymmetry is accidental. The predicted slope is $1/(2\committeelu(\eta))=14.6$ per unit of $\log d$; a least-squares fit to the medians gives $13.7$, and over the factor $32$ in dimension the medians rise from $145$ to $193$. The slope is a parameter-free prediction; the intercept is not, since it depends on the law of $\Gamma$, which the heuristic does not compute. The quartiles show the other prediction: the interquartile range stays between $30$ and $40$ time units at every $d$, as it must if the escape time is $\frac1{2\committeelu}\log d$ plus a fluctuation of order one.

\begin{figure}[t]
\centering
\begin{tikzpicture}
\begin{axis}[hplot,xmin=0,xmax=4.7,ymin=-0.03,ymax=0.75,xlabel={$\eta$},ylabel={decay rate of $\norm{w-w_\star}^2$},legend pos=north west]
\addplot[guide] coordinates {(4.0558,-0.03) (4.0558,0.75)};
\addplot[theory] table[x=eta,y=rate] {figures/data/08-committee-single-theory.dat};
\addplot[simulation] table[x=eta,y=rate] {figures/data/08-committee-single-sim.dat};
\legend{$\eta_{\max}=\pi\sqrt{5/3}$,{$\min(2\eta c_1,2\eta c_2-\eta^2c_3)$},SGD}
\end{axis}
\end{tikzpicture}
\caption{The learning-rate threshold of a single $\erf$ unit (\cref{prop:committee:single}). Curve: the asymptotic decay rate predicted by the linearization, which vanishes at $\eta_{\max}\approx4.056$ (dashed). Points: minus the slope of the median of $\log\norm{w-w_\star}^2$ over $8$ runs of SGD with $d=1000$, fitted where $10^{-40}<\norm{w-w_\star}^2<10^{-6}$, from $R=0$ and $Q=0.25$ (exact chain, seed $10$). Beyond $\eta_{\max}$ the runs do not converge and the fitted rate is zero.}\label{fig:committee:single}
\end{figure}

\Cref{fig:committee:single} tests \cref{prop:committee:single}. The measured rates agree with the predicted rate within $2\%$ for $\eta\le3.75$ and within $4\%$ at $\eta=4$, where the decay is slow and the fitting window short; the largest deviations elsewhere are near the kink at $\eta=2.70$, where the two modes have nearly equal rates. For small $\eta$ the slow mode is the component along $w_\star$; for large $\eta$ it is the orthogonal component, whose contraction $2\eta c_2$ is fought by the noise $\eta^2c_3$. At $\eta=4.25$ and $4.5$ the runs settle at a nonzero error, as the bifurcation in \cref{sec:committee:single} predicts.

\section{What the model gets wrong}\label{sec:committee:wrong}

\paragraph{What was thrown away.}
Isotropic Gaussian inputs, which made the overlaps sufficient statistics and every integral explicit; with structured inputs the state of the dynamics is no longer finite-dimensional. A fixed second layer, which removes the possibility of reweighting units instead of moving them. A width fixed while $d\to\infty$, and a teacher of the same architecture and width, so that the target is realizable with exactly the right number of units. One pass over fresh samples. Small steps of size $\eta/d$, which make the dynamics deterministic at the scale of the theorem and push the decisive fluctuations below it. And small initialization, which is what places the trajectory near the symmetric saddle in the first place; a student started far from $\cS$ need not visit it.

\paragraph{Where it breaks.}
The first failure is internal. The quantity the model is built to explain, the length of the plateau, is set by fluctuations of order $d^{-1/2}$ that \cref{thm:committee:ode} discards: the limit equations predict a plateau of infinite length from every symmetric start, and the heuristic of \cref{sec:committee:plateau} goes beyond what is proved. We do not know a proof of the $\frac1{2\committeelu}\log d$ law for this model. The other failures are external. \emph{Overparameterization.} With more students than teachers the picture changes qualitatively: \textcite{GoldtEtAl2019} prove that the ODE description remains exact and find that the final error of the student grows with the width when only the first layer is trained but stays constant or decreases when both layers are trained; and \textcite{RichertWorschechRosenow2022} show that in the over-realizable case the approach to perfect learning follows a power law, not the exponential convergence of the matched case. \emph{Data.} \textcite{YoshidaOkada2019} show that the plateau depends on the input statistics, and becomes less noticeable when the covariance has small and dispersed eigenvalues; \textcite{GoldtEtAl2020} replace isotropic inputs by the hidden manifold model, inputs that lie near a low-dimensional manifold with labels that depend on the position on it, and show that the learning dynamics then differs from that of i.i.d.\ Gaussian inputs, which the Gaussian model cannot capture. \emph{Optimizers.} Natural gradient descent, which preconditions by the Fisher information, shortens or removes the plateaus \parencite{Amari1998,RattraySaad1999}, so the plateau is a property of plain gradient descent as much as of the network. \emph{Steps and width.} \textcite{VeigaEtAl2022} map the interplay of the learning rate, the time scale, and the number of hidden units, and show that the regime of this chapter, small steps with $m$ fixed, is one corner of a phase diagram whose other corners include the mean-field limit of wide networks and regimes in which the noise dominates the dynamics. \emph{The readout.} \textcite{RieglerBiehl1995} analyze the case where the hidden-to-output weights are trained as well, and compare it with fixed readouts; in the network of a practitioner, the second layer is never frozen.

\paragraph{Fine structure.}
The repairs come in three kinds. \emph{Exact descriptions of more general dynamics.} \textcite{ArnaboldiStephanKrzakalaLoureiro2023} derive one low-dimensional deterministic description that connects the small-step, high-dimensional, and mean-field limits of one-pass SGD in two-layer networks, and \textcite{CollinsWoodfinEtAl2023}, in an arXiv preprint, give ODE deterministic equivalents for streaming SGD on generalized linear and multi-index models, with learning-rate thresholds for stability that generalize \cref{prop:committee:single}; \textcite{BenArousGheissariJagannath2022} prove ballistic and diffusive limit theorems for summary statistics of SGD with constant step size, find a critical step size at which a correction term enters the effective dynamics, the counterpart of our $\eta^2$ term, and show that diffusive limits near fixed points, the regime where the escape from a plateau is decided, can be complicated. \emph{Targets that require several escapes.} For multi-index targets on Gaussian or Boolean inputs, \textcite{AbbeBoixAdseraMisiakiewicz2023} introduce the leap complexity, conjecture that SGD on two-layer networks needs total time about $d^{\max(\mathrm{Leap},2)}$, that is about $d^{\max(\mathrm{Leap}-1,1)}$ samples, prove a version for Gaussian inputs under additional assumptions on how SGD is run, and show that the support of the target is learned in a sequence of saddle-to-saddle transitions; the cascade of plateaus of \cref{obs:committee:plateau} is the matched-teacher ancestor of this picture, and \cref{ch:singleindex,ch:parity} develop the single-direction case. \textcite{BerthierMontanariZhou2023}, in a preprint, describe learning of single-index targets by wide networks as a sequence of separated time scales, partly by non-rigorous derivations. \emph{Landscape and symmetry.} \textcite{FukumizuAmari2000} and \textcite{AmariParkOzeki2006} relate plateaus to the singular regions created by duplicated units, and \textcite{CousseauOzekiAmari2008} analyze the dynamics near them; \textcite{SimsekEtAl2021} count the symmetry-induced critical points of overparameterized networks and show that one extra neuron per layer connects the permutation-equivalent minima into a single manifold. Open: a proof of the logarithmic plateau law and of the law of the escape time; the plateau length when $m>m_\star$; and a description of plateaus with data that are neither Gaussian nor isotropic.

\section{Model card}\label{sec:committee:card}

\modelcard{Plateaus in online training of small sigmoidal networks from small weights: hidden units are redundant copies, then specialize.}
{Teacher and student soft committee machines with $g(z)=\erf(z/\sqrt2)$, Gaussian inputs, online SGD with step $\eta/d$; order parameters $Q,R,T$.}
{As $d\to\infty$ the overlaps follow explicit ODEs with an $\eta^2$ noise term (proved); one unit converges for $\eta<\pi\sqrt{5/3}$ (proved); for $K$ students and teachers, a symmetric saddle with loss $\frac K6-\frac{K^2}\pi\arcsin\frac1{2K}$ and unstable eigenvalue $\frac{2K}{\pi(2K+1)^{3/2}(2K-1)^{1/2}}$ (proved).}
{Plateau level; plateau length $\frac1{2\committeelu}\log d$ in units of $\eta t$, growing linearly in $K$ (derived); optimal and maximal learning rates.}
{The plateau length is decided by fluctuations the ODE omits; overparameterization, structured data, natural gradients, large steps, and trained readouts change or remove the plateau.}
{Overparameterized students; hidden manifold and Gaussian equivalence; dimensionless ODEs and phase diagrams of SGD; leap complexity and saddle-to-saddle dynamics; effective dynamics of summary statistics.}

The committee machine is the matrix version of the single-index model of \cref{ch:singleindex}: Gaussian rotational invariance again reduces the state to overlaps, and the ODE method of \cref{thm:toolkit:ode-method} turns online SGD into a deterministic flow on them. The two chapters produce plateaus for different reasons. In \cref{ch:singleindex} the target is invisible at low degree, and the plateau is a search whose length is a power of $d$; here every unit sees its signal at once, and the plateau is a symmetric saddle whose length is a logarithm. Both are escapes from small initialization, the thread that runs from the deep linear networks of \cref{ch:deeplinear} and the diagonal networks of \cref{ch:diagonal} to the hidden progress of \cref{ch:parity} and the induction heads of \cref{ch:induction}; the escape here has the logistic form of \cref{ch:deeplinear}, with initialization $d^{-1/2}$.

The $\eta^2$ term connects the chapter to the stability of gradient descent in \cref{ch:prologue,ch:stability}: the threshold $\pi\sqrt{5/3}$ is a stability bound in which the variance of the stochastic gradient plays the role that the top curvature plays for full-batch descent. And the plateau is one instance of the abrupt transitions collected in \cref{ch:parity,ch:grokking,ch:induction,ch:scaling}, where a smooth underlying dynamics produces a sudden drop in the loss.

\begin{exercises}
\item \consolidation (Facts about the $\erf$ unit.) (a) Prove \eqref{eq:committee:gprime}, $\E g(z)^2=\frac13$, and $\hat g_1=1/\sqrt\pi$ for $z\sim\cN(0,1)$. (b) Show that the loss at $W=0$ equals $m_\star/6$ for an orthonormal teacher and that the information exponent of $g$ is $1$. (c) Use \cref{lem:committee:tilt} with $S=(a,a)$ to show $\E g'(X_a)^2=\frac2\pi(1+2C_{aa})^{-1/2}$.\label{ex:committee:basics}

\item \consolidation (One unit, explicitly.) For $m=m_\star=1$ and $T=1$, show that the gradient-flow part of the ODE is
\[
  \frac{\dd R}{\dd\tau}=\frac{2}{\pi(1+Q)}\Bigl(\frac{1+Q-R^2}{\sqrt{2(1+Q)-R^2}}-\frac{R}{\sqrt{1+2Q}}\Bigr),
\]
derive the corresponding equation for $Q$, check that $(R,Q)=(1,1)$ is a fixed point, and compute the Jacobian there numerically for $\eta=1$ to confirm the eigenvalues $-c_1$ and $-(2c_2-c_3)$ of \cref{prop:committee:single}.\label{ex:committee:one}

\item \consolidation (The exact chain.) Prove \cref{rem:committee:chain}: write $\norm{Px}^2$ as a quadratic form in $(z,z_\star)$ using the pseudoinverse of $\committeeM(\omega)$, and show that, given $\cF_s$, the law of $\omega_{s+1}$ depends only on $\omega_s$. For $m=m_\star=2$ and $T=I$, give an explicit sampler that uses five scalar random variables per step.\label{ex:committee:chain}

\item \consolidation (The other modes at $P_\star$.) (a) Let $p_j=v$ for every $j$, with $v=(w_{\star1}-w_{\star2})/\sqrt2$: identical students turning toward teacher $1$. Show that the gradient flow gives $\dot p=-(K\cB+\cA)p$ to first order, and that the eigenvalue is
\[
  -\frac{2K}\pi\sqrt{\frac{2K-1}{2K+1}}+\committeelu,
\]
which is $-0.920$ for $K=2$. (b) Explain why the order-parameter Jacobian has zero eigenvalues at $P_\star$ although the weight-space Hessian has none in the corresponding directions except those of \cref{thm:committee:plateau}(d).\label{ex:committee:modes}

\item \variation (Label noise.) Let $y=\sum_kg(z_{\star k})+\sigma\zeta$ with $\zeta\sim\cN(0,1)$ independent. Show that the ODE acquires the term $\eta^2\sigma^2\E[g'(z_j)g'(z_l)]$ in $\dot Q_{jl}$ and nothing else. For one unit, show that $w_\star$ is no longer a fixed point, and that to leading order the excess loss at the new fixed point is $\frac{\eta\sigma^2c_2^2}{2(2c_2-\eta c_3)}$, which is about $\eta\sigma^2/(2\sqrt3\pi)$ for small $\eta$ and diverges at $\eta_{\max}$.\label{ex:committee:noise}

\item \variation (Unequal teachers.) Let $K=2$ and $T=\diag(1,\vartheta^2)$ with $\vartheta\ne1$. Show that $\cS$ is no longer invariant, so that the plateau of \cref{thm:committee:plateau} disappears as an exact object. Integrate the ODE numerically from a symmetric start for $\vartheta=0.9$ and $\vartheta=0.5$, and describe what replaces the plateau: which teacher is learned first, and how the delay depends on $\abs{1-\vartheta}$.\label{ex:committee:unequal}

\item \variation (Other activations.) For an odd activation $\phi$ in place of $g$, with the corresponding symmetric fixed point, show that the Hessian block $\cA$ acts on directions orthogonal to the teachers as multiplication by $-\E[(y-K\phi(z))\phi''(z)]$. Show that this vanishes for every activation with $\phi''(z)=-c\,z\,\phi'(z)$ for a constant $c$, and give an odd activation for which it does not vanish. What does this change in \cref{thm:committee:plateau}(d)?\label{ex:committee:activation}

\item \variation (Specialization for general $K$.) Prove that for $K\ge3$ every vector $v$ in the teacher span with $\ip{v,\sum_kw_{\star k}}=0$ satisfies $\cA v=-\committeelu v$, so that the unstable eigenvalue of $P_\star$ has multiplicity at least $K-1$ in weight space. How many distinct ways to break the symmetry does this give, and how does the plateau length depend on $K$?\label{ex:committee:generalK}

\item \simulation (The committee's learning-rate threshold.) For $K=2$, $T=I$, compute numerically the Jacobian of the full field $F_\eta$ at the specialized fixed point $R=Q=I$ for $\eta$ between $1$ and $3$. You should find that its largest eigenvalue crosses zero at $\eta\approx2.31$, well below the single-unit threshold $4.06$; identify the unstable eigenvector and explain the role of the cross terms $\E[\delta_1\delta_2]$.\label{ex:committee:etamax}

\item \simulation (Plateau length and the number of units.) Integrate the ODE with $\eta=0.1$ from $Q(0)=0.01\,I$ and $R(0)=\epsilon I$ for $\epsilon=10^{-3}$ and $10^{-5}$, for $K=2$ and $K=3$, and record the first time after the approach at which the loss falls below $\Loss_\star/2$. You should find that the escape time grows by about $153$ (for $K=2$) and $217$ (for $K=3$) per factor $e$ decrease of $\epsilon$, within $1\%$ of $1/(\eta\committeelu)$ from \cref{thm:committee:plateau}.\label{ex:committee:Kscaling}

\item \simulation (The chain against SGD.) For $K=2$, $\eta=0.5$, $d=100$, and $w_j(0)\sim\cN(0,\frac{0.01}dI_d)$, run $400$ copies of SGD in $\R^{100}$ and $400$ copies of the chain of \cref{rem:committee:chain} up to $t=80$. Compare the medians and quartiles of the loss at $t=10,20,40,80$; you should observe agreement within $1\%$ in the medians. Then repeat with $\eta=3$, run to $t=200$, and explain with \cref{ex:committee:etamax} why the loss no longer tends to zero.\label{ex:committee:sgdcheck}

\item \research (Plateaus with more students than teachers.) For $m_\star=2$ orthonormal teachers and $m=3$ or $4$ students, find the symmetric fixed points of the gradient-flow field, their unstable eigenvalues, and the predicted plateau length, and compare with simulations. Over-realizable learning ends in a power law rather than an exponential \parencite{RichertWorschechRosenow2022}; how the plateau that precedes it depends on $m/m_\star$ is, as far as we know, not settled.\label{ex:committee:over}

\item \research (A theorem for the escape.) Prove that for $K=2$ and fixed $\eta<\eta_{\max}$, the escape time from random initialization divided by $\log d$ converges in probability to $1/(2\committeelu(\eta))$, and identify the limit law of $t_{\mathrm{esc}}-\frac1{2\committeelu(\eta)}\log d$. The limit theorems of \textcite{BenArousGheissariJagannath2022} for summary statistics near fixed points are a natural starting point; we do not know of a proof for this model.\label{ex:committee:escapethm}
\end{exercises}

\notesandrefs
The committee machine belongs to the statistical mechanics of learning, which \textcite{EngelVandenBroeck2001} survey, from the perceptron to multilayer networks and online learning. The exact dynamics of online learning came in 1995. \Textcite{SaadSolla1995} announced the exact solution for soft committee machines with arbitrary numbers of hidden units in student and teacher, and \textcite{SaadSolla1995b} gave the full derivation of the differential equations for the overlaps, treated the realizable, over-realizable, and unrealizable cases, and studied how the learning rate affects convergence. \Textcite{BiehlSchwarze1995} analyzed online gradient descent in two-layer networks independently in the same year, and \textcite{RieglerBiehl1995} extended the analysis to networks whose hidden-to-output weights are trained.\sidenote{The preprint of the second paper lists Saad at the University of Edinburgh and Solla at the Niels Bohr Institute in Copenhagen; Biehl, Schwarze, Riegler, and W\"ohler circulated their work as W\"urzburg preprints.} \Textcite{BiehlRieglerWohler1996}, in a preprint, found several fixed points of the learning dynamics of the soft committee machine and the cascades of plateaus that connect them.

The plateau problem shaped the next decade of work on optimization. \Textcite{Amari1998} introduced the natural gradient and suggested that it makes the plateau problem of backpropagation less serious, and \textcite{RattraySaad1999} analyzed natural gradient descent in the committee machine with the methods of this chapter, finding shortened or removed plateaus. \Textcite{FukumizuAmari2000} explained the plateaus geometrically, through the lines of critical points that a smaller network embeds in a larger one, and \textcite{AmariParkOzeki2006} and \textcite{CousseauOzekiAmari2008} studied learning near such singularities.\sidenote{The geometric view and the ODE view describe the same object: the set $\cS_0$ of identical students is the image of a network with one hidden unit inside the network with $K$.}

The modern revival began with \textcite{GoldtEtAl2019}, who proved that the ODE description is asymptotically exact and used it to study overparameterized students, and with \textcite{GoldtEtAl2020}, who extended it to structured inputs. \Textcite{VeigaEtAl2022} drew the phase diagram of the regimes of one-pass SGD and gave rigorous convergence rates for the ODE description, and \textcite{ArnaboldiStephanKrzakalaLoureiro2023} unified the high-dimensional and mean-field limits. The proof of \cref{thm:committee:ode} here is a direct application of the elementary ODE method of the appendix, with constants that are explicit but not sharp; the closed-form unstable eigenvalue of \cref{thm:committee:plateau}, derived in weight space with Stein's lemma, is our arrangement of a computation that the 1990s literature performed numerically or in order-parameter coordinates.

\paragraph{Translating the literature.}
\Textcite{SaadSolla1995b} and the papers that follow their conventions, such as \textcite{RichertWorschechRosenow2022}, write $N$ for our $d$, $\xi$ for our input $x$, $J_i$ for our student weights $w_j$, $B_n$ for our teacher weights $w_{\star k}$, $K$ and $M$ for our $m$ and $m_\star$, and update $J_i^{\mu+1}=J_i^\mu+\frac\eta N\delta_i^\mu\xi^\mu$, which is \eqref{eq:committee:sgd} with $\mu$ for our step index $s$; their time $\alpha=\mu/N$ is our $t$, their generalization error $\epsilon_g$ is our $\Loss$, and the overlaps $Q$, $R$, $T$ keep their names, with student indices $i,k$ and teacher indices $n,m$. Their integrals $I_2$, $I_3$, $I_4$ are those of \cref{lem:committee:erf}. \Textcite{GoldtEtAl2019} write $N$ for $d$ and $K$, $M$ for the numbers of student and teacher units, keep the entries of the weights of order one and divide the preactivations by $\sqrt N$, so that their order parameters, normalized by $N$, are ours; they also train the second layer, which our model freezes at one.

\printchapterbib
\end{refsection}

% Chapter 9 of The Hydrogen Atoms of Deep Learning: sparse parities and the limits of gradient methods.
\newcommand{\parityMaj}{\operatorname{Maj}}

\casechapter{Hidden Progress}{Sparse parities and the limits of gradient methods}{Sparse parity on the hypercube}{ch:parity}
\begin{refsection}[chapters/09-parity.bib]

\begin{chaptersummary}
Networks trained by stochastic gradient descent on a sparse parity, the product of $k$ unknown coordinates of a uniformly random sign vector, show no improvement in loss or error for a long time and then learn abruptly, while measurements of the weights reveal progress during the flat stretch. The model of this chapter is the parity itself, learned by one neuron with the activation $z\mapsto z^k$. Its correlation with the target is exactly $k!$ times the product of the $k$ relevant weights, so the population gradient is supported on the relevant coordinates; we prove that one gradient step estimated from about $d^{k-1}\log d$ samples finds them, that every statistical-query algorithm needs a number of queries times an inverse squared tolerance of order $d^k$, and that the population flow spends a time of order $d^{(k-2)/2}$ on a loss plateau while the logarithm of the relevant weights grows steadily from the start. The model is wrong about how networks actually find the coordinates: with random initialization the progress sits in a few lucky neurons, and natural tasks offer low-degree footholds that parities deny.
\end{chaptersummary}

\section{The phenomenon: abrupt learning with steady progress}\label{sec:parity:phenomenon}

Parity is the oldest hard case of neural network theory. \Textcite{MinskyPapert1969} showed that a perceptron that combines predicates on a bounded number of inputs cannot compute the parity of all its inputs; the parity of $n$ bits needs predicates of order $n$.\sidenote{Their \emph{order} of a perceptron is the largest number of inputs that one of its predicates reads.} Modern networks represent parities without difficulty: a single hidden layer of $k+1$ rectified linear units computes the parity of $k$ given bits (\cref{cor:parity:network}). The difficulty that remains is learning: when the $k$ bits are unknown, nothing in the data points to them until they have been found.

Fix a dimension $d$ and an unknown set $S\subseteq[d]=\{1,\dots,d\}$ of size $k$. Inputs are uniform on the cube $\{-1,1\}^d$, and the label is the \term{sparse parity}\index{sparse parity}
\[
  y=\chi_S(x)=\prod_{i\in S}x_i .
\]
Statistically the problem is easy: there are $\binom dk$ candidate sets, so about $k\log d$ samples identify $S$ by exhaustive search, and Gaussian elimination\index{Gaussian elimination} over the field with two elements solves it in polynomial time from about $d$ samples.\sidenote{Write $x_i=(-1)^{z_i}$ with $z_i\in\{0,1\}$; then $\chi_S(x)=(-1)^{\sum_{i\in S}z_i}$, and each sample is a linear equation modulo $2$ for the indicator vector of $S$. About $d$ random equations determine it.} Neither algorithm resembles gradient descent. For a gradient method the task is hard for a reason that the Fourier--Walsh expansion makes precise (\cref{thm:toolkit:fourier-cube}): the label $\chi_S$ is orthogonal to every function of the input that does not depend on all $k$ coordinates of $S$, and in particular to every polynomial of degree less than $k$. No statistic of degree below $k$ sees the target.

\begin{observation}[\textcite{BarakEtAl2022}]\label[observation]{obs:parity:hidden}
Train a network by stochastic gradient descent on fresh samples of a $k$-sparse parity of $d$ bits. Across a range of architectures and initializations, the loss and the error stay at their chance values for a long stretch of training and then fall abruptly; on small instances the number of iterations before the drop grows like $d^{O(k)}$, close to the statistical-query lower bound of \cref{thm:parity:sq} below. During the flat stretch the training makes continual progress that the loss and the error do not show: the authors argue that SGD amplifies the sparse solution gradually, through a Fourier gap in the population gradient, rather than stumbling on it by chance.
\end{observation}

\Cref{fig:parity:mlp} reproduces the observation at a size that runs in a minute. A two-layer network of $256$ rectified linear units, with standard Gaussian initialization, learns $x_1x_2x_3$ in dimension $d=50$ by plain SGD on the hinge loss. In each of four runs the test error sits at $\frac12$ for between $12{,}000$ and $25{,}000$ steps and then falls to zero within about $3{,}000$ steps. The progress measure in the right panel, the largest share of a neuron's squared input weight carried by the three relevant coordinates, starts at a value set by the luck of the initialization and begins to grow several thousand steps before the drop, in every run.

\begin{figure}[t]
\centering
\begin{tikzpicture}
\begin{groupplot}[group style={group size=2 by 1,horizontal sep=1.45cm},hplot,width=0.51\linewidth,height=0.47\linewidth,
  xlabel={SGD step$/1000$},xmin=0,xmax=32]
\nextgroupplot[ymin=-0.02,ymax=0.6,ylabel={test error}]
\pgfplotsinvokeforeach{0,...,3}{\addplot[theory,line width=0.8pt] table[x expr=\thisrow{step}/1000,y=err#1] {figures/data/09-parity-mlp.dat};}
\nextgroupplot[ymode=log,ymin=0.15,ymax=1.2,ylabel={largest share on $S$},ytick={0.2,0.3,0.5,1},yticklabels={$0.2$,$0.3$,$0.5$,$1$}]
\pgfplotsinvokeforeach{0,...,3}{\addplot[theory3,line width=0.8pt] table[x expr=\thisrow{step}/1000,y=prog#1] {figures/data/09-parity-mlp.dat};}
\end{groupplot}
\end{tikzpicture}
\caption{Four runs of SGD for a two-layer ReLU network of width $256$ on the parity $x_1x_2x_3$ of $d=50$ uniform bits: hinge loss, step size $0.02$, fresh minibatches of $32$ samples, Gaussian initialization, all layers trained (seeds $0$ to $3$). Left: error on $4000$ test points. Right: the progress measure $\max_j\sum_{i\in S}W_{ji}^2/\norm{W_j}^2$ over the hidden units $j$, on a logarithmic scale; a unit with no preference has share $k/d=0.06$.}\label{fig:parity:mlp}
\end{figure}

How firm is the observation? The flat-then-sudden curve is easy to reproduce, and it was seen long before the name: \textcite{ShalevShwartzShamirShammah2017} used parities as their first example of a failure of gradient-based learning, arguing that when the gradient carries little information about the target, training by gradients cannot succeed. The scaling $d^{O(k)}$ is measured only on small instances, and the exponent depends on architecture and hyperparameters. The statement about hidden progress depends on the measure: the share of the best neuron in \cref{fig:parity:mlp} grows during the last third of the plateau, not from the first step, while averages over all neurons barely move until the drop.

The observation matters beyond parities because it is the cleanest instance of an \term{emergent}\index{emergence} capability: a skill that is absent for a long time and then present. \Textcite{WeiEtAl2022} catalogued abilities of large language models that appear abruptly as models are scaled, and \textcite{SchaefferMirandaKoyejo2023} argued that many such jumps are produced by the choice of metric, nonlinear or discontinuous metrics turning smooth progress into apparent jumps. Parities sit between these positions: the loss is flat for a real reason, yet a different measurement sees continuous progress. The questions for the model are therefore three. How much signal does a gradient carry about $S$? How many samples, or how much computation, must every gradient method spend? And what progresses while the loss stands still?\index{hidden progress}

\section{The model: one monomial neuron on the cube}\label{sec:parity:model}

We use the Fourier--Walsh expansion of \cref{thm:toolkit:fourier-cube}: every $f\colon\{-1,1\}^d\to\R$ is $\sum_{T\subseteq[d]}\hat f(T)\chi_T$ with $\hat f(T)=\E[f(x)\chi_T(x)]$, and $\E f(x)^2=\sum_T\hat f(T)^2$. Throughout, $x$ is uniform on $\{-1,1\}^d$.

\begin{definition}[The sparse parity model]\label[definition]{def:parity:model}
Let $1\le k\le d$ and let $S\subseteq[d]$ with $\abs S=k$. Inputs are $x\sim\Unif\{-1,1\}^d$ and labels are $y=\chi_S(x)$. The \term{monomial neuron}\index{monomial neuron} with weight $w\in\R^d$ predicts $\ip{w,x}^k$ and pays the correlation loss $-y\ip{w,x}^k$. Its population correlation and loss are
\begin{equation}\label{eq:parity:Psi}
  \Psi(w)=\E\bigl[\chi_S(x)\ip{w,x}^k\bigr],\qquad \Loss(w)=-\Psi(w).
\end{equation}
For a two-layer network $f(x;\theta)=\sum_{j=1}^ma_j\phi(\ip{w_j,x}+b_j)$ the correlation is $\E[yf(x;\theta)]$.
\end{definition}

The activation $z\mapsto z^k$ is chosen because it is the lowest-degree neuron that can see the target at all: a polynomial of degree less than $k$ in $x$ is orthogonal to $\chi_S$, so a neuron with a polynomial activation of degree less than $k$ has $\Psi\equiv0$ and receives no gradient. The correlation loss replaces the hinge and square losses of practice for the same reason as in \cref{ch:singleindex}: it makes the population landscape an explicit polynomial, and the square loss $(y-f)^2$ equals twice the correlation loss plus terms that do not involve the label. On the sphere $\norm w=1$, where we study the dynamics, the neuron is the Boolean counterpart of the single-index student of \cref{ch:singleindex} with information exponent $k$; the difference is that the target depends on $k$ coordinates rather than one direction, and that the cube, unlike the Gaussian, is not rotationally invariant, so the coordinates themselves are the natural basis.

Three algorithms will be studied, in increasing realism. First, \emph{statistical-query algorithms}, the abstraction of every method that learns only from averages over the data, including gradient descent with estimated gradients; this gives the lower bound. Second, \emph{one gradient step} from the symmetric initialization $w_0\propto\mathbf 1_d$, estimated from $n$ samples, followed by a small readout; this gives a matching upper bound. Third, the \emph{population gradient flow} on the sphere from a small start, and its noisy version online SGD; this gives the plateau and the hidden progress.

Each simplification has a reason. Uniform inputs make the characters orthonormal and independent of $S$, so that all $\binom dk$ targets look alike to every algorithm, which is the source of the lower bound.\sidenote{Uniformity is what makes the $\binom dk$ targets an orthonormal family; under a biased product distribution the characters are no longer orthogonal, and \cref{ex:parity:biased} shows how much this helps.} Exact parity labels remove every low-degree correlation: the target has Fourier weight only at the single set $S$. The monomial activation makes the correlation of one neuron with the target an explicit polynomial, \cref{lem:parity:monomial} below. The symmetric initialization $w_0\propto\mathbf 1_d$ removes the luck of a random start, so that what remains is the deterministic signal. What a reader would have had to notice to invent the model is that the population gradient of a neuron is itself a Fourier coefficient: by the chain rule, $\partial_{w_i}\E[y\phi(\ip{w,x})]=\E[\chi_{S\triangle\{i\}}(x)\,\phi'(\ip{w,x})]$, the Fourier coefficient of $\phi'(\ip{w,x})$ at the set $S$ with $i$ removed or added. The question of how much a gradient knows becomes a question about the Fourier spectrum of one function, a theme that \cref{ch:grokking} develops on the group $\Z_p$.

\begin{dictionary}
the set $S$, $\abs S=k$ & the relevant features of the input\\
parity $\chi_S$ & a target with no low-degree footprint: every statistic of degree $<k$ is blind to it\\
monomial neuron $\ip{w,x}^k$ & one hidden unit; its degree is the lowest that can see the target\\
correlation $\Psi(w)$ & minus the loss; the network's agreement with the labels\\
weights $w_i$, $i\in S$ & the progress measure: how much the unit attends to the relevant inputs\\
effective overlap $\rho$ & the alignment of the unit with the target, on a scale where the loss is $1-\rho^k$\\
symmetric start $w_0\propto\mathbf 1_d$ & an initialization with no luck in it\\
statistical query of tolerance $\tau$ & an average over a minibatch of about $\tau^{-2}$ samples\\
loss plateau, $\rho$ near $0$ & the long flat stretch of \cref{fig:parity:mlp}\\
growth of $\log\rho$ during the plateau & hidden progress\\
\end{dictionary}

\section{The solution: what a gradient can see}\label{sec:parity:signal}

Everything rests on one computation, which says that the correlation of a monomial neuron with a parity counts the ways to match the factors of the monomial with the coordinates of $S$.

\begin{lemma}[Monomials against characters]\label[lemma]{lem:parity:monomial}
Let $T\subseteq[d]$, let $j\ge0$ be an integer with $j\le\abs T$, and let $w\in\R^d$. Then
\[
  \E\bigl[\chi_T(x)\ip{w,x}^j\bigr]=\begin{cases}j!\prod_{i\in T}w_i,&\abs T=j,\\0,&\abs T>j.\end{cases}
\]
\end{lemma}

\begin{proof}
Expand the power: $\ip{w,x}^j=\sum w_{i_1}\cdots w_{i_j}\,x_{i_1}\cdots x_{i_j}$, the sum over all sequences $(i_1,\dots,i_j)$ in $[d]^j$. Because $x_i^2=1$, the product $\chi_T(x)x_{i_1}\cdots x_{i_j}$ equals $\chi_U(x)$, where $U$ is the set of indices that occur an odd number of times in the list made of the elements of $T$ and of $i_1,\dots,i_j$; and $\E\chi_U(x)=1$ if $U=\emptyset$ and $0$ otherwise, by \cref{thm:toolkit:fourier-cube}. So a sequence contributes only if every element of $T$ occurs in it an odd number of times, and every other index an even number of times. Such a sequence has at least $\abs T$ terms. If $\abs T>j$, no sequence qualifies and the expectation is $0$. If $\abs T=j$, every element of $T$ occurs exactly once and nothing else occurs, so the qualifying sequences are the $j!$ orderings of $T$, each contributing $\prod_{i\in T}w_i$.
\end{proof}

\begin{theorem}[The exact signal of a monomial neuron]\label{thm:parity:signal}
For every $w\in\R^d$,
\[
  \Psi(w)=k!\prod_{i\in S}w_i,\qquad \partial_{w_i}\Psi(w)=\begin{cases}k!\prod_{j\in S\setminus\{i\}}w_j,&i\in S,\\0,&i\notin S.\end{cases}
\]
In particular, at the symmetric start $w_0=d^{-1/2}\mathbf 1_d$ the population gradient equals $k!\,d^{-(k-1)/2}$ on every coordinate of $S$ and exactly $0$ on every other coordinate.
\end{theorem}

\begin{proof}
The formula for $\Psi$ is \cref{lem:parity:monomial} with $T=S$ and $j=k$. Differentiating the product gives the gradient. For the second statement differentiate under the expectation, which is a finite sum: $\partial_{w_i}\Psi(w)=k\,\E[\chi_S(x)x_i\ip{w,x}^{k-1}]=k\,\E[\chi_{S\triangle\{i\}}(x)\ip{w,x}^{k-1}]$, where $\triangle$ is the symmetric difference. For $i\notin S$ the set $S\triangle\{i\}=S\cup\{i\}$ has $k+1>k-1$ elements, so \cref{lem:parity:monomial} gives $0$; for $i\in S$ it has $k-1$ elements and gives $k\,(k-1)!\prod_{j\in S\setminus\{i\}}w_j$, consistent with the first formula. At $w_0$ every product of $k-1$ coordinates is $d^{-(k-1)/2}$.
\end{proof}

Two features of the theorem drive the rest of the chapter. The gradient off $S$ vanishes identically, not only at $w_0$: in the population, the irrelevant weights of a monomial neuron never move. And the gradient on $S$ is small, of order $d^{-(k-1)/2}$ at a start of unit norm, because it is a product of $k-1$ small weights. The first feature says that the population gradient identifies $S$ perfectly; the second says that a sample estimate of it is drowned in noise unless the sample is large. Before quantifying the second, we check that the first is not an artifact of the monomial.

\begin{proposition}[The Fourier gap of a ReLU neuron]\label[proposition]{prop:parity:gap}
Let $d$ be odd, let $k$ be even with $2\le k<d$, and let $\alpha>0$. The correlation $\Phi(w)=\E[\chi_S(x)\relu(\ip{w,x})]$ is differentiable at $w=\alpha\mathbf 1_d$, and its partial derivative there is $\frac12\widehat{\parityMaj}_d(S\setminus\{i\})$ for $i\in S$ and $\frac12\widehat{\parityMaj}_d(S\cup\{i\})$ for $i\notin S$, where $\parityMaj_d(x)=\sgn(x_1+\dots+x_d)$. These numbers have opposite signs, and the ratio of their absolute values, off $S$ to on $S$, is exactly $(k-1)/(d-k)$. For odd $k$ both vanish.
\end{proposition}

\begin{proof}
Since $d$ is odd, $\ip{\alpha\mathbf 1_d,x}=\alpha(x_1+\dots+x_d)$ is never $0$, so near $\alpha\mathbf 1_d$ each of the finitely many terms of the expectation is a linear function of $w$, and $\partial_{w_i}\Phi(\alpha\mathbf 1_d)=\E[\chi_S(x)x_i\ind\{\ip{\mathbf 1_d,x}>0\}]$. Write $\ind\{\ip{\mathbf 1_d,x}>0\}=\frac12(1+\parityMaj_d(x))$; the set $S\triangle\{i\}$ is not empty, so the constant contributes nothing and the derivative is $\frac12\widehat{\parityMaj}_d(S\triangle\{i\})$. For the values we use the formula for the Fourier coefficients of majority \parencite[Chapter~5]{ODonnell2014}: they vanish on sets of even size, and on a set of odd size $j$,
\[
  \widehat{\parityMaj}_d(T)=(-1)^{(j-1)/2}\,\frac{\binom{(d-1)/2}{(j-1)/2}}{\binom{d-1}{j-1}}\cdot\frac{2}{2^d}\binom{d-1}{(d-1)/2}.
\]
For odd $k$ the sets $S\setminus\{i\}$ and $S\cup\{i\}$ have even sizes, so both derivatives vanish. For even $k$ the sizes are $k-1$ and $k+1$, consecutive odd numbers, so the signs differ. With $\nu=(d-1)/2$, the ratio of absolute values is
\[
  \frac{\binom{\nu}{k/2}}{\binom{\nu}{k/2-1}}\cdot\frac{\binom{d-1}{k-2}}{\binom{d-1}{k}}
  =\frac{\nu-k/2+1}{k/2}\cdot\frac{k(k-1)}{(d-k+1)(d-k)}=\frac{k-1}{d-k},
\]
because $\nu-k/2+1=(d-k+1)/2$.\qedhere
\end{proof}

The ReLU neuron at a symmetric start therefore also separates $S$ from its complement, but with a finite \term{Fourier gap}\index{Fourier gap}\index{majority function}: its gradient off $S$ is smaller by the factor $(k-1)/(d-k)$, not zero. By Stirling's formula the coefficient at level $k-1$ is of order $d^{-(k-1)/2}$, the same order as the monomial's signal. The gradient of a ReLU neuron, read as a Fourier coefficient of the threshold function $\ind\{\ip{w,x}>0\}$, is the mechanism behind \textcite{BarakEtAl2022}'s explanation of hidden progress; the parity of $k$ and the bias matter, since for odd $k$ the weights of a symmetric unit without bias receive no signal at all (\cref{ex:parity:majority}).\sidenote{The bias does receive signal for odd $k$: its derivative is $\frac12\widehat{\parityMaj}_d(S)$, nonzero because $\abs S$ is odd. A symmetric unit learns that the parity is correlated with its threshold, not which coordinates matter.}

\subsection{The statistical-query lower bound}\label{sec:parity:sq}

A gradient estimated from data is a list of averages. The model of computation that keeps only this feature is due to \textcite{Kearns1998}.

\begin{definition}[Statistical queries]\label[definition]{def:parity:sq}
A \term{statistical query}\index{statistical query} is a function $q\colon\{-1,1\}^d\times\{-1,1\}\to\R$ with $\E[q(x,1)^2]\le1$ and $\E[q(x,-1)^2]\le1$; every $q$ with values in $[-1,1]$ qualifies. An \emph{oracle of tolerance} $\tau>0$ for the target $\chi_S$ answers each query with some number $v$ satisfying $\abs{v-\E q(x,\chi_S(x))}\le\tau$. A \emph{statistical-query learner} chooses its queries one at a time, each depending on the previous answers and on internal randomness, and finally outputs a set $\hat S$.
\end{definition}

An average of $q(x_s,y_s)$ over a minibatch of $b$ fresh samples is within about $b^{-1/2}$ of $\E q(x,y)$; a query of tolerance $\tau$ is the idealization of a minibatch of size $\tau^{-2}$, with the sampling error replaced by an adversarial one.

\begin{theorem}[Statistical-query lower bound for sparse parities]\label{thm:parity:sq}
Let $1\le k\le d$, $M=\binom dk$, and $0<\tau\le1$, and let $S$ be uniformly random among the subsets of $[d]$ of size $k$, independent of the learner's internal randomness. There is an oracle of tolerance $\tau$ such that every statistical-query learner that makes at most $Q$ queries outputs $\hat S=S$ with probability at most
\[
  \frac{Q\tau^{-2}+1}{M}.
\]
In particular, a learner that identifies every $S$ with certainty against every oracle of tolerance $\tau$ makes at least $(M-1)\tau^2$ queries.
\end{theorem}

\begin{proof}
The plan is to let the oracle answer every query by its label-independent part, which is correct for all but a few targets, and to count the targets for which it is not. For a query $q$ put
\[
  q_0(x)=\tfrac12\bigl(q(x,1)+q(x,-1)\bigr),\qquad q_1(x)=\tfrac12\bigl(q(x,1)-q(x,-1)\bigr).
\]
Since $\chi_S(x)\in\{-1,1\}$, checking the two cases gives $q(x,\chi_S(x))=q_0(x)+\chi_S(x)q_1(x)$, and therefore
\[
  \E q(x,\chi_S(x))=\E q_0(x)+\hat q_1(S).
\]
\emph{The oracle.} Given the target $S$ and a query $q$, the oracle answers $\E q_0(x)$ if $\abs{\hat q_1(S)}\le\tau$, and the exact value $\E q(x,\chi_S(x))$ otherwise. In both cases the answer is within $\tau$ of the truth.

\emph{Few targets are exposed by one query.} Let $B(q)$ be the set of sets $T$ of size $k$ with $\abs{\hat q_1(T)}>\tau$. Parseval's identity (\cref{thm:toolkit:fourier-cube}) and the inequality $(a-b)^2\le2a^2+2b^2$ give
\[
  \abs{B(q)}\,\tau^2\le\sum_{T}\hat q_1(T)^2=\E q_1(x)^2\le\tfrac12\bigl(\E q(x,1)^2+\E q(x,-1)^2\bigr)\le1,
\]
so $\abs{B(q)}\le\tau^{-2}$.

\emph{Counting.} Fix the learner's internal randomness. Run the learner against the responder that answers $\E q_0(x)$ to every query $q$; this responder ignores $S$, so the run produces queries $q^{(1)},\dots,q^{(Q')}$ with $Q'\le Q$ and an output $\hat S_0$ that do not depend on $S$. Let $B=B(q^{(1)})\cup\dots\cup B(q^{(Q')})$, a set of at most $Q\tau^{-2}$ targets. If $S\notin B$, the oracle for $S$ gives the same answer as the responder at each of these queries, so by induction on the number of queries the learner, facing the oracle, asks the same queries and outputs $\hat S_0$. Hence $\hat S=S$ forces $S\in B$ or $S=\hat S_0$, an event of probability at most $(\abs B+1)/M\le(Q\tau^{-2}+1)/M$ for uniformly random $S$. Averaging over the internal randomness proves the bound. A learner that succeeds with certainty has $1\le(Q\tau^{-2}+1)/M$, which is the last statement.
\end{proof}

The argument goes back to the beginnings of the model: Kearns observed that parities cannot be learned\index{statistical query!lower bound} from statistical queries in polynomial time, and \textcite{BlumEtAl1994} turned the counting of pairwise orthogonal functions into a characterization of weak statistical-query learning by what is now called the statistical-query dimension.\sidenote{For the class of all $2^d$ parities the counting gives $Q\tau^{-2}\gtrsim2^d$; restricting to sets of size $k$ replaces $2^d$ by $\binom dk$.} The same counting bounds learners that only need a weak hypothesis (\cref{ex:parity:weak}).

\begin{remark}[Gradient methods are statistical algorithms]\label[remark]{rem:parity:gd}
For a model $f(x;\theta)$ with $p$ parameters, the population gradient of the correlation loss has coordinates $-\E[y\,\partial_{\theta_j}f(x;\theta)]$, and each is a statistical query after division by $B_j=(\E[\partial_{\theta_j}f(x;\theta)^2])^{1/2}$; for the square loss the gradient differs by $\E[f\,\partial_{\theta_j}f]$, which does not involve the label. So $t$ steps of gradient descent whose gradient coordinates are accurate to $\tau B_j$ are a statistical-query algorithm with $Q=tp$ queries, and if the errors are adversarial, \cref{thm:parity:sq} requires $tp\,\tau^{-2}\gtrsim\binom dk$. With $\tau\approx b^{-1/2}$ for minibatches of size $b$, this reads $tpb\gtrsim\binom dk$: the total computation of a gradient method, parameters times samples, is at least of order $d^k/k!$. This translation is a heuristic, not a theorem about SGD, because sampling errors are random rather than adversarial. The difference can be decisive: \textcite{AbbeSandon2020} showed that SGD with small batches on suitably constructed networks can emulate every polynomial-time learning algorithm, including those for parities that no statistical-query algorithm matches, and proved that this power disappears when gradients are averaged over large batches.
\end{remark}

\section{The solution: one step, and a plateau with progress}\label{sec:parity:upper}

We now show that a single gradient step meets the lower bound up to a constant and a logarithm. The noise of the empirical gradient is controlled by the moments of a sum of independent signs.

\begin{lemma}[Moments of a sum of signs]\label[lemma]{lem:parity:moments}
Let $s=x_1+\dots+x_d$. For every integer $j\ge1$, $\E s^{2j}\le(2j-1)!!\,d^j$, where $(2j-1)!!=1\cdot3\cdots(2j-1)$.
\end{lemma}

\begin{proof}
Expanding, $\E s^{2j}$ is the number of sequences $(i_1,\dots,i_{2j})\in[d]^{2j}$ in which every index occurs an even number of times, because the expectation of a product of signs is $1$ for such sequences and $0$ for the others. Each such sequence arises from at least one pair of a perfect matching of the $2j$ positions and an assignment of one index to each of the $j$ matched pairs: match the occurrences of each index among themselves. There are $(2j-1)!!$ perfect matchings and $d^j$ assignments, which bounds the number of sequences.
\end{proof}

The bound is the corresponding Gaussian moment; for $j=2$ the exact value is $\E s^4=3d^2-2d$. Write the empirical gradient of the correlation at the start $w_0=\mathbf 1_d$, scaled for convenience (the gradient is homogeneous of degree $k-1$ in $w$, so the scale of the start does not matter):
\begin{equation}\label{eq:parity:Ghat}
  \hat G_i=\frac1n\sum_{s=1}^nZ_i(x_s),\qquad Z_i(x)=k\,\chi_S(x)\,\ip{\mathbf 1_d,x}^{k-1}x_i,\qquad i\in[d].
\end{equation}

\begin{theorem}[One gradient step finds the support]\label{thm:parity:onestep}
Let $2\le k\le d$ and $0<\delta<1$, and let $x_1,\dots,x_n$ be independent and uniform on $\{-1,1\}^d$ with labels $\chi_S(x_s)$. If
\[
  n\ge C_k\,d^{k-1}\log\frac{2d}\delta,\qquad C_k=\frac{8k\bigl(k\,(2k-3)!!+k!/3\bigr)}{(k!)^2},
\]
then with probability at least $1-\delta$ every $i\in S$ has $\hat G_i>k!/2$ and every $i\notin S$ has $\hat G_i<k!/2$. In particular the $k$ largest coordinates of the empirical gradient are exactly the coordinates of $S$.
\end{theorem}

\begin{proof}
The plan is to apply Bernstein's inequality to each coordinate and a union bound over the $d$ coordinates. Fix $i$. By \cref{thm:parity:signal} with $w=\mathbf 1_d$, $\E Z_i=\partial_{w_i}\Psi(\mathbf 1_d)$ equals $k!$ for $i\in S$ and $0$ for $i\notin S$. Three bounds on $Z_i$ are needed.
\begin{itemize}
\item \emph{Range.} $\abs{\ip{\mathbf 1_d,x}}\le d$, so $\abs{Z_i}\le kd^{k-1}$; also $k!=k\,(k-1)!\le kd^{k-1}$ because $k\le d$. Hence $\abs{Z_i-\E Z_i}\le2kd^{k-1}$.
\item \emph{Variance.} Since $\chi_S^2=x_i^2=1$, \cref{lem:parity:moments} with $j=k-1\ge1$ gives $\Var Z_i\le\E Z_i^2=k^2\,\E\ip{\mathbf 1_d,x}^{2k-2}\le k^2(2k-3)!!\,d^{k-1}$.
\item \emph{Deviation.} Bernstein's inequality (\cref{thm:toolkit:bernstein}) for the $n$ independent centered variables $Z_i(x_s)-\E Z_i$, with $c=2kd^{k-1}$, $v=nk^2(2k-3)!!\,d^{k-1}$, and threshold $nk!/2$, gives
\[
  \Pr\Bigl(\abs{\hat G_i-\E Z_i}\ge\frac{k!}2\Bigr)\le2\exp\Bigl(-\frac{n^2(k!)^2/4}{2\bigl(v+c\,nk!/6\bigr)}\Bigr),
\]
and inserting $v$ and $c$, the exponent equals $-n(k!)^2/\bigl(8d^{k-1}(k^2(2k-3)!!+k\,k!/3)\bigr)=-n/(C_kd^{k-1})$.
\end{itemize}
By the hypothesis on $n$, the right side is at most $\delta/d$. By the union bound, with probability at least $1-\delta$ every coordinate satisfies $\abs{\hat G_i-\E Z_i}<k!/2$, which is the claim.
\end{proof}

For $k=3$ the constant is $C_3=22/3\approx7.3$.\sidenote{The constant is not optimized. Bernstein's inequality is used because the range of $Z_i$ and its standard deviation times $\sqrt n$ are of the same order $d^{k-1}$ at the sample size of the theorem; Chebyshev's inequality would cost a factor $d$.} The theorem uses a nonlinear step after the gradient, a threshold or a choice of the $k$ largest coordinates. A plain gradient step followed by a linear readout would need more samples: the noise of about $(d^{k-1}/n)^{1/2}$ on each of the $d-k$ irrelevant coordinates adds up, in $\ip{w_1,x}$, to a perturbation of order $(d^k/n)^{1/2}$, so about $d^k$ samples would be needed for the noise to be small compared with the signal (\cref{ex:parity:linear}). With the support found, a small second layer finishes the job.

\begin{corollary}[A network that computes the parity]\label[corollary]{cor:parity:network}
Let $\hat S\subseteq[d]$ have $k$ elements and $s(x)=\sum_{i\in\hat S}x_i$. Then for every $x\in\{-1,1\}^d$,
\[
  \chi_{\hat S}(x)=(-1)^k\Bigl[\relu\bigl(s(x)+k+1\bigr)+\sum_{j=1}^k(-1)^j\,4j\,\relu\bigl(s(x)+k+1-2j\bigr)\Bigr],
\]
a network of $k+1$ ReLU units with input weights $\mathbf 1_{\hat S}$. Consequently, under the hypotheses of \cref{thm:parity:onestep}, one gradient step at $w_0=\mathbf 1_d$, the choice of the $k$ largest coordinates, and this readout output a network equal to $\chi_S$ with probability at least $1-\delta$.
\end{corollary}

\begin{proof}
The value $s(x)$ lies in $\{-k+2i:0\le i\le k\}$, and at $s=-k+2i$ the parity is $(-1)^{k-i}$, since $i$ coordinates equal $+1$ and $k-i$ equal $-1$. Put $c_0=1$ and $c_j=(-1)^j4j$ for $j\ge1$. At $s=-k+2i$ the unit $\relu(s+k+1-2j)$ equals $2(i-j)+1$ for $j\le i$ and $0$ for $j>i$, so the bracket equals $A_i=\sum_{j=0}^ic_j\,(2(i-j)+1)$, and we must show $A_i=(-1)^i$. The partial sums of the coefficients are $P_i=\sum_{j=0}^ic_j=1+4\sum_{j=1}^i(-1)^jj=(-1)^i(2i+1)$, as one checks separately for even and odd $i$. Hence, for $i\ge1$,
\[
  A_i-A_{i-1}=2P_{i-1}+c_i=2(-1)^{i-1}(2i-1)+(-1)^i4i=2(-1)^i,
\]
and since $A_0=1$, $A_i=1+2\sum_{l=1}^i(-1)^l=(-1)^i$. The last statement combines this identity with \cref{thm:parity:onestep}.
\end{proof}

\paragraph{How close to the limit?} Read the $d$ coordinates of the empirical gradient as statistical queries. Normalized to unit second moment, the query for coordinate $i$ is $q_i(x,y)=y\,k\ip{\mathbf 1_d,x}^{k-1}x_i/\beta$ with $\beta^2=k^2(2k-3)!!\,d^{k-1}$, its expectation is $k!/\beta$ on $S$ and $0$ off $S$, and answers of every tolerance below $k!/(2\beta)$ suffice to find $S$. So the one-step algorithm is a statistical-query algorithm with $Q=d$ and, as $\tau$ approaches $k!/(2\beta)$,
\[
  Q\tau^{-2}\to\frac{4k^2(2k-3)!!}{(k!)^2}\,d^k=\frac{4k\,(2k-3)!!}{(k-1)!}\cdot\frac{d^k}{k!},
\]
while \cref{thm:parity:sq} requires $Q\tau^{-2}\ge\binom dk-1\approx d^k/k!$. The algorithm is within the factor $4k(2k-3)!!/(k-1)!$ of the lower bound, $18$ for $k=3$, and its $n\approx d^{k-1}\log d$ samples times $d$ coordinates is the computation $d^k\log d$. This is the precise sense in which gradient methods can learn parities ``near the computational limit'': the limit is a product of samples and queries, and a single step on the right neuron saturates it.

\subsection{The plateau and the hidden progress}\label{sec:parity:hidden}

The one-step algorithm compresses the search into one large average. Training by small steps spreads it over time, and the loss curve acquires a plateau\index{plateau}. We study the population flow of the monomial neuron on the unit sphere, where the loss is bounded. Two constants recur:
\begin{equation}\label{eq:parity:constants}
  c_k=\max_{\norm w=1}\Psi(w)=\frac{k!}{k^{k/2}},\qquad\kappa_k=k\,c_k .
\end{equation}
The maximum is attained when $w_i^2=1/k$ on $S$ and $w_i=0$ off $S$, by the inequality of arithmetic and geometric means: $\prod_{i\in S}w_i^2\le(\frac1k\sum_{i\in S}w_i^2)^k\le k^{-k}$. For $\Psi(w)>0$ define the \term{effective overlap}\index{effective overlap}
\[
  \rho(w)=\bigl(\Psi(w)/c_k\bigr)^{1/k}\in(0,1],
\]
the overlap with $\mathbf 1_S/\sqrt k$ of the vector with equal weights on $S$ that has the same correlation as $w$. In these units the normalized loss $(\Loss(w)-\min\Loss)/(0-\min\Loss)$ is $1-\rho^k$, so the loss sees $\rho^k$ while the weights carry $\rho$. A weight vector is \emph{balanced} if $\abs{w_i}$ is the same for all $i\in S$. The symmetric start $w_0=d^{-1/2}\mathbf 1_d$ is balanced, with $\rho_0=\sqrt{k/d}$; a uniformly random start has $\Psi(w_0)>0$ with probability $\frac12$, and then $\rho_0$ is of order $d^{-1/2}$.

\begin{theorem}[Hidden progress of the population flow]\label{thm:parity:hidden}
Let $k\ge3$ and $w_0\in\Sph^{d-1}$ with $\Psi(w_0)>0$, and let $w(t)$ solve the spherical gradient flow of the correlation,
\begin{equation}\label{eq:parity:flow}
  \dot w=\nabla\Psi(w)-\ip{w,\nabla\Psi(w)}\,w,\qquad w(0)=w_0 .
\end{equation}
\begin{enumerate}
\item[\textup{(a)}] The solution exists for all $t\ge0$ and stays on the sphere, $\Psi(w(t))$ is nondecreasing, and
\[
  \frac{\dd}{\dd t}w_i^2=2\Psi\,(1-kw_i^2)\quad(i\in S),\qquad\frac{\dd}{\dd t}w_i^2=-2k\Psi\,w_i^2\quad(i\notin S).
\]
\item[\textup{(b)}] The effective overlap satisfies
\begin{equation}\label{eq:parity:rhoineq}
  \dot\rho\ge\kappa_k\,\rho^{k-1}(1-\rho^2),
\end{equation}
with equality at all times if $w_0$ is balanced.
\item[\textup{(c)}] Let $t_a$ be the first time with $\rho\ge a$. If $\rho_0\le a<b\le\frac12$, then
\[
  t_b-t_a\le\frac{4\,a^{2-k}}{3(k-2)\kappa_k},
\]
and if $w_0$ is balanced, then also $t_b-t_a\ge(a^{2-k}-b^{2-k})/((k-2)\kappa_k)$.
\item[\textup{(d)}] On $[0,t_{1/2}]$ the logarithm of the effective overlap grows at least linearly, $\frac{\dd}{\dd t}\log\rho\ge\frac34\kappa_k\rho_0^{k-2}$.
\end{enumerate}
\end{theorem}

\begin{proof}
The plan is to compute the evolution of each squared weight, to derive from it a closed differential inequality for $\Psi$, and to integrate that inequality. Write $g=\nabla\Psi$. By \cref{thm:parity:signal}, $g_i=\Psi/w_i$ for $i\in S$ whenever $\Psi\ne0$, and $g_i=0$ for $i\notin S$; and since $\Psi$ is homogeneous of degree $k$, Euler's identity gives $\ip{w,g}=k\Psi$.

\emph{(a)} Read the right side of \eqref{eq:parity:flow} as a polynomial vector field on $\R^d$. Along its solutions $q=\norm w^2-1$ satisfies $\dot q=2\ip{w,g}(1-\norm w^2)=-2k\Psi q$, a linear equation, so $q$ stays $0$: the solution stays on the compact sphere and therefore exists for all time. On the sphere, $\dot\Psi=\ip{g,\dot w}=\norm g^2-\ip{w,g}^2\ge0$ by the Cauchy--Schwarz inequality, so $\Psi\ge\Psi(w_0)>0$ for all $t$, and no weight in $S$ vanishes. For $i\in S$, $\frac{\dd}{\dd t}w_i^2=2w_i(g_i-k\Psi w_i)=2\Psi(1-kw_i^2)$; for $i\notin S$, $\frac{\dd}{\dd t}w_i^2=-2k\Psi w_i^2$.

\emph{(b)} By (a) and the formula for $g$,
\[
  \dot\Psi=\norm g^2-\ip{w,g}^2=\Psi^2\sum_{i\in S}w_i^{-2}-k^2\Psi^2 .
\]
The inequality of arithmetic and geometric means bounds the sum below: $\sum_{i\in S}w_i^{-2}\ge k\bigl(\prod_{i\in S}w_i^{-2}\bigr)^{1/k}=k\,(k!/\Psi)^{2/k}$, using $\prod_{i\in S}w_i^2=(\Psi/k!)^2$. Substitute $\Psi=c_k\rho^k$, so that $\dot\Psi=kc_k\rho^{k-1}\dot\rho$, and use $(k!)^{2/k}c_k^{-2/k}=k$:
\[
  kc_k\rho^{k-1}\dot\rho\ge k\,(k!)^{2/k}(c_k\rho^k)^{2-2/k}-k^2c_k^2\rho^{2k}=k^2c_k^2\rho^{2k-2}(1-\rho^2).
\]
Dividing by $kc_k\rho^{k-1}>0$ gives \eqref{eq:parity:rhoineq}. Equality holds in the inequality of means exactly when all $w_i^2$, $i\in S$, are equal. By (a), $\frac{\dd}{\dd t}(w_i^2-w_j^2)=-2k\Psi(w_i^2-w_j^2)$ for $i,j\in S$, a linear equation, so a balanced start stays balanced and \eqref{eq:parity:rhoineq} is then an equality.

\emph{(c)} Let $h(s)=\kappa_ks^{k-1}(1-s^2)$, which is positive on $(0,1)$. By (b), $\rho$ is nondecreasing, and it is strictly increasing while $\rho<1$. Put $H(r)=\int_a^r\dd s/h(s)$. Suppose first that $t_a$ is finite, which holds for $a=\rho_0$ since then $t_a=0$. On $[t_a,t_b]$ the chain rule and (b) give $\frac{\dd}{\dd t}H(\rho(t))=\dot\rho/h(\rho)\ge1$, so $H(\rho(t))\ge t-t_a$ there; since $H(\rho(t))<H(b)$ for $t<t_b$, the time $t_b$ is finite and $t_b-t_a\le H(b)$. Applied with $a=\rho_0$, this shows that $t_c$ is finite for every $c<1$, so the supposition holds for every $a$. On $[a,b]\subseteq(0,\frac12]$ we have $1-s^2\ge\frac34$, hence
\[
  H(b)\le\frac4{3\kappa_k}\int_a^bs^{1-k}\,\dd s\le\frac{4\,a^{2-k}}{3(k-2)\kappa_k}.
\]
For a balanced start $\frac{\dd}{\dd t}H(\rho(t))=1$, so $t_b-t_a=H(b)\ge\kappa_k^{-1}\int_a^bs^{1-k}\dd s=(a^{2-k}-b^{2-k})/((k-2)\kappa_k)$, using $1-s^2\le1$.

\emph{(d)} By (b), $\frac{\dd}{\dd t}\log\rho=\dot\rho/\rho\ge\kappa_k\rho^{k-2}(1-\rho^2)$, and on $[0,t_{1/2}]$ we have $\rho\ge\rho_0$ and $1-\rho^2\ge\frac34$.
\end{proof}

The theorem is the Boolean version of the search phase of \cref{ch:singleindex}: for a balanced start, $\rho$ obeys exactly the equation $\dot r=(1-r^2)\Psi'(r)$ of that chapter with $\Psi(r)=c_kr^k$. It adds two things. An unbalanced start is never slower than a balanced one with the same correlation, and the imbalance between the relevant weights decays; the irrelevant weights decay at the rate $2k\Psi\ge2k\Psi(w_0)$, so the neuron forgets the coordinates outside $S$ exponentially fast once the correlation is positive. Without the sphere, (a) becomes the statement that $w_i^2-w_j^2$ is conserved for $i,j\in S$, the balancedness of the deep linear networks of \cref{ch:deeplinear}: the monomial neuron on a parity is a depth-$k$ product of scalars (\cref{ex:parity:unconstrained}).

\begin{corollary}[Plateau and drop]\label[corollary]{cor:parity:plateau}
Let $k\ge3$, let $w_0=d^{-1/2}\mathbf 1_d$, and let $0<\eps<2^{-k}$ with $\eps^{1/k}>\rho_0=\sqrt{k/d}$. The normalized loss $1-\rho^k$ stays above $1-\eps$ until the time
\[
  t_{\eps^{1/k}}\ge\frac{(d/k)^{(k-2)/2}-\eps^{-(k-2)/k}}{(k-2)\kappa_k},
\]
and it reaches $1-2^{-k}$ at most $\frac{4}{3(k-2)\kappa_k}\eps^{-(k-2)/k}$ later, a time that does not depend on $d$. During the plateau $[0,t_{\eps^{1/k}}]$ the effective overlap grows by the factor $\eps^{1/k}\sqrt{d/k}$.
\end{corollary}

\begin{proof}
Apply \cref{thm:parity:hidden}(c) twice: with $a=\rho_0$ and $b=\eps^{1/k}$ for the lower bound, using $t_{\rho_0}=0$ and $\rho_0^{2-k}=(d/k)^{(k-2)/2}$; and with $a=\eps^{1/k}$ and $b=\frac12$ for the upper bound. The last statement is the definition of $t_{\eps^{1/k}}$.
\end{proof}

As $d\to\infty$ with $k$ and $\eps$ fixed, the plateau lasts a time of order $d^{(k-2)/2}$ and the drop a time of order one: the loss curve becomes a step. Yet nothing abrupt happens to the weights. The relevant weights are proportional to $\rho$, and by (d) the logarithm of $\rho$ climbs at a rate bounded below from the first instant; most of the climb in $\log\rho$, a fraction $1-O(1/\log d)$ of it, happens while the loss is within $\eps$ of its initial value. The loss is flat because it sees $\rho^k$, the $k$th power of a quantity that grows steadily. This is the mechanism of \textcite{SchaefferMirandaKoyejo2023} in miniature: a nonlinear readout of continuous progress produces a jump.

\begin{heuristic}[Online SGD needs about $d^{k-1}$ samples]
Replace the population gradient by the gradient of one fresh sample, $G_t=-y_t\,k\ip{w_t,x_t}^{k-1}P_{w_t}^\perp x_t$ with $P_w^\perp=I_d-ww^\top$, and renormalize after each step, as in \cref{ch:singleindex}. For a unit vector $w$, $\ip{w,x}$ has bounded moments by \cref{thm:toolkit:bonami}, so the component of $G_t$ along $\mathbf 1_S/\sqrt k$ has variance of order one and $\norm{G_t}^2$ has mean of order $d$. These are the two noise terms of the step-size heuristic of \cref{ch:singleindex}, with the same orders, and the drift is the right side of \eqref{eq:parity:rhoineq}. The heuristic of that chapter then applies verbatim: the step size must be of order $d^{-k/2}$ to keep the noise below the drift at $\rho\approx d^{-1/2}$, and the number of samples is the escape time $d^{(k-2)/2}$ divided by the step, of order $d^{k-1}$. This matches the one-step count of \cref{thm:parity:onestep} up to the logarithm: spreading the average over many small steps costs no extra samples.
\end{heuristic}

\section{Reading the solution}\label{sec:parity:reading}

The solution has three numbers. The \emph{signal}: the population gradient at a start of unit norm is $k!\,d^{-(k-1)/2}$ on $S$ and zero elsewhere, while each coordinate of a single-sample gradient has standard deviation of order $d^0$ in the same units; one sample therefore carries a signal-to-noise ratio of order $d^{-(k-1)/2}$ per coordinate. The \emph{sample size}: averaging $n$ samples multiplies the ratio by $\sqrt n$, and finding the largest of $d$ coordinates needs a ratio of order $\sqrt{\log d}$, hence $n\asymp d^{k-1}\log d$. The \emph{product}: $d$ queries at this precision cost $d^k\log d$, against the lower bound $\binom dk$ of \cref{thm:parity:sq}. Kernel methods, which cannot adapt to $S$, pay instead with the number of features needed to represent all degree-$k$ monomials, also of order $d^k$ (\cref{ch:lazy}); the gain of gradient training over kernels is in memory, not in computation.

\begin{figure}[t]
\centering
\begin{tikzpicture}
\begin{axis}[hplot,xmin=0,xmax=4.1,ymin=0,ymax=1.02,xlabel={$n/(d^2\ln d)$},ylabel={$\Pr(\text{top-}3\text{ coordinates}=S)$},legend pos=south east]
\addplot[theory] table[x=a,y=th16] {figures/data/09-parity-onestep.dat};
\addplot[theory2] table[x=a,y=th32] {figures/data/09-parity-onestep.dat};
\addplot[theory3] table[x=a,y=th64] {figures/data/09-parity-onestep.dat};
\addplot[simulation2,mark=triangle*] table[x=a,y=sim16] {figures/data/09-parity-onestep.dat};
\addplot[simulation2] table[x=a,y=sim32] {figures/data/09-parity-onestep.dat};
\addplot[simulation] table[x=a,y=sim64] {figures/data/09-parity-onestep.dat};
\legend{Gaussian limit $d=16$,Gaussian limit $d=32$,Gaussian limit $d=64$,simulation $d=16$,simulation $d=32$,simulation $d=64$}
\end{axis}
\end{tikzpicture}
\caption{One gradient step of the monomial neuron $\ip{w,x}^3$ at $w_0=\mathbf 1_d$ on the parity $x_1x_2x_3$: the probability that the three largest coordinates of the empirical gradient \eqref{eq:parity:Ghat} are the relevant ones, against the number of samples in units of $d^2\ln d$. Points: $300$ independent data sets for each $(d,n)$ (seed $20261017$). Curves: the central-limit prediction, a Gaussian vector with the exact mean and covariance of $\hat G$, estimated with $20{,}000$ draws. \Cref{thm:parity:onestep} guarantees success probability $1-\delta$ from $n=7.3\,d^2\ln(2d/\delta)$.}\label{fig:parity:onestep}
\end{figure}

\Cref{fig:parity:onestep} tests the sample size. The points are simulations of the one-step algorithm for $k=3$ and $d=16$, $32$, and $64$; the curves are the Gaussian approximation with the exact mean and covariance of $\hat G$, which uses the exact moments $\E\ip{\mathbf 1_d,x}^4=3d^2-2d$ and $\E[\ip{\mathbf 1_d,x}^4x_1x_2]=12d-16$. The two agree to within about two Monte Carlo standard errors (one standard error is about $0.03$), and the curves for the three dimensions are close when $n$ is measured in units of $d^2\ln d$, as the signal-to-noise count predicts. That count also gives the location: the largest of $d-k$ roughly independent centered Gaussians with standard deviation $\sigma=(k^2(2k-3)!!\,d^{k-1}/n)^{1/2}$ is about $\sigma\sqrt{2\ln d}$, and the smallest relevant coordinate is about $k!$, so success becomes likely near $n\approx2k^2(2k-3)!!\,d^{k-1}\ln d/(k!)^2$, which is $1.5\,d^2\ln d$ for $k=3$. The simulated success probability crosses $\frac12$ between about $1.2$ (for $d=16$) and $1.7$ (for $d=64$) in these units. The guarantee of \cref{thm:parity:onestep} at $\delta=0.05$ and $d=64$ asks for about $14\,d^2\ln d$ samples, four times the $3.5\,d^2\ln d$ at which the simulations succeed about $95\%$ of the time. Most of the factor comes from asking every coordinate to stay within $k!/2$ of its mean, half the gap, which costs a factor $4$ in the variance term; the range term of Bernstein's inequality adds about $20\%$ for $k=3$, and $\ln(2d/\delta)$ exceeds $\ln d$.

\begin{figure}[t]
\centering
\begin{tikzpicture}
\begin{groupplot}[group style={group size=2 by 1,horizontal sep=1.45cm},hplot,width=0.51\linewidth,height=0.47\linewidth,
  xlabel={samples$/d^2$},xmin=0,xmax=4.5]
\nextgroupplot[ymode=log,ymin=0.08,ymax=1.2,ylabel={effective overlap $\rho$},ytick={0.125,0.25,0.5,1},
  yticklabels={$\frac18$,$\frac14$,$\frac12$,$1$}]
\pgfplotsinvokeforeach{1,...,8}{\addplot[muted!70,line width=0.35pt] table[x=tau,y=r#1] {figures/data/09-parity-neuron.dat};}
\addplot[theory] table[x=tau,y=ode] {figures/data/09-parity-neuron.dat};
\nextgroupplot[ymin=0,ymax=1.05,ylabel={normalized loss $1-\rho^3$}]
\pgfplotsinvokeforeach{1,...,8}{\addplot[muted!70,line width=0.35pt] table[x=tau,y=l#1] {figures/data/09-parity-neuron.dat};}
\addplot[theory] table[x=tau,y=lode] {figures/data/09-parity-neuron.dat};
\end{groupplot}
\end{tikzpicture}
\caption{Eight runs of spherical online SGD for the monomial neuron $\ip{w,x}^3$ on the parity $x_1x_2x_3$ in dimension $d=256$, from the symmetric start $w_0=d^{-1/2}\mathbf 1_d$, with step size $\eta=0.05\,d^{-3/2}$ and one fresh sample per step (grey; seed $11$), and the population flow of \cref{thm:parity:hidden}, $\dot\rho=\eta\kappa_3\rho^2(1-\rho^2)$ per step (curve). Left: the effective overlap $\rho=(\Psi(w)/c_3)^{1/3}$, on a logarithmic scale. Right: the normalized population loss $1-\rho^3$.}\label{fig:parity:neuron}
\end{figure}

\Cref{fig:parity:neuron} tests the plateau. On the logarithmic scale of the left panel the effective overlap rises from the first step, along the population curve, while the loss in the right panel stays above $0.99$ until $\rho$ has doubled from its start $\rho_0=\sqrt{3/256}\approx0.108$; on the flow this takes $61\%$ of the time to reach $\rho=\frac12$, and the fraction tends to $1$ as $d$ grows. The prediction $t_{1/2}=\kappa_3^{-1}\bigl(\rho_0^{-1}-2+\operatorname{artanh}\frac12-\operatorname{artanh}\rho_0\bigr)$ (\cref{ex:parity:escape}) gives $2.77\,d^2$ samples at this step size, with no fitted constant. The SGD runs reach $\rho=\frac12$ between $0.99$ and $1.51$ times the prediction, with median $1.14$: the noise delays the escape slightly at this step size, as in \cref{ch:singleindex}, and the delay shrinks at smaller steps (\cref{ex:parity:stepsize}). After the escape every run descends in a short, nearly deterministic time.

Placed side by side, \cref{fig:parity:mlp,fig:parity:neuron} show what the model explains and what it does not. The plateau, its abrupt end, and the steady growth of a weight-based measure underneath a flat loss appear in both. But in the network the progress measure is flat for the first half of the plateau and grows only in the best few neurons, and the length of the plateau varies by a factor of two between seeds. Both features come from the random initialization, which the symmetric start of the model removes.

\section{What the model gets wrong}\label{sec:parity:wrong}

\paragraph{What was thrown away.}
Uniform inputs, which make every target look like every other and every character orthogonal to every other; a real input distribution is neither uniform nor a product. Exact parity labels, with all Fourier weight on one set of size $k$ and none below; natural targets have weight spread over many degrees. The monomial activation, chosen so that one neuron's correlation is an exact product; a ReLU unit has the finite Fourier gap of \cref{prop:parity:gap} and, for odd $k$ without bias, no gap at all. The symmetric initialization, which removes luck; with a random start the progress of \cref{fig:parity:mlp} concentrates in the neurons whose initial weights happen to favor $S$. Population dynamics or a single large step, which removes the interaction between noise and the slow escape; only a heuristic connects the theorems to online SGD. One neuron, so that width plays no role. And the correlation loss, which is unbounded below off the sphere.

\paragraph{Where it breaks.}
\emph{Natural tasks have footholds.} A parity is hard because it has no low-degree component, and this is the exception. \Textcite{AbbeEtAl2021} define the staircase property, under which each high-degree Fourier coefficient is reachable from lower-degree ones along a nested chain, and show that such functions are learned in polynomial time by layerwise training of suitable networks, with experiments showing that ResNets trained by SGD learn them too; \textcite{AbbeBoixAdseraMisiakiewicz2022} prove that a merged-staircase property is necessary and nearly sufficient for two-layer networks in the mean-field regime to learn sparse functions with a number of samples of order $d$. For such targets there is no long plateau before every step, and the lower bound of \cref{thm:parity:sq} does not apply. \emph{Non-uniform inputs.} Under a product distribution with $\E x_i=\mu\ne0$, the parity correlates with single coordinates, $\E[\chi_S(x)x_i]=\mu^{k-1}$ for $i\in S$ against $\mu^{k+1}$ for $i\notin S$, and a linear statistic finds $S$ from a number of samples that grows only logarithmically in $d$ (\cref{ex:parity:biased}). \Textcite{CornacchiaMossel2023} turn this into a model of curriculum learning, showing that drawing the training samples from two or more product distributions, rather than from the uniform one alone, substantially lowers the computational cost of learning $k$-parities by SGD, and \textcite{DanielyMalach2020} use non-uniform input distributions to separate two-layer networks trained by gradient descent from linear methods. \emph{Luck, width, and sparsity.} \Textcite{EdelmanEtAl2023} study the trade-offs between data, computation, width, and the luck of the initialization for offline sparse parity learning; they find that wide networks and sparse initialization improve sample efficiency, and they describe width as parallel search for neurons whose initialization is close to $S$. The symmetric model has no luck in it and cannot express these trade-offs. \emph{Statistical queries are not the end.} \Cref{rem:parity:gd} lists the gap between adversarial and random errors: SGD on specially built networks can emulate Gaussian elimination \parencite{AbbeSandon2020}. With label noise, on the other hand, parity learning is believed to be hard for every algorithm, and the best general algorithms, such as that of \textcite{BlumKalaiWasserman2003}, which runs in polynomial time when the parity depends on the first $O(\log d\log\log d)$ bits, are far from gradient methods. \emph{Architecture.} For transformers, \textcite{HahnRofin2024} argue that functions whose output is sensitive to many input bits, parity foremost, occupy isolated regions of parameter space, which biases training toward low-sensitivity functions; and \textcite{KimSuzuki2025} prove that a one-layer transformer learns parity in a single gradient update when trained with intermediate parities as supervision and teacher forcing, while without intermediate supervision finite-precision gradient methods need many iterations. Chain of thought turns one parity of degree $k$ into $k-1$ parities of degree two; the statistical-query cost of \cref{thm:parity:sq} applies to each link, not to the chain.

\paragraph{Fine structure.}
The successors repair one omission at a time. For two-layer networks, \textcite{BarakEtAl2022} prove that a network of width $2^{\Theta(k)}$, trained by SGD with batch size $d^{O(k)}$, learns $k$-sparse parities to error $\eps$ in $2^{O(k)}\poly(1/\eps)$ iterations, using the Fourier gap of the population gradient. \Textcite{KouEtAl2024} show that sign SGD on a two-layer network with $2^{\Theta(k)}$ neurons solves $k$-sparse parity with $\tilde O(d^{k-1})$ samples, matching the statistical-query lower bound in the sense of the product computed above. For $k=2$, the XOR problem, \textcite{Glasgow2024} proves that minibatch SGD on a two-layer ReLU network with logistic loss reaches small error with $d\polylog(d)$ samples, after a signal-finding phase in which neurons evolve independently, the random-start analogue of our hidden progress. Mean-field Langevin dynamics achieves a sample complexity in which $k$ is decoupled from the exponent of $d$ \parencite{SuzukiEtAl2023}. For sparse targets that combine monomials of several degrees, the leap complexity of \textcite{AbbeBoixAdseraMisiakiewicz2023} predicts a sequence of plateaus, one for each new set of coordinates, with lengths set by the size of the leap; this is the multi-index form of \cref{thm:parity:hidden} and the Boolean counterpart of \cref{ch:singleindex}. Open are sharp lower bounds for plain SGD on standard networks beyond the statistical-query heuristic, and a description of hidden progress from a random start: which quantity, computed from the weights at a given time, predicts when the drop will come.

\section{Model card}\label{sec:parity:card}

\modelcard{Loss and error flat for a long time, then an abrupt drop, with progress in the weights during the plateau; training time growing like $d^{O(k)}$ for $k$-sparse parities.}
{$x\sim\Unif\{-1,1\}^d$, $y=\chi_S(x)$ with $\abs S=k$; one neuron $\ip{w,x}^k$ on the correlation loss; statistical queries, one gradient step from $w_0\propto\mathbf 1_d$, and the spherical population flow.}
{$\Psi(w)=k!\prod_{i\in S}w_i$, gradient supported on $S$ (proved); every statistical-query learner needs $Q\tau^{-2}\ge\binom dk-1$ (proved); one step with $O(d^{k-1}\log(d/\delta))$ samples finds $S$ and $k+1$ ReLUs finish (proved); $\dot\rho\ge\kappa_k\rho^{k-1}(1-\rho^2)$ with equality for balanced starts (proved).}
{Sample size $d^{k-1}\log d$ for one step, about $1.5\,d^2\ln d$ at $k=3$; computation $d^k$; plateau of length $d^{(k-2)/2}$ and drop of length $O(1)$; steady growth of $\log\rho$ under a flat loss.}
{No luck: random starts concentrate progress in few neurons; no low-degree structure, which natural tasks have; uniform inputs; statistical queries can be beaten by SGD on special networks; architecture-specific effects for transformers.}
{Staircase and leap complexity; biased inputs and curricula; width, sparsity, and luck; sign SGD matching the SQ bound; chain of thought.}

This chapter is the Boolean partner of \cref{ch:singleindex}. There a single direction hides behind a link whose low Hermite coefficients vanish, here $k$ coordinates hide behind a character whose low-degree Fourier coefficients vanish; the information exponent becomes the degree $k$, the Hermite basis becomes the characters, and the escape equations coincide for balanced starts. Both chapters belong to the thread of escapes from small initialization that starts with the logistic escape of \cref{ch:deeplinear}: without the sphere, the monomial neuron is a depth-$k$ product of scalars whose relevant weights stay balanced, and the plateau of length $d^{(k-2)/2}$ is the higher-order saddle that also produces the plateaus of the committee machine (\cref{ch:committee}) and the sudden formation of induction heads (\cref{ch:induction}).

The statistical-query argument is Fourier analysis on the group $\{-1,1\}^d$; \cref{ch:grokking} uses Fourier analysis on $\Z_p$ to describe what a network computes after an equally abrupt transition, and its progress measures play the role of $\log\rho$ here. Against \cref{ch:lazy}, the chapter gives a second separation between feature learning and kernels: both pay about $d^k$ in computation, but a gradient method stores $d$ weights where a kernel method stores $d^k$ features.

\begin{exercises}
\item \consolidation (Sums of signs.) (a) Compute $\E s^4=3d^2-2d$ for $s=x_1+\dots+x_d$, and check \cref{lem:parity:moments} for $j=2$. (b) Show that $\E[s^4x_1x_2]=12d-16$, and conclude that two coordinates of the empirical gradient \eqref{eq:parity:Ghat} for $k=3$ have correlation of order $1/d$.\label{ex:parity:moments}

\item \consolidation (Majority.) (a) Verify $\parityMaj_3(x)=\frac12(x_1+x_2+x_3)-\frac12x_1x_2x_3$. (b) For $d=5$ and $k=2$, compute the partial derivatives of \cref{prop:parity:gap} directly from the $32$ points of the cube and check the ratio $(k-1)/(d-k)=\frac13$. (c) For odd $k\ge3$, show that the weights of the symmetric unit $\relu(\alpha\ip{\mathbf 1_d,x}+b)$ with $b=0$ receive zero population gradient, while its bias receives $\frac12\widehat{\parityMaj}_d(S)\ne0$.\label{ex:parity:majority}

\item \consolidation (Fewer units.) Show that $k$ ReLU units with input weights $\mathbf 1_S$ and an output bias suffice to compute $\chi_S$. \hint{interpolate the values $(-1)^{k-i}$ at $s=-k+2i$ by a zigzag with breakpoints at the interior nodes, and write the affine part with the unit $\relu(s+k)$.}\label{ex:parity:fewer}

\item \consolidation (Weak learners.) A weak learner outputs a function $h\colon\{-1,1\}^d\to[-1,1]$ and succeeds if $\E[h(x)\chi_S(x)]\ge\gamma$. Adapt the proof of \cref{thm:parity:sq} to show that for uniformly random $S$ the success probability is at most $(Q\tau^{-2}+\gamma^{-2})/\binom dk$.\label{ex:parity:weak}

\item \consolidation (The escape time for $k=3$.) For a balanced start with $\rho_0<\frac12$, solve \eqref{eq:parity:rhoineq} with equality and show $t_{1/2}=\kappa_3^{-1}\bigl(\rho_0^{-1}-2+\operatorname{artanh}\frac12-\operatorname{artanh}\rho_0\bigr)$. Evaluate it for $d=256$ and $\rho_0=\sqrt{3/d}$, and check the value $2.77\,d^2$ samples quoted in \cref{sec:parity:reading} for $\eta=0.05\,d^{-3/2}$.\label{ex:parity:escape}

\item \variation (Without the sphere.) Let $\dot w=\nabla\Psi(w)$ on $\R^d$. (a) Show that $w_i^2-w_j^2$ is conserved for $i,j\in S$ and that $w_i$ is constant for $i\notin S$. (b) For $w_i(0)=u_0>0$ on $S$ and $k\ge3$, show that $u(t)=\bigl(u_0^{2-k}-(k-2)k!\,t\bigr)^{-1/(k-2)}$, which blows up at $t^\star=u_0^{2-k}/((k-2)k!)$, and that the correlation $k!u^k$ stays below $\eps$ until a time $t^\star-O(\eps^{-(k-2)/k})$. (c) For $k=2$ show exponential growth, and compare with the logistic escape of \cref{ch:deeplinear}.\label{ex:parity:unconstrained}

\item \variation (Biased inputs.) Let the coordinates of $x\in\{-1,1\}^d$ be independent with $\E x_i=\mu\in(0,1)$. (a) Show $\E[\chi_S(x)x_i]=\mu^{k-1}$ for $i\in S$ and $\mu^{k+1}$ for $i\notin S$. (b) Using Hoeffding's inequality (\cref{thm:toolkit:hoeffding}), show that the $k$ largest coordinates of $\frac1n\sum_sy_sx_s$ are $S$ with probability at least $1-\delta$ once $n\ge8\log(2d/\delta)/(\mu^{2k-2}(1-\mu^2)^2)$. Why does \cref{thm:parity:sq} not apply?\label{ex:parity:biased}

\item \variation (Label noise.) Flip each label independently with probability $p<\frac12$. Show that every correlation in this chapter is multiplied by $1-2p$, that the sample size of \cref{thm:parity:onestep} grows by the factor $(1-2p)^{-2}$ up to the change in the range term, and that the statistical-query bound keeps its form with $\tau$ replaced by $\tau/(1-2p)$.\label{ex:parity:noise}

\item \variation (Gaussian inputs.) Let $x\sim\cN(0,I_d)$ and $y=\prod_{i\in S}x_i$. Use Isserlis' theorem to show that $\E[y\ip{w,x}^k]=k!\prod_{i\in S}w_i$ and that the gradient vanishes off $S$, as on the cube, and that $\E\ip{\mathbf 1_d,x}^{2k-2}=(2k-3)!!\,d^{k-1}$ exactly. Which conclusions of the chapter survive, and how does the target relate to the multi-index models of \cref{ch:singleindex}?\label{ex:parity:gauss}

\item \variation\hard (A plain step.) Let $w_1=\hat G$ be the empirical gradient \eqref{eq:parity:Ghat} itself, used without thresholding, and let $f(x)=\ip{w_1,x}/k!$. Show that $\E[(f(x)-s_S(x))^2]$, where $s_S=\sum_{i\in S}x_i$, is of order $d^k/n$, so that a readout of $f$ by a few ReLU units needs about $d^k$ samples, a factor $d/\log d$ more than \cref{thm:parity:onestep}. \hint{write $\hat G=k!\,\mathbf 1_S+\xi$; for a fresh input $x$, independent of the data, $\E\ip{\xi,x}^2=\E\norm\xi^2$, and each $\E\xi_i^2$ is the variance bound in the proof of \cref{thm:parity:onestep} divided by $n$.}\label{ex:parity:linear}

\item \simulation (One step for $k=4$.) Repeat the simulation of \cref{fig:parity:onestep} for $k=4$ and $d=12$, $16$, $24$, with $n=a\,d^3\ln d$ and $100$ data sets per point. You should observe a success probability near $\frac12$ for $a$ between $0.5$ and $0.8$, larger for larger $d$,, and above $0.9$ at $a=1.6$, close to the prediction $a\approx2k^2(2k-3)!!/(k!)^2=0.83$ of \cref{sec:parity:reading}.\label{ex:parity:k4}

\item \simulation (Step size and random starts.) (a) Rerun the simulation of \cref{fig:parity:neuron} with $\eta=c\,d^{-3/2}$ for $c=0.02$ and $c=0.1$. You should observe a median escape time within about $15\%$ of the population prediction for $c=0.02$, and a wider spread for $c=0.1$, from about $0.7$ to $1.7$ times the prediction, with the median delayed by about $20\%$; relate this to the noise at the start, as in \cref{ch:singleindex}. (b) Integrate the population flow \eqref{eq:parity:flow} for $d=50$, $k=3$ from several uniformly random starts with $\Psi(w_0)>0$. You should find that $\rho(t)$ always stays above the balanced solution with the same $\rho_0$, as \cref{thm:parity:hidden}(b) asserts, and that the time to $\rho=\frac12$ is shorter than the balanced one, typically by a factor between $0.5$ and $0.95$, and by more for starts with very small $\rho_0$.\label{ex:parity:stepsize}

\item \research (Lower bounds for plain SGD.) Prove that online SGD on a two-layer ReLU network of polynomial width with Gaussian initialization needs at least $d^{k-1-o(1)}$ samples to learn a $k$-sparse parity, for every step size. \Cref{thm:parity:sq} gives this only for adversarial gradient errors, and \textcite{AbbeSandon2020} show that random errors can be exploited by specially constructed networks; whether standard architectures can exploit them is open as far as we know.\label{ex:parity:sgdlower}

\item \research (A progress measure from a random start.) In the setting of \cref{fig:parity:mlp}, find a quantity computed from the weights at time $t$ that predicts the step of the drop with an error that shrinks as $d$ grows, from early in the plateau. The largest share of a neuron on $S$ moves only late; candidates include products of the relevant weights of each neuron, the analogue of $\Psi$, and the population Fourier gap of each neuron. No such predictor with guarantees is known to us.\label{ex:parity:measure}
\end{exercises}

\notesandrefs
Parities entered the theory of neural networks with \textcite{MinskyPapert1969}, whose limitations of bounded-order perceptrons included the parity of all inputs. The statistical-query model was introduced by Kearns in 1993; the journal version \parencite{Kearns1998} shows that statistical-query algorithms tolerate random classification noise and that parities cannot be learned from statistical queries in polynomial time. \Textcite{BlumEtAl1994} characterized weak statistical-query learning through Fourier analysis by the number of nearly orthogonal functions in a class, and \textcite{BlumKalaiWasserman2003} gave a noise-tolerant algorithm for parities on $O(\log d\log\log d)$ bits that no statistical-query algorithm matches, separating the two models. The Fourier analysis used here is in \textcite{ODonnell2014}, including the Fourier coefficients of majority.

\Textcite{ShalevShwartzShamirShammah2017} put parities at the head of their list of failures of gradient-based learning, and \textcite{AbbeSandon2020} delineated what SGD can and cannot do beyond statistical queries. \Textcite{DanielyMalach2020} proved that gradient descent on two-layer networks learns sparse parities under distributions where linear methods fail. The phenomenon of this chapter, abrupt learning with hidden progress near the statistical-query limit, is the subject of \textcite{BarakEtAl2022}; its follow-up on data, compute, width, and luck is \textcite{EdelmanEtAl2023}, and the matching upper bound for sign SGD is \textcite{KouEtAl2024}. The monomial neuron, the one-step theorem, and the comparison argument for the spherical flow are presented here in the form that makes every constant explicit; the doubling-epoch analysis of online SGD that would turn the heuristic of \cref{sec:parity:hidden} into a theorem is the one carried out for single-index models in \cref{ch:singleindex}, following \textcite{BenArousGheissariJagannath2021}.\sidenote{The equality case of \cref{thm:parity:hidden}(b) is exactly the population equation of \cref{ch:singleindex} with $\Psi(r)=c_kr^k$, so the proof of the online SGD theorem there transfers to balanced starts once the noise bounds are redone on the cube.}

For structure that rescues gradient methods see the staircase papers \parencite{AbbeEtAl2021,AbbeBoixAdseraMisiakiewicz2022}, the leap complexity \parencite{AbbeBoixAdseraMisiakiewicz2023}, and curricula with biased inputs \parencite{CornacchiaMossel2023}. For the debate on emergence see \textcite{WeiEtAl2022} and \textcite{SchaefferMirandaKoyejo2023}; for transformers, \textcite{HahnRofin2024} and \textcite{KimSuzuki2025}.

\paragraph{Translating the literature.}
\Textcite{BarakEtAl2022} write $n$ for our dimension $d$, keep $k$ for the sparsity and $S$ for the support, and call the task the $(n,k)$-sparse parity; their labels are in $\{\pm1\}$, as ours. \Textcite{Kearns1998} pairs a query with its tolerance $\tau$, as we do, but his original formulation takes labels and query values in $\{0,1\}$, and later work mostly uses queries with values in $[-1,1]$; the decomposition $q=q_0+yq_1$ in the proof of \cref{thm:parity:sq} is the $\{\pm1\}$ form of the split into a label-independent and a correlational part. \Textcite{ODonnell2014} writes $n$ for our $d$, $\chi_S$ and $\hat f(S)$ as we do, and $\mathrm{Maj}_n$ for our $\parityMaj_d$. \Textcite{KouEtAl2024} write $d$ for the dimension and $k$ for the sparsity, as we do. Our effective overlap $\rho$ is the overlap $r$ of \cref{ch:singleindex} for balanced weights, and our $c_k$ is the coefficient of $r^k$ in that chapter's $\Psi$.

\printchapterbib
\end{refsection}


\casepart{What Networks Represent}{Trained networks build structures that can be read off their weights: Fourier features for modular arithmetic, overcomplete codes for sparse features, a regular simplex for classes. Each structure is the optimum of a small problem that can be solved by hand.}{60}
% Chapter 10 of The Hydrogen Atoms of Deep Learning: grokking and the Fourier circuit.
\newcommand{\grokZp}{\Z_p}
\newcommand{\grokmem}{t_{\mathrm{mem}}}
\newcommand{\grokt}{t_{\mathrm{grok}}}
\newcommand{\grokcw}{\bar\gamma}
\newcommand{\grokii}{\mathrm{i}}

\casechapter{Grokking}{Modular addition and its Fourier circuit}{Modular addition, quadratic two-layer network}{ch:grokking}
\begin{refsection}[chapters/10-grokking.bib]

\begin{chaptersummary}
Small networks trained on part of the addition table of $\Z_p$ fit their training data early and predict the rest of the table only much later: this is grokking. For a two-layer network with quadratic activation the generalizing solution can be written down. Each neuron carries one Fourier frequency, and trigonometric identities turn the sum of the neurons into clocks $\cos(\omega(a+b-c))$ that peak at the correct answer. We prove that margin maximization forces this circuit: every neuron has a single frequency, its readout phase is the sum of its input phases, every frequency is present, and each frequency needs at least six neurons. On a diagonal linear network, weight decay acts exactly as an initialization scale shrinking like $e^{-\lambda t}$, which predicts a delay $\lambda^{-1}\log(\alpha/\alpha_c)$ that simulations confirm to within one percent. The model says nothing about transformers that use a handful of frequencies or a different algorithm, about grokking without weight decay, or about networks trained far from the margin limit.
\end{chaptersummary}

\section{The phenomenon: memorize now, generalize later}\label{sec:grokking:phenomenon}

In 2021 a group at OpenAI trained small transformers on tables of binary operations: addition, subtraction and division modulo $97$, polynomials such as $x^2+xy+y^2$, and composition of permutations in $S_5$. Each network saw a random part of a table and had to fill in the rest \parencite{PowerEtAl2022}.\sidenote{The work was presented at an ICLR 2021 workshop on mathematical reasoning and posted to arXiv in January 2022. The verb \emph{grok}, to understand something completely, comes from Robert Heinlein's novel \emph{Stranger in a Strange Land} (1961).} The networks were two-layer decoder-only transformers of width $128$ with four attention heads and about $4\times10^5$ parameters, trained with AdamW at learning rate $10^{-3}$ and weight decay $1$. On division modulo $97$ with half of the table for training, training accuracy was nearly perfect after about $10^3$ steps, while validation accuracy stayed near chance and approached $100\%$ only after close to $10^6$ steps. Power and coauthors called this delayed generalization \term{grokking}\index{grokking}. They also found that the delay grows steeply as the training set shrinks: at fractions of $25$ to $30\%$ of the table, removing $1\%$ of the data raised the median time to generalization by $40$ to $50\%$, and some operations never generalized. Among the interventions they tried, weight decay improved data efficiency the most.

\begin{observation}[\textcite{PowerEtAl2022}; \textcite{NandaEtAl2023}]\label[observation]{obs:grokking:delay}
A small network trained with AdamW and weight decay on a random fraction of the addition table of $\Z_p$ (between about $25\%$ and $50\%$ of it, for $p$ near $100$) reaches $100\%$ training accuracy early. It stays near chance on the rest of the table for one to three orders of magnitude longer, and then its test accuracy rises to about $100\%$ within a comparatively short interval. The delay grows steeply as the training fraction decreases, and weight decay shortens it.
\end{observation}

What does the network compute once it generalizes? \Textcite{NandaEtAl2023} reverse-engineered a one-layer transformer trained on $30\%$ of the addition table modulo $113$, with full-batch AdamW and weight decay $1$. Its embedding of a number $a$ is concentrated on five ``key frequencies'' $\omega_k=2\pi k/113$, $k\in\{14,35,41,42,52\}$: the network represents $a$ by $\cos(\omega_ka)$ and $\sin(\omega_ka)$. Attention and the MLP combine these into $\cos(\omega_k(a+b))$ and $\sin(\omega_k(a+b))$ by the addition formulas, and the unembedding produces logits close to $\sum_k\cos(\omega_k(a+b-c))$, which are largest at $c=a+b$. Nanda and coauthors also defined progress measures, losses computed after keeping or removing the key-frequency components of the logits, and found three phases: memorization (epochs $0$ to $1.4$ thousand), circuit formation (to $9.4$ thousand), during which the Fourier circuit grows while test loss stays flat, and cleanup (to $14$ thousand), during which weight decay removes the memorizing components and test loss drops. The jump in test accuracy is the visible end of a gradual process, the same kind of hidden progress as in \cref{ch:parity}.

How robust is this? Grokking on algorithmic tasks is reproducible across transformers and multilayer perceptrons and across group operations. \Textcite{LiuEtAl2022} mapped phase diagrams over learning rate and weight decay with four phases (comprehension, grokking, memorization, confusion) and read grokking as a symptom of hyperparameters that sit between fast comprehension and pure memorization. \Textcite{LiuEtAl2023} induced grokking on MNIST, with $1000$ training images and an initialization scaled up nine-fold, and found the delay roughly proportional to the inverse weight decay; on modular addition, a smaller initialization with the weight norm held fixed almost removed it. Delayed generalization has since been reported for convolutional networks on CIFAR-10 \parencite{HumayunEtAl2024}, Gaussian processes and linear regression \parencite{MillerEtAl2024}, and kernel machines with learned features \parencite{MallinarEtAl2025}. Two conclusions are safe. Grokking is a property of a training configuration (initialization scale, weight decay, data fraction) as much as of a task. And on group operations the generalizing networks compute with characters of the group.

These observations raise two questions, and the model of this chapter answers both in a minimal setting. \emph{Which} computation does a generalizing network perform, and why that one rather than another of the many functions that fit the training data? And \emph{why so late}: what sets the delay between fitting the training set and generalizing?

\section{The model: a quadratic network on the addition table}\label{sec:grokking:model}

In this chapter $p\ge3$ is an odd prime, the modulus of the arithmetic; this overrides, locally, the global use of $p$ for a number of features. The elements of $\grokZp=\{0,1,\dots,p-1\}$ are written $a$, $b$, $c$, and arithmetic on them is modulo $p$.\sidenote{The proofs use only that $p$ is odd. Primality matters elsewhere: for composite $p$, a frequency that shares a factor with $p$ cannot separate all the classes on its own.}

\paragraph{The task.} The data are the $p^2$ correct triples $(a,b,a+b)$. A training set $D$ is a random set of pairs $(a,b)$, a fraction of the table, labeled by $a+b$; the test error is the error on the remaining pairs. The input is one-hot: the pair $(a,b)$ is the vector $x=(e_a,e_b)\in\R^{2p}$, which says nothing about the arithmetic. Whatever structure the network finds, it must find in the labels.

\paragraph{The network.} A two-layer network with $m$ hidden neurons, quadratic activation, and $p$ outputs, the logits. The first-layer weight vector $w_r\in\R^{2p}$ of neuron $r$ splits into two functions $u_r,u'_r\colon\grokZp\to\R$, so that $\ip{w_r,x}=u_r(a)+u'_r(b)$, and the readout of neuron $r$ is a third function $v_r\colon\grokZp\to\R$. The logit of class $c$ is
\begin{equation}\label{eq:grokking:net}
  f_c(a,b;\theta)=\sum_{r=1}^m v_r(c)\bigl(u_r(a)+u'_r(b)\bigr)^2 .
\end{equation}
The parameters are $\theta=(\theta_1,\dots,\theta_m)$ with $\theta_r=(u_r,u'_r,v_r)\in(\R^p)^3$; we write $\norm{u}^2=\sum_au(a)^2$ and $\norm{\theta_r}^2=\norm{u_r}^2+\norm{u'_r}^2+\norm{v_r}^2$. The network predicts the class with the largest logit. Training minimizes the cross-entropy of the softmax of the logits over $D$ plus the weight decay $\frac\lambda2\norm\theta^2$, by gradient flow in the analysis and by AdamW in the experiments.

\paragraph{Two structural facts.} The network is \term{homogeneous of degree three} neuron by neuron\index{homogeneous network}: replacing $\theta_r$ by $t\theta_r$ with $t>0$ multiplies the contribution of neuron $r$ to every logit by $t^3$. Homogeneity makes margins the right currency, as in \cref{ch:margin}. And the task is invariant under translations: if $(a,b,c)$ is a correct triple, then so is $(a+g,b+h,c+g+h)$ for all $g,h\in\grokZp$. A problem invariant under a commutative group is diagonalized by the characters of the group, here the functions $a\mapsto e^{2\pi\grokii ka/p}$, so translation invariance is the reason to expect Fourier analysis.

\paragraph{What was thrown away.} The networks in which grokking was discovered are transformers: tokens are embedded, attention mixes the two positions, an MLP with ReLU activations follows, and an unembedding produces logits; training uses Adam with decoupled weight decay on a fraction of the table. The model keeps the two ingredients that the solution needs and nothing else. One-hot inputs with free first-layer weights are a learned embedding. The quadratic activation is the simplest nonlinearity that multiplies, since $(x+y)^2=x^2+2xy+y^2$, and the trigonometric identities behind the circuit need exactly one multiplication; it also makes every logit a polynomial in the weights, so that the Fourier analysis is exact rather than approximate. The margin analysis is done on the full table, which describes what a generalizing network must look like but not why training on part of the table produces one; \cref{sec:grokking:wrong} returns to this. For the delay we use a second, even simpler model, the diagonal linear network of \cref{ch:diagonal} with weight decay, in which the passage from the lazy to the rich regime can be computed exactly.

\begin{dictionary}
one-hot $(e_a,e_b)$ and first-layer functions $u_r,u'_r$ & token embeddings of $a$ and $b$, read by an MLP neuron\\
quadratic activation; the cross term $2u_r(a)u'_r(b)$ & the multiplication done by attention and by the MLP nonlinearity\\
frequency $\zeta$ and phases $\varphi,\psi$ of a Fourier neuron & a key frequency, and the cosine and sine directions of the embedding \parencite{NandaEtAl2023}\\
readout $v_r(c)\propto\cos(\omega_\zeta c+\varphi+\psi)$ & rows of the unembedding; logits $\propto\sum_\zeta\cos(\omega_\zeta(a+b-c))$\\
normalized margin, norm $\norm{\theta}_{2,3}$ & the long-time limit of cross-entropy training with small weight decay\\
average margin over all wrong classes & gap between the correct logit and the mean wrong logit\\
initialization scale $\alpha$; effective scale $\alpha e^{-\lambda t}$ & parameter norm, which decays under weight decay after memorization\\
critical scale $\alpha_c$ & the norm below which training leaves the lazy regime\\
\end{dictionary}

\paragraph{What one had to notice.} Three things. First, that a network trained on part of a group's table can only generalize by finding the group structure, and that for a cyclic group the structure is carried by characters, which a network can store as cosines in its embedding. Second, that one multiplication suffices: $\cos(\omega(a+b))$ is a sum of products of a function of $a$ and a function of $b$, and a square supplies products. Third, that among all networks that classify the table correctly, training selects one by a norm, so the question ``which circuit?'' becomes an extremal problem, and the extremal problem is diagonalized by the same characters.

\section{The solution: a Fourier circuit for addition}\label{sec:grokking:circuit}

\paragraph{Fourier analysis on $\grokZp$.} For $u\colon\grokZp\to\R$ let $\hat u(k)=\sum_{a\in\grokZp}u(a)e^{-2\pi\grokii ka/p}$ for $k\in\grokZp$. In this normalization the inversion formula, Plancherel's identity, and the symmetry of real functions read
\[
  u(a)=\frac1p\sum_{k\in\grokZp}\hat u(k)e^{2\pi\grokii ka/p},\qquad \sum_au(a)^2=\frac1p\sum_k\abs{\hat u(k)}^2,\qquad \hat u(-k)=\overline{\hat u(k)}
\]
(\cref{thm:toolkit:fourier-zp}). All three follow from the orthogonality relation
\begin{equation}\label{eq:grokking:orth}
  \sum_{a\in\grokZp}e^{2\pi\grokii qa/p}=p\,\ind\{q\equiv0\pmod p\}\qquad\text{for every integer }q .
\end{equation}
Because $p$ is odd, the nonzero frequencies pair up as $\pm\zeta$ with $\zeta\in\{1,\dots,\frac{p-1}2\}$; we write $\omega_\zeta=2\pi\zeta/p$. A real function whose transform vanishes outside $\pm\zeta$ is $a\mapsto\rho\cos(\omega_\zeta a+\varphi)$ with $\rho=\frac2p\abs{\hat u(\zeta)}$ and $\varphi=\arg\hat u(\zeta)$.

\begin{lemma}[Sum of cosines]\label[lemma]{lem:grokking:cosines}
For every integer $q$, $\displaystyle\sum_{\zeta=1}^{(p-1)/2}\cos\frac{2\pi\zeta q}{p}=\frac{p\,\ind\{q\equiv0\}-1}{2}$.
\end{lemma}

\begin{proof}
In \eqref{eq:grokking:orth} the term $a=0$ equals $1$, and the terms $a=\zeta$ and $a=p-\zeta$ are complex conjugates with real part $\cos(2\pi\zeta q/p)$. Hence \eqref{eq:grokking:orth} reads $1+2\sum_{\zeta=1}^{(p-1)/2}\cos(2\pi\zeta q/p)=p\,\ind\{q\equiv0\}$.
\end{proof}

With $q=a+b-c$ the lemma says that the logits
\begin{equation}\label{eq:grokking:target}
  f_c(a,b)=\sum_{\zeta=1}^{(p-1)/2}\cos\bigl(\omega_\zeta(a+b-c)\bigr)=\frac{p\,\ind\{c=a+b\}-1}{2}
\end{equation}
are $\frac{p-1}2$ on the correct class and $-\frac12$ on every other class. Each term is a \term{clock}\index{clock}: $\cos(\omega_\zeta(a+b-c))$ is largest when the angle $\omega_\zeta c$ equals $\omega_\zeta a+\omega_\zeta b$, so it compares the sum of two rotations with a third. One frequency already peaks only at the correct class, because $p$ is prime; summing all of them makes every wrong class tie. The question is how a sum of squares of first-layer activations produces one clock.

\paragraph{One frequency at a time.} Fix a frequency $\zeta$ and write $\omega=\omega_\zeta$. Call the neuron with $u(a)=\cos(\omega a+\varphi)$, $u'(b)=\cos(\omega b+\psi)$ and $v(c)=\cos(\omega c+\varphi+\psi)$ the \term{Fourier neuron}\index{Fourier neuron} of frequency $\zeta$ and phases $(\varphi,\psi)$, and write
\[
  n_{\varphi,\psi}(a,b,c)=\bigl(\cos(\omega a+\varphi)+\cos(\omega b+\psi)\bigr)^2\cos(\omega c+\varphi+\psi)
\]
for its contribution to the logits. The readout phase is the sum of the input phases; \cref{thm:grokking:neuron} will explain why. It is convenient to use the readout phase $\chi=\varphi+\psi$ and the difference $\delta=\varphi-\psi$ as coordinates; the function $n_{\varphi,\psi}$ depends on $(\varphi,\psi)$ only through $\chi$ and $\delta$ modulo $2\pi$.

\begin{lemma}[Cancellation within one frequency]\label[lemma]{lem:grokking:block}
Let $\beta_1,\dots,\beta_K$ be real, let $(\varphi_r,\psi_r)$ be phases, and let $\chi_r=\varphi_r+\psi_r$ and $\delta_r=\varphi_r-\psi_r$. Then
\begin{equation}\label{eq:grokking:block}
  \sum_{r=1}^K\beta_r\,n_{\varphi_r,\psi_r}(a,b,c)=\frac12\Bigl(\sum_{r=1}^K\beta_r\Bigr)\cos\bigl(\omega(a+b-c)\bigr)\qquad\text{for all }a,b,c\in\grokZp
\end{equation}
if and only if, for all integers $j$ and $s$ with $\abs j\le2$, $\abs s\le1$ and $(j,s)\ne(0,0)$,
\begin{equation}\label{eq:grokking:moments}
  M_{j,s}\defeq\sum_{r=1}^K\beta_re^{\grokii(j\chi_r+s\delta_r)}=0 .
\end{equation}
\end{lemma}

\begin{proof}
The plan is to expand one neuron into cosines of linear forms in $(a,b,c)$ and to read off the coefficient of each. Put $A=\omega a+\varphi$, $B=\omega b+\psi$ and $C=\omega c+\chi$. Using $\cos^2x=\frac12(1+\cos2x)$ and $4\cos A\cos B\cos C=\cos(A+B-C)+\cos(A+B+C)+\cos(A-B+C)+\cos(A-B-C)$,
\begin{align*}
  n_{\varphi,\psi}&=\tfrac12\cos(A+B-C)+\cos C+\tfrac14\bigl[\cos(2A\pm C)+\cos(2B\pm C)\bigr]\\
  &\qquad+\tfrac12\bigl[\cos(A+B+C)+\cos(A-B\pm C)\bigr],
\end{align*}
each $\pm$ standing for the sum of the two terms. Because $\chi=\varphi+\psi$, the first term is $\frac12\cos(\omega(a+b-c))$ for every neuron. Every other term is $\cos(\omega\ell+j\chi+s\delta)$ with a linear form $\ell$ in $(a,b,c)$ and integers $(j,s)$, as follows: $C$ has $(\ell;j,s)=(c;1,0)$; $2A+C$ has $(2a+c;2,1)$; $2A-C$ has $(2a-c;0,1)$; $2B+C$ has $(2b+c;2,-1)$; $2B-C$ has $(2b-c;0,-1)$; $A+B+C$ has $(a+b+c;2,0)$; $A-B+C$ has $(a-b+c;1,1)$; and $A-B-C$ has $(a-b-c;-1,1)$. Up to sign, these pairs $(j,s)$ are the seven nonzero pairs with $\abs j\le2$ and $\abs s\le1$. Weighting by $\beta_r$ and summing over $r$, the term with $(\ell;j,s)$ contributes its coefficient times $\operatorname{Re}\bigl(e^{\grokii\omega\ell}M_{j,s}\bigr)$, and $M_{-j,-s}=\overline{M_{j,s}}$ because the weights are real. If \eqref{eq:grokking:moments} holds, then every such contribution vanishes, and \eqref{eq:grokking:block} follows.

Conversely, the eighteen functions $e^{\pm\grokii\omega(a+b-c)}$ and $e^{\pm\grokii\omega\ell}$ are characters of $\grokZp^3$ with distinct frequency vectors: two of these vectors differ in some coordinate by $1$, $2$ or $4$, which multiplied by $\zeta$ is not divisible by $p$. Characters with distinct frequency vectors are orthogonal by \eqref{eq:grokking:orth}. So if \eqref{eq:grokking:block} holds, then the coefficient of each $e^{\grokii\omega\ell}$ on its left side vanishes, and that coefficient is a nonzero multiple of the corresponding $M_{j,s}$.
\end{proof}

In the language of numerical integration, \eqref{eq:grokking:moments} says that the points $(\chi_r,\delta_r)$ with weights $\beta_r$ integrate every trigonometric polynomial of degree at most $2$ in $\chi$ and at most $1$ in $\delta$ exactly, as the uniform measure of total mass $\sum_r\beta_r$ does. Such a set of weighted points is a \term{cubature rule}\index{cubature rule} on the torus, and every cubature rule of this kind gives a circuit.

\begin{theorem}[A Fourier network for modular addition]\label{thm:grokking:fourier}
For each frequency $\zeta\in\{1,\dots,\frac{p-1}2\}$ and each pair of phases $(\varphi,\psi)\in\{0,\frac{2\pi}3,\frac{4\pi}3\}^2$, take the neuron with $u(a)=\cos(\omega_\zeta a+\varphi)$, $u'(b)=\cos(\omega_\zeta b+\psi)$ and $v(c)=\frac29\cos(\omega_\zeta c+\varphi+\psi)$. The resulting network of width $\frac92(p-1)$ has logits
\[
  f_c(a,b)=\frac{p\,\ind\{c=a+b\}-1}2\qquad\text{for all }a,b,c\in\grokZp .
\]
In particular it classifies the whole table correctly, and every margin $f_{a+b}(a,b)-f_c(a,b)$ with $c\ne a+b$ equals $\frac p2$.
\end{theorem}

\begin{proof}
Fix $\zeta$. Its nine neurons are the Fourier neurons $n_{\varphi,\psi}$ with weights $\beta=\frac29$. To compute the moments, write $j\chi+s\delta=(j+s)\varphi+(j-s)\psi$; the sum over the grid factorizes,
\[
  M_{j,s}=\frac29\Bigl(\sum_{\varphi}e^{\grokii(j+s)\varphi}\Bigr)\Bigl(\sum_{\psi}e^{\grokii(j-s)\psi}\Bigr),
\]
and by \eqref{eq:grokking:orth} with modulus $3$ the first factor vanishes unless $3$ divides $j+s$, the second unless $3$ divides $j-s$. If $3$ divides both, then it divides $2j$ and $2s$, hence $j$ and $s$, and $\abs j\le2$, $\abs s\le1$ force $j=s=0$. So \eqref{eq:grokking:moments} holds, and \cref{lem:grokking:block} shows that the nine neurons of frequency $\zeta$ contribute $\frac12\cdot9\cdot\frac29\cos(\omega_\zeta(a+b-c))=\cos(\omega_\zeta(a+b-c))$. Summing over $\zeta$ gives \eqref{eq:grokking:target}.
\end{proof}

The proof is a bookkeeping of phases. The squares $u(a)^2$ and $u'(b)^2$ create terms that ignore $b$ or $a$; the cross term creates the wanted clock together with the unwanted $\cos(\omega(a-b\pm c)+\cdots)$ and $\cos(\omega(a+b+c)+\cdots)$; and three phases per input make every unwanted term a sum of three unit vectors spaced by $2\pi/3$. We checked the identity numerically for all odd $p\le23$, where it holds to rounding error.

\begin{remark}[Other designs]\label[remark]{rem:grokking:designs}
A pair of Fourier neurons with phases $(\varphi,\psi)$ and $(\varphi,\psi+\pi)$ turns squares into a product, $n_{\varphi,\psi}+n_{\varphi,\psi+\pi}=\bigl[(x+y)^2-(x-y)^2\bigr]z=4xyz$, with $x=\cos(\omega a+\varphi)$, $y=\cos(\omega b+\psi)$ and $z=\cos(\omega c+\varphi+\psi)$. \Textcite{MorwaniEtAl2024} use four such pairs, one for each product in $\cos(A+B-C)=\cos A\cos B\cos C-\sin A\sin B\cos C+\sin A\cos B\sin C+\cos A\sin B\sin C$: eight neurons per frequency. Three pairs suffice: for $\theta\in\{0,\frac\pi3,\frac{2\pi}3\}$ take $(\varphi,\psi)=(\theta,\theta)$ and $(\theta,\theta+\pi)$. Expanding $4xyz$ as in the proof of \cref{lem:grokking:block}, the three unwanted terms carry the phases $4\theta$, $2\theta$ and $-2\theta$, whose exponentials sum to zero over $\theta$, so
\begin{align*}
  \frac13\sum_\theta\bigl(n_{\theta,\theta}+n_{\theta,\theta+\pi}\bigr)&=\frac43\sum_\theta\cos(\omega a+\theta)\cos(\omega b+\theta)\cos(\omega c+2\theta)\\
  &=\cos\bigl(\omega(a+b-c)\bigr).
\end{align*}
These six neurons per frequency are the ``order-6'' block of \textcite{Tian2025}, who assembles global minimizers of a squared loss from such blocks; they give a network of width $3(p-1)$. \Textcite{Gromov2023} gives a construction with independent uniformly random phases, readout weight $2/N$ and $N$ neurons per frequency; there each $M_{j,s}$ is a sum of $N$ independent terms of mean zero, of size about $2/\sqrt N$, so the unwanted terms cancel only as the width grows.
\end{remark}

The six-neuron identity is a three-term decomposition of the trilinear form $\cos(A+B-C)=\operatorname{Re}(e^{\grokii A}e^{\grokii B}e^{-\grokii C})$ into products of cosines, and three is the least possible: it is the number of real multiplications needed to multiply two complex numbers, as in Gauss's trick.\sidenote{Multiplying $x_1+\grokii x_2$ by $y_1+\grokii y_2$ needs the four products $x_jy_k$ in the obvious way, but three suffice: $x_1y_1$, $x_2y_2$, and $(x_1+x_2)(y_1+y_2)$. Over the reals, three is the minimum.} Can a frequency do with fewer than six Fourier neurons? Not if the weights are positive, which is the case that margins will select.

\begin{lemma}[Six neurons per frequency]\label[lemma]{lem:grokking:six}
If $\beta_1,\dots,\beta_K>0$ and the moment conditions \eqref{eq:grokking:moments} hold, then the points $(\chi_r,\delta_r)$, taken modulo $2\pi$, include at least six distinct points. In particular $K\ge6$.
\end{lemma}

\begin{proof}
Neurons with the same point $(\chi_r,\delta_r)$ modulo $2\pi$ contribute the same summand to every moment, so merging them and adding their weights changes neither the hypothesis nor the positivity; assume that the $K$ points are distinct. Consider the six-dimensional space of trigonometric polynomials
\[
  P(\chi,\delta)=\sum_{\abs j\le1}\,\sum_{s\in\{0,1\}}c_{js}e^{\grokii(j\chi+s\delta)} .
\]
If $K<6$, then the $K$ linear conditions $P(\chi_r,\delta_r)=0$ on six unknowns have a nonzero solution $(c_{js})$. Expanding $\abs{P}^2$, every product of two of its exponentials has frequency $(j-j',s-s')$ with $\abs{j-j'}\le2$ and $\abs{s-s'}\le1$, so the moment conditions kill every term with $(j,s)\ne(j',s')$:
\[
  0=\sum_{r=1}^K\beta_r\abs{P(\chi_r,\delta_r)}^2=\sum_{j,s,j',s'}c_{js}\overline{c_{j's'}}\,M_{j-j',s-s'}=M_{0,0}\sum_{j,s}\abs{c_{js}}^2 .
\]
But $M_{0,0}=\sum_r\beta_r>0$ and the coefficients are not all zero, a contradiction.
\end{proof}

The argument is the classical lower bound for cubature rules with positive weights: a rule that is exact on all products $P\bar Q$ of functions from a space of dimension $N$ needs at least $N$ points.\sidenote{The bound is attained by the six-point rule of \cref{rem:grokking:designs}. A rule that attains it is called tight, and a tight rule has equal weights (\cref{ex:grokking:phases}).} Without the positivity the bound fails: a block of four neurons with unequal input amplitudes and readout weights of both signs produces the clock up to a term that depends on $c$ alone, and \textcite{Tian2025} cancels such terms across frequencies to build exact circuits of width $2p$. Those neurons are not Fourier neurons, and the next section shows that margins exclude them.

\section{The solution: margins select the circuit}\label{sec:grokking:margin}

\Cref{thm:grokking:fourier} exhibits a circuit; it does not say why training finds it. \Textcite{MorwaniEtAl2024} proposed that margin maximization alone determines it. Recall from \cref{ch:margin} why margins are the natural target. For a homogeneous network trained by cross-entropy on data that it can classify, the minimizers of the regularized loss converge in direction, as the weight decay tends to zero, to a network that maximizes the normalized margin \parencite{WeiEtAl2019}, and gradient flow on the unregularized loss converges in direction to a Karush--Kuhn--Tucker point of the margin problem \parencite{LyuLi2020}.\sidenote{The first statement is about global minimizers, for the norm that is penalized. The second gives a stationary point of the margin problem, not necessarily a global maximizer.} Weight decay penalizes $\norm\theta^2=\sum_r\norm{\theta_r}^2$. Morwani and coauthors normalize instead by
\[
  \norm{\theta}_{2,3}=\Bigl(\sum_{r=1}^m\norm{\theta_r}^3\Bigr)^{1/3},
\]
the norm matched to neurons that are homogeneous of degree three: with it the margin problem decouples into one problem per neuron and becomes solvable. With the Euclidean norm it does not decouple (\cref{ex:grokking:l2}), and we know of no complete solution in that case. The analysis is on the full table, all $p^2$ pairs.

\begin{definition}[Margins]\label[definition]{def:grokking:margin}
The \term{margin}\index{margin!of a network} of $\theta$ at $(a,b)$ against $c\ne a+b$ is $f_{a+b}(a,b;\theta)-f_c(a,b;\theta)$. The margin $\gamma(\theta)$ of the network is the smallest of these $p^2(p-1)$ numbers, and its \term{average margin}\index{margin!average} $\grokcw(\theta)$ is their average. The \term{normalized margin} of $\theta\ne0$ is $\gamma(\theta)/\norm\theta_{2,3}^3$, which is invariant under $\theta\mapsto t\theta$ for $t>0$. A network of width $m$ is a \term{max-margin network}\index{max-margin network} if it maximizes the normalized margin among networks of width $m$.
\end{definition}

The maximum exists, because $\gamma$ is continuous on the compact set $\{\norm{\theta}_{2,3}=1\}$. The average margin is the class-weighted margin of Morwani and coauthors: the margin averaged over a uniform pair $(a,b)$ and a uniform wrong class. Its virtue is linearity: $\grokcw(\theta)=\sum_r\grokcw(\theta_r)$, the sum of the average margins of the one-neuron networks $\theta_r$.

\begin{lemma}[Average margin of one neuron]\label[lemma]{lem:grokking:cw}
For one neuron $(u,u',v)$, with $\bar v=\frac1p\sum_cv(c)$ and $(a,b)$ uniform on $\grokZp^2$,
\begin{align}
  \grokcw(u,u',v)&=\frac{2p}{p-1}\,\E_{a,b}\bigl[u(a)u'(b)\bigl(v(a+b)-\bar v\bigr)\bigr]\nonumber\\
  &=\frac{2}{(p-1)p^2}\sum_{k\ne0}\hat u(k)\,\hat u'(k)\,\overline{\hat v(k)} .\label{eq:grokking:cw}
\end{align}
\end{lemma}

\begin{proof}
For fixed $(a,b)$ the average of $v(c)$ over the $p-1$ wrong classes is $(p\bar v-v(a+b))/(p-1)$, so the average margin at $(a,b)$ is $\frac p{p-1}\bigl(u(a)+u'(b)\bigr)^2\bigl(v(a+b)-\bar v\bigr)$. Expand the square. The term $u(a)^2$ averages to zero, because for fixed $a$ the class $a+b$ is uniform when $b$ is, so that $\E_b[v(a+b)]=\bar v$; the term $u'(b)^2$ vanishes in the same way. The cross term $2u(a)u'(b)$ remains, which proves the first equality. For the second, insert $v(c)=\frac1p\sum_k\hat v(k)e^{2\pi\grokii kc/p}$ and sum over $a$ and $b$ separately:
\[
  \sum_{a,b}u(a)u'(b)v(a+b)=\frac1p\sum_k\hat v(k)\,\hat u(-k)\,\hat u'(-k).
\]
The term $k=0$ equals $\hat u(0)\hat u'(0)\hat v(0)/p=\sum_{a,b}u(a)u'(b)\bar v$, which the subtraction of $\bar v$ removes. Divide by $p^2$, multiply by $\frac{2p}{p-1}$, replace $k$ by $-k$, and use $\hat v(-k)=\overline{\hat v(k)}$.
\end{proof}

The cancellation of the squares is the reason for averaging uniformly over wrong classes: the terms $u(a)^2v(c)$ and $u'(b)^2v(c)$, which no single neuron can avoid, carry no information about $a+b$ and contribute equally to the correct logit and to the average wrong logit. What remains is a cubic form that the Fourier transform diagonalizes.

\begin{theorem}[The best single neuron]\label{thm:grokking:neuron}
Every neuron satisfies $\grokcw(\theta_r)\le\gamma^\star\norm{\theta_r}^3$, with
\[
  \gamma^\star=\sqrt{\frac2{27}}\,\frac1{\sqrt p\,(p-1)} .
\]
A neuron $\theta_r\ne0$ attains equality if and only if it is a positive multiple of a Fourier neuron: $u(a)=\rho\cos(\omega_\zeta a+\varphi)$, $u'(b)=\rho\cos(\omega_\zeta b+\psi)$ and $v(c)=\rho\cos(\omega_\zeta c+\varphi+\psi)$ for some $\rho>0$, frequency $\zeta\in\{1,\dots,\frac{p-1}2\}$, and phases $\varphi,\psi$.
\end{theorem}

\begin{proof}
The plan is to pair conjugate frequencies, to bound each pair by the arithmetic--geometric mean inequality, and then to use that $s\mapsto s^{3/2}$ grows faster than linearly to put all the weight on one pair. Both sides are homogeneous of degree three, so assume $\norm{\theta_r}=1$.

\emph{Pairing.} Because $u$, $u'$ and $v$ are real, the terms $k$ and $-k$ of \eqref{eq:grokking:cw} are complex conjugates, and $k\ne-k$ because $p$ is odd. With $x_\zeta=\abs{\hat u(\zeta)}$, $y_\zeta=\abs{\hat u'(\zeta)}$ and $z_\zeta=\abs{\hat v(\zeta)}$,
\[
  \grokcw(\theta_r)=\frac4{(p-1)p^2}\sum_{\zeta=1}^{(p-1)/2}\operatorname{Re}\Bigl(\hat u(\zeta)\hat u'(\zeta)\overline{\hat v(\zeta)}\Bigr)\le\frac4{(p-1)p^2}\sum_{\zeta=1}^{(p-1)/2}x_\zeta y_\zeta z_\zeta .
\]
\emph{The constraint.} Plancherel's identity, with the terms $\pm\zeta$ grouped and the term $k=0$ dropped, gives $1=\norm{u}^2+\norm{u'}^2+\norm{v}^2\ge\frac2p\sum_\zeta s_\zeta$ with $s_\zeta=x_\zeta^2+y_\zeta^2+z_\zeta^2$; equality requires $\hat u(0)=\hat u'(0)=\hat v(0)=0$.

\emph{One frequency.} By the arithmetic--geometric mean inequality $x_\zeta y_\zeta z_\zeta\le(s_\zeta/3)^{3/2}$, with equality if and only if $x_\zeta=y_\zeta=z_\zeta$. And $\sum_\zeta s_\zeta^{3/2}\le(\max_\zeta s_\zeta)^{1/2}\sum_\zeta s_\zeta\le\bigl(\sum_\zeta s_\zeta\bigr)^{3/2}$, with equality if and only if at most one $s_\zeta$ is nonzero. With $\sum_\zeta s_\zeta\le p/2$ these give
\[
  \grokcw(\theta_r)\le\frac4{(p-1)p^2}\Bigl(\frac{p}{6}\Bigr)^{3/2}=\sqrt{\frac2{27}}\,\frac1{\sqrt p\,(p-1)} .
\]
\emph{Equality.} Equality holds throughout if and only if the transforms vanish at $0$ and at every frequency except one pair $\pm\zeta$, the three moduli at $\zeta$ are equal (to $\sqrt{p/6}$), and $\hat u(\zeta)\hat u'(\zeta)\overline{\hat v(\zeta)}$ is a positive real, that is, $\arg\hat v(\zeta)=\arg\hat u(\zeta)+\arg\hat u'(\zeta)$. By the remark before \cref{lem:grokking:cosines}, these are exactly the functions $\rho\cos(\omega_\zeta a+\varphi)$, $\rho\cos(\omega_\zeta b+\psi)$ and $\rho\cos(\omega_\zeta c+\varphi+\psi)$ with $\rho=\frac2p\sqrt{p/6}=\sqrt{2/(3p)}$. Conversely, such a neuron makes every inequality of the chain an equality.
\end{proof}

Two features of the optimum deserve their names. The best neuron is \emph{single-frequency}: mixing frequencies wastes norm, because the objective is cubic and the constraint quadratic, and a cubic gains more from concentration than a quadratic loses. And its phases \emph{add}: the readout phase is the sum of the two input phases, so that the cross term $\cos(\omega a+\varphi)\cos(\omega b+\psi)$ contains $\frac12\cos(\omega(a+b)+\chi)$, which the readout $\cos(\omega c+\chi)$ turns into the clock $\frac14\cos(\omega(a+b-c))$. A single neuron cannot be the whole network, though: its logits $(u(a)+u'(b))^2v(c)$ peak at the correct class only for some pairs. The next theorem assembles optimal neurons into a network all of whose margins are equal.

\begin{theorem}[Max-margin networks; \textcite{MorwaniEtAl2024}, with sharp width]\label{thm:grokking:maxmargin}
\leavevmode
\begin{enumerate}
\item Every network of width $m$ satisfies $\gamma(\theta)\le\gamma^\star\norm{\theta}_{2,3}^3$.
\item A network $\theta\ne0$ attains equality if and only if its logits are
\[
  f_c(a,b;\theta)=\gamma^\star\norm\theta_{2,3}^3\Bigl(\ind\{c=a+b\}-\frac1p\Bigr)\qquad\text{for all }a,b,c .
\]
Then every nonzero neuron is a positive multiple of a Fourier neuron, and each frequency $\zeta\in\{1,\dots,\frac{p-1}2\}$ is carried by at least six neurons.
\item Equality is attained at width $m$ if and only if $m\ge3(p-1)$.
\end{enumerate}
\end{theorem}

\begin{proof}
The plan is to compare the minimum margin with the average margin, which decomposes over neurons, and then to read the equality cases in Fourier space; part (iii) rests on the cubature lemmas of \cref{sec:grokking:circuit}.

\emph{The bound.} A minimum is at most an average, so by \cref{thm:grokking:neuron},
\begin{equation}\label{eq:grokking:chain}
  \gamma(\theta)\le\grokcw(\theta)=\sum_{r=1}^m\grokcw(\theta_r)\le\gamma^\star\sum_{r=1}^m\norm{\theta_r}^3=\gamma^\star\norm\theta_{2,3}^3 .
\end{equation}
\emph{Equality.} If the logits have the stated form, then every margin equals $\gamma^\star\norm\theta^3_{2,3}$, and equality holds. Conversely, suppose that equality holds, and write $\gamma'=\gamma(\theta)$. The first inequality in \eqref{eq:grokking:chain} is then an equality, so all margins equal $\gamma'$: for every $(a,b)$ the wrong logits share a common value $g(a,b)$, and $f_c(a,b)=g(a,b)+\gamma'\ind\{c=a+b\}$. The second inequality is an equality term by term, so by \cref{thm:grokking:neuron} every nonzero neuron is a positive multiple of a Fourier neuron. A Fourier readout has mean zero over $c$, so $\sum_cf_c(a,b)=0$, which gives $g=-\gamma'/p$.

\emph{Every frequency, six times.} Group the neurons by frequency. The neurons of frequency $\zeta$ are $\rho_r$ times Fourier neurons with phases $(\varphi_r,\psi_r)$, and they contribute $B_\zeta=\sum_r\rho_r^3\,n_{\varphi_r,\psi_r}$, whose Fourier transform in $c$, for fixed $(a,b)$, lives on $\pm\zeta$. By \cref{lem:grokking:cosines} the logits are $\gamma'(\ind\{c=a+b\}-\frac1p)=\frac{2\gamma'}p\sum_\zeta\cos(\omega_\zeta(a+b-c))$, and the transform in $c$ of this sum at $\pm\zeta$ is that of its $\zeta$ term alone. Comparing transforms in $c$, frequency by frequency,
\[
  B_\zeta(a,b,c)=\frac{2\gamma'}p\cos\bigl(\omega_\zeta(a+b-c)\bigr)\qquad\text{for every }\zeta .
\]
The right side is not zero, so some neuron carries frequency $\zeta$; by \cref{lem:grokking:block}, applied with the positive weights $\beta_r=\rho_r^3$, the moment conditions hold; and by \cref{lem:grokking:six} at least six neurons carry $\zeta$.

\emph{Width.} Six neurons for each of the $\frac{p-1}2$ frequencies make $m\ge3(p-1)$ necessary. For $m\ge3(p-1)$, use the six-neuron blocks of \cref{rem:grokking:designs} at every frequency and set the other neurons to zero. Multiplying $u$ and $u'$ by $t$ and $v$ by $t^{-2}$ does not change a neuron's function, so the neuron $(\cos(\omega a+\varphi),\cos(\omega b+\psi),\frac13\cos(\omega c+\chi))$ may be replaced by the Fourier neuron $3^{-1/3}(\cos(\omega a+\varphi),\cos(\omega b+\psi),\cos(\omega c+\chi))$. The logits are \eqref{eq:grokking:target}, so all margins are equal and the first inequality in \eqref{eq:grokking:chain} is an equality; every neuron is a Fourier neuron, so the second is one as well.
\end{proof}

So the max-margin networks form a large family, one positive cubature rule per frequency with arbitrary rotations and scalings, but they all compute one function, the clock sum \eqref{eq:grokking:target} up to scale, and they share three visible features: single-frequency neurons, a readout phase equal to the sum of the input phases, and all $\frac{p-1}2$ frequencies present. \Textcite{MorwaniEtAl2024} prove parts (i) and (ii) for width $m\ge4(p-1)$, using their eight-neuron blocks; the six-neuron blocks and \cref{lem:grokking:six} make the width threshold exact. The same paper treats sparse parities and finite groups, where irreducible representations replace the characters.

\section{The solution: why the circuit comes late}\label{sec:grokking:delay}

The max-margin theorem describes the end of training. Why is the end reached so long after the training set is fit? Two groups proposed the same answer in 2023 \parencite{LyuEtAl2024,KumarEtAl2024}: training starts in the lazy regime of \cref{ch:lazy}, where the network fits the training set with its initial tangent kernel and memorizes, and reaches the rich regime, where margin-maximizing features form, only later.

\begin{theorem}[\textcite{LyuEtAl2024}; stated without proof]\label{thm:grokking:lyu}
Let $f(\theta;x)$ be $L$-homogeneous and twice continuously differentiable in $\theta$, let $f(\bar\theta;x)=0$ for every input $x$, and run gradient flow on $\Lhat(\theta)+\frac\lambda2\norm\theta^2$ from $\theta(0)=\alpha\bar\theta$, with weight decay $\lambda=\lambda(\alpha)$ of order $\alpha^{-q}$ for a constant $q>0$. Fix a constant $c\in(0,1)$ and let $\alpha\to\infty$.
\begin{enumerate}
\item At time $(1-c)\lambda^{-1}\log\alpha$ the predictor converges pointwise, after division by a normalizing factor in the case of classification, to a kernel predictor for the tangent features $\nabla_\theta f(\bar\theta;x)$: the maximum-$\ell_2$-margin classifier, if these features separate the training data, or the minimum-norm interpolant, for regression with linearly independent features.
\item At suitable times between $\lambda^{-1}\log\alpha$ and $(1+c)\lambda^{-1}\log\alpha$, every limit point of $\theta/\norm\theta$ (classification) or of $\theta$ (regression) is a Karush--Kuhn--Tucker point of $\min\frac12\norm\theta^2$ subject to $y_if(\theta;x_i)\ge1$ for all $i$, respectively $f(\theta;x_i)=y_i$ for all $i$.
\end{enumerate}
\end{theorem}

For diagonal linear networks the first problem is the maximum-$\ell_2$-margin classifier and the second the maximum-$\ell_1$-margin classifier, which needs far fewer samples for a sparse target; that gap is grokking. \Textcite{KumarEtAl2024} reach the same picture from the other side. They observe grokking in a two-layer network that learns $y=\frac12\ip{w_\star,x}^2$ from Gaussian inputs by plain gradient descent with no weight decay, control it by an output scale that sets the laziness of the dynamics, and find that it needs a large scale and a training set of intermediate size; they derive no formula for the transition time. For modular addition itself, \textcite{MohamadiEtAl2024} prove that permutation-equivariant kernel methods, a model of the lazy phase, need a constant fraction of the table to generalize: the lazy phase can only memorize.

The diagonal network of \cref{ch:diagonal} with weight decay makes the two phases explicit. Let $X\in\R^{n\times d}$ with $n<d$, let $y=Xw_\star$ with a $k$-sparse target $w_\star$, let $w=u\odot u-v\odot v$ with $u(0)=v(0)=\alpha\mathbf1_d$, and run gradient flow on
\[
  \Lhat(w)+\frac\lambda2\bigl(\norm u^2+\norm v^2\bigr),\qquad \Lhat(w)=\frac1{2n}\norm{Xw-y}^2 .
\]
Recall from \cref{ch:diagonal} the hyperbolic-entropy potential $\Phi_\alpha(w)=\frac{\alpha^2}4\sum_iq(w_i/\alpha^2)$ with $q(z)=2-\sqrt{4+z^2}+z\arcsinh(z/2)$, whose gradient is $\nabla\Phi_\alpha(w)_i=\frac14\arcsinh(w_i/2\alpha^2)$. Without weight decay the flow converges to the interpolator $w^\Phi_\alpha=\argmin\{\Phi_\alpha(w):Xw=y\}$, which is close to the minimum-$\ell_2$-norm interpolator when $\alpha$ is large and to the minimum-$\ell_1$-norm interpolator when $\alpha$ is small.

\begin{proposition}[Weight decay shrinks the initialization]\label[proposition]{prop:grokking:decay}
Let $\alpha_t=\alpha e^{-\lambda t}$ and $g(t)=\nabla\Lhat(w(t))$. For every $t\ge0$ and every coordinate $i$:
\begin{enumerate}
\item $u_i(t)v_i(t)=\alpha_t^2$;
\item $\frac{\mathrm d}{\mathrm dt}\nabla\Phi_{\alpha_t}(w(t))=-g(t)$;
\item $w(t)=\argmin\{\Phi_{\alpha_t}(w'):Xw'=Xw(t)\}$.
\end{enumerate}
\end{proposition}

\begin{proof}
The flow is $\dot u=-2u\odot g-\lambda u$ and $\dot v=2v\odot g-\lambda v$. It exists for all time, because the objective decreases along it and is at least $\frac\lambda2(\norm u^2+\norm v^2)$, so the trajectory stays bounded.

\emph{Part (i).} The product rule gives $\frac{\mathrm d}{\mathrm dt}(u_iv_i)=(-2g_i-\lambda)u_iv_i+(2g_i-\lambda)u_iv_i=-2\lambda u_iv_i$, and $u_i(0)v_i(0)=\alpha^2$.

\emph{Part (ii).} Let $\sigma_i=u_i^2+v_i^2$. By part (i), $\sigma_i^2-w_i^2=4u_i^2v_i^2=4\alpha_t^4$, so $\sigma_i=\sqrt{w_i^2+4\alpha_t^4}$, and the flow gives $\dot w_i=2u_i\dot u_i-2v_i\dot v_i=-4\sigma_ig_i-2\lambda w_i$. Let $z_i=w_i/(2\alpha_t^2)$. Since $\frac{\mathrm d}{\mathrm dt}\alpha_t^2=-2\lambda\alpha_t^2$, the quotient rule gives $\dot z_i=(\dot w_i+2\lambda w_i)/(2\alpha_t^2)=-2\sigma_ig_i/\alpha_t^2$, while $\sqrt{1+z_i^2}=\sigma_i/(2\alpha_t^2)$. Hence
\[
  \frac{\mathrm d}{\mathrm dt}\,\frac14\arcsinh z_i=\frac{\dot z_i}{4\sqrt{1+z_i^2}}=-g_i .
\]
\emph{Part (iii).} At $t=0$, $\nabla\Phi_\alpha(w(0))=0$, so by part (ii) $\nabla\Phi_{\alpha_t}(w(t))=-\int_0^tg=-\frac1nX^\top\int_0^t(Xw-y)$ lies in the row space of $X$. The potential $\Phi_{\alpha_t}$ is strictly convex, because $q''>0$, so a point $w'$ minimizes it on the affine set $\{Xw'=Xw(t)\}$ if and only if $\nabla\Phi_{\alpha_t}(w')$ lies in the row space; $w(t)$ is such a point, and the minimizer is unique.
\end{proof}

In words: on a diagonal network, weight decay acts exactly as an initialization scale that shrinks like $e^{-\lambda t}$, and at every moment the network sits at the predictor that the flow of \cref{ch:diagonal} would have chosen from the scale $\alpha_t$, among the predictors with its own training outputs. Nothing in the proof is asymptotic. The rest of the story needs one approximation.

\begin{heuristic}[Heuristic derivation of the grokking time]
Assume that the training residual stays negligible after an initial transient, so that $Xw(t)\approx y$. Then \cref{prop:grokking:decay} gives $w(t)\approx w^\Phi_{\alpha_t}$, and the test error at time $t$ is a fixed function of $\log\alpha-\lambda t$, read off from the static family $a\mapsto w^\Phi_a$.

\emph{Memorization.} At first $\sigma_i\approx2\alpha^2$, so $\dot w\approx-8\alpha^2\nabla\Lhat(w)$, and the residual $Xw-y$ decays at rate at least $8\alpha^2\lmin(XX^\top/n)$. The training set is fit by the time $\grokmem\approx\log(1/\eps)/\bigl(8\alpha^2\lmin(XX^\top/n)\bigr)$, while $\alpha_t\approx\alpha$ and $w\approx w^\Phi_\alpha$ is close to the minimum-norm interpolator, whose test error for a sparse target is near the null risk.

\emph{Generalization.} The static family passes from $\ell_2$-like to $\ell_1$-like behavior when $a$ crosses a critical scale $\alpha_c$, so the test error drops when $\alpha e^{-\lambda t}=\alpha_c$, that is, at
\begin{equation}\label{eq:grokking:tgrok}
  \grokt\approx\frac1\lambda\log\frac\alpha{\alpha_c} .
\end{equation}
\emph{The critical scale.} For $\abs{w_i}\ll a^2$ the potential is quadratic, $\Phi_a(w)\approx\sum_iw_i^2/(16a^2)$; for $\abs{w_i}\gg a^2$ it is $\Phi_a(w)\approx\frac14\sum_i\abs{w_i}(\log(\abs{w_i}/a^2)-1)$. The minimum-norm interpolator $w_2=X^+y$ is spread over all coordinates, with $\norm{w_2}^2\approx\frac nd\norm{w_\star}^2$, while the target is concentrated on $k$ of them, and the minimizer switches from one to the other when their potentials are comparable. For a target with entries $\pm1$ this gives, to logarithmic accuracy, $\norm{w_2}^2/(16a^2)\approx\frac14\norm{w_\star}_1\log(1/a^2)$, that is,
\[
  \alpha_c^2\log\frac1{\alpha_c^2}\approx\frac{\norm{w_2}^2}{4\norm{w_\star}_1}\approx\frac{n}{4d}.\qedhere
\]
\end{heuristic}

The ratio $\grokt/\grokmem$ grows like $\alpha^2\log(\alpha/\alpha_c)/\lambda$: a large initialization and a small weight decay make grokking dramatic. The prediction \eqref{eq:grokking:tgrok} has the form of \cref{thm:grokking:lyu}, with $\log\alpha$ refined to $\log(\alpha/\alpha_c)$; the refinement matters at the moderate $\alpha$ of experiments. The estimate of $\alpha_c$ is cruder than the rest: on random instances it is within a factor of about $1.5$ of the scale at which the static test error halves, an error that enters $\grokt$ only through a logarithm.

\section{Reading the solution}\label{sec:grokking:reading}

The model makes five quantitative predictions.
\begin{enumerate}
\item The logits of every max-margin network are the clock sum $\gamma(\ind\{c=a+b\}-\frac1p)$: the wrong classes tie. As vectors in $\R^p$, the logit vectors of the $p$ possible sums are the vertices of a regular simplex centered at the origin, the geometry of neural collapse (\cref{ch:collapse}).
\item Every neuron of a max-margin network has a single frequency, equal norms $\norm u=\norm{u'}=\norm v$, and phase defect $\chi-\varphi-\psi=0$, where $\varphi$, $\psi$ and $\chi$ are the phases of $u$, $u'$ and $v$ at their common frequency.
\item All $\frac{p-1}2$ frequencies appear, each in at least six neurons, and the maximal normalized margin is $\gamma^\star=\sqrt{2/27}/(\sqrt p(p-1))$, for instance $6.1\times10^{-4}$ for $p=59$.
\item In the diagonal caricature, the training set is fit at $\grokmem\propto\alpha^{-2}$ and the test error drops at $\grokt\approx\lambda^{-1}\log(\alpha/\alpha_c)$; the test error is a function of $\log\alpha-\lambda t$ alone.
\item The delay $\grokt/\grokmem$ grows like $\alpha^2/\lambda$, up to a logarithm.
\end{enumerate}

\begin{figure}[t]
\begin{fullwidth}
\centering
\begin{tikzpicture}
\begin{groupplot}[group style={group size=2 by 1,horizontal sep=17mm},hplot,width=0.475\linewidth,height=0.37\linewidth]
\nextgroupplot[xmode=log,xmin=1e-5,xmax=3000,ymin=-0.03,ymax=1.03,xlabel={time $t$},ylabel={error},
  xtick={1e-4,1e-2,1,1e2},legend style={at={(0.03,0.04)},anchor=south west}]
\addplot[theory3] table[x=t,y=trainrel] {figures/data/10-grokking-diagonal.dat};
\addplot[simulation,mark repeat=3] table[x=t,y=test] {figures/data/10-grokking-diagonal.dat};
\addplot[theory] table[x=t,y=theory] {figures/data/10-grokking-diagonal.dat};
\addplot[guide] coordinates {(429,-0.03) (429,1.03)};
\legend{training error,test error (flow),test error of $w^\Phi_{\alpha e^{-\lambda t}}$}
\nextgroupplot[xmode=log,xmin=0.8,xmax=125,ymin=1.5,ymax=7,xlabel={initialization scale $\alpha$},ylabel={$\lambda\,t_{1/2}$},
  legend style={at={(0.03,0.97)},anchor=north west}]
\addplot[theory,domain=0.9:110,samples=50] {ln(x/0.1369)};
\addplot[simulation] table[x=alpha,y=sim0003] {figures/data/10-grokking-delay.dat};
\addplot[simulation2] table[x=alpha,y=sim001] {figures/data/10-grokking-delay.dat};
\addplot[only marks,mark=triangle*,mark size=2pt,color=sea,mark options={fill=seamist}] table[x=alpha,y=sim003] {figures/data/10-grokking-delay.dat};
\legend{$\log(\alpha/\alpha_c)$,$\lambda=0.003$,$\lambda=0.01$,$\lambda=0.03$}
\end{groupplot}
\end{tikzpicture}
\caption{Grokking in a diagonal linear network with weight decay ($d=100$, $n=30$, a $3$-sparse target with entries $1$, Gaussian design, seed $0$). The flow of $(u,v)$ is integrated with a stiff solver. Left: $\alpha=10$, $\lambda=0.01$. The training error $\norm{Xw-y}^2/\norm y^2$ falls below $10^{-4}$ by $t=0.005$; the relative test error $\norm{w-w_\star}^2/\norm{w_\star}^2$ stays at $0.72$, the value for the minimum-norm interpolator, until $t\approx300$ and falls below $0.01$ by $t\approx690$. The curve is the quasi-static prediction of \cref{prop:grokking:decay}, the test error of $\argmin\{\Phi_{\alpha_t}(w):Xw=y\}$; the dashed line is $\lambda^{-1}\log(\alpha/\alpha_c)=429$, where $\alpha_c=0.137$ is the scale at which that static test error halves. Right: the time $t_{1/2}$ at which the simulated test error halves, for three weight decays and seven initialization scales, against $\log(\alpha/\alpha_c)$; all $21$ points lie within $0.06$ of the line.}\label{fig:grokking:diagonal}
\end{fullwidth}
\end{figure}

\Cref{fig:grokking:diagonal} tests the delay. The left panel is grokking in miniature: the training error falls at $t\approx10^{-3}$, the test error at $t\approx4\times10^2$, and between them the test error sits at the null-like value $0.72$ for five decades of time. After the fast transient the simulated test error is indistinguishable from the quasi-static curve, which is the test error of the static interpolator $w^\Phi_a$ evaluated at $a=\alpha e^{-\lambda t}$, computed without reference to the dynamics. The right panel shows the collapse that the proposition predicts: $\lambda t_{1/2}$ depends on $\alpha$ only through $\log(\alpha/\alpha_c)$, with the same $\alpha_c$ for every $\lambda$. The small upward offset grows with $\lambda$, from $0.005$ at $\lambda=0.003$ to $0.05$ at $\lambda=0.03$; it is the training residual of order $\lambda$ that the quasi-static assumption neglects. The heuristic estimate $\alpha_c^2\log(1/\alpha_c^2)\approx\norm{w_2}^2/(4\norm{w_\star}_1)$ gives $\alpha_c=0.131$ here and $\grokt=434$.

\begin{figure}[t]
\begin{fullwidth}
\centering
\begin{tikzpicture}
\begin{groupplot}[group style={group size=2 by 1,horizontal sep=17mm},hplot,width=0.475\linewidth,height=0.37\linewidth]
\nextgroupplot[xmode=log,xmin=1,xmax=10000,ymin=-0.03,ymax=1.03,xlabel={AdamW step},ylabel={accuracy or concentration},
  legend style={at={(0.03,0.97)},anchor=north west}]
\addplot[theory3] table[x=step,y=train] {figures/data/10-grokking-modadd-curves.dat};
\addplot[simulation,mark repeat=2] table[x=step,y=test] {figures/data/10-grokking-modadd-curves.dat};
\addplot[theory2] table[x=step,y=conc] {figures/data/10-grokking-modadd-curves.dat};
\legend{training accuracy,test accuracy,Fourier concentration}
\nextgroupplot[xmin=0,xmax=30,ymin=-3.3,ymax=3.3,xlabel={frequency $\zeta$ of the neuron},ylabel={phase defect $\chi-\varphi-\psi$},
  ytick={-3.1416,-1.5708,0,1.5708,3.1416},yticklabels={$-\pi$,$-\frac\pi2$,$0$,$\frac\pi2$,$\pi$},legend style={at={(0.03,0.97)},anchor=north west}]
\addplot[guide] coordinates {(0,0) (30,0)};
\addplot[simulation,restrict expr to domain={\thisrow{relnorm}}{0.3:1.1}] table[x=zeta,y=defect] {figures/data/10-grokking-modadd-neurons.dat};
\addplot[only marks,mark=o,mark size=1.4pt,color=muted,restrict expr to domain={\thisrow{relnorm}}{0:0.3}] table[x=zeta,y=defect] {figures/data/10-grokking-modadd-neurons.dat};
\legend{theory,large neurons,small neurons}
\end{groupplot}
\end{tikzpicture}
\caption{The quadratic network \eqref{eq:grokking:net} trained on modular addition: $p=59$, width $m=256$, $30\%$ of the table for training, full-batch AdamW (learning rate $10^{-3}$, $\beta=(0.9,0.98)$, weight decay $1$), cross-entropy, seed $0$. Left: training accuracy reaches $99\%$ at step $177$, test accuracy at step $5960$. The Fourier concentration is the share of each neuron's first-layer energy at its dominant frequency $\pm\zeta$, averaged with weights $\norm{\theta_r}^3$; it rises steadily from $0.10$ before the test accuracy moves. Right: after $10^4$ steps, the phase defect of every neuron at its dominant frequency. The $140$ large neurons (norm above $0.3$ times the largest) have median concentration $0.91$, defects below $0.13$ in absolute value (median $0.015$), and cover all $29$ frequencies, with $3$ to $8$ of them per frequency.}\label{fig:grokking:modadd}
\end{fullwidth}
\end{figure}

\Cref{fig:grokking:modadd} tests the features. The trained network groks, with a delay of a factor of about $34$, and at the end of training its neurons have the form that \cref{thm:grokking:neuron} prescribes: the large ones are concentrated on single frequencies, and their readout phases equal the sum of their input phases to within a few hundredths of a radian, although nothing in training refers to phases. All $29$ frequencies are present, as \cref{thm:grokking:maxmargin} predicts for max-margin networks and as \textcite{MorwaniEtAl2024} found in their own experiments with quadratic and ReLU networks. The Fourier concentration rises long before the test accuracy, a progress measure in the sense of \textcite{NandaEtAl2023}.

Two predictions fail, and they fail informatively. The network is far from the margin limit: at step $10^4$ its smallest margin is $0.22$ times its average margin, whereas a max-margin network has all margins equal, and its normalized margin is $0.145\,\gamma^\star$. Six frequencies are carried by fewer than six large neurons, so the unwanted terms of \cref{lem:grokking:block} do not cancel exactly; the network classifies correctly anyway, because correct classification needs only the largest logit to be right, not the wrong logits to tie. The features of the max-margin solution are in place long before its margin is.

What about the delay in this network? With weight decay $0.5$, $1$ and $2$ the test accuracy reaches $99\%$ after about $15000$, $6000$ and $4000$ steps, so the product of weight decay and delay stays between $6000$ and $8000$ steps, as $\grokt\propto\lambda^{-1}$ suggests. But with weight decay $4$ the network has not generalized after $40000$ steps, and with $0.3$ it needs more than $30000$ (\cref{ex:grokking:wdsweep}). The diagonal model cannot see the upper end, where weight decay competes with fitting, the confusion phase of \textcite{LiuEtAl2022}.

\section{What the model gets wrong}\label{sec:grokking:wrong}

\paragraph{What was thrown away.} Four simplifications, each with a price. \emph{The task} is an abelian group, so the characters are known in advance and diagonalize everything; natural data have no such basis, and the analysis has nothing to say about them. \emph{The architecture} is a two-layer network with a polynomial activation, which makes the Fourier analysis exact; there is no attention, no factorization into embedding and unembedding, no normalization, and no ReLU. \emph{The margin theory} describes networks at the end of training on the full table: it uses the limit of vanishing weight decay, the $\ell_{2,3}$ norm instead of the $\ell_2$ norm that weight decay penalizes, and all $p^2$ pairs as training data. It therefore explains which features a generalizing network uses, but not why training on a fraction of the table generalizes, which is the phenomenon. \emph{The delay model} is a reparameterized linear model on a sparse regression task, whose only ``feature'' is a choice of coordinates; it shows how weight decay moves a network from the lazy to the rich regime, but in a model with no circuit to form.

\paragraph{Where it breaks.} \emph{Different networks find different algorithms.} \Textcite{ZhongEtAl2023} found that transformers with weak attention learn a ``pizza'' algorithm, which averages the embeddings of $a$ and $b$ instead of multiplying them, so that all pairs with the same sum lie on a line through the origin; its logits carry an extra factor $\abs{\cos(\omega(a-b)/2)}$ and are not clocks. Which algorithm appears depends on the architecture, and margin maximization within one architecture does not predict it.\sidenote{Zhong and coauthors interpolate between constant attention, which learns the pizza, and standard attention, which learns the clock, by an attention-rate parameter; the transition between the two is continuous.} \emph{Transformers use few frequencies.} The transformer of \textcite{NandaEtAl2023} uses five key frequencies, and other runs between two and nine, while every max-margin network of \cref{thm:grokking:maxmargin} uses all $56$; our MLP agrees with the theory, the transformers do not. \emph{Grokking happens without weight decay.} \Textcite{KumarEtAl2024} observe it with plain gradient descent, controlled by laziness; \textcite{ThilakEtAl2022} (unrefereed) observe it with Adam and no regularization, at the onset of cyclic instabilities in which the norm of the last layer grows abruptly; and \textcite{PrietoEtAl2025} show that without regularization the logits grow until floating-point absorption makes the softmax saturate and the gradient vanish exactly, halting learning, while the gradient aligns with a direction that scales the logits without changing predictions. In these settings the clock that sets the delay is not $\lambda^{-1}\log(\alpha/\alpha_c)$. \emph{Grokking happens without circuits.} \Textcite{LeviEtAl2024} show that a linear estimator on Gaussian data groks in thresholded accuracy, with a delay set by the Marchenko--Pastur spectrum (\cref{ch:doubledescent}) while the losses evolve smoothly; \textcite{VarmaEtAl2023} (unrefereed) explain grokking by the relative efficiency of memorizing and generalizing circuits, which predicts a critical dataset size and ``ungrokking'' below it, an axis the initialization scale does not capture. \emph{Non-abelian groups.} For finite groups the characters give way to irreducible representations, and \textcite{ChughtaiEtAl2023} found networks composing representation matrices, with the set of representations learned varying from seed to seed; but \textcite{StanderEtAl2024}, reverse-engineering networks trained on $S_5$ and $S_6$, found neurons that detect membership in cosets of subgroups and little evidence that the embeddings store representation matrices. Even the group-theoretic generalization of the clock can misdescribe what a network computes.

\paragraph{Fine structure.} The repairs come in four directions. \emph{Representation theory} replaces Fourier analysis for finite groups: \textcite{MorwaniEtAl2024} characterize max-margin networks for general group operations, and \textcite{ChughtaiEtAl2023} describe the algorithm by matrix products. \emph{Algebraic structure of optima}: \textcite{Tian2025} decomposes the squared loss of quadratic networks on abelian groups frequency by frequency and builds global minimizers by composing small blocks, including the six-neuron block above and narrower mixtures of width $2p$ with non-Fourier neurons, and reports that most gradient-descent solutions match these constructions. \emph{Dynamics and phase transitions}: \textcite{RubinEtAl2024} treat grokking in two-layer networks, including the quadratic network on modular addition, as a first-order phase transition in an adaptive-kernel theory of the Bayesian posterior; \textcite{ZunkovicIlievski2024} compute exact critical exponents in teacher--student models with local rules; and \textcite{BeckEtAl2025} place grokking in logistic regression at the edge of linear separability. \emph{Fourier features elsewhere}: recursive feature machines, kernel machines with learned feature metrics, grok modular arithmetic with block-circulant features that implement the same Fourier algorithm \parencite{MallinarEtAl2025}; GPT-2-XL fine-tuned on addition represents numbers with sparse Fourier features, low frequencies in the MLPs for magnitude and high frequencies in attention for the residue \parencite{ZhouEtAl2024}; and language models represent days of the week and months of the year on circles \parencite{EngelsEtAl2025}, a cross-link to \cref{ch:superposition}. What is open is a dynamical theory: a proof that gradient descent with weight decay on a fraction of the table finds the circuit, with the time it takes and the set of frequencies it selects.

\section{Model card}\label{sec:grokking:card}

\modelcard{A network fits part of a group's table early and generalizes much later; once it generalizes, it computes with characters of the group.}{Two-layer network with quadratic activation and one-hot inputs on $a+b\bmod p$; for the delay, a diagonal linear network with weight decay and large initialization.}{An explicit Fourier circuit; max-$\ell_{2,3}$-margin networks consist of Fourier neurons with additive phases, use every frequency with at least six neurons each, and exist exactly at width $m\ge3(p-1)$; weight decay acts as an initialization scale $\alpha e^{-\lambda t}$.}{Logits $\propto\sum_\zeta\cos(\omega_\zeta(a+b-c))$; $\gamma^\star=\sqrt{2/27}/(\sqrt p(p-1))$; phase defect $0$; $\grokt\approx\lambda^{-1}\log(\alpha/\alpha_c)$ and $\grokmem\propto\alpha^{-2}$.}{Transformers use few frequencies or the pizza algorithm; trained networks are far from max margin; grokking without weight decay or without circuits; coset circuits for permutation groups.}{Representation theory of finite groups; algebraic composition of optima; lazy-to-rich dynamics; phase-transition theories; Fourier features in language models.}

The chapter sits on two threads of the book. The first is implicit bias: the max-margin argument is the one of \cref{ch:margin}, transported to a network that is homogeneous of degree three, and the diagonal model of \cref{ch:diagonal} supplies both the potential and the exact effect of weight decay. The second is lazy versus rich training: the memorization phase is the kernel regime of \cref{ch:lazy}, and grokking is the crossing from that regime to the rich one, which weight decay forces at a predictable time. The search phase of \cref{ch:singleindex} and the hidden progress of \cref{ch:parity} are the same story told with small initialization, where the plateau comes before fitting rather than after.

Fourier analysis on groups links this chapter to the Boolean cube of \cref{ch:parity}, where the characters are parities, and to the circular features of \cref{ch:superposition}. The equal-margin logits of a max-margin network form a regular simplex, the configuration that \cref{ch:collapse} derives for the last layer of a classifier.

\begin{exercises}
\item \consolidation (The normalized margin of the circuit.) (a) Show that the network of \cref{thm:grokking:fourier} has normalized margin $0.39\,\gamma^\star$, to two digits. (b) Rescale each neuron, multiplying $u$ and $u'$ by $t=(2/9)^{1/3}$ and $v$ by $t^{-2}$, and show that the normalized margin becomes exactly $\gamma^\star$. Why does rescaling change the normalized margin but not the logits? \label{ex:grokking:normalized}

\item \consolidation\hard (Six phases suffice, and tight rules have equal weights.) (a) Verify the six-neuron identity of \cref{rem:grokking:designs}, and check that its points $(\chi_r,\delta_r)$ are $(\frac{\pi r}3,\pi r)$, $r=0,\dots,5$, up to order. (b) Suppose that six points with positive weights satisfy \eqref{eq:grokking:moments}. Show that every weight equals $\frac16\sum_r\beta_r$. (c) Conclude that the nine phases of \cref{thm:grokking:fourier} are not minimal, and that six is the least number of Fourier neurons per frequency. \hint{for (b), use the space of \cref{lem:grokking:six}: the evaluation map at the six points is a bijection, so for each point there is a $P_r$ that is $1$ there and $0$ at the others. Compute $\sum_s\beta_s\abs{P_r(\chi_s,\delta_s)}^2$ in two ways, and use Cauchy--Schwarz in the form $1=P_r(\chi_r,\delta_r)\le\sqrt6\,(\sum\abs{c_{js}}^2)^{1/2}$.}\label{ex:grokking:phases}

\item \consolidation (Why the full table.) (a) Repeat \cref{lem:grokking:cw} with the wrong class drawn from a weighting $\tau(a,b)\in\Delta(\grokZp\setminus\{a+b\})$ that is not uniform, and show that the squares $u(a)^2$ and $u'(b)^2$ no longer cancel in general. (b) Repeat it with $(a,b)$ uniform on a training set $D$ instead of on $\grokZp^2$, and show that the squares cancel if $D$ is a union of full rows $\{a\}\times\grokZp$, but not for a typical random $D$. What does this say about using the margin theory to explain generalization from a fraction of the table? \label{ex:grokking:subset}

\item \consolidation (Weight decay and the Lasso.) (a) Show that the minimum of $\norm u^2+\norm v^2$ subject to $u\odot u-v\odot v=w$ is $\norm w_1$, so that the regularized objective of \cref{sec:grokking:delay} has the same minimum value as the Lasso $\Lhat(w)+\frac\lambda2\norm w_1$, and its global minimizers map to Lasso solutions. (b) For gradient descent with step $\eta$, show that $u_iv_i$ is multiplied by $(1-\eta\lambda)^2-4\eta^2g_i^2$ at each step, so part (i) of \cref{prop:grokking:decay} holds up to a relative error of order $\eta^2$ per step. (c) Derive the memorization time of the heuristic derivation and evaluate it for the run of \cref{fig:grokking:diagonal}. \label{ex:grokking:lasso}

\item \consolidation (A simplex of logits.) Let $\theta$ be a max-margin network. (a) Show that its logit vector $f(a,b)\in\R^p$ depends on $(a,b)$ only through $a+b$, and that the $p$ distinct vectors have equal norms, sum to zero, and have pairwise inner products $-\frac1{p-1}$ times their squared norm. (b) Compare with the simplex equiangular tight frame of \cref{ch:collapse}. \label{ex:grokking:simplex}

\item \variation\hard (Multiplication and even moduli.) (a) Let $g$ generate the cyclic group $\grokZp^\times$ of order $p-1$. Show that $c=ab$ in $\grokZp^\times$ if and only if $\log_gc=\log_ga+\log_gb$ in $\Z_{p-1}$, and build a Fourier network for multiplication. (b) The modulus $N=p-1$ is even, so the frequency $N/2$, whose character $(-1)^x$ is the Legendre symbol after the change of variables, is real and has no conjugate partner. Redo \cref{thm:grokking:neuron} on $\Z_N$ and show that the best single neuron sits at frequency $N/2$, with average margin $\frac{2}{3\sqrt3(N-1)\sqrt N}$, a factor $\sqrt2$ above the best paired frequency. (c) Show that a network whose neurons all sit at frequency $N/2$ gives every wrong class of the same parity as $a+b$ a margin of zero, and conclude that the proof of \cref{thm:grokking:maxmargin} fails for even moduli. \label{ex:grokking:even}

\item \variation (A non-abelian group.) Let $G$ be a finite group with irreducible characters $\chi_\rho$ of dimensions $d_\rho$. (a) Use the orthogonality of characters to show that $\sum_\rho d_\rho\chi_\rho(g)=\abs G\,\ind\{g=e\}$. (b) Write the character table of $S_3$ and check (a). (c) Deduce that the logits $f_c(a,b)=\sum_\rho d_\rho\chi_\rho(c^{-1}ab)-1$ give every margin $\abs G$, the analogue of \cref{thm:grokking:fourier}. (d) Write $\chi_\rho(c^{-1}ab)$ as a trace of a product of three matrices, and explain why a network needs products of matrix entries rather than of numbers; compare with the algorithm of \textcite{ChughtaiEtAl2023}. \label{ex:grokking:s3}

\item \variation (The Euclidean norm.) (a) Show that $\min_{t>0}\bigl(t^2(\norm u^2+\norm{u'}^2)+t^{-4}\norm v^2\bigr)=3\cdot2^{-2/3}(\norm u^2+\norm{u'}^2)^{2/3}\norm v^{2/3}$, so that the weight-decay cost of a neuron grows like the power $\frac23$ of the size of its contribution to the logits. (b) Show that under $\norm\theta_2\le1$ the average margin is maximized only by networks with a single nonzero neuron, and that for $p\ge5$ such a network does not have equal margins. Conclude that the proof of \cref{thm:grokking:maxmargin} does not apply to the norm that weight decay penalizes. \label{ex:grokking:l2}

\item \variation (Unbalanced initialization.) In the diagonal network start from $u(0)=\alpha\mathbf1_d$ and $v(0)=\alpha'\mathbf1_d$ with $\alpha\ne\alpha'$. Show that $u_iv_i=\alpha\alpha'e^{-2\lambda t}$, that part (ii) of \cref{prop:grokking:decay} holds with $\alpha_t^2=\alpha\alpha'e^{-2\lambda t}$, and that $w(t)$ minimizes $\Phi_{\alpha_t}(w')-\frac14\log(\alpha/\alpha')\sum_iw'_i$ on $\{Xw'=Xw(t)\}$. Which sign of target coordinates does a larger $\alpha$ favor? \label{ex:grokking:unbalanced}

\item \simulation (Weight decay and the delay in the network.) Run the script of \cref{fig:grokking:modadd} with weight decay $0.3$, $0.5$, $1$, $2$ and $4$ and $40000$ steps, and record the first step at which the test accuracy reaches $99\%$. The reader should observe about $15000$, $6000$ and $4000$ steps for $0.5$, $1$ and $2$, consistent with $\grokt\propto\lambda^{-1}$ up to a slowly varying logarithm; more than $30000$ for $0.3$; and no generalization within $40000$ steps for $4$, where the network memorizes later and weight decay competes with fitting. Plot the weight norm against time in each run and compare with the effective scale $\alpha e^{-\lambda t}$ of the diagonal model. \label{ex:grokking:wdsweep}

\item \simulation (The critical scale.) Using the script of \cref{fig:grokking:diagonal}, (a) set $\lambda=0$ and check that the test error stays at that of $w^\Phi_\alpha$ forever; (b) for $n/d\in\{0.2,0.3,0.4\}$ and several random designs, compute the scale $\alpha_{1/2}$ at which the static test error halves and compare it with the root of $\alpha_c^2\log(1/\alpha_c^2)=\norm{w_2}^2/(4\norm{w_\star}_1)$. The reader should find agreement within a factor of about $1.5$, and an error in $\lambda\grokt$ of about $\log1.5\approx0.4$, small against $\log(\alpha/\alpha_c)$ when $\alpha$ is large. \label{ex:grokking:alphac}

\item \research (Dynamics on a fraction of the table.) Prove that gradient descent with small weight decay, started from a small random initialization and trained on a random fraction of the table, converges to a network whose neurons are close to Fourier neurons, and estimate the time. \Textcite{MohamadiEtAl2024} prove generalization bounds for regularized global minimizers with large $\ell_\infty$ margin, and \textcite{Tian2025} describes the global minimizers of a squared loss; a proof for the dynamics from random initialization, and a prediction of which frequencies are selected and when, are open. \label{ex:grokking:dynamics}

\item \research (The narrowest circuit.) What is the smallest width of a quadratic network \eqref{eq:grokking:net} whose logits equal the clock sum \eqref{eq:grokking:target} exactly? Fourier neurons need $3(p-1)$, and \textcite{Tian2025} reaches width $2p$ with other neurons. Numerical searches suggest that one frequency cannot be represented exactly by four neurons of that frequency, and that with five the error can be made small only with very large weights. A matching lower bound, and the analogous question for the $\ell_2$ max-margin problem of \cref{ex:grokking:l2}, are open. \label{ex:grokking:narrow}
\end{exercises}

\notesandrefs
Grokking was reported by \textcite{PowerEtAl2022}, and its explanation in terms of Fourier features comes from the mechanistic study of \textcite{NandaEtAl2023}. \Textcite{Gromov2023} trained the two-layer quadratic network of this chapter, with a squared loss, observed the same features, and wrote down an approximate solution with random phases. The effective theory of \textcite{LiuEtAl2022} and the weight-norm picture of \textcite{LiuEtAl2023} first tied the delay to hyperparameters. \Textcite{DaviesEtAl2023} (a workshop paper) relate grokking to epoch-wise double descent, and \textcite{DoshiEtAl2024} study it with corrupted labels, where memorizing and generalizing neurons can be told apart.

The margin analysis is that of \textcite{MorwaniEtAl2024}, with the class-weighted margin as the device that linearizes the problem; their paper treats sparse parities and general finite groups as well. The general margin results behind it are those of \textcite{WeiEtAl2019} and \textcite{LyuLi2020}, discussed in \cref{ch:margin}. The eight-neuron blocks are Morwani and coauthors'; the six-neuron block appears in \textcite{Tian2025}, whose algebraic framework also produces the narrower mixtures mentioned in \cref{sec:grokking:circuit}. The nine-phase grid of \cref{thm:grokking:fourier} is the most symmetric design, and the cubature view of \cref{lem:grokking:block,lem:grokking:six} organizes all of them.

For the delay, \textcite{LyuEtAl2024} prove the dichotomy between early kernel and late margin biases for homogeneous networks, with diagonal linear networks and matrix completion as examples, and they also exhibit ``misgrokking'', in which the early bias generalizes and the late one does not. \Textcite{KumarEtAl2024} argue for the lazy-to-rich transition without weight decay, building on the lazy regime of \textcite{ChizatEtAl2019} and the kernel and rich regimes of diagonal networks of \textcite{WoodworthEtAl2020}. \Textcite{MohamadiEtAl2024} prove the kernel lower bound for modular addition and generalization bounds for quadratic networks with large margin. The weight-decay identity of \cref{prop:grokking:decay} is a short computation in the setting of \cref{ch:diagonal}. AdamW is due to \textcite{LoshchilovHutter2019}; its decoupled weight decay multiplies the parameters by $1-\eta\lambda$ at each step, so that $\eta\lambda$ is the decay rate per step in our experiments. A related account explains grokking by the competition of a sparse and a dense subnetwork \parencite{MerrillEtAl2023}.

\paragraph{Translating the literature.} \Textcite{MorwaniEtAl2024} write a neuron as $(u_a+v_b)^2w$: their $u$, $v$, $w$ are our $u$, $u'$, $v$; their frequency $\zeta$ and phases $\theta^*_u$, $\theta^*_v$, $\theta^*_w$ are our $\zeta$, $\varphi$, $\psi$, $\chi$; their class-weighted margin with uniform weighting is our average margin $\grokcw$; their $L_{2,3}$ norm is our $\norm\cdot_{2,3}$; their $\gamma^*$ is our $\gamma^\star$; their Fourier transform is unnormalized, as ours is. \Textcite{NandaEtAl2023} write $w_k=2\pi k/p$ for the key frequencies, our $\omega_\zeta$ with $\zeta=k$. \Textcite{Gromov2023} writes $N$ for the width (our $m$), $n$ and $m$ for the two inputs (our $a$ and $b$), $q$ for the output (our $c$), and $\phi^{(1)}$, $\phi^{(2)}$, $\phi^{(3)}$ for the phases (our $\varphi$, $\psi$, $\chi$). \Textcite{Tian2025} writes $d$ for the order of the group (our $p$), and the order of a block is our number of neurons per frequency. \Textcite{LyuEtAl2024} use our $\alpha$ and $\lambda$, write $\bar\theta_{\mathrm{init}}$ for our $\bar\theta$, and write $T^\pm_c(\alpha)=(1\pm c)\lambda^{-1}\log\alpha$ for the two times of \cref{thm:grokking:lyu}.

\printchapterbib
\end{refsection}

% Chapter 11 of The Hydrogen Atoms of Deep Learning: superposition in the ReLU toy autoencoder.
\newcommand{\supv}{v(S)}
\newcommand{\supOmega}{\Omega}
\newcommand{\supD}{\mathrm{D}}

\casechapter{Superposition}{More features than dimensions}{ReLU toy autoencoder}{ch:superposition}
\begin{refsection}[chapters/11-superposition.bib]

\begin{chaptersummary}
Neurons of trained networks often respond to several unrelated things, and the activations of language models decompose into sparse combinations of many more directions than there are dimensions. The toy model of this chapter stores $p$ sparse features in $m<p$ dimensions with a linear map and reads them out with a ReLU. We prove that without the ReLU the optimum is principal component analysis at every sparsity, so a linear network never stores more features than dimensions; that with it, two features share one dimension as an antipodal pair once they are sparse enough, with exact losses and a provable region of the phase diagram; that sparse features fit without first-order interference exactly when their directions have pairwise nonpositive inner products, which is possible for at most $2m$ of them; and that random directions and a thresholded readout store $e^{cm}$ features if few are active at a time. The model knows nothing of correlated or multidimensional features, of computation, or of how sparse autoencoders choose among decompositions.
\end{chaptersummary}

\section{The phenomenon: polysemantic neurons and sparse dictionaries}\label{sec:superposition:phenomenon}

Look at what makes a single neuron of a trained vision network fire, and the answer is often a list of things with nothing in common. In their survey of the circuits of InceptionV1, \textcite{OlahEtAl2020} found many \term{polysemantic neurons}\index{polysemantic neuron}, neurons that respond to several unrelated inputs, and named them a main obstacle to reading networks neuron by neuron.\sidenote{Neuroscience met the same fact earlier under the name \emph{mixed selectivity}: \textcite{RigottiEtAl2013} found that neurons in monkey prefrontal cortex respond to combinations of task variables, and that every variable can be decoded from the population although single cells mix them.} Word embeddings tell the same story in a form one can compute with. \Textcite{AroraEtAl2018} showed that the embedding of a polysemous word such as \emph{tie} lies close to a linear combination of vectors for its senses, and that sparse coding of the embeddings against a dictionary of about $2000$ ``discourse atoms'' recovers those senses.

\Textcite{ElhageEtAl2022} turned these observations into a hypothesis. A network may represent more \emph{features}, properties of the input that it uses, than it has dimensions, by assigning each feature a direction, with the directions almost but not exactly orthogonal. The cost is \emph{interference}: reading one feature picks up a little of every other. If features are \emph{sparse}, so that few are present in a given input, the interference is usually small, and a nonlinearity can remove what remains. They called this \term{superposition}\index{superposition}, studied it in the toy model of this chapter, and found that whether a feature is stored in superposition changes abruptly with its sparsity and importance, and that stored features arrange themselves in regular polytopes.

If the hypothesis is right, then the right units for reading a network are directions, not neurons, and they can be found by sparse dictionary learning: write each activation vector $h\in\R^m$ as $h\approx\sum_{i}a_i(h)\,d_i$ with many more dictionary vectors $d_i$ than dimensions and few nonzero coefficients $a_i(h)\ge0$, as \textcite{OlshausenField1997} did for natural images. A \term{sparse autoencoder}\index{sparse autoencoder} computes the coefficients with one ReLU layer and learns the dictionary. \Textcite{CunninghamEtAl2024} trained such autoencoders on the residual stream of small language models of the Pythia family and found directions more interpretable than neurons. \Textcite{BrickenEtAl2023} decomposed the $512$ neurons of a one-layer transformer's MLP into more than $4000$ features, among them features for DNA sequences, legal language, HTTP requests, and Hebrew text, invisible in single neurons. \Textcite{TempletonEtAl2024} scaled the method to a production model, Claude 3 Sonnet, with dictionaries of up to $34$ million features, many of them multilingual, multimodal, and usable to steer the model, and \textcite{GaoEtAl2025} trained a $16$-million-latent autoencoder on GPT-4 activations and measured scaling laws in the dictionary size.

\begin{observation}[\textcite{OlahEtAl2020}; \textcite{ElhageEtAl2022}; \textcite{BrickenEtAl2023}; \textcite{TempletonEtAl2024}]\label[observation]{obs:superposition:poly}
Individual neurons of trained networks are often polysemantic. The activation vectors of a layer of width $m$ are approximated, with a reconstruction error that is small but not negligible, by sparse nonnegative combinations of $p\gg m$ learned directions, and many of these directions are interpretable.
\end{observation}

The status of the observation should be stated carefully. Polysemanticity is robust and old. That activations \emph{are} sparse combinations of a dictionary is a modeling claim, supported by the interpretability of the recovered directions, and challenged by three facts we return to in \cref{sec:superposition:wrong}: the reconstructions leave a structured error, the dictionaries depend on the random seed, and some features are not one-dimensional. The question for this chapter is the one that comes before all of this. \emph{When and why does a network store more features than it has dimensions, how many can it store, and at what cost?}

\section{The model: a ReLU autoencoder for sparse features}\label{sec:superposition:model}

\paragraph{The features.} The input is a vector $x\in[0,1]^p$ of $p$ features with independent coordinates, each equal to $0$ with probability $S$ and uniform on $[0,1]$ otherwise. The number $S\in[0,1)$ is the \term{sparsity}\index{sparsity}; it is local notation, as is the name $p$ for the number of features, which matches the global use of $p$ for the number of features of a linear model. A feature is \emph{active} when it is nonzero. One feature has mean, second moment, and variance
\begin{equation}\label{eq:superposition:moments}
  \E x_i=\frac{1-S}2,\qquad \E x_i^2=\frac{1-S}3,\qquad \supv\defeq\Var x_i=\frac{(1-S)(1+3S)}{12}.
\end{equation}
Each feature carries an \term{importance}\index{importance of a feature} $\omega_i>0$, and $\supOmega=\sum_i\omega_i$.

\paragraph{The network.} A matrix $W\in\R^{m\times p}$ with $m<p$ compresses the features to $h=Wx\in\R^m$, and the transpose decodes them through a ReLU with bias $b\in\R^p$:
\begin{equation}\label{eq:superposition:model}
  x'=\relu\bigl(W^\top Wx+b\bigr),\qquad \Loss(W,b)=\sum_{i=1}^p\omega_i\,\E\bigl(x_i-x'_i\bigr)^2 .
\end{equation}
We write $w_i\in\R^m$ for the $i$-th column of $W$, the direction of feature $i$, and $G=W^\top W$ for the Gram matrix of the directions. Coordinate $i$ of the output reads
\[
  x'_i=\relu\Bigl(\norm{w_i}^2x_i+\sum_{j\ne i}\ip{w_i,w_j}\,x_j+b_i\Bigr):
\]
the signal, the \term{interference}\index{interference} from the other features, and a bias. The \emph{linear model} is the same network without the ReLU, $x'=W^\top Wx+b$.

\begin{definition}[Representation and superposition]\label[definition]{def:superposition:sup}
A configuration $(W,b)$ \emph{represents} feature $i$ if $w_i\ne0$, and it \emph{drops} feature $i$ if $w_i=0$. It is in \term{superposition} if it represents more than $m$ features.
\end{definition}

Represented directions in superposition cannot all be orthogonal, so superposition means interference. \Textcite{ElhageEtAl2022} also measure superposition by $\norm{W}_F^2/m$, the number of features per dimension, which equals $1$ when $W^\top W$ is an orthogonal projection of rank $m$; we use it in \cref{sec:superposition:reading}.

\paragraph{What was thrown away.} The features are given rather than learned from data, independent, identically distributed except for their importances, and of one kind (a magnitude in $[0,1]$). The task is to reconstruct them, not to compute with them. The decoder is tied to the encoder, there is one layer, and the loss is the population loss. Each simplification removes something that would hide the trade-off we want to see: with independent features, interference is pure noise; with reconstruction, the target is the identity and every error is visible; tied weights make the geometry of $W$ the whole story.

\begin{dictionary}
features $x_i$, sparsity $S$, importance $\omega_i$ & properties of the input a network uses (a word sense, a language, a concept), present rarely, of unequal value to the loss\\
directions $w_i\in\R^m$, $m<p$ & directions in a layer's activation space: the rows of a sparse autoencoder's decoder\\
$h=Wx$ & the activation vector of the layer, of width $m$\\
interference $\ip{w_i,w_j}x_j$ & polysemanticity: a neuron or direction picks up unrelated features\\
ReLU with bias $b_i<0$ & a nonlinearity that thresholds out small interference; the encoder of a sparse autoencoder\\
loss $\Loss(W,b)$ & what the layer must preserve for the rest of the network\\
\end{dictionary}

\paragraph{What one had to notice.} That a sparse feature is usually absent, so that interference with it costs only when it is present together with another one; that a ReLU with a slightly negative bias discards small positive interference for free and that negative interference costs nothing at all, because the ReLU clips it; and that in high dimension there are exponentially many almost orthogonal directions. The rest is computation.

\section{The solution: what the ReLU buys}\label{sec:superposition:two}

\paragraph{No superposition without a nonlinearity.} The linear model is principal component analysis, and so it never superposes. The theorem holds for every distribution with uncorrelated coordinates; we use only the mean $\mu\in\R^p$ and the diagonal covariance $\diag(v_1,\dots,v_p)$ of $x$.

\begin{theorem}[The linear model keeps the top $m$ features]\label{thm:superposition:linear}
Let $x$ have mean $\mu$ and covariance $\diag(v_1,\dots,v_p)$ with all $v_i>0$, let $m<p$, and let $c_{(1)}\ge\dots\ge c_{(p)}$ be the numbers $c_i=\omega_iv_i$ in decreasing order. For every $W\in\R^{m\times p}$ and $b\in\R^p$,
\[
  \Loss_{\mathrm{lin}}(W,b)\defeq\sum_{i=1}^p\omega_i\E\bigl(x_i-(W^\top Wx+b)_i\bigr)^2\ \ge\ \sum_{k=m+1}^pc_{(k)} .
\]
Equality holds when, for a set $K$ of $m$ indices carrying the $m$ largest values $c_i$, the columns $(w_i)_{i\in K}$ are orthonormal, $w_i=0$ for $i\notin K$, and $b=(I_p-W^\top W)\mu$. If $c_{(m)}>c_{(m+1)}$, then these are the only minimizers; in particular every minimizer drops exactly $p-m$ features, whatever the sparsity.
\end{theorem}

\begin{proof}
The plan is to remove the bias, to rewrite the loss as the distance from a diagonal matrix to a matrix of rank at most $m$, and to bound that distance by a linear program on the diagonal of a projection. Put $P=W^\top W$, $D=\diag(\omega_i)$ and $\Sigma=\diag(v_i)$.

\emph{The bias.} Split the error as $x-Px-b=(I_p-P)(x-\mu)+\bigl((I_p-P)\mu-b\bigr)$. The first part has mean zero, so the cross term vanishes in expectation, and
\[
  \Loss_{\mathrm{lin}}(W,b)=\sum_i\omega_i\bigl[(I_p-P)\Sigma(I_p-P)\bigr]_{ii}+\sum_i\omega_i\bigl((I_p-P)\mu-b\bigr)_i^2 .
\]
The second sum is nonnegative and vanishes exactly for $b=(I_p-P)\mu$. The first equals $\norm{A-B}_F^2$ with $A=D^{1/2}\Sigma^{1/2}=\diag(\sqrt{c_i})$ and $B=D^{1/2}P\Sigma^{1/2}$, a matrix of rank at most $m$.

\emph{A projection.} Because $\rank B\le m<p$, there is an orthogonal projection $Q$ of rank $m$ whose range contains the range of $B$, so that $(I_p-Q)B=0$. The matrices $(I_p-Q)A$ and $Q(A-B)$ have columns in orthogonal subspaces and sum to $A-B$, so
\begin{equation}\label{eq:superposition:proj}
  \norm{A-B}_F^2=\norm{(I_p-Q)A}_F^2+\norm{Q(A-B)}_F^2\ \ge\ \sum_ic_i\norm{(I_p-Q)e_i}^2=\sum_ic_i(1-q_i),
\end{equation}
with $q_i=\norm{Qe_i}^2=Q_{ii}$. The numbers $q_i$ lie in $[0,1]$ and add up to $\tr Q=m$.

\emph{The linear program.} For every real $t$, $\sum_ic_iq_i=tm+\sum_i(c_i-t)q_i$. Take $t=c_{(m)}$: the terms with $c_i<t$ are at most $0$, and those with $c_i\ge t$ are at most $c_i-t$, since $q_i\le1$; the indices with $c_i>t$ are among the top $m$. Hence $\sum_ic_iq_i\le tm+\sum_{k\le m}(c_{(k)}-t)=\sum_{k\le m}c_{(k)}$, and \eqref{eq:superposition:proj} gives the bound.

\emph{Attainment.} For the configuration in the statement, $P$ is the diagonal projection $\Pi_K$ onto the coordinates in $K$, so $B=\diag(\sqrt{c_i}\,\ind\{i\in K\})$ and $\norm{A-B}_F^2=\sum_{i\notin K}c_i=\sum_{k>m}c_{(k)}$; the bias term vanishes by the choice of $b$.

\emph{Uniqueness.} Let $c_{(m)}>c_{(m+1)}$, so that $K$ is unique, and take $t$ strictly between the two values. Every term $(c_i-t)q_i$ is then at most $c_i-t$ for $i\in K$ and at most $0$ for $i\notin K$, with equality only if $q_i=1$, respectively $q_i=0$. So equality forces $\norm{(I_p-Q)e_i}=0$ for $i\in K$, and since $\rank Q=m=\abs K$, $Q=\Pi_K$. Equality in \eqref{eq:superposition:proj} also forces $Q(A-B)=0$, that is $B=\Pi_KA$, and then $P=D^{-1/2}B\Sigma^{-1/2}=\Pi_K$. The identity $W^\top W=\Pi_K$ says $\ip{w_i,w_j}=\ind\{i=j\in K\}$, and the bias must be $(I_p-\Pi_K)\mu$.
\end{proof}

Sparsity enters the linear model only through the variances, so it cannot make superposition pay. The reason is visible in the proof: a linear decoder pays for interference in every input, whether the interfering feature is present or not, through the variance it adds. When importances tie the minimizer is not unique, and rotations of $W$ within a tied block may represent more than $m$ features; but $W^\top W$ is still an orthogonal projection, one feature per dimension, and the slightest difference in importances removes the rotation (\cref{ex:superposition:ties}).\sidenote{That a linear autoencoder trained by least squares recovers the principal subspace goes back to \textcite{BourlardKamp1988} and \textcite{BaldiHornik1989}, who also showed that its other critical points are saddles. The low-rank approximation behind it is that of \textcite{EckartYoung1936}.}

\paragraph{Two facts about the ReLU.} Everything that follows uses two inequalities: for all real $z,z'$ and every $y\ge0$,
\begin{equation}\label{eq:superposition:relu}
  \abs{\relu(z)-\relu(z')}\le\abs{z-z'},\qquad\abs{y-\relu(z)}\le\abs{y-z}.
\end{equation}
The first says that the ReLU is $1$-Lipschitz. For the second, if $z\ge0$ both sides agree, and if $z<0$, then $\abs{y-0}=y\le y-z$. The ReLU is the projection onto $[0,\infty)$, and projecting onto a convex set that contains the target never increases the distance to it. This is the whole advantage of the nonlinearity: negative interference is free.

\begin{lemma}[The cost of dropping a feature]\label[lemma]{lem:superposition:drop}
If $w_i=0$, then the term $i$ of $\Loss(W,b)$ is at least $\omega_i\,\supv$, with equality if and only if $b_i=\frac{1-S}2$.
\end{lemma}

\begin{proof}
If $w_i=0$, then $x'_i=\relu(b_i)$ is a constant $c\ge0$, and $\E(x_i-c)^2=\supv+(\E x_i-c)^2$ is smallest exactly at $c=\E x_i=\frac{1-S}2>0$, that is, at $b_i=\frac{1-S}2$.
\end{proof}

A dropped feature is replaced by its mean, and the bias makes this possible; without a bias the cost of dropping would be $\omega_i\E x_i^2$, which changes the comparisons below (\cref{ex:superposition:nobias}).

\paragraph{Two features in one dimension.} Take $p=2$, $m=1$, so that $W=(w_1,w_2)$ is a row of two numbers, and importances $\omega=(1,r)$ with $0<r\le1$. The loss is a function of four numbers, and we compare three explicit configurations: the \term{dedicated} configuration $w=(1,0)$, $b=(0,\frac{1-S}2)$, which keeps feature $1$ exactly and replaces feature $2$ by its mean; the \term{antipodal pair} $w=(1,-1)$, $b=(0,0)$; and the antipodal tilt $w=(1,-c)$ of the dedicated configuration. The tilt needs one more function. For $\tau\in\R$ and a feature $x$, let
\begin{equation}\label{eq:superposition:f}
  f_S(\tau)=S\,\E\relu(\tau-x)^2+(1-S)\,\E(\tau-x)^2,\qquad \kappa(S)=\min_{\tau\in\R}f_S(\tau),
\end{equation}
and let $\tau_S$ be the minimizer (\cref{ex:superposition:kappa} shows that it exists, is unique, lies in $(0,1)$, and solves a quadratic equation). Numerically, $\kappa(S)/\supv$ equals $1$ at $S=0$, $0.60$ at $S=\frac12$, and $0.11$ at $S=0.9$, and $\kappa(S)\le f_S(0)=\frac{(1-S)^2}3$.

\begin{theorem}[Two features in one dimension]\label{thm:superposition:two}
Let $p=2$, $m=1$, $\omega=(1,r)$ with $0<r\le1$, and $0\le S<1$.
\begin{enumerate}
\item Every configuration that drops a feature has $\Loss\ge r\,\supv$. The dedicated configuration attains $r\,\supv$.
\item The antipodal pair has $\Loss=\frac{1+r}6(1-S)^2$.
\item For every $c>0$, the configuration $w=(1,-c)$, $b=\bigl(c\tau_S,\,(1+c-c^2)\frac{1-S}2\bigr)$ has
\[
  \Loss\le r\,\supv-c^2\bigl(r\,\supv-\kappa(S)\bigr)+c^4r\,\supv .
\]
\item Consequently, if $r\,\supv>\min\bigl\{\kappa(S),\frac{1+r}6(1-S)^2\bigr\}$, then the infimum of $\Loss$ is strictly smaller than the loss of every configuration that drops a feature: no optimal or nearly optimal configuration drops one. This holds in particular whenever $S>\frac{2+r}{2+5r}$.
\end{enumerate}
\end{theorem}

\begin{proof}
The plan is to evaluate the antipodal pair exactly and to bound the tilt with \eqref{eq:superposition:relu}; part (iv) then compares numbers.

\emph{Part (i).} If $w_2=0$, then \cref{lem:superposition:drop} bounds the second term by $r\,\supv$, and the first is nonnegative. If $w_1=0$, then the first term alone is at least $\supv\ge r\,\supv$. For the dedicated configuration, $x'_1=\relu(x_1)=x_1$, and the second term equals $r\,\supv$ by \cref{lem:superposition:drop}.

\emph{Part (ii).} Here $x'_1=\relu(x_1-x_2)$. If $x_2=0$, then $x'_1=x_1$. If $x_1=0$, then $x'_1=\relu(-x_2)=0=x_1$: the interference is negative and the ReLU clips it. If both are active, which happens with probability $(1-S)^2$, then $x_1-\relu(x_1-x_2)=\min(x_1,x_2)$, and for independent uniform $U,V$ the minimum has density $2(1-t)$ on $[0,1]$, so $\E\min(U,V)^2=\int_0^12t^2(1-t)\,\dd t=\frac16$. The first term is therefore $\frac{(1-S)^2}6$, and by symmetry the second is $r$ times the same.

\emph{Part (iii), feature 1.} Now $x'_1=\relu(x_1+c(\tau_S-x_2))$. If $x_1>0$, then $x_1=\relu(x_1)$ and the Lipschitz bound in \eqref{eq:superposition:relu} gives $\abs{x_1-x'_1}\le c\abs{\tau_S-x_2}$. If $x_1=0$, then $(x_1-x'_1)^2=c^2\relu(\tau_S-x_2)^2$, because $c>0$. Since $x_1$ and $x_2$ are independent, the first term is at most $c^2f_S(\tau_S)=c^2\kappa(S)$.

\emph{Part (iii), feature 2.} Now $x'_2=\relu(z)$ with $z=-cx_1+c^2x_2+b_2$. By the second inequality in \eqref{eq:superposition:relu}, $(x_2-x'_2)^2\le(x_2-z)^2=\bigl((1-c^2)x_2+cx_1-b_2\bigr)^2$. Our $b_2$ is the mean of $(1-c^2)x_2+cx_1$, so the expectation is its variance, $(1-c^2)^2\supv+c^2\supv=\supv(1-c^2+c^4)$. Adding $r$ times this to the bound for feature $1$ gives (iii).

\emph{Part (iv).} If $r\,\supv>\kappa(S)$, take $c^2=(r\supv-\kappa(S))/(2r\supv)$ in (iii); the loss is then at most
\[
  r\,\supv-\frac{\bigl(r\,\supv-\kappa(S)\bigr)^2}{4r\,\supv}<r\,\supv .
\]
If $r\,\supv>\frac{1+r}6(1-S)^2$, then (ii) gives a loss below $r\,\supv$. In both cases (i) shows that every configuration dropping a feature is worse than one of these. Finally, by \eqref{eq:superposition:moments} and because $1-S>0$, the inequality
\[
  \frac{1+r}6(1-S)^2<r\,\frac{(1-S)(1+3S)}{12}
\]
is equivalent to $2(1+r)(1-S)<r(1+3S)$, that is, to $(2+5r)S>2+r$.
\end{proof}

Both constructions superpose by making the interference negative. The antipodal pair is the extreme case: each feature is read out exactly whenever it is alone, and the error $\min(x_1,x_2)$ appears only when both are present, an event of probability $(1-S)^2$, while dropping a feature costs order $1-S$. The tilt is the infinitesimal version. Feature $2$ gains $c^2r\supv$ because its readout now sees a little of $x_2$; feature $1$ pays $c^2\kappa(S)$, and since negative interference is clipped when $x_1=0$, $\kappa(S)$ is smaller than the variance $\supv$ that a linear decoder would pay, which is what makes the trade profitable. At $S=0$ the two coincide, $\kappa(0)=\frac1{12}=v(0)$, and at $r=1$ the tilt is then neutral to second order in $c$, the rotational degeneracy of \cref{thm:superposition:linear}; at higher order the ReLU, which clips the rare negative outputs of a rotated readout, breaks the tie in favor of superposition.

Part (iii) also shows that for $r>\kappa(S)/\supv$ the dedicated configuration is not a local minimum: the loss decreases along the tilt. The converse is what a second-order computation gives.

\begin{heuristic}[Heuristic derivation of the spinodal]
Perturb the dedicated configuration to $w=(1+\delta,\pm c)$ and $b=(b_1,\frac{1-S}2+\beta)$ with all four perturbations small. For feature $2$, whose ReLU stays in its linear region, the best readout of $x_2$ from $\pm c(1+\delta)x_1+c^2x_2$ plus a constant gains at most $c^2r\supv+O(c^4)$. For feature $1$, the error is $-(2\delta x_1\pm cx_2+b_1)$ when $x_1>0$, up to a set of measure $O(c)$ near $x_1=0$, and $\relu(\pm cx_2+b_1)$ when $x_1=0$; the term in $\delta$ only adds a positive quadratic, the sign $+$ is worse than $-$ (it is not clipped), and with $b_1=c\tau$ the cost is $c^2f_S(\tau)\ge c^2\kappa(S)$. So the second variation is $c^2(\kappa(S)-r\supv)$ to leading order: the dedicated configuration is a strict local minimum for $r<\kappa(S)/\supv$ and a saddle for $r>\kappa(S)/\supv$.
\end{heuristic}

We checked the heuristic numerically: from small perturbations of the dedicated configuration with a negative tilt $w_2<0$, gradient descent on the exact loss returns to it at $(S,r)=(0.7,0.3)$, $(0.5,0.55)$ and $(0.3,0.75)$, just below the curve $r=\kappa(S)/\supv$, and leaves it at $(0.7,0.4)$, $(0.5,0.65)$ and $(0.3,0.85)$, just above. In the language of phase transitions, the curve $r=\kappa(S)/\supv$ is a \term{spinodal}\index{spinodal}, the edge of metastability of the dedicated phase. \Cref{sec:superposition:reading} shows that the true transition comes before it and is discontinuous.

\section{The solution: geometry and capacity}\label{sec:superposition:geometry}

\paragraph{The sparse limit.} With many features, the cleanest statements are about $S\to1$. Condition on the set $A\subseteq\{1,\dots,p\}$ of active features, which has probability $(1-S)^{\abs A}S^{p-\abs A}$, and let $\ell_A(W,b)$ be the conditional expected loss. For $A=\emptyset$ and $A=\{i\}$, with $U$ uniform on $[0,1]$,
\begin{align*}
  \ell_\emptyset&=\sum_j\omega_j\relu(b_j)^2,\\
  \ell_{\{i\}}&=\omega_i\E\bigl(U-\relu(\norm{w_i}^2U+b_i)\bigr)^2+\sum_{j\ne i}\omega_j\E\relu\bigl(\ip{w_j,w_i}U+b_j\bigr)^2 .
\end{align*}
Since $x\in[0,1]^p$, every $\ell_A$ is bounded by a constant $C_{W,b}$ that depends on the configuration, and the probability that two or more features are active is at most $\binom p2(1-S)^2$. Hence
\begin{equation}\label{eq:superposition:expansion}
  \Loss(W,b)=S^p\ell_\emptyset+(1-S)S^{p-1}\sum_{i=1}^p\ell_{\{i\}}+R,\qquad 0\le R\le\tbinom p2C_{W,b}(1-S)^2 .
\end{equation}
To first order in the density $1-S$, the loss sees only inputs with one active feature: the problem becomes one of \emph{one-hot} inputs.

\begin{definition}[Obtuse configurations]\label[definition]{def:superposition:obtuse}
A configuration is \term{obtuse}\index{obtuse configuration} if its directions are unit vectors with pairwise nonpositive inner products, $\norm{w_i}=1$ and $\ip{w_i,w_j}\le0$ for $i\ne j$, and its bias is $b=0$.
\end{definition}

\begin{lemma}[No first-order interference]\label[lemma]{lem:superposition:firstorder}
A configuration has $\ell_\emptyset=0$ and $\ell_{\{i\}}=0$ for every $i$ if and only if it is obtuse.
\end{lemma}

\begin{proof}
The term $\omega_i\E(U-\relu(\norm{w_i}^2U+b_i))^2$ vanishes if and only if $\relu(\norm{w_i}^2u+b_i)=u$ for almost every $u\in[0,1]$; the left side is then a linear function equal to $u$ on an interval, so $\norm{w_i}^2=1$ and $b_i=0$. Given $b=0$, the term $\E\relu(\ip{w_j,w_i}U)^2$ vanishes if and only if $\ip{w_j,w_i}\le0$, and $\ell_\emptyset=0$ holds. Conversely, an obtuse configuration makes every one of these terms zero.
\end{proof}

How many directions can be pairwise obtuse? The answer is classical.

\begin{lemma}[At most $2m$ obtuse directions]\label[lemma]{lem:superposition:2m}
If $u_1,\dots,u_k\in\R^m$ are nonzero and $\ip{u_i,u_j}\le0$ for all $i\ne j$, then $k\le2m$. The $2m$ vectors $\pm e_1,\dots,\pm e_m$ attain the bound.
\end{lemma}

\begin{proof}
We use induction on $m$. For $m=1$ the vectors are nonzero numbers, and two of the same sign have a positive product, so $k\le2$. Let $m\ge2$, and project $u_1,\dots,u_{k-1}$ onto the hyperplane orthogonal to $u_k$: $u'_i=u_i-\frac{\ip{u_i,u_k}}{\norm{u_k}^2}u_k$. For $i\ne j$,
\[
  \ip{u'_i,u'_j}=\ip{u_i,u_j}-\frac{\ip{u_i,u_k}\ip{u_j,u_k}}{\norm{u_k}^2}\le0,
\]
because the subtracted product of two nonpositive numbers is nonnegative. A projection $u'_i$ vanishes only if $u_i$ is a multiple of $u_k$, necessarily a negative one; two such $u_i$ would have a positive inner product, so at most one $u'_i$ vanishes. The remaining $k-2$ or more nonzero vectors lie in an $(m-1)$-dimensional space with pairwise nonpositive inner products, so $k-2\le2(m-1)$ by the induction hypothesis.
\end{proof}

\begin{theorem}[Sparse features are superposed]\label{thm:superposition:sparse}
Let $m<p$ and $0\le S<1$, and let $\omega_{\min}$ be the smallest importance.
\begin{enumerate}
\item A configuration $(W,b)$ satisfies $\Loss(W,b)=O\bigl((1-S)^2\bigr)$ as $S\to1$ if and only if it is obtuse. If it is not, then $\Loss(W,b)\ge2^{1-p}(1-S)\bigl(\ell_\emptyset+\sum_i\ell_{\{i\}}\bigr)>0$ for $S\ge\frac12$.
\item Obtuse configurations exist if and only if $p\le2m$. Every obtuse configuration has
\[
  \Loss(W,0)\le\frac{\supOmega}3(1-S)\bigl(1-S^{p-1}\bigr)\le\frac{(p-1)\supOmega}3(1-S)^2 .
\]
\item If $p\le2m$ and $4(1-S^{p-1})\supOmega<\omega_{\min}(1+3S)$, then the loss of every configuration that drops a feature exceeds the loss of every obtuse configuration. In particular, for $S$ close enough to $1$, no optimal or nearly optimal configuration drops a feature, and all are in superposition.
\end{enumerate}
\end{theorem}

\begin{proof}
The plan is to read part (i) from the expansion \eqref{eq:superposition:expansion}, part (ii) from \cref{lem:superposition:2m} and the clipping of negative interference, and part (iii) from \cref{lem:superposition:drop}.

\emph{Part (i).} By \eqref{eq:superposition:expansion}, $\Loss(W,b)=O((1-S)^2)$ if and only if $S^p\ell_\emptyset+(1-S)S^{p-1}\sum_i\ell_{\{i\}}=O((1-S)^2)$, which forces $\ell_\emptyset=0$ (the coefficient $S^p$ tends to $1$) and then $\sum_i\ell_{\{i\}}=0$; by \cref{lem:superposition:firstorder} this means obtuse. Conversely, for an obtuse configuration the first two terms vanish and $\Loss=R=O((1-S)^2)$. For the lower bound, drop $R\ge0$ and use $S^p\ge(1-S)S^{p-1}$ and $S^{p-1}\ge2^{1-p}$, both valid for $S\ge\frac12$.

\emph{Part (ii), existence.} By \cref{lem:superposition:2m}, $p$ unit vectors with pairwise nonpositive inner products exist only if $p\le2m$. If $p\le2m$, then $\lceil p/2\rceil\le m$, and the directions $w_{2k-1}=e_k$, $w_{2k}=-e_k$ for $k\le\lfloor p/2\rfloor$, together with $w_p=e_{(p+1)/2}$ when $p$ is odd, are obtuse.

\emph{Part (ii), the bound.} Fix $i$ and an obtuse configuration. If $x_i=0$, then $x'_i=\relu(\sum_{j\ne i}\ip{w_i,w_j}x_j)=0=x_i$, since every term of the sum is at most $0$. If $x_i>0$, then $x'_i=\relu(x_i+\sum_{j\ne i}\ip{w_i,w_j}x_j)$ lies in $[0,x_i]$, so $\abs{x_i-x'_i}\le x_i$, and $x'_i=x_i$ when no other feature is active. By independence, $\E(x_i-x'_i)^2\le\E x_i^2\cdot\Pr\{\text{some }x_j\ne0,\ j\ne i\}=\frac{1-S}3(1-S^{p-1})$. Summing with weights $\omega_i$ gives the first bound, and $1-S^{p-1}\le(p-1)(1-S)$ gives the second.

\emph{Part (iii).} By \cref{lem:superposition:drop} and \eqref{eq:superposition:moments}, a configuration that drops a feature has loss at least $\omega_{\min}\frac{(1-S)(1+3S)}{12}$, and this exceeds the bound of part (ii) exactly under the stated inequality. As $S\to1$ its left side tends to $0$ and its right side to $4\omega_{\min}$. A configuration that drops no feature represents $p>m$ of them.
\end{proof}

For $p=2$ and $m=1$ the obtuse configuration is the antipodal pair, and part (ii) overestimates its exact loss of \cref{thm:superposition:two} by a factor of two. For $p\le2m$ in general, sparse features pair up into antipodal pairs in orthogonal planes, the cross-polytope, or form other obtuse sets such as a regular simplex of $m+1$ directions with pairwise inner products $-\frac1m$; at first order all of them are free. With $p>2m$ features first-order interference cannot be avoided. Then some inner products are positive, and a negative bias must cut off the interference, at the price of also cutting off the bottom of every true signal; part (i) says that the loss of every configuration is then of order $1-S$. Its constant is positive configuration by configuration; that it is bounded below uniformly is the content of the next remark.

\begin{remark}[More than $2m$ features]\label[remark]{rem:superposition:uniform}
For $p>2m$ the infimum of $\ell_\emptyset+\sum_i\ell_{\{i\}}$ over all configurations is positive, so the optimal loss is at least a constant times $1-S$ for $S\ge\frac12$.
\end{remark}

\begin{proofsketch}
Along a sequence of configurations on which this sum tends to $0$, $\relu(b_j)\to0$, and the functions $u\mapsto\relu(\norm{w_i}^2u+b_i)$ tend to $u\mapsto u$ in $L^2[0,1]$; one checks that this forces $\norm{w_i}\to1$ and $b_i\to0$, so a subsequence of $W$ converges, and its limit is obtuse, which \cref{lem:superposition:2m} forbids. We omit the routine verification of the second step; the lower bound then follows from the inequality of part (i) of \cref{thm:superposition:sparse}.
\end{proofsketch}

\paragraph{The shape of the solutions.} In the configurations that training finds (\cref{fig:superposition:pentagon}), the represented directions form regular polygons and polytopes, and the share of a dimension that a feature receives takes simple fractional values. \Textcite{ElhageEtAl2022} define the \term{feature dimensionality}\index{feature dimensionality} of a represented feature,
\[
  \supD_i=\frac{\norm{w_i}^4}{\sum_{j}\ip{w_i,w_j}^2},
\]
which is $1$ for a feature with a dimension of its own and smaller under interference; the sum runs over all $j$, including $j=i$.

\begin{proposition}[Tight frames share dimensions equally]\label[proposition]{prop:superposition:frame}
Let $w_1,\dots,w_N$ have equal norms and span a subspace $V$ of dimension $k$, and suppose that they form a \term{tight frame} of $V$: $\sum_jw_jw_j^\top=\gamma\Pi_V$ for a number $\gamma$, with $\Pi_V$ the orthogonal projection onto $V$. Then $\supD_i=k/N$ for every $i$.
\end{proposition}

\begin{proof}
Taking traces, $N\norm{w_1}^2=\gamma k$. For each $i$, $w_i\in V$ gives $\sum_j\ip{w_i,w_j}^2=w_i^\top\bigl(\sum_jw_jw_j^\top\bigr)w_i=\gamma\norm{w_i}^2=\frac Nk\norm{w_i}^4$, so $\supD_i=k/N$.
\end{proof}

An antipodal pair is a tight frame of a line, so each of its features has dimensionality $\frac12$; the vertices of a regular $N$-gon with $N\ge3$ are a tight frame of the plane (\cref{ex:superposition:frames}), with dimensionality $\frac2N$: $\frac23$ for a triangle, $\frac25$ for a pentagon; a regular tetrahedron in $\R^3$ gives $\frac34$. When the represented features split into tight frames of mutually orthogonal subspaces, each frame keeps its own value.

\paragraph{Exponentially many almost orthogonal directions.} Obtuse configurations stop at $2m$ features. Beyond that, interference is unavoidable but can be made small, and in high dimension small costs little room.

\begin{theorem}[Packing almost orthogonal directions]\label{thm:superposition:packing}
For every $\eps\in(0,1)$ and every integer $m\ge1$, there are $N=\lfloor e^{\eps^2m/4}\rfloor$ unit vectors $u_1,\dots,u_N\in\R^m$ with $\abs{\ip{u_i,u_j}}\le\eps$ for all $i\ne j$.
\end{theorem}

\begin{proof}
The plan is to choose the vectors at random and to show that the expected number of bad pairs is less than one. Let $u_i=s_i/\sqrt m$ with independent uniformly random sign vectors $s_i\in\{\pm1\}^m$; each $u_i$ is a unit vector. For $i\ne j$, $m\ip{u_i,u_j}=\sum_{k=1}^ms_{ik}s_{jk}$ is a sum of $m$ independent uniform signs, because the products $s_{ik}s_{jk}$ are independent and uniform on $\{\pm1\}$. Hoeffding's inequality (\cref{thm:toolkit:hoeffding}) for $m$ variables with values in $[-1,1]$ gives
\[
  \Pr\bigl\{\abs{\ip{u_i,u_j}}>\eps\bigr\}\le\Pr\Bigl\{\Bigl\lvert\sum_{k=1}^ms_{ik}s_{jk}\Bigr\rvert\ge\eps m\Bigr\}\le2\exp\Bigl(-\frac{2(\eps m)^2}{4m}\Bigr)=2e^{-\eps^2m/2}.
\]
Summing over the $\binom N2$ pairs, the expected number of bad pairs is at most $N(N-1)e^{-\eps^2m/2}<N^2e^{-\eps^2m/2}\le1$, where the last step uses $N\le e^{\eps^2m/4}$. A nonnegative integer random variable with mean less than $1$ is $0$ with positive probability, so some choice of signs has no bad pair.
\end{proof}

The theorem is a form of the lemma of \textcite{JohnsonLindenstrauss1984}, and its exponent is the right order up to a logarithm.\sidenote{The Welch bound, \cref{ex:superposition:welch}, shows that $N>m$ unit vectors always have some $\abs{\ip{u_i,u_j}}\ge\sqrt{(N-m)/(m(N-1))}$ \parencite{Welch1974}, so $\eps$ cannot go below about $(2m)^{-1/2}$ once $N\ge2m$. That $\log N$ can be at most of order $\eps^2m\log(1/\eps)$ is a theorem of \textcite{Alon2003}.} For $m=1000$ and $\eps=0.25$ it guarantees $6.1$ million directions with all cosines at most $0.25$ in absolute value; in two samples of $3000$ random sign vectors the largest cosines were $0.16$ and $0.18$.

\paragraph{Reading sparse features through a ReLU.} Almost orthogonal directions are useful only if the features can be read back. For the toy model with unit directions this is one ReLU with a negative bias.

\begin{theorem}[Thresholded readout]\label{thm:superposition:readout}
Let the directions $w_1,\dots,w_p\in\R^m$ be unit vectors with \term{coherence}\index{coherence} $\mu=\max_{i\ne j}\abs{\ip{w_i,w_j}}$, let $s\ge1$ be an integer, let $\beta\ge s\mu$, and set $b=-\beta\mathbf 1_p$ in \eqref{eq:superposition:model}. If $x\in[0,1]^p$ has at most $s$ nonzero coordinates, then for every $i$:
\begin{enumerate}
\item if $x_i=0$, then $x'_i=0$;
\item if $x_i>0$, then $x_i-\beta-(s-1)\mu\le x'_i\le x_i$; with $\beta=s\mu$, the error is at most $(2s-1)\mu$.
\end{enumerate}
\end{theorem}

\begin{proof}
Write $x'_i=\relu(x_i+\iota_i-\beta)$ with the interference $\iota_i=\sum_{j\ne i}\ip{w_i,w_j}x_j$, which has at most as many nonzero terms as there are active features $j\ne i$, each of absolute value at most $\mu$. If $x_i=0$, there are at most $s$ such terms, so $\iota_i-\beta\le s\mu-\beta\le0$ and $x'_i=0$. If $x_i>0$, there are at most $s-1$, so $\abs{\iota_i}\le(s-1)\mu$. The ReLU is nondecreasing, hence $x'_i\le\relu(x_i+(s-1)\mu-\beta)\le\relu(x_i-\mu)\le x_i$, using $\beta\ge s\mu$ and $x_i\ge0$; and $x'_i\ge x_i+\iota_i-\beta\ge x_i-\beta-(s-1)\mu$.
\end{proof}

\begin{corollary}[Capacity of the toy model]\label[corollary]{cor:superposition:capacity}
Let $p\ge2$, $s\ge1$ and $\delta\in(0,1)$. If $m\ge4(2s-1)^2\log p/\delta^2$, then there is a configuration of \eqref{eq:superposition:model} that outputs $x'_i=0$ for every inactive feature and $\abs{x'_i-x_i}\le\delta$ for every active one, for every $x\in[0,1]^p$ with at most $s$ active features.
\end{corollary}

\begin{proof}
Let $\eps=\delta/(2s-1)\in(0,1)$. The hypothesis gives $\eps^2m/4\ge\log p$, so \cref{thm:superposition:packing} provides at least $p$ unit vectors with coherence at most $\eps$; take $p$ of them as columns of $W$. Apply \cref{thm:superposition:readout} with $\mu\le\eps$ and $\beta=s\eps\ge s\mu$: the error on active features is at most $\beta+(s-1)\mu\le(2s-1)\eps=\delta$.
\end{proof}

The number of features that fit in $m$ dimensions is therefore exponential in $m\delta^2/s^2$. This is compressed sensing done by a single layer.\sidenote{Compressed sensing recovers every $s$-sparse vector from $m$ of order $s\log(p/s)$ random linear measurements by $\ell_1$ minimization \parencite{CandesEtAl2006,Donoho2006}. A one-step thresholded readout pays a factor of about $s$ in the number of dimensions, and the Welch bound shows that no guarantee based on coherence alone can do better (\cref{ex:superposition:welch}).} The guarantee is worst case: it holds for every input with at most $s$ active features, while for random sparse features the interference is a sum of terms with random signs and is typically of order $\mu\sqrt s$ rather than $\mu s$.

\section{Reading the solution}\label{sec:superposition:reading}

The theorems make four kinds of prediction. (i) A linear network never superposes, at every sparsity. (ii) Two features share one dimension when they are sparse and nearly equally important; superposition provably wins for $S>\frac{2+r}{2+5r}$ and for $r>\kappa(S)/\supv$, and the dedicated configuration is a saddle in the second region. (iii) Features fit in $m$ dimensions without first-order interference exactly up to $p=2m$, in obtuse configurations with zero bias; beyond $2m$, positive interference appears and must be cut off by a negative bias, and the loss jumps from order $(1-S)^2$ to order $1-S$. (iv) With few active features, the number of storable features is exponential in the dimension.

\begin{figure}[t]
\begin{fullwidth}
\centering
\begin{tikzpicture}
\begin{groupplot}[group style={group size=2 by 1,horizontal sep=17mm},hplot,width=0.475\linewidth,height=0.37\linewidth]
\nextgroupplot[ymode=log,xmin=0,xmax=1,ymin=1e-4,ymax=0.5,xlabel={sparsity $S$},ylabel={loss},
  legend style={at={(0.02,0.02)},anchor=south west}]
\addplot[theory] table[x=S,y=dedicated] {figures/data/11-superposition-loss.dat};
\addplot[theory2] table[x=S,y=antipodal] {figures/data/11-superposition-loss.dat};
\addplot[theory3] table[x=S,y=optimum] {figures/data/11-superposition-loss.dat};
\addplot[simulation] table[x=S,y=seed0] {figures/data/11-superposition-loss.dat};
\addplot[simulation,forget plot] table[x=S,y=seed1] {figures/data/11-superposition-loss.dat};
\addplot[simulation,forget plot] table[x=S,y=seed2] {figures/data/11-superposition-loss.dat};
\legend{dedicated $r\,v(S)$,antipodal $\frac{1+r}6(1-S)^2$,numerical optimum,trained (three seeds)}
\nextgroupplot[xmin=0,xmax=1,ymin=0,ymax=1,xlabel={importance ratio $r$},ylabel={sparsity $S$},
  legend style={at={(0.02,0.02)},anchor=south west}]
\addplot[theory] table[x=r,y=kappa] {figures/data/11-superposition-phasecurves.dat};
\addplot[theory2] table[x=r,y=antipodal] {figures/data/11-superposition-phasecurves.dat};
\addplot[simulation2] table[x=r,y=numeric] {figures/data/11-superposition-phase.dat};
\node[font=\footnotesize,text=ink] at (axis cs:0.72,0.86) {superposition};
\node[font=\footnotesize,text=ink] at (axis cs:0.4,0.4) {feature 2 dropped};
\legend{$r=\kappa(S)/v(S)$,$S=\frac{2+r}{2+5r}$,numerical boundary}
\end{groupplot}
\end{tikzpicture}
\caption{Two features in one dimension. Left: importance ratio $r=\frac12$. The curves are the exact losses of the dedicated configuration and of the antipodal pair (\cref{thm:superposition:two}) and the minimum of the exact loss over $(w_1,w_2,b_1,b_2)$, computed by L-BFGS from three starting points. The points are the exact loss of parameters trained by Adam on fresh batches of $2048$ samples ($1500$ steps, learning rate $10^{-2}$, then $2\times10^{-3}$ after step $1000$; initialization $\cN(0,\frac14)$; seeds $0$, $1$, $2$). Right: in the region above the dots the numerical optimum represents both features, and below them it drops feature $2$; superposition is proved above each of the two curves.}\label{fig:superposition:two}
\end{fullwidth}
\end{figure}

\Cref{fig:superposition:two} tests prediction (ii). On the left, at $r=\frac12$, the optimal loss follows the dedicated curve $r\supv$ exactly up to $S\approx0.52$, then breaks away, and above $S\approx0.7$ it is within $4\%$ of the antipodal curve, which it approaches as $S\to1$. Training finds the optimum to within $0.1\%$ whenever it does not stall, and the stalls are informative. At $S=0.5$ one seed ends in a superposed configuration $3\%$ worse than the dedicated optimum: below the transition the superposed branch is metastable. At $S=0.55$ two of the three seeds end at the dedicated configuration, which there is a strict local minimum: $0.55$ lies between the true transition, $S_c=0.516$, and the spinodal, $S=0.586$, at which $\kappa(S)/\supv=\frac12$. For $S\ge0.75$ one seed still sits near the dedicated configuration, now a saddle, but a flat one: the negative curvature $2(r\supv-\kappa(S))$ is about $0.02$ at $S=0.9$, and the escape from a saddle with a small starting tilt takes a time of order $\log(1/c_0)$ divided by the curvature, the logistic escape met in \cref{ch:deeplinear,ch:singleindex}.

On the right, the numerical boundary $S_c(r)$ lies below both proven curves for every $r$. The transition is first order: along the left panel, the optimal configuration jumps at $S_c=0.516$ from $w=(1,0)$ to $w\approx(1.04,-0.77)$, with $\norm{W}_F^2$ jumping from $1$ to $1.67$, and then moves toward the antipodal pair, reaching $w\approx(1.04,-0.91)$ at $S=0.6$; it does not grow continuously from the dedicated one. Near $r=1$ the boundary falls to $S=0$ and meets the spinodal, and at $r=1$ every $S\ge0$ superposes, by a small rotation when $S$ is small (at $S=0$ the optimal loss is $0.08326$, against $\frac1{12}\approx0.08333$ for the dedicated configuration). For $r\to0$ both proven curves and the numerical boundary tend to $S=1$; a feature much less important than the other is superposed only when it is very sparse. The proven curve $S=\frac{2+r}{2+5r}$ is close to the true boundary at small $r$, because the transition goes to a nearly antipodal configuration with nearly zero bias, and the spinodal is close to it at large $r$.

\begin{figure}[t]
\begin{fullwidth}
\centering
\begin{tikzpicture}
\begin{groupplot}[group style={group size=4 by 1,horizontal sep=4mm},hplot,width=0.29\linewidth,height=0.29\linewidth,
  xmin=-1.45,xmax=1.45,ymin=-1.45,ymax=1.45,axis equal image,xtick=\empty,ytick=\empty,axis line style={muted!50}]
\nextgroupplot[title={$S=0$}]
\addplot[quiver={u=\thisrow{x},v=\thisrow{y}},-stealth,sea,line width=0.9pt,restrict expr to domain={\thisrow{panel}}{0:0}] table[x expr=0,y expr=0] {figures/data/11-superposition-pentagon.dat};
\addplot[only marks,mark=none,nodes near coords,point meta=explicit symbolic,every node near coord/.style={font=\scriptsize,text=ink,anchor=center},restrict expr to domain={\thisrow{panel}}{0:0}] table[x expr=1.22*\thisrow{x},y expr=1.22*\thisrow{y},meta=feature] {figures/data/11-superposition-pentagon.dat};
\nextgroupplot[title={$S=0.5$}]
\addplot[quiver={u=\thisrow{x},v=\thisrow{y}},-stealth,sea,line width=0.9pt,restrict expr to domain={\thisrow{panel}}{1:1}] table[x expr=0,y expr=0] {figures/data/11-superposition-pentagon.dat};
\addplot[only marks,mark=none,nodes near coords,point meta=explicit symbolic,every node near coord/.style={font=\scriptsize,text=ink,anchor=center},restrict expr to domain={\thisrow{panel}}{1:1}] table[x expr=1.22*\thisrow{x},y expr=1.22*\thisrow{y},meta=feature] {figures/data/11-superposition-pentagon.dat};
\nextgroupplot[title={$S=0.8$}]
\addplot[quiver={u=\thisrow{x},v=\thisrow{y}},-stealth,sea,line width=0.9pt,restrict expr to domain={\thisrow{panel}}{2:2}] table[x expr=0,y expr=0] {figures/data/11-superposition-pentagon.dat};
\addplot[only marks,mark=none,nodes near coords,point meta=explicit symbolic,every node near coord/.style={font=\scriptsize,text=ink,anchor=center},restrict expr to domain={\thisrow{panel}}{2:2}] table[x expr=1.22*\thisrow{x},y expr=1.22*\thisrow{y},meta=feature] {figures/data/11-superposition-pentagon.dat};
\nextgroupplot[title={$S=0.95$}]
\addplot[quiver={u=\thisrow{x},v=\thisrow{y}},-stealth,sea,line width=0.9pt,restrict expr to domain={\thisrow{panel}}{3:3}] table[x expr=0,y expr=0] {figures/data/11-superposition-pentagon.dat};
\addplot[only marks,mark=none,nodes near coords,point meta=explicit symbolic,every node near coord/.style={font=\scriptsize,text=ink,anchor=center},restrict expr to domain={\thisrow{panel}}{3:3}] table[x expr=1.15*\thisrow{x},y expr=1.15*\thisrow{y},meta=feature] {figures/data/11-superposition-pentagon.dat};
\end{groupplot}
\end{tikzpicture}
\caption{Five features in two dimensions: the directions $w_i$ of norm above $0.1$ in trained models (importances $\omega_i=0.8^{i-1}$; Adam, $3000$ steps on fresh batches of $4096$ samples, learning rate $10^{-2}$, then $2\times10^{-3}$ after step $2000$; initialization $\cN(0,0.09)$; the better of seeds $0$ and $1$). Labels are feature indices. Dense: the two most important features, orthogonal. $S=0.5$: feature $1$ alone and features $2$, $3$ as an antipodal pair. $S=0.8$: two antipodal pairs, biases below $0.01$. $S=0.95$: all five, nearly a regular pentagon, with norms $1.09$ to $1.14$ and biases $-0.22$ to $-0.24$.}\label{fig:superposition:pentagon}
\end{fullwidth}
\end{figure}

\Cref{fig:superposition:pentagon} tests prediction (iii) with $p=5$ and $m=2$, so that $2m=4$. In the dense model the two most important features get orthogonal directions and the other three are dropped, with biases within $0.001$ of their means $\frac{1-S}2$, as \cref{thm:superposition:linear,lem:superposition:drop} predict. As $S$ grows, features enter in order of importance: three features at $S=0.5$, as a dedicated direction plus an antipodal pair orthogonal to it, and four at $S=0.8$, as two antipodal pairs, the obtuse configuration of \cref{thm:superposition:sparse} with zero bias. The fifth feature cannot join without positive interference, and at $S=0.95$ it has joined: the five directions form a nearly regular pentagon, with consecutive angles between $68.5^\circ$ and $75^\circ$, neighbors have inner products between $0.34$ and $0.45$, and every bias is about $-0.23$, the negative threshold that \cref{thm:superposition:readout} uses. The features per dimension, $\norm{W}_F^2/m$, are $1.0$, $1.4$, $2.0$ and $3.2$, and the dimensionalities are close to those of \cref{prop:superposition:frame}: $1$, $0.54$ and $0.46$ at $S=0.5$ (a dedicated feature and an unequal antipodal pair), $0.49$ to $0.51$ at $S=0.8$, and $0.38$ to $0.41$ at $S=0.95$, against $\frac25$ for a regular pentagon.

\paragraph{Limits of the predictions.} The phase boundary of the two-feature model is numerical: we prove which side of each curve superposes, not where the true transition is, and we know no closed form for it. The capacity bound is worst case and loose for random inputs. And all predictions concern optimal configurations; the stalls in \cref{fig:superposition:two} show that training can miss them.

\section{What the model gets wrong}\label{sec:superposition:wrong}

\paragraph{What was thrown away.} \emph{The features are given.} In a network, features are not inputs but functions of the input, learned along with the weights that store them; the model says how a fixed set would be stored, not which set would be learned, and so it cannot say what a polysemantic neuron means. \emph{The features are independent} and identically distributed apart from importance. Correlations change everything the geometry depends on: features that never co-occur can be stored antipodally for free (\cref{ex:superposition:anticorrelated}), and features that always co-occur interfere in every input. \emph{The values are magnitudes} in $[0,1]$; with binary features the thresholds move (\cref{ex:superposition:binary}), and features with signs cannot use the clipping of negative interference. \emph{The task is reconstruction} through tied weights and one layer, so the model has nothing to say about computing with features stored in superposition. \emph{The loss is the population loss}, so there is no memorization and no finite-sample effect.

\paragraph{Where it breaks.} \emph{Computation in superposition.} \Textcite{HanniEtAl2024} (unrefereed) construct one-layer networks that compute pairwise ANDs of many sparse Boolean features, with the inputs themselves in superposition, using far fewer neurons than there are features; there superposition is a resource for computation, not only a way through a bottleneck, and the reconstruction model cannot express this. \emph{Features that are not directions.} \Textcite{EngelsEtAl2025} found, with sparse autoencoders, multidimensional features in language models: days of the week and months of the year lie on circles in GPT-2 and Mistral~7B, and the circles are used to compute with modular arithmetic. A circular feature is a two-dimensional irreducible object, the Fourier feature of \cref{ch:grokking}, and the toy model, in which every feature is a single direction, cannot produce it. \emph{The geometry is not Euclidean.} \Textcite{ParkEtAl2024} formalize the hypothesis that concepts are linear directions, and argue that the inner product in which concept directions are orthogonal is a ``causal'' inner product determined by the model, not the Euclidean one; the toy model's choice of $\ip{\cdot,\cdot}$ is a convention. \emph{The decomposition is not unique.} Sparse autoencoders do not recover a unique dictionary: \textcite{PauloBelrose2025} find that autoencoders differing only in their random seed share only about $30\%$ of their latents in one large configuration, and \textcite{EngelsSmithTegmark2025} show that the reconstruction error, the ``dark matter'', is structured, with about half of the error vector linearly predictable from the input activation. In the toy model the true features are given and every optimal configuration stores them; in a network one can find many sparse dictionaries, and the model offers no criterion for choosing among them. \emph{Data, not features.} With a finite training set, \textcite{HenighanEtAl2023} found that the same toy model stores individual training examples in superposition when the data set is small and features when it is large, with double descent in between (\cref{ch:doubledescent}); the infinite-data model sees only the second regime.

\paragraph{Fine structure.} The geometry of optimal configurations with $p>2m$ is richer than \cref{thm:superposition:sparse} suggests: \textcite{ElhageEtAl2022} report that the directions organize into polytopes whose dimensionalities cluster at fractions such as $\frac34$, $\frac23$ and $\frac12$, that several configurations match solutions of the Thomson problem of placing charges on a sphere, and that the transition of a feature into superposition is abrupt. \Textcite{ScherlisEtAl2022} (unrefereed) recast superposition in a related toy model as an allocation of capacity, the dimensionality $\supD_i$, among features: the most important features get dimensions of their own, less important ones share, the least important are dropped, and sharing is more common for sparser and heavier-tailed inputs. \Textcite{LiuLiuGore2025} vary the strength of superposition in the toy model with weight decay and find that in the strong-superposition regime the loss scales inversely with the dimension $m$ for a broad class of feature frequencies, a mechanism for the scaling laws of \cref{ch:scaling}. On the side of dictionaries, \textcite{GaoEtAl2025} find clean scaling laws for $k$-sparse autoencoders. What is open: the dynamics by which gradient descent reaches these polytopes, the theory of computation in superposition beyond constructions, and a criterion that selects one sparse dictionary among many.

\section{Model card}\label{sec:superposition:card}

\modelcard{Polysemantic neurons; activations that decompose into sparse combinations of many more directions than dimensions.}{$p$ independent sparse features in $[0,1]$, stored in $m<p$ dimensions by $W$ and read out by $\relu(W^\top Wx+b)$, importance-weighted squared loss.}{Linear readout: principal components at every sparsity. ReLU: exact losses for two features in one dimension, superposition provable for $S>\frac{2+r}{2+5r}$ or $r>\kappa(S)/v(S)$; zero first-order interference if and only if obtuse, possible if and only if $p\le2m$; $e^{\eps^2m/4}$ directions with coherence $\eps$ and thresholded readout.}{A first-order transition into superposition with a metastable region; antipodal pairs, cross-polytopes, then polygons with negative bias; loss $O((1-S)^2)$ up to $2m$ features and order $1-S$ beyond; capacity exponential in $m\delta^2/s^2$.}{Features are given, independent, one-dimensional; reconstruction instead of computation; Euclidean geometry; unique true dictionary.}{Computation in superposition; multidimensional and circular features; causal inner products; memorization in superposition; scaling laws from superposition.}

The chapter stands at a junction of the book. Its linear half is \cref{ch:deeplinear}: a linear autoencoder learns principal components, and in stages, one per singular value, so that under gradient flow features enter one at a time, as they do here when the sparsity grows. Its nonlinear half is a phase transition with metastability, and the stalls of \cref{fig:superposition:two} are escapes from a saddle, as in \cref{ch:singleindex,ch:committee}. The almost orthogonal directions of \cref{thm:superposition:packing} are the random features of \cref{ch:lazy} seen from the other side: there random directions are fixed and only a readout is learned; here learned directions are as orthogonal as random ones would be.

The circular features of \textcite{EngelsEtAl2025} are the Fourier features of \cref{ch:grokking}, now found in language models: a two-dimensional irreducible representation where this chapter has one direction. The many directions sharing few dimensions are the opposite of the simplex of \cref{ch:collapse}, in which $K$ class means sit in $K-1$ dimensions, as far apart as possible, because classes are never present together; sparse features are rarely present together, and that is why superposition works. Finally, the dependence of the loss on the dimension through the number of features stored links superposition to the scaling laws of \cref{ch:scaling}.

\begin{exercises}
\item \consolidation (Moments and the antipodal pair with a bias.) (a) Derive \eqref{eq:superposition:moments}. (b) For the antipodal pair with $b=(\beta,\beta)$ and $r=1$, write the loss of feature $1$ as $S^2\beta^2+S(1-S)(\beta^2+\frac{\beta^3}3)+(1-S)^2C(\beta)$ for $0\le\beta\le1$, with $C(\beta)=\E\bigl(U-\relu(U-V+\beta)\bigr)^2$. (c) Show that $C'(0)=-\frac13$, and conclude that a small positive bias improves the antipodal pair at every $S<1$. Why does a positive bias help, although it produces false positives? \label{ex:superposition:antibias}

\item \consolidation (Ties in the linear model.) (a) With equal importances and equal variances, show that every $W$ with $W^\top W$ an orthogonal projection of rank $m$ minimizes $\Loss_{\mathrm{lin}}$, so that for $m=1$, $p=2$ every unit vector $w=(\cos\theta,\sin\theta)$ is optimal, and represents both features when $\sin2\theta\ne0$. (b) Show that $\norm W_F^2=m$ for all these minimizers. (c) Show that if $\omega_1>\omega_2$, then only $\theta\in\{0,\pi\}$ remains. \label{ex:superposition:ties}

\item \consolidation (The function $\kappa$.) (a) Show that for $\tau\in[0,1]$, $f_S(\tau)=S^2\tau^2+\frac{S(1-S)}3\tau^3+(1-S)\bigl(\tau^2-(1-S)\tau+\frac{1-S}3\bigr)$, that $f_S$ is decreasing on $(-\infty,0]$ and increasing on $[1,\infty)$, and that $f_S$ is strictly convex on $[0,1]$. (b) Conclude that for $0<S<1$ the minimizer is the root in $(0,1)$ of $S(1-S)\tau^2+2(S^2+1-S)\tau-(1-S)^2=0$, and that $\tau_0=\frac12$, $\kappa(0)=\frac1{12}$. (c) Compute $\kappa(S)/\supv$ at $S=\frac12$ and find the $S$ at which it equals $\frac12$; compare with the spinodal in \cref{fig:superposition:two}. \label{ex:superposition:kappa}

\item \consolidation\hard (The Welch bound.) Let $u_1,\dots,u_N\in\R^m$ be unit vectors with $N>m$, and let $\mu=\max_{i\ne j}\abs{\ip{u_i,u_j}}$. (a) Let $G$ be their Gram matrix. Show that $\norm G_F^2\ge(\tr G)^2/\rank G\ge N^2/m$, and deduce $\mu^2\ge\frac{N-m}{m(N-1)}$. (b) Conclude that if $N\ge2m$, then $(2s-1)\mu<1$ requires $m>(2s-1)^2/2$, so that \cref{thm:superposition:readout} needs $m$ of order $s^2$ to say anything. \hint{in (a), the nonzero eigenvalues of $G$ add up to $N$, and there are at most $m$ of them; use the Cauchy--Schwarz inequality.}\label{ex:superposition:welch}

\item \consolidation (Tight frames.) (a) Show that the vertices $w_j=(\cos\frac{2\pi j}N,\sin\frac{2\pi j}N)$ of a regular $N$-gon, $N\ge3$, satisfy $\sum_jw_jw_j^\top=\frac N2I_2$, and that this fails for $N=2$. (b) Compute $\supD_i$ for the configurations of \cref{fig:superposition:pentagon} from the data file, and compare with \cref{prop:superposition:frame}. \label{ex:superposition:frames}

\item \variation (No bias.) Force $b=0$ in the two-feature model. (a) Show that the best configuration that drops feature $2$ has loss $r\frac{1-S}3$, and that the antipodal pair still has loss $\frac{1+r}6(1-S)^2$. (b) Conclude that the antipodal pair beats every configuration dropping a feature when $S>\frac{1-r}{1+r}$, so for $r=1$ at every $S>0$, and that at $S=0$, $r=1$ the two tie. Explain why the bias, which lets a dropped feature be replaced by its mean, makes dropping more attractive. \label{ex:superposition:nobias}

\item \variation (Binary features.) Let each feature equal $1$ with probability $1-S$ and $0$ otherwise. (a) Show that dropping feature $2$ costs at least $rS(1-S)$ and that the antipodal pair costs $(1+r)(1-S)^2$. (b) Conclude that superposition beats dropping when $S>\frac{1+r}{1+2r}$, for instance $S>\frac23$ at $r=1$, and compare with $\frac37$ for uniform values. Why do binary features superpose less readily? \label{ex:superposition:binary}

\item \variation (Anticorrelated features.) Let $x_1$ and $x_2$ never be active together: $x=(U,0)$, $(0,U)$, or $(0,0)$ with probabilities $\frac{1-S}2$, $\frac{1-S}2$, $S$. (a) Show that the antipodal pair has loss $0$ for every $S$. (b) More generally, show that if at most one of $p$ features is active at a time, then every obtuse configuration has loss $0$, so $p\le2m$ features fit exactly at every sparsity. (c) Let instead $x_2=x_1$ always. Show that $w=(1,1)/\sqrt2$ with $b=0$ has loss $0$. Correlations decide the geometry. \label{ex:superposition:anticorrelated}

\item \variation\hard (Random sparsity.) Use the configuration of \cref{cor:superposition:capacity} with $\eps=\delta/(2s-1)$ on the random features of \cref{sec:superposition:model}, and let $K$ be the number of active features, a binomial variable with parameters $p$ and $1-S$. (a) Show that, for every $x$, every coordinate has error at most $\max(K-s,0)\eps$ if it is inactive and at most $(s+K-1)\eps$ if it is active. (b) Deduce $\Loss\le\supOmega\bigl[(1-S)\delta^2+\E\bigl((s+K)^2\eps^2\ind\{K>s\}\bigr)\bigr]$, and choose $s$ for $p(1-S)=10$ and $p=10^6$. \label{ex:superposition:random}

\item \simulation (The pentagon and its successors.) Run the script of \cref{fig:superposition:pentagon} for $S\in\{0,0.1,\dots,0.9,0.95,0.98\}$ and record $\norm W_F^2$ and the number of directions of norm above $0.1$. With the better of seeds $0$ and $1$, the reader should observe two features up to $S=0.3$, three at $S=0.4$ and $0.5$, four from $S=0.6$ to $S=0.8$, and five from $S=0.9$ on; the counts depend on the seed (seed $0$ gives four at $S=0.5$, seed $1$ only four at $S=0.9$: training ends in different local optima). The biases of represented features should be above $-0.05$ with up to four features, and near $-0.2$ with five. Repeat with equal importances and describe what changes. \label{ex:superposition:pentasim}

\item \simulation (Escape from the dedicated saddle.) In the two-feature model with $r=\frac12$ and $S=0.9$, start gradient descent on the exact loss (the function \texttt{loss} of the script of \cref{fig:superposition:two}) from $w=(1,-c_0)$, $b=(0,\frac{1-S}2)$ with $c_0\in\{10^{-2},10^{-4},10^{-6}\}$ and step size $\eta=0.1$. Record the number of steps until $\abs{w_2}>\frac12$. The reader should observe about $1200$, $3100$ and $5000$ steps: increments of about $1900$ steps per factor $100$ in $c_0$, as predicted by the linearization, $\log(100)/(2\eta(r\supv-\kappa(S)))\approx1900$, the logistic escape of \cref{ch:singleindex}. Explain why the linearization predicts the increments but not the counts themselves. \label{ex:superposition:escape}

\item \research (The phase boundary.) Prove that the dedicated configuration is a global minimizer of the two-feature loss for $S<S_c(r)$ for some explicit function $S_c$, and find the exact transition. The optimum on the superposed side has nonzero biases and unequal norms, and we know of no closed form for it; part (iii) of \cref{thm:superposition:two} and the heuristic spinodal give one side only. \label{ex:superposition:boundary}

\item \research (Computation in superposition.) Replace reconstruction by a computation, for instance the outputs $\abs{x_i-x_j}$ or $x_i\wedge x_j$ for selected pairs, with a hidden ReLU layer of width $m$. How many pairs can a width-$m$ network compute when the inputs are $s$-sparse? \Textcite{HanniEtAl2024} give constructions; matching lower bounds, the role of training, and the geometry of learned solutions are open. \label{ex:superposition:computation}
\end{exercises}

\notesandrefs
Overcomplete sparse representations are older than deep learning: \textcite{OlshausenField1997} learned dictionaries with more atoms than dimensions for natural images and argued that cortex uses such codes, and compressed sensing \parencite{CandesEtAl2006,Donoho2006} made precise when a sparse vector can be recovered from few linear measurements. The packing argument of \cref{thm:superposition:packing} is the probabilistic method in the form used for the lemma of \textcite{JohnsonLindenstrauss1984}, with Hoeffding's inequality \parencite{Hoeffding1963}, and \cref{lem:superposition:2m} is a classical fact about obtuse sets of vectors. \Cref{thm:superposition:linear} is the Eckart--Young theorem \parencite{EckartYoung1936} specialized to a diagonal matrix, with the bias handled first; its proof by a linear program on the diagonal of a projection is Ky Fan's maximum principle in disguise.

The superposition hypothesis grew out of the circuits program \parencite{OlahEtAl2020} and was made concrete by \textcite{ElhageEtAl2022}, whose toy model this chapter solves in its simplest cases; they also studied the phase diagram of two features in one dimension and the polytope geometry in larger cases numerically, and compared the former with the losses of a few configurations. \Textcite{HenighanEtAl2023} extended the toy model to finite data, and \textcite{ScherlisEtAl2022} to an allocation of capacity. Word senses as sparse combinations of directions are due to \textcite{AroraEtAl2018}. The dictionary-learning program followed: \textcite{CunninghamEtAl2024} and \textcite{BrickenEtAl2023} at small scale, \textcite{TempletonEtAl2024} and \textcite{GaoEtAl2025} at large scale.\sidenote{The Transformer Circuits Thread is the web publication of Anthropic's interpretability team; its articles are not refereed. \Textcite{ElhageEtAl2022} is also on arXiv.}

The two-feature analysis of \cref{thm:superposition:two}, the spinodal, and \cref{thm:superposition:sparse} were derived for this chapter. The comparison of explicit configurations is in the spirit of \textcite{ElhageEtAl2022}, but we have not compared our formulas with theirs line by line, and the bounds on the side of superposition are new as stated here as far as we know.

\paragraph{Translating the literature.} \Textcite{ElhageEtAl2022} write $n$ for the number of features (our $p$), $m$ for the number of dimensions (our $m$), $W\in\R^{m\times n}$ with columns $W_i$ (our $w_i$), $S$ for the sparsity, as here, $I_i$ for the importance (our $\omega_i$; we keep $I$ for the identity), and $x'=\relu(W^\top Wx+b)$ for the reconstruction, as here; their feature dimensionality $D_i=\norm{W_i}^2/\sum_j(\hat W_i\cdot W_j)^2$, with $\hat W_i$ the unit vector, equals our $\supD_i$, and their dimensions per feature $m/\norm W_F^2$ is the inverse of our features per dimension. In \textcite{BrickenEtAl2023} and \textcite{TempletonEtAl2024} the activation vector of the layer being decomposed plays the role of our $h$, the number of dictionary features plays the role of our $p$, and the rows of the decoder play the role of our directions $w_i$; the encoder's ReLU with bias is our thresholded readout.

\printchapterbib
\end{refsection}

% Chapter 12 of The Hydrogen Atoms of Deep Learning: neural collapse.
\newcommand{\collapseone}{\mathbf 1}
\newcommand{\collapseYt}{\widetilde Y}
\newcommand{\collapseHb}{\bar H}
\newcommand{\collapsenuc}[1]{\norm{#1}_*}

\casechapter{Neural Collapse}{The geometry of the last layer}{Unconstrained features}{ch:collapse}
\begin{refsection}[chapters/12-collapse.bib]

\begin{chaptersummary}
Train a deep classifier long past zero training error and its last layer becomes rigid: the features of each class shrink to their mean, the $K$ class means spread out into a regular simplex, the classifier vectors line up with the means, and the network classifies by the nearest class mean. The hydrogen atom of this neural collapse treats the last-layer features as free variables, constrained only in size. We prove that every global minimizer of the cross-entropy with norm constraints is a simplex, solve the weight-decay and square-loss versions in closed form, find a transition at a single dimensionless strength $\kappa=K\sqrt{n\lambda_W\lambda_H}$ of the regularization, and show that the max-margin problem reached by unregularized gradient descent has the same solution. The model explains the geometry and predicts how it breaks under class imbalance; it is silent on test data, on depth, and on whether collapse is good for anything.
\end{chaptersummary}

\section{The phenomenon: four regularities at the end of training}\label{sec:collapse:phenomenon}

\Cref{ch:margin} followed a linear classifier on fixed features into the regime past zero training error, where the cross-entropy keeps falling and the weights keep growing. In a network the features below the last layer keep moving as well. What do they converge to?

\Textcite{PapyanHanDonoho2020} measured it. They trained networks from three families (VGG, ResNet and DenseNet) on seven image classification data sets (MNIST, FashionMNIST, CIFAR10, CIFAR100, SVHN, STL10 and ImageNet), by stochastic gradient descent with momentum and without data augmentation, far into what they called the \term{terminal phase of training}\index{terminal phase of training}: the epochs after the training error first reaches zero. Throughout that phase they recorded the penultimate-layer features $h\in\R^m$ of the training examples and the last-layer weights. Write $K$ for the number of classes, $h_{k,i}$ for the features of the $i$th training example of class $k$, $\bar h_k$ for the mean of class $k$, $h_G$ for the global mean, and $w_k$ for the $k$th row of the last weight matrix $W\in\R^{K\times m}$, so that the logits of an input with features $h$ are $Wh+b$.

\begin{observation}[\textcite{PapyanHanDonoho2020}]\label[observation]{obs:collapse:nc}
During the terminal phase of training, with the cross-entropy loss and weight decay, four quantities tend toward zero together:
\begin{description}
\item[(NC1) variability collapse:] the within-class scatter, relative to the between-class scatter of the means;
\item[(NC2) convergence to a simplex:] the spread of the norms $\norm{\bar h_k-h_G}$, and the deviation of every cosine between two centered means from $-1/(K-1)$;
\item[(NC3) self-duality:] the distance between $W/\norm W_F$ and the matrix of centered means, normalized the same way;
\item[(NC4) nearest class center:] the fraction of training points on which the network's prediction differs from the class of the nearest mean.
\end{description}
\end{observation}

The name \term{neural collapse}\index{neural collapse} covers all four. NC1 is measured by $\tr(\Sigma_W\Sigma_B^+)/K$, with $\Sigma_W$ the average within-class covariance, $\Sigma_B$ the covariance of the class means and $\Sigma_B^+$ its pseudo-inverse. NC2 says that the centered means become the vertices of a regular simplex centered at the origin, the most spread-out configuration of $K$ points of equal norm; NC3 says that the classifier is that same simplex, rescaled; NC4 follows from the other three (\cref{ex:collapse:ncc}). \Textcite{PapyanHanDonoho2020} argued that the collapsed geometry brings benefits in generalization, robustness and interpretability, and their paper started a literature: \textcite{HanPapyanDonoho2022} reported the same four regularities for networks trained with the square loss, across three networks and five data sets, and \textcite{HeSu2023} found a regularity one level down, a \term{law of equi-separation}\index{law of equi-separation}: in a trained network each layer improves a measure of class separation by roughly the same factor, so that collapse builds up geometrically with depth.\sidenote{\textcite{RangamaniEtAl2023} report a related ``intermediate neural collapse'': the within-class covariance, relative to the between-class one, decreases with depth.}

\begin{figure}[t]
\centering
\begin{tikzpicture}[scale=1.1,font=\small]
\pgfmathsetseed{7}
\foreach \a/\c in {90/sea,210/seadeep,330/sandline!55!black}{
  \foreach \j in {1,...,24}{
    \pgfmathsetmacro{\rx}{0.42*rand}\pgfmathsetmacro{\ry}{0.42*rand}
    \fill[\c,opacity=0.7] ({1.55*cos(\a)+\rx},{1.55*sin(\a)+\ry}) circle (1.1pt);}
  \foreach \j in {1,...,14}{
    \pgfmathsetmacro{\rx}{0.06*rand}\pgfmathsetmacro{\ry}{0.06*rand}
    \fill[\c] ({4.9+1.55*cos(\a)+\rx},{1.55*sin(\a)+\ry}) circle (1.1pt);}
  \draw[->,\c,line width=0.8pt] (4.9,0) -- ({4.9+1.05*cos(\a)},{1.05*sin(\a)});
}
\draw[muted,densely dashed] ({4.9+1.55*cos(90)},{1.55*sin(90)}) -- ({4.9+1.55*cos(210)},{1.55*sin(210)}) -- ({4.9+1.55*cos(330)},{1.55*sin(330)}) -- cycle;
\fill[ink] (0,0) circle (0.9pt); \fill[ink] (4.9,0) circle (0.9pt);
\node[anchor=north] at (0,-1.25) {early in training};
\node[anchor=north] at (4.9,-1.25) {terminal phase};
\node[anchor=west,text=sea] at ({4.9+0.12},0.75) {$w_1$};
\node[anchor=west] at ({4.9+1.55*cos(90)+0.15},{1.55*sin(90)}) {$\bar h_1$};
\end{tikzpicture}
\caption{Neural collapse for $K=3$ classes, drawn in the plane of the class means. Left: the penultimate features of the training examples (dots, one color per class) form clouds. Right: in the terminal phase the clouds shrink to their means (NC1), the means form an equilateral triangle centered at the global mean, which is the simplex of NC2, and the classifier vectors $w_k$ (arrows) point at the means (NC3).}\label{fig:collapse:cartoon}
\end{figure}

The status of \cref{obs:collapse:nc} is good on training data and weaker elsewhere. On the training set the four metrics decrease steadily in a wide range of architectures and data sets, but they do not reach zero in finite time, and how close they get depends on the training time, the weight decay and the learning-rate schedule. In a small experiment of our own (\cref{fig:collapse:network}) NC1 fell by a factor of thirty and then stalled. On test data collapse is much weaker \parencite{HuiBelkinNakkiran2022}, and NC2 needs at least $K-1$ feature dimensions, so it cannot hold literally when the classes outnumber the dimensions, as they do for the vocabulary of a language model (\cref{sec:collapse:wrong}).

The regularity is still striking. A network has millions of parameters and its features are produced by a long nonlinear computation, yet at the end of training the last layer settles into the most symmetric configuration available, the same for every architecture. Why a simplex? And why does the classifier copy it?

\section{The model: free features of bounded size}\label{sec:collapse:model}

The data are $K\ge2$ classes with $n$ training examples each; deviating from the global convention, $n$ counts the examples per class, so the training set has $Kn$ points. The model keeps only the last layer. The classifier is $W\in\R^{K\times m}$ with rows $w_1,\dots,w_K$, and the features are $Kn$ vectors $h_{k,i}\in\R^m$, collected as the columns of $H\in\R^{m\times Kn}$. The features are \emph{free variables}: they are optimized directly, as if the network below the last layer could produce whatever features it likes. This is the \term{unconstrained features model}\index{unconstrained features model}, introduced under that name by \textcite{MixonParshallPi2020} and as the \term{layer-peeled model}\index{layer-peeled model} by \textcite{FangEtAl2021}. The loss of an example with logits $z\in\R^K$ and class $k$ is the cross-entropy
\[
  \ell(z,k)=\log\sum_{j=1}^Ke^{z_j-z_k},
\]
and the training loss is $\Lhat(W,H)=\frac1{Kn}\sum_{k=1}^K\sum_{i=1}^n\ell(Wh_{k,i},k)$, without a bias; adding one changes nothing in the main theorem (\cref{ex:collapse:bias}). Write $\bar w=\frac1K\sum_kw_k$ for the mean classifier and $\bar h_k=\frac1n\sum_ih_{k,i}$ for the class means.

The loss alone has no minimizer: scaling up $W$ and $H$ in a separating configuration drives it to zero. Something must bound the size, and the literature uses three devices, each a stand-in for a feature of real training.

\begin{assumption}[Size of the last layer]\label[assumption]{ass:collapse:size}
One of the following holds.
\begin{enumerate}
\item (Norm constraints.) For given $E_W,E_H>0$, the parameters satisfy $\frac1K\norm W_F^2\le E_W$ and $\frac1{Kn}\norm H_F^2\le E_H$.
\item (Weight decay.) The objective is $\Lhat(W,H)+\frac{\lambda_W}2\norm W_F^2+\frac{\lambda_H}2\norm H_F^2$ with $\lambda_W,\lambda_H>0$.
\item (Margins.) No regularization; the object of study is the direction of $(W,H)$ along gradient descent, as in \cref{ch:margin}.
\end{enumerate}
\end{assumption}

Norm constraints are the cleanest for proofs; weight decay is what practice uses, on all layers, and a penalty on the features stands in for the penalty on the layers that produce them; and without regularization the margin story of \cref{ch:margin} applies, because the logits $Wh$ are homogeneous of degree two in $(W,H)$. All three give the same geometry.

\begin{dictionary}
classifier $W$, rows $w_k$ & last weight matrix\\
free features $h_{k,i}\in\R^m$ & penultimate-layer activations of the training examples\\
bound on $\norm H_F$ & weight decay on all the layers that produce the features\\
$m\ge K-1$ & penultimate width at least the number of classes minus one\\
class means $\bar h_k$, $h_{k,i}=\bar h_k$ & NC1\\
simplex equiangular tight frame & NC2\\
$w_k\propto\bar h_k$ & NC3\\
global minimizer, or max-margin direction & the end of the terminal phase of training\\
regularization strength $\kappa=K\sqrt{n\lambda_W\lambda_H}$ & weight decay times the square root of the data per class\\
\end{dictionary}

\paragraph{What was thrown away.} Everything that makes the features depend on the inputs. The model assumes that the network can realize every configuration of last-layer features at a cost that depends only on their norms. Gone too are the test set, the layers below the last, the optimizer, the bias, and every asymmetry between classes. Expressiveness is the right assumption to make first, because a network that drives the training error to zero has already shown that it can place its features almost where it likes; and the question of \cref{obs:collapse:nc} is about the training features only.

\paragraph{What one had to notice.} That NC3 and NC4 make the last layer look like the solution of an optimization problem in which the features and the classifier play symmetric roles; and that a problem with that much symmetry should have a symmetric ground state. The problem is invariant under every rotation $Q\in O(m)$ acting on both sides, $(W,H)\mapsto(WQ^\top,QH)$, and under every relabeling of the classes. The simplex is the configuration of $K$ points that relabelings map to itself by rotations. Its minimizers will form a single orbit of the rotation group: the permutation symmetry is preserved, and the rotational symmetry is broken, as in a ferromagnet below its Curie point, where the magnetization picks a direction.

\section{The solution: the simplex is the ground state}\label{sec:collapse:solution}

We first describe the simplex and the property that makes it optimal.

\begin{definition}[Simplex equiangular tight frame]\label[definition]{def:collapse:etf}
Vectors $v_1,\dots,v_K\in\R^m$ form a \term{simplex equiangular tight frame}\index{simplex equiangular tight frame} of squared norm $r>0$ if
\[
  \ip{v_j,v_k}=r\Bigl(\delta_{jk}-\frac{1-\delta_{jk}}{K-1}\Bigr)\qquad\text{for all }j,k ,
\]
that is, if their Gram matrix is $\frac{rK}{K-1}(I_K-\frac1K\collapseone_K\collapseone_K^\top)$.
\end{definition}

The Gram matrix is $\frac{rK}{K-1}$ times the orthogonal projection onto the complement of $\collapseone_K$, so it is positive semidefinite of rank $K-1$, and a simplex equiangular tight frame exists in $\R^m$ if and only if $m\ge K-1$. The vectors sum to zero, because $\norm{\sum_kv_k}^2=\collapseone^\top G\collapseone=0$ for the Gram matrix $G$, and every two frames with the same $r$ differ by a rotation of $\R^m$. For $K=2$ the frame is a pair of antipodal points, for $K=3$ an equilateral triangle, for $K=4$ a regular tetrahedron. ``Tight frame'' refers to $\sum_kv_kv_k^\top$ being a multiple of the projection onto the span of the $v_k$; ``equiangular'' to the equal angles.\sidenote{Equiangular tight frames are a classical subject in frame theory and coding, where they solve packing problems for lines and points on spheres \parencite{StrohmerHeath2003}.}

\begin{lemma}[The simplex is the best code]\label[lemma]{lem:collapse:simplex}
Let $u_1,\dots,u_K$ be unit vectors in $\R^m$ with $K\ge2$. Then
\[
  \max_{j\ne k}\ip{u_j,u_k}\ge-\frac1{K-1},
\]
with equality if and only if $u_1,\dots,u_K$ form a simplex equiangular tight frame of squared norm $1$.
\end{lemma}

\begin{proof}
Write $\beta=\max_{j\ne k}\ip{u_j,u_k}$. Expanding a square,
\[
  0\le\Bigl\|\sum_ku_k\Bigr\|^2=K+\sum_{j\ne k}\ip{u_j,u_k}\le K+K(K-1)\beta ,
\]
which gives the inequality. Equality holds if and only if both steps are equalities: every $\ip{u_j,u_k}$ with $j\ne k$ equals $\beta=-1/(K-1)$, which is the definition of the frame, and then $\sum_ku_k=0$ holds automatically.
\end{proof}

Among $K$ points on a sphere, the simplex makes the largest pairwise correlation as small as possible, and therefore the smallest pairwise distance as large as possible. That is already the shape of the answer. What remains is to show that the loss and the classifier force the features onto a sphere and push them apart, and this is the job of one inequality.

\begin{lemma}[Alignment]\label[lemma]{lem:collapse:align}
For every $W\in\R^{K\times m}$ and $H\in\R^{m\times Kn}$,
\begin{equation}\label{eq:collapse:align}
  \frac1{Kn}\sum_{k,i}\ip{w_k-\bar w,h_{k,i}}\le\Bigl(\frac1K\norm W_F^2\Bigr)^{1/2}\Bigl(\frac1{Kn}\norm H_F^2\Bigr)^{1/2}.
\end{equation}
If equality holds and the left side is positive, then $\bar w=0$, $h_{k,i}=\bar h_k$ for all $k$ and $i$, and $w_k=c\,\bar h_k$ for all $k$, with $c=\sqrt n\norm W_F/\norm H_F>0$.
\end{lemma}

\begin{proof}
The plan is to average over the examples of each class, apply the Cauchy--Schwarz inequality in $\R^{K\times m}$, and then drop two nonnegative terms. Averaging over $i$ turns the left side of \eqref{eq:collapse:align} into $\frac1K\sum_k\ip{w_k-\bar w,\bar h_k}$, and the Cauchy--Schwarz inequality bounds this by
\[
  \Bigl(\frac1K\sum_k\norm{w_k-\bar w}^2\Bigr)^{1/2}\Bigl(\frac1K\sum_k\norm{\bar h_k}^2\Bigr)^{1/2}.
\]
Two variance identities finish the proof:
\[
  \sum_k\norm{w_k-\bar w}^2=\norm W_F^2-K\norm{\bar w}^2,\qquad
  \norm{\bar h_k}^2=\frac1n\sum_i\norm{h_{k,i}}^2-\frac1n\sum_i\norm{h_{k,i}-\bar h_k}^2 .
\]
Suppose now that equality holds in \eqref{eq:collapse:align} with a positive left side. Then both dropped terms vanish, so $\bar w=0$ and $h_{k,i}=\bar h_k$, and equality holds in the Cauchy--Schwarz inequality with a positive inner product, so $w_k-\bar w=c\,\bar h_k$ for a constant $c>0$. With $\bar w=0$ this reads $w_k=c\,\bar h_k$, and then $\norm W_F^2=c^2\sum_k\norm{\bar h_k}^2=c^2\norm H_F^2/n$ determines $c$.
\end{proof}

The left side of \eqref{eq:collapse:align} is the average over examples of the gap between the mean logit and the logit of the correct class, with the sign reversed. The lemma says that this average gap is at most the product of the root-mean-square sizes of the classifier and the features, and that it can reach that bound only in a configuration that already satisfies NC1 and NC3. The theorem adds NC2.

\begin{theorem}[Neural collapse under norm constraints]\label{thm:collapse:main}
Let $K\ge2$, $m\ge K-1$, and $E_W,E_H>0$, and minimize $\Lhat(W,H)$ subject to $\frac1K\norm W_F^2\le E_W$ and $\frac1{Kn}\norm H_F^2\le E_H$. The minimum value is
\begin{equation}\label{eq:collapse:minval}
  \Lhat^\star=\log\Bigl(1+(K-1)\exp\Bigl(-\frac K{K-1}\sqrt{E_WE_H}\Bigr)\Bigr),
\end{equation}
and $(W,H)$ is a minimizer if and only if
\begin{enumerate}
\item (NC1) $h_{k,i}=\bar h_k$ for all $k$ and $i$;
\item (NC2) $\bar h_1,\dots,\bar h_K$ form a simplex equiangular tight frame of squared norm $E_H$;
\item (NC3) $w_k=\sqrt{E_W/E_H}\,\bar h_k$ for all $k$.
\end{enumerate}
\end{theorem}

\begin{proof}
The plan is to bound the loss below in three steps, by the arithmetic--geometric mean inequality inside each example, by Jensen's inequality across examples, and by \cref{lem:collapse:align}, and then to read off the configuration from the cases of equality. Fix a feasible $(W,H)$ and write $z^{k,i}=Wh_{k,i}$ for the logits of example $(k,i)$.

\emph{Step 1: within an example.} The $K-1$ numbers $e^{z_j-z_k}$ with $j\ne k$ have geometric mean $\exp\bigl(\frac1{K-1}\sum_{j\ne k}(z_j-z_k)\bigr)$, and $\frac1{K-1}\sum_{j\ne k}(z_j-z_k)=\frac K{K-1}(\bar z-z_k)$ with $\bar z=\frac1K\sum_jz_j$. Put $\psi(u)=\log(1+(K-1)e^u)$. The arithmetic--geometric mean inequality gives
\[
  \ell(z^{k,i},k)=\log\Bigl(1+\sum_{j\ne k}e^{z^{k,i}_j-z^{k,i}_k}\Bigr)\ge\psi(u_{k,i}),\qquad u_{k,i}=\frac K{K-1}\ip{\bar w-w_k,h_{k,i}},
\]
because $\bar z-z_k=\ip{\bar w-w_k,h_{k,i}}$ and $\log(1+x)$ is increasing.

\emph{Step 2: across examples.} The function $\psi$ is strictly convex and strictly increasing, since $\psi'(u)=(K-1)e^u/(1+(K-1)e^u)$ is positive and strictly increasing. By Jensen's inequality, $\Lhat(W,H)\ge\psi(\bar u)$ with $\bar u=\frac1{Kn}\sum_{k,i}u_{k,i}$.

\emph{Step 3: the alignment bound.} By \cref{lem:collapse:align} and the constraints, $\bar u\ge-\frac K{K-1}\sqrt{E_WE_H}$, and since $\psi$ is increasing, $\Lhat(W,H)\ge\psi(-\frac K{K-1}\sqrt{E_WE_H})$, which is the right side of \eqref{eq:collapse:minval}.

\emph{Attainment.} Let $(W,H)$ satisfy (i)--(iii); such configurations exist because $m\ge K-1$, and they are feasible with equality. Put $c=\sqrt{E_W/E_H}$. For $j\ne k$, $z^{k,i}_j-z^{k,i}_k=c(\ip{\bar h_j,\bar h_k}-\norm{\bar h_k}^2)=-cE_H\frac K{K-1}=-\frac K{K-1}\sqrt{E_WE_H}$, so every example has loss equal to the right side of \eqref{eq:collapse:minval}. This proves \eqref{eq:collapse:minval} and that (i)--(iii) describe minimizers.

\emph{Equality.} Conversely, let $(W,H)$ be a minimizer, so that all three steps hold with equality. In Step 3, since $\psi$ is strictly increasing, the left side of \eqref{eq:collapse:align} equals $\sqrt{E_WE_H}$, which is positive and at least the right side of \eqref{eq:collapse:align} by the constraints. So equality holds in \cref{lem:collapse:align} and in both constraints: $\bar w=0$, NC1 holds, and $w_k=c\bar h_k$ with $c=\sqrt n\norm W_F/\norm H_F=\sqrt{E_W/E_H}$, which is NC3. In Step 2, strict convexity forces all $u_{k,i}$ to be equal; here $u_{k,i}=-\frac K{K-1}c\norm{\bar h_k}^2$, so all the means have the same squared norm $r$. In Step 1, equality in the arithmetic--geometric mean inequality forces the numbers $z_j-z_k=c(\ip{\bar h_j,\bar h_k}-r)$, $j\ne k$, to be equal for each $k$, so $\ip{\bar h_j,\bar h_k}=\beta_k$ for all $j\ne k$; by symmetry of the inner product $\beta_k=\beta_j$, so all the off-diagonal inner products equal one number $\beta$. Finally $\sum_k\bar h_k=\bar w/c=0$, and
\[
  0=\Bigl\|\sum_k\bar h_k\Bigr\|^2=Kr+K(K-1)\beta
\]
gives $\beta=-r/(K-1)$. The constraint on $H$ holds with equality and $h_{k,i}=\bar h_k$, so $r=\frac1K\sum_k\norm{\bar h_k}^2=E_H$, which is NC2.
\end{proof}

The proof follows \textcite{FangEtAl2021} and \textcite{LuSteinerberger2022} in outline, with the equality cases made explicit; \textcite{GrafEtAl2021} proved a version with spherical constraints at about the same time. Each inequality produces one collapse property: Jensen's inequality within a class gives NC1, the Cauchy--Schwarz step gives NC3, and the arithmetic--geometric mean inequality inside each example gives NC2. The minimizer is unique up to rotations of $\R^m$, the orbit predicted in \cref{sec:collapse:model}, and every logit gap of every example equals $\frac K{K-1}\sqrt{E_WE_H}$: the configuration of \cref{lem:collapse:simplex}.

\paragraph{Weight decay.} The penalized problem reduces to the constrained one at the norms the minimizer chooses, and the choice of norms is a problem in one variable. Write
\[
  \kappa=K\sqrt{n\lambda_W\lambda_H},\qquad\rho(W,H)=\Bigl(\frac1K\norm W_F^2\cdot\frac1{Kn}\norm H_F^2\Bigr)^{1/2}=\frac{\norm W_F\norm H_F}{K\sqrt n}.
\]
The number $\kappa$ is the only combination of $\lambda_W$, $\lambda_H$ and $n$ that matters, and $\rho$ is the root-mean-square size of the logits that the last layer can produce.

\begin{corollary}[Neural collapse under weight decay]\label[corollary]{cor:collapse:wd}
Let $K\ge2$, $m\ge K-1$ and $\lambda_W,\lambda_H>0$, and minimize $F(W,H)=\Lhat(W,H)+\frac{\lambda_W}2\norm W_F^2+\frac{\lambda_H}2\norm H_F^2$.
\begin{enumerate}
\item If $\kappa\ge1$, the unique minimizer is $W=0$, $H=0$.
\item If $\kappa<1$, the minimizers are exactly the configurations of \cref{thm:collapse:main} with $\lambda_W\norm W_F^2=\lambda_H\norm H_F^2$ and
\begin{equation}\label{eq:collapse:rhostar}
  \rho=\rho^\star=\frac{K-1}K\log\Bigl(1+\frac{K(1-\kappa)}\kappa\Bigr).
\end{equation}
At every minimizer the training loss is $\log\frac K{K-(K-1)\kappa}$, and every training example is assigned probability $1-\frac{K-1}K\kappa$ to its class.
\end{enumerate}
\end{corollary}

\begin{proof}
The plan is to bound $F$ below by a convex function $g$ of $\rho$ alone, using \cref{thm:collapse:main} for the loss and the arithmetic--geometric mean inequality for the penalty, and then to minimize $g$. Put
\[
  g(\rho)=\psi\Bigl(-\frac K{K-1}\rho\Bigr)+\kappa\rho,\qquad\rho\ge0,
\]
with $\psi$ as in the proof of \cref{thm:collapse:main}. A minimizer of $F$ exists, because $F$ is continuous and $F(W,H)\ge\frac{\lambda_W}2\norm W_F^2+\frac{\lambda_H}2\norm H_F^2$ tends to infinity with $(W,H)$.

\emph{The lower bound.} Fix $(W,H)$ and set $E_W=\frac1K\norm W_F^2$, $E_H=\frac1{Kn}\norm H_F^2$, so that $\rho=\sqrt{E_WE_H}$. If $\rho>0$, then $(W,H)$ is feasible for \cref{thm:collapse:main} with these budgets, and $\Lhat(W,H)\ge\psi(-\frac K{K-1}\rho)$. If $\rho=0$, then $W=0$ or $H=0$, all logits vanish, and $\Lhat=\log K=\psi(0)$. For the penalty, the arithmetic--geometric mean inequality gives
\[
  \frac{\lambda_W}2\norm W_F^2+\frac{\lambda_H}2\norm H_F^2=\frac K2\bigl(\lambda_WE_W+n\lambda_HE_H\bigr)\ge K\sqrt{n\lambda_W\lambda_HE_WE_H}=\kappa\rho ,
\]
with equality if and only if $\lambda_WE_W=n\lambda_HE_H$, that is, $\lambda_W\norm W_F^2=\lambda_H\norm H_F^2$. Hence $F(W,H)\ge g(\rho(W,H))$.

\emph{The function $g$.} It is strictly convex, as a strictly convex function of a linear function plus a linear function. Its derivative is $g'(\rho)=\kappa-\frac K{K-1}\psi'(-\frac K{K-1}\rho)$, and $\psi'(0)=\frac{K-1}K$ gives $g'(0)=\kappa-1$. If $\kappa\ge1$, then $g'\ge g'(0)\ge0$ on $[0,\infty)$, and $g$ is uniquely minimized at $\rho=0$. If $\kappa<1$, then $g'(0)<0$ and $g'(\rho)\to\kappa>0$ as $\rho\to\infty$, so $g$ is uniquely minimized at the root of $g'$. With $s=\frac K{K-1}\rho$ the equation $g'=0$ reads $Ke^{-s}=\kappa\bigl(1+(K-1)e^{-s}\bigr)$, so $e^{-s}=\kappa/(K-(K-1)\kappa)$, which is \eqref{eq:collapse:rhostar}.

\emph{Minimizers.} The bound $F\ge g(\rho)\ge\min g$ is attained: for $\kappa\ge1$ by $(0,0)$, and for $\kappa<1$ by a configuration of \cref{thm:collapse:main} with $\rho=\rho^\star$ and balanced norms, which exists because $m\ge K-1$. So $\min F=\min g$, and $(W,H)$ is a minimizer if and only if $\rho(W,H)$ minimizes $g$ and both inequalities of the lower bound are equalities. For $\kappa\ge1$ this means $\rho=0$ and $\lambda_W\norm W_F^2=\lambda_H\norm H_F^2$, so $W=0$ and $H=0$. For $\kappa<1$ it means $\rho=\rho^\star>0$, balanced norms, and $\Lhat(W,H)$ equal to the minimum of \cref{thm:collapse:main} at the budgets $E_W,E_H$, which by that theorem holds exactly for the configurations it describes. The loss is then $\psi(-s)=\log\bigl(1+(K-1)e^{-s}\bigr)=\log\frac K{K-(K-1)\kappa}$, and each example's probability of its class is $1/(1+(K-1)e^{-s})=1-\frac{K-1}K\kappa$.
\end{proof}

\paragraph{The square loss.} \Textcite{HanPapyanDonoho2022} found neural collapse with the square loss, and in the unconstrained features model the square loss can be solved in closed form. Let $Y\in\R^{K\times Kn}$ be the matrix of one-hot labels, whose column $(k,i)$ is $e_k$, and add an unpenalized bias $b\in\R^K$:
\begin{equation}\label{eq:collapse:mse}
  F_{\mathrm{sq}}(W,H,b)=\frac1{2Kn}\bigl\|WH+b\collapseone_{Kn}^\top-Y\bigr\|_F^2+\frac{\lambda_W}2\norm W_F^2+\frac{\lambda_H}2\norm H_F^2 .
\end{equation}
We use the nuclear norm $\collapsenuc Z$, the sum of the singular values of $Z$, and two facts about it: it is the dual of the operator norm, $\collapsenuc Z=\max\{\ip{G,Z}:\norm G_{\op}\le1\}$ with $\ip{G,Z}=\tr(G^\top Z)$; and consequently $\collapsenuc{WX}\le\norm W_F\norm X_F$, because $\ip{G,WX}=\ip{W^\top G,X}\le\norm{W^\top G}_F\norm X_F\le\norm W_F\norm G_{\op}\norm X_F$.

\begin{proposition}[Square loss: closed form]\label[proposition]{prop:collapse:mse}
Let $K\ge2$, $m\ge K-1$ and $\lambda_W,\lambda_H>0$, and let $\kappa=K\sqrt{n\lambda_W\lambda_H}$.
\begin{enumerate}
\item If $\kappa\ge1$, the unique minimizer of \eqref{eq:collapse:mse} is $W=0$, $H=0$, $b=\frac1K\collapseone_K$.
\item If $\kappa<1$, every minimizer satisfies NC1, $\sum_k\bar h_k=0$, $w_k=\frac\kappa{K\lambda_W}\bar h_k$ (NC3), $b=\frac1K\collapseone_K$, and
\[
  W\bar h_k+b=(1-\kappa)e_k+\frac\kappa K\collapseone_K\qquad\text{for every }k ,
\]
so that the means form a simplex equiangular tight frame (NC2) and the correct-class output is $1-\frac{K-1}K\kappa$. The minimum value is $\frac{K-1}{2K}\kappa(2-\kappa)$.
\end{enumerate}
\end{proposition}

\begin{proof}
The plan is to eliminate the bias, to bound $F_{\mathrm{sq}}$ below by a convex function of $Z=WH$, to solve that convex problem with a dual certificate, and then to read off the configuration from the stationarity conditions. Put $\mu=\sqrt{\lambda_W\lambda_H}$, let $P=I_{Kn}-\frac1{Kn}\collapseone\collapseone^\top$ be the projection that centers the columns, and let $\collapseYt=YP$, whose column $(k,i)$ is $e_k-\frac1K\collapseone_K$.

\emph{The bias.} For fixed $A=WH-Y$, the minimum of $\norm{A+b\collapseone^\top}_F^2$ over $b$ is $\norm{AP}_F^2$, attained only at $b=-\frac1{Kn}A\collapseone$, and $AP=WHP-\collapseYt$.

\emph{The lower bound.} By the arithmetic--geometric mean inequality, $\norm{HP}_F\le\norm H_F$, and the second fact about the nuclear norm,
\begin{equation}\label{eq:collapse:pen}
  \frac{\lambda_W}2\norm W_F^2+\frac{\lambda_H}2\norm H_F^2\ge\mu\norm W_F\norm H_F\ge\mu\norm W_F\norm{HP}_F\ge\mu\collapsenuc{WHP}.
\end{equation}
Hence $F_{\mathrm{sq}}(W,H,b)\ge\Phi(WHP)$ with $\Phi(Z)=\frac1{2Kn}\norm{Z-\collapseYt}_F^2+\mu\collapsenuc Z$.

\emph{The convex problem.} Let $Q=\collapseYt/\sqrt n$. A direct count gives $QQ^\top=I_K-\frac1K\collapseone\collapseone^\top$, an orthogonal projection of rank $K-1$, so all the nonzero singular values of $Q$ equal $1$: $\norm Q_{\op}=1$ and $\collapsenuc Q=\ip{Q,Q}=K-1$. Let $G=Q/\max(1,\kappa)$, so that $\norm G_{\op}\le1$ and $\collapsenuc Z\ge\ip{G,Z}$ for every $Z$. Completing the square,
\[
  \Phi(Z)\ge\frac1{2Kn}\bigl\|Z-\collapseYt+Kn\mu G\bigr\|_F^2+\mathrm{const},
\]
a strictly convex function of $Z$ minimized only at $Z_0=\collapseYt-Kn\mu G=\sqrt n\bigl(1-\kappa/\max(1,\kappa)\bigr)Q$, because $Kn\mu=\sqrt n\kappa$. At $Z_0$ the bound is an equality: $Z_0$ is a nonnegative multiple of $Q$, so $\collapsenuc{Z_0}=\ip{G,Z_0}$ when $\kappa<1$, and $Z_0=0$ when $\kappa\ge1$. So $Z_0$ is the unique minimizer of $\Phi$, with $Z_0=\sqrt n(1-\kappa)Q=(1-\kappa)\collapseYt$ if $\kappa<1$ and $Z_0=0$ otherwise.

\emph{Minimizers.} A minimizer of $F_{\mathrm{sq}}$ exists by coercivity in $(W,H)$ and by the formula for the optimal bias. The configurations described below attain $\Phi(Z_0)$ (\cref{ex:collapse:mseattain}), so $\min F_{\mathrm{sq}}=\Phi(Z_0)$ and every minimizer has $WHP=Z_0$, equality throughout \eqref{eq:collapse:pen}, and the optimal bias. If $\kappa\ge1$, equality in \eqref{eq:collapse:pen} with $\collapsenuc{Z_0}=0$ forces $W=0$ and $H=0$, and then $b=\frac1K\collapseone$; this is part (i). If $\kappa<1$, the residual is
\[
  R=WH+b\collapseone^\top-Y=WHP-\collapseYt=-\kappa\collapseYt .
\]
The objective is differentiable, so at a minimizer its gradients in $H$ and $W$ vanish:
\[
  \frac1{Kn}W^\top R+\lambda_HH=0,\qquad\frac1{Kn}RH^\top+\lambda_WW=0 .
\]
The first gives $h_{k,i}=\frac\kappa{Kn\lambda_H}W^\top(e_k-\frac1K\collapseone)=\frac\kappa{Kn\lambda_H}(w_k-\bar w)$, which depends only on $k$ (NC1) and sums to zero over $k$. The second, with $\sum_{k,i}h_{k,i}=0$, gives $W=\frac\kappa{Kn\lambda_W}\sum_{k,i}(e_k-\frac1K\collapseone)h_{k,i}^\top=\frac\kappa{K\lambda_W}\sum_ke_k\bar h_k^\top$, which is NC3. Since the features are centered, the optimal bias is $b=\frac1K\collapseone-\frac1{Kn}WH\collapseone=\frac1K\collapseone$, and the column of $WH=WHP=Z_0$ for class $k$ is $(1-\kappa)(e_k-\frac1K\collapseone)$; adding $b$ gives the displayed logits. Reading the $j$th entry with NC3, $\frac\kappa{K\lambda_W}\ip{\bar h_j,\bar h_k}=(1-\kappa)(\delta_{jk}-\frac1K)$, so the Gram matrix of the means is a positive multiple of $I_K-\frac1K\collapseone\collapseone^\top$, which is NC2. The minimum value is $\Phi(Z_0)=\frac{\kappa^2}{2Kn}\norm{\collapseYt}_F^2+\mu(1-\kappa)\sqrt n(K-1)=\frac{K-1}{2K}\kappa(2-\kappa)$, using $\norm{\collapseYt}_F^2=n(K-1)$ and $\mu\sqrt n=\kappa/K$.
\end{proof}

The bias matters. Without it the target is $Y$ itself, whose $K$ singular values all equal $\sqrt n$, and the same argument gives class means that are orthogonal rather than a simplex, and needs $m\ge K$ (\cref{ex:collapse:nobias}). The cross-entropy needs no bias, because the softmax ignores a common shift of the logits.

\section{The solution: margins, landscapes, and imbalance}\label{sec:collapse:dynamics}

Practice does not minimize a constrained problem; it runs gradient descent, often with little weight decay. Three questions remain: where gradient descent goes without regularization, whether it can get stuck, and what happens when the classes are not balanced.

\paragraph{The margin problem.} Without regularization the logits $Wh$ are homogeneous of degree two in $\theta=(W,H)$, and the cross-entropy has the exponential tail of \cref{ch:margin} in every logit gap. By the results described in that chapter for homogeneous models, once the training loss is small the direction of $\theta$ along gradient flow is governed by the max-margin problem
\begin{equation}\label{eq:collapse:margin}
  \min\ \frac12\bigl(\norm W_F^2+\norm H_F^2\bigr)\quad\text{subject to}\quad\ip{w_k-w_j,h_{k,i}}\ge1\ \ \text{for all }k,i,\ j\ne k .
\end{equation}

\begin{theorem}[The max-margin solution is neural collapse]\label{thm:collapse:margin}
Let $K\ge2$ and $m\ge K-1$. The minimum value of \eqref{eq:collapse:margin} is $\sqrt n(K-1)$, and $(W,H)$ is a minimizer if and only if NC1 holds, the means form a simplex equiangular tight frame of squared norm $\frac{K-1}{K\sqrt n}$, and $w_k=\sqrt n\,\bar h_k$ for every $k$.
\end{theorem}

\begin{proof}
The plan is to average the constraints into a single inequality for the left side of \eqref{eq:collapse:align}, and then to apply \cref{lem:collapse:align}. Let $(W,H)$ be feasible. Summing the constraints of example $(k,i)$ over $j\ne k$ gives $\ip{Kw_k-K\bar w,h_{k,i}}\ge K-1$, that is,
\begin{equation}\label{eq:collapse:sumconstr}
  \ip{w_k-\bar w,h_{k,i}}\ge\frac{K-1}K\qquad\text{for every }(k,i).
\end{equation}
Averaging over the $Kn$ examples and applying \cref{lem:collapse:align}, $\frac{K-1}K\le\frac{\norm W_F\norm H_F}{K\sqrt n}$, so $\norm W_F\norm H_F\ge\sqrt n(K-1)$, and
\[
  \frac12\bigl(\norm W_F^2+\norm H_F^2\bigr)\ge\norm W_F\norm H_F\ge\sqrt n(K-1).
\]
Now let $(W,H)$ attain the bound. Then $\norm W_F=\norm H_F$, equality holds in \cref{lem:collapse:align} with a positive left side, and every inequality \eqref{eq:collapse:sumconstr} is an equality. The lemma gives $\bar w=0$, NC1, and $w_k=c\bar h_k$ with $c=\sqrt n\norm W_F/\norm H_F=\sqrt n$. Equality in \eqref{eq:collapse:sumconstr} means that the $K-1$ constraints of each example, each at least $1$, sum to $K-1$, so each equals $1$: $c(\norm{\bar h_k}^2-\ip{\bar h_j,\bar h_k})=1$ for all $j\ne k$. Exchanging $j$ and $k$ shows that all the means have the same squared norm $r$, and then all the off-diagonal inner products equal $r-1/c$. As in the proof of \cref{thm:collapse:main}, $\sum_k\bar h_k=0$ forces $r-1/c=-r/(K-1)$, so $r=\frac{K-1}{Kc}=\frac{K-1}{K\sqrt n}$. Conversely, a configuration of this form exists because $m\ge K-1$; every constraint is active, since $\sqrt n(\norm{\bar h_k}^2-\ip{\bar h_j,\bar h_k})=\sqrt n\,r\frac K{K-1}=1$; and $\norm H_F^2=nKr=\sqrt n(K-1)=\norm W_F^2$, so its value is $\sqrt n(K-1)$.
\end{proof}

\Cref{thm:collapse:margin} identifies the global minimizers of \eqref{eq:collapse:margin}. The problem is not convex, and the theory of \cref{ch:margin} for homogeneous models gives convergence to its Karush--Kuhn--Tucker points, not to its minimizers. For the unconstrained features model, the landscape closes this gap.

\begin{theorem}[Benign landscapes; stated with a sketch]\label{thm:collapse:landscape}
\begin{enumerate}
\item \parencite{JiEtAl2022} For the cross-entropy without regularization, gradient flow on $(W,H)$ converges in direction to Karush--Kuhn--Tucker points of \eqref{eq:collapse:margin}, and for the cross-entropy the critical points of this problem that are not global minimizers are strict saddles.
\item \parencite{ZhuEtAl2021} For the cross-entropy with weight decay on $W$, $H$ and a bias, and with a feature dimension large enough relative to $K$, every critical point is either a global minimizer, which is neural collapse, or a strict saddle, at which the Hessian has a negative eigenvalue.
\item \parencite{ZhouEtAl2022} The same holds for the square loss \eqref{eq:collapse:mse}.
\end{enumerate}
\end{theorem}

\begin{proofsketch}
We describe a mechanism common to parts (ii) and (iii), and omit the case analysis, for which we refer to the three papers. The objective depends on $W$ and $H$ through $Z=WH$ and through the penalty, and by \eqref{eq:collapse:pen} and its equality case (\cref{ex:collapse:mseattain}), minimizing the penalty over factorizations of a fixed $Z$ gives $\mu\collapsenuc Z$. So the problem in $(W,H)$ is a factorized form of the convex problem $\min_Z\Lhat(Z)+\mu\collapsenuc Z$. A critical point of the factorized problem that is not a global minimizer either has an unbalanced factorization or a product $Z$ that is not optimal for the convex problem; when $m$ exceeds the rank of the optimal $Z$, there is room in $\R^m$ to move $W$ and $H$ in a new direction along which the objective decreases to second order, which is the negative curvature. For part (i), the convergence to Karush--Kuhn--Tucker points is the homogeneous-model theory of \textcite{LyuLi2020} described in \cref{ch:margin}, and \textcite{JiEtAl2022} analyze the critical points of \eqref{eq:collapse:margin} directly.
\end{proofsketch}

Strict saddles do not attract gradient descent from random initialization, so the theorem predicts that gradient methods find neural collapse in the model, and they do: in all our runs of \cref{fig:collapse:phase}, quasi-Newton descent from random starting points reached the predicted minimum to six digits.

\paragraph{Imbalanced classes.} Suppose now that $K_A$ majority classes have $n_A$ examples each and $K_B\ge2$ minority classes have $n_B$ each, with imbalance ratio $R=n_A/n_B$, $K=K_A+K_B$, and loss averaged over all examples. Following \textcite{FangEtAl2021}, constrain $\frac1K\sum_k\norm{w_k}^2\le E_W$ and $\frac1K\sum_k\frac1{n_k}\sum_i\norm{h_{k,i}}^2\le E_H$, where $n_k$ is the size of class $k$. Imbalance breaks the permutation symmetry, and with it the simplex. \Textcite{FangEtAl2021} proved that it breaks in a definite way.

\begin{theorem}[Minority collapse; \textcite{FangEtAl2021}; stated with a sketch]\label{thm:collapse:minority}
With $K_A,K_B,E_W,E_H$ fixed, let $(W^R,H^R)$ be a minimizer of the imbalanced problem for each $R$. Then $\lim_{R\to\infty}(w^R_k-w^R_{k'})=0$ for every two minority classes $k$ and $k'$.
\end{theorem}

The minority classifiers become identical, so that the network cannot tell the minority classes apart, however their features are arranged. The key step is an averaging lemma.

\begin{lemma}[Averaging the minority classifiers]\label[lemma]{lem:collapse:average}
Fix the features and the majority classifiers, and replace each minority classifier $w_j$, $j\in B$, by their mean $\bar w_B=\frac1{K_B}\sum_{j\in B}w_j$. Then the loss of every majority example does not increase, and $\sum_{j\in B}\norm{w_j}^2$ strictly decreases unless the $w_j$, $j\in B$, are already equal.
\end{lemma}

\begin{proof}
Let $(k,i)$ be an example of a majority class, with logits $z=Wh_{k,i}$. Its loss is
\[
  \ell(z,k)=\log\Bigl(\sum_{j\notin B}e^{z_j}+\sum_{j\in B}e^{\ip{w_j,h_{k,i}}}\Bigr)-z_k ,
\]
an increasing function of the sum over $B$, and the averaging changes nothing else. By convexity of the exponential, $\sum_{j\in B}e^{\ip{w_j,h}}\ge K_Be^{\ip{\bar w_B,h}}$, which is the value after averaging. The norm statement is strict convexity of $\norm\cdot^2$: $\sum_{j\in B}\norm{w_j}^2=K_B\norm{\bar w_B}^2+\sum_{j\in B}\norm{w_j-\bar w_B}^2$.
\end{proof}

\begin{heuristic}[Why the minority classifiers merge]
Take $K_A=K_B=2$ for definiteness. The minority examples carry a fraction $1/(1+R)$ of the loss, so as $R\to\infty$ the problem tends to that of minimizing the majority loss alone. In that limit the minority features cost budget and earn nothing, so they shrink to zero, and the minority classifiers matter only through the logits they give to majority examples, which they should make as small as possible. By \cref{lem:collapse:average}, merging the two minority classifiers into their mean never hurts the majority examples and frees classifier budget, which can be spent on the majority margins; so at the limiting optimum the two minority classifiers coincide, both pointing away from the majority means. What this argument does not show is that the freed budget strictly lowers the loss at every limiting optimum; \textcite{FangEtAl2021} give the full proof.
\end{heuristic}

The numerical solutions of \cref{fig:collapse:minority} show something stronger than the limit: the two minority classifiers become exactly equal at a finite imbalance ratio, which grows quickly with the budget. With the square loss and no bias the threshold can be computed: a class whose size $n_k$ satisfies $n_k\le N^2\lambda_W\lambda_H$, where $N$ is the total number of examples, gets the classifier $w_k=0$ and zero features, so all such minority classes collapse to the same point (\cref{ex:collapse:imbalance}).

\section{Reading the solution}\label{sec:collapse:reading}

The solution makes four predictions about the last layer of a classifier trained to convergence.

\emph{The geometry is universal.} With norm constraints, weight decay, or no regularization at all (in the max-margin limit), with the cross-entropy or with the square loss and a bias, the class means form a simplex, the classifier copies it, and the features collapse. Only the scale depends on the details.

\emph{A transition at $\kappa=1$.} \Cref{cor:collapse:wd} and \cref{prop:collapse:mse} have the same threshold. For $\kappa=K\sqrt{n\lambda_W\lambda_H}\ge1$ the regularization wins and the last layer is zero; for $\kappa<1$ it is a simplex whose scale grows as $\kappa$ decreases. Near the threshold both orders of the transition can be read off: $\rho^\star\approx(K-1)(1-\kappa)$ for the cross-entropy and $\rho=\frac{K-1}K(1-\kappa)$ for the square loss, and since $\rho$ is a product of two norms, $\norm W_F$ and $\norm H_F$ vanish like $(1-\kappa)^{1/2}$: a continuous transition with the mean-field exponent $\frac12$ of Landau theory.\sidenote{In Landau's theory of a continuous transition the order parameter vanishes like the square root of the distance to the critical point; here $\norm W_F^2\propto\rho$ plays the role of the squared order parameter.} With two or more linear layers on top of the features the transition becomes discontinuous (\cref{ex:collapse:deep}).

\emph{The confidence is set by $\kappa$ alone.} The probability that the cross-entropy solution assigns to the correct class and the correct-class output of the square-loss solution are the same number, $1-\frac{K-1}K\kappa$. Small weight decay gives confident training predictions, and the training loss of the cross-entropy solution, $\log\frac K{K-(K-1)\kappa}\approx\frac{K-1}K\kappa$ for small $\kappa$, is proportional to $\kappa$.

\emph{Logits grow like $\log(1/\lambda)$.} As $\kappa\to0$, $\rho^\star=\frac{K-1}K\log\frac1\kappa+O(1)$: the logits grow logarithmically in the inverse weight decay. This is the regularized counterpart of the $\log t$ growth of \cref{ch:margin}, where time plays the role of $1/\lambda$.

\begin{figure}[t]
\begin{fullwidth}
\centering
\begin{tikzpicture}
\begin{groupplot}[group style={group size=2 by 1,horizontal sep=17mm},hplot,width=0.46\linewidth,height=0.33\linewidth,
  xmin=0,xmax=1.4,xlabel={$\kappa=K\sqrt{n\lambda_W\lambda_H}$}]
\nextgroupplot[ymin=0,ymax=3,ylabel={$\rho=\norm W_F\norm H_F/(K\sqrt n)$},legend style={at={(0.97,0.97)},anchor=north east}]
\addplot[theory] table[x=kappa,y=rhoce] {figures/data/12-collapse-phase-theory.dat};
\addplot[theory2] table[x=kappa,y=rhomse] {figures/data/12-collapse-phase-theory.dat};
\addplot[simulation] table[x=kappa,y=rhoce] {figures/data/12-collapse-phase-sim.dat};
\addplot[simulation2] table[x=kappa,y=rhomse] {figures/data/12-collapse-phase-sim.dat};
\legend{cross-entropy \eqref{eq:collapse:rhostar},square loss,cross-entropy (opt.),square loss (opt.)}
\nextgroupplot[ymin=0,ymax=1.05,ylabel={correct-class probability or output},legend style={at={(0.97,0.97)},anchor=north east}]
\addplot[theory] table[x=kappa,y=p] {figures/data/12-collapse-phase-theory.dat};
\addplot[simulation] table[x=kappa,y=pce] {figures/data/12-collapse-phase-sim.dat};
\addplot[simulation2] table[x=kappa,y=pmse] {figures/data/12-collapse-phase-sim.dat};
\addplot[guide] coordinates {(1,0) (1,1.05)};
\legend{$1-\frac{K-1}K\kappa$,cross-entropy,square loss}
\end{groupplot}
\end{tikzpicture}
\caption{The transition at $\kappa=1$ in the unconstrained features model with weight decay, $K=4$, $n=5$, $m=6$, $\lambda_W=\lambda_H$. Curves: \cref{cor:collapse:wd} for the cross-entropy and \cref{prop:collapse:mse} for the square loss with a bias. Points: the regularized objective minimized numerically by L-BFGS from three random initializations (seeds $0,1,2$), keeping the best. Left: the scale $\rho$ of the last layer. Right: the probability of the correct class (cross-entropy) and the correct-class output (square loss), which the theory predicts to be equal.}\label{fig:collapse:phase}
\end{fullwidth}
\end{figure}

\Cref{fig:collapse:phase} tests the first three predictions. The points are numerical minimizers of the regularized objectives from random initializations, and they fall on the curves to the precision of the optimizer: at $\kappa\in\{0.1,0.5,0.9,0.95\}$ the discrepancy in $\rho$ is below $10^{-7}$ for both losses, the within-class scatter is below $10^{-13}$ of the size of the means, and the cosines between the centered means are within $10^{-6}$ of $-\frac13$. For $\kappa>1$ the optimizer returns $W=0$, $H=0$, as \cref{cor:collapse:wd} and \cref{prop:collapse:mse} predict. In the left panel the cross-entropy scale diverges logarithmically as $\kappa\to0$, while the square-loss scale saturates at $\frac{K-1}K$, where the outputs equal their targets. The model is exactly solvable, and a figure of this kind tests the algebra and the landscape (the optimizer was never trapped), not the model's relevance to networks.

\begin{figure}[t]
\centering
\begin{tikzpicture}
\begin{axis}[hplot,xmode=log,ymode=log,xmin=1,xmax=20000,ymin=0.04,ymax=3,
  xlabel={gradient step},ylabel={collapse metric},legend style={at={(0.97,0.97)},anchor=north east,font=\scriptsize},legend columns=2]
\addplot[sea,line width=0.9pt,mark=*,mark size=1.2pt] table[x=step,y=nc1] {figures/data/12-collapse-network.dat};
\addplot[sea,line width=0.9pt,densely dashed,mark=o,mark size=1.2pt] table[x=step,y=nc1test] {figures/data/12-collapse-network.dat};
\addplot[seadeep,line width=0.9pt,mark=square*,mark size=1.1pt] table[x=step,y=nc2] {figures/data/12-collapse-network.dat};
\addplot[sandline!55!black,line width=0.9pt,mark=triangle*,mark size=1.4pt] table[x=step,y=nc3] {figures/data/12-collapse-network.dat};
\addplot[muted,line width=0.5pt,densely dotted] coordinates {(35,0.04) (35,3)};
\legend{NC1 train,NC1 test,NC2,NC3}
\end{axis}
\end{tikzpicture}
\caption{Collapse metrics of a ReLU network $32\to256\to256\to32\to4$ trained by full-batch gradient descent (step $0.1$, momentum $0.9$, weight decay $10^{-3}$ on all parameters, cross-entropy, seed $0$) on a Gaussian mixture with $K=4$ classes in $\R^{32}$ (class means $\cN(0,\frac14I)$, unit noise, $64$ training examples per class). NC1 is $\tr(\Sigma_W\Sigma_B^+)/K$ on the training and on $2000$ test examples; NC2 is the Frobenius distance between the normalized Gram matrix of the centered means and that of a simplex; NC3 is the distance between $W/\norm W_F$ and the centered means normalized the same way. The training error reaches zero near step $35$ (dotted line). The unconstrained features model predicts that all the metrics tend to zero.}\label{fig:collapse:network}
\end{figure}

\Cref{fig:collapse:network} turns to a network. It is a small experiment, but it shows the typical shape. The training error reaches zero near step $35$; at that point NC1 has fallen from $2.3$ to $0.2$, mostly as a by-product of learning to classify. In the terminal phase the metrics decrease again, slowly: NC2 from $0.35$ to $0.09$ and NC3 from $0.32$ to $0.09$ over the next twenty thousand steps, and NC4 is exactly zero throughout. The prediction of the model is that they tend to zero. They do not, within our budget: NC1 stalls near $0.08$ after step $2000$, and the test NC1 stalls near $0.21$, more than twice the training value; the test error wanders between $7.5\%$ and $8.7\%$ during the terminal phase, ending at its highest value. The geometry of the model is the direction in which the network moves; the network stops short of it.

\begin{figure}[t]
\centering
\begin{tikzpicture}
\begin{axis}[hplot,xmode=log,xmin=1,xmax=1000,ymin=-0.45,ymax=1.5,ytick={-0.5,0,0.5,1},
  xlabel={imbalance ratio $R=n_A/n_B$},ylabel={$\cos\angle(w_3,w_4)$},legend columns=3,legend style={at={(0.5,0.97)},anchor=north,font=\scriptsize}]
\addplot[guide] coordinates {(1,-0.3333) (1000,-0.3333)};
\addplot[guide] coordinates {(1,1) (1000,1)};
\addplot[sea,line width=0.9pt,mark=*,mark size=1.5pt] table[x=R,y=cosE1] {figures/data/12-collapse-minority.dat};
\addplot[seadeep,line width=0.9pt,mark=square*,mark size=1.4pt] table[x=R,y=cosE3] {figures/data/12-collapse-minority.dat};
\addplot[sandline!55!black,line width=0.9pt,mark=triangle*,mark size=1.6pt] table[x=R,y=cosE10] {figures/data/12-collapse-minority.dat};
\legend{,,$E_W=E_H=1$,$E_W=E_H=3$,$E_W=E_H=10$}
\end{axis}
\end{tikzpicture}
\caption{Minority collapse in the norm-constrained model with $K_A=2$ majority and $K_B=2$ minority classes, $m=4$: the cosine between the two minority classifiers at the numerical optimum (SLSQP from four random starts, seeds $0$ to $3$, best kept) against the imbalance ratio. At $R=1$ the classifiers form a simplex (cosine $-\frac13$, lower dashed line); as $R$ grows they merge (cosine $1$, upper dashed line), at a ratio that grows quickly with the budget.}\label{fig:collapse:minority}
\end{figure}

\Cref{fig:collapse:minority} shows \cref{thm:collapse:minority} at work. With budget $E_W=E_H=1$ the two minority classifiers are identical from $R\approx5$ on; with budget $3$, from $R\approx300$; with budget $10$ their cosine has only risen from $-0.33$ to $-0.17$ at $R=1000$. A larger budget means a smaller effective regularization, and the merging is a consequence of the regularization: with a small budget the classifier norm is precious, and the minority classes are not worth it.

\section{What the model gets wrong}\label{sec:collapse:wrong}

\paragraph{What was thrown away.} \emph{The inputs.} The features are free, so the model cannot distinguish a training set that is easy to collapse from one that is hard, and it cannot predict how fast collapse happens or whether it happens at all for a network of given width and depth. \emph{The test set.} With no population, there is nothing to say about the features of unseen examples, which is where collapse would matter for generalization. \emph{The layers below.} A penalty on $\norm H_F$ replaces the weight decay on all the layers that produce the features; the replacement is exact for no architecture. \emph{Balance and room.} The simplex needs equal class sizes and $m\ge K-1$; both fail in common settings. \emph{The optimizer.} The theorems concern minimizers and limits; the time scale of collapse is absent.

\paragraph{Where it breaks.} \emph{Test features do not collapse.} \Textcite{HuiBelkinNakkiran2022}, in an unrefereed preprint, separate collapse on training data, which they find, from collapse on test data, which they find much weaker, and conclude that neural collapse is primarily an optimization phenomenon whose relation to generalization is unclear; in \cref{fig:collapse:network} the test NC1 stalls at more than twice the training value. \emph{Collapse can hurt transfer.} \Textcite{KornblithEtAl2021} compared nine training objectives on ImageNet and found that those producing more class separation in the penultimate layer gave higher accuracy but frozen features that transferred worse; \textcite{HuiBelkinNakkiran2022} report that training longer can likewise hurt transfer, while \textcite{GalantiGyorgyHutter2022} argue that collapse helps few-shot learning on new classes. The model, which knows only the training classes, can decide neither. \emph{Depth changes the optimum.} With several free layers between the features and the classifier, a deep unconstrained features model, the global optimum exhibits neural collapse layer after layer for binary classification \parencite{SukenikMondelliLampert2023}; but beyond two layers or two classes, solutions of lower rank have lower cost, because a product of penalized layers favors low rank \parencite{SukenikLampertMondelli2024}. \Cref{ex:collapse:deep} shows the mechanism: depth makes the penalty concave in the singular values. \emph{More classes than dimensions.} The classes of a language model are tokens, and the vocabulary exceeds the embedding dimension, so no simplex fits: the geometry is frustrated, as for many charges on a sphere. \Textcite{JiangEtAl2024} show that the minimizers of a spherically constrained model then maximize the smallest one-versus-rest margin, under conditions on $m$ and $K$; \textcite{WuPapyan2024} find that in causal language models, where the vocabulary is also imbalanced and training lasts a few epochs, collapse properties develop with scale and are linked to generalization; and \textcite{ZhaoEtAl2024} show that next-token prediction drives the logits toward a sparse-plus-low-rank structure in which contexts with the same set of possible successors collapse onto a subspace rather than a point. \emph{Imbalance.} For moderate imbalance the geometry is not a simplex: \textcite{ThrampoulidisEtAl2022} characterize it, for vanishing regularization, by the singular value decomposition of a simplex-encoded label matrix, and find that ridge regularization matters for the minority classes, as the dependence on the budget in \cref{fig:collapse:minority} suggests.

\paragraph{Fine structure.} \Textcite{TirerBruna2022} add a layer of weights and a ReLU below the free features and find that, for the regularized square loss, the minimizers can have a more delicate structure than the simplex. \Textcite{JacotEtAl2025} remove the free features: for networks ending in at least two linear layers, trained with weight decay and with a wide first layer, they prove that gradient descent reaches approximately collapsed last-layer features, under conditions that they state. For dynamics, \textcite{HanPapyanDonoho2022} observe that the classifier stays close to the least-squares classifier of the current features, and analyze a renormalized gradient flow on the resulting central path, which predicts collapse. \Textcite{EWojtowytsch2022} explain the simplex through toy models of very expressive networks, and \textcite{HeSu2023} and \textcite{RangamaniEtAl2023} document partial collapse that builds up with depth. \emph{Open:} the rate of collapse in a network of given architecture; a model with data in which test features collapse partially; and the geometry that gradient descent selects when $K>m+1$.

\section{Model card}\label{sec:collapse:card}

\modelcard{In the terminal phase of training, last-layer features collapse to their class means (NC1), the means form a simplex equiangular tight frame (NC2), the classifier aligns with them (NC3), and prediction becomes nearest class center (NC4).}{Unconstrained features: free $h_{k,i}\in\R^m$, classifier $W$, $K$ balanced classes, $m\ge K-1$; cross-entropy or square loss with norm constraints, weight decay, or no regularization (max-margin limit).}{Every minimizer is neural collapse, with closed-form scale; a transition at $\kappa=K\sqrt{n\lambda_W\lambda_H}=1$; the max-margin problem has the same solution; benign landscapes (sketched); minority collapse under imbalance (sketched).}{The simplex geometry; correct-class confidence $1-\frac{K-1}K\kappa$; logits growing like $\log(1/\kappa)$; minority classifiers merging at large imbalance, at a threshold that grows with the budget.}{Silent on test data and transfer; partial collapse in networks; deeper free models prefer lower rank; no simplex when $K>m+1$ (language models); the geometry under moderate imbalance is not a simplex.}{Deep and extended unconstrained features models; generalized neural collapse; SELI geometry; central-path dynamics; collapse in end-to-end training of wide networks.}

The chapter ends the thread on implicit bias that starts in \cref{ch:diagonal} and runs through \cref{ch:margin}. In \cref{ch:margin} the features were fixed and gradient descent found the max-margin classifier; here the features are free as well, the margin problem \eqref{eq:collapse:margin} is solved over both, and the solution is the simplex: the max-margin attractor of a model that is homogeneous of degree two. The Fourier features of \cref{ch:grokking} are another solution of a margin problem for a homogeneous network, found the same way, by symmetry. The superposition of \cref{ch:superposition} is the opposite regime: there more features must be stored than there are dimensions, and the geometry is a compromise; here there is room, and the geometry is perfect.

The unconstrained features model is a mean-field theory: it replaces the layers below by a single effective cost on the features, and it gets the symmetry of the ground state right while missing the test set and the finite training time.

\begin{exercises}
\item \consolidation (The simplex.) Let $U\in\R^{m\times K}$ have orthonormal columns, $m\ge K$, and put $M=\sqrt{\frac{rK}{K-1}}\,U(I_K-\frac1K\collapseone\collapseone^\top)$. (a) Show that the columns of $M$ form a simplex equiangular tight frame of squared norm $r$. (b) Show that $MM^\top$ is $\frac{rK}{K-1}$ times an orthogonal projection of rank $K-1$, which is the tight-frame property. (c) For $m=K-1$, construct a frame explicitly by projecting the standard basis vectors of $\R^K$ onto the hyperplane orthogonal to $\collapseone_K$ and choosing an orthonormal basis of that hyperplane. \label{ex:collapse:etf}

\item \consolidation (NC4 from NC1--NC3.) Suppose that the training features satisfy NC1--NC3 with means of equal norm centered at the origin and zero bias. Show that $\argmax_k\ip{w_k,h}=\argmin_k\norm{h-\bar h_k}$ for every $h\in\R^m$. Show by an example with $K=2$ that the conclusion fails if the means have unequal norms. \label{ex:collapse:ncc}

\item \consolidation (A bias does not help.) Add a bias $b\in\R^K$ to the logits in \cref{thm:collapse:main}, unconstrained. Show that the lower bound of the proof still holds, because $\frac1{Kn}\sum_{k,i}(\bar b-b_k)=0$ for balanced classes, and deduce that the minimum value is unchanged and that at every minimizer $b$ is a multiple of $\collapseone_K$. \label{ex:collapse:bias}

\item \consolidation (Checking \cref{cor:collapse:wd}.) (a) Show that at a minimizer with $\kappa<1$, $\norm W_F^2=K\rho^\star\sqrt{n\lambda_H/\lambda_W}$ and $\norm H_F^2=K\rho^\star\sqrt{n\lambda_W/\lambda_H}$. (b) Show that $\rho^\star=(K-1)(1-\kappa)+O((1-\kappa)^2)$ as $\kappa\uparrow1$ and $\rho^\star=\frac{K-1}K\log\frac1\kappa+\frac{K-1}K\log K+o(1)$ as $\kappa\to0$. \label{ex:collapse:wdcheck}

\item \consolidation (Attaining the square-loss bound.) Let $\kappa<1$. Take a simplex frame $\bar h_1,\dots,\bar h_K$ of squared norm $r$, put $h_{k,i}=\bar h_k$, $w_k=c\,\bar h_k$ and $b=\frac1K\collapseone$. (a) Find $c$ and $r$ such that $\lambda_W\norm W_F^2=\lambda_H\norm H_F^2$ and $WH=(1-\kappa)\collapseYt$. (b) Show that then $\norm W_F\norm H_F=\collapsenuc{WH}$, so that every inequality in \eqref{eq:collapse:pen} is an equality and $F_{\mathrm{sq}}=\Phi(Z_0)$. \label{ex:collapse:mseattain}

\item \variation (The square loss without a bias.) Drop $b$ from \eqref{eq:collapse:mse} and assume $m\ge K$. Show that the unique minimizer of the analogue of $\Phi$ is $(1-\kappa)_+Y$, and that for $\kappa<1$ every minimizer has collapsed features whose means are orthogonal with equal norms: an orthogonal frame, not a simplex. What does NC4 look like for this geometry? \label{ex:collapse:nobias}

\item \variation (Minority collapse, exactly solvable.) In the square loss without a bias, let class $k$ have $n_k$ examples, $N=\sum_kn_k$, and minimize $\frac1{2N}\norm{WH-Y}_F^2+\frac{\lambda_W}2\norm W_F^2+\frac{\lambda_H}2\norm H_F^2$ with $m\ge K$. (a) Show that $YY^\top=\diag(n_1,\dots,n_K)$ and that the analogue of $\Phi$ is minimized by $Z_0=\sum_k(1-N\mu/\sqrt{n_k})_+e_k\collapseone_{(k)}^\top$, where $\collapseone_{(k)}$ is the indicator of the examples of class $k$ and $\mu=\sqrt{\lambda_W\lambda_H}$. \hint{the singular value decomposition of $Y$ is explicit, and the dual certificate of \cref{prop:collapse:mse} is replaced by a different multiple of each singular pair.} (b) Show that a class with $n_k\le N^2\lambda_W\lambda_H$ has $w_k=0$ and $h_{k,i}=0$ at every minimizer, so that all such classes become indistinguishable. (c) With $n=(20,20,3)$ and $\lambda_W=\lambda_H=\sqrt8/43$, check that the correct-class output is $0.3675$ for the two large classes and $0$ for the small one. \label{ex:collapse:imbalance}

\item \variation\harder (Depth makes the transition discontinuous.) Replace $WH$ in \eqref{eq:collapse:mse} by $W_2W_1H$ with $W_1\in\R^{m\times m}$, penalize $\frac\lambda2(\norm{W_2}_F^2+\norm{W_1}_F^2+\norm H_F^2)$, and you may use the fact that the minimum of $\frac13(\norm A_F^2+\norm B_F^2+\norm C_F^2)$ over factorizations $Z=ABC$ with large enough inner dimensions is $\sum_i\sigma_i(Z)^{2/3}$. (a) Assuming that the optimal $Z$ shares its singular vectors with $\collapseYt$, reduce the problem to minimizing $(t-1)^2+at^{2/3}$ over $t\ge0$ for each nonzero singular value $\sqrt nt$ of $Z$, with $a=3\lambda Kn^{1/3}$. (b) Show that the minimizer jumps from $t=\frac12$ to $t=0$ at $a_c=\frac32\cdot2^{-1/3}$. (c) Compare with the continuous transition of \cref{prop:collapse:mse}, and relate to \textcite{SukenikLampertMondelli2024}. \hint{at the jump, the value at the nonzero critical point equals the value $1$ at $t=0$, and the critical point satisfies $a=3(1-t)t^{1/3}$.}\label{ex:collapse:deep}

\item \simulation (Collapse in the model under gradient descent.) Run gradient descent with step $0.5$ on the unregularized cross-entropy of the unconstrained features model with $K=4$, $n=5$, $m=6$, from small random initialization ($\cN(0,0.1^2)$ entries, seed $0$) for $2\times10^5$ steps. Record the within-class scatter $\sum_{k,i}\norm{h_{k,i}-\bar h_k}^2/\sum_k\norm{\bar h_k}^2$, the largest deviation of the cosines between centered means from $-\frac13$, and the normalized margin $\min_{k,i,j\ne k}\ip{w_k-w_j,h_{k,i}}/\bigl(\frac12(\norm W_F^2+\norm H_F^2)\bigr)$. The reader should observe the normalized margin within $0.2\%$ of the value $1/(\sqrt n(K-1))=0.1491$ of \cref{thm:collapse:margin} after $10^3$ steps, while the scatter and the cosine deviation are still about $0.02$ after $2\times10^5$ steps and decrease slowly. Explain why: the margin is flat to second order at its maximizer, so a near-optimal margin allows a visible deviation from the simplex. \label{ex:collapse:gdsim}

\item \simulation (Thresholds of minority collapse.) With the script of \cref{fig:collapse:minority}, locate the smallest ratio $R$ at which the two minority classifiers coincide, for $E_W=E_H\in\{0.5,1,2,3\}$, and plot it on a logarithmic scale. The reader should find thresholds near $R=2$, $4$, $30$ and $300$, growing roughly exponentially in $E$ over this range. Explain qualitatively why a larger budget delays the merging. \label{ex:collapse:threshold}

\item \research (Dynamics of collapse.) In a network trained past zero error, NC1 decreases slowly and may stall (\cref{fig:collapse:network}). Find a solvable model with data, for example a two-layer network on a Gaussian mixture, in which the rate of decrease of NC1 and its limiting value can be computed. \Textcite{HanPapyanDonoho2022} and \textcite{JiEtAl2022} analyze dynamics with free features; a model in which the features are produced from inputs is open, beyond the results of \textcite{JacotEtAl2025}. \label{ex:collapse:dynamics}

\item \research (More classes than dimensions.) For $K>m+1$ the minimizers of \cref{thm:collapse:main} are no longer simplices. Characterize them, for example for $m=2$ and all $K$, where a regular $K$-gon is the natural guess, and decide when gradient descent finds them. \Textcite{LuSteinerberger2022} study the limit of many points, and \textcite{JiangEtAl2024} prove one-versus-rest margin maximization under conditions; a general description is open. \label{ex:collapse:many}
\end{exercises}

\notesandrefs
Neural collapse was named and measured by \textcite{PapyanHanDonoho2020}, whose paper appeared on arXiv in August 2020 and in PNAS that October; \textcite{HanPapyanDonoho2022} extended it to the square loss, an ICLR 2022 outstanding paper, and introduced the central path. The unconstrained features model appeared within months in several groups: \textcite{MixonParshallPi2020} gave it its name and studied its loss landscape; \textcite{LuSteinerberger2022} solved the cross-entropy on the sphere and studied the limit of many classes; \textcite{EWojtowytsch2022} proved simplex symmetry in toy models of expressive networks; \textcite{FangEtAl2021} introduced the layer-peeled model, proved \cref{thm:collapse:main} in their formulation, and discovered minority collapse; and \textcite{GrafEtAl2021} proved that both the cross-entropy and the supervised contrastive loss are minimized by simplex configurations. \Textcite{ZhuEtAl2021} proved the benign landscape with weight decay and used it to fix the last layer to a simplex in practice, saving memory; \textcite{ZhouEtAl2022} did the same for the square loss; and \textcite{JiEtAl2022} connected the unregularized model to the margin theory of \cref{ch:margin}. \Textcite{Kothapalli2023} reviews this literature.

We isolated the alignment inequality of \cref{lem:collapse:align} because \cref{thm:collapse:main}, \cref{cor:collapse:wd} and \cref{thm:collapse:margin} all reduce to it. The closed form of \cref{cor:collapse:wd} in the single parameter $\kappa$, the dual-certificate proof of \cref{prop:collapse:mse}, and the confidence $1-\frac{K-1}K\kappa$ shared by both losses are our own presentation; we have not found them stated in this form.

Equiangular tight frames, and the bound of \cref{lem:collapse:simplex}, belong to frame theory and coding \parencite{StrohmerHeath2003}; the margin view of collapse comes from \textcite{SoudryEtAl2018} and \textcite{LyuLi2020}, through \textcite{JiEtAl2022}. Language models have made the regime $K>m+1$ the practically important one.

\paragraph{Translating the literature.} \Textcite{PapyanHanDonoho2020} write $C$ for our number of classes $K$, $N$ for our examples per class $n$, $p$ for our feature dimension $m$, $\boldsymbol h_{i,c}$ for our $h_{c,i}$, $\boldsymbol\mu_c$ and $\boldsymbol\mu_G$ for our $\bar h_c$ and $h_G$, and $\Sigma_W$, $\Sigma_B$ as we do; their NC1 metric is $\tr(\Sigma_W\Sigma_B^\dagger)/C$, our $\tr(\Sigma_W\Sigma_B^+)/K$. \Textcite{FangEtAl2021} use $K$, $n_k$, $\boldsymbol h_{k,i}$ and the budgets $E_W$, $E_H$ as we do, write the classifier as $\boldsymbol W=[\boldsymbol w_1,\dots,\boldsymbol w_K]^\top$, and denote the majority and minority counts by $K_A$, $K_B$ and the sizes by $n_A$, $n_B$, with ratio $R=n_A/n_B$. \Textcite{ZhuEtAl2021} and \textcite{ZhouEtAl2022} write $d$ for our feature dimension $m$ (we keep $d$ for the input dimension), $\boldsymbol H\in\R^{d\times N}$ with $N=nK$ as we do, and $\lambda_{\boldsymbol W},\lambda_{\boldsymbol H},\lambda_{\boldsymbol b}$ for the weight decays; our $\kappa$ is not used there, and their thresholds can be rewritten in terms of it. Several papers, including \textcite{FangEtAl2021}, call the configuration of \cref{def:collapse:etf} a simplex ETF and normalize it with $\sqrt{K/(K-1)}$ as in \cref{ex:collapse:etf}.

\printchapterbib
\end{refsection}


\casepart{Transformers}{Two solvable models of learning in context: a linear attention layer that performs one step of gradient descent, and a two-layer circuit that counts bigrams. Each is exact in its own world and partial in ours.}{74}
% Chapter 13 of The Hydrogen Atoms of Deep Learning: in-context learning by linear attention.
\newcommand{\iclKQ}{W^{KQ}}
\newcommand{\iclPV}{W^{PV}}
\newcommand{\iclShat}{\hat\Sigma}
\newcommand{\iclc}{\hat c}

\casechapter{Learning in Context}{Linear attention as one step of gradient descent}{Linear self-attention on regression tasks}{ch:icl}
\begin{refsection}[chapters/13-icl.bib]

\begin{chaptersummary}
Large language models learn a task from a few examples placed in the prompt, without changing a weight, and transformers trained on prompts of linear regression problems predict nearly as well as least squares. A single layer of linear self-attention, trained on such prompts, is the solvable model: its prediction is one step of preconditioned gradient descent on the prompt, and we prove that the best such layer, among all its parameters, takes the step $\Gamma^\star=N\bigl((N+1)\Sigma+(\tr\Sigma+\sigma^2)I_d\bigr)^{-1}$, that gradient flow finds it, and that for isotropic inputs and a small balanced initialization it does so after a plateau, along a logistic curve. The model gets the mechanism of one trained layer right and the rate wrong: one step loses a factor $(d+1+\sigma^2)/\sigma^2$ against least squares, deeper trained transformers do much better, softmax attention computes different estimators, and in language models in-context learning is in part the recognition of a task already learned.
\end{chaptersummary}

\section{The phenomenon: learning from the prompt}\label{sec:icl:phenomenon}

A trained network is a fixed function. Yet \textcite{BrownEtAl2020} found that GPT-3, a transformer language model with 175 billion parameters, could be taught a new task by its input alone: shown a few demonstrations of a task, written as text (pairs of English and French words, or arithmetic problems with their answers), followed by a new instance, it often completed the instance correctly, and more demonstrations usually helped. They called this \term{in-context learning}\index{in-context learning}, distinguished zero-, one- and few-shot prompts, and stressed that no gradient update takes place: whatever learning happens is computed by the forward pass, from the prompt.\sidenote{The idea that a network can learn to learn inside its activations is older: \textcite{HochreiterEtAl2001} trained LSTM networks whose recurrent dynamics, with fixed weights, learned new functions from the examples presented to them.}

Language is a poor laboratory: the tasks are vague, and one cannot tell whether a model learns from the demonstrations or recognizes a task it has seen in training. \textcite{GargEtAl2022} replaced language with a well-posed problem. Each training sequence is a prompt $(x_1,f(x_1),\dots,x_N,f(x_N),x_q)$ with a function $f$ drawn afresh from a class, and the model is trained to predict $f(x_q)$. For linear functions $f(x)=\ip{w_\star,x}$ with $x$ and $w_\star$ isotropic Gaussian in dimension $20$, a transformer trained from scratch reached an error comparable to that of least squares on the same prompt, and it did so on functions it had never seen. Sparse linear functions, decision trees and two-layer networks could be learned in context too.

What does the trained transformer compute? \textcite{AkyurekEtAl2023} showed by construction that transformers can implement gradient descent and closed-form ridge regression, and found that trained models track the predictions of least squares and ridge. \textcite{VonOswaldEtAl2023} showed that one layer of linear self-attention can implement exactly one step of gradient descent on the squared loss of the prompt, and found that a trained single layer matches that step closely, in prediction and in its weights. \textcite{RaventosEtAl2023} varied the number of distinct tasks $w_\star$ seen in training and found a threshold: below it, the transformer behaves like the Bayes estimator for the finite set of training tasks and fails on new ones; above it, it behaves like ridge regression, the Bayes estimator for a Gaussian prior over all tasks.

\begin{observation}[\textcite{GargEtAl2022}; \textcite{VonOswaldEtAl2023}; \textcite{AkyurekEtAl2023}]\label[observation]{obs:icl:icl}
Transformers trained on prompts of noisy linear regression problems with random coefficients, to predict the label of a query input from the labeled examples before it, learn to predict nearly as well as least squares or ridge regression on the prompt, for coefficients never seen in training. A single trained layer of linear self-attention predicts like one step of gradient descent on the squared loss of the prompt.
\end{observation}

How robust is this? On synthetic regression the observation has been reproduced many times, but it describes models trained on nothing else, with many distinct tasks. In language models the evidence is mixed. \textcite{MinEtAl2022} found that replacing the labels of the demonstrations by random ones barely hurts many classification tasks, and \textcite{PanEtAl2023} separated \emph{task recognition}, which works without correct labels and does not improve much with scale, from \emph{task learning}, which emerges with scale; \textcite{WeiEtAl2023} (a preprint) found that large models, unlike small ones, can follow demonstrations whose labels are flipped. In-context learning also depends on the training data \parencite{ChanEtAl2022} and can be transient, appearing and then fading during training \parencite{SinghEtAl2023}.

The question for this chapter: when a single attention layer is trained on regression prompts, which algorithm does it learn, how well does that algorithm predict, and how does training find it?

\section{The model: one layer of linear attention}\label{sec:icl:model}

\paragraph{Prompts.} A task is a coefficient vector $w_\star\in\R^d$; a prompt shows $N$ labeled examples of the task and one query.

\begin{assumption}[Linear regression prompts]\label[assumption]{ass:icl:prompts}
For each prompt, independently: the task is $w_\star\sim\cN(0,I_d)$; the inputs $x_1,\dots,x_N,x_q$ are independent with law $\cN(0,\Sigma)$, $\Sigma$ positive definite; the labels are $y_i=\ip{w_\star,x_i}+\eps_i$ and $y_q=\ip{w_\star,x_q}+\eps_q$ with noises $\eps_1,\dots,\eps_N,\eps_q\sim\cN(0,\sigma^2)$; all these variables are independent.
\end{assumption}

The prompt is stored as a matrix whose columns are tokens, one per example, with the query's label left blank:
\begin{equation}\label{eq:icl:Z}
  Z=\begin{pmatrix} x_1 & x_2 & \cdots & x_N & x_q\\ y_1 & y_2 & \cdots & y_N & 0\end{pmatrix}\in\R^{(d+1)\times(N+1)} .
\end{equation}
We write $z_i=(x_i,y_i)$ for the $i$th column and $z_q=(x_q,0)$ for the last. Three statistics of the prompt will appear everywhere: the sample covariance, the label-weighted input mean, and the mean square label,
\[
  \iclShat=\frac1N\sum_{i=1}^Nx_ix_i^\top,\qquad s=\frac1N\sum_{i=1}^Ny_ix_i,\qquad \iclc=\frac1N\sum_{i=1}^Ny_i^2 .
\]
Their relevance is no accident: the empirical loss of the prompt, $\Lhat(w)=\frac1{2N}\sum_i(y_i-\ip{w,x_i})^2$, has gradient $\nabla\Lhat(w)=\iclShat w-s$, so $\iclShat$ and $s$ are the sufficient statistics of least squares, and $\nabla\Lhat(0)=-s$.

\paragraph{Linear self-attention.} A softmax attention head \parencite{VaswaniEtAl2017} maps $Z$ to $Z+W^PW^VZ\,\softmax\bigl((W^KZ)^\top W^QZ\bigr)$, the softmax acting on each column. Removing the softmax and merging the products of matrices that appear only together gives \term{linear self-attention}\index{linear self-attention} \parencite{SchlagEtAl2021,VonOswaldEtAl2023}, with parameters $\theta=(\iclKQ,\iclPV)$, two $(d+1)\times(d+1)$ matrices:
\begin{equation}\label{eq:icl:lsa}
  f(Z;\theta)=Z+\frac1N\,\iclPV Z\,M\,Z^\top\iclKQ Z,\qquad M=\begin{pmatrix}I_N&0\\0&0\end{pmatrix}\in\R^{(N+1)\times(N+1)} .
\end{equation}
The mask $M$ lets every token attend to the $N$ examples but not to the query, whose blank label carries no information. The prediction is read from the bottom entry of the query column, $\hat y_q=f(Z;\theta)_{d+1,N+1}$; since $Z_{d+1,N+1}=0$, the residual term does not contribute. The model is trained on the \term{pretraining loss}\index{pretraining loss}, the population risk over prompts,
\[
  \Loss(\theta)=\E\bigl(y_q-\hat y_q(\theta)\bigr)^2 ,
\]
with the expectation over \cref{ass:icl:prompts}. A good parameter is a learning algorithm: it must predict well for every task, using only the prompt.

\paragraph{Blocks.} Only some blocks of $\theta$ matter. Write
\[
  \iclKQ=\begin{pmatrix}W_{11}&w_{12}\\ w_{21}^\top&w_{22}\end{pmatrix},\qquad
  \iclPV=\begin{pmatrix}V_{11}&v_{12}\\ v_{21}^\top&v_{22}\end{pmatrix},
\]
with $W_{11},V_{11}\in\R^{d\times d}$, column vectors $w_{12},w_{21},v_{12},v_{21}\in\R^d$, and scalars $w_{22},v_{22}$. The prediction is the last row of $\iclPV$ applied to the last column of $\frac1NZMZ^\top\iclKQ Z$, and $\frac1NZMZ^\top=\frac1N\sum_{i\le N}z_iz_i^\top$ has blocks $\iclShat$, $s$, $s^\top$, $\iclc$. So $V_{11}$, $v_{12}$ (the rows that write into the inputs) and $w_{12}$, $w_{22}$ (which multiply the blank label) never affect $\hat y_q$; they matter only in deeper models.

\begin{dictionary}
prompt of $N$ pairs $(x_i,y_i)$ and a query & few-shot prompt: demonstrations followed by a new input\\
task $w_\star$, drawn afresh for each prompt & the task the demonstrations specify\\
law of $w_\star$ & the diversity of tasks implicit in the pretraining corpus\\
one token $z_i=(x_i,y_i)$ per example & an input and its label in separate tokens, interleaved\\
$\iclKQ$, $\iclPV$ & key--query and value--output products of one head\\
linear attention, one layer & softmax attention, many heads and layers, MLPs, normalization\\
preconditioner $\Gamma$ & the learning rule implemented by the forward pass\\
pretraining loss $\Loss(\theta)$ & next-token loss on documents\\
Bayes estimator (ridge) & the best a learner could do with the prompt\\
\end{dictionary}

\paragraph{What was thrown away, and why first.} The softmax goes, because without it attention is a cubic polynomial in the prompt and every expectation is a Gaussian moment. Depth, MLPs and normalization go, because one layer already learns a nontrivial algorithm. Language goes, replaced by a task family where the optimal learner is known, so that ``learning in context'' can be measured against ridge regression. Each example is one token, rather than two, which spares the model the job of pairing inputs with labels, the job of the induction heads of \cref{ch:induction}. \Cref{sec:icl:wrong} says what each choice costs.

\paragraph{What one had to notice.} Attention computes $\sum_iz_iz_i^\top$, a sum of outer products over the prompt, and so does the gradient of the squared loss: $-\nabla\Lhat(0)=\frac1N\sum_iy_ix_i$ is a block of $\frac1N\sum_iz_iz_i^\top$. A layer whose update is a sum over the examples can therefore take a gradient step on them.\sidenote{\Textcite{SchlagEtAl2021} had shown that linear attention is a \emph{fast weight programmer}: it writes the outer products $v_ik_i^\top$ of values and keys into a weight matrix, which then multiplies the query.}

\section{The solution: the optimal single layer}\label{sec:icl:solution}

The first result reads off what one layer computes; the second finds the best learner of that form; the third shows that nothing in the full parameter space does better.

\begin{proposition}[What one layer computes]\label[proposition]{prop:icl:express}
For every $\theta$, with $\Gamma(\theta)=v_{22}W_{11}^\top+w_{21}v_{21}^\top$ and $b(\theta)=W_{11}^\top\iclShat\,v_{21}+\iclc\,v_{22}w_{21}$,
\begin{equation}\label{eq:icl:predict}
  \hat y_q=x_q^\top\Gamma(\theta)\,s+x_q^\top b(\theta) .
\end{equation}
In particular, if $v_{21}=0$ and $w_{21}=0$, then $\hat y_q=x_q^\top\Gamma s$ with $\Gamma=v_{22}W_{11}^\top$, which is the prediction $\ip{w_1,x_q}$ of one step of preconditioned gradient descent $w_1=w_0-\Gamma\nabla\Lhat(w_0)$ from $w_0=0$. Every $\Gamma\in\R^{d\times d}$ arises this way, for instance with $v_{22}=1$ and $W_{11}=\Gamma^\top$.
\end{proposition}

\begin{proof}
The query column of $Z^\top\iclKQ Z$ is $Z^\top\iclKQ z_q$. The last entry of $z_q$ is $0$, so $\iclKQ z_q$ has blocks $W_{11}x_q$ and $w_{21}^\top x_q$. Applying $\frac1NZMZ^\top$ and then the last row $(v_{21}^\top,v_{22})$ of $\iclPV$ gives
\begin{align*}
  \hat y_q&=(v_{21}^\top,\ v_{22})\begin{pmatrix}\iclShat&s\\ s^\top&\iclc\end{pmatrix}\begin{pmatrix}W_{11}x_q\\ w_{21}^\top x_q\end{pmatrix}\\
  &=v_{21}^\top\iclShat W_{11}x_q+(v_{21}^\top s)(w_{21}^\top x_q)+v_{22}\,s^\top W_{11}x_q+v_{22}\,\iclc\,w_{21}^\top x_q .
\end{align*}
Each term is a scalar; transposing the first and third and regrouping gives \eqref{eq:icl:predict}. If $v_{21}=w_{21}=0$, then $b=0$ and $\Gamma=v_{22}W_{11}^\top$, and $w_1=0-\Gamma(\iclShat\cdot0-s)=\Gamma s$.
\end{proof}

The odd and even parts of \eqref{eq:icl:predict} will be the key. Flipping the sign of the task and of all the noises maps $y_i\mapsto-y_i$, so $s\mapsto-s$ while $\iclShat$ and $\iclc$ are unchanged: the first term of $\hat y_q$ is odd in the labels and the second, $x_q^\top b$, is even. The target $y_q$ is odd. An even term can only add error, as the proof of \cref{thm:icl:global} will show.

The risk of the one-step learner needs the fourth moments of a Gaussian vector.

\begin{lemma}[Gaussian fourth moments]\label[lemma]{lem:icl:wick}
Let $x\sim\cN(0,\Sigma)$ in $\R^d$. For every $A\in\R^{d\times d}$,
\[
  \E\bigl[xx^\top Axx^\top\bigr]=\Sigma A\Sigma+\Sigma A^\top\Sigma+\tr(A\Sigma)\,\Sigma .
\]
Consequently, if $x_1,\dots,x_N$ are independent $\cN(0,\Sigma)$ and $\iclShat=\frac1N\sum_ix_ix_i^\top$, then $\E[\iclShat^2]=\frac1N\bigl((N+1)\Sigma^2+\tr(\Sigma)\Sigma\bigr)$.
\end{lemma}

\begin{proof}
\emph{The standard case.} Let $g\sim\cN(0,I_d)$. Its coordinates are independent with $\E g_i^2=1$, $\E g_i^4=3$, and vanishing odd moments, so $\E[g_ig_jg_kg_l]$ is $1$ if the indices form two equal pairs and differ across pairs, $3$ if all four are equal, and $0$ otherwise; in every case $\E[g_ig_jg_kg_l]=\delta_{ij}\delta_{kl}+\delta_{ik}\delta_{jl}+\delta_{il}\delta_{jk}$. For $B\in\R^{d\times d}$, the $(i,l)$ entry of $gg^\top Bgg^\top$ is $\sum_{j,k}g_ig_jB_{jk}g_kg_l$, so
\[
  \E\bigl[gg^\top Bgg^\top\bigr]_{il}=\sum_{j,k}B_{jk}\bigl(\delta_{ij}\delta_{kl}+\delta_{ik}\delta_{jl}+\delta_{il}\delta_{jk}\bigr)=B_{il}+B_{li}+\delta_{il}\tr B .
\]
\emph{General $\Sigma$.} Write $x=\Sigma^{1/2}g$ and $B=\Sigma^{1/2}A\Sigma^{1/2}$. Then $xx^\top Axx^\top=\Sigma^{1/2}gg^\top Bgg^\top\Sigma^{1/2}$, whose mean is $\Sigma^{1/2}(B+B^\top+\tr(B)I)\Sigma^{1/2}=\Sigma A\Sigma+\Sigma A^\top\Sigma+\tr(A\Sigma)\Sigma$.

\emph{The sample covariance.} Expanding $\iclShat^2=\frac1{N^2}\sum_{i,j}x_ix_i^\top x_jx_j^\top$, the $N(N-1)$ terms with $i\ne j$ have mean $\Sigma^2$ by independence, and the $N$ terms with $i=j$ have mean $2\Sigma^2+\tr(\Sigma)\Sigma$ by the first part with $A=I_d$. Adding, $\E[\iclShat^2]=\frac1{N^2}\bigl(N(N+1)\Sigma^2+N\tr(\Sigma)\Sigma\bigr)$.
\end{proof}

\begin{theorem}[The optimal one-step learner]\label{thm:icl:onestep}
Under \cref{ass:icl:prompts}, let $\cR(\Gamma)=\E\bigl(y_q-x_q^\top\Gamma s\bigr)^2$ for $\Gamma\in\R^{d\times d}$, and let
\[
  H=\frac1N\Bigl((N+1)\Sigma^2+(\tr\Sigma+\sigma^2)\Sigma\Bigr).
\]
Then
\begin{equation}\label{eq:icl:risk}
  \cR(\Gamma)=\sigma^2+\tr\Sigma-2\tr(\Sigma\Gamma\Sigma)+\tr(\Sigma\Gamma H\Gamma^\top),
\end{equation}
and $\cR$ has the unique minimizer
\begin{equation}\label{eq:icl:gammastar}
  \Gamma^\star=N\Bigl((N+1)\Sigma+(\tr\Sigma+\sigma^2)I_d\Bigr)^{-1},
\end{equation}
with $\cR(\Gamma)-\cR(\Gamma^\star)=\tr\bigl(\Sigma(\Gamma-\Gamma^\star)H(\Gamma-\Gamma^\star)^\top\bigr)$. If $\Sigma$ has eigenvalues $\lambda_1,\dots,\lambda_d$, the minimal risk is
\begin{equation}\label{eq:icl:riskstar}
  \cR^\star=\sigma^2+\sum_{k=1}^d\frac{\lambda_k(\lambda_k+\tr\Sigma+\sigma^2)}{(N+1)\lambda_k+\tr\Sigma+\sigma^2} .
\end{equation}
For $\Sigma=I_d$: $\Gamma^\star=\gamma^\star I_d$ with $\gamma^\star=\frac{N}{N+d+1+\sigma^2}$, and $\cR^\star=\sigma^2+\frac{d(d+1+\sigma^2)}{N+d+1+\sigma^2}$.
\end{theorem}

\begin{proof}
The plan is to average over the query first, then over the prompt with \cref{lem:icl:wick}, and to complete the square.

\emph{The query.} The pair $(x_q,\eps_q)$ is independent of $(w_\star,s)$, with $\E x_qx_q^\top=\Sigma$ and $\E\eps_q^2=\sigma^2$. Since $y_q-x_q^\top\Gamma s=\eps_q+x_q^\top(w_\star-\Gamma s)$ is a sum of two uncorrelated terms given $(w_\star,s)$,
\[
  \cR(\Gamma)=\sigma^2+\E\bigl[(w_\star-\Gamma s)^\top\Sigma(w_\star-\Gamma s)\bigr].
\]
\emph{Two moments of $s$.} Write $s=\iclShat w_\star+\bar\xi$ with $\bar\xi=\frac1N\sum_i\eps_ix_i$. The task is independent of the inputs and noises, with $\E w_\star w_\star^\top=I_d$, and $\E[\bar\xi\mid X]=0$, so $\E[sw_\star^\top]=\E[\iclShat]=\Sigma$. For the second moment, the cross terms vanish because the noises have mean zero given everything else; conditioning on the inputs gives $\E[\iclShat w_\star w_\star^\top\iclShat]=\E[\iclShat^2]$; and independence of the noises gives $\E[\bar\xi\bar\xi^\top]=\frac{\sigma^2}{N^2}\sum_i\E[x_ix_i^\top]=\frac{\sigma^2}N\Sigma$. With \cref{lem:icl:wick},
\[
  \E[ss^\top]=\frac1N\bigl((N+1)\Sigma^2+\tr(\Sigma)\Sigma\bigr)+\frac{\sigma^2}N\Sigma=H .
\]
\emph{The risk.} Expanding the quadratic form, $\E[w_\star^\top\Sigma w_\star]=\tr\Sigma$, $\E[w_\star^\top\Sigma\Gamma s]=\tr(\Sigma\Gamma\,\E[sw_\star^\top])=\tr(\Sigma\Gamma\Sigma)$, and $\E[s^\top\Gamma^\top\Sigma\Gamma s]=\tr(\Sigma\Gamma H\Gamma^\top)$, which is \eqref{eq:icl:risk}.

\emph{Completing the square.} Let $K=\frac1N\bigl((N+1)\Sigma+(\tr\Sigma+\sigma^2)I_d\bigr)$, so that $H=\Sigma K=K\Sigma$ and $\Gamma^\star=K^{-1}$, a symmetric matrix with $H\Gamma^\star=\Sigma$. Expanding $\tr\bigl(\Sigma(\Gamma-\Gamma^\star)H(\Gamma-\Gamma^\star)^\top\bigr)$ and using $H\Gamma^{\star\top}=\Sigma$ turns its cross term into $-2\tr(\Sigma\Gamma\Sigma)$ and its last term into $\tr(\Sigma\Gamma^\star\Sigma)$; comparing with \eqref{eq:icl:risk},
\[
  \cR(\Gamma)=\sigma^2+\tr\Sigma-\tr(\Sigma\Gamma^\star\Sigma)+\tr\bigl(\Sigma(\Gamma-\Gamma^\star)H(\Gamma-\Gamma^\star)^\top\bigr).
\]
The matrices $\Sigma$ and $H$ are positive definite ($H$ is a product of two commuting positive definite matrices), so the last term equals $\norm{\Sigma^{1/2}(\Gamma-\Gamma^\star)H^{1/2}}_F^2$, which vanishes only at $\Gamma=\Gamma^\star$. Hence $\Gamma^\star$ is the unique minimizer and $\cR^\star=\sigma^2+\tr\Sigma-\tr(\Sigma^2\Gamma^\star)$. In an eigenbasis of $\Sigma$, the matrix $\Gamma^\star$ is diagonal with entries $N/\bigl((N+1)\lambda_k+\tr\Sigma+\sigma^2\bigr)$, and $\lambda_k-N\lambda_k^2/\bigl((N+1)\lambda_k+\tr\Sigma+\sigma^2\bigr)$ simplifies to the $k$th term of \eqref{eq:icl:riskstar}. For $\Sigma=I_d$, set $\lambda_k=1$ and $\tr\Sigma=d$.
\end{proof}

The proof used only $\E w_\star w_\star^\top=I_d$, independence, and Gaussian inputs: the law of the task enters only through its second moment.\sidenote{We checked \eqref{eq:icl:risk} and \eqref{eq:icl:riskstar} by Monte Carlo with $5\times10^5$ prompts for $d=5$, $N=12$, $\sigma^2=0.7$, and $\Sigma=\diag(2,1,0.5,0.3,0.2)$, at $\Gamma^\star$ and at a random $\Gamma$; the deviations were within $2.5$ standard errors.} The one-step learner is restricted, but the restriction is not ours: the full layer can do no better.

\begin{theorem}[The global minimum of one layer]\label{thm:icl:global}
Under \cref{ass:icl:prompts}, every $\theta=(\iclKQ,\iclPV)$ satisfies
\[
  \Loss(\theta)=\cR\bigl(\Gamma(\theta)\bigr)+\E\bigl(x_q^\top b(\theta)\bigr)^2\ \ge\ \cR^\star ,
\]
with $\Gamma(\theta)$ and $b(\theta)$ as in \cref{prop:icl:express}. Equality holds if and only if $\Gamma(\theta)=\Gamma^\star$ and $b(\theta)=0$ almost surely. In particular, $\min_\theta\Loss(\theta)=\cR^\star$, and every global minimizer predicts $\hat y_q=x_q^\top\Gamma^\star s$.
\end{theorem}

\begin{proof}
The plan is to split the prediction \eqref{eq:icl:predict} into its odd and even parts and to show that they do not interact.

\emph{The symmetry.} Let $\tau$ map $(w_\star,\eps_1,\dots,\eps_N,\eps_q)$ to its negative and leave the inputs unchanged. Under \cref{ass:icl:prompts} these variables are independent of the inputs and have centered Gaussian laws, so $\tau$ preserves the joint law of all the variables. It maps every label to its negative, hence $s\mapsto-s$ and $y_q\mapsto-y_q$, while $\iclShat$, $\iclc$, and $x_q$ are unchanged.

\emph{No interaction.} Write $\hat y_q=x_q^\top\Gamma s+e$ with $e=x_q^\top b(\theta)$, and let $F=(y_q-x_q^\top\Gamma s)\,e$. All the variables are polynomials in Gaussian variables, so $F$ is integrable. By the symmetry, $F\circ\tau=-F$ and $F\circ\tau$ has the law of $F$, so $\E F=-\E F=0$. Expanding the square,
\[
  \Loss(\theta)=\E\bigl(y_q-x_q^\top\Gamma s\bigr)^2-2\E F+\E e^2=\cR(\Gamma)+\E e^2 .
\]
\emph{The bound.} \Cref{thm:icl:onestep} gives $\cR(\Gamma)\ge\cR^\star$, with equality only at $\Gamma=\Gamma^\star$, and $\E e^2\ge0$, with equality only if $e=0$ almost surely. Since $x_q$ is independent of $b(\theta)$ and has a positive definite covariance, $\E e^2=\E\bigl[b^\top\Sigma b\bigr]$ vanishes exactly when $b(\theta)=0$ almost surely. The parameter with $v_{22}=1$, $W_{11}=\Gamma^{\star\top}$, and all other blocks zero attains $\cR^\star$ by \cref{prop:icl:express}.
\end{proof}

The theorem says more than that the optimum is attained: \emph{whatever} global minimum training reaches, the trained layer runs one step of gradient descent with preconditioner $\Gamma^\star$. The minimizers themselves are not unique (rescaling $v_{22}$ by $c\neq0$ and $W_{11}$ by $1/c$ changes nothing), but their predictions are.

How good is this learner? The best predictor of $y_q$ from the prompt and $x_q$ is the posterior mean, which, for Gaussian prior and noise, is ridge regression.

\begin{corollary}[One step against the Bayes estimator]\label[corollary]{cor:icl:bayes}
Let $\Sigma=I_d$ in \cref{ass:icl:prompts}, and let $X\in\R^{N\times d}$ have rows $x_i$. Then:
\begin{enumerate}
\item the Bayes-optimal prediction is $x_q^\top\hat w$ with $\hat w=(X^\top X+\sigma^2I_d)^{-1}X^\top y$, and its risk is $\sigma^2+\E\tr\bigl(\sigma^2(X^\top X+\sigma^2I_d)^{-1}\bigr)$;
\item if $d,N\to\infty$ with $d/N\to\gamma\in(0,\infty)$ and $\sigma^2/d\to\varsigma>0$, then the excess risks per dimension converge:
\begin{equation}\label{eq:icl:proportional}
  \frac{\cR^\star-\sigma^2}d\to\frac{\gamma(1+\varsigma)}{1+\gamma(1+\varsigma)},\qquad
  \frac{\cR_{\mathrm{B}}-\sigma^2}d\to\frac{\sqrt{b^2+4\varsigma}-b}2,
\end{equation}
where $\cR_{\mathrm{B}}$ is the Bayes risk and $b=\frac1\gamma-1+\varsigma$;
\item as $\gamma\to0$ (many examples per dimension), the ratio of the two limits tends to $(1+\varsigma)/\varsigma$.
\end{enumerate}
\end{corollary}

\begin{proof}
\emph{Part (i).} Given $(X,y)$, the density of $w_\star$ at $w$ is proportional to
\[
  \exp\Bigl(-\frac1{2\sigma^2}\norm{y-Xw}^2-\frac12\norm w^2\Bigr),
\]
a Gaussian whose precision matrix is $(X^\top X+\sigma^2I_d)/\sigma^2$ and whose mean $\hat w$ is the stationary point $X^\top(y-Xw)=\sigma^2w$. The squared loss is minimized by the conditional mean of $y_q$ given the prompt and $x_q$, which is $x_q^\top\hat w$, because $x_q$ and $\eps_q$ are independent of the rest. Its risk is $\sigma^2+\E\norm{w_\star-\hat w}^2$, and the conditional mean of $\norm{w_\star-\hat w}^2$ is the trace of the posterior covariance $\sigma^2(X^\top X+\sigma^2I_d)^{-1}$.

\emph{Part (ii).} For the one-step learner, \cref{thm:icl:onestep} gives $(\cR^\star-\sigma^2)/d=(d+1+\sigma^2)/(N+d+1+\sigma^2)$; divide numerator and denominator by $N$. For the Bayes estimator, write $\hat\Sigma_X=X^\top X/N$ and $\lambda=\sigma^2/N$; then $\frac1d\tr\bigl(\sigma^2(X^\top X+\sigma^2I)^{-1}\bigr)=\frac1d\tr\bigl(\lambda(\hat\Sigma_X+\lambda I)^{-1}\bigr)$, which is the function $v_p(\lambda)$ of \cref{ch:doubledescent} with $p=d$ features and $n=N$ samples. Here $\lambda=\lambda_N=\varsigma_N\,d/N\to\varsigma\gamma$ with $\varsigma_N=\sigma^2/d$. The proof of \cref{thm:doubledescent:ridge} shows that $\E v_p(\lambda)\to v(\lambda)$ at each fixed $\lambda>0$; since $\abs{\partial_\lambda v_p(\lambda)}\le1/(4\lambda)$, the convergence also holds along $\lambda_N\to\varsigma\gamma>0$. The limit $v(\varsigma\gamma)$ is the positive root of $\gamma v^2+(1-\gamma+\varsigma\gamma)v-\varsigma\gamma=0$, that is, of $v^2+bv-\varsigma=0$, which is the stated limit.

\emph{Part (iii).} As $\gamma\to0$, the first limit is $\gamma(1+\varsigma)+O(\gamma^2)$, and the second is $\frac{2\varsigma}{\sqrt{b^2+4\varsigma}+b}=\varsigma\gamma+O(\gamma^2)$, because $b=1/\gamma+O(1)$.
\end{proof}

\section{The solution: how training finds the step}\label{sec:icl:dynamics}

\Cref{thm:icl:global} is about the landscape; does gradient descent reach its minimum? Two facts make the dynamics solvable. First, the set
\[
  \cM=\{\theta:\ v_{21}=0,\ w_{21}=0\}
\]
is invariant: on it the odd--even symmetry kills the gradient in the off-manifold directions. Second, on $\cM$ the model is a two-layer linear network, $\Gamma=v_{22}W_{11}^\top$, a scalar times a matrix, with the conservation law of \cref{ch:deeplinear}. Gradient flow is $\dot\theta=-\nabla\Loss(\theta)$.

\begin{theorem}[Gradient flow finds the optimal step]\label{thm:icl:flow}
Under \cref{ass:icl:prompts}, let $\theta(0)\in\cM$ with $v_{22}(0)>\norm{W_{11}(0)}_F$. Then the gradient flow exists for all $t\ge0$, stays in $\cM$, and conserves $c=v_{22}^2-\norm{W_{11}}_F^2>0$; and as $t\to\infty$, $\Gamma(\theta(t))\to\Gamma^\star$ and $\Loss(\theta(t))\to\cR^\star$.
\end{theorem}

\begin{proof}
The plan is to show invariance by the odd--even symmetry, to obtain the conservation law and boundedness on $\cM$, and then to identify every limit point as a global minimum.

\emph{Invariance.} The loss does not depend on $V_{11}$, $v_{12}$, $w_{12}$, $w_{22}$, so their gradients vanish everywhere. By \cref{prop:icl:express}, $\partial\hat y_q/\partial v_{21}=\iclShat W_{11}x_q+(w_{21}^\top x_q)s$ and $\partial\hat y_q/\partial w_{21}=(v_{21}^\top s+v_{22}\iclc)x_q$. At a point of $\cM$ these are $\iclShat W_{11}x_q$ and $v_{22}\iclc\,x_q$, which are even under the map $\tau$ of the proof of \cref{thm:icl:global}, while $y_q-\hat y_q=y_q-x_q^\top\Gamma s$ is odd. So $\partial\Loss/\partial v_{21}=-2\E[(y_q-\hat y_q)\,\partial\hat y_q/\partial v_{21}]=0$ on $\cM$, and likewise for $w_{21}$. The vector field is polynomial, hence locally Lipschitz: the solution of the flow restricted to $\cM$, in the variables $(W_{11},v_{22})$, solves the full equation, and by uniqueness it is the full solution, as long as it exists.

\emph{The flow on $\cM$.} On $\cM$, $\Loss(\theta)=\cR(\Gamma)$ with $\Gamma=v_{22}A$ and $A=W_{11}^\top$. Let $G=\nabla\cR(\Gamma)=2(\Sigma\Gamma H-\Sigma^2)$, from \eqref{eq:icl:risk}. The chain rule gives
\[
  \dot A=-v_{22}\,G,\qquad \dot v_{22}=-\ip{G,A}_F ,
\]
so $\frac{\dd}{\dd t}v_{22}^2=-2v_{22}\ip{G,A}_F=\frac{\dd}{\dd t}\norm A_F^2$, and $c=v_{22}^2-\norm A_F^2$ is constant. In particular $v_{22}^2\ge c>0$, so $v_{22}$ never vanishes and stays positive.

\emph{Boundedness.} The loss does not increase along the flow, and by \cref{thm:icl:onestep}, $\cR(\Gamma)\ge\cR^\star+\mu\norm{\Gamma-\Gamma^\star}_F^2$ with $\mu=\lmin(\Sigma)\lmin(H)>0$. So $\norm{\Gamma(t)}_F\le B$ for a constant $B$. Since $\norm\Gamma_F^2=v_{22}^2\norm A_F^2=(c+\norm A_F^2)\norm A_F^2\ge\norm A_F^4$, also $\norm A_F^2\le B$ and $v_{22}^2\le c+B$. A bounded solution of a locally Lipschitz equation exists for all time.

\emph{Limit points.} Integrating $\frac{\dd}{\dd t}\Loss=-\norm{\nabla\Loss}^2$ gives $\int_0^\infty\norm{\nabla\Loss(\theta(t))}^2\dd t\le\Loss(\theta(0))-\cR^\star<\infty$, so $\nabla\Loss(\theta(t_k))\to0$ along some sequence $t_k\to\infty$. By boundedness, a subsequence of $\theta(t_k)$ converges to some $\bar\theta\in\cM$, where $\nabla\Loss(\bar\theta)=0$ by continuity and $\bar v_{22}^2\ge c>0$. The $A$-component of the gradient, $\bar v_{22}\,G(\bar\Gamma)$, vanishes, so $\nabla\cR(\bar\Gamma)=0$, and $\bar\Gamma=\Gamma^\star$, the only critical point of the strictly convex $\cR$. Hence $\Loss(\theta(t_k))\to\cR^\star$ along the subsequence; since $\Loss(\theta(t))$ is nonincreasing and bounded below by $\cR^\star$, $\Loss(\theta(t))\to\cR^\star$. Finally $\mu\norm{\Gamma(t)-\Gamma^\star}_F^2\le\cR(\Gamma(t))-\cR^\star\to0$.
\end{proof}

The hypothesis $v_{22}(0)>\norm{W_{11}(0)}_F$ keeps the flow away from the saddle at $\Gamma=0$. From a small, balanced initialization, $c=0$, and the flow must escape that saddle. For isotropic inputs this escape can be solved in closed form.

\begin{proposition}[Escape from the saddle: a logistic curve]\label[proposition]{prop:icl:logistic}
Let $\Sigma=I_d$, and start the gradient flow on $\cM$ at $W_{11}(0)=a_0I_d$ and $v_{22}(0)=\sqrt d\,a_0$ with $a_0>0$, so that $\Gamma(0)=\gamma_0I_d$ with $\gamma_0=\sqrt d\,a_0^2$. Then $\Gamma(t)=\gamma(t)I_d$ for all $t$, where
\begin{equation}\label{eq:icl:logistic}
  \dot\gamma=4\sqrt d\,\gamma\Bigl(1-\frac\gamma{\gamma^\star}\Bigr),\qquad
  \gamma(t)=\frac{\gamma^\star}{1+(\gamma^\star/\gamma_0-1)e^{-4\sqrt d\,t}},
\end{equation}
with $\gamma^\star=N/(N+d+1+\sigma^2)$. The step reaches half its final value at $t_{1/2}=\log(\gamma^\star/\gamma_0-1)/(4\sqrt d)$.
\end{proposition}

\begin{proof}
With $\Sigma=I_d$ we have $H=hI_d$ with $h=(N+d+1+\sigma^2)/N=1/\gamma^\star$, and $G=2(h\Gamma-I_d)$. Try $A=a(t)I_d$ and $v_{22}=\sqrt d\,a(t)$, so $\Gamma=\gamma I_d$ with $\gamma=\sqrt d\,a^2$ and $G=2(h\gamma-1)I_d$. The flow equations of the proof of \cref{thm:icl:flow} become
\[
  \dot a\,I_d=-2\sqrt d\,a(h\gamma-1)I_d,\qquad \dot v_{22}=-\ip{G,aI_d}_F=-2d\,a(h\gamma-1)=\sqrt d\,\dot a ,
\]
so the ansatz is consistent, with the single equation $\dot a=-2\sqrt d\,a(h\gamma-1)$; by uniqueness it gives the solution. Then $\dot\gamma=2\sqrt d\,a\dot a=-4d\,a^2(h\gamma-1)=4\sqrt d\,\gamma(1-h\gamma)$, a logistic equation with carrying capacity $1/h=\gamma^\star$, whose solution with $\gamma(0)=\gamma_0$ is \eqref{eq:icl:logistic}. Setting $\gamma=\gamma^\star/2$ gives $e^{4\sqrt d\,t}=\gamma^\star/\gamma_0-1$.
\end{proof}

For general $\Sigma$, \textcite{ZhangFreiBartlett2024} prove that gradient flow on the population loss, from a suitable random initialization, converges to a global minimum; by \cref{thm:icl:global}, the limit predicts with $\Gamma^\star$.

\begin{proofsketch}
What is omitted is the case $c=0$ for a nonisotropic $\Sigma$: one must show that a balanced initialization whose matrix part is positive semidefinite moves away from the saddle at $\Gamma=0$, where $\nabla\cR(0)=-2\Sigma^2$ pulls $\Gamma$ toward positive definite matrices, and that it is not attracted to the critical set $\{v_{22}=0,\ \tr(\Sigma^2W_{11}^\top)=0\}$. Once $\Gamma$ is bounded away from $0$, the argument of \cref{thm:icl:flow} applies. The details, including the exact form of the initialization, are in \textcite{ZhangFreiBartlett2024}; we have not checked their constants.
\end{proofsketch}

\paragraph{Task diversity.} \Cref{ass:icl:prompts} draws a fresh task for every prompt. If instead the pretraining tasks are a finite set $\{w_1,\dots,w_M\}$, then the derivation of \cref{thm:icl:onestep} goes through with the identity replaced by the task second-moment matrix $\frac1M\sum_kw_kw_k^\top$ (\cref{ex:icl:diversity}): a single linear layer cannot memorize tasks, it can only adapt its preconditioner to their span. \Textcite{LuEtAl2025} solve a linear-attention model of this kind in a proportional limit in which $N$ and the number of tasks grow linearly with $d$ and the number of pretraining prompts quadratically; they find a double descent in the number of pretraining prompts and a transition from memorizing the pretraining tasks to generalizing beyond them as the diversity grows, and they report the same behavior in experiments with linear attention and with nonlinear transformers.

\section{Reading the solution}\label{sec:icl:reading}

\Cref{fig:icl:risk} compares the theory with trained layers. Each point is a layer with all $2(d+1)^2=162$ parameters trained from a random initialization; its risk lies within $1.8$ standard errors of $\cR^\star$ at every $N$, its learned step $\tr(\Gamma)/d$ within $0.004$ of $\gamma^\star$, the off-diagonal part of $\Gamma$ is below $0.011$ in normalized Frobenius norm, and the even part $b(\theta)$ has coefficients of order $10^{-2}$, the level of the optimizer's noise. Training did not get stuck in this experiment, although \cref{thm:icl:flow} covers only initializations on $\cM$.

\begin{figure}[t]
\begin{fullwidth}
\centering
\begin{tikzpicture}
\begin{groupplot}[group style={group size=2 by 1,horizontal sep=17mm},hplot,
  width=0.5\linewidth,height=0.37\linewidth,xmode=log,log basis x=2,xmin=0.8,xmax=320,
  xtick={1,4,16,64,256},xticklabels={$1$,$4$,$16$,$64$,$256$},unbounded coords=jump]
\nextgroupplot[xlabel={context length $N$},ylabel={excess risk per dimension},ymode=log,ymin=0.0015,ymax=1.5,
  title={(a) risk},legend pos=south west,
  ytick={0.001,0.01,0.1,1},yticklabels={$10^{-3}$,$10^{-2}$,$10^{-1}$,$1$}]
\addplot[guide,forget plot] coordinates {(8,0.0015) (8,1.5)};
\addplot[theory] table[x=N,y=onestep] {figures/data/13-icl-risk-theory.dat};
\addlegendentry{optimal one step}
\addplot[theory2] table[x=N,y=bayes] {figures/data/13-icl-risk-theory.dat};
\addlegendentry{Bayes (ridge)}
\addplot[theory3] table[x=N,y=ols] {figures/data/13-icl-risk-theory.dat};
\addlegendentry{least squares}
\addplot[simulation] table[x=N,y=risk] {figures/data/13-icl-risk-sim.dat};
\addlegendentry{trained layer}
\nextgroupplot[xlabel={context length $N$},ylabel={step $\tr(\Gamma)/d$},ymin=0,ymax=1.05,
  title={(b) learned step},legend pos=south east]
\addplot[theory] table[x=N,y=gamma] {figures/data/13-icl-risk-theory.dat};
\addlegendentry{$\gamma^\star=N/(N+d+1+\sigma^2)$}
\addplot[simulation] table[x=N,y=gammahat] {figures/data/13-icl-risk-sim.dat};
\addlegendentry{trained layer}
\end{groupplot}
\end{tikzpicture}
\caption{In-context regression with $d=8$, $\Sigma=I_d$, and $\sigma^2=0.5$; excess risk is $(\text{risk}-\sigma^2)/d$. (a)~Curves: the optimal one-step learner (\cref{thm:icl:onestep}), the Bayes estimator (\cref{cor:icl:bayes}(i), Monte Carlo over $4000$ designs per $N$), and minimum-norm least squares (\cref{thm:doubledescent:exact} with $\E\norm{w_\star}^2=d$; infinite for $\abs{N-d}\le1$; the dashed line marks $N=d$). Points: one layer of linear self-attention with all parameters trained by Adam ($1500$ steps, batches of $1024$ fresh prompts, cosine-decayed rate from $10^{-2}$, initialization $\cN(0,0.01)$), risk estimated on $10^5$ fresh prompts (seed $20261017$). (b)~The learned step $\tr\Gamma(\theta)/d$ of the same layers.}
\label{fig:icl:risk}
\end{fullwidth}
\end{figure}

\paragraph{The learned step.} The optimal preconditioner \eqref{eq:icl:gammastar} is a step size that knows how much to trust the prompt. For $\Sigma=I_d$ it is $\gamma^\star=N/(N+d+1+\sigma^2)$: a short or noisy prompt gives a short step. The denominator adds $N$ (the signal in $s$, whose mean is $w_\star$) to $d+1+\sigma^2$, the fluctuation of $s$ around $w_\star$: $\E\norm{s-w_\star}^2=\frac{d(d+1+\sigma^2)}N$, of which $\sigma^2d/N$ comes from the noise and $d(d+1)/N$ from the fluctuation of $\iclShat$ around $I_d$. As $N\to\infty$, $\Gamma^\star\to\Sigma^{-1}$, so for long prompts the one step becomes a Newton step, and for short ones it is a ridge-like shrinkage, $\Gamma^\star\approx\frac N{\tr\Sigma+\sigma^2}I_d$, which does not depend on the shape of $\Sigma$ at all.

\paragraph{The rate.} At fixed $d$ the excess risk of the optimal step is $d(d+1+\sigma^2)/N+O(N^{-2})$, and that of least squares is $\sigma^2d/(N-d-1)$ (\cref{thm:doubledescent:exact}); the one-step learner is consistent but pays a factor $(d+1+\sigma^2)/\sigma^2$ for not inverting $\iclShat$. Without noise least squares is exact as soon as $N\ge d$, while one step never is: its excess risk is $d(d+1)/(N+d+1)$. In \cref{fig:icl:risk}(a) the gap at $N=32$ is a factor of about $10.5$, against the asymptotic $1+(d+1)/\sigma^2=19$. Near $N=d$ the one-step learner beats the interpolating least-squares estimator, whose variance explodes, but not the Bayes estimator, which is ridge; in the proportional limit at $\gamma=d/N=1$ its excess risk per dimension is $(1+\varsigma)/(2+\varsigma)$, against $\frac12\bigl(\sqrt{\varsigma^2+4\varsigma}-\varsigma\bigr)$ for Bayes (\cref{cor:icl:bayes}): $0.52$ against $0.27$ at $\varsigma=0.1$.\sidenote{We checked both limits in \eqref{eq:icl:proportional} at $d=200$ for $(N,\varsigma)\in\{(200,0.1),(400,0.1),(100,0.5)\}$: the finite-size values differ from the limits by less than $0.002$.}

\paragraph{Plateau, then a logistic jump.} \Cref{fig:icl:dynamics} shows the dynamics of \cref{prop:icl:logistic}. The origin $\Gamma=0$ is a saddle of the loss in the parameters, as in the deep linear networks of \cref{ch:deeplinear}: the product $v_{22}W_{11}^\top$ has a vanishing gradient when both factors are small, so a small initialization spends a time $\log(\gamma^\star/\gamma_0)/(4\sqrt d)$ on a plateau before the step jumps to $\gamma^\star$, the time growing like the logarithm of the inverse initialization. In-context learning appears abruptly in training, but in this model the abruptness is the generic escape from a saddle, not a new mechanism; the induction heads of \cref{ch:induction} give a second instance.

\begin{figure}[t]
\centering
\begin{tikzpicture}
\begin{axis}[hplot,height=0.5\linewidth,xlabel={training time $t=\eta\times\text{step}$},ylabel={step $\tr(\Gamma)/d$},
  xmin=0,xmax=2,ymin=0,ymax=0.75,legend pos=south east]
\addplot[guide,domain=0:2,samples=2,forget plot] {0.677966};
\addplot[theory] table[x=t,y=theoryA] {figures/data/13-icl-dynamics.dat};
\addlegendentry{$\gamma_0=10^{-2}$}
\addplot[theory2] table[x=t,y=theoryB] {figures/data/13-icl-dynamics.dat};
\addlegendentry{$\gamma_0=10^{-4}$}
\addplot[simulation,forget plot] table[x=t,y=simA] {figures/data/13-icl-dynamics.dat};
\addplot[simulation2,forget plot] table[x=t,y=simB] {figures/data/13-icl-dynamics.dat};
\node[font=\scriptsize,text=muted,anchor=north east] at (axis cs:1.98,0.67) {$\gamma^\star$};
\end{axis}
\end{tikzpicture}
\caption{Escape from the saddle. Curves: the logistic solution \eqref{eq:icl:logistic} for $d=8$, $N=20$, $\sigma^2=0.5$ ($\gamma^\star=0.678$, dashed line) and two initial steps $\gamma_0$. Points: gradient descent with step $\eta=0.005$ on all $162$ parameters, each step on a batch of $4096$ fresh prompts, from the balanced initialization of \cref{prop:icl:logistic} plus independent $\cN(0,10^{-8})$ perturbations of every parameter (seed $1017$). The largest deviation from the curves is $0.017$, on the steep part of the second curve.}
\label{fig:icl:dynamics}
\end{figure}

\paragraph{What shifts it tolerates.} The prediction $x_q^\top\Gamma^\star s$ is linear in the labels, and its risk depends on the law of the task only through $\E w_\star w_\star^\top$; for a fixed task it depends on $w_\star$ only through $\norm{w_\star}$ when $\Sigma=I_d$ (\cref{ex:icl:taskshift}). So the learner transfers to new task distributions with the same second moment, including tasks never seen in training. It does not transfer across input distributions: $\Gamma^\star$ depends on $\Sigma$, and the prediction is cubic in the inputs. If the inputs of a test prompt are scaled by $a$, the noiseless part of the target scales by $a$ and that of the prediction by $a^3$. If the input covariance is multiplied by $c=2$, with the parameters of \cref{fig:icl:dynamics}, the excess risk rises from $2.6$ to $15.6$, nearly the null risk $16$, while least squares does not notice the scaling (\cref{ex:icl:covshift}). \Textcite{ZhangFreiBartlett2024} found the trained layer brittle under covariate shift in this sense, and reported that larger nonlinear transformers are more robust.

\section{What the model gets wrong}\label{sec:icl:wrong}

\paragraph{What was thrown away.} Without the softmax, attention cannot select examples: every example enters the prediction with a weight that is a quadratic form, so the layer cannot implement nearest neighbors, sparse regression, or a data-dependent step. With one layer, the forward pass can take one step, and the learner cannot be consistent at the rate of least squares. Putting $x_i$ and $y_i$ in one token removes the hardest part of in-context learning in language, finding which earlier token goes with which, the job of induction heads (\cref{ch:induction}). Linear tasks with a Gaussian prior make the Bayes estimator ridge regression, a linear function of the labels; for other priors the Bayes estimator is nonlinear in the labels, and no single linear layer can follow it. Drawing a fresh task for every prompt hides the memorization regime of \textcite{RaventosEtAl2023}, and training on the population loss hides the number of pretraining prompts needed \parencite{WuEtAl2024,LuEtAl2025}.

\paragraph{Where it breaks.} \emph{Deeper models are not one step.} The transformers of \textcite{GargEtAl2022} reach the error of least squares, far below the one-step risk $d(d+1)/(N+d+1)$ of the noiseless model, so they cannot be computing one gradient step. \Textcite{AhnEtAl2023} proved that for $L$ layers of linear attention certain critical points of the loss implement $L$ steps of preconditioned gradient descent, but not that training finds them; \textcite{FuEtAl2024} found that the predictions of successive layers of trained transformers track the iterates of an iterative Newton method, which converges exponentially faster than gradient descent, and proved that $k$ Newton iterations can be implemented with $k+O(1)$ layers; \textcite{VonOswaldEtAl2023} themselves reported that deeper trained models improve on plain gradient descent with a learned curvature correction. ``One layer is one gradient step'' survives; ``in-context learning is gradient descent'' does not follow from it.

\emph{Softmax attention computes something else.} \Textcite{CollinsEtAl2024} showed that a softmax attention layer trained on in-context regression learns a nearest-neighbor-like predictor whose window adapts to the Lipschitz constant and the noise of the pretraining functions, which linear attention cannot do. A softmax layer is not a gradient step on the squared loss, and its optimum depends on the function class in a way the linear model does not see.

\emph{Language models are not trained on regression prompts.} \Textcite{ShenEtAl2024} compared in-context learning in pretrained language models (LLaMa-7B among them) with gradient descent fine-tuning on the same demonstrations and found that the two respond differently, for instance to the order of the demonstrations; they argued that the evidence for the equivalence came from models trained specifically for in-context regression, with hand-built or specially trained weights. Much of in-context learning in language models is task recognition rather than learning \parencite{PanEtAl2023,MinEtAl2022}, a retrieval of something learned in the weights, which no model with a fresh task per prompt can represent. And in-context learning in transformers can be transient, giving way to in-weights learning as training continues \parencite{SinghEtAl2023}, while in this model the global minimum is stable.

\emph{The data decide whether there is in-context learning at all.} \Textcite{ChanEtAl2022} found that in-context learning emerges only for training data with particular properties, bursty sequences and many rare classes, and that it trades off against in-weights learning. In \cref{ass:icl:prompts} the task changes with every prompt, so in-context learning is the only option.

\paragraph{Fine structure.} The Bayesian view of \textcite{XieEtAl2022} explains in-context learning in language as inference of a latent concept shared by a document, of which the Gaussian posterior of \cref{cor:icl:bayes} is the simplest case. Task diversity was quantified by \textcite{RaventosEtAl2023} empirically and by \textcite{WuEtAl2024} and \textcite{LuEtAl2025} in solvable linear-attention models: \textcite{WuEtAl2024} bound the number of independent pretraining tasks needed and show that the pretrained model is then nearly Bayes optimal at a fixed context length. Transformers can select among algorithms in context, for instance between ridge with different penalties, and reach near-Bayes risk on mixtures of noise levels \parencite{BaiEtAl2023}. Looped transformers, which apply one block repeatedly, learn in-context regression and other function classes about as well as standard transformers with under a tenth of the parameters \parencite{YangEtAl2024}. For softmax attention, \textcite{HuangChengLiang2024} proved convergence of gradient descent for a one-layer transformer on in-context linear regression with features drawn from a finite set, with stage-wise dynamics when the features are imbalanced. The solvable model of how a multi-layer softmax transformer comes to learn in context, and of which algorithm it then runs, is open; the induction-head models of \cref{ch:induction} solve one piece of it.

\section{Model card}\label{sec:icl:card}

\modelcard{Transformers learn a task from the examples in the prompt, without weight updates; trained on regression prompts, they predict nearly as well as least squares.}
{One layer of linear self-attention on prompts $(x_1,y_1,\dots,x_N,y_N,x_q)$ with Gaussian inputs, a Gaussian task drawn per prompt, and Gaussian noise; population pretraining loss.}
{The layer computes $x_q^\top\Gamma s+x_q^\top b$; the global optimum is $b=0$, $\Gamma^\star=N((N+1)\Sigma+(\tr\Sigma+\sigma^2)I)^{-1}$ (\cref{thm:icl:onestep,thm:icl:global}); gradient flow finds it (\cref{thm:icl:flow}), along a logistic curve from a small initialization (\cref{prop:icl:logistic}).}
{One preconditioned gradient step, a Newton step for long prompts; excess risk $d(d+1+\sigma^2)/N$, a factor $(d+1+\sigma^2)/\sigma^2$ above least squares; a plateau of length $\log(1/\gamma_0)/(4\sqrt d)$; tolerance of task shift, brittleness under covariate shift.}
{Deeper and softmax transformers run other algorithms (Newton-like iterations, adaptive nearest neighbors); language models also recognize tasks; data properties and transience; no memorization regime.}
{Multi-layer linear attention; task diversity in proportional limits; algorithm selection; looped transformers; softmax dynamics; induction heads.}

The Gaussian design returns from \cref{ch:doubledescent}: the Bayes estimator of \cref{cor:icl:bayes} is the optimally tuned ridge of \cref{prop:doubledescent:optimal}, and its proportional limit comes from the Marchenko--Pastur law. Rotational and sign symmetries of Gaussian data did the work here as in \cref{ch:singleindex,ch:committee}. The escape from the saddle at $\Gamma=0$ is that of the deep linear networks of \cref{ch:deeplinear}, with the same conservation law and the same logistic curve, and it recurs in the induction heads of \cref{ch:induction}, where a second layer must compose with a first before either can learn.

\begin{exercises}
\item \consolidation (Gaussian fourth moments.) (a) Prove $\E[g_ig_jg_kg_l]=\delta_{ij}\delta_{kl}+\delta_{ik}\delta_{jl}+\delta_{il}\delta_{jk}$ for $g\sim\cN(0,I_d)$ by counting cases. (b) Deduce from \cref{lem:icl:wick} that $\E[\iclShat A\iclShat]=\frac1N\bigl((N-1)\Sigma A\Sigma+\Sigma A\Sigma+\Sigma A^\top\Sigma+\tr(A\Sigma)\Sigma\bigr)$. (c) Show that $\E\norm{s-w_\star}^2=d(d+1+\sigma^2)/N$ when $\Sigma=I_d$.\label{ex:icl:wick}

\item \consolidation (The best scalar step.) Under \cref{ass:icl:prompts}, restrict to $\Gamma=\eta I_d$. Show that the best $\eta$ is $\eta^\star=N\tr(\Sigma^2)/\bigl((N+1)\tr(\Sigma^3)+(\tr\Sigma+\sigma^2)\tr(\Sigma^2)\bigr)$, and that it coincides with $\Gamma^\star$ exactly when $\Sigma$ is a multiple of $I_d$. For $\Sigma=\diag(1,\epsilon)$ and $\sigma=0$, show that as $N\to\infty$ the risk of $\Gamma^\star$ tends to $0$, while that of $\eta^\star I_2$ tends to $\epsilon(1-\epsilon)^2/(1+\epsilon^3)$.\label{ex:icl:scalar}

\item \consolidation (Logistic escape.) (a) For the setting of \cref{fig:icl:dynamics}, compute $t_{1/2}$ for $\gamma_0=10^{-2}$ and $\gamma_0=10^{-4}$ (answers: $0.37$ and $0.78$). (b) Show that halving $\gamma_0$ delays the escape by $\log2/(4\sqrt d)+O(\gamma_0/\gamma^\star)$ as $\gamma_0\to0$, and that the time to reach $(1-\delta)\gamma^\star$ is $\log\bigl((\gamma^\star/\gamma_0-1)(1-\delta)/\delta\bigr)/(4\sqrt d)$. (c) Show that along the flow of \cref{prop:icl:logistic} the loss is $\sigma^2+d-2d\gamma+d\gamma^2/\gamma^\star$, and sketch it against $t$.\label{ex:icl:logistic}

\item \consolidation (Task shift.) Let $\Sigma=I_d$ and fix $w_\star$. Show that the risk of the predictor $x_q^\top\gamma s$, averaged over the inputs and noises only, is $\sigma^2+\norm{w_\star}^2\bigl(1-2\gamma+\gamma^2\frac{N+d+1}N\bigr)+\gamma^2\frac{\sigma^2d}N$. Deduce that the trained layer has the same risk on every task distribution with $\E\norm{w_\star}^2=d$, for instance tasks uniform on the sphere of radius $\sqrt d$, or a single fixed task of that norm.\label{ex:icl:taskshift}

\item \consolidation (The proportional limit.) (a) Check part (iii) of \cref{cor:icl:bayes}, and show that at $\gamma=1$ the limits are $(1+\varsigma)/(2+\varsigma)$ and $\frac12(\sqrt{\varsigma^2+4\varsigma}-\varsigma)$. (b) Show that at $\varsigma\to0$ the Bayes limit is $(1-1/\gamma)_+$, the fraction of directions the prompt does not see, while the one-step limit is $\gamma/(1+\gamma)$.\label{ex:icl:proportional}

\item \variation (Covariate shift.) Train at $\Sigma=I_d$, so that $\Gamma=\gamma^\star I_d$, and test on prompts whose inputs (context and query) have law $\cN(0,cI_d)$. (a) Show that the excess risk is $dc-2d\gamma^\star c^2+d\gamma^{\star2}c^2\bigl((N+1+d)c+\sigma^2\bigr)/N$, which grows like $c^3$. (b) With $d=8$, $N=20$, $\sigma^2=0.5$, evaluate it at $c\in\{\frac12,1,2\}$ and compare with the null risk $dc$. (c) Show that for $N\ge d+2$ the risk of least squares does not depend on $c$.\label{ex:icl:covshift}

\item \variation (Two layers, two steps.) Stack $L$ layers \eqref{eq:icl:lsa}, each acting on all columns of the output of the previous one, with $\iclKQ=\begin{pmatrix}\eta I_d&0\\0&0\end{pmatrix}$ and $\iclPV=\begin{pmatrix}0&0\\0&-1\end{pmatrix}$ in every layer, and read the prediction as minus the bottom entry of the query column. Show by induction that after $\ell$ layers the label entry of column $i\le N$ is $y_i-\ip{w_\ell,x_i}$ and that of the query column is $-\ip{w_\ell,x_q}$, where $w_\ell$ are the gradient descent iterates $w_{\ell+1}=w_\ell-\eta\nabla\Lhat(w_\ell)$ from $w_0=0$ \parencite{VonOswaldEtAl2023}. Replace $\eta I_d$ by $\Gamma^\top$ to get preconditioned steps.\label{ex:icl:twolayers}

\item \variation (Finite task diversity.) Let the task be $\pm w_k$, with $k$ uniform on $\{1,\dots,M\}$ and an independent uniform sign, let $\Sigma=I_d$, and let $W=\frac1M\sum_kw_kw_k^\top$. (a) Show that \cref{thm:icl:onestep,thm:icl:global} hold with
\[
  \Gamma^\star=NW\bigl((N+1)W+(\tr W+\sigma^2)I_d\bigr)^{-1}.
\] (b) If $M<d$ and the $w_k$ are linearly independent, show that as $N\to\infty$ the trained layer tested on fresh tasks $w_\star\sim\cN(0,I_d)$ has excess risk tending to $d-M$: it has learned the span of the training tasks, not the tasks themselves. \hint{the sign flips keep the law symmetric, which is what \cref{thm:icl:global} needs.}\label{ex:icl:diversity}

\item \variation (Other symmetric laws.) (a) Show that \cref{thm:icl:onestep} holds for every task law with $\E w_\star w_\star^\top=I_d$ independent of the inputs and noises. (b) Show that \cref{thm:icl:global} holds whenever, in addition, the laws of $w_\star$ and of the noises are symmetric and have finite fourth moments. (c) Let the query attend to itself, replacing $M$ by $I_{N+1}$ in \eqref{eq:icl:lsa}; show that the extra terms are even in the labels, so that the global minimum is unchanged.\label{ex:icl:symmetric}

\item \simulation (Anisotropic inputs.) Take $d=4$, $\Sigma=\diag(4,2,1,\frac12)$, $N=10$, and $\sigma^2=\frac14$. Train all $50$ parameters of the layer with Adam on fresh prompts, as in \cref{fig:icl:risk}. Observe that the learned $\Gamma(\theta)$ is close to $\diag(0.193,0.336,0.533,0.755)$, that the coefficients of $b(\theta)$ are near zero, and that the risk approaches \eqref{eq:icl:riskstar}. Then test the trained layer with $\Sigma'=\diag(\frac12,1,2,4)$ and explain the result with \eqref{eq:icl:risk}.\label{ex:icl:aniso}

\item \simulation (A softmax head.) Replace linear attention by one softmax head with the same embedding \eqref{eq:icl:Z}, train it on the prompts of \cref{fig:icl:risk} with $d=8$ and $N\in\{8,32\}$, and compare its risk with $\cR^\star$. Then train both models on prompts with labels $y_i=\sgn\ip{w_\star,x_i}$ and compare again. Report which model wins in each case; \cref{thm:icl:global} predicts only the linear model's risk, and the outcome for the softmax head is the point of the exercise.\label{ex:icl:softmax}

\item \research\harder (Global minima of deeper linear transformers.) For $L\ge2$ layers of linear self-attention on the prompts of \cref{ass:icl:prompts}, \textcite{AhnEtAl2023} identify critical points that implement $L$ steps of preconditioned gradient descent. Open, as far as we know: whether these are global minima for general $L$ and $\Sigma$, what the minimal risk is as a function of $(N,d,\sigma^2,L)$, and whether gradient flow from small initialization reaches it, with plateaus as in \cref{prop:icl:logistic} between steps.\label{ex:icl:deep}

\item \research (A solvable softmax model.) Find a task family and a softmax-attention model for which the global minimum of the pretraining loss and the gradient-flow dynamics can be computed as explicitly as in \cref{thm:icl:global,prop:icl:logistic}, and which reproduces the adaptive window of \textcite{CollinsEtAl2024}. Results exist for special data \parencite{HuangChengLiang2024}; a model as complete as the linear one is open.\label{ex:icl:researchsoftmax}
\end{exercises}

\notesandrefs
\Textcite{BrownEtAl2020} described in-context learning in GPT-3 and measured it against the number of demonstrations; the idea of a network that learns in its activations goes back at least to the meta-learning experiments of \textcite{HochreiterEtAl2001}. \Textcite{GargEtAl2022} turned the phenomenon into a controlled experiment on function classes, and within a few months \textcite{AkyurekEtAl2023} and \textcite{VonOswaldEtAl2023}, first posted in November and December 2022, connected trained transformers with gradient descent and ridge regression. Linear attention had been introduced for efficiency \parencite{KatharopoulosEtAl2020} and connected with fast weights \parencite{SchlagEtAl2021}.

The landscape of one layer was analyzed in 2023 by three groups. \Textcite{AhnEtAl2023} showed that the global minimum of a single layer of linear attention implements one step of preconditioned gradient descent, with a preconditioner that adapts to the input distribution and to the finite number of examples, and studied critical points of deeper models. \Textcite{MahankaliEtAl2024} proved that one step of gradient descent is the optimal one-layer learner for isotropic Gaussian inputs, preconditioned for nonisotropic ones, and that it remains one step of gradient descent on the least-squares loss even when the responses are nonlinear functions. \Textcite{ZhangFreiBartlett2024} proved convergence of gradient flow and computed the limit; they analyzed robustness to distribution shift. The odd--even argument of \cref{thm:icl:global} is a short route to the global minimum for the masked layer; we do not know whether it appears in this form in the literature.

The connection with ridge regression and the Bayes estimator is due to \textcite{AkyurekEtAl2023} and \textcite{RaventosEtAl2023}; asymptotic analyses of task diversity are in \textcite{WuEtAl2024} and \textcite{LuEtAl2025}. Algorithm selection in context is analyzed by \textcite{BaiEtAl2023}. The mechanistic study of in-context learning in language models began with induction heads \parencite{OlssonEtAl2022}, the subject of \cref{ch:induction}.

\paragraph{Translating the literature.} \Textcite{ZhangFreiBartlett2024} write $N$, as we do, for the number of examples in a training prompt, and $M$ for the length of a test prompt, which may differ from $N$; their trained layer predicts with our $\Gamma^\star$ for their label model and training length $N$, applied to test prompts of length $M$. \Textcite{AhnEtAl2023} write $n$ for our $N$ and $P$ and $Q$ for the two matrices of a layer, our $\iclPV$ and $\iclKQ$, and call the matrix that multiplies the gradient the preconditioner, our $\Gamma$. \Textcite{VonOswaldEtAl2023} write $e_j=(x_j,y_j)$ for our tokens $z_j$ and keep separate key, query, value, and projection matrices, whose products are our $\iclKQ$ and $\iclPV$; their construction of gradient steps is \cref{ex:icl:twolayers}. \Textcite{GargEtAl2022} write $k$ for the number of in-context examples, our $N$, and report squared errors divided by the dimension, as in \cref{fig:icl:risk}.

\printchapterbib
\end{refsection}

% Chapter 14 of The Hydrogen Atoms of Deep Learning: induction heads and in-context Markov chains.
\newcommand{\indcnt}{\nu}
\newcommand{\indpb}{p^{\mathrm{B}}}
\newcommand{\indAU}{A_{UU}}
\newcommand{\indAX}{A_{UB}}
\newcommand{\indAB}{A_{BB}}
\newcommand{\indLU}{\Loss_{U}}
\newcommand{\indDir}{\operatorname{Dir}}
\newcommand{\indH}{\mathrm{H}}

\casechapter{Induction Heads}{In-context Markov chains and a two-layer circuit}{Attention-only transformer on random Markov chains}{ch:induction}
\begin{refsection}[chapters/14-induction.bib]

\begin{chaptersummary}
Small transformers trained on text acquire, in a short burst early in training, attention heads that complete the pattern $[A][B]\dots[A]\to[B]$, and their ability to use long contexts jumps at the same moment. We study the phenomenon on sequences drawn from a random Markov chain, a new chain for every sequence, so that next-token prediction can only be learned in context. The Bayes-optimal predictor is an add-$\kappa$ bigram estimator, and we prove that a two-layer attention-only circuit (a previous-token head, an induction head, and an attention sink that carries the prior) computes it with an error exponentially small in the attention scales. In a three-parameter caricature with linear attention we prove that gradient flow first learns the in-context unigram statistics, sits on a plateau, and leaves it after a time $\log(1/\alpha)/(2G)$ set by the initialization scale $\alpha$ and an explicit induction gain $G$; a trained softmax transformer shows the same plateau and the same logarithmic dependence. The model is silent on semantics, on heads that copy abstract patterns, and on why induction heads in language models coincide with other changes.
\end{chaptersummary}

\section{The phenomenon: a phase change in context use}\label{sec:induction:phenomenon}

In December 2021 a group at Anthropic published a study of transformers stripped down to attention: no MLPs, no layer normalization, one to three layers \parencite{ElhageEtAl2021}. A one-layer model of this kind is a sum of terms that each read one earlier token, and it learns what the authors called skip-trigrams. Two-layer models did something new. Some heads in the second layer attended, from the current token $A$, to the position \emph{just after} an earlier occurrence of $A$, and copied the token found there: on a sequence $\dots[A][B]\dots[A]$ they predicted $[B]$. Elhage and coauthors named these \term{induction heads}\index{induction head}. Such a head cannot work alone. Its key at position $s$ must contain the token at position $s-1$, which a \term{previous-token head}\index{previous-token head} in the first layer writes there; the second head's keys read the first head's output, a mechanism they called \term{K-composition}\index{K-composition}.\sidenote{In their language a head is described by two low-rank matrices, $W_{QK}=W_Q^\top W_K$, which decides where to attend, and $W_{OV}=W_OW_V$, which decides what to copy. K-composition means that $W_{QK}$ of a later head reads a subspace written by $W_{OV}$ of an earlier one.}

\Textcite{OlssonEtAl2022} followed the heads through training. They measured context use by an \term{in-context learning score}\index{in-context learning score}: the loss on the $500$th token of a $512$-token context minus the loss on the $50$th token, averaged over text. In models of many sizes, with and without MLPs, induction heads appear in a short window early in training, the score improves abruptly in the same window, and the training loss shows a visible bump there. They argued, from several lines of evidence, that induction heads may be the mechanism behind most in-context learning; the evidence is causal for small attention-only models, where removing the heads removes the effect, and correlational for large models with MLPs.

\begin{observation}[\textcite{ElhageEtAl2021}; \textcite{OlssonEtAl2022}]\label[observation]{obs:induction:phase}
In transformers with at least two layers trained on text, induction heads form during a short period early in training. During it the in-context learning score improves abruptly, and the training loss curve shows a bump. One-layer attention-only models do not form induction heads.
\end{observation}

The observation has been reproduced in controlled settings, which also show that it is fragile. \Textcite{ChanEtAl2022} trained transformers on sequences of image--label pairs and found that in-context learning emerges only when the training data are \emph{bursty}, with classes appearing in clusters within a sequence, and contain many rare classes; with a Zipfian class distribution in-context learning and learning in the weights coexist. \Textcite{SinghEtAl2023} found that in-context learning can be \emph{transient}: on such data it emerges and then fades as in-weights learning takes over, although the training loss keeps decreasing. \Textcite{Reddy2024} reduced an in-context classification task to a two-parameter model of an induction head, in which the abrupt transition comes from nested nonlinearities learned in sequence. \Textcite{SinghEtAl2024} dissected the formation of induction heads in a synthetic task and found that several subcircuits must form and interact before the head works.\sidenote{Their tool is a causal intervention on the training dynamics, modifying activations throughout training, rather than on the trained network.}

The cleanest setting is the one we adopt. \Textcite{EdelmanEtAl2024} trained small transformers on sequences, each generated by its own random Markov chain, so that the next token can be predicted well only from the statistics of the current sequence. They observed training in stages: uniform predictions, then a plateau at a suboptimal strategy that uses only the in-context frequencies of single tokens, then a rapid jump to predictions that use the in-context frequencies of pairs, computed by what they called \emph{statistical induction heads}. They also gave evidence that the simpler unigram solution delays the bigram one.

Three questions follow. \emph{What} should an induction head compute, statistically? \emph{How} does a two-layer network compute it, and why two layers? And \emph{why} does it form late and abruptly, after a plateau, rather than gradually from the start?

\section{The model: Markov chains drawn afresh}\label{sec:induction:model}

\paragraph{The data.} We study \term{in-context Markov chains}\index{Markov chain!in-context}. Tokens take values in $[K]=\{1,\dots,K\}$; we use $K$, the global symbol for a number of classes, because next-token prediction is a $K$-class problem. Fix a concentration $\kappa>0$.\sidenote{Papers on this task often write $\alpha$ for the Dirichlet parameter; we reserve $\alpha$ for the initialization scale, as everywhere in the book.} Each sequence is generated in three steps.

\begin{assumption}[In-context Markov chains]\label[assumption]{ass:induction:data}
(i) A transition matrix $P$ is drawn with independent rows $P_{i\cdot}\sim\indDir(\kappa\mathbf 1_K)$, $i\in[K]$. (ii) The first token $x_1$ is drawn from a law $\mu$ on $[K]$ that does not depend on $P$. (iii) Given $P$ and $x_1,\dots,x_s$, the next token is $x_{s+1}=j$ with probability $P_{x_sj}$.
\end{assumption}

Recall that the Dirichlet law $\indDir(\kappa_1,\dots,\kappa_K)$ on the simplex has density proportional to $\prod_jq_j^{\kappa_j-1}$ and mean $\kappa_j/\sum_l\kappa_l$. With probability one every entry of $P$ is positive, so $P$ has a unique stationary law $\pi=\pi(P)$, with $\pi P=\pi$. A model sees a context $x_{1:t}=(x_1,\dots,x_t)$ and outputs a probability vector $q_t\in\R^K$ for $x_{t+1}$; it is trained on the cross-entropy $\E[-\log q_t(x_{t+1})]$ averaged over the positions $t=1,\dots,N$ of contexts of length $N$. A fresh $P$ for every sequence makes all in-weights knowledge about transitions useless: averaged over $P$, every row is uniform.\sidenote{In the language of \cref{ch:icl}, the chain $P$ is the task and the context is its training set; the network must learn a learning algorithm.} The counts in the context are everything,
\begin{gather*}
  \indcnt_{ij}(t)=\#\{s<t:\ x_s=i,\ x_{s+1}=j\},\qquad \indcnt_i(t)=\sum_j\indcnt_{ij}(t),\\
  \indcnt_j^U(t)=\#\{s\le t:\ x_s=j\},
\end{gather*}
the transition counts out of state $i$ and the unigram counts.

\paragraph{The network.} We use an attention-only transformer with two layers, one head each, in which every layer writes its output to a fresh subspace of the residual stream, so that later layers can read each part separately.\sidenote{This ``disentangled'' form is used by \textcite{NichaniEtAl2024}; a standard transformer with a wide residual stream can embed it, by placing token, position and head outputs in orthogonal subspaces.} Position $s\ge1$ carries the one-hot vector $e_{x_s}\in\R^K$ of its token; an extra position $s=0$, the \term{attention sink}\index{attention sink}, carries $e_{x_0}\defeq0$. Layer~1 is a previous-token head of inverse temperature $\beta_1$: position $s\ge1$ attends to $s'\in\{0,\dots,s\}$ with weights
\begin{equation}\label{eq:induction:layer1}
  a^{(1)}_{ss'}=\frac{e^{\beta_1\ind\{s'=s-1\}}}{\sum_{r=0}^se^{\beta_1\ind\{r=s-1\}}},\qquad h_s=\sum_{s'=0}^sa^{(1)}_{ss'}\,e_{x_{s'}}\in\R^K .
\end{equation}
Its logits depend on positions only (one-hot positions and a shift matrix as $W_{QK}$), and its output $h_s$ is written next to $e_{x_s}$. Layer~2 is an induction head of inverse temperature $\beta_2$ with a sink. From the last position $t$ it attends to $s\in\{0,\dots,t\}$ with logits
\begin{equation}\label{eq:induction:layer2}
  \ell_0=\beta_2+\log(K\kappa),\qquad \ell_s=\beta_2\ip{e_{x_t},h_s}\quad(1\le s\le t),
\end{equation}
and weights $a^{(2)}_s=e^{\ell_s}/\sum_re^{\ell_r}$; the query reads the token at $t$ and the key reads the subspace written by layer~1, which is K-composition. Its values are $e_{x_s}$ for $s\ge1$ and the uniform vector $\frac1K\mathbf1_K$ for the sink, so the output
\begin{equation}\label{eq:induction:output}
  \hat y_t=a^{(2)}_0\,\tfrac1K\mathbf 1_K+\sum_{s=1}^ta^{(2)}_s\,e_{x_s}
\end{equation}
is a probability vector with positive entries, and the network predicts $q_t=\hat y_t$. Reading the output in this way, as a mixture of the values, replaces the unembedding and MLP of a real transformer, which would have to compute logarithms of the mixture; in the experiments of \cref{sec:induction:reading} the logits are $\gamma\log\hat y_t$ with a learned $\gamma$.

\paragraph{The caricature for the dynamics.} Gradient flow on softmax attention is not solvable in closed form. For the dynamics we keep the composition and drop everything else: linear attention, square loss, three scalar parameters. Layer~1 has positional score $a\ind\{s'=s-1\}$, so that $h_s=a\,e_{x_{s-1}}$ (and $h_1=0$); layer~2 has score $c+b\ip{e_{x_t},h_s}$, values $e_{x_s}$, and normalization $1/t$. Its output is
\begin{equation}\label{eq:induction:caricature}
  \hat y_t=c\,\varphi_U+ab\,\varphi_B,\qquad \varphi_U=\frac1t\sum_{s=1}^te_{x_s},\qquad \varphi_B=\frac1t\sum_{s=2}^t\ind\{x_{s-1}=x_t\}\,e_{x_s},
\end{equation}
that is, $\varphi_U(j)=\indcnt^U_j(t)/t$ and $\varphi_B(j)=\indcnt_{x_tj}(t)/t$. The parameter $c$ is a unigram strength, $a$ the strength of the previous-token head, and $b$ that of the induction head; the bigram feature enters through the product $ab$ because it needs both layers. We train at one position $t=T$ on the square loss $\Loss(c,a,b)=\E\norm{\hat y_T-e_{x_{T+1}}}^2$ by gradient flow from $c=0$ and $a=b=\alpha$, with a small initialization scale $\alpha$.

\begin{dictionary}
random transition matrix $P$, fresh for each sequence & a document with its own statistics; patterns that recur within it\\
transition counts $\indcnt_{ij}(t)$ & repeated bigrams in a context, $[A][B]\dots[A][B]$\\
layer~1, logits $\beta_1\ind\{s'=s-1\}$ & a previous-token head\\
layer~2, logits $\beta_2\ip{e_{x_t},h_s}$ & an induction head, K-composed with the previous-token head\\
sink of logit $\beta_2+\log(K\kappa)$, value $\frac1K\mathbf1_K$ & an attention sink, here carrying the Dirichlet prior\\
unigram strength $c$ & uniform attention over the context, the first thing learned\\
product $ab$ of layer strengths & the induction circuit, useful only when both layers are in place\\
induction gain $G$ & how much the context's bigram statistics improve on its unigram statistics\\
\end{dictionary}

\paragraph{What was thrown away.} Semantics, MLPs, layer normalization, multiple heads, learned embeddings, and every in-weights regularity: the data are designed so that only in-context statistics help. The circuit model fixes the form of both heads and leaves two scales; the caricature fixes the form of the heads and leaves three scalars, linear attention and a square loss. These are the right first simplifications because the phenomenon is a fact about \emph{composition} (a computation that needs two heads, one feeding the other) and about \emph{statistics} (what an optimal predictor should do with a context), and both survive them.

\paragraph{What one had to notice.} That an induction head is a statistical estimator: copying the token after an earlier match, averaged over all matches, is counting bigrams. That a sequence model trained on a mixture of tasks predicts with the posterior of the task, so that the right benchmark is a Bayes predictor. And that a product of two small strengths is a saddle, the object of \cref{ch:deeplinear}, so that composition alone predicts a plateau whose length grows like the logarithm of the inverse initialization.

\section{The solution: what an induction head computes}\label{sec:induction:circuit}

\subsection{The Bayes predictor}

The \term{Bayes predictor}\index{Bayes predictor} of the next token is its conditional law given the context.

\begin{theorem}[Bayes-optimal next-token prediction]\label{thm:induction:bayes}
Under \cref{ass:induction:data}, for every $t\ge1$ and every context $x_{1:t}$ of positive probability, the rows of $P$ are conditionally independent given $x_{1:t}$, with $P_{i\cdot}\sim\indDir(\kappa\mathbf1_K+\indcnt_{i\cdot}(t))$, and
\begin{equation}\label{eq:induction:bayes}
  \Pr(x_{t+1}=j\mid x_{1:t})=\indpb_t(j)\defeq\frac{\indcnt_{x_tj}(t)+\kappa}{\indcnt_{x_t}(t)+K\kappa}\qquad\text{for every }j\in[K].
\end{equation}
Every predictor $q_t$, a function of $x_{1:t}$ with values in the simplex, satisfies $\E[-\log q_t(x_{t+1})]\ge\E[-\log\indpb_t(x_{t+1})]$, with equality if and only if $q_t=\indpb_t$ almost surely.
\end{theorem}

\begin{proof}
The plan is to write the joint density of $P$ and the context, recognize a product of Dirichlet densities, and take the posterior mean. Let $Z=\Gamma(\kappa)^K/\Gamma(K\kappa)$ be the normalizing constant of $\indDir(\kappa\mathbf1_K)$. By \cref{ass:induction:data}, the joint density of $P$ and $x_{1:t}$ is
\[
  \mu(x_1)\prod_{s=1}^{t-1}P_{x_sx_{s+1}}\prod_{i=1}^K\frac1Z\prod_{j=1}^KP_{ij}^{\kappa-1}=\frac{\mu(x_1)}{Z^K}\prod_{i=1}^K\prod_{j=1}^KP_{ij}^{\kappa+\indcnt_{ij}(t)-1},
\]
because the product over $s$ contains the factor $P_{ij}$ exactly $\indcnt_{ij}(t)$ times. Here we used (ii): the law of $x_1$ does not involve $P$. As a function of $P$, the right side is a product over rows of unnormalized $\indDir(\kappa\mathbf1_K+\indcnt_{i\cdot}(t))$ densities, which proves the first claim. By (iii), $\Pr(x_{t+1}=j\mid x_{1:t},P)=P_{x_tj}$, so $\Pr(x_{t+1}=j\mid x_{1:t})=\E[P_{x_tj}\mid x_{1:t}]$, the mean of the $j$th coordinate of $\indDir(\kappa\mathbf1_K+\indcnt_{x_t\cdot}(t))$, which is \eqref{eq:induction:bayes}.

For the second claim, condition on $x_{1:t}$. The difference of the conditional cross-entropies of $q_t$ and $\indpb_t$ is $\sum_j\indpb_t(j)\log(\indpb_t(j)/q_t(j))$, the Kullback--Leibler divergence of $q_t$ from $\indpb_t$, which is nonnegative and vanishes only when $q_t=\indpb_t$ (Gibbs' inequality). Take expectations.
\end{proof}

Two cases deserve a word. At $t=1$, and whenever the current token has not been seen before position $t$, all counts out of $x_t$ vanish and the prediction is uniform. Hypothesis (ii) is used only in the first display; if the first token is drawn from the stationary law $\pi(P)$, the formula fails, because the first token then carries information about $P$ (\cref{ex:induction:stationary}). The estimator \eqref{eq:induction:bayes} adds $\kappa$ pseudo-counts to each transition: $\kappa=1$ is Laplace's rule of succession and $\kappa=\frac12$ the estimator of \textcite{KrichevskyTrofimov1981}.\sidenote{Laplace derived his rule for two outcomes: after $s$ successes in $n$ trials, the probability of a success next is $(s+1)/(n+2)$ under a uniform prior. Krichevsky and Trofimov showed that $\kappa=\frac12$ is nearly minimax for the cumulative log loss of universal coding.}

Two cruder predictors serve as baselines: the \emph{uniform} predictor $\frac1K\mathbf1_K$, with loss $\log K$ at every position, and the \emph{unigram} predictors $(\indcnt^U_j(t)+\kappa')/(t+K\kappa')$, which use the frequencies of single tokens in the context but not the current token. For long contexts their losses have simple limits. Write $\indH(q)=-\sum_jq_j\log q_j$ for entropy.

\begin{proposition}[Long-context losses]\label[proposition]{prop:induction:limits}
Under \cref{ass:induction:data}, as $t\to\infty$, the expected loss of $\indpb_t$ converges to the mean entropy rate $\indH_B=\E\sum_i\pi_i\indH(P_{i\cdot})$, and that of the unigram predictor, for every fixed $\kappa'>0$, to the mean stationary entropy $\indH_U=\E\,\indH(\pi)$. Moreover $\indH_B\le\indH_U\le\log K$.
\end{proposition}

\begin{proofsketch}
Fix $P$; its entries are positive, so the chain is ergodic, and by the ergodic theorem $\indcnt^U_j(t)/t\to\pi_j$ and $\indcnt_{ij}(t)/t\to\pi_iP_{ij}$ almost surely. Hence $\indpb_t(j)\to P_{x_tj}$ and the unigram predictor tends to $\pi$, uniformly over the next token. The law of $(x_t,x_{t+1})$ given $P$ converges to $\pi_iP_{ij}$. The losses are bounded by $\log((t+K\kappa)/\kappa)$, which grows too slowly to matter after a truncation argument, and the limits follow by dominated convergence; we omit the routine details. For the inequalities: $\indH(\pi)=\indH(\sum_i\pi_iP_{i\cdot})\ge\sum_i\pi_i\indH(P_{i\cdot})$ by concavity of entropy and stationarity, and $\indH(\pi)\le\log K$.
\end{proofsketch}

For $K=2$ the rows are $(1-u,u)$ and $(w,1-w)$ with $u,w$ independent uniform, and the three limits are $\log2=0.693$, $\indH_U=0.564$ and $\indH_B=0.481$ (two-dimensional quadrature, \cref{ex:induction:k2}). For $K=3$, Monte Carlo over $2\times10^5$ matrices gives $\log3=1.099$, $\indH_U=0.979$ and $\indH_B=0.819$. These are the three levels at which training will pause.

\subsection{A two-layer circuit}

\begin{figure}[t]
\centering
\begin{tikzpicture}[x=9mm,y=10mm,>={Stealth[length=2mm]}]
  \tikzset{tok/.style={draw=sea,fill=seafoam,line width=0.5pt,minimum width=7mm,minimum height=6mm,inner sep=0pt,font=\small},
           snk/.style={draw=sandline,fill=sand,line width=0.5pt,minimum width=7mm,minimum height=6mm,inner sep=0pt,font=\small}}
  \node[snk] (s0) at (0,0) {$\varnothing$};
  \foreach \i/\t in {1/C,2/A,3/B,4/C,5/A,6/C,7/A,8/B,9/A}{\node[tok] (x\i) at (\i,0) {$\t$};}
  \node[font=\scriptsize,text=muted] at (0,0.5) {sink};
  \foreach \i in {1,...,9}{\node[font=\scriptsize,text=muted] at (\i,0.5) {$\i$};}
  \foreach \i/\j in {2/3,5/6,7/8}{\draw[->,muted,line width=0.5pt] (x\i.south east) to[bend right=70] (x\j.south west);}
  \foreach \j in {3,6,8}{\draw[->,sea,line width=0.8pt] (9,-0.32) to[bend left=30] (\j,-0.32);}
  \draw[->,sandline!55!black,line width=0.8pt] (9,-0.32) to[bend left=25] (0,-0.32);
  \node[font=\scriptsize,text=muted,anchor=west] at (-0.4,-1.75) {grey: layer 1 copies $x_{s-1}$ to position $s$;};
  \node[font=\scriptsize,text=sea,anchor=west] at (4.3,-1.75) {blue: layer 2 attends from $t$ to $s$ with $x_{s-1}=x_t$};
\end{tikzpicture}
\caption{The circuit of \cref{thm:induction:circuit} on the context $CABCACABA$, with $t=9$ and $x_t=A$. Layer~1 writes $x_{s-1}$ at each position $s$ (three of the grey arrows drawn). Layer~2 attends from position $9$ to the positions $3,6,8$ that follow an earlier $A$, whose tokens are $B,C,B$, and to the sink. With $\kappa=1$ and $K=3$ the output is $\bigl(\frac{0+1}{3+3},\frac{2+1}{3+3},\frac{1+1}{3+3}\bigr)$ for $(A,B,C)$, the Bayes predictor.}\label{fig:induction:circuit}
\end{figure}

The network \eqref{eq:induction:layer1}--\eqref{eq:induction:output} computes $\indpb_t$ in the limit of large inverse temperatures (\cref{fig:induction:circuit}). We need one fact about softmax.

\begin{lemma}[Softmax is Lipschitz in ratio]\label[lemma]{lem:induction:softmax}
Let $\ell,\ell'$ be real vectors indexed by the same finite set, with $\abs{\ell_s-\ell'_s}\le\eps$ for every $s$, and let $w$ and $w'$ be their softmax weights. Then $e^{-2\eps}\le w_s/w'_s\le e^{2\eps}$ for every $s$.
\end{lemma}

\begin{proof}
Write $w_s/w'_s=e^{\ell_s-\ell'_s}\cdot\sum_re^{\ell'_r}/\sum_re^{\ell_r}$. The first factor lies in $[e^{-\eps},e^{\eps}]$, and so does the second, because each term $e^{\ell_r}$ lies between $e^{-\eps}e^{\ell'_r}$ and $e^{\eps}e^{\ell'_r}$.
\end{proof}

\begin{theorem}[A two-layer circuit computes the Bayes predictor]\label{thm:induction:circuit}
Let $\kappa,\beta_1,\beta_2>0$. For every $t\ge1$, every sequence $x_1,\dots,x_t\in[K]$, and every $j\in[K]$, the output \eqref{eq:induction:output} satisfies
\begin{equation}\label{eq:induction:circuit}
  \abs{\log\hat y_t(j)-\log\indpb_t(j)}\le2\beta_2te^{-\beta_1}+\frac{t}{\kappa}e^{-\beta_2} .
\end{equation}
In particular, if $0<\eps<\min\{1,2t/\kappa\}$, $\beta_2=\log(2t/(\kappa\eps))$ and $\beta_1=\log(4\beta_2t/\eps)$, then the cross-entropy of $\hat y_t$ exceeds that of the Bayes predictor by at most $\eps$, for every data distribution.
\end{theorem}

\begin{proof}
The plan is to compare the network with an ideal one whose first layer is exact, using \cref{lem:induction:softmax}, and then to compare the ideal network at finite $\beta_2$ with its limit. Write $i=x_t$ and set $\ind\{x_0=i\}\defeq0$ for the sink.

\emph{Layer 1.} For $s\ge1$ the softmax \eqref{eq:induction:layer1} puts weight $e^{\beta_1}/(e^{\beta_1}+s)$ on $s-1$, so the remaining weight is $\delta_s=s/(e^{\beta_1}+s)\le te^{-\beta_1}$. Since $\ip{e_i,e_{x_{s'}}}=\ind\{x_{s'}=i\}\in\{0,1\}$,
\begin{align*}
  \ip{e_i,h_s}&=(1-\delta_s)\ind\{x_{s-1}=i\}+\sum_{s'\ne s-1}a^{(1)}_{ss'}\ind\{x_{s'}=i\}\\
  &\in\bigl[(1-\delta_s)\ind\{x_{s-1}=i\},\,(1-\delta_s)\ind\{x_{s-1}=i\}+\delta_s\bigr],
\end{align*}
so $\abs{\ip{e_i,h_s}-\ind\{x_{s-1}=i\}}\le\delta_s\le te^{-\beta_1}$.

\emph{An ideal second layer.} Let $\ell^*_0=\ell_0$ and $\ell^*_s=\beta_2\ind\{x_{s-1}=i\}$ for $1\le s\le t$, and let $w$ be their softmax weights. By the bound on layer~1, $\abs{\ell_s-\ell^*_s}\le\beta_2te^{-\beta_1}$ for every $s$, so by \cref{lem:induction:softmax} each weight $a^{(2)}_s$ lies within a factor $e^{\pm2\beta_2te^{-\beta_1}}$ of $w_s$. The values in \eqref{eq:induction:output} have nonnegative coordinates, so the same factor bounds $\hat y_t(j)$ against $\tilde y(j)\defeq w_0/K+\sum_{s\ge1}w_s\ind\{x_s=j\}$:
\begin{equation}\label{eq:induction:step1}
  \abs{\log\hat y_t(j)-\log\tilde y(j)}\le2\beta_2te^{-\beta_1}.
\end{equation}

\emph{The ideal network against the Bayes predictor.} The positions $s\in\{2,\dots,t\}$ with $x_{s-1}=i$ are exactly the ends of the transitions out of $i$ within $x_{1:t}$: there are $\indcnt_i(t)$ of them, and $\indcnt_{ij}(t)$ of them carry the token $j$. Each has unnormalized weight $e^{\beta_2}$; the sink has $K\kappa e^{\beta_2}$; the other $m=t-\indcnt_i(t)$ positions have weight $1$, and $r_j\in[0,m]$ of them carry $j$. Hence
\[
  \tilde y(j)=\frac{e^{\beta_2}(\indcnt_{ij}(t)+\kappa)+r_j}{e^{\beta_2}(\indcnt_i(t)+K\kappa)+m} .
\]
Dividing by $\indpb_t(j)=(\indcnt_{ij}(t)+\kappa)/(\indcnt_i(t)+K\kappa)$ and using $0\le r_j\le m\le t$, $\indcnt_{ij}(t)+\kappa\ge\kappa$ and $\indcnt_i(t)+K\kappa\ge\kappa$,
\[
  \Bigl(1+\frac t\kappa e^{-\beta_2}\Bigr)^{-1}\le\frac{\tilde y(j)}{\indpb_t(j)}\le1+\frac t\kappa e^{-\beta_2},
\]
and $\log(1+u)\le u$ gives $\abs{\log\tilde y(j)-\log\indpb_t(j)}\le te^{-\beta_2}/\kappa$. Together with \eqref{eq:induction:step1} this proves \eqref{eq:induction:circuit}.

\emph{The choice of scales.} With the stated $\beta_2$ the second term of \eqref{eq:induction:circuit} is $\eps/2$, and with the stated $\beta_1$ the first is $\eps/2$. The cross-entropies of $\hat y_t$ and $\indpb_t$ are expectations of $-\log\hat y_t(x_{t+1})$ and $-\log\indpb_t(x_{t+1})$, whose difference is at most $\eps$ pointwise.
\end{proof}

The bound holds sequence by sequence: the circuit is an exact algorithm for the add-$\kappa$ bigram estimator, and the Markov assumption is needed only to make that estimator optimal. We checked \eqref{eq:induction:circuit} on $300$ random sequences and scales; the observed error never exceeded $0.65$ times the bound. Three features of the construction will matter later. The \emph{sink} carries the prior: its logit offset $\log(K\kappa)$ adds $\kappa$ pseudo-counts to every transition, and without it the circuit outputs the empirical bigram distribution, whose cross-entropy is infinite whenever a transition is unseen (\cref{ex:induction:nosink}). When the current token has not occurred before, every non-sink logit is near $0$ against $\beta_2+\log(K\kappa)$ at the sink, and the output is nearly uniform, as the Bayes predictor requires. And when $\beta_2=0$ the second layer ignores the first and returns a mixture of the in-context unigram frequencies and the uniform vector: an untrained induction head is a unigram head.\sidenote{The same remark explains the first stage of training in \cref{fig:induction:transformer}: uniform attention is already a unigram estimator, and only its readout must be learned.}

Why two layers? The key at position $s$ must depend on $x_{s-1}$, information that sits at another position, and in a one-layer transformer keys are computed from one position's embedding alone. \Textcite{SanfordEtAl2024} make this precise: by a communication-complexity argument, a one-layer transformer that solves an induction task must be much larger than a two-layer one. Higher-order chains need more composition: \textcite{RajaramanEtAl2024} show that attention-only transformers with $O(\log k)$ layers represent the in-context conditional distribution of $k$th-order Markov sources, and that layer normalization brings this to a constant number of layers.

\section{The solution: why the circuit forms late}\label{sec:induction:dynamics}

We now turn to the caricature \eqref{eq:induction:caricature}. Its loss depends on $(a,b)$ only through the product $v=ab$, and it is a quadratic in $(c,v)$.

\begin{lemma}[The loss is quadratic in $(c,ab)$]\label[lemma]{lem:induction:quadratic}
Let $\indAU=\E\norm{\varphi_U}^2$, $\indAX=\E\ip{\varphi_U,\varphi_B}$, $\indAB=\E\norm{\varphi_B}^2$, $g_U=\E\ip{\varphi_U,e_{x_{T+1}}}$ and $g_B=\E\ip{\varphi_B,e_{x_{T+1}}}$, at the training position $T$. Then
\[
  \Loss(c,a,b)=1-2g_Uc-2g_Bv+\indAU c^2+2\indAX cv+\indAB v^2,\qquad v=ab .
\]
The moments satisfy $\indAX\ge0$, $g_U>0$, and $\indAU\indAB\ge\indAX^2$ with $\indAU,\indAB>0$.
\end{lemma}

\begin{proof}
Expand $\norm{c\varphi_U+v\varphi_B-e_{x_{T+1}}}^2$ and use $\norm{e_{x_{T+1}}}^2=1$. The features have nonnegative coordinates, so $\indAX\ge0$; the inequality is Cauchy--Schwarz in $L^2$. Finally $g_U>0$ and $\indAU>0$ because $\varphi_U\ne0$, and $\indAB>0$ because $\varphi_B\ne0$ with positive probability.
\end{proof}

The gradient flow $\dot c=-\partial_c\Loss$, $\dot a=-\partial_a\Loss$, $\dot b=-\partial_b\Loss$ reads
\begin{equation}\label{eq:induction:flow}
  \begin{aligned}
  \dot c&=2\bigl(g_U-\indAU c-\indAX v\bigr),\qquad \dot a=b\,R,\qquad \dot b=a\,R,\\
  R&=2\bigl(g_B-\indAX c-\indAB v\bigr).
  \end{aligned}
\end{equation}
At the origin of the $(a,b)$ plane the gradients in $a$ and $b$ vanish, while that in $c$ is $-2g_U\ne0$: the unigram direction is learned at once, and the induction directions sit at a saddle, the saddle of a product $ab$ met in \cref{ch:deeplinear}.\sidenote{There the product of the two layers of a linear network multiplies a singular value of the data; here the product of two attention strengths multiplies a feature of the context. With $a(0)\ne b(0)$ the conserved quantity $a^2-b^2$ plays the role of the balancedness of \cref{ch:deeplinear} (\cref{ex:induction:unbalanced}).} On the unigram manifold $v=0$ the loss is minimized at $c^\star=g_U/\indAU$, with value $\indLU=1-g_U^2/\indAU$. The rate at which the induction circuit grows there is governed by the \term{induction gain}\index{induction gain}\index{plateau}
\begin{equation}\label{eq:induction:gain}
  G\defeq g_B-\indAX c^\star=-\tfrac12\,\partial_v\Loss(c^\star,0),
\end{equation}
the correlation of the bigram feature with the target that the unigram feature does not already explain. If $G>0$ and $\indAU\indAB>\indAX^2$, the quadratic has the unique minimizer
\[
  v^{\star\star}=\frac{\indAU G}{\indAU\indAB-\indAX^2}>0,\qquad c^{\star\star}=\frac{g_U\indAB-g_B\indAX}{\indAU\indAB-\indAX^2},
\]
with value $\Loss^{\star\star}$ below the plateau by $\indLU-\Loss^{\star\star}=\indAU G^2/(\indAU\indAB-\indAX^2)$.

\begin{lemma}[Long-context moments]\label[lemma]{lem:induction:longcontext}
Under \cref{ass:induction:data}, as $T\to\infty$, $\indAU$ and $g_U$ both converge to $\E\norm{\pi}^2$, so that $c^\star\to1$, and
\begin{gather*}
  \indAX\to\E\sum_i\pi_i^2\ip{\pi,P_{i\cdot}},\qquad \indAB\to\E\sum_i\pi_i^3\norm{P_{i\cdot}}^2,\\
  G\to\E\sum_i\pi_i^2\ip{P_{i\cdot}-\pi,P_{i\cdot}} .
\end{gather*}
\end{lemma}

\begin{proofsketch}
By the ergodic theorem, $\varphi_U\to\pi$ and $\frac1T\indcnt_{i\cdot}(T)\to\pi_iP_{i\cdot}$ almost surely given $P$, so $\varphi_B=\sum_i\ind\{x_T=i\}\frac1T\indcnt_{i\cdot}(T)$ is $\pi_{x_T}P_{x_T\cdot}+o(1)$. The law of $(x_T,x_{T+1})$ given $P$ tends to $\pi_iP_{ij}$, and all features are bounded by $1$, so bounded convergence gives the limits; for instance $g_U\to\E\sum_{i}\pi_i\ip{\pi,P_{i\cdot}}=\E\norm\pi^2$ by stationarity, $\sum_i\pi_iP_{i\cdot}=\pi$.
\end{proofsketch}

In the long-context limit the plateau predictor is $\hat y=\varphi_U$, the in-context unigram frequencies, and $G=\E\sum_i\pi_i^2\ip{P_{i\cdot}-\pi,P_{i\cdot}}$ measures how far the rows of the chain are from its stationary law. If all rows are equal, the tokens are independent, $G=0$, and there is nothing for an induction head to learn. For $\kappa=1$ Monte Carlo gives $G=0.034$ for $K=2$, $0.037$ for $K=3$ and $0.030$ for $K=4$, against $\indAU=0.61$, $0.41$ and $0.30$: the unigram direction is learned about ten times faster than the induction circuit grows.

\begin{theorem}[Unigram plateau and escape time]\label{thm:induction:escape}
Assume the conclusions of \cref{lem:induction:quadratic}, $\indAU\indAB>\indAX^2$, and $G>0$. Run \eqref{eq:induction:flow} from $c(0)=0$ and $a(0)=b(0)=\alpha>0$. Let
\[
  \eta_0=\min\Bigl\{\frac G{2\indAB},\ \frac{G\indAU}{4\indAX^2}\Bigr\}
\]
(the second term is $+\infty$ if $\indAX=0$), fix $\eta\in(0,\eta_0]$ with $\alpha^2<\eta$, and let $\tau_\eta=\inf\{t\ge0:\ v(t)\ge\eta\}$. Then:
\begin{enumerate}
\item \emph{(The plateau.)} For $0\le t\le\tau_\eta$, $\abs{c(t)-c^\star(1-e^{-2\indAU t})}\le\indAX\eta/\indAU$ and $\abs{\Loss(t)-\indLU}\le2\indAU{c^\star}^2e^{-4\indAU t}+4g_B\eta$.
\item \emph{(The escape time.)} $\tau_\eta$ is finite, and
\[
  \frac{\log(\eta/\alpha^2)-2\indAX c^\star/\indAU}{4\bigl(G+\indAX^2\eta/\indAU\bigr)}\le\tau_\eta\le\frac{\log(\eta/\alpha^2)}{4\bigl(G-\indAB\eta\bigr)} .
\]
\item \emph{(Convergence.)} If moreover $2\indAU{c^\star}^2e^{-4\indAU\tau_\eta}<G\eta$, which holds for every small enough $\alpha$ by (ii), then $(c,a,b)\to(c^{\star\star},\sqrt{v^{\star\star}},\sqrt{v^{\star\star}})$ and $\Loss(t)\to\Loss^{\star\star}<\indLU$.
\end{enumerate}
\end{theorem}

\begin{proof}
The plan is to reduce the flow to two scalar equations, solve the linear one for $c$ in terms of $v$, and then integrate the growth rate of $v$ between two comparison bounds. Write $A=\indAU$, $B=\indAX$, $C=\indAB$ for brevity.

\emph{Reduction.} The flow \eqref{eq:induction:flow} has polynomial right side, so a solution exists on a maximal interval. Along it $\Loss$ is nonincreasing, since $\frac{\mathrm d}{\mathrm dt}\Loss=-\norm{\nabla\Loss}^2$. The difference $a-b$ satisfies $\frac{\mathrm d}{\mathrm dt}(a-b)=-R\,(a-b)$ and vanishes at $t=0$, so $a=b$ for all $t$, and $v=a^2$. The quadratic form of \cref{lem:induction:quadratic} is positive definite because $AC>B^2$, so $\Loss$ tends to infinity as $\abs c+v\to\infty$; since $\Loss$ is nonincreasing, $c$ and $a$ stay bounded and the solution exists for all $t\ge0$. Finally $\dot a=aR$ gives $a(t)=\alpha\exp\int_0^tR>0$ and
\begin{equation}\label{eq:induction:vgrowth}
  v(t)=\alpha^2\exp\Bigl(\int_0^t2R(s)\,\mathrm ds\Bigr),\qquad 2R=4\bigl(g_B-Bc-Cv\bigr).
\end{equation}
\emph{The unigram coordinate.} The equation for $c$ is linear with source $-2Bv$. Variation of constants, with $c(0)=0$ and $g_U=Ac^\star$, gives
\begin{equation}\label{eq:induction:cformula}
  c(t)=c^\star\bigl(1-e^{-2At}\bigr)-2B\int_0^te^{-2A(t-s)}v(s)\,\mathrm ds .
\end{equation}
For $t\le\tau_\eta$ we have $0<v\le\eta$, so the integral term lies between $0$ and $2B\eta\int_0^\infty e^{-2As}\mathrm ds=B\eta/A$. This is the first claim of (i), and it gives in particular
\begin{equation}\label{eq:induction:cbounds}
  c^\star\bigl(1-e^{-2At}\bigr)-\frac{B\eta}A\le c(t)\le c^\star\qquad(0\le t\le\tau_\eta).
\end{equation}
\emph{Upper bound on $\tau_\eta$.} For $t<\tau_\eta$, \eqref{eq:induction:cbounds} and $v<\eta$ give $2R\ge4(g_B-Bc^\star-C\eta)=4(G-C\eta)$, and $G-C\eta\ge G/2>0$ because $\eta\le G/(2C)$. By \eqref{eq:induction:vgrowth}, $v(t)\ge\alpha^2e^{4(G-C\eta)t}$ for $t<\tau_\eta$. The right side reaches $\eta$ at the time on the right of (ii), so $\tau_\eta$ is finite and at most that time.

\emph{Lower bound on $\tau_\eta$.} For $t\le\tau_\eta$, drop the term $-Cv\le0$ and use the lower bound of \eqref{eq:induction:cbounds}:
\[
  2R\le4\bigl(g_B-Bc\bigr)\le4\Bigl(G+Bc^\star e^{-2At}+\frac{B^2\eta}A\Bigr).
\]
Integrating from $0$ to $\tau_\eta$, with $\int_0^\infty4Bc^\star e^{-2As}\mathrm ds=2Bc^\star/A$, \eqref{eq:induction:vgrowth} at $t=\tau_\eta$, where $v=\eta$ by continuity, gives $\log(\eta/\alpha^2)\le4(G+B^2\eta/A)\tau_\eta+2Bc^\star/A$, which is the left inequality of (ii).

\emph{The loss on the plateau.} Completing the square in $c$ and then expanding in $v$,
\begin{equation}\label{eq:induction:lossdecomp}
  \Loss=\indLU+A(c-c^\star)^2-2v\bigl(g_B-Bc\bigr)+Cv^2 .
\end{equation}
For $t\le\tau_\eta$, \eqref{eq:induction:cbounds} gives $\abs{c-c^\star}\le c^\star e^{-2At}+B\eta/A$, hence $A(c-c^\star)^2\le2A{c^\star}^2e^{-4At}+2B^2\eta^2/A$, and $0\le g_B-Bc\le g_B+B^2\eta/A$. With $v\le\eta$, $C\eta\le G/2$, $4B^2\eta/A\le G$ (both from $\eta\le\eta_0$) and $G\le g_B$,
\[
  \abs{\Loss-\indLU}\le2A{c^\star}^2e^{-4At}+2g_B\eta+\eta^2\Bigl(\frac{4B^2}A+C\Bigr)\le2A{c^\star}^2e^{-4At}+2g_B\eta+\tfrac32G\eta,
\]
which is at most the bound in (i).

\emph{Convergence.} At $t=\tau_\eta$ we have $v=\eta$ and, by \eqref{eq:induction:cbounds}, $g_B-Bc\ge G$; so \eqref{eq:induction:lossdecomp} and the estimates of the previous step give
\[
  \Loss(\tau_\eta)-\indLU\le2A{c^\star}^2e^{-4A\tau_\eta}+\frac{2B^2\eta^2}A-2G\eta+C\eta^2\le2A{c^\star}^2e^{-4A\tau_\eta}-G\eta<0
\]
under the hypothesis of (iii). Since $\int_0^\infty\norm{\nabla\Loss}^2\mathrm dt=\Loss(0)-\lim_t\Loss(t)<\infty$ and $\norm{\nabla\Loss}^2$ has a bounded time derivative on the bounded trajectory, $\nabla\Loss\to0$ along the flow. Every limit point is therefore a critical point with $\Loss\le\Loss(\tau_\eta)<\indLU$ and $a=b\ge0$. A critical point with $a=b$ has $\partial_c\Loss=0$ and either $a=0$ or $\partial_v\Loss=0$. If $a=0$, then $v=0$, $c=c^\star$, and $\Loss=\indLU$, which is excluded. Otherwise $(c,v)$ is the unique minimizer $(c^{\star\star},v^{\star\star})$ of the quadratic, and $a=b=\sqrt{v^{\star\star}}$. The limit point is unique, so the flow converges to it.
\end{proof}

\begin{corollary}[Logarithmic escape]\label[corollary]{cor:induction:log}
Under the hypotheses of \cref{thm:induction:escape}, for every fixed $\eta\in(0,\eta_0]$,
\begin{gather*}
  \frac1{2(G+\indAX^2\eta/\indAU)}\le\liminf_{\alpha\to0}\frac{\tau_\eta}{\log(1/\alpha)},\\
  \limsup_{\alpha\to0}\frac{\tau_\eta}{\log(1/\alpha)}\le\frac1{2(G-\indAB\eta)} .
\end{gather*}
\end{corollary}

\begin{proof}
Divide the bounds of (ii) by $\log(1/\alpha)$, note $\log(\eta/\alpha^2)/\log(1/\alpha)\to2$, and let $\alpha\to0$.
\end{proof}

The plateau thus lasts $\log(1/\alpha)/(2G)$, up to a factor $1+O(\eta)$ and an additive constant, while the unigram coordinate settles in time $O(1/\indAU)$. The escape is the logistic escape of \cref{ch:deeplinear}: two factors that each need the other grow together as $v\propto e^{4Gt}$ while small, the composite strength $v$ at rate $4G$, until the quadratic term $\indAB v^2$ stops them. The induction gain $G$ plays the role of a singular value of the input--output correlation, and the unigram feature the role of a mode that has already been learned and must be projected out of the bigram signal; this projection is the competition that \textcite{EdelmanEtAl2024} observed.

\paragraph{Rigorous results for softmax models.} \Textcite{NichaniEtAl2024} analyze gradient descent on a simplified two-layer transformer for in-context learning of a latent causal structure; they show that the first attention layer learns the causal graph because the gradient of its attention matrix reflects mutual information between tokens, and for Markov chains the trained network implements an induction head. \Textcite{BiettiEtAl2023} study a two-layer transformer on data mixing global bigrams with context-specific ones, describe weight matrices as associative memories, and show that global bigrams are learned quickly and the induction head more slowly, with theoretical insight into how gradient steps build these memories. Neither result gives the plateau length as a function of the initialization; that is the piece our caricature supplies, at the cost of linear attention.

\section{Reading the solution}\label{sec:induction:reading}

The model makes five quantitative predictions. (i) Loss levels: a trained network should pass through the uniform level $\log K$, then the unigram level, then the Bayes level of \cref{thm:induction:bayes}, whose long-context values are $1.099$, $0.979$ and $0.819$ for $K=3$, $\kappa=1$. (ii) The Bayes predictor is computable by two heads and a sink, with error $O(\beta_2te^{-\beta_1}+te^{-\beta_2}/\kappa)$, so a two-layer network can reach the Bayes level exactly. (iii) After the unigram coordinate settles, the loss sits at $\indLU$ for a time $\log(1/\alpha)/(2G)$: each factor $10$ in the initialization adds $\log(10)/(2G)$ to the plateau, and the drop from it is a logistic curve of rate $4G$. (iv) Weaker rows mean longer plateaus: $G$ falls as $\kappa$ grows, like $(K-1)^2/(K^3(K\kappa+1))$ for large $\kappa$ (\cref{ex:induction:kappa}); for $K=3$ and $\kappa=\frac14,\frac12,1,2,4$ Monte Carlo gives $G=0.070$, $0.055$, $0.037$, $0.021$ and $0.0115$. (v) Even the Bayes predictor learns slowly from a finite context. The add-$\kappa$ estimator of one row from $n$ transitions has expected excess loss $(K-1)/(2n)$ to leading order; with $n\approx\pi_it$ transitions out of state $i$ and the current state distributed as $\pi$, the excess loss at position $t$ is
\begin{equation}\label{eq:induction:excess}
  \E\bigl[-\log\indpb_t(x_{t+1})\bigr]-\indH_B\approx\sum_i\pi_i\frac{K-1}{2\pi_it}=\frac{K(K-1)}{2t},
\end{equation}
a heuristic that Monte Carlo confirms slowly ($0.048$, $0.027$ and $0.0141$ at $t=50,100,200$ for $K=3$, against $0.060$, $0.030$ and $0.015$). This is why the Bayes level averaged over a short context lies well above $\indH_B$.

\begin{figure}[t]
\begin{fullwidth}
\centering
\begin{tikzpicture}
\begin{groupplot}[group style={group size=2 by 1,horizontal sep=17mm},hplot,width=0.475\linewidth,height=0.37\linewidth]
\nextgroupplot[ymode=log,xmin=0,xmax=200,ymin=1e-6,ymax=1,xlabel={time $t$},ylabel={$\Loss(t)-\Loss^{\star\star}$},
  ytick={1e-6,1e-4,1e-2,1},legend style={at={(0.97,0.97)},anchor=north east}]
\addplot[theory] table[x=t,y=ode] {figures/data/14-induction-caricature-curve.dat};
\addplot[simulation,mark repeat=4] table[x=sgdt,y=sgd] {figures/data/14-induction-caricature-curve.dat};
\addplot[guide] coordinates {(0,0.03136) (200,0.03136)};
\legend{gradient flow,SGD,$\indLU-\Loss^{\star\star}$}
\nextgroupplot[xmin=0.8,xmax=5.2,ymin=0,ymax=180,xlabel={$\log_{10}(1/\alpha)$},ylabel={escape time},
  legend style={at={(0.03,0.97)},anchor=north west}]
\addplot[theory] table[x=log10inv,y=ode] {figures/data/14-induction-caricature-escape.dat};
\addplot[simulation] table[x=log10inv,y=sgdmean] {figures/data/14-induction-caricature-escape.dat};
\addplot[guide] table[x=log10inv,y=asym] {figures/data/14-induction-caricature-escape.dat};
\legend{gradient flow,SGD (mean of 5),$\log(1/\alpha)/(2G)$}
\end{groupplot}
\end{tikzpicture}
\caption{The caricature \eqref{eq:induction:caricature} with $K=3$, $\kappa=1$, uniform first token and training position $T=24$. The moments are computed on a pool of $4\times10^5$ sequences (seed $0$): $\indAU=0.433$, $\indAX=0.185$, $\indAB=0.115$, $g_U=0.410$, $g_B=0.209$, hence $c^\star=0.947$, $G=0.0337$, $\indLU=0.6116$ and $\Loss^{\star\star}=0.5802$. Curves: the gradient flow \eqref{eq:induction:flow}, integrated numerically. Points: minibatch SGD on $(c,a,b)$ with batches of $256$ sequences from the pool and step $0.05$, plotted against time $0.05\times$step. Left: excess loss for $\alpha=10^{-3}$; after the unigram coordinate settles ($t\approx5$) the loss stays on the plateau until $t\approx70$ and then drops logistically; the SGD points level off near $10^{-5}$, the noise floor of the step size. Right: the time at which $v=ab$ reaches $v^{\star\star}/2$; flow and SGD (five seeds, spread below $1.2$) agree, and both are parallel to $\log(1/\alpha)/(2G)$, with slope $34.2$ per decade.}\label{fig:induction:caricature}
\end{fullwidth}
\end{figure}

\Cref{fig:induction:caricature} tests (iii). The plateau is flat to four digits, at the value $\indLU$ predicted from the moments, and its length grows by $34.2$ time units per decade of $\alpha$, the slope $\log(100)/(4G)$ of \cref{cor:induction:log}; the curves sit about $5$ units below the asymptote, the additive constant of (ii). The bounds of (ii) bracket the flow: at $\eta=\eta_0=0.107$ and $\alpha=10^{-3}$ they give $64\le\tau_\eta\le135$, and the flow takes $81$. The second drop is small, $\indLU-\Loss^{\star\star}=0.031$ against the first drop $0.39$: the square loss of the caricature rewards the bigram feature less than cross-entropy does, and only the time scale, not the size of the gain, should be compared with experiments.

\begin{figure}[t]
\centering
\begin{tikzpicture}
\begin{axis}[hplot,xmin=0,xmax=3000,ymin=0.935,ymax=1.118,xlabel={SGD step},ylabel={validation loss},
  legend style={at={(0.97,0.2)},anchor=south east},legend cell align=left]
\addplot[sea,line width=1.0pt] table[x=step,y=l1] {figures/data/14-induction-transformer.dat};
\addplot[seadeep,line width=1.0pt,dashed] table[x=step,y=l2] {figures/data/14-induction-transformer.dat};
\addplot[sandline!55!black,line width=1.0pt] table[x=step,y=l3] {figures/data/14-induction-transformer.dat};
\addplot[guide] table[x=step,y=uniform] {figures/data/14-induction-transformer.dat};
\addplot[guide] table[x=step,y=unigram] {figures/data/14-induction-transformer.dat};
\addplot[guide] table[x=step,y=bigram] {figures/data/14-induction-transformer.dat};
\legend{init $10^{-1}$,init $10^{-2}$,init $10^{-3}$}
\node[font=\scriptsize,text=muted,anchor=south west] at (axis cs:1500,1.0986) {uniform, $\log3$};
\node[font=\scriptsize,text=muted,anchor=south west] at (axis cs:1900,1.0322) {best add-$\kappa'$ unigram};
\node[font=\scriptsize,text=muted,anchor=north east] at (axis cs:2980,0.9465) {Bayes, \cref{thm:induction:bayes}};
\end{axis}
\end{tikzpicture}
\caption{A softmax transformer on in-context Markov chains ($K=3$, $\kappa=1$, context $N=32$, loss averaged over positions, validation set of $4000$ sequences, seed $0$). Layer~1: one head with learned relative-position logits and value matrix; layer~2: one head with a learned bilinear score on (token, layer-1 output) and a sink with learned logit; logits $\gamma\log\hat y_t$ with learned $\gamma$; all weights initialized as the stated scale times a standard Gaussian, $\gamma$ and the sink logit at $0$. Plain SGD, learning rate $1$, batches of $64$ fresh sequences. Dashed lines: the three levels, computed on the validation set. The loss crosses the midpoint between the unigram and Bayes levels at steps $680$, $1200$ and $1740$, and ends at $0.952$; the previous-token weight $a^{(1)}_{t,t-1}$ ends near $0.9$.}\label{fig:induction:transformer}
\end{figure}

\Cref{fig:induction:transformer} tests (i) and the qualitative content of (iii) on a softmax network that can represent the circuit of \cref{thm:induction:circuit}. The three stages appear: the readout $\gamma$, whose gradient is of order one, turns the uniform attention into a unigram estimator within about $300$ steps; the loss then sits at $1.029$, slightly below the best add-$\kappa'$ unigram estimator because the readout can also sharpen it; and it then drops to within $0.006$ of the Bayes level. The plateau grows by about $530$ steps per decade of initialization, linearly in $\log(1/\alpha)$ as \cref{cor:induction:log} predicts. The slope itself is not predicted: the softmax network has more than two factors in its induction path (relative-position logits, a value matrix, a bilinear score and the readout), cross-entropy instead of square loss, and stochastic gradients. With the previous-token head frozen in place the plateau disappears, and the loss crosses the midpoint by step $60$ (\cref{ex:induction:freeze}), which confirms that composition causes the delay.

\section{What the model gets wrong}\label{sec:induction:wrong}

\paragraph{What was thrown away.} \emph{The data} are designed so that in-context statistics are the only signal: a fresh chain per sequence, a symmetric prior, no vocabulary structure, no semantics. Real text mixes in-context and in-weights regularities, and their competition is a large part of the phenomenon. \emph{The circuit} is fixed by hand, with one head per layer and the positional pattern of a previous-token head built in; real induction heads match on fuzzy, learned features and copy through learned subspaces. \emph{The caricature} keeps one scalar per layer, linear attention and a square loss; it captures the saddle created by composition and nothing else. In particular it does not explain why gradient descent finds the \emph{previous-token} pattern among all positional patterns, a symmetry-breaking question that \textcite{NichaniEtAl2024} answer with mutual information and that our three-parameter ansatz assumes away. \emph{The readout} of our experiments is a logarithm of the attention output, which spares the network the MLP a real transformer needs to turn mixtures of one-hot vectors into calibrated logits.

\paragraph{Where it breaks.} \emph{Language models may do more at the phase change.} For large models with MLPs the evidence of \textcite{OlssonEtAl2022} that induction heads cause in-context learning is correlational: the heads and the jump in context use appear in the same window of training, but other mechanisms may form in that window too, and the abrupt improvement is not cleanly attributable to a two-head circuit. \emph{In-context learning in large models goes beyond copying.} \Textcite{ToddEtAl2024} find a small set of heads whose outputs, averaged over demonstrations, form \emph{function vectors} that trigger a task even in zero-shot contexts; \textcite{HendelEtAl2023} find that in-context learning compresses the demonstrations into a single \emph{task vector}. These are abstractions of a task, not counts of repeated tokens. \emph{Data control whether the head forms at all.} \Textcite{ChanEtAl2022} show that in-context learning requires bursty data with many rare classes, and \textcite{SinghEtAl2023} that it can disappear with further training as in-weights learning takes over; in our model $G=0$ for data without in-context structure (\cref{ex:induction:iid}), which captures the first finding, but nothing in it can capture the second, because in-weights knowledge is useless by construction. \emph{One layer gets stuck.} \Textcite{MakkuvaEtAl2024} (unrefereed in the version we cite) characterize the loss landscape of single-layer transformers on Markov data and find bad local minima that depend on the data and the architecture; our caricature has no local minima at all, only a saddle. \emph{The plateau is not always one plateau.} \Textcite{SinghEtAl2024} find that induction-head formation depends on several interacting subcircuits, and \textcite{Reddy2024} that the abrupt transition comes from a sequence of nested nonlinearities; the caricature, with one product of two strengths, has a single escape.

\paragraph{Fine structure.} \emph{Rigorous dynamics.} \Textcite{NichaniEtAl2024} prove that gradient descent on a simplified two-layer transformer learns latent causal graphs, and induction heads for Markov chains, and \textcite{BiettiEtAl2023} explain the order in which global and in-context bigrams are learned through associative memories. \emph{Higher-order chains.} For $k$th-order Markov sources the Bayes predictor counts $(k+1)$-grams; \textcite{EdelmanEtAl2024} extend their experiments to such sources, \textcite{AkyurekEtAl2024} find ``$n$-gram heads'', higher-order induction heads, in transformers trained on in-context language learning tasks, and \textcite{RajaramanEtAl2024} show that depth logarithmic in $k$ suffices, and constant depth with layer normalization. \emph{Lower bounds.} \Textcite{SanfordEtAl2024} prove that one-layer transformers need much larger size than two-layer ones for induction tasks. \emph{Attention sinks.} \Textcite{XiaoEtAl2024} found that language models put large attention on the first tokens regardless of content and that keeping these ``sinks'' stabilizes streaming inference; in our circuit a sink has a statistical role, carrying the prior. Open: a theorem for the softmax network of \cref{fig:induction:transformer} that predicts the slope of the plateau length against $\log(1/\alpha)$, and an account of how the unigram stage, which our proof treats as fast, changes the escape when it is not.

\section{Model card}\label{sec:induction:card}

\modelcard{Induction heads form abruptly after a plateau, and context use jumps at the same time.}{Sequences from a random Markov chain with $\indDir(\kappa\mathbf1_K)$ rows, fresh per sequence; a two-layer attention-only circuit with a sink; for the dynamics, linear attention with three scalar strengths $c$, $a$, $b$.}{Bayes predictor $(\indcnt_{x_tj}+\kappa)/(\indcnt_{x_t}+K\kappa)$; a two-layer circuit computes it with error $O(\beta_2te^{-\beta_1}+te^{-\beta_2}/\kappa)$; gradient flow learns the unigram first and escapes the induction saddle at time $\log(1/\alpha)/(2G)$.}{Plateaus at $\log K$, $\indH_U$ and $\indH_B$; plateau length linear in $\log(1/\alpha)$; $G$ and escape rate set by the spread of the rows, $G\sim(K-1)^2/(K^3(K\kappa+1))$ as $\kappa\to\infty$.}{No semantics or in-weights competition; no transience; language-model phase changes involve more than one circuit; function and task vectors; bad local minima of one-layer models.}{Causal-structure dynamics; associative memories; $n$-gram heads and higher-order chains; depth lower bounds; attention sinks.}

The chapter closes the book's thread on escape from saddles. The escape time $\log(1/\alpha)/(2G)$ is that of a deep linear network (\cref{ch:deeplinear}), with $G$ in the role of a singular value; it is a cousin of the search phase of single-index models (\cref{ch:singleindex}), of the plateaus of the committee machine (\cref{ch:committee}), where the student must break a symmetry, and of the hidden progress of sparse parities (\cref{ch:parity}). As in grokking (\cref{ch:grokking}), a circuit can be written down exactly and the hard part is the dynamics. The comparison with \cref{ch:icl} is instructive: there, linear attention on regression tasks implements one step of gradient descent with a single layer; here, the in-context task is discrete and the optimal algorithm needs two composed heads, which is where the plateau comes from. The question whether such plateaus add up to the smooth loss curves of large models is taken up in \cref{ch:scaling}.

\begin{exercises}
\item \consolidation (When the first token knows the chain.) (a) Redo \cref{thm:induction:bayes} with rows $P_{i\cdot}\sim\indDir(\kappa_{i1},\dots,\kappa_{iK})$ and show $\Pr(x_{t+1}=j\mid x_{1:t})=(\indcnt_{x_tj}+\kappa_{x_tj})/(\indcnt_{x_t}+\sum_l\kappa_{x_tl})$. (b) Let $K=2$, $\kappa=1$, and draw $x_1$ from the stationary law $\pi(P)$. Show that $\Pr(x_2=x_1\mid x_1)=2\E[w(1-u)/(u+w)]$ with $u,w$ independent uniform on $[0,1]$, and evaluate it numerically (about $0.591$), while \eqref{eq:induction:bayes} gives $\frac12$. Explain which step of the proof fails. \label{ex:induction:stationary}

\item \consolidation (Three levels for two states.) For $K=2$ and $\kappa=1$, write the stationary law and the entropy rate of $P=\bigl(\begin{smallmatrix}1-u&u\\w&1-w\end{smallmatrix}\bigr)$, and compute $\indH_U=\E\,\indH(\pi)$ and $\indH_B$ by two-dimensional quadrature. You should find $0.564$ and $0.481$. Show that $\indH_B<\indH_U$ unless $u+w=1$, and interpret $u+w=1$. \label{ex:induction:k2}

\item \consolidation (The $1/t$ excess.) Let $n$ draws from $q\in\R^K$, with all $q_j>0$, give counts $n_j$, and let $\hat q_j=(n_j+\kappa)/(n+K\kappa)$. Show that $\E\,\mathrm{KL}(q\,\|\,\hat q)=(K-1)/(2n)+O(n^{-2})$ by expanding $\log\hat q_j$ to second order around $q_j$. Derive \eqref{eq:induction:excess}, and compare with the Monte Carlo values in \cref{sec:induction:reading}. \label{ex:induction:excess}

\item \consolidation (The quadratic.) (a) Verify the formulas for $v^{\star\star}$, $c^{\star\star}$ and $\indLU-\Loss^{\star\star}$ after \eqref{eq:induction:gain}. (b) Check them against the moments in the caption of \cref{fig:induction:caricature}. (c) Show that $c^\star-c^{\star\star}=\indAX G/(\indAU\indAB-\indAX^2)$: once the induction circuit is on, the unigram strength decreases. \label{ex:induction:quadratic}

\item \consolidation (Without a sink.) Remove the sink from the circuit, so that layer~2 attends only to $s\in\{1,\dots,t\}$. (a) Show that as $\beta_1,\beta_2\to\infty$ the output tends to the empirical distribution $\indcnt_{x_tj}/\indcnt_{x_t}$ when $\indcnt_{x_t}>0$, and to the in-context unigram frequencies when $\indcnt_{x_t}=0$. (b) Show that the expected cross-entropy at position $t\ge3$ tends to $+\infty$ as $\beta_1,\beta_2\to\infty$, and that for fixed large $\beta_1$ it grows linearly in $\beta_2$. \label{ex:induction:nosink}

\item \variation (Second-order chains.) Let $x_{s+1}$ depend on $(x_{s-1},x_s)$ through a random $K^2\times K$ matrix with $\indDir(\kappa\mathbf1_K)$ rows. (a) Derive the Bayes predictor. (b) Build a two-layer circuit with two heads in layer~1, writing $x_{s-1}$ and $x_{s-2}$ into two subspaces, and a layer-2 head whose logit is $\beta_2(\ip{e_{x_t},h_s}+\ip{e_{x_{t-1}},h'_s}-1)$, and prove the analogue of \cref{thm:induction:circuit}. (c) What does the analogue of the caricature predict about the order in which unigram, bigram and trigram statistics are learned? \label{ex:induction:trigram}

\item \variation (Unbalanced initialization.) Start the caricature from $a(0)=\alpha$, $b(0)=0$. (a) Show that $a^2-b^2$ is conserved and that $\dot v=\sqrt{\alpha^4+4v^2}\,R$. (b) Show that $v$ first grows linearly and then exponentially at rate $4G$, and that the escape time is still $\log(1/\alpha)/(2G)+O(1)$ as $\alpha\to0$. Compare with the unbalanced deep linear networks of \cref{ch:deeplinear}. \label{ex:induction:unbalanced}

\item \variation\hard (No in-context structure.) Suppose all rows of $P$ are equal to one random vector $\pi\sim\indDir(\kappa\mathbf1_K)$, so that tokens are independent given $\pi$. (a) Show that in the long-context limit $G=0$. (b) Assume $G\le0$ and $\indAU\indAB>\indAX^2$. Using \eqref{eq:induction:vgrowth} and \eqref{eq:induction:cformula}, show that $v(t)\le\alpha^2e^{2\indAX c^\star/\indAU}$ for all $t$, and that the flow converges to $(c^\star,0,0)$: the induction circuit never forms. \hint{at a time when $v$ reaches its running maximum $V$, bound $c^\star-c\le c^\star e^{-2\indAU t}+\indAX V/\indAU$ and use $\indAX^2<\indAU\indAB$ to show that $\frac{\mathrm d}{\mathrm dt}\log v\le4\indAX c^\star e^{-2\indAU t}$ there.} Relate this to the role of burstiness in \textcite{ChanEtAl2022}. \label{ex:induction:iid}

\item \variation\hard (The induction gain for nearly uniform rows.) Write $P_{i\cdot}=\frac1K\mathbf1_K+\Delta_i$. (a) Show that $\pi=\frac1K\mathbf1_K+\frac1K\sum_i\Delta_i+O(\norm\Delta^2)$. (b) Using \cref{lem:induction:longcontext}, show that $G=\frac{K-1}{K^2}\E\norm{\Delta_1}^2+O(\E\norm\Delta^3)$, and with $\E\norm{\Delta_1}^2=(K-1)/(K(K\kappa+1))$ conclude $G\sim(K-1)^2/(K^3(K\kappa+1))$ as $\kappa\to\infty$. (c) Compare with the Monte Carlo values $0.021$ and $0.0115$ for $K=3$, $\kappa=2,4$. \hint{the terms of first order in $\Delta$ cancel because $\ip{\Delta_i,\mathbf1_K}=0$.} \label{ex:induction:kappa}

\item \simulation (Composition causes the delay.) In the script of \cref{fig:induction:transformer}, at initialization $10^{-2}$, (a) freeze layer~1 as a previous-token head (relative-position logit $6$ at offset $1$ and $0$ elsewhere, value matrix $I_K$). The reader should observe that the loss crosses the midpoint between the unigram and Bayes levels by step $60$ instead of $1200$: with one factor fixed, the induction path loses its saddle. (b) Remove the sink. The reader should observe that the loss stalls near $0.98$ and fluctuates, because unseen transitions receive near-zero probability (\cref{ex:induction:nosink}). \label{ex:induction:freeze}

\item \simulation (Plateau length and the prior.) Run the script of \cref{fig:induction:caricature} with $\kappa\in\{\frac12,1,2\}$ and measure the slope of the escape time against $\log_{10}(1/\alpha)$. The reader should find slopes close to $\log(100)/(4G(\kappa))$ with $G$ computed from the pool moments at $T=24$, increasing with $\kappa$ roughly as $1/G\propto K\kappa+1$ for large $\kappa$. \label{ex:induction:kappasim}

\item \simulation (The Bayes level is reached.) Train the transformer of \cref{fig:induction:transformer} for $10^4$ steps and record the learned sink logit $\sigma$ and the layer-2 score of a matching position. \Cref{thm:induction:circuit} predicts that their difference approaches $\log(K\kappa)=\log3$ when the circuit is sharp. Report how close the trained network comes, and why the readout $\gamma$, which is not part of the theorem, can absorb part of the discrepancy. \label{ex:induction:sink}

\item \research (Dynamics of softmax induction heads.) Prove a version of \cref{thm:induction:escape} for the softmax network of \cref{fig:induction:transformer}, from a small Gaussian initialization: identify the unigram plateau, show that its length is $\Theta(\log(1/\alpha))$, and compute the constant. \Textcite{NichaniEtAl2024} and \textcite{BiettiEtAl2023} analyze related models with staged training or a few gradient steps; a theorem that predicts the plateau length of a trained softmax network is open. \label{ex:induction:research}

\item \research (Stages for higher-order chains.) For $k$th-order chains, do transformers learn the $1$-gram, $2$-gram, \dots, $(k+1)$-gram statistics in successive stages with their own plateaus, and are the plateau lengths predicted by a caricature with $k$ composed strengths? \Textcite{EdelmanEtAl2024} and \textcite{AkyurekEtAl2024} study such sources experimentally; a quantitative theory of the stage lengths is open. \label{ex:induction:stages}
\end{exercises}

\notesandrefs
Induction heads were named by \textcite{ElhageEtAl2021}, in a study of attention-only transformers that introduced the residual-stream view, the $W_{QK}$ and $W_{OV}$ circuits and K-composition; \textcite{OlssonEtAl2022} connected their formation to an abrupt improvement of in-context learning. Both appeared in the Transformer Circuits Thread, and the second also on arXiv. The data-dependence of in-context learning is from \textcite{ChanEtAl2022} and \textcite{SinghEtAl2023}, and mechanistic studies of formation from \textcite{Reddy2024} and \textcite{SinghEtAl2024}.

The in-context Markov chain task is that of \textcite{EdelmanEtAl2024}, who observed the uniform, unigram and bigram stages, studied the effect of the prior over chains, and extended the task to $n$-grams; \textcite{NichaniEtAl2024} and \textcite{BiettiEtAl2023} proved results on gradient dynamics for related tasks, and \textcite{MakkuvaEtAl2024} and \textcite{RajaramanEtAl2024} studied the loss landscape and the depth needed for Markov data. \Cref{thm:induction:bayes} is classical Dirichlet--multinomial conjugacy, and the add-$\frac12$ estimator and its $\frac{K-1}2\log n$ redundancy go back to \textcite{KrichevskyTrofimov1981}. The construction of \cref{thm:induction:circuit} follows the circuits of \textcite{ElhageEtAl2021} and \textcite{EdelmanEtAl2024}; the sink that carries the prior and the explicit error bound are, as far as we know, new here, as is the caricature with its escape theorem, which transfers the analysis of \cref{ch:deeplinear} to a composition of attention layers.

\Textcite{SanfordEtAl2024} (unrefereed) prove the separation between one and two layers by communication complexity. \Textcite{XiaoEtAl2024} introduced the term attention sink for the first tokens that absorb attention in language models. Function vectors and task vectors \parencite{ToddEtAl2024,HendelEtAl2023} were found independently and posted to arXiv simultaneously in October 2023.

\paragraph{Translating the literature.} \Textcite{ElhageEtAl2021} write $W_{QK}^h=W_Q^{h\top}W_K^h$ and $W_{OV}^h=W_O^hW_V^h$ for the circuits of head $h$; in our circuit layer~1 has $W_{QK}=\beta_1\times$(the shift on one-hot positions) and $W_{OV}$ copying the token into a fresh subspace, and layer~2 has $W_{QK}=\beta_2\times$(the identity from the token subspace of the query to the layer-1 subspace of the key) and $W_{OV}$ copying the token. Their pattern $[A][B]\dots[A]\to[B]$ is our $x_{s-1}=x_t=A$, $x_s=B$. Papers on in-context Markov chains commonly write $k$ for our number of states $K$ and $\alpha$ for the Dirichlet parameter that we call $\kappa$; the ``statistical induction head'' of \textcite{EdelmanEtAl2024} is our layer~2, and their unigram and bigram strategies are our unigram predictors and $\indpb_t$. \Textcite{OlssonEtAl2022} measure context use by the loss difference between tokens $500$ and $50$, which for the Bayes predictor of our setting is about $\frac{K(K-1)}2(\frac1{50}-\frac1{500})$ by \eqref{eq:induction:excess}. \Textcite{NichaniEtAl2024} call our disentangled layout a disentangled transformer. Our $\alpha$ is the initialization scale of \cref{ch:deeplinear}, and our $G$ plays the role of its singular value $s$.

\printchapterbib
\end{refsection}


\casepart{Training at Scale}{The last two models concern large steps and large models. A product of two numbers shows why gradient descent can run at the edge of stability, and power-law spectra and Zipf's law show where power-law scaling can come from.}{88}
% Chapter 15 of The Hydrogen Atoms of Deep Learning: the edge of stability.
\newcommand{\stabce}{c_{\mathrm{e}}}
\newcommand{\stabcf}{c_{4}}

\casechapter{The Edge of Stability}{Large step sizes and the catapult}{The two-parameter product model}{ch:stability}
\begin{refsection}[chapters/15-stability.bib]

\begin{chaptersummary}
Gradient descent on a quadratic diverges when the step size exceeds two over the largest curvature, yet networks trained by full-batch gradient descent spend most of their training at exactly that curvature, and large first steps make the loss spike and the curvature fall. The product of two numbers, $\Loss(u,v)=\frac12(uv-y)^2$, is the smallest model whose minimizers differ in curvature. We prove that gradient descent almost never converges to a minimizer sharper than $2/\eta$; that it leaves a sharp region only through steps of large loss, because gradient flow conserves $u^2-v^2$; and that an invariant ellipse makes $4/\eta$ the exact boundary between the catapult and divergence. A heuristic derivation, checked against the exact map, predicts where a small catapult lands. The model has no progressive sharpening, so it shows the catapult but not the sustained edge of stability.
\end{chaptersummary}

\section{The phenomenon: sharpening to the edge}\label{sec:stability:phenomenon}

Gradient descent on a quadratic with step size $\eta$ converges if and only if every curvature is below $2/\eta$; above that threshold the stiffest mode flips sign at every step and grows (\cref{ch:prologue}). The largest eigenvalue of the Hessian of the training loss,
\[
  S(\theta)=\lmax\bigl(\nabla^2\Lhat(\theta)\bigr),
\]
is called the \term{sharpness}\index{sharpness}, and the classical theory of optimization treats $\eta S<2$ as a constraint on the user: if the step size is below $2$ divided by the largest curvature met along the way, then the descent lemma guarantees that every step decreases the loss.\sidenote{Sharpness entered deep learning through generalization rather than optimization: \textcite{HochreiterSchmidhuber1997} argued that flat minima generalize better, and \textcite{KeskarEtAl2017} found that large-batch training lands in sharper minima and generalizes worse. This chapter is about optimization only.}

For neural networks the constraint describes neither practice nor what training does. \Textcite{CohenEtAl2021} measured the sharpness along full-batch gradient descent for fully connected, convolutional and transformer networks, with the square loss and the cross-entropy loss, and found the same two-stage pattern almost everywhere.

\begin{observation}[Progressive sharpening and the edge of stability; \textcite{CohenEtAl2021}]\label{obs:stability:eos}
Train a neural network by full-batch gradient descent with a fixed step size $\eta$. (i) While $S(\theta_t)<2/\eta$, the sharpness tends to increase. (ii) Once it reaches $2/\eta$, it stops increasing and hovers at or slightly above $2/\eta$ for the rest of training, while the training loss continues to decrease over long time scales, although not monotonically.
\end{observation}

The first stage is \term{progressive sharpening}\index{progressive sharpening}, the second the \term{edge of stability}\index{edge of stability}. The observation is robust for full-batch gradient descent on standard architectures and data sets, and its scope is known in three directions. Progressive sharpening is weaker for wider networks and with the cross-entropy loss, whose sharpness can also fall late in training, when the margins grow \parencite{CohenEtAl2021}. With mini-batches the sharpness settles below $2/\eta$, at a level that depends on the batch size \parencite{JastrzebskiEtAl2020}. With adaptive optimizers it is the preconditioned sharpness that sits at an edge \parencite{CohenEtAl2022}. \Cref{fig:stability:mlp} shows the observation in a network small enough to train in seconds.

\begin{figure}[t]
\begin{fullwidth}
\centering
\begin{tikzpicture}
\begin{groupplot}[group style={group size=2 by 1,horizontal sep=17mm},hplot,width=82mm,height=54mm,xmin=0,xmax=900]
\nextgroupplot[xlabel={time $\eta t$},ylabel={sharpness $S(\theta_t)$},ymin=0,ymax=15,legend style={at={(0.5,1.03)},anchor=south},legend columns=3]
\addplot[theory] table[x=time,y=sA] {figures/data/15-stability-mlp.dat};
\addplot[theory3] table[x=time,y=sB] {figures/data/15-stability-mlp.dat};
\addplot[theory2] table[x=time,y=sC] {figures/data/15-stability-mlp.dat};
\addplot[guide,sea,domain=0:900] {2/0.15};
\addplot[guide,sandline!55!black,domain=0:900] {2/0.25};
\addplot[guide,seadeep,domain=0:900] {2/0.4};
\legend{$\eta=0.15$,$\eta=0.25$,$\eta=0.4$}
\nextgroupplot[xlabel={time $\eta t$},ylabel={training loss $\Lhat(\theta_t)$},ymode=log,ymin=4e-4,ymax=2]
\addplot[theory] table[x=time,y=lA] {figures/data/15-stability-mlp.dat};
\addplot[theory3] table[x=time,y=lB] {figures/data/15-stability-mlp.dat};
\addplot[theory2] table[x=time,y=lC] {figures/data/15-stability-mlp.dat};
\end{groupplot}
\end{tikzpicture}
\end{fullwidth}
\caption{Progressive sharpening and the edge of stability in a small network. A tanh network with two hidden layers of width $64$ is trained by full-batch gradient descent on the square loss, on $500$ Gaussian inputs in $\R^{10}$ labelled by a random two-layer teacher (seeds fixed in the script \texttt{15-stability-mlp.py}), with three step sizes. Left: the sharpness, computed every $5$ units of time by power iteration with Hessian--vector products; the dashed lines are at $2/\eta$. All three runs follow one curve until the sharpness reaches their own $2/\eta$; the run with $\eta=0.15$ is still sharpening at the end. Right: the training loss keeps falling at the edge, with occasional spikes.}
\label{fig:stability:mlp}
\end{figure}

A second observation concerns the first steps of training with a large step size.

\begin{observation}[The catapult; \textcite{LewkowyczEtAl2020}]\label{obs:stability:catapult}
Let $\lambda_{\mathrm{init}}$ be the curvature at initialization. If $\eta<2/\lambda_{\mathrm{init}}$, wide networks train much as under gradient flow. If $2/\lambda_{\mathrm{init}}<\eta<\eta_{\max}$, the loss first grows for some steps, the curvature falls below $2/\eta$, and training then converges in a flatter region. Above $\eta_{\max}$ training diverges. In a two-layer linear model that Lewkowycz and coauthors solved, $\eta_{\max}=4/\lambda_{\mathrm{init}}$; for wide ReLU networks they measured $\eta_{\max}\approx12/\lambda_{\mathrm{init}}$.
\end{observation}

The \term{catapult}\index{catapult} is the name they gave to the loss spike that throws the parameters out of a sharp region.\sidenote{The idea that the step size selects among minima by their sharpness is older. \Textcite{WuMaE2018} showed that gradient descent can be linearly stable only at minima with $\eta S\le2$, and proposed this dynamical stability as the reason that gradient methods find flat minima. \Textcite{JastrzebskiEtAl2019} observed that the sharpness along the trajectory adapts to the step size early in training.} The two observations pose one question in two parts. The prologue's theorem says that gradient descent with $\eta S>2$ must diverge, and it does not. Where can gradient descent end up, given its step size, and how does it get there from a region that is too sharp for it? The model of this chapter answers both parts exactly, and the part of the phenomenon that it cannot reproduce, the sharpness that stays at $2/\eta$ for thousands of steps, points to what it leaves out.

\section{The model: a product of two numbers}\label{sec:stability:model}

\begin{definition}[The product model]\label[definition]{def:stability:model}
Fix a label $y>0$ and a step size $\eta>0$. The \term{product model}\index{product model} is the loss
\[
  \Loss(u,v)=\tfrac12\,(uv-y)^2,\qquad \theta=(u,v)\in\R^2,
\]
trained by gradient descent $\theta_{t+1}=G(\theta_t)$, where $G(u,v)=(u-\eta zv,\;v-\eta zu)$ and $z=uv-y$.
\end{definition}

The model is a linear network $f(x)=vux$ with one hidden unit, input weight $u$ and output weight $v$, trained on the single example $(x,y)=(1,y)$; the number $z$ is its residual. The gradient is $\nabla\Loss=z\,(v,u)$, and the Hessian is
\[
  \nabla^2\Loss(u,v)=\begin{pmatrix} v^2 & uv+z\\ uv+z & u^2\end{pmatrix}.
\]
The minimizers form the hyperbola $\cM=\{uv=y\}$, with one branch in the positive and one in the negative quadrant. At a minimizer the Hessian $\bigl(\begin{smallmatrix}v^2&y\\y&u^2\end{smallmatrix}\bigr)$ has the eigenvalue $0$ along the tangent $(u,-v)$ of $\cM$ and the eigenvalue $u^2+v^2$ along the normal $(v,u)$, so the sharpness of a minimizer is $u^2+v^2$.

Two quantities organize everything that follows: the \term{imbalance}\index{imbalance} $c=u^2-v^2$ and the trace $\tau=u^2+v^2=\tr\nabla^2\Loss$. Because $(u^2+v^2)^2=(u^2-v^2)^2+4(uv)^2$, they are related by
\begin{equation}\label{eq:stability:tau}
  \tau^2=c^2+4(z+y)^2 .
\end{equation}
On $\cM$ we have $z=0$, so the minimizer with imbalance $c$ has sharpness
\begin{equation}\label{eq:stability:Lambda}
  \Lambda(c)=\sqrt{c^2+4y^2}.
\end{equation}
Each branch of $\cM$ contains exactly one minimizer with each imbalance, the sharpness grows with $\abs c$, and the flattest minimizers are the balanced ones, $u=v=\pm\sqrt y$, with $\Lambda(0)=2y$. A quadratic analysis at a minimizer calls it stable for gradient descent when $\eta\Lambda(c)<2$, that is, when $\abs c<\stabce$, with the \term{edge imbalance}
\[
  \stabce=2\sqrt{\eta^{-2}-y^2}.
\]
Stable minimizers exist only when $\eta y<1$. Two more facts about the scales: away from $\cM$ the sharpness is not $\tau$ but $S(\theta)=\frac12\bigl(\tau+\sqrt{c^2+4(uv+z)^2}\bigr)$ (Exercise~\ref{ex:stability:hessian}); and the rescaling $(u,v)\mapsto(u,v)/\sqrt y$ turns the model with label $y$ and step $\eta$ into the model with label $1$ and step $\eta y$, so the dynamics depends on $\eta$ and $y$ only through the dimensionless step size $\eta y$ and on the initialization only through $\sqrt\eta\,(u,v)$.

\begin{dictionary}
the loss $\frac12(uv-y)^2$ & the square loss of a network on its training set\\
the hyperbola $\cM$ of minimizers & the manifold of parameters that interpolate the data\\
the imbalance $c=u^2-v^2$ & the differences $W_{\ell+1}^\top W_{\ell+1}-W_\ell W_\ell^\top$ between consecutive layers of a linear network, and their traces in a ReLU network, which gradient flow conserves\\
$\Lambda(c)$, the sharpness of the minimizer with imbalance $c$ & the sharpness of the minimizer that gradient flow would reach from the current parameters\\
the trace $\tau=u^2+v^2$ & the neural tangent kernel $\norm{\nabla_\theta f}^2$, the Gauss--Newton part of the Hessian\\
the residual $z$, flipping sign at every step & the component of the iterate along the top eigenvector of the Hessian at the edge of stability\\
the dimensionless step size $\eta y$ & the step size in units of the scale of the outputs
\end{dictionary}

\paragraph{What was thrown away.} Almost everything: the data (one example), the width and the depth (one hidden unit, two layers), the nonlinearity, and all directions of parameter space but two. What is kept is the smallest structure that the phenomenon needs. A quadratic has a constant Hessian, so its sharpness cannot respond to the dynamics; a one-parameter model has isolated minimizers, so there is nowhere flatter to go. Two parameters give a curve of minimizers along which the sharpness varies, and a Hessian that changes along the oscillation, because the third derivative of $\Loss$ is not zero. Both ingredients are present in every network with a continuum of interpolating parameters, which is every overparameterized network.

\paragraph{What one had to notice.} Gradient flow on this loss conserves $c$, so it lands at the minimizer with the imbalance it started with, sharp or not. Gradient descent does not conserve $c$, and the only thing that changes it is the residual: each step multiplies $c$ by $1-\eta^2z^2$. To invent the model one had to see that an oscillation of the residual, which a quadratic would treat as the prelude to divergence, is the mechanism that moves the parameters along the valley to a flatter place.

\section{The solution: where gradient descent can land}\label{sec:stability:landing}

We begin with the baseline, a quadratic loss, in the form we shall need. Let $\Loss(\theta)=\frac12(\theta-\theta_\star)^\top H(\theta-\theta_\star)$, where $H$ is symmetric positive semidefinite with eigenvalues $\lambda_1\ge\dots\ge\lambda_p\ge0$ and orthonormal eigenvectors $e_1,\dots,e_p$.

\begin{proposition}[Linear stability]\label[proposition]{prop:stability:quadratic}
Gradient descent with step size $\eta$ on this quadratic converges from $\theta_0$ if and only if $\ip{\theta_0-\theta_\star,e_k}=0$ for every $k$ with $\eta\lambda_k\ge2$. Consequently it converges from every initialization if $\eta\lambda_1<2$; if $\eta\lambda_1\ge2$, it converges only from a proper affine subspace of $\R^p$, which is a Lebesgue-null set.
\end{proposition}

\begin{proof}
By the normal-mode formula of \cref{ch:prologue}, $\ip{\theta_t-\theta_\star,e_k}=(1-\eta\lambda_k)^t\ip{\theta_0-\theta_\star,e_k}$. A geometric sequence with a nonzero first term converges if and only if its ratio lies in $(-1,1]$, which here means $0\le\eta\lambda_k<2$. So the iterates converge exactly when the components with $\eta\lambda_k\ge2$ vanish, that is, when $\theta_0$ lies in $\theta_\star+\operatorname{span}\{e_k\colon\eta\lambda_k<2\}$, an affine subspace that is proper when $\eta\lambda_1\ge2$.
\end{proof}

At every minimizer of the product model the Hessian has the eigenvalue $\Lambda(c)$ in the direction normal to $\cM$, so the quadratic approximation predicts that gradient descent is stable near the minimizer exactly when $\eta\Lambda(c)<2$. To see what gradient descent actually does we need its exact form, which is simple in the right coordinates.

\begin{lemma}[Gradient descent in product coordinates]\label[lemma]{lem:stability:coordinates}
Let $\theta=(u,v)$ and $\theta'=G(\theta)$, and write $a=u+v$ and $b=u-v$, so that $c=ab$. Then
\begin{enumerate}
\item $a'=(1-\eta z)\,a$ and $b'=(1+\eta z)\,b$;
\item $c'=(1-\eta^2z^2)\,c=\bigl(1-2\eta^2\Loss(\theta)\bigr)\,c$;
\item $z'=z\bigl(1-\eta\tau+\eta^2z(z+y)\bigr)$;
\item $\tau'=\tau+\eta z^2(\eta\tau-4)-4\eta yz$.
\end{enumerate}
\end{lemma}

\begin{proof}
Adding and subtracting the two coordinates of $G(\theta)$ gives (i): $u'+v'=(u+v)-\eta z(v+u)$ and $u'-v'=(u-v)+\eta z(u-v)$. Multiplying the two identities of (i) gives (ii). For (iii) and (iv), expand $u'v'=(u-\eta zv)(v-\eta zu)=uv-\eta z\tau+\eta^2z^2uv$ and $u'^2+v'^2=\tau-4\eta zuv+\eta^2z^2\tau$, and substitute $uv=z+y$.
\end{proof}

Each identity carries a message. By (i), the sum and the difference of the weights evolve by multiplication. By (ii), the imbalance changes only through the loss: the factor $\abs{1-2\eta^2\Loss(\theta)}$ is less than $1$ whenever $0<\Loss(\theta)<1/\eta^2$, so $\abs c$ can grow only at steps with $\Loss(\theta)>1/\eta^2$. By (iii), the residual is multiplied at each step by a factor that, near $\cM$, is close to $1-\eta\Lambda(c)$: it decays when $\eta\Lambda(c)<2$ and flips sign and grows when $\eta\Lambda(c)>2$. By (iv), during a large oscillation the trace falls if $\eta\tau<4$ and rises if $\eta\tau>4$; for $y=0$ this is the whole update, and it is the equation from which \textcite{LewkowyczEtAl2020} read off their threshold $4/\lambda_{\mathrm{init}}$. In the limit of small steps the factor in (ii) tends to $1$, and the imbalance is conserved.

\begin{proposition}[Gradient flow conserves the imbalance]\label[proposition]{prop:stability:flow}
Along gradient flow $\dot\theta=-\nabla\Loss(\theta)$ the imbalance $c$ is constant. If $u(0)+v(0)\ne0$, then $\theta(t)$ converges to the minimizer with imbalance $c(0)$ on the branch of $\cM$ on which $u+v$ has the sign of $u(0)+v(0)$, and the sharpness of the limit is $\Lambda(c(0))$. If $u(0)+v(0)=0$, then $\theta(t)$ converges to the origin, a saddle point.
\end{proposition}

\begin{proof}
The plan is to show that $c$ is constant and $z$ decays, and then to identify the limit. Differentiating, $\dot c=2u\dot u-2v\dot v=-2uzv+2vzu=0$, $\dot z=\dot uv+u\dot v=-z\tau$, $\dot a=-za$ and $\dot b=zb$. Suppose first that $c\ne0$. Then $\tau\ge\abs c$ by \eqref{eq:stability:tau}, so $\abs{z(t)}\le\abs{z(0)}e^{-\abs ct}$, and $\tau(t)=(c^2+4(z(t)+y)^2)^{1/2}$ stays bounded. The speed $\abs{\dot\theta}=\abs z\sqrt\tau$ is therefore integrable, so $\theta(t)$ has a limit, at which $z=0$ and $u^2-v^2=c$. The two minimizers with imbalance $c$ differ in the sign of $a$, and $a(t)=a(0)\exp(-\int_0^tz)$ keeps its sign. Suppose next that $c=0$ and $a(0)\ne0$. Then $b=c/a=0$ for all $t$, so $u=v=a/2$, $z=a^2/4-y$, and the scalar equation $\dot a=-(a^2/4-y)a$ drives $a$ to $\pm2\sqrt y$, a balanced minimizer. Finally, if $a(0)=0$, then $a$ stays $0$, so $z=-b^2/4-y\le-y$ and $\dot b=zb$ gives $\abs{b(t)}\le\abs{b(0)}e^{-yt}$.
\end{proof}

So gradient flow ends wherever its imbalance points it, and when $\eta\Lambda(c(0))>2$ it ends at a minimizer where gradient descent is linearly unstable. To show that gradient descent almost never ends there, we need one fact from measure theory.

\begin{lemma}[Preimages of null sets]\label[lemma]{lem:stability:null}
Let $P\colon\R^n\to\R^n$ be a polynomial map whose Jacobian determinant $J=\det DP$ is not the zero polynomial. Then $P^{-1}(N)$ is a Lebesgue-null set for every null set $N\subseteq\R^n$.
\end{lemma}

\begin{proof}
The plan is to discard the zero set of $J$, which is null, and to cover the rest by countably many open sets on which $P$ has a $C^1$ inverse, since such an inverse maps null sets to null sets. We use three facts.

\emph{(a) The zero set of a nonzero polynomial $q$ on $\R^n$ is null.} For $n=1$ it is finite. For $n>1$, write $q(x',x_n)=\sum_kq_k(x')x_n^k$ with some $q_k$ not the zero polynomial. For $x'$ outside the zero set of $q_k$, which is null by induction, $x_n\mapsto q(x',x_n)$ is a nonzero polynomial with finitely many zeros, and Fubini's theorem shows that the zero set of $q$ is null.

\emph{(b) A $C^1$ map $F$ from an open set $W\subseteq\R^n$ into $\R^n$ maps null sets to null sets.} Write $W$ as a countable union of closed cubes. On each such cube $Q$ the map $F$ is Lipschitz with constant $K_Q=\max_Q\norm{DF}_{\op}$, by the mean value inequality on the convex set $Q$. Let $N\subseteq W$ be null and $\delta>0$, and cover $N\cap Q$ by cubes $C_i$ of total volume less than $\delta$. Each $F(C_i\cap Q)$ has diameter at most $K_Q\sqrt n$ times the side of $C_i$, so it lies in a cube of volume $(2K_Q\sqrt n)^n\operatorname{vol}(C_i)$, and $F(N\cap Q)$ is covered by cubes of total volume less than $(2K_Q\sqrt n)^n\delta$. Hence $F(N\cap Q)$ is null, and so is $F(N)$, the countable union of these sets.

\emph{(c) Local inverses.} Let $U=\{J\ne0\}$. By the inverse function theorem each point of $U$ has an open neighbourhood $W\subseteq U$ that $P$ maps diffeomorphically onto an open set. Because $\R^n$ has a countable base, countably many such sets $W_1,W_2,\dots$ cover $U$; write $\psi_i$ for the inverse of $P$ on $W_i$.

Now $P^{-1}(N)\subseteq\{J=0\}\cup\bigcup_i\psi_i\bigl(N\cap P(W_i)\bigr)$. The first set is null by (a), and each $\psi_i(N\cap P(W_i))$ is null by (b).
\end{proof}

\begin{theorem}[Where gradient descent can land]\label{thm:stability:landing}
Let $y>0$ and $\eta>0$, and let $\theta_t=G^t(\theta_0)$ be gradient descent on the product model.
\begin{enumerate}
\item For every $\theta_0$ and $t$, the imbalance is $c_t=c_0\prod_{s<t}\bigl(1-2\eta^2\Loss(\theta_s)\bigr)$.
\item For every $\theta_0$ outside a Lebesgue-null set, either the iterates do not converge, or they converge to a minimizer $\theta_\infty$ with $\eta S(\theta_\infty)\le2$, that is, with $\abs{c_\infty}\le\stabce$.
\item Suppose that $\eta y\le1$ and $\eta\Lambda(c_0)>2$. If $\theta_0$ lies outside the null set of (ii) and the iterates converge, then
\[
  \prod_{t\ge0}\bigl\lvert1-2\eta^2\Loss(\theta_t)\bigr\rvert=\frac{\abs{c_\infty}}{\abs{c_0}}\le\frac{\stabce}{\abs{c_0}}<1 .
\]
\end{enumerate}
\end{theorem}

\begin{proof}
Part (i) follows from \cref{lem:stability:coordinates}(ii) by induction on $t$.

Part (ii) carries the weight. The plan is to show that convergence to a fixed point at which gradient descent is linearly unstable forces the iterates to land exactly on one of two curves after finitely many steps, and that the initializations whose orbits ever meet those curves form a null set. If $\theta_t\to\theta_\infty$, then $\theta_\infty=G(\theta_\infty)$ by continuity, so $z_\infty\cdot(v_\infty,u_\infty)=0$: the limit is a minimizer or the origin.

\emph{A sharp minimizer.} Let $\theta_\infty\in\cM$ with $\eta S(\theta_\infty)>2$. By \cref{lem:stability:coordinates}(iii), $z_{t+1}=m_tz_t$ with $m_t=1-\eta\tau_t+\eta^2z_t(z_t+y)$. As $t\to\infty$, $z_t\to0$ and $\tau_t\to\tau(\theta_\infty)=S(\theta_\infty)$, so $m_t\to1-\eta S(\theta_\infty)<-1$, and there is a $T$ with $\abs{m_t}\ge1$ for all $t\ge T$. If $z_T$ were nonzero, then $\abs{z_t}\ge\abs{z_T}>0$ for all $t\ge T$, contradicting $z_t\to0$. Hence $z_T=0$, that is, $\theta_T\in\cM$.

\emph{The origin.} Let $\theta_\infty=0$. By \cref{lem:stability:coordinates}(i), $a_{t+1}=(1-\eta z_t)a_t$, and $1-\eta z_t\to1+\eta y>1$ because $z_t\to-y$. The landing argument for a sharp minimizer, applied to $a_t$ in place of $z_t$, gives $a_T=0$ for some $T$: the iterate $\theta_T$ lies on the line $\cD=\{u+v=0\}$.

\emph{The exceptional set is null.} The curves $\cM$ and $\cD$ are null, being graphs of continuous functions. The map $G$ is polynomial, and its Jacobian determinant $\det\bigl(I-\eta\nabla^2\Loss(u,v)\bigr)$ equals $1-\eta^2y^2-\eta u^2$ at $v=0$, so it is not the zero polynomial. By \cref{lem:stability:null} and induction on $T$, each set $G^{-T}(\cM\cup\cD)$ is null, and so is their union over $T\ge0$, which contains every initialization whose iterates converge to the origin or to a minimizer with $\eta S>2$. At the remaining limits $\eta S(\theta_\infty)\le2$, and since $S=\Lambda(c)$ on $\cM$, this means $\abs{c_\infty}\le\stabce$.

For part (iii), part (ii) says that the limit is a minimizer with $\abs{c_\infty}\le\stabce$. By part (i) and the continuity of $c=u^2-v^2$, the partial products $\prod_{s<t}\abs{1-2\eta^2\Loss(\theta_s)}=\abs{c_t}/\abs{c_0}$ converge to $\abs{c_\infty}/\abs{c_0}$. Finally, $\eta\Lambda(c_0)>2$ means $\abs{c_0}>\stabce$.
\end{proof}

If $\eta y>1$, then every minimizer has $\eta S\ge2\eta y>2$, and part (ii) says that gradient descent from almost every initialization never converges. Part (ii) is a nonlinear version of \cref{prop:stability:quadratic}: the minimizers too sharp for the step size are reached only from a null set. For general smooth losses the same conclusion follows from the stable manifold theorem when the step size is small enough that $G$ is a diffeomorphism \parencite{LeeEtAl2016,LeeEtAl2019}; here no such condition is needed, because $G$ is polynomial and the landing argument replaces the stable manifold theorem.

Parts (i) and (iii) say why the catapult is forced. Gradient flow from $\theta_0$ would end at sharpness $\Lambda(c_0)$. If $\eta\Lambda(c_0)>2$, gradient descent cannot end there, and the only way it can move its destination is through the factors $1-2\eta^2\Loss(\theta_t)$, which differ from $1$ only when the loss is not small. A run that starts near a sharp minimizer and converges must pass through a phase of large loss. \Cref{lem:stability:coordinates}(ii) also fixes the direction of travel: while $0<\Loss(\theta_t)<1/\eta^2$, the sharpness $\Lambda(c_t)$ of the flow's destination strictly decreases.\sidenote{\Textcite{KreislerEtAl2023} call $\Lambda(c_t)$ the sharpness of the gradient flow solution, and prove that it decreases monotonically along gradient descent for scalar linear networks of every depth, under conditions on the step size and the initialization.} Below the edge, by contrast, gradient descent lands almost where the flow lands.

\begin{proposition}[Below the edge, gradient descent follows the flow]\label[proposition]{prop:stability:lazy}
Let $\eta y<1$ and let $\theta_0$ satisfy $\eta\Lambda(c_0)<2$. Put $q=\max\{\eta\Lambda(c_0)-1,\,1-2\eta y\}$, which lies in $[0,1)$, and suppose that $\abs{z_0}\le(1-q)/(8\eta)$. Then $\abs{z_t}\le\bigl(\frac{1+q}2\bigr)^t\abs{z_0}$ for every $t$, $\abs{c_t}$ is nonincreasing, the iterates converge to a minimizer, and
\[
  \abs{c_0}\ \ge\ \abs{c_\infty}\ \ge\ \abs{c_0}\Bigl(1-\frac{2\eta^2z_0^2}{1-q}\Bigr).
\]
\end{proposition}

\begin{proof}
The plan is an induction showing that the residual contracts by the factor $\rho=(1+q)/2$ at every step, which keeps the iterates where the imbalance cannot grow. The number $q$ is at least $\abs{1-2\eta y}\ge0$ because $\Lambda(c_0)\ge2y$, and it is less than $1$ because $\eta\Lambda(c_0)<2$ and $\eta y>0$.

Suppose that $\abs{c_t}\le\abs{c_0}$ and $\abs{z_t}\le\abs{z_0}$, as holds at $t=0$. Compare the factor $m_t$ of \cref{lem:stability:coordinates}(iii) with $m_t^0=1-\eta\Lambda(c_t)$. Since $2y\le\Lambda(c_t)\le\Lambda(c_0)$, we have $1-\eta\Lambda(c_0)\le m_t^0\le1-2\eta y$, so $\abs{m_t^0}\le q$. By \eqref{eq:stability:tau}, $\tau_t$ and $\Lambda(c_t)$ are the lengths of the vectors $(c_t,2(z_t+y))$ and $(c_t,2y)$, which differ by $(0,2z_t)$, so $\abs{\tau_t-\Lambda(c_t)}\le2\abs{z_t}$. Hence, using $\eta\abs{z_t}\le\eta\abs{z_0}\le\frac18$ and $\eta y<1$ in the last step,
\[
  \abs{m_t-m_t^0}\le\eta\abs{\tau_t-\Lambda(c_t)}+\eta^2\abs{z_t}\bigl(\abs{z_t}+y\bigr)\le\eta\abs{z_t}\bigl(2+\eta\abs{z_t}+\eta y\bigr)<4\eta\abs{z_t}\le\frac{1-q}{2}.
\]
Therefore $\abs{m_t}<q+\frac{1-q}2=\rho$, and $\abs{z_{t+1}}\le\rho\abs{z_t}\le\abs{z_0}$. Moreover $\eta^2z_t^2\le\frac1{64}$, so the factor $1-\eta^2z_t^2$ of \cref{lem:stability:coordinates}(ii) lies in $[0,1]$ and $\abs{c_{t+1}}\le\abs{c_t}\le\abs{c_0}$. This completes the induction.

Consequently $\abs{z_t}\le\rho^t\abs{z_0}$. The step length $\abs{\theta_{t+1}-\theta_t}=\eta\abs{z_t}\sqrt{\tau_t}$ is at most a constant times $\rho^t$, since $\tau_t\le(c_0^2+4(\abs{z_0}+y)^2)^{1/2}$, so $(\theta_t)$ is a Cauchy sequence, and its limit has $z=0$. Finally, the factors $1-\eta^2z_t^2$ lie in $[0,1]$, and the inequality $\prod_t(1-x_t)\ge1-\sum_tx_t$ for numbers $x_t\in[0,1]$ gives
\[
  \abs{c_\infty}\ge\abs{c_0}\Bigl(1-\sum_{t\ge0}\eta^2z_t^2\Bigr)\ge\abs{c_0}\Bigl(1-\frac{\eta^2z_0^2}{1-\rho^2}\Bigr)\ge\abs{c_0}\Bigl(1-\frac{2\eta^2z_0^2}{1-q}\Bigr),
\]
because $1-\rho^2\ge1-\rho=(1-q)/2$.
\end{proof}

\section{The solution: the catapult}\label{sec:stability:catapult}

\Cref{thm:stability:landing} says where gradient descent can end, not whether it gets there. Start at a minimizer with $\eta\Lambda(c_0)$ above $2$ and perturb it slightly. The residual flips sign and grows, and every step with nonzero loss shrinks $\abs c$, as long as the loss stays below $1/\eta^2$. Does the shrinking win before the oscillation runs away? The answer has an exact part, which locates the boundary between the catapult and divergence, and an approximate part, which predicts where a small catapult lands.

\begin{theorem}[The catapult ellipse]\label{thm:stability:ellipse}
Let $0<\eta y<2$, and define
\[
  \Psi(\theta)=\eta\,(u^2+v^2)-\eta^2y\,(uv-y)-4=\eta\tau-\eta^2yz-4 .
\]
Then for every $\theta$,
\begin{equation}\label{eq:stability:psi}
  \Psi\bigl(G(\theta)\bigr)=\bigl(1+\eta^2z(z+y)\bigr)\,\Psi(\theta),\qquad\text{and}\qquad 1+\eta^2z(z+y)\ge1-\tfrac14\eta^2y^2>0 .
\end{equation}
Consequently, with $\cE=\{\Psi=0\}$ and $\cK=\{\Psi<0\}$:
\begin{enumerate}
\item the map $G$ sends $\cE$, $\cK$ and the exterior $\{\Psi>0\}$ into themselves;
\item $\cK$ is the interior of an ellipse with axes along the diagonals $u=\pm v$, and every iterate started in $\cK$ satisfies $\eta\abs{z_t}<2$, that is, $\Loss(\theta_t)<2/\eta^2$;
\item a minimizer $\theta$ lies in $\cK$ if and only if $\eta S(\theta)<4$, and on $\cE$ if and only if $\eta S(\theta)=4$;
\item for every $\theta_0$ with $\Psi(\theta_0)\ge0$ outside the null set of \cref{thm:stability:landing}, the iterates do not converge.
\end{enumerate}
\end{theorem}

\begin{proof}
The plan is to verify the identity with \cref{lem:stability:coordinates} and then read off the consequences. Write $Q=1+\eta^2z(z+y)$. By \cref{lem:stability:coordinates}(iv), $\eta\tau'-4=(1+\eta^2z^2)(\eta\tau-4)-4\eta^2yz$, and by \cref{lem:stability:coordinates}(iii), $\eta^2yz'=\eta^2yz(1-\eta\tau)+\eta^4yz^2(z+y)$. Subtract the second identity from the first, and collect the terms that contain $\eta\tau-4$ (among them $\eta^2yz\cdot\eta\tau-4\eta^2yz$):
\[
  \Psi(\theta')=(1+\eta^2z^2+\eta^2yz)(\eta\tau-4)-\eta^2yz\bigl(1+\eta^2z(z+y)\bigr)=Q\,(\eta\tau-4)-Q\,\eta^2yz=Q\,\Psi(\theta).
\]
Completing the square, $Q=(\eta z+\frac12\eta y)^2+1-\frac14\eta^2y^2$, which is positive because $\eta y<2$.

(i) By induction, $\Psi(\theta_t)=\Psi(\theta_0)\prod_{s<t}Q(\theta_s)$, and every factor is positive.

(ii) In the coordinates $a=u+v$ and $b=u-v$ we have $u^2+v^2=\frac12(a^2+b^2)$ and $uv=\frac14(a^2-b^2)$, so the inequality $\Psi<0$ reads
\[
  \frac{\eta a^2}{4(2+\eta y)}+\frac{\eta b^2}{4(2-\eta y)}<1 ,
\]
the interior of an ellipse. For the bound on $z$, use $\tau\ge2\abs{uv}=2\abs{z+y}$. If $z\ge0$, then $\Psi<0$ gives $2\eta(z+y)<4+\eta^2yz$, that is, $\eta z\,(2-\eta y)<2\,(2-\eta y)$, so $\eta z<2$. If $z<-y$, then it gives $-2\eta(z+y)<4+\eta^2yz$, that is, $-\eta z\,(2+\eta y)<2\,(2+\eta y)$, so $\eta z>-2$. If $-y\le z<0$, then $\eta\abs z\le\eta y<2$.

(iii) At a minimizer $z=0$ and $\tau=S$, so $\Psi=\eta S-4$.

(iv) If $\Psi(\theta_0)\ge0$ and $\theta_t\to\theta_\infty$, then $\Psi(\theta_\infty)\ge0$ by (i) and continuity. The limit is a fixed point of $G$, so it is the origin, where $\Psi=\eta^2y^2-4<0$, or a minimizer, which by (iii) has $\eta S(\theta_\infty)\ge4>2$. The first is impossible, and \cref{thm:stability:landing}(ii) confines the second to the null set.
\end{proof}

Together with \cref{thm:stability:landing}, the ellipse gives an exact phase diagram for gradient descent started near a minimizer $\theta_\star$ with imbalance $c_0$, for $\eta y<1$ (\cref{fig:stability:ellipse}). If $\eta\Lambda(c_0)<2$, gradient descent converges close to $\theta_\star$ (\cref{prop:stability:lazy}). If $2<\eta\Lambda(c_0)<4$, then $\theta_\star$ lies in $\cK$, and so does a neighbourhood of it; gradient descent started there never leaves $\cK$, its loss stays below $2/\eta^2$, and for almost every start, if it converges, then it lands at a minimizer with $\eta S\le2$, below the edge. If $\eta\Lambda(c_0)>4$, then from almost every start near $\theta_\star$ it never converges. The threshold $4/\lambda_{\mathrm{init}}$ of \cref{obs:stability:catapult} is therefore exact in the product model for every $\eta y<2$, and the same identity holds for the width-$m$ model of \textcite{LewkowyczEtAl2020} with $\eta^2/m$ in place of $\eta^2$ (Exercise~\ref{ex:stability:width}).\sidenote{We found the identity \eqref{eq:stability:psi} while looking for the reason why, in simulations, the boundary between catapult and divergence sat at $\eta S=4$ to four digits for every $\eta y$. We have not seen it stated in the literature.} Two things the theorem does not say should be said. That gradient descent started in $\cK$ actually converges is a numerical observation (every one of $10^5$ random initializations in $\cK$ converged for $\eta y\in\{0.1,0.4,0.7,0.9\}$), not a theorem. And that it diverges outside is again numerical, except for $y=0$, where it is an exercise (Exercise~\ref{ex:stability:yzero}).

The ellipse is also where the dynamics is most complicated. In the coordinates $a=2\sqrt{(2+\eta y)/\eta}\,\sin\varphi$ and $b=2\sqrt{(2-\eta y)/\eta}\,\cos\varphi$ on $\cE$, gradient descent is the map $\varphi\mapsto\pi-3\varphi$ of the circle (Exercise~\ref{ex:stability:tripling}), an expanding map with a dense set of periodic orbits, whose only fixed points are the four minimizers with $\eta S=4$. Inside $\cK$, large oscillations behave chaotically for a while before they die out, and the landing point of a large catapult depends sensitively on where it started (\cref{fig:stability:catapult}).

\begin{figure}[t]
\begin{fullwidth}
\centering
\begin{tikzpicture}
\begin{groupplot}[group style={group size=2 by 1,horizontal sep=18mm},hplot]
\nextgroupplot[width=74mm,height=74mm,axis equal image,xlabel={$u$},ylabel={$v$},xmin=-4.3,xmax=4.3,ymin=-4.3,ymax=4.3,xtick={-4,-2,0,2,4},ytick={-4,-2,0,2,4}]
\addplot[muted,line width=0.6pt,densely dotted,domain=0.2326:0.3178,samples=20] {1/x};
\addplot[muted,line width=0.6pt,densely dotted,domain=3.1463:4.3,samples=20] {1/x};
\addplot[muted,line width=0.6pt,densely dotted,domain=-0.3178:-0.2326,samples=20] {1/x};
\addplot[muted,line width=0.6pt,densely dotted,domain=-4.3:-3.1463,samples=20] {1/x};
\addplot[sea,line width=1.2pt,domain=0.3178:0.4569,samples=20] {1/x};
\addplot[sea,line width=1.2pt,domain=2.1889:3.1463,samples=40] {1/x};
\addplot[sea,line width=1.2pt,domain=-0.4569:-0.3178,samples=20] {1/x};
\addplot[sea,line width=1.2pt,domain=-3.1463:-2.1889,samples=40] {1/x};
\addplot[sandline!55!black,line width=2.2pt,domain=0.4569:2.1889,samples=60] {1/x};
\addplot[sandline!55!black,line width=2.2pt,domain=-2.1889:-0.4569,samples=60] {1/x};
\addplot[seadeep,line width=1.0pt,domain=0:360,samples=181] ({sqrt(6)*sin(x)+2*cos(x)},{sqrt(6)*sin(x)-2*cos(x)});
\addplot[only marks,mark=*,mark size=0.55pt,color=seadeep] table[x=u,y=v] {figures/data/15-stability-ellipse-in.dat};
\addplot[muted,line width=0.3pt,mark=square*,mark size=0.9pt,mark options={fill=muted}] table[x=u,y=v] {figures/data/15-stability-ellipse-out.dat};
\node[font=\footnotesize,text=seadeep,anchor=south west] at (axis cs:-3.6,2.9) {$\cE$};
\node[font=\footnotesize,text=ink,anchor=west] at (axis cs:1.2,1.6) {$\eta S\le2$};
\nextgroupplot[width=74mm,height=74mm,xlabel={$\eta y$},ylabel={$\eta\Lambda(c_0)$},xmin=0,xmax=2,ymin=0,ymax=5,xtick={0,0.5,1,1.5,2},ytick={0,1,2,3,4,5}]
\fill[seafoam] (0,0) -- (1,2) -- (0,2) -- cycle;
\fill[seamist] (0,2) -- (1,2) -- (1,4) -- (0,4) -- cycle;
\fill[sand] (1,2) -- (2,4) -- (1,4) -- cycle;
\fill[muted!12] (0,0) -- (2,4) -- (2,0) -- cycle;
\draw[seadeep,line width=0.9pt] (0,4) -- (2,4);
\draw[seadeep,line width=0.9pt,dashed] (0,2) -- (1,2);
\draw[muted,line width=0.6pt] (0,0) -- (2,4);
\draw[sandline!55!black,line width=0.6pt,densely dotted] (0.4,0.8) -- (0.4,4);
\node[font=\footnotesize,text=ink] at (axis cs:0.2,1.65) {lazy};
\node[font=\footnotesize,text=ink] at (axis cs:0.7,3.0) {catapult};
\node[font=\footnotesize,text=ink,align=center] at (axis cs:1.32,3.45) {oscillates};
\node[font=\footnotesize,text=ink] at (axis cs:1.0,4.5) {diverges};
\node[font=\footnotesize,text=muted,align=center] at (axis cs:1.45,1.0) {no minimizer};
\end{groupplot}
\end{tikzpicture}
\end{fullwidth}
\caption{The geometry of the product model. Left: the plane of parameters for $y=1$ and $\eta=0.4$. The hyperbola $uv=1$ is drawn thick where its minimizers are stable ($\eta S\le2$), solid inside the catapult ellipse $\cE$ ($2<\eta S<4$), and dotted outside it. Circles: a catapult, gradient descent started next to the minimizer with $\eta S=2.3$; it zigzags across the valley and lands at $\eta S=1.67$. Squares: a run started next to the minimizer with $\eta S=4.3$, outside $\cE$, which leaves the picture after eight steps. Right: the phase diagram for a start near a minimizer, in the dimensionless step size $\eta y$ and the start's sharpness; the lines $\eta\Lambda=2$ (dashed) and $\eta\Lambda=4$ are the exact boundaries of \cref{thm:stability:landing,thm:stability:ellipse}, no minimizer is flatter than $2\eta y$, and in the band marked ``oscillates'' no minimizer is stable but the iterates stay bounded. The dotted line marks $\eta y=0.4$, used in \cref{fig:stability:catapult,fig:stability:trajectory}.}
\label{fig:stability:ellipse}
\end{figure}

How far does a catapult throw? Near the edge the answer comes from a slow drift of the sharpness through a period-doubling instability, and it is only derived.\sidenote{The expansion below is the normal form of a flip (period-doubling) bifurcation, here with a parameter that drifts because of the oscillation itself \parencite{Kuznetsov2004}.}

\begin{heuristic}[Heuristic derivation of the reflection law]
Let $\eta y<1$, and start near a minimizer just beyond the edge: $\eps_0=\eta\Lambda(c_0)-2$ is small and positive, and $z_0$ is much smaller than $\eps_0$. Follow two numbers, the squared residual $A_t=z_t^2=2\Loss(\theta_t)$ and the excess $\eps_t=\eta\Lambda(c_t)-2$ of the sharpness of the flow's destination over the edge. Both stay small, $A$ of order $\eps_0^2$ and $\eps$ of order $\eps_0$, and in each recursion we keep the leading terms.

\emph{The fast variable.} Expand $\tau=(\Lambda(c)^2+8yz+4z^2)^{1/2}$ in powers of $z$ and substitute it in \cref{lem:stability:coordinates}(iii), with $\eta\Lambda(c)=2+\eps$ and with the coefficients of $z^2$ and $z^3$ evaluated at the edge:
\[
  z'=-(1+\eps)\,z-\eta^2y\,z^2+\eta^4y^2z^3+\cdots .
\]
The residual flips sign at every step, and its size is multiplied by $1+\eps$. The quadratic term changes sign with $z$ and nearly cancels over two steps: composing $g(z)=-(1+\eps)z+k_2z^2+k_3z^3$ with itself gives $g(g(z))=(1+\eps)^2z-2(k_2^2+k_3)z^3$ plus terms of order $\eps z^2$ and $z^4$. Here $k_2=-\eta^2y$ and $k_3=\eta^4y^2$, so $k_2^2+k_3=2\eta^4y^2$, and over one step, on average,
\[
  A'=(1+2\eps)A-\beta A^2+\cdots,\qquad \beta=4\eta^4y^2 .
\]

\emph{The slow variable.} By \cref{lem:stability:coordinates}(ii), $c'=(1-\eta^2A)c$, and $d\eps/dc=\eta c/\Lambda(c)=\eta^2c/(2+\eps)$. At the edge $c^2=\stabce^2$, so
\[
  \eps'=\eps-\nu A+\cdots,\qquad \nu=\tfrac12\eta^4\stabce^2=2\eta^2(1-\eta^2y^2) .
\]

\emph{The flow.} In one step $A$ changes by a factor $1+O(\eps)$ and $\eps$ by $O(A)$, so both vary slowly, and we replace the recursions by differential equations in the step count:
\begin{equation}\label{eq:stability:slowfast}
  \frac{dA}{dt}=2\eps A-\beta A^2,\qquad \frac{d\eps}{dt}=-\nu A .
\end{equation}
The oscillation grows while $\eps>0$ and decays once $\eps<0$, and its growth is what drives $\eps$ down. Relative to $2\eps A$, the term $\beta A^2$ is of order $\beta\eps_0/\nu$, a correction of relative order $\eps_0$, and to leading order we drop it. Then $\frac{d}{dt}(\nu A+\eps^2)=2\nu\eps A-2\eps\nu A=0$, so $\nu A+\eps^2$ is constant during the catapult. At the start $A\approx0$ and $\eps=\eps_0$; at the end $A\to0$ again, and $\eps$, which only decreases, has become negative. Hence $\eps_\infty=-\eps_0$, which is the \term{reflection law}\index{reflection law}
\begin{equation}\label{eq:stability:reflection}
  \eta S(\theta_\infty)=4-\eta\Lambda(c_0).
\end{equation}
The same first integral gives the height of the loss spike: $A$ is largest when $\eps=0$, where $\nu A=\eps_0^2$, so the largest loss is about $\eps_0^2/(2\nu)=\eps_0^2/\bigl(4\eta^2(1-\eta^2y^2)\bigr)$.
\end{heuristic}

\begin{remark}[The next order]\label[remark]{rem:stability:second}
Each approximation in the derivation (dropping $\beta A^2$, replacing the growth factor $(1+\eps)^2$ by $1+2\eps$, freezing $d\eps/dc$ at its value at the edge, and replacing the recursions by a flow) changes the dynamics by a relative amount of order $\eps$, and so each moves $\eps_\infty$ at order $\eps_0^2$. Carrying all four to the next order (Exercise~\ref{ex:stability:second}) gives
\begin{equation}\label{eq:stability:second}
  \eps_\infty=-\eps_0+\Bigl(\frac{2\eta^2y^2}{1-\eta^2y^2}-\frac23\Bigr)\eps_0^2+O(\eps_0^3).
\end{equation}
With $s=\beta/\nu=2\eta^2y^2/(1-\eta^2y^2)$, and with $\log A$ and $\log c$ as the variables, the four contributions are $\frac23s$ from the saturation, $\frac13$ from the growth factor, $\frac13(1+s)$ from the curvature of $\eps$ as a function of $\log c$, and $-\frac43$ from the discreteness: the map updates $A$ and $c$ simultaneously, as the explicit Euler method does, and the first integral drifts. The exact map confirms \eqref{eq:stability:second} to six digits for every $\eta y$ from $0.05$ to $0.8$ (\cref{fig:stability:catapult}, right). Keeping the saturation alone turns \eqref{eq:stability:slowfast} into a system solvable with the Lambert $W$ function, and gives the coefficient $\frac23s$, which has the wrong sign whenever $\eta y<\frac12$.
\end{remark}

\section{Reading the solution}\label{sec:stability:reading}

\paragraph{The catapult curve.} \Cref{fig:stability:catapult} plots where gradient descent lands against where it starts, for $\eta y=0.4$. Below the edge the landing sharpness equals the starting one, as \cref{prop:stability:lazy} requires. Just above the edge the landing point is the mirror image of the start across $\eta S=2$: a start at $\eta S=2.08$ lands at $1.918$, where the reflection law \eqref{eq:stability:reflection} predicts $1.92$ and the second-order law \eqref{eq:stability:second} $1.9182$. The reflection is a statement about a first integral, and it has a physical reading: the oscillation is an instability that feeds on the excess sharpness and destroys it, and since it must die out exactly as it grew, it overshoots the edge by as much as it started above it.\sidenote{Lotka and Volterra's predators and prey, and the epidemics of \textcite{KermackMcKendrick1927}, have first integrals of the same kind; the reflection law is the analogue of the final-size relation of an epidemic, with the oscillation as the infection and the excess sharpness as the susceptible population.} The landing lies slightly below the mirror image, by the second-order term, and beyond $\eta\Lambda(c_0)\approx2.4$ the gap widens, because a large oscillation shrinks $c$ by the factor $\abs{1-\eta^2z^2}$, which is smaller than the $e^{-\eta^2z^2}$ of the expansion. From about $2.9$ on, most runs land near the flattest minimizer, $\eta S=2\eta y=0.8$, with a scatter that depends on the initial perturbation at the level of $10^{-6}$: in these runs the oscillation grows large enough to behave chaotically before it dies out. At $\eta\Lambda(c_0)=4$ the scatter ends in divergence, exactly where \cref{thm:stability:ellipse} puts it.

\begin{figure}[t]
\begin{fullwidth}
\centering
\begin{tikzpicture}
\begin{groupplot}[group style={group size=2 by 1,horizontal sep=17mm},hplot,width=82mm,height=56mm]
\nextgroupplot[xlabel={start: $\eta\Lambda(c_0)$},ylabel={landing: $\eta S(\theta_\infty)$},xmin=0.8,xmax=4.2,ymin=0.6,ymax=2.25,legend style={at={(0.5,1.03)},anchor=south,font=\scriptsize},legend columns=2]
\fill[seamist] (4,0.6) rectangle (4.2,2.25);
\addplot[guide] coordinates {(0.8,2) (4.2,2)};
\addplot[guide] coordinates {(0.8,0.8) (4.2,0.8)};
\addplot[theory] table[x=el0,y=refl] {figures/data/15-stability-catapult.dat};
\addplot[theory2] table[x=el0,y=second] {figures/data/15-stability-catapult.dat};
\addplot[simulation,mark size=1.1pt] table[x=el0,y=rb] {figures/data/15-stability-catapult.dat};
\addplot[simulation2,mark size=0.9pt] table[x=el0,y=rc] {figures/data/15-stability-catapult.dat};
\legend{,,reflection law,second order,{exact, $z_0=10^{-6}$},{exact, $z_0=-10^{-6}$}}
\node[font=\scriptsize,text=seadeep,rotate=90] at (axis cs:4.1,1.4) {diverges};
\nextgroupplot[xlabel={$\eta y$},ylabel={$\lim\,(\eps_\infty+\eps_0)/\eps_0^2$},xmin=0,xmax=0.82,ymin=-1,ymax=3,legend style={at={(0.5,1.03)},anchor=south,font=\scriptsize},legend columns=2]
\addplot[guide] coordinates {(0,0) (0.82,0)};
\addplot[theory,domain=0:0.8,samples=120] {2*x^2/(1-x^2)-2/3};
\addplot[theory3,domain=0:0.8,samples=120] {(4/3)*x^2/(1-x^2)};
\addplot[simulation] table[x=h,y=c2exact] {figures/data/15-stability-second.dat};
\legend{,{second-order law \eqref{eq:stability:second}},saturation only,exact map}
\end{groupplot}
\end{tikzpicture}
\end{fullwidth}
\caption{Where a catapult lands. Left: the sharpness at landing against the sharpness of the start, for $y=1$ and $\eta=0.4$. Each point is one run of the exact map, started on the level set of $c_0$ at residual $z_0$ and iterated to convergence; the grid of starts has spacing $0.0125$. Curves: the identity below the edge (\cref{prop:stability:lazy}), the reflection law \eqref{eq:stability:reflection}, and the second-order law \eqref{eq:stability:second}, both cut at the flattest minimizer $2\eta y=0.8$ (dashed). No run converges beyond $\eta\Lambda(c_0)=4$. Right: the second-order coefficient of the landing law, measured from the exact map by extrapolating runs with $\eps_0=0.01,0.005,0.0025$ and $z_0=10^{-9}$ (points), against \eqref{eq:stability:second} and against the saturation term alone. Script: \texttt{15-stability-catapult.py}.}
\label{fig:stability:catapult}
\end{figure}

\paragraph{Sizes and times.} Three more numbers come out of the solution. The loss never exceeds $2/\eta^2$ during a catapult that starts inside $\cE$ (\cref{thm:stability:ellipse}), a bound that is attained in the limit $\eta\Lambda(c_0)\to4$. Near the edge the spike has height about $\eps_0^2/(4\eta^2(1-\eta^2y^2))$: for $\eta y=0.4$ and $\eps_0=0.08$ the prediction is $0.0119$, and the exact map gives $0.0125$. And the spike comes after about $\log(A_{\mathrm{peak}}/A_0)/(2\eps_0)$ steps, the time that an oscillation growing by the factor $(1+\eps_0)^2$ per step needs to rise from its initial size $A_0$ to $\eps_0^2/\nu$. The delay is logarithmic in the initial perturbation and inversely proportional to the distance from the edge, the signature of an escape from an unstable fixed point that recurs throughout this book (\cref{ch:deeplinear,ch:singleindex}). \Cref{fig:stability:trajectory} shows one catapult in time.

\begin{figure}[t]
\begin{fullwidth}
\centering
\begin{tikzpicture}
\begin{groupplot}[group style={group size=2 by 1,horizontal sep=17mm},hplot,width=82mm,height=52mm,xmin=0,xmax=100]
\nextgroupplot[xlabel={step $t$},ylabel={loss $\Loss(\theta_t)$},ymode=log,ymin=1e-17,ymax=1]
\addplot[theory,mark=*,mark size=0.7pt] table[x=t,y=loss] {figures/data/15-stability-trajectory.dat};
\nextgroupplot[xlabel={step $t$},ylabel={sharpness},ymin=1.5,ymax=2.5,legend pos=north east,legend style={font=\scriptsize},legend columns=2]
\addplot[guide] coordinates {(0,2) (100,2)};
\addplot[simulation,mark size=1.1pt] table[x=t,y=sharp] {figures/data/15-stability-trajectory.dat};
\addplot[theory] table[x=t,y=gfs] {figures/data/15-stability-trajectory.dat};
\addplot[theory3,domain=0:100] {1.7};
\addplot[theory2,domain=0:100] {1.6743};
\legend{,{$\eta S(\theta_t)$},{$\eta\Lambda(c_t)$},reflection law,second order}
\end{groupplot}
\end{tikzpicture}
\end{fullwidth}
\caption{One catapult, for $y=1$, $\eta=0.4$, started on the level set of the minimizer with $\eta S=2.3$ at residual $z_0=10^{-6}$. Left: the loss grows geometrically for about fifty steps, peaks near $0.14$, and then decays geometrically at the rate set by the new minimizer. Right: the sharpness $\eta S(\theta_t)$ (points) and the sharpness $\eta\Lambda(c_t)$ of the minimizer that gradient flow would reach from $\theta_t$ (curve). The second decreases monotonically, as \cref{lem:stability:coordinates}(ii) requires while the loss is below $1/\eta^2$, from $2.3$ to $1.670$; the lines mark the predictions $1.7$ of \eqref{eq:stability:reflection} and $1.674$ of \eqref{eq:stability:second}. Script: \texttt{15-stability-trajectory.py}.}
\label{fig:stability:trajectory}
\end{figure}

\paragraph{From the catapult to the edge of stability.} The model has a catapult but no edge of stability: after one oscillation it lands below $2/\eta$ and stays there. The reason is that its gradient flow conserves $c$, so nothing pushes the sharpness of the flow's destination back up. In a network that is trained at the edge, progressive sharpening does exactly that. Add it to \eqref{eq:stability:slowfast} as a constant rate $\gamma>0$:
\begin{equation}\label{eq:stability:lv}
  \frac{dA}{dt}=2\eps A,\qquad \frac{d\eps}{dt}=\gamma-\nu A .
\end{equation}
This is the Lotka--Volterra system. The function $\eps^2+\nu A-\gamma\log A$ is constant along its solutions, which are closed curves around the point $(A,\eps)=(\gamma/\nu,0)$. Over one period $A$ and $\eps$ return to their initial values, so the time averages of $d\log A/dt=2\eps$ and of $d\eps/dt=\gamma-\nu A$ vanish:
\[
  \langle\eps\rangle=0,\qquad \langle A\rangle=\gamma/\nu .
\]
The sharpness of the flow's destination oscillates around $2/\eta$, its time average is exactly $2/\eta$, and the oscillation carries exactly the energy needed to cancel the sharpening. With the saturation term $-\beta A^2$ restored, the closed curves become spirals into the fixed point $A=\gamma/\nu$, $\eps=\beta\gamma/(2\nu)$, so the sharpness settles slightly above $2/\eta$, as in \cref{obs:stability:eos}. This is a two-variable caricature of the mechanism that \textcite{DamianNichaniLee2023} called \term{self-stabilization}\index{self-stabilization} and analysed for general losses.

\paragraph{When no minimizer is stable.} For $1<\eta y<2$ every minimizer is sharper than $2/\eta$, and gradient descent started in $\cK$ stays bounded (\cref{thm:stability:ellipse}) but almost never converges (\cref{thm:stability:landing}). Numerically, its imbalance tends to zero, and it settles on the balanced line $u=v$, on an orbit of period two around the balanced minimizer for $\eta y$ up to about $1.24$. The orbit then period-doubles, and for most $\eta y$ above $1.3$ the motion is chaotic; at $\eta y=2$ the map on the balanced line is again a tripling map. In this band the loss never goes to zero: training is stuck at the edge, without the decrease of the loss that \cref{obs:stability:eos} reports.

\section{What the model gets wrong}\label{sec:stability:wrong}

\paragraph{What was thrown away.} The product model keeps one example, one hidden unit and two layers, and so the top eigenvector of its Hessian is pinned to the normal of a single curve. In a network the top eigenvector rotates during training, and several eigenvalues can reach $2/\eta$ at once. The model uses the square loss, whose minimum has a curvature that does not change once it is reached; with the cross-entropy loss, the curvature of a minimizing sequence tends to zero. It uses full-batch gradient descent with a fixed step, with no noise, momentum, preconditioning, schedule or warmup. Above all, its gradient flow conserves the imbalance, so the sharpness of the flow's destination never drifts upward: the model has no progressive sharpening. It reproduces the catapult and the selection of minimizers flatter than $2/\eta$, and it cannot reproduce the long phase at the edge, which needs sharpening to balance. Finally, the exact identities of \cref{lem:stability:coordinates,thm:stability:ellipse} come from the bilinear structure, and nothing guarantees them for deeper or nonlinear models.

\paragraph{Where it breaks.} Six results mark its limits. (1) The product model does not explain why the sharpness rises in networks. \Textcite{CohenEtAl2021} left progressive sharpening unexplained, and the partial explanations since then use other models: second-order regression models \parencite{AgarwalaPedregosaPennington2023}, and an analysis that ties the rise to the growth of the output layer \parencite{WangLiLi2022}. (2) With the cross-entropy loss the sharpness eventually decreases \parencite{CohenEtAl2021}, and for logistic regression on separable data gradient descent converges with every step size, after an initial phase in which the loss oscillates \parencite{WuBravermanLee2023,WuEtAl2024}; there the edge has no fixed location (\cref{ch:margin}). (3) With minibatches the quantity that settles at $2/\eta$ is not the full-batch sharpness: the sharpness that minibatch training reaches early is lower for larger step sizes and smaller batches \parencite{JastrzebskiEtAl2020}, and the edge of stability has a stochastic counterpart, measured with a sharpness adapted to the minibatch \parencite{LeeJang2023,AndreyevBeneventano2024}. (4) With Adam the preconditioned sharpness equilibrates near $38/\eta$ for $\beta_1=0.9$, the stability threshold of momentum of the kind Adam uses (Exercise~\ref{ex:stability:momentum}), rather than $2/\eta$ \parencite{CohenEtAl2022}. (5) The model's divergence threshold $4/\lambda_{\mathrm{init}}$ is not universal: for wide ReLU networks \textcite{LewkowyczEtAl2020} measured a largest step size of about $12/\lambda_{\mathrm{init}}$. (6) Linear stability concerns limits, and \cref{thm:stability:landing} says nothing about the thousands of steps that a network spends at the edge without converging, during which the loss keeps falling; that phase is where the interesting optimization happens.

\paragraph{Fine structure.} Three successors repair the main failure. \Textcite{DamianNichaniLee2023} expand the loss to third order around a point on the edge and show that the oscillation along the top eigenvector produces, through the third derivative, a gradient of the sharpness that pushes the iterates back below $2/\eta$; gradient descent then tracks gradient descent projected onto the set $\{S\le2/\eta\}$, and \eqref{eq:stability:lv} is the two-variable shadow of their analysis. \Textcite{CohenEtAl2025} derive a \term{central flow}\index{central flow}, a differential equation for the time-averaged trajectory at the edge that predicts the loss and the sharpness, including the case of several unstable directions. A series of minimalist models give complete proofs of edge-of-stability behaviour in low dimension: a two-parameter quartic loss \parencite{ZhuEtAl2023}, a model of learning threshold neurons \parencite{AhnEtAl2023}, normalized gradient descent and its sharpness-reduction flow \parencite{AroraLiPanigrahi2022}, scalar networks, for which the sharpness of the gradient-flow solution decreases monotonically \parencite{KreislerEtAl2023}, and a bifurcation analysis of the trajectory \parencite{SongYun2023}. Large step sizes balance the layers of two-layer models \parencite{WangEtAl2022}, and drive diagonal linear networks toward sparser solutions \parencite{EvenEtAl2023}, the subject of \cref{ch:diagonal}. Still open are a solvable model in which progressive sharpening is derived and the long phase at the edge is solved (Exercise~\ref{ex:stability:sharpening}), convergence from the interior of the catapult ellipse (Exercise~\ref{ex:stability:converge}), and a theory of the stochastic and adaptive edges.

\section{Model card}\label{sec:stability:card}

\modelcard{Full-batch gradient descent raises the sharpness to $2/\eta$ and then trains at that edge; large initial steps cause a loss spike and a drop in curvature, the catapult.}
{$\Loss(u,v)=\frac12(uv-y)^2$ trained by gradient descent with step size $\eta$; imbalance $c=u^2-v^2$; minimizer sharpness $\Lambda(c)=(c^2+4y^2)^{1/2}$.}
{Exact: $c'=(1-2\eta^2\Loss)c$; convergence, for almost every start, only to minimizers with $\eta S\le2$; an invariant ellipse $\Psi'=(1+\eta^2z(z+y))\Psi$. Derived: the reflection law.}
{Phase boundaries at $\eta\Lambda=2$ and $\eta\Lambda=4$ (exact); landing at $\eta S\approx4-\eta\Lambda(c_0)$, with correction $(\frac{2\eta^2y^2}{1-\eta^2y^2}-\frac23)\eps_0^2$; spike height $\eps_0^2/(4\eta^2(1-\eta^2y^2))$; loss below $2/\eta^2$.}
{No progressive sharpening, hence no sustained edge of stability; square loss, full batch, plain gradient descent; the threshold $4$ is special to bilinear models.}
{Self-stabilization; central flows; minimalist models with proofs; stochastic and adaptive edges; large-step logistic regression.}

The model is the prologue's quadratic with one more parameter, and the extra parameter changes the verdict: the threshold $2/\lambda_1$ of \cref{ch:prologue} becomes a rule about which minimizers can be reached, not a constraint on the step size. The product $uv$ is also a diagonal linear network with one coordinate, and its conserved imbalance is the quantity whose size sets the implicit bias in \cref{ch:diagonal}; large steps act on that same quantity. The conservation law is the scalar case of the balancedness of deep linear networks in \cref{ch:deeplinear}. The delay before the catapult, logarithmic in the initial perturbation, is the time to escape an unstable fixed point, as in \cref{ch:deeplinear,ch:singleindex,ch:committee}. In the committee machine of \cref{ch:committee} the step size enters the order-parameter equations through terms of order $\eta^2$, and a large enough step destabilizes the learning dynamics there too. With the logistic loss, large steps lead to the margin-maximizing solutions of \cref{ch:margin} after an unstable phase.

\begin{exercises}
\item \consolidation (The Hessian off the valley.) (a) Show that the eigenvalues of $\nabla^2\Loss(u,v)$ are $\frac12\bigl(\tau\pm\sqrt{c^2+4(uv+z)^2}\bigr)$, so that $S(\theta)\ge\tau/2$, with $S=\tau$ on $\cM$. (b) On the balanced line $u=v$, show that $S=\max\{3u^2-y,\;y-u^2\}$. Conclude that along gradient flow from a small balanced initialization the sharpness first falls from $y$ to $y/2$ and then rises to $2y$, a two-parameter caricature of progressive sharpening that ends at the flattest minimizer.\label{ex:stability:hessian}

\item \consolidation (Edges in the valley.) Let $\eta=0.4$ and $y=1$. (a) Show that the minimizers on the positive branch of $\cM$ with $2<\eta S<4$ are the points $(u,1/u)$ with $u\in(0.3178,0.4569)\cup(2.1889,3.1463)$. (b) Show that $\cE$ meets the line $u=v$ where $\eta u^2=2+\eta y$, and the line $u=-v$ where $\eta u^2=2-\eta y$. (c) For general $\eta y<2$, show that $\cK$ contains every point with $\eta(u^2+v^2)<4-2\eta y$, and that $\cK$ is contained in $\{\eta(u^2+v^2)<4+2\eta y\}$.\label{ex:stability:edges}

\item \consolidation (The tripling map on the ellipse.) Let $\eta y<2$, and parameterize $\cE$ by $a=2\sqrt{(2+\eta y)/\eta}\,\sin\varphi$ and $b=2\sqrt{(2-\eta y)/\eta}\,\cos\varphi$. (a) Show that $\eta z=-2\cos2\varphi$ on $\cE$. (b) Using $\sin3\varphi=\sin\varphi\,(1+2\cos2\varphi)$ and $\cos3\varphi=-\cos\varphi\,(1-2\cos2\varphi)$, show that $G$ acts on $\cE$ as $\varphi\mapsto\pi-3\varphi$. (c) Deduce that $G$ has exactly four fixed points on $\cE$, the minimizers with $\eta S=4$, and exactly $3^k-(-1)^k$ points of $\cE$ fixed by $G^k$ for each $k\ge1$. (d) Show that the two orbits of period two are $\{\theta,-\theta\}$ for the points $\theta$ where $\cE$ meets the diagonals, with $\eta z=2$ on $u=v$ and $\eta z=-2$ on $u=-v$.\label{ex:stability:tripling}

\item \consolidation (Volterra's averages.) (a) Show that $\eps^2+\nu A-\gamma\log A$ is constant along the solutions of \eqref{eq:stability:lv}, and that every solution with $A(0)>0$ is periodic. (b) Show that small oscillations around $(\gamma/\nu,0)$ have period $2\pi/\sqrt{2\gamma}$. (c) With the term $-\beta A^2$ added to the first equation, show that the fixed point moves to $(\gamma/\nu,\beta\gamma/(2\nu))$ and is a stable focus.\label{ex:stability:volterra}

\item \consolidation\harder (The second-order reflection law.) Let $\ell=\log A$ and $w=\log c$, and model one step of gradient descent near the edge by
\[
  \ell'=\ell+F(w)-\beta e^{\ell},\qquad w'=w-\eta^2e^{\ell},\qquad F(w)=2\log\bigl(1+\eps(w)\bigr),
\]
with $\eps(w)=\eta(e^{2w}+4y^2)^{1/2}-2$. (a) Without the term $\beta e^\ell$, show that $H=\eta^2e^\ell+\int^wF$ is constant along the flow $\dot\ell=F(w)$, $\dot w=-\eta^2e^\ell$, and that one step of the map increases it by $\frac12G^2F'(w)+\frac12GF(w)^2$ up to higher-order terms, with $G=\eta^2e^\ell$. (b) Sum these increments over a catapult, converting sums into integrals with $dw=-G$ and using the leading-order solution $\nu A=\eps_0^2-\eps^2$; add the effect of $\beta e^\ell$, and derive \eqref{eq:stability:second}. (c) Check with the exact map that $(\eps_\infty+\eps_0)/\eps_0^2\to-2/7$ for $\eta y=0.4$. \hint{at the edge $d\eps/dw=2(1-\eta^2y^2)$ and $d^2\eps/dw^2=2(1-\eta^4y^4)$; the increments of (a) contribute $-\frac43$, and the expansion $F=2\eps-\eps^2$ contributes $\frac13$.}\label{ex:stability:second}

\item \variation (No label.) Let $y=0$, so that $\Loss=\frac12u^2v^2$ and the minimizers are the two axes, with sharpness $\tau$. (a) Show that $\tau'=\tau+\eta z^2(\eta\tau-4)$ exactly, and that $\cE$ is the circle $\eta\tau=4$. (b) If $\eta\tau_0<4$, show that $\Loss(\theta_t)\to0$. (c) If $\eta\tau_0>4$, show that $\tau_t\to\infty$ for almost every $\theta_0$. \hint{in (c), if $\tau_t$ stayed bounded it would converge, and $\sum_t\eta z_t^2(\eta\tau_t-4)<\infty$ would force $z_t\to0$; now argue as in \cref{thm:stability:landing}.}\label{ex:stability:yzero}

\item \variation (Three factors.) Let $\Loss(u,v,w)=\frac12(uvw-y)^2$ and $z=uvw-y$. (a) Show that gradient flow conserves $u^2-v^2$ and $v^2-w^2$, and that gradient descent multiplies $u^2-v^2$ by $1-\eta^2z^2w^2$. (b) Show that the sharpness of a minimizer is $y^2(u^{-2}+v^{-2}+w^{-2})$, and that the flattest minimizers are balanced, with sharpness $3y^{4/3}$. (c) Adapt the proof of \cref{thm:stability:landing}(ii) to show that the initializations from which gradient descent converges to a minimizer with $\eta S>2$ form a null set. Deduce that if $\eta>2/(3y^{4/3})$, then gradient descent from almost every initialization converges to no minimizer at all.\label{ex:stability:three}

\item \variation (Every width.) In the model of \textcite{LewkowyczEtAl2020}, $u,v\in\R^m$, $f=\ip{u,v}/\sqrt m$, $\Loss=\frac12(f-y)^2$, and the curvature is $\lambda=(\norm u^2+\norm v^2)/m$. (a) With $z=f-y$, show that gradient descent satisfies $z'=z\bigl(1-\eta\lambda+\frac{\eta^2}mz(z+y)\bigr)$ and $\lambda'=\lambda+\frac\eta mz^2(\eta\lambda-4)-\frac{4\eta}myz$. (b) Show that $\Psi_m=\eta\lambda-\frac{\eta^2y}mz-4$ satisfies $\Psi_m'=\bigl(1+\frac{\eta^2}mz(z+y)\bigr)\Psi_m$. (c) Deduce that if $\eta y<2\sqrt m$ and $\Psi_m<0$ at initialization, then the iterates stay bounded; so the threshold $4/\lambda_{\mathrm{init}}$ is exact up to a correction of order $1/m$.\label{ex:stability:width}

\item \variation (Momentum and the adaptive edge.) Gradient descent with momentum in the form used by Adam updates $m_t=\beta_1m_{t-1}+(1-\beta_1)\nabla\Loss(\theta_t)$ and $\theta_{t+1}=\theta_t-\eta m_t$. (a) On a quadratic with curvature $\lambda_1$ in some direction, show that this iteration is stable in that direction if and only if $\eta\lambda_1<2(1+\beta_1)/(1-\beta_1)$. (b) Evaluate the threshold for $\beta_1=0.9$, and compare with the value $38/\eta$ at which \textcite{CohenEtAl2022} observed the preconditioned sharpness of Adam to equilibrate.\label{ex:stability:momentum}

\item \simulation (Other step sizes.) Run \texttt{15-stability-catapult.py} with $\eta y=0.2$ and with $\eta y=0.8$. Check that near the edge the landing follows \eqref{eq:stability:second}, with coefficients $-0.583$ and $2.889$; that the region of smooth landing ends, and the scatter begins, at a start sharpness that decreases as $\eta y$ grows; and that no run converges beyond $\eta\Lambda(c_0)=4$. For $\eta y=0.8$, explain why the reflection law reaches the flattest minimizer already at $\eta\Lambda(c_0)=2.4$.\label{ex:stability:sim1}

\item \simulation (A network at the edge.) Using \texttt{15-stability-mlp.py}: (a) Train with $\eta\in\{0.3,0.5\}$, and check that the sharpness levels off within a few percent of $2/\eta$. (b) Plot the time $\eta t$ at which the sharpness first reaches $2/\eta$ against $2/\eta$ for several step sizes, and check that the points lie on the sharpness curve of the run with the smallest step: before the edge, gradient descent follows gradient flow. (c) Replace the full gradient by a minibatch gradient with $50$ examples and $\eta=0.25$; the full-batch sharpness should level off well below $2/\eta=8$ (near $6.6$ in our runs), as minibatch training is known to do.\label{ex:stability:sim2}

\item \research (Convergence inside the ellipse.) For $0<\eta y<1$, every one of the $10^5$ random initializations in $\cK$ that we tried converged. Prove that gradient descent converges from almost every point of $\cK$. The case $y=0$ is Exercise~\ref{ex:stability:yzero}(b), where $\tau$ is monotone; for $y>0$ no monotone quantity is known to us, and the orbits that pass near $\cE$ are chaotic for a while. A second open question is the law of the landing sharpness for starts near $\cE$, where it is effectively random.\label{ex:stability:converge}

\item \research (Sharpening from first principles.) The product model must be given progressive sharpening by hand, as the rate $\gamma$ of \eqref{eq:stability:lv}. Find a solvable model in which the rate is derived, and which then trains at the edge of stability with a decreasing loss. \Textcite{AgarwalaPedregosaPennington2023} derive progressive sharpening in quadratic regression models; what is open is a model in which both the rate and the long phase at the edge are solved, and a derivation of the rate for a network.\label{ex:stability:sharpening}
\end{exercises}


\notesandrefs
%%NOTES%%

\printchapterbib
\end{refsection}

% Chapter 16 of The Hydrogen Atoms of Deep Learning: scaling laws from spectra and from Zipf's law.
\casechapter{Scaling Laws}{Power laws from spectra and from Zipf's law}{Power-law linear regression and quanta}{ch:scaling}
\begin{refsection}[chapters/16-scaling.bib]

\begin{chaptersummary}
The test loss of language models falls as a power law in the number of parameters, the number of training tokens, and the compute, over many orders of magnitude, with small exponents that differ between parameters and data. Two solvable models produce such laws. In linear regression on Gaussian features whose target energy decays as $j^{-(1+b)}$ we compute the learning curve of least squares exactly at every model size $N$ and sample size $D$, and read off a parameter exponent $b$ and a data exponent $b/(1+b)$; ridge regression and the deterministic equivalent show that the spectrum of the features only converts a regularization or a training time into a number of fitted features. A memorizing learner on items with Zipf frequencies $p_k\propto k^{-(1+b)}$ gives the same two exponents, with an exact constant, and an elementary theorem allocates compute between parameters and data. The models get the mechanism of power laws but not the exponents: those are inputs, and fitted neural exponents violate the models' prediction that the data exponent is the smaller.
\end{chaptersummary}

\section{The phenomenon: smooth power laws over many decades}\label{sec:scaling:phenomenon}

How much better does a model get when it is made ten times larger, or trained on ten times more data? For most of the history of machine learning the honest answer was ``it depends''. \Textcite{HestnessEtAl2017} measured it systematically, in several domains among them machine translation and language modeling, and found that the generalization error falls as a power law in the size of the training set, with an exponent that depends on the task more than on the architecture, and that the model size needed to fit the data grows sublinearly with it.

\Textcite{KaplanEtAl2020} made the observation quantitative for transformer language models. Measured in nats per token, the test loss of a model with $N$ non-embedding parameters trained to convergence on a large corpus followed $\Loss(N)=(N_c/N)^{\alpha_N}$ with $\alpha_N\approx0.076$ and $N_c\approx8.8\times10^{13}$, and the loss of a large model trained with early stopping on $D$ tokens followed $\Loss(D)=(D_c/D)^{\alpha_D}$ with $\alpha_D\approx0.095$ and $D_c\approx5.4\times10^{13}$. Some of the trends spanned more than seven orders of magnitude, and architectural details such as depth and width at fixed $N$ mattered little.\sidenote{Kaplan et al.\ estimate the training compute as $C\approx6ND$ floating-point operations: about $2N$ per token for the forward pass and twice that for the backward pass. We use $C=6ND$ throughout.} From these fits they concluded that a growing compute budget should be spent mostly on larger models, with the optimal number of parameters growing like $C^{0.73}$.

\Textcite{HoffmannEtAl2022} revisited the allocation with more than $400$ models and three estimation procedures. The third fitted the parametric form
\begin{equation}\label{eq:scaling:chinchilla}
  \Loss(N,D)=\Loss_\infty+\frac A{N^{\alpha_N}}+\frac B{D^{\alpha_D}}
\end{equation}
(they write $E$, $\alpha$, $\beta$ for our $\Loss_\infty$, $\alpha_N$, $\alpha_D$), and reported $\alpha_N\approx0.34$ and $\alpha_D\approx0.28$. All three procedures said that parameters and tokens should grow at about the same rate, and a $70$-billion-parameter model trained on more tokens with the compute of a $280$-billion-parameter one outperformed it. The two studies disagree about the allocation, and later work traced much of the disagreement to how parameters were counted and to the small scale of the earlier fits \parencite{PearceSong2024}, and to the cost of the last layer, the warmup, and the tuning of the optimizer \parencite{PorianEtAl2024}.

\begin{observation}[\textcite{HestnessEtAl2017}; \textcite{KaplanEtAl2020}; \textcite{HoffmannEtAl2022}]\label[observation]{obs:scaling:powerlaw}
For language models trained on large corpora, the test loss, minus an irreducible part, decreases as a power law in the number of parameters when data are plentiful, and in the number of training tokens when the model is large, over several orders of magnitude, with exponents well below $1$. The loss of compute-optimal models is a power law in the compute.
\end{observation}

How robust is this? The power-law form is robust; the numbers are not. \Textcite{BesirogluEtAl2024} reconstructed the data behind the third procedure of \textcite{HoffmannEtAl2022} from its plots and refitted it: the published estimates fit the reconstructed data poorly, were inconsistent with the other two procedures, and came with confidence intervals far too narrow for a few hundred runs, while their own fit had a larger data exponent and was consistent with the scaling policy actually used for Chinchilla. Curves of other quantities bend, and \textcite{CaballeroEtAl2023} fit them with smoothly broken power laws. Repeating data breaks the law in $D$ slowly: up to about four epochs of repeated tokens are almost as good as fresh ones, and then the value of more compute decays \parencite{MuennighoffEtAl2023}. Choosing the data instead of sampling it can beat the law \parencite{SorscherEtAl2022}. And while the loss is smooth, some downstream abilities appear abruptly with scale \parencite{WeiEtAl2022}, although part of the abruptness comes from the choice of metric \parencite{SchaefferEtAl2023}.

Three questions follow. Why a power law at all, rather than an exponential or a logarithm? What sets the exponents, and why do parameters and data have different ones? And how should compute be split between them? The models of this chapter answer the first and third exactly and the second in a form that can be checked.

\section{The models: power-law features and Zipf-distributed quanta}\label{sec:scaling:model}

A power law in $N$ needs a hierarchy: a model of size $N$ must capture the $N$ most important pieces of the task and miss the rest, and the missing importance must decay as a power. Two hierarchies are natural. In regression the pieces are directions of feature space, and their importance is the variance they explain; in language the pieces are items of knowledge, and their importance is how often they are used. In this chapter $N$ and $D$ follow the scaling-law literature: $N$ counts parameters, the $p$ of \cref{ch:doubledescent}, and $D$ counts training examples, the $n$ of that chapter.

\begin{assumption}[Power-law features]\label[assumption]{ass:scaling:features}
The input is a sequence of independent features $x_j\sim\cN(0,\lambda_j)$, $j=1,\dots,P$, with $P\in\{1,2,\dots\}\cup\{\infty\}$ and variances $\lambda_1\ge\lambda_2\ge\dots>0$. The label is $y=\sum_jw_{\star,j}x_j+\eps$, with $\sum_j\lambda_jw_{\star,j}^2<\infty$ and noise $\eps$ of mean $0$ and variance $\sigma^2$, independent of $x$. In the \term{power-law case}\index{power-law spectrum} $P=\infty$ and, for exponents $a>1$ and $b>0$,
\begin{equation}\label{eq:scaling:powerlaw}
  \lambda_j=j^{-a},\qquad \lambda_jw_{\star,j}^2=j^{-(1+b)} .
\end{equation}
\end{assumption}

The \term{model of size $N$}\index{model size} predicts with the first $N$ features, $f(x)=\sum_{j\le N}w_jx_j$, and is fitted by least squares on $D$ independent examples $(x_i,y_i)$; its data matrix $X_N\in\R^{D\times N}$ has rows $(x_{i1},\dots,x_{iN})$, and for $N\le D$ the estimator is $\hat w_N=(X_N^\top X_N)^{-1}X_N^\top y$. Because the features are independent, the population loss of a weight vector $w$ supported on the first $N$ coordinates is
\begin{equation}\label{eq:scaling:loss}
  \Loss(w)=\E\bigl(y-f(x)\bigr)^2=\sigma^2+T(N)+\sum_{j\le N}\lambda_j(w_j-w_{\star,j})^2,\qquad T(N)=\sum_{j>N}\lambda_jw_{\star,j}^2 .
\end{equation}
The \term{approximation error}\index{approximation error} $T(N)$ is what the best model of size $N$ still misses; the \term{excess loss}\index{excess loss} is $\Loss-\sigma^2$. In the power-law case $T(N)$ is a tail of $\zeta(1+b)$, and $b$ is called the \term{source exponent}: it measures how fast the importance of the features decays, and it does not involve $a$.

The second model replaces regression by recall.

\begin{assumption}[Zipf-distributed quanta]\label[assumption]{ass:scaling:quanta}
A task consists of items $k=1,2,\dots$, the \term{quanta}\index{quanta}. Each example, at training or test time, requires one item, drawn independently with probability
\begin{equation}\label{eq:scaling:zipf}
  p_k=\frac{k^{-(1+b)}}{\zeta(1+b)},\qquad b>0 .
\end{equation}
A learner answers correctly exactly when it has learned the item required. A learner of \term{capacity} $N$ can hold $N$ items; a \term{memorizing learner} learns every item that occurs in its $D$ training examples.
\end{assumption}

Frequencies of words follow \term{Zipf's law}\index{Zipf's law}, $p_k\propto k^{-s}$ with $s$ close to $1$ \parencite{Zipf1949}; \textcite{Hutter2021} built the memorizing learner to show that power-law learning curves with every exponent arise from Zipf-distributed data, and \textcite{MichaudEtAl2023} proposed that a language model's skills are such quanta, learned in order of decreasing frequency.\sidenote{The exponent is written $1+b$ so that $b>0$ is the exponent of the tail $\sum_{k>N}p_k$, as $b$ is the exponent of $T(N)$ in \eqref{eq:scaling:powerlaw}; \textcite{Hutter2021} writes $\alpha$ for our $b$.}

\begin{dictionary}
feature $j$, or quantum $k$ & a direction of the representation, or a skill, ranked by importance\\
model size $N$ & number of parameters\\
sample size $D$ & number of training tokens, each seen once\\
variances $\lambda_j=j^{-a}$ & spectrum of the data in the network's representation, e.g.\ of a tangent kernel\\
target energy $j^{-(1+b)}$, frequency $p_k$ & share of the loss that each direction or skill carries\\
approximation error $T(N)$ & loss of the best model of size $N$ on infinite data\\
noise variance $\sigma^2$ & irreducible loss $\Loss_\infty$, the entropy of the text\\
ridge parameter $\lambda$, training time $t$ & weight decay, training time: how many directions are fitted\\
memorizing an item seen once & learning a skill from few examples\\
$C=6ND$ & training compute in floating-point operations\\
\end{dictionary}

\paragraph{What was thrown away, and why first.} The features are fixed and independent, so the model learns no representation and its exponents are properties of the data in a given basis; but the question ``why a power law'' is already nontrivial for fixed features, and independence makes the learning curve exact. The model of size $N$ uses the top $N$ features, which assumes that a network spends its parameters on the most important directions first, the staged learning of \cref{ch:deeplinear}. The noise is white, and it plays the role of the irreducible loss. Every example is fresh, as in one-epoch training of language models. The quanta are independent and are learned all at once; nothing is composed of anything.

\paragraph{What one had to notice.} Two things. A power law in the loss needs a power-law hierarchy of importance, and real data supply them: kernel spectra of image data decay as powers \parencite{SpiglerEtAl2020}, and word frequencies follow Zipf's law. And a finite sample makes the hierarchy visible only down to a threshold: a direction or item worth less than about $1/D$ of the loss cannot be learned from $D$ examples. The data exponent comes from where this threshold falls.

\section{The solution: exact learning curves of power-law regression}\label{sec:scaling:regression}

\begin{lemma}[Power-law tails]\label[lemma]{lem:scaling:tail}
For $b>0$ and every integer $N\ge1$,
\[
  \frac{(N+1)^{-b}}{b}\le\sum_{j>N}j^{-(1+b)}\le\frac{N^{-b}}{b} .
\]
In particular, in the power-law case $T(N)=\frac{N^{-b}}b\,(1+O(1/N))$.
\end{lemma}

\begin{proof}
The function $t\mapsto t^{-(1+b)}$ decreases, so $\int_j^{j+1}t^{-(1+b)}\dd t\le j^{-(1+b)}\le\int_{j-1}^jt^{-(1+b)}\dd t$ for $j\ge2$. Summing over $j>N$ gives $\int_{N+1}^\infty t^{-(1+b)}\dd t\le\sum_{j>N}j^{-(1+b)}\le\int_N^\infty t^{-(1+b)}\dd t$, and the integrals equal $(N+1)^{-b}/b$ and $N^{-b}/b$. Finally $(1+1/N)^{-b}=1+O(1/N)$.
\end{proof}

The main theorem says that, for least squares on the top $N$ features, the variances $\lambda_j$ do not matter at all: the learning curve depends on the target only through the tail $T(N)$.

\begin{theorem}[Exact learning curve]\label{thm:scaling:exact}
Under \cref{ass:scaling:features}, let $1\le N\le\min(P,D-2)$. The least-squares estimator on the first $N$ features from $D$ examples has test loss
\begin{equation}\label{eq:scaling:exact}
  \E\Loss(\hat w_N)=\bigl(\sigma^2+T(N)\bigr)\frac{D-1}{D-N-1} .
\end{equation}
If $\sigma^2+T(N)>0$ and $N\in\{D-1,D\}$, then $\E\Loss(\hat w_N)=+\infty$.
\end{theorem}

\begin{proof}
The plan is to treat the omitted features as noise, as in the weak-features model of \cref{ch:doubledescent}, and to whiten the retained ones.

\emph{The omitted features are noise.} For example $i$ write $y_i=\sum_{j\le N}w_{\star,j}x_{ij}+\xi_i$ with $\xi_i=\sum_{j>N}w_{\star,j}x_{ij}+\eps_i$; the series converges in $L^2$ because $\sum_j\lambda_jw_{\star,j}^2<\infty$. The features of example $i$ are independent and $\eps_i$ is independent of them, so $\xi_i$ is independent of $(x_{i1},\dots,x_{iN})$, has mean $0$, and has variance $T(N)+\sigma^2$; across examples the $\xi_i$ are independent. Hence the vector $\xi=(\xi_1,\dots,\xi_D)$ is independent of $X_N$ and satisfies $\E\xi\xi^\top=(\sigma^2+T(N))I_D$.

\emph{Whitening.} Let $\Lambda=\diag(\lambda_1,\dots,\lambda_N)$ and $Z=X_N\Lambda^{-1/2}$, a $D\times N$ matrix with independent $\cN(0,1)$ entries, of rank $N$ almost surely (\cref{thm:toolkit:inv-wishart}). The estimator is $\hat w_N=w_{\star,\le N}+(X_N^\top X_N)^{-1}X_N^\top\xi$, and
\[
  \Lambda^{1/2}(\hat w_N-w_{\star,\le N})=\Lambda^{1/2}\bigl(\Lambda^{1/2}Z^\top Z\Lambda^{1/2}\bigr)^{-1}\Lambda^{1/2}Z^\top\xi=(Z^\top Z)^{-1}Z^\top\xi .
\]
By \eqref{eq:scaling:loss} the last sum in $\Loss(\hat w_N)$ is the squared norm of this vector.

\emph{Averaging.} Conditionally on $Z$, the squared norm has mean $(\sigma^2+T(N))\tr\bigl((Z^\top Z)^{-1}Z^\top Z(Z^\top Z)^{-1}\bigr)=(\sigma^2+T(N))\tr\bigl((Z^\top Z)^{-1}\bigr)$. By the inverse Wishart mean (\cref{thm:toolkit:inv-wishart}, or \cref{lem:doubledescent:wishart}), $\E\tr\bigl((Z^\top Z)^{-1}\bigr)=N/(D-N-1)$ for $N\le D-2$, and it is $+\infty$ for $N\in\{D-1,D\}$. Adding $\sigma^2+T(N)$ from \eqref{eq:scaling:loss} gives $(\sigma^2+T(N))\bigl(1+\frac N{D-N-1}\bigr)$, which is \eqref{eq:scaling:exact}.
\end{proof}

The degenerate size $N=0$ (the zero predictor) has loss $\sigma^2+T(0)$, which \eqref{eq:scaling:exact} also gives. The formula says that a model of size $N$ pays its approximation error, and pays it again, together with the noise, inflated by the variance factor $N/(D-N-1)$: the unlearned part of the target is noise for the learner. This self-consistent noise returns in the deterministic equivalent below. In the power-law case \cref{thm:scaling:exact} turns into scaling laws.

\begin{corollary}[Two scaling laws and the optimal size]\label[corollary]{cor:scaling:laws}
Under \cref{ass:scaling:features}, suppose that $T(N)\sim cN^{-b}$ as $N\to\infty$, for some $c>0$ and $b>0$ (in the power-law case $c=1/b$), and that $\sigma>0$. Let $\cE(N,D)=\E\Loss(\hat w_N)-\sigma^2$ and $\cE^\star(D)=\min_{0\le N\le D-2}\cE(N,D)$.
\begin{enumerate}
\item \emph{Parameter-limited.} For fixed $N$, $\cE(N,D)\to T(N)\sim cN^{-b}$ as $D\to\infty$.
\item \emph{Data-limited.} As $D\to\infty$,
\begin{equation}\label{eq:scaling:datalaw}
  \cE^\star(D)\sim(1+b)\,c^{\frac1{1+b}}\Bigl(\frac{\sigma^2}{bD}\Bigr)^{\frac b{1+b}},
\end{equation}
and every minimizer $N_D$ satisfies $N_D\sim(bcD/\sigma^2)^{1/(1+b)}$.
\end{enumerate}
\end{corollary}

\begin{proof}
Part (i) follows from \eqref{eq:scaling:exact}, since $(D-1)/(D-N-1)\to1$. For part (ii), the plan is to sandwich $\cE(N,D)$ between multiples of $h_D(N)=cN^{-b}+\sigma^2N/D$, whose minimum we compute first.

\emph{The comparison function.} Let $M_D=(bcD/\sigma^2)^{1/(1+b)}$, the critical point of $h_D$, and $\epsilon_D=cM_D^{-b}$. Since $\sigma^2M_D/D=\sigma^2M_D^{1+b}/(DM_D^b)=b\,\epsilon_D$, substituting $N=sM_D$ gives $h_D(sM_D)=\epsilon_D\,g(s)$ with $g(s)=s^{-b}+bs$. The function $g$ is strictly convex on $(0,\infty)$, with $g'(s)=b(1-s^{-1-b})$, so its minimum is $g(1)=1+b$, and for every $\eta\in(0,1)$ the number $g_\eta=\min\{g(1-\eta),g(1+\eta)\}$ exceeds $1+b$ and bounds $g$ from below off $(1-\eta,1+\eta)$. The right side of \eqref{eq:scaling:datalaw} is $(1+b)\epsilon_D$, by the definition of $M_D$.

\emph{Upper bound.} Take $N=\lceil M_D\rceil$. Since $M_D\to\infty$ and $M_D/D\to0$, we have $T(N)=cN^{-b}(1+o(1))=\epsilon_D(1+o(1))$ and $(\sigma^2+T(N))\frac N{D-N-1}=\sigma^2\frac{M_D}D(1+o(1))=b\epsilon_D(1+o(1))$. Hence $\cE^\star(D)\le\cE(N,D)=(1+b)\epsilon_D(1+o(1))$.

\emph{Lower bound.} Since $\frac N{D-N-1}\ge\frac ND$ and $\sigma^2+T(N)\ge\sigma^2$, every $N\le D-2$ has $\cE(N,D)\ge T(N)+\sigma^2N/D$. Fix $\delta\in(0,1)$ and choose $N_0$ with $T(N)\ge(1-\delta)cN^{-b}>0$ for all $N\ge N_0$. For $N\ge N_0$, $\cE(N,D)\ge(1-\delta)h_D(N)=(1-\delta)\epsilon_Dg(N/M_D)\ge(1-\delta)(1+b)\epsilon_D$. For $N<N_0$, $\cE(N,D)\ge T(N_0)$, because $T$ is nonincreasing, and $T(N_0)>0$ exceeds $(1+b)\epsilon_D$ once $D$ is large, because $\epsilon_D\to0$. Hence $\liminf\cE^\star(D)/\epsilon_D\ge(1-\delta)(1+b)$ for every $\delta$, which with the upper bound proves \eqref{eq:scaling:datalaw}.

\emph{The minimizer.} Fix $\eta\in(0,1)$, and choose $\delta>0$ so small that $(1-\delta)g_\eta>(1+\delta)(1+b)$. For large $D$, the upper bound gives $\cE^\star(D)<(1+\delta)(1+b)\epsilon_D$, while every $N\ge N_0$ with $\abs{N/M_D-1}\ge\eta$ has $\cE(N,D)\ge(1-\delta)g_\eta\epsilon_D$, and every $N<N_0$ has $\cE(N,D)\ge T(N_0)$; neither can be a minimizer. So $\abs{N_D/M_D-1}<\eta$ for large $D$.
\end{proof}

The proof shows where the data exponent comes from. At the optimum the approximation error $cN^{-b}$ and the variance $\sigma^2N/D$ are of the same order, so the model keeps exactly the features whose share of the loss, about $N^{-(1+b)}$ per feature, exceeds the estimation cost $\sigma^2/D$ per feature; this gives $N_D\asymp D^{1/(1+b)}$ and $\cE^\star\asymp N_D^{-b}=D^{-b/(1+b)}$. Without noise the variance has nothing to amplify except $T(N)$ itself, the best model has a size proportional to $D$, and the data exponent rises to $b$ (\cref{ex:scaling:noiseless}).\sidenote{This is the distinction of \textcite{BahriEtAl2024} between variance-limited and resolution-limited scaling, here in one formula: the factor $(D-1)/(D-N-1)$ is the variance, and $T(N)$ the resolution.}

\section{The solution: the spectrum, ridge, and Zipf's law}\label{sec:scaling:spectrum}

Least squares on the top $N$ features never sees the variances $\lambda_j$. A learner that uses all features and controls its complexity by a ridge parameter $\lambda$, or by its training time, does: then the spectrum decides how many features are fitted.

\paragraph{The sequence model.} Strip the regression to its normal modes. For each $j$, observe $z_j=\theta_j+\frac\sigma{\sqrt D}g_j$, with $\theta_j=\sqrt{\lambda_j}\,w_{\star,j}$ the target in units of prediction and independent $g_j\sim\cN(0,1)$: the per-feature estimation noise of \cref{thm:scaling:exact} is $\sigma^2/D$ when $N\ll D$. Ridge regression multiplies the $j$th coordinate by $\lambda_j/(\lambda_j+\lambda)$, as in \cref{ch:doubledescent}; the estimator $\hat\theta_j=\frac{\lambda_j}{\lambda_j+\lambda}z_j$ has excess loss
\begin{equation}\label{eq:scaling:seq}
  \cE_{\mathrm{seq}}(\lambda)=\E\sum_j(\hat\theta_j-\theta_j)^2=\underbrace{\sum_j\lambda_jw_{\star,j}^2\Bigl(\frac\lambda{\lambda_j+\lambda}\Bigr)^2}_{\text{bias}}+\underbrace{\frac{\sigma^2}D\sum_j\Bigl(\frac{\lambda_j}{\lambda_j+\lambda}\Bigr)^2}_{\text{variance}} .
\end{equation}

\begin{proposition}[Ridge on a power-law spectrum]\label[proposition]{prop:scaling:ridge}
In the power-law case \eqref{eq:scaling:powerlaw}, let $\frac12<a$ and $0<b<2a$, and let $J=\lambda^{-1/a}$. There are constants $0<c_1<c_2$, depending only on $a$ and $b$, such that for $0<\lambda\le1$
\[
  c_1J^{-b}\le\text{bias}\le c_2J^{-b},\qquad c_1J\le\sum_j\Bigl(\frac{\lambda_j}{\lambda_j+\lambda}\Bigr)^2\le c_2J .
\]
Consequently, if $\sigma^2\le D$, then $\min_{\lambda>0}\cE_{\mathrm{seq}}(\lambda)\asymp(\sigma^2/D)^{b/(1+b)}$, attained up to constants at $J=(D/\sigma^2)^{1/(1+b)}$.
\end{proposition}

\begin{proof}
The plan is to split the sums at $j=J$, where $\lambda_j=\lambda$, and compare each part with a power.

\emph{Two regimes of the filter.} Write $r_j=\lambda j^a=(j/J)^a$, so that $\frac\lambda{\lambda_j+\lambda}=\frac{r_j}{1+r_j}$ and $\frac{\lambda_j}{\lambda_j+\lambda}=\frac1{1+r_j}$. For $j\le J$ we have $r_j\le1$, so $\frac{r_j}{1+r_j}\in[\frac{r_j}2,r_j]$ and $\frac1{1+r_j}\in[\frac12,1]$; for $j>J$ we have $r_j>1$, so $\frac{r_j}{1+r_j}\in[\frac12,1]$ and $\frac1{1+r_j}\in[\frac1{2r_j},\frac1{r_j}]$. Since $\lambda\le1$, $J\ge1$.

\emph{The bias.} By the first regime, the bias lies between $\frac14$ and $1$ times
\[
  \sum_{j\le J}j^{-(1+b)}\Bigl(\frac jJ\Bigr)^{2a}+\sum_{j>J}j^{-(1+b)}=J^{-2a}\sum_{j\le J}j^{2a-1-b}+\sum_{j>J}j^{-(1+b)} .
\]
The exponent $2a-1-b$ exceeds $-1$ because $b<2a$, so $\sum_{j\le J}j^{2a-1-b}$ lies between two constant multiples of $J^{2a-b}$ for $J\ge1$ (compare with $\int_0^{J+1}t^{2a-1-b}\dd t$ and with the single term $j=1$ when $J<2$); and $\sum_{j>J}j^{-(1+b)}$ lies between two constant multiples of $J^{-b}$ by \cref{lem:scaling:tail} applied at $\lfloor J\rfloor$. Both parts are of order $J^{-b}$.

\emph{The variance sum.} By the two regimes it lies between $\frac14$ and $1$ times $\#\{j\le J\}+\sum_{j>J}(J/j)^{2a}$. The first term is $\lfloor J\rfloor\in[J/2,J]$, and the second is at most $J^{2a}\bigl(J^{-2a}+\int_J^\infty t^{-2a}\dd t\bigr)=1+\frac J{2a-1}$, finite because $2a>1$.

\emph{The optimum.} With $\epsilon=\sigma^2/D\le1$, the bounds give $\cE_{\mathrm{seq}}\asymp J^{-b}+\epsilon J$ on $\lambda\le1$. At $J=\epsilon^{-1/(1+b)}\ge1$ both terms equal $\epsilon^{b/(1+b)}$. Conversely, for every $J\ge1$, either $J\le\epsilon^{-1/(1+b)}$ and $J^{-b}\ge\epsilon^{b/(1+b)}$, or $J>\epsilon^{-1/(1+b)}$ and $\epsilon J\ge\epsilon^{b/(1+b)}$. For $\lambda>1$ every factor $\frac\lambda{\lambda_j+\lambda}$ is at least $\frac12$, so the bias is at least $\frac14\zeta(1+b)$, which is at least a constant times $\epsilon^{b/(1+b)}$.
\end{proof}

So ridge on all features behaves like least squares on the top $J=\lambda^{-1/a}$ features: the spectrum exponent $a$ only converts the regularization into an effective model size, and the exponents are those of \cref{cor:scaling:laws}. The same holds for training time. Gradient flow on the population loss from zero fits the $j$th coordinate up to the factor $1-e^{-\lambda_jt}$, so it has learned the features with $\lambda_jt\gtrsim1$, an effective size $J=t^{1/a}$, and its excess loss decays like $t^{-b/a}$ (\cref{ex:scaling:time}). Spectral bias, the order in which kernel gradient descent learns eigenfunctions (\cref{ch:lazy}), becomes a law in time with exponent $b/a$. If $b>2a$, the bias of ridge saturates at order $\lambda^2$ and the exponent drops to $2a/(1+2a)$ (\cref{ex:scaling:saturation}): the target is smoother than ridge can exploit.

\paragraph{From the sequence model to Gaussian designs.} With a finite Gaussian sample the spectrum of $\frac1DX^\top X$ is distorted, as the Marchenko--Pastur law of \cref{ch:doubledescent} showed for $\lambda_j\equiv1$. The distortion is captured by a renormalized ridge parameter.

\begin{heuristic}[The deterministic equivalent]
Under \cref{ass:scaling:features} with $\sum_j\lambda_j<\infty$, let $\hat w_\lambda=(X^\top X+D\lambda I)^{-1}X^\top y$ use all features. Let $\kappa>\lambda$ solve $\kappa=\lambda+\frac\kappa D\sum_j\frac{\lambda_j}{\lambda_j+\kappa}$ (for $\lambda=0$ and $D<P$: $\sum_j\frac{\lambda_j}{\lambda_j+\kappa}=D$), and let $\gamma=\frac1D\sum_j\bigl(\frac{\lambda_j}{\lambda_j+\kappa}\bigr)^2<1$. The excess loss $R$ of $\hat w_\lambda$ is then close to the solution of
\begin{equation}\label{eq:scaling:de}
  R=\sum_j\lambda_jw_{\star,j}^2\Bigl(\frac\kappa{\lambda_j+\kappa}\Bigr)^2+\gamma\,(\sigma^2+R),\qquad\text{that is,}\qquad R=\frac{\text{bias}(\kappa)+\gamma\sigma^2}{1-\gamma}.
\end{equation}
In words: the Gaussian design behaves like the sequence model \eqref{eq:scaling:seq} with ridge $\kappa$ instead of $\lambda$ and with noise $\sigma^2+R$ instead of $\sigma^2$, the unlearned signal acting as noise exactly as $T(N)$ does in \eqref{eq:scaling:exact}. For $\lambda_j\equiv1$ on $P$ features, $\kappa$ is the inverse of the Stieltjes-transform variable of \cref{thm:doubledescent:ridge}, and \eqref{eq:scaling:de} reproduces that theorem (\cref{ex:scaling:isotropic}). \Textcite{ChengMontanari2024} prove nonasymptotic bounds of the form $(1\pm\Delta)$ for the bias and the variance under conditions on effective ranks that hold for power laws; \textcite{AtanasovEtAl2026} derive \eqref{eq:scaling:de} as a renormalization, and \textcite{BordelonEtAl2020} obtained it for kernel regression with methods of statistical physics.
\end{heuristic}

For the power law, the equation for $\kappa$ reads $\kappa-\lambda\asymp\kappa\,\kappa^{-1/a}/D$, so $\kappa\asymp\max(\lambda,D^{-a})$, and $\gamma$ stays bounded away from $1$; inserting \cref{prop:scaling:ridge} gives an excess loss of order $D^{-b/(1+b)}$ for optimally tuned ridge with noise, and of order $\kappa_0^{b/a}\asymp D^{-b}$ for the noiseless minimum-norm interpolator, the two exponents of \cref{cor:scaling:laws} and \cref{ex:scaling:noiseless}. \Textcite{CuiEtAl2021} describe the crossover between them as $D$ grows, and \textcite{CaponnettoDeVito2007} proved the noisy rate, with matching lower bounds, for kernel ridge regression under source and capacity conditions.

\paragraph{Zipf's law.} The second model gives the same exponents with no regression at all.

\begin{lemma}[Integral of a unimodal sequence]\label[lemma]{lem:scaling:unimodal}
Let $g\colon\R\to[0,\infty)$ be integrable, nondecreasing on $(-\infty,m]$ and nonincreasing on $[m,\infty)$. Then $\bigl\lvert\sum_{k\in\Z}g(k)-\int g\bigr\rvert\le2\sup g$.
\end{lemma}

\begin{proof}
Let $k_0=\lfloor m\rfloor$. For $k\le k_0-1$ the interval $[k,k+1]$ lies in $(-\infty,m]$, where $g$ increases, so $g(k)\le\int_k^{k+1}g$; for $k\ge k_0+2$ the interval $[k-1,k]$ lies in $[m,\infty)$, so $g(k)\le\int_{k-1}^kg$. These intervals do not overlap $(k_0,k_0+1)$, and the two remaining terms are at most $\sup g$ each, so $\sum_kg(k)\le\int g+2\sup g$. In the other direction, $g(k)\ge\int_{k-1}^kg$ for $k\le k_0$ and $g(k)\ge\int_k^{k+1}g$ for $k\ge k_0+1$; these intervals cover $\R$ except $(k_0,k_0+1)$, whose integral is at most $\sup g$.
\end{proof}

\begin{theorem}[Zipf learning curves]\label{thm:scaling:zipf}
Under \cref{ass:scaling:quanta}, the memorizing learner trained on $D$ examples errs with probability
\begin{equation}\label{eq:scaling:zipfexact}
  M(D)=\sum_{k\ge1}p_k(1-p_k)^D\sim C_b\,D^{-\frac b{1+b}},\qquad C_b=\frac{\Gamma\bigl(\frac b{1+b}\bigr)}{(1+b)\,\zeta(1+b)^{\frac1{1+b}}},
\end{equation}
as $D\to\infty$. The learner of capacity $N$ that holds the items $1,\dots,N$ errs with probability $\sum_{k>N}p_k\sim N^{-b}/\bigl(b\,\zeta(1+b)\bigr)$.
\end{theorem}

\begin{proof}
The plan is to replace $(1-p)^D$ by $e^{-Dp}$, up to a factor that is harmless where the sum lives, and to evaluate the resulting sum as an integral.

\emph{The exact formula.} The test item is independent of the training items, and the learner errs exactly when the test item is $k$ and none of the $D$ training items is $k$; this has probability $p_k(1-p_k)^D$. Summing over $k$ gives $M(D)$.

\emph{The exponential sum.} Let $s=1+b$, $Z=\zeta(s)$, and $S(c)=\sum_kp_ke^{-cp_k}$ for $c>0$. Then $S(c)=\sum_kg(k)$ with $g(x)=\frac{x^{-s}}Ze^{-cx^{-s}/Z}$ for $x>0$ and $g(x)=0$ for $x\le0$. In the variable $u=cx^{-s}/Z$, which decreases from $\infty$ to $0$ as $x$ increases, $g=\frac uce^{-u}$, which increases for $u>1$ and decreases for $u<1$; so $g$ is unimodal in $x$, with $\sup g=\frac1{ec}$. Changing variables, with $x=(c/(Zu))^{1/s}$ and $\dd x=-\frac1s(c/Z)^{1/s}u^{-1/s-1}\dd u$,
\[
  \int_0^\infty g(x)\dd x=\frac1s\Bigl(\frac cZ\Bigr)^{1/s}\frac1c\int_0^\infty u^{-1/s}e^{-u}\dd u=\frac{\Gamma(1-\frac1s)}{s\,Z^{1/s}}\,c^{\frac1s-1}=C_b\,c^{-\frac b{1+b}},
\]
since $1-\frac1s=\frac b{1+b}$. By \cref{lem:scaling:unimodal}, $\abs{S(c)-C_bc^{-b/(1+b)}}\le\frac2{ec}$, so $S(c)\sim C_bc^{-b/(1+b)}$ as $c\to\infty$.

\emph{Comparing $(1-p)^D$ with $e^{-Dp}$.} For $0\le p<1$, $\log(1-p)\le-p$ and $-\log(1-p)=\sum_{i\ge1}p^i/i\le p/(1-p)$, so
\[
  e^{-Dp/(1-p)}\le(1-p)^D\le e^{-Dp} .
\]
The upper bound gives $M(D)\le S(D)$. For the lower bound fix $\delta\in(0,1)$ and let $K$ be the number of items with $p_k>\delta$. Every $k>K$ has $e^{-Dp_k/(1-p_k)}\ge e^{-Dp_k/(1-\delta)}$, and the items $k\le K$ contribute at most $e^{-Dp_K/(1-\delta)}\le e^{-D\delta}$ to $S(D/(1-\delta))$ (if $K\ge1$); hence $M(D)\ge S\bigl(\frac D{1-\delta}\bigr)-e^{-D\delta}$. With the asymptotics of $S$, $\liminf_DD^{b/(1+b)}M(D)\ge C_b(1-\delta)^{b/(1+b)}$, and letting $\delta\to0$ proves \eqref{eq:scaling:zipfexact}.

\emph{Capacity.} The learner holding the items $1,\dots,N$ errs exactly on the items $k>N$, with probability $\sum_{k>N}p_k=Z^{-1}\sum_{k>N}k^{-(1+b)}$; apply \cref{lem:scaling:tail}.
\end{proof}

The constants are explicit: $C_{1/2}\approx0.942$, $C_1=\sqrt6/(2\sqrt\pi)\approx0.691$, and $C_2\approx0.425$. Once more the parameter exponent is $b$ and the data exponent is $b/(1+b)$, and for the same reason: item $k$ is learned once it has been seen, which happens when $Dp_k\gtrsim1$, so the effective capacity after $D$ examples is $N\asymp D^{1/(1+b)}$, and the error is the tail beyond it, $N^{-b}=D^{-b/(1+b)}$.\sidenote{\Textcite{CabannesEtAl2024} derive scaling laws of this kind for associative memories built from outer products of embeddings, where capacity and data limit the recall of Zipf-distributed associations together.}

\section{The solution: spending compute}\label{sec:scaling:compute}

If the loss has the additive form \eqref{eq:scaling:chinchilla} and compute costs $C=\kappa ND$ (with $\kappa=6$ for transformers), how should a budget $C$ be split?

\begin{theorem}[Compute-optimal allocation]\label{thm:scaling:compute}
Let $A,B,\alpha_N,\alpha_D,\kappa>0$, and for $C>0$ minimize $\Loss(N,D)=\Loss_\infty+AN^{-\alpha_N}+BD^{-\alpha_D}$ over $N,D>0$ subject to $\kappa ND=C$. Let $\nu_N=\frac{\alpha_D}{\alpha_N+\alpha_D}$, $\nu_D=\frac{\alpha_N}{\alpha_N+\alpha_D}$, and $G=\bigl(\frac{\alpha_NA}{\alpha_DB}\bigr)^{1/(\alpha_N+\alpha_D)}$. The minimizer is unique and equals
\begin{equation}\label{eq:scaling:alloc}
  N^\star=G\Bigl(\frac C\kappa\Bigr)^{\nu_N},\qquad D^\star=G^{-1}\Bigl(\frac C\kappa\Bigr)^{\nu_D}.
\end{equation}
At the minimum the two reducible terms are in the ratio $AN^{\star-\alpha_N}:BD^{\star-\alpha_D}=\alpha_D:\alpha_N$, and
\begin{equation}\label{eq:scaling:lossC}
  \Loss(N^\star,D^\star)-\Loss_\infty=\frac{\alpha_N+\alpha_D}{\alpha_D}\,A\,G^{-\alpha_N}\Bigl(\frac C\kappa\Bigr)^{-\frac{\alpha_N\alpha_D}{\alpha_N+\alpha_D}} .
\end{equation}
\end{theorem}

\begin{proof}
The plan is to eliminate $D$ and minimize a strictly convex function of $u=\log N$. On the constraint, $D=C/(\kappa N)$, and the reducible loss is
\[
  \varphi(u)=Ae^{-\alpha_Nu}+B\Bigl(\frac\kappa C\Bigr)^{\alpha_D}e^{\alpha_Du} .
\]
Each term is strictly convex in $u$, and $\varphi(u)\to\infty$ as $u\to\pm\infty$, so $\varphi$ has a unique minimizer, the unique zero of $\varphi'(u)=-\alpha_NAN^{-\alpha_N}+\alpha_DBD^{-\alpha_D}$ (with $N=e^u$ and $D=C/(\kappa N)$). The condition $\alpha_NAN^{-\alpha_N}=\alpha_DBD^{-\alpha_D}$ gives the ratio of the two terms. Substituting $D=C/(\kappa N)$ turns it into $N^{\alpha_N+\alpha_D}=\frac{\alpha_NA}{\alpha_DB}(C/\kappa)^{\alpha_D}$, which is the first formula of \eqref{eq:scaling:alloc}, and $D^\star=C/(\kappa N^\star)$ is the second, because $1-\nu_N=\nu_D$. Finally the loss is $AN^{\star-\alpha_N}(1+\frac{\alpha_N}{\alpha_D})$, and $N^{\star-\alpha_N}=G^{-\alpha_N}(C/\kappa)^{-\alpha_N\alpha_D/(\alpha_N+\alpha_D)}$.
\end{proof}

Four readings. The allocation exponents $\nu_N$ and $\nu_D$ depend only on the ratio of the exponents: the faster-decaying term gets the smaller share of the budget's growth. Parameters and data grow at the same rate exactly when $\alpha_N=\alpha_D$; otherwise the number of tokens per parameter, $D^\star/N^\star=G^{-2}(C/\kappa)^{(\alpha_N-\alpha_D)/(\alpha_N+\alpha_D)}$, drifts with the budget. The loss exponent in compute, $\frac{\alpha_N\alpha_D}{\alpha_N+\alpha_D}$, is half the harmonic mean of the two exponents, so it is smaller than both. And the constants $A$ and $B$ enter only $G$: the exponents of the allocation do not depend on them, but the number of tokens per parameter does, which is why it is so sensitive to the fit. With the rounded exponents $\alpha_N=0.34$ and $\alpha_D=0.28$ of \textcite{HoffmannEtAl2022}, $\nu_N=0.28/0.62\approx0.45$, $\nu_D\approx0.55$, the tokens per parameter grow like $C^{0.10}$, and the compute-optimal loss falls like $C^{-0.154}$. A data exponent equal to the parameter exponent, closer to the refit of \textcite{BesirogluEtAl2024}, gives $\nu_N=\nu_D=\frac12$.

\section{Reading the solution}\label{sec:scaling:reading}

\Cref{fig:scaling:curves}(a) compares \cref{thm:scaling:exact} with simulations. Each fixed-size model follows a curve that falls like $1/D$ and then flattens at its approximation error $T(N)$: it is parameter-limited. The lower envelope of the curves, the best size at each $D$, is the data-limited law \eqref{eq:scaling:datalaw}, with slope $-\frac b{1+b}=-\frac12$ here; at $D=2048$ the best size is $N=96$, against $M_D=101$ from \cref{cor:scaling:laws}. Below the envelope no model of the family can go, and above it each model wastes either parameters or data. The smooth law is thus the envelope of a family of curves that each saturate, which is the shape of the compute-efficient frontier in \textcite{KaplanEtAl2020}.

\begin{figure}[t]
\begin{fullwidth}
\centering
\begin{tikzpicture}
\begin{groupplot}[group style={group size=2 by 1,horizontal sep=16mm},hplot,
  width=0.5\linewidth,height=0.4\linewidth,xmode=log,ymode=log,unbounded coords=jump,
  log basis x=2]
\nextgroupplot[xlabel={sample size $D$},ylabel={excess loss},xmin=6,xmax=2300,ymin=0.008,ymax=1.2,
  title={(a) least squares on the top $N$ features}]
\addplot[muted,line width=0.8pt] table[x=D,y=N4] {figures/data/16-scaling-ols-theory.dat};
\addplot[sandline!55!black,line width=0.9pt,dashdotted] table[x=D,y=N16] {figures/data/16-scaling-ols-theory.dat};
\addplot[theory2] table[x=D,y=N64] {figures/data/16-scaling-ols-theory.dat};
\addplot[seadeep,line width=0.8pt,densely dotted] table[x=D,y=N256] {figures/data/16-scaling-ols-theory.dat};
\addplot[theory,line width=1.6pt] table[x=D,y=best] {figures/data/16-scaling-ols-theory.dat};
\addplot[guide] table[x=D,y=asym] {figures/data/16-scaling-ols-theory.dat};
\addplot[simulation] table[x=D,y=N4] {figures/data/16-scaling-ols-sim.dat};
\addplot[simulation] table[x=D,y=N16] {figures/data/16-scaling-ols-sim.dat};
\addplot[simulation] table[x=D,y=N64] {figures/data/16-scaling-ols-sim.dat};
\addplot[simulation] table[x=D,y=N256] {figures/data/16-scaling-ols-sim.dat};
\node[font=\scriptsize,text=ink,anchor=south east] at (axis cs:2000,0.225) {$N=4$};
\node[font=\scriptsize,text=ink,anchor=south east] at (axis cs:2000,0.064) {$16$};
\node[font=\scriptsize,text=ink,anchor=north east] at (axis cs:1500,0.026) {$64$};
\node[font=\scriptsize,text=ink,anchor=west] at (axis cs:560,0.2) {$256$};
\node[font=\scriptsize,text=sea,anchor=north east] at (axis cs:300,0.03) {best $N$};
\nextgroupplot[xlabel={sample size $D$},xmin=6,xmax=1200,ymin=0.0015,ymax=0.6,
  title={(b) ridge on all features},legend pos=south west]
\addplot[theory] table[x=D,y=optimal] {figures/data/16-scaling-ridge-theory.dat};
\addlegendentry{optimal ridge, $\sigma^2=0.2$}
\addplot[theory2] table[x=D,y=noiseless] {figures/data/16-scaling-ridge-theory.dat};
\addlegendentry{minimum norm, $\sigma=0$}
\addplot[simulation,forget plot] table[x=D,y=optimal] {figures/data/16-scaling-ridge-sim.dat};
\addplot[simulation2,forget plot] table[x=D,y=noiseless] {figures/data/16-scaling-ridge-sim.dat};
\addplot[guide,domain=8:1024,samples=2,forget plot] {0.9*x^(-0.5)};
\addplot[guide,domain=8:1024,samples=2,forget plot] {1.6*x^(-1)};
\node[font=\scriptsize,text=muted,anchor=south west] at (axis cs:128,0.085) {$D^{-1/2}$};
\node[font=\scriptsize,text=muted,anchor=north east] at (axis cs:100,0.012) {$D^{-1}$};
\end{groupplot}
\end{tikzpicture}
\caption{Learning curves on power-law features with $a=\frac32$ and $b=1$. (a)~Least squares on the top $N\in\{4,16,64,256\}$ of $P=1024$ features, $\sigma^2=0.2$. Thin curves: \cref{thm:scaling:exact}; thick curve: its minimum over $N$; dashed: the asymptote \eqref{eq:scaling:datalaw} for $P=\infty$. Points: averages over $150$ training sets (seed $1617$), scored by their exact test loss; the largest deviation is $1.4$ standard errors. Points with $D<2N+4$, where the simulated average has a large variance, are omitted. (b)~Ridge on all $P=4000$ features: minimum-norm interpolation without noise (squares), and ridge with the $\lambda$ that minimizes the deterministic equivalent, $\sigma^2=0.2$ (circles); curves: the deterministic equivalent \eqref{eq:scaling:de}; $60$ training sets per point. The equivalent is asymptotic: at $D\le16$ the simulations lie up to $18\%$ above it, and from $D=64$ on within $4\%$.}
\label{fig:scaling:curves}
\end{fullwidth}
\end{figure}

\Cref{fig:scaling:curves}(b) shows the two exponents for ridge on all features: $D^{-b/(1+b)}$, here $D^{-1/2}$, with noise and optimal ridge, and $D^{-b}$, here $D^{-1}$, for the noiseless interpolator, as predicted by \cref{prop:scaling:ridge} and the deterministic equivalent. With noise, the local slope is steeper than $\frac12$ at moderate $D$, because the unlearned signal adds to the noise in \eqref{eq:scaling:de} and shrinks as $D$ grows; for smaller $b$ this transient lasts longer (\cref{ex:scaling:simridge}).

\paragraph{What sets the exponents.} In both models the exponents are $b$ for parameters and $b/(1+b)$ for data, and $b$ is a property of the data: how the importance of directions or items decays. Neither the model class nor the spectrum exponent $a$ appears, except through saturation when $b>2a$ and through the time exponent $b/a$. This is the solvable version of the observation of \textcite{HestnessEtAl2017} that exponents depend on the task more than on the architecture. Three consequences can be checked. First, the data exponent is smaller than the parameter exponent, and determined by it: $\alpha_D=\alpha_N/(1+\alpha_N)$. Second, data that are less heavy-tailed (larger $b$) scale faster. Third, small exponents mean $b$ close to $0$, a nearly flat hierarchy: the frequencies of words, $p_k\propto k^{-s}$ with $s\approx1$, correspond to $b\approx0$ \parencite{Hutter2021}, and a model limited by rare words would scale very slowly.

\paragraph{Emergence on a smooth curve.} \Cref{fig:scaling:zipf}(b) follows single quanta. Item $k$ is known after $D$ examples with probability $1-(1-p_k)^D\approx1-e^{-Dp_k}$, which rises from $0.1$ to $0.9$ as $Dp_k$ grows from $0.105$ to $2.30$, a factor of about $22$ in data, at $D\approx1/p_k$. Every skill turns on over the same window in $\log D$, at a position set by its frequency, while their sum, the loss, is a power law without a feature. A downstream task that needs a few specific rare quanta shows a jump; the loss does not. This is the quanta picture of emergence of \textcite{MichaudEtAl2023}.

\begin{figure}[t]
\begin{fullwidth}
\centering
\begin{tikzpicture}
\begin{groupplot}[group style={group size=2 by 1,horizontal sep=15mm},hplot,
  width=0.49\linewidth,height=0.4\linewidth,xmode=log]
\nextgroupplot[ymode=log,xlabel={training examples $D$},ylabel={error $M(D)$},xmin=1,xmax=1e6,ymin=1e-5,ymax=1,
  title={(a) memorizing learner},legend pos=south west]
\addplot[theory] table[x=D,y=exact0] {figures/data/16-scaling-zipf-theory.dat};
\addlegendentry{$b=\frac12$}
\addplot[theory2] table[x=D,y=exact1] {figures/data/16-scaling-zipf-theory.dat};
\addlegendentry{$b=1$}
\addplot[theory3] table[x=D,y=exact2] {figures/data/16-scaling-zipf-theory.dat};
\addlegendentry{$b=2$}
\addplot[guide,forget plot] table[x=D,y=asym0] {figures/data/16-scaling-zipf-theory.dat};
\addplot[guide,forget plot] table[x=D,y=asym1] {figures/data/16-scaling-zipf-theory.dat};
\addplot[guide,forget plot] table[x=D,y=asym2] {figures/data/16-scaling-zipf-theory.dat};
\addplot[simulation,forget plot] table[x=D,y=sim0] {figures/data/16-scaling-zipf-sim.dat};
\addplot[simulation,forget plot] table[x=D,y=sim1] {figures/data/16-scaling-zipf-sim.dat};
\addplot[simulation,forget plot] table[x=D,y=sim2] {figures/data/16-scaling-zipf-sim.dat};
\nextgroupplot[xlabel={training examples $D$},ylabel={probability},xmin=1,xmax=1e7,ymin=0,ymax=1.02,
  title={(b) single quanta and the loss, $b=1$}]
\addplot[theory3] table[x=D,y=k1] {figures/data/16-scaling-zipf-items.dat};
\addplot[theory3] table[x=D,y=k10] {figures/data/16-scaling-zipf-items.dat};
\addplot[theory3] table[x=D,y=k100] {figures/data/16-scaling-zipf-items.dat};
\addplot[theory3] table[x=D,y=k1000] {figures/data/16-scaling-zipf-items.dat};
\addplot[theory] table[x=D,y=loss] {figures/data/16-scaling-zipf-items.dat};
\node[font=\scriptsize,text=ink,anchor=west] at (axis cs:2.5,0.75) {$k=1$};
\node[font=\scriptsize,text=ink,anchor=east] at (axis cs:70,0.55) {$10$};
\node[font=\scriptsize,text=ink,anchor=east] at (axis cs:7e3,0.55) {$100$};
\node[font=\scriptsize,text=ink,anchor=east] at (axis cs:7e5,0.55) {$1000$};
\node[font=\scriptsize,text=sea,anchor=west] at (axis cs:1.3,0.13) {error $M(D)$};
\end{groupplot}
\end{tikzpicture}
\caption{The quanta model. (a)~Error of the memorizing learner for $b\in\{\frac12,1,2\}$, items truncated at $K=10^7$. Curves: the exact sum \eqref{eq:scaling:zipfexact}; dashed: the asymptote $C_bD^{-b/(1+b)}$ of \cref{thm:scaling:zipf} for the untruncated law; points: $20$ simulated training sets per point (seed $1618$), each scored by the exact mass of its unseen items, within $2.2$ standard errors of the curves. (b)~For $b=1$: the probability that item $k\in\{1,10,100,1000\}$ has been seen (sigmoids, each spanning a factor of about $22$ in $D$ from $0.1$ to $0.9$), and the error $M(D)$, which is smooth.}
\label{fig:scaling:zipf}
\end{fullwidth}
\end{figure}

\paragraph{Allocation in the solvable models.} The form \eqref{eq:scaling:chinchilla} is additive; \cref{thm:scaling:exact} is not, because the variance $\sigma^2N/D$ couples $N$ and $D$. Minimizing $cN^{-b}+\sigma^2N/D$ under $6ND=C$ gives $N^\star\propto C^{1/(2+b)}$ and $D^\star\propto C^{(1+b)/(2+b)}$: with $b=0.34$, parameters grow like $C^{0.43}$ (\cref{ex:scaling:olscompute}). For a model trained by gradient flow on fresh data, whose cost is proportional to $N$ times the training time, \cref{ex:scaling:time} gives $N^\star\propto C^{1/(1+a)}$ instead, and the spectrum exponent enters the allocation.

\section{What the model gets wrong}\label{sec:scaling:wrong}

\paragraph{What was thrown away.} \emph{Feature learning.} The features, their spectrum, and the target's decay in their basis are fixed in advance, so the exponent $b$ is an input; a network chooses its representation, and with it the basis in which importance is measured. \emph{Optimization.} The estimators are closed-form; compute enters only through the accounting $C=6ND$, not through the dynamics of stochastic gradient descent, its noise, its step size, or its schedule. \emph{Architecture.} Only the number of parameters matters, and they are spent on the top features in order. \emph{The irreducible loss.} White noise of variance $\sigma^2$ stands for the entropy of the data, which in language is structured. \emph{Repetition.} Each example is fresh and used once. \emph{Composition.} Quanta are independent and are learned from one example; real skills share parts and need many examples (\cref{ex:scaling:sightings}).

\paragraph{Where it breaks.} \emph{The exponents are not derived.} The models say how a decay of importance becomes a decay of loss, not what the decay is. For kernels it can be measured: \textcite{SpiglerEtAl2020} estimated learning-curve exponents from the kernel spectra of MNIST and CIFAR-10 that came close to the measured ones. For trained networks there is no such recipe, and feature learning changes the exponents: in the solvable model of \textcite{BordelonEtAl2025}, feature learning nearly doubles the exponent in time and compute for targets outside the reproducing-kernel space of the initial tangent kernel, and leaves it unchanged for easier targets.

\emph{The ratio of the exponents.} Both models predict $\alpha_D=\alpha_N/(1+\alpha_N)<\alpha_N$. With $\alpha_N=0.076$, this gives $0.071$, against the $\alpha_D\approx0.095$ of \textcite{KaplanEtAl2020}; with $\alpha_N=0.34$ it gives $0.25$, against the $0.28$ of \textcite{HoffmannEtAl2022}, and the refit of \textcite{BesirogluEtAl2024} has a data exponent that is larger, not smaller. The comparison is not clean, since the exponents of \textcite{KaplanEtAl2020} were measured in different regimes (convergence in $N$, early stopping in $D$), but the prediction fails in the direction that matters: data are worth more, relative to parameters, than the hierarchy alone allows. A learner that improves on many items from each example, rather than on one, would do this; memorization cannot.

\emph{The fits.} The exponents of \eqref{eq:scaling:chinchilla} are fitted, and the fits are fragile: the published estimates of \textcite{HoffmannEtAl2022} did not survive refitting \parencite{BesirogluEtAl2024}, and the allocations of \textcite{KaplanEtAl2020} and \textcite{HoffmannEtAl2022} disagree because of parameter counting, scale, warmup, and tuning \parencite{PearceSong2024,PorianEtAl2024}. \Cref{thm:scaling:compute} is exact for the additive form, but the solvable models are not additive, and the tokens per parameter depend on constants that the models do not predict.

\emph{Emergence.} The quanta model makes each skill turn on over a factor of about $22$ in data. Abilities that appear abruptly \parencite{WeiEtAl2022} may be sharper than this, or may look sharp because of the metric \parencite{SchaefferEtAl2023}; the model cannot tell which, since it has no notion of a task built from several quanta.

\emph{The data distribution.} The law in $D$ is a law of random sampling. If the training set contains each item once instead of sampling with repetition, the memorizing learner's error after $D$ examples is $\sum_{k>D}p_k\asymp D^{-b}$: curation raises the data exponent from $b/(1+b)$ to $b$ (\cref{ex:scaling:pruning}). \Textcite{SorscherEtAl2022} showed in theory and experiment that pruning data can beat power-law scaling, and \textcite{MuennighoffEtAl2023} that repeated tokens are worth less than fresh ones after a few epochs; neither effect exists in a model of fresh, independent samples. And real curves bend: a spectrum that is a power law only over a range gives a broken power law \parencite{CaballeroEtAl2023}, with plateaus where the range ends \parencite{MaloneyEtAl2022}.

\paragraph{Fine structure.} \Textcite{MaloneyEtAl2022} solved random-feature regression on data with a power-law latent spectrum: nonlinear random features extend the power law, the test loss scales in both $N$ and $D$, and it plateaus when either exceeds the size of the latent spectrum. \Textcite{BahriEtAl2024} separate variance-limited scaling, with exponent $1$, from resolution-limited scaling with data-dependent exponents, in both $N$ and $D$. Dynamics: \textcite{BordelonEtAl2024} solved a random-feature model trained by gradient descent, in which time and model size have different exponents and the gap between training and test loss builds up through reuse of data; \textcite{PaquetteEtAl2024} found four phases and three subphases of compute-optimal scaling for one-pass stochastic gradient descent in the plane of data and target complexity, with boundaries set by model capacity, optimizer noise, and the embedding of the features; and \textcite{LinEtAl2024} proved that for sketched linear regression trained by one-pass stochastic gradient descent the reducible loss is $\Theta(N^{-(a-1)}+D^{-(a-1)/a})$ for a Gaussian prior with $\lambda_j\asymp j^{-a}$, our $b=a-1$, with no variance term growing with $N$: implicit regularization removes it. \Textcite{SharmaKaplan2022} derive exponents of about $4/d$ from the intrinsic dimension $d$ of the data manifold, and \textcite{CabannesEtAl2024} treat associative memories. Open: a model in which the exponent $b$ of the learned representation is an output, computed from the data and the architecture, and checked against language models.

\section{Model card}\label{sec:scaling:card}

\modelcard{Test loss falls as a power law in parameters, data, and compute, with small exponents; compute-optimal training scales parameters and data together.}
{Least squares and ridge on independent Gaussian features with target energy $j^{-(1+b)}$; a memorizing learner on quanta with Zipf frequencies $p_k\propto k^{-(1+b)}$; an additive loss with $C=6ND$.}
{Exact learning curve $(\sigma^2+T(N))\frac{D-1}{D-N-1}$ (\cref{thm:scaling:exact}); exponents $b$ and $b/(1+b)$ (\cref{cor:scaling:laws,prop:scaling:ridge}); Zipf curve with constant $C_b$ (\cref{thm:scaling:zipf}); allocation (\cref{thm:scaling:compute}).}
{Optimal size $N\propto D^{1/(1+b)}$; data exponent $b/(1+b)$ with noise, $b$ without; time exponent $b/a$; allocation exponents $\frac{\alpha_D}{\alpha_N+\alpha_D}$ and $\frac{\alpha_N}{\alpha_N+\alpha_D}$; skills emerging over a factor $22$ in data.}
{Exponents are inputs; the predicted $\alpha_D<\alpha_N$ fails for fitted language-model exponents; no feature learning, optimizer, repetition, or curation.}
{Random-feature models; dynamical models and the 4+3 phases; one-pass stochastic gradient descent; feature learning that changes exponents.}

This chapter closes the thread of spectra that began with the Marchenko--Pastur law of \cref{ch:doubledescent}: \cref{thm:scaling:exact} is the weak-features corollary of that chapter with an anisotropic design, and the deterministic equivalent is its ridge theorem with a renormalized penalty. The spectral bias of kernel gradient descent from \cref{ch:lazy}, which learns eigenfunctions in order of their eigenvalues, becomes here a scaling law in time with exponent $b/a$, and the staged learning of \cref{ch:deeplinear}, one singular mode after another, is the mechanism that the model of size $N$ assumes.

The quanta connect smooth loss curves with the sudden transitions of earlier chapters: the plateaus of the committee machine (\cref{ch:committee}), hidden progress on parities (\cref{ch:parity}), grokking (\cref{ch:grokking}), and the formation of induction heads (\cref{ch:induction}) are each one quantum, abrupt on its own scale, and a power law is what many such transitions look like when added up with Zipf weights. The capacity to hold more items than dimensions is the subject of \cref{ch:superposition}.

\begin{exercises}
\item \consolidation (Tails and constants.) (a) Show that $\sum_{k>N}p_k\sim N^{-b}/(b\zeta(1+b))$ and that $T(N)=N^{-b}/b+O(N^{-b-1})$ in the power-law case. (b) Show that $C_1=\sqrt6/(2\sqrt\pi)$, and compute $C_{1/2}$ and $C_2$ to three digits. (c) Show that $\sum_{k\ge1}p_ke^{-cp_k}$ is within $2/(ec)$ of $C_bc^{-b/(1+b)}$ for every $c>0$, and deduce that $D^{b/(1+b)}M(D)-C_b=O(D^{-1/(1+b)})$.\label{ex:scaling:tails}

\item \consolidation (Noiseless learning curves.) Under \cref{ass:scaling:features} with $\sigma=0$ and $T(N)\sim cN^{-b}$, show that $\min_{N\le D-2}\E\Loss(\hat w_N)\sim c\,\frac{(1+b)^{1+b}}{b^b}D^{-b}$, attained at $N\sim\frac b{1+b}D$. Check numerically that for $b=1$ and $T(N)=\sum_{j>N}j^{-2}$ the product $D\cdot\min_N\E\Loss(\hat w_N)$ is $4.000$ at $D=1000$. Why does the optimal model grow linearly with $D$ here?\label{ex:scaling:noiseless}

\item \consolidation (Chinchilla arithmetic.) (a) Derive \eqref{eq:scaling:alloc} with a Lagrange multiplier instead of the substitution $u=\log N$. (b) For $\alpha_N=0.34$ and $\alpha_D=0.28$, compute $\nu_N$, $\nu_D$, the compute exponent of the loss, and the exponent of $D^\star/N^\star$. (c) Show that at the optimum the parameter term carries the fraction $\nu_N$ of the reducible loss, and that a mistaken allocation $N=\rho N^\star$, $D=D^\star/\rho$ raises the reducible loss by the factor $\nu_N\rho^{-\alpha_N}+\nu_D\rho^{\alpha_D}$. Evaluate it for $\rho=4$.\label{ex:scaling:chinchilla}

\item \consolidation (The isotropic case.) Take $\lambda_j=1$ for $j\le P$ and $P/D\to\gamma$. Show that the equation for $\kappa$ in the deterministic equivalent becomes $\kappa=\lambda+\gamma\kappa/(1+\kappa)$, and that with $v=\lambda/\kappa$ it is the equation $\gamma v^2+(1-\gamma+\lambda)v-\lambda=0$ of \cref{ch:doubledescent}. Check that \eqref{eq:scaling:de} then gives the bias $r^2(v-\lambda v')$ and the variance $\sigma^2\gamma v'$ of \cref{thm:doubledescent:ridge}. \hint{differentiate the quadratic equation in $\lambda$.}\label{ex:scaling:isotropic}

\item \consolidation (Exact constants for ridge.) In \cref{prop:scaling:ridge}, show that as $\lambda\to0$ the bias is asymptotic to $\lambda^{b/a}\frac1aB\bigl(\frac{2a-b}a,\frac ba\bigr)$ and the variance sum to $\lambda^{-1/a}\frac1aB\bigl(\frac1a,2-\frac1a\bigr)$, with $B$ the Beta function, and compute the optimal $\lambda$ and the constant in $\min_\lambda\cE_{\mathrm{seq}}\sim c'(\sigma^2/D)^{b/(1+b)}$. \hint{$\int_0^\infty\frac{u^{s-1}}{(1+u^a)^2}\dd u=\frac1aB\bigl(\frac sa,2-\frac sa\bigr)$; use \cref{lem:scaling:unimodal}-type comparisons on the two monotone parts of each summand.}\label{ex:scaling:beta}

\item \variation (Training time.) Gradient flow on the population loss of the model of size $N$, from $w=0$, gives $w_j(t)=(1-e^{-\lambda_jt})w_{\star,j}$, so the excess loss is $\sum_{j\le N}j^{-(1+b)}e^{-2\lambda_jt}+T(N)$ in the power-law case. (a) Show that for $t\ge1$ and $N\to\infty$ the first sum is of order $t^{-b/a}$, and that the excess loss is of order $t^{-b/a}+N^{-b}$ when $N\ge t^{1/a}$ and of order $N^{-b}$ when $N\le t^{1/a}$. (b) If one unit of compute buys one unit of time for one feature, $C=Nt$, show that the compute-optimal size is $N\asymp C^{1/(1+a)}$ and the loss is of order $C^{-b/(1+a)}$. Compare with \textcite{BordelonEtAl2024} and \textcite{PaquetteEtAl2024}, where the model is also trained with finite data.\label{ex:scaling:time}

\item \variation (Saturation.) In \cref{prop:scaling:ridge}, let $b>2a$. Show that the bias is of order $\lambda^2$, that the best excess loss is of order $(\sigma^2/D)^{2a/(1+2a)}$, and that least squares on the top $N$ features, with $N$ chosen well, does better. What happens at $b=2a$?\label{ex:scaling:saturation}

\item \variation (Several sightings.) Suppose an item is learned only after it has been seen $m$ times. Show that the error is $\sum_kp_k\Pr\bigl(\mathrm{Bin}(D,p_k)<m\bigr)\sim C_{b,m}D^{-b/(1+b)}$ with
\[
  C_{b,m}=\frac1{(1+b)\,\zeta(1+b)^{1/(1+b)}}\sum_{i=0}^{m-1}\frac{\Gamma\bigl(i+\frac b{1+b}\bigr)}{i!} .
\]
For $b=1$, check numerically that $C_{1,2}\approx1.036$ and $C_{1,3}\approx1.296$. The exponent does not change; what does?\label{ex:scaling:sightings}

\item \variation (Capacity and data together.) A learner of capacity $N$ trained on $D$ examples holds the items it has seen among $1,\dots,N$. (a) Show that its error is $\sum_{k>N}p_k+\sum_{k\le N}p_k(1-p_k)^D$. (b) Show that this is of order $N^{-b}+D^{-b/(1+b)}$ uniformly in $N,D\ge1$, and that the crossover lies at $N\asymp D^{1/(1+b)}$. (c) With $C=6ND$, find the compute-optimal allocation and compare it with \cref{thm:scaling:compute} for $\alpha_N=b$ and $\alpha_D=b/(1+b)$.\label{ex:scaling:combined}

\item \variation (Curated data.) (a) In \cref{ass:scaling:quanta}, a curator shows the learner the items $1,\dots,D$ once each. Show that the error is $\sum_{k>D}p_k\sim D^{-b}/(b\zeta(1+b))$. (b) Now the curator removes repeated items from a sample of size $D'$ and keeps the first $D$ distinct ones. Express the error in terms of $D'$, and show that the exponent in $D'$ is still $b/(1+b)$ while the exponent in the number $D$ of kept examples is $b$. Relate this to \textcite{SorscherEtAl2022}.\label{ex:scaling:pruning}

\item \variation (Compute for least squares.) Minimize $cN^{-b}+\sigma^2N/D$ under $6ND=C$, and show that $N^\star\propto C^{1/(2+b)}$, $D^\star\propto C^{(1+b)/(2+b)}$, and that the loss falls like $C^{-b/(2+b)}$. Why does the variance term push the allocation toward data?\label{ex:scaling:olscompute}

\item \simulation (Ridge exponents.) With $a=\frac32$, $\sigma^2=0.2$, and $P=10^6$ features, compute the deterministic equivalent of optimally tuned ridge at $D=10^3$ and $D=10^4$ for $b\in\{\frac12,1,4\}$, and the slope $-\log_{10}$ of their ratio. Observe slopes of about $0.39$, $0.52$, and $0.74$, against $\frac13$, $\frac12$, and the saturated $\frac34$. Explain why the slope for $b=\frac12$ is too steep, using the effective noise $\sigma^2+R$ in \eqref{eq:scaling:de}, and check that the sequence model \eqref{eq:scaling:seq}, whose noise does not grow with $R$, gives $0.335$. Confirm one point by simulating ridge at $D=10^3$ with $4000$ features.\label{ex:scaling:simridge}

\item \simulation (Fitting a law.) For $b=1$, compute the local slope $-\dd\log M/\dd\log D$ of the exact Zipf curve at $D=3$, $10$, $100$, and $10^4$; observe about $0.48$, $0.44$, $0.50$, and $0.50$. Then draw $D$ samples for $D=10,20,\dots,320$, record the error of the memorizing learner, and fit $\log M$ against $\log D$ by least squares on $20$ independent repetitions. How far is the fitted exponent from $\frac12$, and how wide is the spread? Relate the result to the replication of \textcite{BesirogluEtAl2024}.\label{ex:scaling:fit}

\item \research (Predicting an exponent.) \Textcite{SpiglerEtAl2020} predicted kernel learning-curve exponents from kernel spectra of image data. Open: for a network that learns its features, define and measure the analogue of $b$ in its learned representation, and predict its parameter and data exponents; a solvable model in which feature learning changes $b$ during training, as in \textcite{BordelonEtAl2025}, and which yields $\alpha_D>\alpha_N$, would be a first step.\label{ex:scaling:research}
\end{exercises}

\notesandrefs
Learning curves are older than scaling laws. In the statistical mechanics of learning, \textcite{AmariEtAl1992} classified the asymptotic learning curves of parametric learning machines into four types, among them $D^{-1}$ for a deterministic machine and a noiseless teacher and $D^{-1/2}$ for a deterministic machine and a noisy teacher; such universal exponents are set by the number of parameters of a model that contains the target. The exponents of this chapter, which depend on the data, appear when the target is not representable by a finite model. In statistics, \textcite{CaponnettoDeVito2007} proved the minimax rates of kernel ridge regression under a source condition, the decay of the target in the eigenbasis, and a capacity condition, the decay of the eigenvalues; \cref{prop:scaling:ridge} is the sequence-model version of their rate. \Textcite{BordelonEtAl2020} and \textcite{SpiglerEtAl2020} brought these rates to wide networks and real data, and \textcite{CuiEtAl2021} mapped the crossover from the noiseless rate $D^{-b}$ to the noisy one.

\Textcite{HestnessEtAl2017} established power laws in data across domains for deep networks, \textcite{KaplanEtAl2020} in parameters, data, and compute for language models, and \textcite{HoffmannEtAl2022} the compute-optimal allocation that is used today, after correcting the earlier one. \Cref{thm:scaling:compute} is their allocation formula for the parametric fit, which they write as $N_{\mathrm{opt}}\propto C^a$ and $D_{\mathrm{opt}}\propto C^b$ with $a=\beta/(\alpha+\beta)$ and $b=\alpha/(\alpha+\beta)$. \Textcite{BesirogluEtAl2024} is an unrefereed preprint, and our account of it rests on its abstract.

Solvable models of the exponents came in several waves: memorization of Zipf-distributed data \parencite{Hutter2021} and quanta \parencite{MichaudEtAl2023}; random features \parencite{MaloneyEtAl2022,BahriEtAl2024}; dynamics \parencite{BordelonEtAl2024,PaquetteEtAl2024,LinEtAl2024}; and feature learning \parencite{BordelonEtAl2025}. The exact finite-sample curve of \cref{thm:scaling:exact} is an elementary consequence of the inverse Wishart mean; we did not find it stated in this form in the scaling literature.

\paragraph{Translating the literature.} \Textcite{HoffmannEtAl2022} write $E$, $\alpha$, and $\beta$ for our $\Loss_\infty$, $\alpha_N$, and $\alpha_D$, and $a$, $b$ for our allocation exponents $\nu_N$, $\nu_D$, which are not our spectral exponents. \Textcite{KaplanEtAl2020} use our $\alpha_N$ and $\alpha_D$, but without an irreducible term, and count non-embedding parameters. \Textcite{LinEtAl2024} write $M$ for our $N$ and $N$ for our $D$; their spectral exponent is our $a$, and their Gaussian prior on the target makes our $b=a-1$, so their rate $M^{-(a-1)}+N^{-(a-1)/a}$ has our two exponents $b$ and $b/(1+b)$. \Textcite{Hutter2021} writes $\theta_i\propto i^{-(1+\alpha)}$ and $n$ for our $p_k\propto k^{-(1+b)}$ and $D$. In the kernel literature \parencite{CaponnettoDeVito2007,CuiEtAl2021} the capacity and source conditions are written with other exponents; with $\lambda_j\asymp j^{-a}$ and $\lambda_jw_{\star,j}^2\asymp j^{-(1+b)}$, the noisy rate is always $D^{-b/(1+b)}$ below saturation, and translating a paper's exponents into $a$ and $b$ is the quickest way to compare results.

\printchapterbib
\end{refsection}


\bookmarksetup{startatroot}
% Epilogue of The Hydrogen Atoms of Deep Learning.
\specialchapter{Epilogue}{Fine structure: what the solvable models get wrong together}{100}\label{ch:epilogue}
\begin{refsection}[chapters/17-epilogue.bib]

Sixteen models, each solved in its own chapter, each with its own list of failures. Read one at a time, they look like sixteen separate tricks. Read side by side, they share a small number of mechanisms, and they also share their failures. A failure that one model has is a fact about that model. A failure that every model has is a fact about the method, and it is the fine structure that the next generation of models has to explain. This epilogue collects both, lists phenomena that still have no hydrogen atom, and ends with a recipe for building one.

\addsec{Mechanisms that recur}

The table below lists the mechanisms that carry more than one chapter, the chapters in which they do the main work, and what each one computes. None of them is new to this book; most are older than deep learning. What is striking is how few there are.

\begin{fullwidth}
\small
\renewcommand{\arraystretch}{1.22}
\begin{tabularx}{\linewidth}{@{}>{\raggedright}p{0.2\linewidth}>{\raggedright}p{0.14\linewidth}X@{}}
\toprule
\textsc{mechanism} & \textsc{chapters} & \textsc{what it computes}\\
\midrule
Escape from a saddle & 2, 3, 7, 8, 9, 10, 13, 14 & The length of a plateau: the logarithm of the initial distance from the saddle, divided by the unstable rate\\
Conservation laws from rescaling symmetries & 2, 3, 13, 15 & Which minimizer training reaches, how sharp it is, and which initializations make the dynamics closed\\
Reduction to order parameters & 1, 5, 7, 8, 13 & A closed recursion or ordinary differential equation in a few overlaps, from rotational invariance of Gaussian data\\
Harmonic analysis & 5, 6, 7, 9, 10 & What the gradient sees at each degree: Hermite coefficients, spherical harmonics, characters of $\Z_p$ and of the cube\\
Margins as attractors & 4, 10, 12 & The direction that an exponentially tailed loss selects after the training error reaches zero\\
Spectra of random and structured matrices & 1, 6, 13, 16 & Risk as a functional of a spectrum: Marchenko--Pastur, kernel eigenvalues, power-law tails\\
A limit that closes the equations & 1, 5, 6, 7, 8, 16 & Width, dimension or sample size taken to infinity in a fixed ratio or a fixed order\\
\bottomrule
\end{tabularx}
\end{fullwidth}

The first row deserves a comment, because it is the most reliable prediction in the book. Near a saddle the dynamics is linear, and an unstable mode with rate $\lambda$ grows like $e^{\lambda t}$. Starting at distance $\eps$, it needs a time $\lambda^{-1}\log(1/\eps)$ to reach order one. This is the normal-mode theorem of \hyperref[ch:prologue]{the prologue} with the sign of one eigenvalue reversed. It gives the half-learning time $\log(s/a_0)/(2s)$ of a two-layer linear network (\cref{ch:deeplinear}), the time $\log(1/\alpha)/(2\abs{w_{\star,i}})$ at which a diagonal network learns coordinate $i$ (\cref{ch:diagonal}), the $\log d$ plateau of the committee machine (\cref{ch:committee}), the delay $\lambda^{-1}\log(\alpha/\alpha_c)$ of grokking under weight decay (\cref{ch:grokking}), and the escape times of linear attention from small initialization (\cref{ch:icl,ch:induction}). The single-index model and the sparse parity (\cref{ch:singleindex,ch:parity}) are the same story at a degenerate saddle, where the unstable direction grows like a power rather than an exponential, and the logarithm becomes a power of the dimension. Whenever training shows a plateau followed by a drop, this is the first model to try, and the two numbers to measure are the initial distance and the unstable rate.

The second row is less visible but as useful. A rescaling symmetry of the parameterization, such as $(W_{\ell+1},W_\ell)\mapsto(W_{\ell+1}A^{-1},AW_\ell)$ for a linear network or $(u,v)\mapsto(cu,v/c)$ for a product, gives a conserved quantity under gradient flow: the balancedness of \cref{ch:deeplinear}, the products $u_iv_i=\alpha^2$ of \cref{ch:diagonal}. Under gradient descent with a finite step the same quantity is no longer conserved but evolves by an exact law, and in \cref{ch:stability} that law is what forces the catapult. In every case the conserved or slowly varying quantity selects among minimizers that the loss cannot tell apart.

The third and fourth rows are the reason the models are solvable at all. Gaussian data are rotationally invariant, so a high-dimensional state collapses to a few overlaps; uniform data on a group are invariant under translation, so the right basis is the group's characters. The fifth row is the mechanism of implicit bias in classification: once the data are separated, an exponentially tailed loss is dominated by its smallest margins, and the trajectory follows the maximum-margin problem. Three chapters of very different appearance, on logistic regression, modular addition and the last layer of a classifier, are three solutions of that problem.

\addsec{What the models get wrong together}

Each chapter ends with what its model gets wrong. Putting the lists side by side, six failures recur.

\paragraph{The data are isotropic, independent, or uniform.} Most of the sixteen models assume Gaussian inputs with identity covariance, uniform inputs on a group or a cube, or independent features. This is what makes the order parameters close. Real data have structure at every scale, and in the models that structure enters, if at all, only through a spectrum (\cref{ch:doubledescent,ch:scaling}) or through a target that is given rather than learned. Where a chapter tests what happens without isotropy, the prediction changes: anisotropic inputs change the information exponent (\cref{ch:singleindex}), and structured inputs change or remove the plateau of the committee machine (\cref{ch:committee}).

\paragraph{One mechanism at a time.} Every model isolates one phenomenon and removes everything that could compete with it. In a trained network the mechanisms run together. The delay of grokking in \cref{ch:grokking} had to be modelled by a different network from the one that explains the Fourier circuit; the induction head of \cref{ch:induction} is one circuit among many that form during the same phase change; implicit bias, edge-of-stability dynamics and feature learning act on the same weights at the same time. No model in this book solves an interaction between two of its own mechanisms.

\paragraph{Limits in a convenient order.} Solvability comes from a limit, and the limit is usually taken in the order that makes the equations close: width to infinity at fixed depth (\cref{ch:signalprop,ch:lazy}), dimension to infinity before time (\cref{ch:singleindex,ch:committee}), time to infinity for a margin (\cref{ch:margin}). The interesting behavior often lives where two limits compete. Signal propagation fails when depth is comparable to width; the committee machine's plateau length is set by fluctuations of order $d^{-1/2}$ that its own ordinary differential equation omits; the margin converges at a rate in $1/\log t$ that no training run reaches.

\paragraph{Gradient flow, or plain gradient descent.} Almost every theorem is proved for gradient flow, online stochastic gradient descent with vanishing steps, or full-batch gradient descent with a small step. Only \cref{ch:stability} takes a large step seriously, and no chapter treats an adaptive optimizer. Yet the chapters themselves record that the optimizer changes the answer: step size and label noise change the implicit bias of diagonal networks (\cref{ch:diagonal}), steepest descent and Adam change the margin that classification selects (\cref{ch:margin}), and reusing samples beats the information exponent (\cref{ch:singleindex}).

\paragraph{Special initializations.} Several solutions require an initialization chosen to make the dynamics solvable: aligned and balanced for deep linear networks, symmetric for the committee machine and the sparse parity, small for diagonal networks and attention. The chapters show that random initializations keep the scaling and change the constants, but the exact solutions are exact only on a set of measure zero.

\paragraph{Minimal architectures.} The models are linear, diagonal, two-layer, or one layer of attention. Residual connections, normalization layers and softmax attention appear only in the sections on what the models get wrong. The transformer simulation of \cref{ch:induction} did not reach the bigram level with a standard linear readout and needed a modified one.

\medskip

The book's own simulations are a fair test of these failures, because every chapter compares its model with a finite simulation, and several of them disagreed with the model in an instructive way. A small ReLU network trained in \cref{ch:collapse} collapsed only partially. The networks of \cref{ch:grokking} are far from the maximum-margin solutions that explain their Fourier features. The fitted exponents of language models violate the relation between data and parameter exponents that both models of \cref{ch:scaling} predict. And the transformers that match least squares on regression prompts cannot be performing the single gradient step that \cref{ch:icl} proves optimal for one layer.

The pattern across these comparisons is consistent. What the models predict well are scalings and orders of events: a time that grows like $\log(1/\alpha)$ or a power of $d$, a threshold at $\eta\Lambda=2$, the order in which modes are learned, the geometry of an optimum. What they predict badly are constants, and which mechanism is active in a given network. This is the same division as in physics, where the hydrogen atom gets the structure of the spectrum right and the fine structure wrong, and where the corrections are small. In deep learning the corrections are not always small.

\addsec{Missing hydrogen atoms}

Several phenomena that matter in practice have no model in this book, either because none exists or because the existing ones do not yet pass the four tests of \hyperref[ch:prologue]{the prologue}.

\emph{Adaptive optimizers.} Almost every large network is trained with Adam \parencite{KingmaBa2015} or a relative of it. Sign descent and steepest descent in other norms have partial theories (\cref{ch:margin}), but there is no solvable model that predicts when Adam trains a network that gradient descent cannot, or the step-size schedules used in practice.

\emph{Progressive sharpening.} The model of \cref{ch:stability} explains why gradient descent at a large step does not diverge, but not why the sharpness rises to $2/\eta$ in the first place, which \textcite{CohenEtAl2021} observed in many networks. A model that produces sustained sharpening from its own dynamics is missing.

\emph{Normalization layers.} Batch normalization \parencite{IoffeSzegedy2015} and layer normalization \parencite{BaKirosHinton2016} change the geometry of the loss, make it invariant to the scale of the weights, and interact with weight decay and the step size. They appear in this book only as things the models leave out.

\emph{Data mixtures and curricula.} Real training sets are mixtures of sources with different difficulty and redundancy, and the order in which examples arrive can matter \parencite{BengioEtAl2009}. The models here draw every example from one fixed distribution.

\emph{Learning from feedback.} Reinforcement learning from human preferences \parencite{ChristianoEtAl2017,OuyangEtAl2022} trains a pretrained model on a reward that is itself learned. Its dynamics, including how far the model moves from its pretrained state, have no solvable model of the kind collected here.

\emph{Chain of thought.} Prompting a model to write intermediate steps improves its answers on reasoning tasks \parencite{WeiEtAl2022}. The sparse parity of \cref{ch:parity} suggests why intermediate steps can help a learner, since they break a high-degree target into low-degree pieces, but a model that predicts the size of the effect in a trained transformer does not exist yet.

\emph{Long-context retrieval.} Language models use information at the start and end of a long input better than information in the middle \parencite{LiuEtAl2024}. The models of \cref{ch:icl,ch:induction} treat every position in the prompt alike.

\addsec{A recipe for a new hydrogen atom}

The chapters suggest a procedure, which a reader can apply to each of the phenomena above.

\begin{enumerate}
\item \emph{Choose a phenomenon with a number.} An exponent, a time scale, a threshold or a geometry, observed robustly in several networks. Without a number the model cannot fail, and a model that cannot fail cannot teach.
\item \emph{Find the symmetry.} Look for the invariance that reduces the state to a few variables: rotations of Gaussian inputs, translations on a group, a rescaling of the parameters. If no symmetry is visible, build it into the data.
\item \emph{Choose the limit and the initialization.} Take the limit (dimension, width, time, initialization scale) and choose the starting point that makes the reduced equations close. Write down what this choice throws away; it will be the first section of ``What the model gets wrong''.
\item \emph{Solve the smallest instance completely.} Find the saddle, the conserved quantity or the margin problem. Prove the main theorem, and label as heuristic whatever you cannot prove.
\item \emph{Predict a number and simulate with no fitted constant.} Compare the prediction with a finite simulation of the model, and then with a network that the model does not describe. The disagreement with the second is the result.
\item \emph{Break one assumption on purpose.} Change the data, the optimizer, the architecture or the limit, one at a time, and solve again. This is what the variation exercises of every chapter ask for, and it is how one model becomes the next.
\item \emph{Write the model card.} If the six lines cannot be filled honestly, in particular the line on what the model gets wrong, the model is not ready.
\end{enumerate}

The hydrogen atom became the foundation of atomic physics not because it was right, but because every way in which it was wrong pointed to something real: relativity, the electron's spin, the quantized field. The models in this book are wrong in ways that are known and specific. The hope is that their failures, studied as carefully as their solutions, point to what deep learning theory needs next.

\notesandrefs

The table of mechanisms condenses material from the chapters, where each mechanism is proved and referenced. The observation that the time to escape a saddle scales like $\log(1/\eps)/\lambda$ is classical in the theory of dynamical systems, and its role in the chapters is discussed in the notes to \cref{ch:deeplinear,ch:diagonal,ch:committee}. Adam is due to \textcite{KingmaBa2015}; batch and layer normalization to \textcite{IoffeSzegedy2015} and \textcite{BaKirosHinton2016}; curriculum learning was named by \textcite{BengioEtAl2009}; learning from human preferences is due to \textcite{ChristianoEtAl2017} and was applied to language models by \textcite{OuyangEtAl2022}; chain-of-thought prompting is due to \textcite{WeiEtAl2022}; the position dependence of long-context retrieval was measured by \textcite{LiuEtAl2024}; progressive sharpening was named by \textcite{CohenEtAl2021}.

\printchapterbib
\end{refsection}


\appendix
% Appendix A of The Hydrogen Atoms of Deep Learning: the shared mathematical toolkit.
\casechapter{A Toolkit}{The mathematics the models share}{}{ch:toolkit}
\begin{refsection}[chapters/18-toolkit.bib]

\begin{chaptersummary}
Five pieces of mathematics recur across the models of this book: Gaussian integration by parts and Hermite polynomials, which turn expectations over Gaussian data into power series in an overlap; spherical harmonics, which do the same on the sphere; the Marchenko--Pastur law for the spectrum of a sample covariance matrix; martingale inequalities and the method of ordinary differential equations, which reduce a stochastic recursion in high dimension to a deterministic flow; and Fourier analysis on finite groups. This appendix states what the chapters use and proves what can be proved in a page.
\end{chaptersummary}

The appendix is a reference, not a chapter to be read in order. Hermite polynomials (\cref{sec:toolkit:hermite}) are used in \cref{ch:signalprop,ch:singleindex,ch:committee,ch:lazy}; spherical harmonics (\cref{sec:toolkit:sphere}) in \cref{ch:lazy}; the Marchenko--Pastur law (\cref{sec:toolkit:mp}) in \cref{ch:doubledescent,ch:scaling}; martingales and the ODE method (\cref{sec:toolkit:ode}) in \cref{ch:singleindex,ch:committee}; Fourier analysis (\cref{sec:toolkit:fourier}) in \cref{ch:parity,ch:grokking}; and concentration inequalities (\cref{sec:toolkit:conc}) throughout.

\section{Gaussian integration by parts and Hermite polynomials}\label{sec:toolkit:hermite}

Throughout this section $z$ is a standard Gaussian random variable, with density $\varphi(z)=(2\pi)^{-1/2}e^{-z^2/2}$, and $L^2(\gamma)$ is the space of functions $f\colon\R\to\R$ with $\E[f(z)^2]<\infty$, with inner product $\ip{f,g}=\E[f(z)g(z)]$.

\begin{theorem}[Gaussian integration by parts]\label{thm:toolkit:stein}\index{Gaussian integration by parts}\index{Stein's lemma}
Let $f\colon\R\to\R$ be absolutely continuous with $\E\abs{f'(z)}<\infty$. Then $\E\abs{zf(z)}<\infty$ and
\[
  \E[zf(z)]=\E[f'(z)].
\]
More generally, let $x\sim\cN(0,\Sigma)$ in $\R^d$ and let $f\colon\R^d\to\R$ be continuously differentiable with $\E\abs{f(x)}$ and $\E\norm{\nabla f(x)}$ finite. Then $\E[x f(x)]=\Sigma\,\E[\nabla f(x)]$.
\end{theorem}

\begin{proof}
The plan is to write $f'$ as an integral against the Gaussian tail and exchange the order of integration. The density satisfies $\varphi'(u)=-u\varphi(u)$, so $\varphi(t)=\int_t^\infty u\varphi(u)\dd u$ for $t\ge0$ and $\varphi(t)=\int_{-\infty}^t(-u)\varphi(u)\dd u$ for $t\le0$. On the positive half-line we substitute the first formula and exchange the order of integration, which Fubini's theorem allows because the same computation with $\abs{f'}$ in place of $f'$ gives the finite value $\int_0^\infty\abs{f'}\varphi$:
\begin{align*}
  \int_0^\infty f'(t)\varphi(t)\dd t&=\int_0^\infty f'(t)\int_t^\infty u\varphi(u)\dd u\dd t\\
  &=\int_0^\infty u\varphi(u)\big(f(u)-f(0)\big)\dd u .
\end{align*} The negative half-line gives in the same way $\int_{-\infty}^0f'\varphi=\int_{-\infty}^0u\varphi(u)(f(u)-f(0))\dd u$. Adding the two, $\E[f'(z)]=\E[z(f(z)-f(0))]=\E[zf(z)]$, since $\E z=0$; the integrability of $zf(z)$ follows from the bound $\E\abs{z}\abs{f(z)-f(0)}\le\E\abs{f'(z)}$, which the same computation with absolute values provides.

For the vector form, write $x=\Sigma^{1/2}g$ with $g\sim\cN(0,I_d)$ and $F(g)=f(\Sigma^{1/2}g)$, so that $\nabla F(g)=\Sigma^{1/2}\nabla f(x)$. Conditionally on the other coordinates, the scalar statement applied to the $i$th coordinate of $g$ gives $\E[g_iF(g)]=\E[\partial_iF(g)]$, by Fubini's theorem and the integrability assumptions. Hence $\E[xf(x)]=\Sigma^{1/2}\E[gF(g)]=\Sigma^{1/2}\E[\nabla F(g)]=\Sigma\,\E[\nabla f(x)]$.
\end{proof}

Hermite polynomials are the orthogonal polynomials of the Gaussian measure. The quickest definition is through their generating function: for real $s$ and $z$,
\begin{equation}\label{eq:toolkit:genfun}
  G(s,z)\defeq e^{sz-s^2/2}=\sum_{k\ge0}\He_k(z)\,\frac{s^k}{k!}.
\end{equation}
Expanding $e^{sz}e^{-s^2/2}$ as a product of two power series shows that $\He_k$ is a monic polynomial of degree $k$: $\He_0=1$, $\He_1(z)=z$, $\He_2(z)=z^2-1$, $\He_3(z)=z^3-3z$, $\He_4(z)=z^4-6z^2+3$. The normalized Hermite polynomials are $h_k=\He_k/\sqrt{k!}$, and the \term{Hermite coefficients}\index{Hermite coefficient} of $f\in L^2(\gamma)$ are $\hat f_k=\E[f(z)h_k(z)]$.

\begin{lemma}[Hermite polynomials]\label{lem:toolkit:hermite}\index{Hermite polynomials}
\begin{enumerate}
\item $\He_k'=k\He_{k-1}$ and $\He_{k+1}(z)=z\He_k(z)-k\He_{k-1}(z)$; equivalently $h_k'=\sqrt k\,h_{k-1}$ and $zh_k=\sqrt{k+1}\,h_{k+1}+\sqrt k\,h_{k-1}$.
\item $\E[\He_j(z)\He_k(z)]=k!\,\delta_{jk}$, so the $h_k$ are orthonormal in $L^2(\gamma)$; they form an orthonormal basis.
\item If $f$ has $k$ derivatives with $\E\abs{f^{(i)}(z)}\abs{\He_{k-i}(z)}<\infty$ for $i\le k$, then $\E[f(z)\He_k(z)]=\E[f^{(k)}(z)]$, so $\hat f_k=\E[f^{(k)}(z)]/\sqrt{k!}$.
\end{enumerate}
\end{lemma}

\begin{proof}
For (i), differentiate the generating function: $\partial_zG=sG$ and $\partial_sG=(z-s)G$; comparing the coefficients of $s^k/k!$ gives $\He_k'=k\He_{k-1}$ and $\He_{k+1}=z\He_k-k\He_{k-1}$. Dividing by the factorials gives the normalized forms. For (ii), compute with the moment generating function of $z$:
\[
  \E[G(s,z)G(t,z)]=e^{-(s^2+t^2)/2}\,\E[e^{(s+t)z}]=e^{st}=\sum_{k\ge0}\frac{(st)^k}{k!}.
\]
The series in \eqref{eq:toolkit:genfun} converges absolutely, and $\E[e^{(\abs s+\abs t)\abs z}]$ is finite, so expectation and summation can be exchanged; comparing the coefficients of $s^jt^k$ gives $\E[\He_j\He_k]/(j!\,k!)=\delta_{jk}/k!$. That the $h_k$ are complete, that is, that polynomials are dense in $L^2(\gamma)$, is classical \parencite{Bogachev1998}. For (iii), the recurrence and $\He_{k-1}'=(k-1)\He_{k-2}$ give $\He_k=z\He_{k-1}-\He_{k-1}'$, and \cref{thm:toolkit:stein} applied to $f\He_{k-1}$ gives $\E[zf\He_{k-1}]=\E[f'\He_{k-1}]+\E[f\He_{k-1}']$. Subtracting $\E[f\He_{k-1}']$ from both sides,
\[
  \E[f\He_k]=\E[zf\He_{k-1}]-\E[f\He_{k-1}']=\E[f'\He_{k-1}],
\]
and induction on $k$ finishes the proof.
\end{proof}

For example, the rectified linear unit $\relu(z)=\max(z,0)$ has $\hat\relu_0=\E[\relu(z)]=1/\sqrt{2\pi}$, $\hat\relu_1=\E[\ind\{z>0\}]=\frac12$ by part (iii), and $\hat\relu_2=\E[(z^3-z)\ind\{z>0\}]/\sqrt2=\varphi(0)/\sqrt2=1/(2\sqrt\pi)$ by direct computation.

The basic fact that makes Gaussian data tractable is that Hermite polynomials of correlated Gaussians are orthogonal across degrees.

\begin{theorem}[Hermite polynomials of correlated Gaussians]\label{thm:toolkit:hermite-corr}\index{Hermite polynomials!of correlated Gaussians}
Let $(u,v)$ be jointly Gaussian with $\E u=\E v=0$, $\E u^2=\E v^2=1$ and $\E uv=\rho$. Then $\E[\He_j(u)\He_k(v)]=\delta_{jk}\,k!\,\rho^k$, that is, $\E[h_j(u)h_k(v)]=\delta_{jk}\rho^k$. In particular, for unit vectors $a,b\in\R^d$ and $x\sim\cN(0,I_d)$,
\[
  \E\big[h_j(\ip{a,x})\,h_k(\ip{b,x})\big]=\delta_{jk}\,\ip{a,b}^k.
\]
\end{theorem}

\begin{proof}
The moment generating function of $su+tv$ is $e^{(s^2+t^2+2st\rho)/2}$, so $\E[G(s,u)G(t,v)]=e^{st\rho}=\sum_k(st\rho)^k/k!$. Comparing coefficients of $s^jt^k$, as in the proof of \cref{lem:toolkit:hermite}, gives the first claim. The pair $(\ip{a,x},\ip{b,x})$ is jointly Gaussian with unit variances and covariance $\ip{a,b}$.
\end{proof}

\begin{corollary}[Dual activations]\label{cor:toolkit:dual}\index{dual activation}
Let $f,g\in L^2(\gamma)$ and let $(u,v)$ be as in \cref{thm:toolkit:hermite-corr}. Then
\[
  \E[f(u)g(v)]=\sum_{k\ge0}\hat f_k\,\hat g_k\,\rho^k,
\]
and the series converges absolutely for $\abs\rho\le1$. In particular $\rho\mapsto\E[f(u)f(v)]=\sum_k\hat f_k^2\rho^k$ is a power series with nonnegative coefficients.
\end{corollary}

\begin{proof}
The bilinear form $B(f,g)=\E[f(u)g(v)]$ satisfies $\abs{B(f,g)}\le\norm f\norm g$ by the Cauchy--Schwarz inequality, because $u$ and $v$ are standard Gaussians; so $B$ is continuous on $L^2(\gamma)\times L^2(\gamma)$. For the truncations $f_K=\sum_{k\le K}\hat f_kh_k$ and $g_K$, \cref{thm:toolkit:hermite-corr} gives $B(f_K,g_K)=\sum_{k\le K}\hat f_k\hat g_k\rho^k$. As $K\to\infty$, $f_K\to f$ and $g_K\to g$ in $L^2(\gamma)$ by \cref{lem:toolkit:hermite}(ii), so the left side tends to $B(f,g)$, while the right side converges absolutely because $\sum_k\abs{\hat f_k\hat g_k}\le\norm f\norm g$.
\end{proof}

\begin{lemma}[Price's theorem]\label{lem:toolkit:price}\index{Price's theorem}
Let $f$ and $g$ be absolutely continuous with $f,f',g,g'\in L^2(\gamma)$, and let $(u,v)$ be as in \cref{thm:toolkit:hermite-corr}. Then $\rho\mapsto\E[f(u)g(v)]$ is differentiable on $(-1,1)$, with
\[
  \frac{\partial}{\partial\rho}\E[f(u)g(v)]=\E[f'(u)g'(v)].
\]
\end{lemma}

\begin{proof}
The Hermite coefficients of $f'$ are $\widehat{(f')}_{k-1}=\sqrt k\,\hat f_k$: by \cref{lem:toolkit:hermite}(iii) with one derivative, $\E[f'\He_{k-1}]=\E[f\He_k]$, and dividing by $\sqrt{(k-1)!}$ gives the claim. \Cref{cor:toolkit:dual} applied to $f'$ and $g'$ gives $\E[f'(u)g'(v)]=\sum_{k\ge1}k\hat f_k\hat g_k\rho^{k-1}$, which is the termwise derivative of the power series $\sum_k\hat f_k\hat g_k\rho^k$. A power series with radius of convergence at least $1$ may be differentiated termwise on $(-1,1)$.
\end{proof}

\section{Spherical harmonics}\label{sec:toolkit:sphere}

Let $\sigma$ be the uniform probability measure on the unit sphere $\Sph^{d-1}\subset\R^d$, $d\ge3$. A \term{spherical harmonic}\index{spherical harmonics} of degree $k$ is the restriction to the sphere of a homogeneous polynomial of degree $k$ on $\R^d$ that is harmonic (its Laplacian vanishes). They form a finite-dimensional space $\cH_k$, and the following facts are classical \parencite[Chapter~IV]{SteinWeiss1971}; see also \textcite{AtkinsonHan2012} and \textcite{DaiXu2013}.

\begin{theorem}[Spherical harmonics]\label{thm:toolkit:addition}\index{addition theorem}\index{Funk--Hecke formula}
\begin{enumerate}
\item $L^2(\Sph^{d-1},\sigma)$ is the orthogonal direct sum of the spaces $\cH_k$, $k\ge0$, and
\[
  N_{d,k}\defeq\dim\cH_k=\binom{d+k-1}{k}-\binom{d+k-3}{k-2},
\]
with the second binomial coefficient read as $0$ for $k<2$. For fixed $k$, $N_{d,k}=\frac{d^k}{k!}\big(1+O(1/d)\big)$ as $d\to\infty$.
\item \emph{(Addition theorem.)} Let $Y_1,\dots,Y_{N_{d,k}}$ be an orthonormal basis of $\cH_k$ in $L^2(\sigma)$. The function $Z_k(x,y)=\sum_jY_j(x)Y_j(y)$ does not depend on the choice of basis, is the kernel of the orthogonal projection onto $\cH_k$, and depends only on $\ip{x,y}$: $Z_k(x,y)=N_{d,k}\,P_{d,k}(\ip{x,y})$, where $P_{d,k}$ is the Gegenbauer polynomial of degree $k$ for dimension $d$, normalized by $P_{d,k}(1)=1$. In particular $Z_k(x,x)=N_{d,k}$ for every $x$.
\item \emph{(Funk--Hecke formula.)} For every bounded measurable $F\colon[-1,1]\to\R$ and every $Y\in\cH_k$,
\[
  \int F(\ip{x,y})\,Y(y)\dd\sigma(y)=\mu_k(F)\,Y(x),
\]
with $\mu_k(F)=\frac{\abs{\Sph^{d-2}}}{\abs{\Sph^{d-1}}}\int_{-1}^1F(t)P_{d,k}(t)(1-t^2)^{(d-3)/2}\dd t$.
Consequently every kernel of the form $K(x,y)=F(\ip{x,y})$ on the sphere has the spherical harmonics as eigenfunctions, with the eigenvalue $\mu_k(F)$ shared by all of $\cH_k$.
\end{enumerate}
\end{theorem}

\begin{proof}[Proof of the parts used in this book]
We take the decomposition and the dimension formula in (i) from the references. The asymptotics follow from the formula: $\binom{d+k-1}k=\frac{d^k}{k!}(1+O(1/d))$, while $\binom{d+k-3}{k-2}=O(d^{k-2})$.

For (ii), the function $Z_k(\cdot,y)$ satisfies $\ip{Y,Z_k(\cdot,y)}=Y(y)$ for every $Y\in\cH_k$, which characterizes it as the reproducing kernel of the finite-dimensional space $\cH_k$; so it does not depend on the basis. If $U$ is an orthogonal matrix, then $Y\mapsto Y\circ U$ maps $\cH_k$ onto itself and preserves $L^2(\sigma)$ inner products, because harmonicity, homogeneity and $\sigma$ are rotation invariant; so $\{Y_j\circ U\}$ is another orthonormal basis, and $Z_k(Ux,Uy)=Z_k(x,y)$. Since every pair of points with the same inner product is mapped to every other such pair by some rotation, $Z_k(x,y)$ depends only on $\ip{x,y}$. In particular $Z_k(x,x)$ is a constant $c$, and integrating $\sum_jY_j(x)^2$ against $\sigma$ gives $c=\sum_j\norm{Y_j}^2=N_{d,k}$. The identification of $Z_k(x,y)/N_{d,k}$ with the Gegenbauer polynomial and part (iii) are in \textcite[Chapter~IV]{SteinWeiss1971}.
\end{proof}

For $d=3$, the polynomial $P_{3,k}$ is the Legendre polynomial of degree $k$, and $N_{3,k}=2k+1$.

\section{Wishart matrices and the Marchenko--Pastur law}\label{sec:toolkit:mp}

Let $X\in\R^{n\times p}$ have independent $\cN(0,1)$ entries, with rows $x_1,\dots,x_n\in\R^p$, and let $\hat\Sigma=X^\top X/n$ be the sample covariance matrix; its population counterpart is $I_p$. The first result is exact at every size.

\begin{theorem}[Mean of the inverse Wishart matrix]\label{thm:toolkit:inv-wishart}\index{Wishart matrix}
If $n\ge p+2$, then $\E[(X^\top X)^{-1}]=I_p/(n-p-1)$. If $n\in\{p,p+1\}$, then $X^\top X$ is invertible almost surely and $\E\tr[(X^\top X)^{-1}]=+\infty$.
\end{theorem}

\begin{proof}
The plan is to use rotational invariance to reduce to one diagonal entry, and then to compute that entry with a Schur complement. When $n\ge p$, the columns of $X$ are linearly independent almost surely, so $X^\top X$ is invertible. For every orthogonal $U\in\R^{p\times p}$ the matrix $XU$ has the same law as $X$, so $M=\E[(X^\top X)^{-1}]$, if finite, satisfies $U^\top MU=M$ for all $U$, and is therefore a multiple $cI_p$ of the identity; also all diagonal entries of $(X^\top X)^{-1}$ have the same law. Let $X_1$ be the first column of $X$ and $P$ the orthogonal projection onto the span of the other $p-1$ columns. By the formula for the inverse of a block matrix,
\[
  \big[(X^\top X)^{-1}\big]_{11}=\frac{1}{X_1^\top X_1-X_1^\top P X_1}=\frac1{\norm{(I-P)X_1}^2}.
\]
Conditionally on the other columns, $I-P$ is a projection of rank $n-p+1$ independent of $X_1\sim\cN(0,I_n)$, so $\norm{(I-P)X_1}^2$ has the $\chi^2$ distribution with $k=n-p+1$ degrees of freedom. For $k\ge3$, $\E[1/\chi^2_k]=\frac{\Gamma(k/2-1)}{2\Gamma(k/2)}=\frac1{k-2}$, and for $k\in\{1,2\}$ the expectation is infinite, because the $\chi^2_k$ density does not vanish at $0$. With $k-2=n-p-1$ this gives $c=1/(n-p-1)$ when $n\ge p+2$, and an infinite expected trace when $n-p\in\{0,1\}$.
\end{proof}

In the proportional regime, $p/n\to\gamma\in(0,\infty)$, the eigenvalues of $\hat\Sigma$ do not concentrate at $1$ but spread over an interval. We describe their distribution through the \term{Stieltjes transform}\index{Stieltjes transform}
\[
  s_p(z)=\frac1p\tr\big(\hat\Sigma-zI_p\big)^{-1}=\frac1p\sum_{j=1}^p\frac1{\lambda_j(\hat\Sigma)-z},
\]
which for real $z=-\lambda<0$ is a positive, decreasing, convex function of $\lambda$. Its derivative gives the other trace that appears in ridge regression: $\frac{d}{d\lambda}s_p(-\lambda)=-\frac1p\tr(\hat\Sigma+\lambda I)^{-2}$.

\begin{theorem}[Marchenko--Pastur law]\label{thm:toolkit:mp}\index{Marchenko--Pastur law}
Let $n,p\to\infty$ with $p/n\to\gamma\in(0,\infty)$. For every $\lambda>0$, almost surely,
\[
  \frac1p\tr\big(\hat\Sigma+\lambda I_p\big)^{-1}\to m(-\lambda),\qquad \frac1p\tr\big(\hat\Sigma+\lambda I_p\big)^{-2}\to m'(-\lambda),
\]
where $m(-\lambda)$ is the unique positive solution $m$ of
\begin{equation}\label{eq:toolkit:mp}
  \gamma zm^2-(1-\gamma-z)m+1=0\qquad\text{at }z=-\lambda,
\end{equation}
namely $m(-\lambda)=\big(\gamma-1-\lambda+\sqrt{(\gamma-1-\lambda)^2+4\gamma\lambda}\big)/(2\gamma\lambda)$, and $m'$ denotes the derivative of $m(z)$ in $z$. The function $m$ is the Stieltjes transform of the Marchenko--Pastur law $\mu_\gamma$, which has density $\sqrt{(b-x)(x-a)}/(2\pi\gamma x)$ on $[a,b]=[(1-\sqrt\gamma)^2,(1+\sqrt\gamma)^2]$ and, when $\gamma>1$, an atom of mass $1-1/\gamma$ at $0$.
\end{theorem}

\begin{proof}
The plan is to derive an approximate quadratic equation for $\E s_p$ from a resolvent identity, using that each row of $X$ is independent of the matrix built from the other rows, and then to show that $s_p$ concentrates around its mean. Fix $\lambda>0$, write $z=-\lambda$, $R=(\hat\Sigma+\lambda I)^{-1}$, $s_p=\frac1p\tr R$ and $\gamma_n=p/n$. Every resolvent below is positive definite with operator norm at most $1/\lambda$.

\emph{A resolvent identity.} Since $\hat\Sigma R=I+zR$, taking normalized traces gives $1+zs_p=\frac1p\tr(\hat\Sigma R)=\frac1{pn}\sum_ix_i^\top Rx_i$. Let $R_{(i)}=(\hat\Sigma-\frac1nx_ix_i^\top+\lambda I)^{-1}$ be the resolvent without the $i$th row and $q_i=\frac1nx_i^\top R_{(i)}x_i\ge0$. The Sherman--Morrison formula gives $Rx_i=R_{(i)}x_i/(1+q_i)$, hence $x_i^\top Rx_i=nq_i/(1+q_i)$, and
\begin{equation}\label{eq:toolkit:resolvent}
  1+zs_p=\frac1p\sum_{i=1}^n\frac{q_i}{1+q_i}=\frac1{\gamma_n}\Big(1-\frac1n\sum_{i=1}^n\frac1{1+q_i}\Big).
\end{equation}

\emph{Removing one row changes little.} Again by Sherman--Morrison, $R_{(i)}-R=\frac1nR_{(i)}x_ix_i^\top R_{(i)}/(1+q_i)$, a positive semidefinite matrix of rank one with trace $\frac1nx_i^\top R_{(i)}^2x_i/(1+q_i)\le\frac1\lambda\cdot\frac{q_i}{1+q_i}\le\frac1\lambda$. Hence $0\le\tr R_{(i)}-\tr R\le1/\lambda$. As a consequence, replacing the row $x_i$ by an independent copy changes $s_p$ by at most $1/(p\lambda)$.

\emph{Concentration.} By the rank-one bound, $s_p$ is a function of $n$ independent rows with bounded differences $1/(p\lambda)$, so McDiarmid's inequality (\cref{thm:toolkit:mcdiarmid}) gives $\Pr(\abs{s_p-\E s_p}\ge\eps)\le2\exp(-2\eps^2p^2\lambda^2/n)$. Since $p^2/n\to\infty$ at a linear rate, these probabilities are summable, and the Borel--Cantelli lemma gives $s_p-\E s_p\to0$ almost surely. The Efron--Stein inequality \parencite[Theorem~3.1]{BoucheronLugosiMassart2013} gives in the same way $\Var(s_p)\le n/(2p^2\lambda^2)$.

\emph{The quadratic form concentrates.} The row $x_i\sim\cN(0,I_p)$ is independent of $R_{(i)}$, and for a fixed symmetric matrix $A$ we have $\E[x^\top Ax]=\tr A$ and $\Var(x^\top Ax)=2\tr(A^2)$ when $x\sim\cN(0,I_p)$. Hence, with $s_{(i)}=\frac1p\tr R_{(i)}$, conditionally on $R_{(i)}$ the quadratic form $q_i$ has mean $\gamma_ns_{(i)}$ and variance $\frac2{n^2}\tr R_{(i)}^2\le\frac{2p}{n^2\lambda^2}$. We split $q_i-\gamma_n\E s_p$ into $q_i-\gamma_ns_{(i)}$, $\gamma_n(s_{(i)}-s_p)$ and $\gamma_n(s_p-\E s_p)$, and bound the three in mean by the conditional variance, the rank-one bound, and the variance bound:
\[
  \E\abs{q_i-\gamma_n\E s_p}\le\frac{\sqrt{2\gamma_n/n}}{\lambda}+\frac{\gamma_n}{p\lambda}+\gamma_n\sqrt{\Var(s_p)}\eqqcolon\epsilon_n,
\]
and $\epsilon_n\to0$.

\emph{The limiting equation.} Take expectations in \eqref{eq:toolkit:resolvent} and write $\bar s=\E s_p\in[0,1/\lambda]$. Because $q_i\ge0$ and $\gamma_n\bar s\ge0$, the function $a\mapsto1/(1+a)$ is $1$-Lipschitz on the relevant range, so $\E[1/(1+q_i)]=1/(1+\gamma_n\bar s)+O(\epsilon_n)$, and
\[
  1+z\bar s=\frac1{\gamma_n}\Big(1-\frac1{1+\gamma_n\bar s}\Big)+O(\epsilon_n)=\frac{\bar s}{1+\gamma_n\bar s}+O(\epsilon_n).
\]
Multiplying by $1+\gamma_n\bar s\le1+\gamma_n/\lambda$ gives $(1+z\bar s)(1+\gamma_n\bar s)-\bar s\to0$. Along every subsequence on which $\bar s$ converges, its limit $s^*\ge0$ satisfies $(1+zs^*)(1+\gamma s^*)=s^*$, which expands to \eqref{eq:toolkit:mp}. At $z=-\lambda$ the quadratic $\gamma zm^2-(1-\gamma-z)m+1$ has roots whose product is $1/(\gamma z)<0$, so exactly one root is positive; and $0$ is not a root. Hence every subsequential limit of $\bar s$ equals the positive root $m(-\lambda)$, so $\bar s\to m(-\lambda)$, and with the concentration step, $s_p\to m(-\lambda)$ almost surely.

\emph{Derivatives.} The functions $\lambda\mapsto\frac1p\tr(\hat\Sigma+\lambda I)^{-1}$ are convex on $(0,\infty)$. Almost surely they converge on a countable dense set of $\lambda$ simultaneously, hence, being monotone with a continuous limit, at every $\lambda>0$. Pointwise convergence of convex functions on an open interval implies convergence of their derivatives at every point where the limit is differentiable \parencite[Theorem~24.5]{Rockafellar1970}, and $\lambda\mapsto m(-\lambda)$ is differentiable. The derivative of $\frac1p\tr(\hat\Sigma+\lambda I)^{-1}$ in $\lambda$ is $-\frac1p\tr(\hat\Sigma+\lambda I)^{-2}$, and the derivative of $m(-\lambda)$ in $\lambda$ is $-m'(-\lambda)$, which gives the second limit.

The identification of $m$ with the Stieltjes transform of $\mu_\gamma$ is a computation with the explicit density, which we omit; see \textcite{MarchenkoPastur1967} and \textcite[Chapter~3]{BaiSilverstein2010}. A probability measure on $[0,\infty)$ is determined by its Stieltjes transform on the negative half-line, because the transform extends analytically to $\C\setminus[0,\infty)$, where the Stieltjes inversion formula recovers the measure.
\end{proof}

\begin{remark}\label{rem:toolkit:mp-edges}
The edges of the Marchenko--Pastur law explain the double descent peak of \cref{ch:doubledescent}: the smallest nonzero eigenvalue of $\hat\Sigma$ tends to $(1-\sqrt\gamma)^2$, which vanishes at $\gamma=1$, so the minimum-norm least squares solution, which inverts $\hat\Sigma$ on its range, becomes unstable exactly at the interpolation threshold. (The convergence of the extreme eigenvalues to the edges of the support needs an argument beyond the Stieltjes transform; see \textcite[Chapter~5]{BaiSilverstein2010}.)
\end{remark}

\section{Martingales, Gr\"onwall, and the ODE method}\label{sec:toolkit:ode}

Several chapters analyze a stochastic recursion in which each step moves the state by an amount of order $1/d$, with random fluctuations of the same order, and claim that over $d$ steps the state follows a deterministic ordinary differential equation. This section proves such a claim in the generality the chapters need. We write $\cF_t$ for the $\sigma$-algebra generated by the first $t$ steps.

\begin{lemma}[Doob's maximal inequality]\label{lem:toolkit:doob}\index{Doob's maximal inequality}
Let $(M_t)_{t=0}^T$ be a martingale with values in $\R^D$ and $\E\norm{M_T}^2<\infty$. Then for every $a>0$,
\[
  \Pr\Big(\max_{t\le T}\norm{M_t}\ge a\Big)\le\frac{\E\norm{M_T}^2}{a^2},\qquad \E\Big[\max_{t\le T}\norm{M_t}^2\Big]\le4\,\E\norm{M_T}^2 .
\]
If the increments satisfy $\E[\norm{M_{t+1}-M_t}^2\mid\cF_t]\le v_t$, then $\E\norm{M_T}^2\le\norm{M_0}^2+\sum_{t<T}\E v_t$.
\end{lemma}

\begin{proof}
The process $\norm{M_t}$ is a nonnegative submartingale, by Jensen's inequality for conditional expectations, and the two inequalities are Doob's maximal and $L^2$ inequalities for nonnegative submartingales \parencite{Williams1991}. For the last claim, increments of a square-integrable martingale are orthogonal: $\E\ip{M_{t+1}-M_t,M_s-M_{s-1}}=0$ for $s\le t$, by conditioning on $\cF_t$; so $\E\norm{M_T-M_0}^2=\sum_t\E\norm{M_{t+1}-M_t}^2$, and the same conditioning shows $\E\ip{M_T-M_0,M_0}=0$.
\end{proof}

\begin{theorem}[Freedman's inequality]\label{thm:toolkit:freedman}\index{Freedman's inequality}
Let $(M_t)$ be a real martingale with $M_0=0$ whose increments satisfy $M_{t+1}-M_t\le c$, and let $V_t=\sum_{s<t}\E[(M_{s+1}-M_s)^2\mid\cF_s]$. Then for all $a,v>0$,
\[
  \Pr\big(M_t\ge a\text{ and }V_t\le v\text{ for some }t\big)\le\exp\Big(-\frac{a^2}{2(v+ca)}\Big).
\]
In particular (Azuma--Hoeffding), if $\abs{M_{t+1}-M_t}\le c_t$ almost surely, then $\Pr(\abs{M_T}\ge a)\le2\exp\big(-a^2/(2\sum_{t<T}c_t^2)\big)$.
\end{theorem}

Freedman's inequality is proved by \textcite{Freedman1975}. The Azuma--Hoeffding inequality follows from the exponential method; see \textcite{BoucheronLugosiMassart2013}.

\begin{theorem}[Gr\"onwall's inequality]\label{thm:toolkit:gronwall}\index{Gr\"onwall's inequality}
\begin{enumerate}
\item If $u\colon[0,T]\to[0,\infty)$ is continuous and $u(t)\le a+\int_0^tb\,u(s)\dd s$ for constants $a,b\ge0$, then $u(t)\le ae^{bt}$.
\item If $u_0,u_1,\dots$ are nonnegative and $u_t\le a+b\sum_{s<t}u_s$ for all $t$, then $u_t\le a(1+b)^t\le ae^{bt}$.
\end{enumerate}
\end{theorem}

\begin{proof}
For (i), let $U(t)=a+\int_0^tbu$; then $U'=bu\le bU$, so $(e^{-bt}U)'\le0$ and $u\le U(t)\le e^{bt}U(0)=ae^{bt}$. For (ii), let $S_t=a+b\sum_{s<t}u_s$; then $u_t\le S_t$ and $S_{t+1}=S_t+bu_t\le(1+b)S_t$, so $u_t\le S_t\le(1+b)^tS_0=(1+b)^ta$.
\end{proof}

\begin{theorem}[The ODE method]\label{thm:toolkit:ode-method}\index{ODE method}
Let $U\subseteq\R^D$ be open and let $F\colon U\to\R^D$ be $L$-Lipschitz and bounded by $B$ on $U$. Let $(Z_t)_{t\ge0}$ be adapted random vectors with
\[
  Z_{t+1}=Z_t+\frac1d\big(F(Z_t)+\xi_{t+1}+\epsilon_{t+1}\big),
\]
where, on the event $\{Z_t\in U\}$, the noise satisfies $\E[\xi_{t+1}\mid\cF_t]=0$ and $\E[\norm{\xi_{t+1}}^2\mid\cF_t]\le\varsigma^2$, and the error satisfies $\norm{\epsilon_{t+1}}\le\epsilon$. Let $z\colon[0,T]\to U$ solve $\dot z=F(z)$, and suppose that every point within distance $\rho$ of the path $z([0,T])$ lies in $U$. For $a>0$ put
\[
  \Delta=\big(\norm{Z_0-z(0)}+a+T\epsilon+TLB/(2d)\big)e^{LT}.
\]
If $\Delta<\rho$, then
\[
  \Pr\Big(\max_{0\le t\le Td}\norm{Z_t-z(t/d)}>\Delta\Big)\le\frac{\varsigma^2T}{a^2d}.
\]
\end{theorem}

\begin{proof}
The plan is to compare the recursion with the Euler scheme for the ODE, to control the accumulated noise with Doob's inequality, and to close the argument with Gr\"onwall's inequality, using a stopping time to stay inside $U$. Write $D_t=Z_t-z(t/d)$. Integrating the ODE over one step gives $z(\frac{t+1}d)=z(\frac td)+\frac1dF(z(\frac td))+r_t$ with the Euler error
\[
  \norm{r_t}\le\int_{t/d}^{(t+1)/d}\norm[\big]{F(z(u))-F\big(z(\tfrac td)\big)}\dd u\le\frac{LB}{2d^2},
\]
because $\norm{z(u)-z(t/d)}\le B(u-t/d)$. Let $\tau$ be the first $t$ with $\norm{D_t}>\rho$; for $t<\tau$ the point $Z_t$ lies within $\rho$ of the path, hence in $U$. Let $M_t=\sum_{s<t\wedge\tau}\xi_{s+1}$, a martingale because $\tau$ is a stopping time. Subtracting the two recursions and summing, for $t\le\tau\wedge\lfloor Td\rfloor$,
\[
  D_t=D_0+\frac1d\sum_{s<t}\big(F(Z_s)-F(z(\tfrac sd))\big)+\frac1dM_t+\frac1d\sum_{s<t}\epsilon_{s+1}-\sum_{s<t}r_s,
\]
Taking norms and using the Lipschitz bound,
\[
  \norm{D_t}\le\Big(\norm{D_0}+\frac1d\max_{u\le Td}\norm{M_u}+T\epsilon+\frac{TLB}{2d}\Big)+\frac Ld\sum_{s<t}\norm{D_s}.
\]
On the event $E=\{\max_{u\le Td}\norm{M_u}\le ad\}$, the discrete Gr\"onwall inequality (\cref{thm:toolkit:gronwall}(ii)) gives $\norm{D_t}\le\Delta$ for every $t\le\tau\wedge\lfloor Td\rfloor$. Since $\Delta<\rho$, this rules out $\tau\le\lfloor Td\rfloor$, so on $E$ the bound holds for all $t\le Td$. Finally, the increments of $M$ have conditional second moment at most $\varsigma^2$, so by \cref{lem:toolkit:doob}, $\Pr(E^c)\le\E\norm{M_{\lfloor Td\rfloor}}^2/(ad)^2\le\varsigma^2Td/(a^2d^2)=\varsigma^2T/(a^2d)$.
\end{proof}

In the applications, $\epsilon$ and $1/d$ tend to zero as the dimension grows and $a$ is chosen as $d^{-1/4}$, say; then the recursion stays within $O(d^{-1/4})$ of the ODE on the time interval $[0,T]$ with probability $1-O(d^{-1/2})$, and choosing $a$ a little larger than $d^{-1/2}$ shows that the fluctuations are of order $d^{-1/2}$.

\section{Fourier analysis on finite groups}\label{sec:toolkit:fourier}

\paragraph{The cyclic group.} Let $p\ge2$ be an integer, $\Z_p=\{0,1,\dots,p-1\}$ with addition modulo $p$, and $e_p(x)=e^{2\pi ix/p}$. For $f\colon\Z_p\to\C$ define the Fourier transform $\hat f(\omega)=\frac1p\sum_{a\in\Z_p}f(a)\,e_p(-\omega a)$ for $\omega\in\Z_p$, and the convolution $(f*g)(c)=\frac1p\sum_af(a)g(c-a)$.

\begin{theorem}[Fourier analysis on $\Z_p$]\label{thm:toolkit:fourier-zp}\index{Fourier analysis!on $\Z_p$}
For all $f,g,u,u',v\colon\Z_p\to\C$:
\begin{enumerate}
\item (inversion) $f(a)=\sum_\omega\hat f(\omega)e_p(\omega a)$;
\item (Plancherel) $\frac1p\sum_af(a)\overline{g(a)}=\sum_\omega\hat f(\omega)\overline{\hat g(\omega)}$;
\item (convolution) $\widehat{f*g}=\hat f\,\hat g$;
\item (triple correlation) $\frac1{p^2}\sum_{a,b}u(a)u'(b)v(a+b)=\sum_\omega\hat u(-\omega)\,\hat u'(-\omega)\,\hat v(\omega)$.
\end{enumerate}
If $f$ is real, then $\hat f(-\omega)=\overline{\hat f(\omega)}$, and $f(a)=\hat f(0)+\sum_{\omega}\abs{\hat f(\omega)}\cos(2\pi\omega a/p+\arg\hat f(\omega))$, where for odd $p$ the sum runs over $\omega=1,\dots,(p-1)/2$ with each term doubled.
\end{theorem}

\begin{proof}
Everything follows from the orthogonality relation
\[
  \frac1p\sum_{a\in\Z_p}e_p(\omega a)=\ind\{\omega=0\},
\]
which holds because for $\omega\ne0$ the sum is a geometric series of $p$th roots of unity whose ratio is not $1$. For (i), substitute the definition of $\hat f$ into the right side and sum over $\omega$ first. For (ii), substitute (i) for $f$ and $g$ and use orthogonality. For (iii), $\widehat{f*g}(\omega)=\frac1{p^2}\sum_{a,c}f(a)g(c-a)e_p(-\omega a)e_p(-\omega(c-a))=\hat f(\omega)\hat g(\omega)$, substituting $b=c-a$. For (iv), expand $v$ by (i): $\sum_{a,b}u(a)u'(b)v(a+b)=\sum_\omega\hat v(\omega)\sum_au(a)e_p(\omega a)\sum_bu'(b)e_p(\omega b)=p^2\sum_\omega\hat v(\omega)\hat u(-\omega)\hat u'(-\omega)$. The statements for real $f$ follow from conjugating the definition, and from pairing the terms $\omega$ and $-\omega$ in (i).
\end{proof}

\paragraph{The Boolean cube.} For $S\subseteq[d]=\{1,\dots,d\}$ the \term{parity}\index{parity function} or character $\chi_S(x)=\prod_{i\in S}x_i$ is a function on $\{\pm1\}^d$, with $\chi_\emptyset=1$. Let $x$ be uniform on $\{\pm1\}^d$ and $\ip{f,g}=\E[f(x)g(x)]$.

\begin{theorem}[Fourier--Walsh expansion]\label{thm:toolkit:fourier-cube}\index{Fourier analysis!on the Boolean cube}
The $2^d$ characters form an orthonormal basis of the functions on $\{\pm1\}^d$. Every $f$ equals $\sum_S\hat f(S)\chi_S$ with $\hat f(S)=\ip{f,\chi_S}$, and $\E[f(x)^2]=\sum_S\hat f(S)^2$.
\end{theorem}

\begin{proof}
Since $x_i^2=1$, $\chi_S\chi_T=\chi_{S\triangle T}$, and by independence of the coordinates $\E[\chi_{S\triangle T}(x)]=\prod_{i\in S\triangle T}\E[x_i]$, which is $1$ if $S=T$ and $0$ otherwise. So the characters are orthonormal, and $2^d$ orthonormal vectors in a $2^d$-dimensional space form a basis.
\end{proof}

The \term{degree} of $f$ is the largest $\abs S$ with $\hat f(S)\ne0$. Low-degree functions cannot have heavy tails, a fact known as hypercontractivity. We prove the fourth-moment version, which suffices for most purposes, and state the general one.

\begin{theorem}[Bonami's lemma and hypercontractivity]\label{thm:toolkit:bonami}\index{hypercontractivity}\index{Bonami's lemma}
If $f\colon\{\pm1\}^d\to\R$ has degree at most $k$, then $\E[f^4]\le9^k\big(\E[f^2]\big)^2$. More generally, for every $q\ge2$, $\big(\E\abs f^q\big)^{1/q}\le(q-1)^{k/2}\big(\E f^2\big)^{1/2}$, and consequently, for $t\ge(2e)^{k/2}$,
\[
  \Pr\big(\abs{f(x)}\ge t\,(\E f^2)^{1/2}\big)\le\exp\Big(-\frac k{2e}\,t^{2/k}\Big).
\]
The same inequalities hold for polynomials of degree at most $k$ in a standard Gaussian vector.
\end{theorem}

\begin{proof}[Proof of the fourth-moment bound]
The plan is induction on the number $d$ of variables. For $d=0$, $f$ is constant and the bound holds with equality. For $d\ge1$, write $f(x)=x_dg(x')+h(x')$, where $x'=(x_1,\dots,x_{d-1})$, $g$ has degree at most $k-1$ and $h$ degree at most $k$. Expanding $(x_dg+h)^4$ and averaging over $x_d$ removes the odd powers of $x_d$, and $\E[f^2]=\E[g^2]+\E[h^2]$ in the same way. Hence $\E f^4=\E g^4+6\,\E[g^2h^2]+\E h^4$. We bound the outer terms by the induction hypothesis and the middle one by the Cauchy--Schwarz inequality and the induction hypothesis again:
\[
  \E f^4\le9^{k-1}(\E g^2)^2+6\cdot3^{k-1}3^k\,\E g^2\,\E h^2+9^k(\E h^2)^2 .
\]
Since $6\cdot3^{2k-1}=2\cdot9^k$ and $9^{k-1}\le9^k$, the right side is at most $9^k(\E g^2+\E h^2)^2=9^k(\E f^2)^2$.
\end{proof}

The general moment bound is the Bonami--Beckner hypercontractive inequality, and the tail bound follows from it by Markov's inequality with an optimized $q$; see \textcite[Chapters~9 and~11]{ODonnell2014}, where the Gaussian case is also treated.

\section{Concentration inequalities}\label{sec:toolkit:conc}

We collect the standard inequalities used in the chapters; proofs are in \textcite{BoucheronLugosiMassart2013} and \textcite{Vershynin2018}.

\begin{theorem}[Hoeffding's inequality]\label{thm:toolkit:hoeffding}\index{Hoeffding's inequality}
If $X_1,\dots,X_n$ are independent with $X_i\in[a_i,b_i]$ almost surely, then for every $t>0$,
\[
  \Pr\Big(\Big|\sum_i(X_i-\E X_i)\Big|\ge t\Big)\le2\exp\Big(-\frac{2t^2}{\sum_i(b_i-a_i)^2}\Big).
\]
\end{theorem}

\begin{theorem}[Bernstein's inequality]\label{thm:toolkit:bernstein}\index{Bernstein's inequality}
If $X_1,\dots,X_n$ are independent with $\E X_i=0$, $\abs{X_i}\le c$ almost surely and $\sum_i\E X_i^2\le v$, then for every $t>0$,
\[
  \Pr\Big(\Big|\sum_iX_i\Big|\ge t\Big)\le2\exp\Big(-\frac{t^2}{2(v+ct/3)}\Big).
\]
\end{theorem}

\begin{theorem}[McDiarmid's bounded differences inequality]\label{thm:toolkit:mcdiarmid}\index{bounded differences inequality}
Let $X_1,\dots,X_n$ be independent and let $F$ satisfy $\abs{F(x)-F(x')}\le c_i$ whenever $x$ and $x'$ differ only in the $i$th coordinate. Then $\Pr(\abs{F(X)-\E F(X)}\ge t)\le2\exp\big(-2t^2/\sum_ic_i^2\big)$.
\end{theorem}

\begin{theorem}[Gaussian concentration]\label{thm:toolkit:gauss-conc}\index{Gaussian concentration}
If $x\sim\cN(0,I_d)$ and $f\colon\R^d\to\R$ is $L$-Lipschitz, then $\Pr(\abs{f(x)-\E f(x)}\ge t)\le2e^{-t^2/(2L^2)}$ for every $t>0$.
\end{theorem}

Gaussian concentration is due to \textcite{TsirelsonIbragimovSudakov1976}; it shows, for example, that $\norm x$ concentrates around $\sqrt d$ with fluctuations of order one, since the norm is $1$-Lipschitz, and that $\norm{x}^2/d=1+O(d^{-1/2})$ with high probability.

\begin{exercises}
\item \consolidation Use \cref{lem:toolkit:hermite}(iii) to compute the Hermite coefficients of $\phi(z)=\abs z$ and of $\phi(z)=\sgn(z)$, and check that $\sum_k\hat\phi_k^2=\E[\phi(z)^2]$ numerically for the first twenty coefficients.\label{ex:toolkit:coeffs}
\item \consolidation Let $u$ and $v$ be standard Gaussians with correlation $\rho$. Show that
\[
  \E[\relu(u)\relu(v)]=\tfrac1{2\pi}\big(\sqrt{1-\rho^2}+(\pi-\arccos\rho)\rho\big).
\] \hint{By \cref{lem:toolkit:price} the $\rho$-derivative of the left side is $\E[\ind\{u>0\}\ind\{v>0\}]=\frac14+\frac{\arcsin\rho}{2\pi}$; integrate from $\rho=0$.}
\item \consolidation Prove that $N_{d,1}=d$, $N_{d,2}=d(d+1)/2-1$, and identify $\cH_1$ and $\cH_2$ explicitly.
\item \consolidation\hard Show from the explicit formula for $m(-\lambda)$ that $\lambda\,m(-\lambda)\to1-1/\gamma$ as $\lambda\downarrow0$ when $\gamma>1$, and $\lambda\,m(-\lambda)\to0$ when $\gamma<1$; interpret the limit through the atom of $\mu_\gamma$ at $0$. \hint{$\lambda m(-\lambda)=\int\frac{\lambda}{x+\lambda}\dd\mu_\gamma(x)$.}
\item \consolidation Check \cref{thm:toolkit:inv-wishart} by simulation for $n=30$ and $p\in\{5,15,25,27\}$, and explain why the empirical average of $\tr(X^\top X)^{-1}$ is unstable at $p=29$.
\item \variation Prove the variance formula $\Var(x^\top Ax)=2\tr(A^2)$ used in the proof of \cref{thm:toolkit:mp}, and generalize it to $x\sim\cN(0,\Sigma)$.
\item \consolidation\hard Show that the error bound of \cref{thm:toolkit:ode-method} is sharp in its dependence on $d$: for $D=1$, $F=0$ and $\pm1$ noise, the maximum of $\abs{Z_t-Z_0}$ over $t\le d$ is of order $d^{-1/2}$.
\end{exercises}

\notesandrefs

Gaussian integration by parts is often called Stein's lemma, after its role in Stein's method for normal approximation \parencite{Stein1981}; Hermite polynomials and their generating function are classical, and \textcite{ODonnell2014} develops them alongside Fourier analysis on the cube in a way that matches this book's use. The expansion of $\E[\phi(u)\phi(v)]$ in powers of the correlation, which \textcite{DanielyFrostigSinger2016} called the dual activation, is a form of Mehler's formula. Price's theorem is from \textcite{Price1958}.

The Marchenko--Pastur law appeared in \textcite{MarchenkoPastur1967}; the resolvent argument above follows the standard Stieltjes-transform proof \parencite{BaiSilverstein2010}, restricted to the negative real axis, where every estimate is real and elementary. The ODE method for stochastic approximation goes back to \textcite{Ljung1977} and \textcite{KushnerYin2003}; its high-dimensional form, in which the step size is $1/d$ and the state is a finite set of overlaps, is the method of \textcite{SaadSolla1995} made rigorous, for example by \textcite{WangMattinglyLu2017} and \textcite{GoldtEtAl2019}.

\printchapterbib
\end{refsection}


\backmatter
\printindex
\specialchapter{How This Book Was Made}{Colophon}{34}

\addsec{Roles}

\emph{The Hydrogen Atoms of Deep Learning} was written by Claude Opus 5.5, an AI model made by Anthropic, in October 2026. Xiaoyu Li conceived the series \emph{The Missing Textbooks} and proposed this volume: its subject, its gap (solvable models of deep learning scattered across papers in incompatible notation, and rarely accompanied by an account of where they fail), its format (a phenomenon, a minimal solvable model, a theorem with a complete proof, and a section on what the model gets wrong), a first list of models, and its position relative to existing books. Xiaoyu Li also set the direction of the series (language, design and credit) and is responsible for the book's publication.

\addsec{Process}

A written plan fixed the thesis, the sixteen models, the global notation, and a specification for each chapter, including the editor's plan for each main theorem. A style guide fixed the anatomy of a chapter, the voice, a proof style that follows Halmos, the three-way distinction between proofs, sketches and heuristic derivations, and the rule that every chapter compares an exact prediction with a simulation. Separate instances of the same model then drafted the chapters, at most five at a time, each re-deriving its results, checking them numerically, verifying its references with web search, and writing the scripts behind its figures. Five chapters were begun by one set of instances and completed by another; two of these received a separate pass to finish and check the draft before refereeing. A coordinating instance wrote the prologue, preface, notation, epilogue and this colophon, designed the volume's additions to the series typography, and integrated the whole.

Each chapter, the prologue and the appendix then received one independent referee pass from a separate instance, which re-derived the proofs, re-ran the figure scripts, checked constants and exercises numerically, and checked references and historical claims by web search. The referees corrected what they found directly. They found few errors in the main theorems, and several gaps in proofs, which they closed or marked; most of their corrections were to numbers quoted from simulations, to the statements and answers of exercises, and to claims about what other papers prove, many of which were softened or removed because they could not be confirmed. The epilogue was written last and was not refereed.

\addsec{Limits}

No human has yet checked every proof. Several results are presented as heuristic derivations, marked with $\Diamond$, because no proof is known to the author or because a proof would not fit; they are checked against simulations but not proved. Web search was available while the book was written, but the full texts of most papers could not be opened, so references were checked against search results, abstracts and secondary accounts. Some historical details, and some translations of the notation of source papers, rest on these or on the author's memory, and the chapters say so where it matters. One referee pass per chapter finds many errors but not all of them. Results from 2024 to 2026 are cited from preprints where no refereed version could be found, and are marked as such where a claim rests on them. Corrections are welcome and will be recorded in the repository of \emph{The Missing Textbooks}.

\medskip
{\small\noindent Typeset with pdf\LaTeX{} in \texttt{missingbook}, the series class built on \textsc{koma}-Script, with the volume style \texttt{hydrogen}; Palatino (newpx) for text and mathematics, biblatex and biber for references, Ti\emph{k}Z and pgfplots for figures, whose data were computed with Python, NumPy, SciPy and PyTorch. The lines on the title page are the Balmer series of hydrogen, drawn at the wavelengths given by Rydberg's formula; the wave at the foot of each opening page rises from chapter to chapter.\par}

\end{document}
